% Auto-generated from data/segments.json + layout.json (+ transforms) — do not edit by hand.
% volume: 25Khu08
% mode: reading
% segments: 2218
% transform_rules: 0
% toc_marks: 82

% Volume layout defaults (layout.json ``layout``).
\csromanlayoutapply{1}{1.5}{21.6pt}{5pt}{6.3pt}{65pt}{2.5em}%

% Reading physical-page layout (layout.json ``page_layout_reading_mode``).
\csromanreadingpagelayoutvolume{1}{1.5}{21.6pt}{5pt}{6.3pt}{65pt}{2.5em}%
\csromanreadingpagelayoutenable

\csromanheader{2}{\textbf{ขุทฺทกนิกาย}}
\csromanmatikaanchor{mtk.0}
\csromantocmarknopage{chapter}{จูฬนิทฺเทสปาฬิ}{mtk.0}
\bookmark[dest={mtk.0},level=0]{จูฬนิทฺเทสปาฬิ}
\gambhira{\textbf{จูฬนิทฺเทสปาฬิ}}
\namakkaram{นโม ตสฺส ภควโต อรหโต สมฺมาสมฺพุทฺธสฺส.}
\csromanmatikaanchor{mtk.1}
\csromantocmarknopage{section}{ปารายนวคฺค}{mtk.1}
\bookmark[dest={mtk.1},level=1]{ปารายนวคฺค}
\csromanheader{1}{\textbf{ปารายนวคฺค}}
\csromanmatikaanchor{mtk.2}
\csromantocmarkref{subsection}{1. วตฺถุคาถา}{mtk.2}
\bookmark[dest={mtk.2},level=2]{1. วตฺถุคาถา}
\proseitem{1}{\csromanfolio{1}โกสลานํ ปุรา รมฺมา, อคมา ทกฺขิณาปถํ.}

\prose{อากิญฺจญฺญํ ปตฺถยาโน, พฺราหฺมโณ มนฺตปารคู.}

\proseitem{2}{โส อสฺสกสฺส วิสเย, มฬกสฺส\footnote{อฬกสฺส (ขุ 1. 429 ปิฏฺเฐ), มุฬกสฺส (สฺยา), มูฬฺหกสฺส (ก)} สมาสเน\footnote{สมาสนฺเน (ก)}.}

\prose{วสิ โคธาวรีกูเล, อุญฺเฉน จ ผเลน จ.}

\proseitem{3}{ตสฺเสว\footnote{ตํเยว (ก) อฏฺฐกถา โอโลเกตพฺพา.} อุปนิสฺสาย, คาโม จ วิปุโล อหุ.}

\prose{ตโต ชาเตน อาเยน, มหายญฺญมกปฺปยิ.}

\proseitem{4}{มหายญฺญํ ยชิตฺวาน, ปุน ปาวิสิ อสฺสมํ.}

\prose{ตสฺมิํ ปฏิปวิฏฺฐมฺหิ, อญฺโญ อาคญฺฉิ พฺราหฺมโณ.}

\proseitem{5}{อุคฺฆฏฺฏปาโท ตสิโต\footnote{ตสฺสิโต (ก)}, ปงฺกทนฺโต รชสฺสิโร.}

\csromangathameasure{โส จ นํ อุปสงฺกมฺม, สตานิ ปญฺจ ยาจติ.}
\csromangathagroup{\csromangathastanza{โส จ นํ อุปสงฺกมฺม, สตานิ ปญฺจ ยาจติ.}}

\proseitem{6}{ตเมนํ พาวรี ทิสฺวา, อาสเนน นิมนฺตยิ.}

\csromangathameasure{สุขญฺจ กุสลํ ปุจฺฉิ, อิทํ วจนมพฺรวิ.}
\csromangathagroup{\csromangathastanza{สุขญฺจ กุสลํ ปุจฺฉิ, อิทํ วจนมพฺรวิ\footnote{วจนมพฺรุวิ (สี)}.}}

\proseitem{7}{ยํ โข มม เทยฺยธมฺมํ, สพฺพํ วิสชฺชิตํ มยา.}

\csromangathameasure{อนุชานาหิ เม พฺรหฺเม, นตฺถิ ปญฺจสตานิ เม.}
\csromangathagroup{\csromangathastanza{อนุชานาหิ เม พฺรหฺเม, นตฺถิ ปญฺจสตานิ เม.}}

\csromanhangingbegin
\hangingheaditem{8}{\csromanfolio{2}สเจ เม ยาจมานสฺส, ภวํ นานุปทสฺสติ\footnote{ปเทสฺสติ (ก)}.}
\hangingbody{สตฺตเม ทิวเส ตุยฺหํ, มุทฺธา ผลตุ สตฺตธา.}
\csromanhangingend

\csromanhangingbegin
\hangingheaditem{9}{อภิสงฺขริตฺวา กุหโก, เภรวํ โส อกิตฺตยิ.}
\hangingbody{ตสฺส ตํ วจนํ สุตฺวา, พาวรี ทุกฺขิโต อหุ.}
\csromanhangingend

\csromanhangingbegin
\hangingheaditem{10}{อุสฺสุสฺสติ อนาหาโร, โสกสลฺลสมปฺปิโต.}
\hangingbody{อโถปิ เอวํ จิตฺตสฺส, ฌาเน น รมตี มโน.}
\csromanhangingend

\csromanhangingbegin
\hangingheaditem{11}{อุตฺรสฺตํ ทุกฺขิตํ ทิสฺวา, เทวตา อตฺถกามินี.}
\hangingbody{พาวริํ อุปสงฺกมฺม, อิทํ วจนมพฺรวิ.}
\csromanhangingend

\csromanhangingbegin
\hangingheaditem{12}{น โส มุทฺธํ ปชานาติ, กุหโก โส ธนตฺถิโก.}
\hangingbody{มุทฺธนิ มุทฺธปาเต2 วา, ญาณํ ตสฺส น วิชฺชติ.}
\csromanhangingend

\csromanhangingbegin
\hangingheaditem{13}{โภตี\footnote{โภติ (ก)} จรหิ ชานาติ, ตํ เม อกฺขาหิ ปุจฺฉิตา.}
\hangingbody{มุทฺธํ มุทฺธาธิปาตญฺจ4, ตํ สุโณม วโจ ตว.}
\csromanhangingend

\csromanhangingbegin
\hangingheaditem{14}{อหมฺเปตํ น ชานามิ, ญาณํ เมตฺถ น วิชฺชติ.}
\hangingbody{มุทฺธนิ มุทฺธาธิปาเต จ, ชินานํ เหตฺถ5 ทสฺสนํ.}
\csromanhangingend

\csromanhangingbegin
\hangingheaditem{15}{อถ โก จรหิ\footnote{โย จรติ (ก)} ชานาติ, อสฺมิํ ปถวิมณฺฑเล\footnote{ปุถวิมณฺฑเล (สี)}.}
\hangingbody{มุทฺธํ มุทฺธาธิปาตญฺจ, ตํ เม อกฺขาหิ เทวเต.}
\csromanhangingend

\csromanhangingbegin
\hangingheaditem{16}{ปุรา กปิลวตฺถุมฺหา, นิกฺขนฺโต โลกนายโก.}
\hangingbody{อปจฺโจ โอกฺกากราชสฺส, สกฺยปุตฺโต ปภงฺกโร.}
\csromanhangingend

\csromanhangingbegin
\hangingheaditem{17}{โส หิ พฺราหฺมณ สมฺพุทฺโธ, สพฺพธมฺมาน ปารคู.}
\hangingbody{สพฺพาภิญฺญาพลปฺปตฺโต8, สพฺพธมฺเมสุ จกฺขุมา.}
\hangingbody{สพฺพกมฺมกฺขยํ ปตฺโต, วิมุตฺโต อุปธิกฺขเย.}
\csromanhangingend

\csromanhangingbegin
\hangingheaditem{18}{พุทฺโธ โส ภควา โลเก, ธมฺมํ เทเสติ จกฺขุมา.}
\hangingbody{ตํ ตฺวํ คนฺตฺวาน ปุจฺฉสฺสุ, โส เต ตํ พฺยากริสฺสติ.}
\csromanhangingend

\csromanhangingbegin
\hangingheaditem{19}{สมฺพุทฺโธติ วโจ สุตฺวา, อุทคฺโค พาวรี อหุ.}
\hangingbody{โสก’สฺส ตนุโก อาสิ, ปีติญฺจ วิปุลํ ลภิ.}
\csromanhangingend

\csromanhangingbegin
\hangingheaditem{20}{\csromanfolio{3}โส พาวรี อตฺตมโน อุทคฺโค,}
\hangingbody{ตํ เทวตํ ปุจฺฉติ เวทชาโต.}
\hangingbody{กตมมฺหิ คาเม นิคมมฺหิ วา ปน,}
\hangingbody{กตมมฺหิ วา ชนปเท โลกนาโถ.}
\hangingbody{ยตฺถ คนฺตฺวาน ปสฺเสมุ, สมฺพุทฺธํ ทฺวิปทุตฺตมํ.}
\csromanhangingend

\csromanhangingbegin
\hangingheaditem{21}{สาวตฺถิยํ โกสลมนฺทิเร ชิโน,}
\hangingbody{ปหูตปญฺโญ วรภูริเมธโส.}
\hangingbody{โส สกฺยปุตฺโต วิธุโร อนาสโว,}
\hangingbody{มุทฺธาธิปาตสฺส วิทู นราสโภ.}
\csromanhangingend

\csromanhangingbegin
\hangingheaditem{22}{ตโต อามนฺตยี สิสฺเส, พฺราหฺมเณ มนฺตปารคู\footnote{ปารเค (สฺยา)}.}
\hangingbody{“เอถ มาณวา อกฺขิสฺสํ, สุณาถ วจนํ มม.}
\csromanhangingend

\csromanhangingbegin
\hangingheaditem{23}{ยสฺเสโส ทุลฺลโภ โลเก, ปาตุภาโว อภิณฺหโส.}
\hangingbody{สฺวาชฺช โลกมฺหิ อุปฺปนฺโน, สมฺพุทฺโธ อิติ วิสฺสุโต.}
\hangingbody{ขิปฺปํ คนฺตฺวาน สาวตฺถิํ, ปสฺสโวฺห ทฺวิปทุตฺตมํ”.}
\csromanhangingend

\csromanhangingbegin
\hangingheaditem{24}{กถํ จรหิ ชาเนมุ, ทิสฺวา “พุทฺโธ”ติ พฺราหฺมณ.}
\hangingbody{อชานตํ โน ปพฺรูหิ, ยถา ชาเนมุ ตํ มยํ.}
\csromanhangingend

\csromanhangingbegin
\hangingheaditem{25}{อาคตานิ หิ มนฺเตสุ, มหาปุริสลกฺขณา.}
\hangingbody{ทฺวตฺติํสานิ จ พฺยากฺขาตา, สมตฺตา อนุปุพฺพโส.}
\csromanhangingend

\csromanhangingbegin
\hangingheaditem{26}{ยสฺเสเต โหนฺติ คตฺเตสุ, มหาปุริสลกฺขณา.}
\hangingbody{เทฺวเยว ตสฺส คติโย, ตติยา หิ น วิชฺชติ.}
\csromanhangingend

\csromanhangingbegin
\hangingheaditem{27}{สเจ อคารํ อาวสติ, วิเชยฺย ปถวิํ อิมํ.}
\hangingbody{อทณฺเฑน อสตฺเถน, ธมฺเมน อนุสาสติ.}
\csromanhangingend

\csromanhangingbegin
\hangingheaditem{28}{สเจ จ โส ปพฺพชติ, อคารา อนคาริยํ.}
\hangingbody{วิวฏฺฏจฺฉโท2 สมฺพุทฺโธ, อรหา ภวติ อนุตฺตโร.}
\csromanhangingend

\csromanhangingbegin
\hangingheaditem{29}{ชาติํ โคตฺตญฺจ ลกฺขณํ, มนฺเต สิสฺเส ปุนาปเร.}
\hangingbody{มุทฺธํ มุทฺธาธิปาตญฺจ, มนสาเยว ปุจฺฉถ.}
\csromanhangingend

\csromanhangingbegin
\hangingheaditem{30}{อนาวรณทสฺสาวี, ยทิ พุทฺโธ ภวิสฺสติ.}
\hangingbody{มนสา ปุจฺฉิเต ปเญฺห, วาจาย วิสชฺชิสฺสติ3.}
\csromanhangingend

\csromanhangingbegin
\hangingheaditem{31}{\csromanfolio{4}พาวริสฺส วโจ สุตฺวา, สิสฺสา โสฬส พฺราหฺมณา.}
\hangingbody{อชิโต ติสฺสเมตฺเตยฺโย, ปุณฺณโก อถ เมตฺตคู.}
\csromanhangingend

\csromanhangingbegin
\hangingheaditem{32}{โธตโก อุปสีโว จ, นนฺโท จ อถ เหมโก.}
\hangingbody{โตเทยฺย-กปฺปา ทุภโย, ชตุกณฺณี จ ปณฺฑิโต.}
\csromanhangingend

\csromanhangingbegin
\hangingheaditem{33}{ภทฺราวุโธ อุทโย จ, โปสาโล จาปิ พฺราหฺมโณ.}
\hangingbody{โมฆราชา จ เมธาวี, ปิงฺคิโย จ มหา-อิสิ.}
\csromanhangingend

\csromanhangingbegin
\hangingheaditem{34}{ปจฺเจกคณิโน สพฺเพ, สพฺพโลกสฺส วิสฺสุตา.}
\hangingbody{ฌายี ฌานรตา ธีรา, ปุพฺพวาสนวาสิตา.}
\csromanhangingend

\csromanhangingbegin
\hangingheaditem{35}{พาวริํ อภิวาเทตฺวา, กตฺวา จ นํ ปทกฺขิณํ.}
\hangingbody{ชฏาชินธรา สพฺเพ, ปกฺกามุํ อุตฺตรามุขา.}
\csromanhangingend

\csromanhangingbegin
\hangingheaditem{36}{มฬกสฺส ปติฏฺฐานํ, ปุรมาหิสฺสติํ\footnote{ปุรมาหิยติ (ก)} ตทา\footnote{สทา (ก)}.}
\hangingbody{อุชฺเชนิญฺจาปิ โคนทฺธํ, เวทิสํ วนสวฺหยํ.}
\csromanhangingend

\csromanhangingbegin
\hangingheaditem{37}{โกสมฺพิญฺจาปิ สาเกตํ, สาวตฺถิญฺจ ปุรุตฺตมํ.}
\hangingbody{เสตพฺยํ กปิลวตฺถุํ, กุสินารญฺจ มนฺทิรํ.}
\csromanhangingend

\csromanhangingbegin
\hangingheaditem{38}{ปาวญฺจ โภคนครํ, เวสาลิํ มาคธํ ปุรํ.}
\hangingbody{ปาสาณกํ เจติยญฺจ, รมณียํ มโนรมํ.}
\csromanhangingend

\csromanhangingbegin
\hangingheaditem{39}{ตสิโตวุทกํ สีตํ, มหาลาภํว วาณิโช.}
\hangingbody{ฉายํ ฆมฺมาภิตตฺโตว, ตุริตา ปพฺพตมารุหุํ.}
\csromanhangingend

\csromanhangingbegin
\hangingheaditem{40}{ภควา ตมฺหิ สมเย, ภิกฺขุสํฆปุรกฺขโต.}
\hangingbody{ภิกฺขูนํ ธมฺมํ เทเสติ, สีโหว นทตี วเน.}
\csromanhangingend

\csromanhangingbegin
\hangingheaditem{41}{อชิโต อทฺทส พุทฺธํ, ปีตรํสิํว\footnote{ชิตรํสิํ สีตรํสิํ (ก), วีตรํสิํ (สี, สฺยา)} ภาณุมํ.}
\hangingbody{จนฺทํ ยถา ปนฺนรเส, ปริปูรํ4 อุปาคตํ.}
\csromanhangingend

\csromanhangingbegin
\hangingheaditem{42}{อถสฺส คตฺเต ทิสฺวาน, ปริปูรญฺจ พฺยญฺชนํ.}
\hangingbody{เอกมนฺตํ ฐิโต หฏฺโฐ, มโนปเญฺห อปุจฺฉถ.}
\csromanhangingend

\csromanhangingbegin
\hangingheaditem{43}{อาทิสฺส ชมฺมนํ พฺรูหิ, โคตฺตํ พฺรูหิ สลกฺขณํ.}
\hangingbody{มนฺเตสุ ปารมิํ พฺรูหิ, กติ วาเจติ พฺราหฺมโณ.}
\csromanhangingend

\csromanhangingbegin
\hangingheaditem{44}{\csromanfolio{5}วีสํ วสฺสสตํ อายุ, โส จ โคตฺเตน พาวรี.}
\hangingbody{ตีณิสฺส ลกฺขณา คตฺเต, ติณฺณํ เวทาน ปารคู.}
\csromanhangingend

\csromanhangingbegin
\hangingheaditem{45}{ลกฺขเณ อิติหาเส จ, สนิฆณฺฑุสเกฏุเภ.}
\hangingbody{ปญฺจสตานิ วาเจติ, สธมฺเม ปารมิํ คโต.}
\csromanhangingend

\csromanhangingbegin
\hangingheaditem{46}{ลกฺขณานํ ปวิจยํ, พาวริสฺส นรุตฺตม.}
\hangingbody{ตณฺหจฺฉิท1 ปกาเสหิ, มา โน กงฺขายิตํ อหุ.}
\csromanhangingend

\csromanhangingbegin
\hangingheaditem{47}{มุขํ ชิวฺหาย ฉาเทติ, อุณฺณ’สฺส ภมุกนฺตเร.}
\hangingbody{โกโสหิตํ วตฺถคุยฺหํ, เอวํ ชานาหิ มาณว.}
\csromanhangingend

\csromanhangingbegin
\hangingheaditem{48}{ปุจฺฉญฺหิ กิญฺจิ อสุณนฺโต, สุตฺวา ปเญฺห วิยากเต.}
\hangingbody{วิจินฺเตติ ชโน สพฺโพ, เวทชาโต กตญฺชลี.’}
\csromanhangingend

\csromanhangingbegin
\hangingheaditem{49}{โก นุ เทโว วา พฺรหฺมา วา, อินฺโท วาปิ สุชมฺปติ.}
\hangingbody{มนสา ปุจฺฉิเต ปเญฺห, กเมตํ ปฏิภาสติ.}
\csromanhangingend

\csromanhangingbegin
\hangingheaditem{50}{มุทฺธํ มุทฺธาธิปาตญฺจ, พาวรี ปริปุจฺฉติ.}
\hangingbody{ตํ พฺยากโรหิ ภควา, กงฺขํ วินย โน อิเส.}
\csromanhangingend

\csromanhangingbegin
\hangingheaditem{51}{อวิชฺชา มุทฺธาติ ชานาหิ, วิชฺชา มุทฺธาธิปาตินี.}
\hangingbody{สทฺธาสติสมาธีหิ, ฉนฺทวีริเยน สํยุตา.}
\csromanhangingend

\csromanhangingbegin
\hangingheaditem{52}{ตโต เวเทน มหตา, สนฺถมฺเภตฺวาน มาณโว.}
\hangingbody{เอกํสํ อชินํ กตฺวา, ปาเทสุ สิรสา ปติ.}
\csromanhangingend

\csromanhangingbegin
\hangingheaditem{53}{พาวรี พฺราหฺมโณ โภโต, สห สิสฺเสหิ มาริส.}
\hangingbody{อุทคฺคจิตฺโต สุมโน, ปาเท วนฺทติ จกฺขุม.}
\csromanhangingend

\csromanhangingbegin
\hangingheaditem{54}{สุขิโต พาวรี โหตุ, สห สิสฺเสหิ พฺราหฺมโณ.}
\hangingbody{ตฺวญฺจาปิ สุขิโต โหหิ, จิรํ ชีวาหิ มาณว.}
\csromanhangingend

\csromanhangingbegin
\hangingheaditem{55}{พาวริสฺส จ ตุยฺหํ วา, สพฺเพสํ สพฺพสํสยํ.}
\hangingbody{กตาวกาสา ปุจฺฉโวฺห, ยํ กิญฺจิ มนสิจฺฉถ.}
\csromanhangingend

\csromanhangingbegin
\hangingheaditem{56}{สมฺพุทฺเธน กโตกาโส, นิสีทิตฺวาน ปญฺชลี.}
\hangingbody{อชิโต ปฐมํ ปญฺหํ, ตตฺถ ปุจฺฉิ ตถาคตํ.}
\csromanhangingend

\nitthitam{วตฺถุคาถา นิฏฺฐิตา.}
\csromanmatikaanchor{mtk.3}
\csromantocmarkref{subsection}{1. อชิตมาณวปุจฺฉา}{mtk.3}
\bookmark[dest={mtk.3},level=2]{1. อชิตมาณวปุจฺฉา}
\csromanheader{2}{1. อชิตมาณวปุจฺฉา}
\csromanhangingbegin
\hangingheaditem{57}{\csromanfolio{6}เกนสฺสุ นิวุโต โลโก, (อิจฺจายสฺมา อชิโต,)}
\hangingbody{เกนสฺสุ นปฺปกาสติ.}
\hangingbody{กิสฺสาภิเลปนํ พฺรูสิ, กิํสุ ตสฺส มหพฺภยํ.}
\csromanhangingend

\csromanhangingbegin
\hangingheaditem{58}{อวิชฺชาย นิวุโต โลโก, (อชิตาติ ภควา,)}
\hangingbody{เววิจฺฉา ปมาทา นปฺปกาสติ,}
\hangingbody{ชปฺปาภิเลปนํ พฺรูมิ, ทุกฺข’มสฺส มหพฺภยํ.}
\csromanhangingend

\csromanhangingbegin
\hangingheaditem{59}{สวนฺติ สพฺพธิ โสตา, (อิจฺจายสฺมา อชิโต,)}
\hangingbody{โสตานํ กิํ นิวารณํ,}
\hangingbody{โสตานํ สํวรํ พฺรูหิ, เกน โสตา ปิธิยฺยเร.’}
\csromanhangingend

\csromanhangingbegin
\hangingheaditem{60}{ยานิ โสตานิ โลกสฺมิํ, (อชิตาติ ภควา,)}
\hangingbody{สติ เตสํ นิวารณํ.}
\hangingbody{โสตานํ สํวรํ พฺรูมิ, ปญฺญาเยเต ปิธิยฺยเร.}
\csromanhangingend

\csromanhangingbegin
\hangingheaditem{61}{ปญฺญา เจว สติ จาปิ\footnote{สตี เจว (สี)}, (อิจฺจายสฺมา อชิโต,)}
\hangingbody{นามรูปญฺจ มาริส.}
\hangingbody{เอตํ เม ปุฏฺโฐ ปพฺรูหิ, กตฺเถตํ อุปรุชฺฌติ.}
\csromanhangingend

\csromanhangingbegin
\hangingheaditem{62}{ยเมตํ ปญฺหํ อปุจฺฉิ, อชิต ตํ วทามิ เต.}
\hangingbody{ยตฺถ นามญฺจ รูปญฺจ, อเสสํ อุปรุชฺฌติ.}
\hangingbody{วิญฺญาณสฺส นิโรเธน, เอตฺเถตํ อุปรุชฺฌติ.}
\csromanhangingend

\csromanhangingbegin
\hangingheaditem{63}{เย จ สงฺขาตธมฺมาเส, เย จ เสขา\footnote{เสกฺขา (ก)} ปุถู อิธ.}
\hangingbody{เตสํ เม นิปโก อิริยํ, ปุฏฺโฐ ปพฺรูหิ มาริส.}
\csromanhangingend

\csromanhangingbegin
\hangingheaditem{64}{กาเมสุ นาภิคิชฺเฌยฺย, มนสา’นาวิโล สิยา.}
\hangingbody{กุสโล สพฺพธมฺมานํ, สโต ภิกฺขุ ปริพฺพเชติ.}
\csromanhangingend

\vspace{\tipitakaparskip}
\csromancenter{อชิตมาณวปุจฺฉา ปฐมา.}
\csromanmatikaanchor{mtk.4}
\csromanmatikaanchor{mtk.5}
\csromantocmarkref{subsection}{2. ติสฺสเมตฺเตยฺยมาณวปุจฺฉา}{mtk.4}
\bookmark[dest={mtk.4},level=2]{2. ติสฺสเมตฺเตยฺยมาณวปุจฺฉา}
\csromantocmarkref{subsection}{3. ปุณฺณกมาณวปุจฺฉา}{mtk.5}
\bookmark[dest={mtk.5},level=2]{3. ปุณฺณกมาณวปุจฺฉา}
\csromanheader{2}{3. ปุณฺณกมาณวปุจฺฉา \textbf{2. ติสฺสเมตฺเตยฺยมาณวปุจฺฉา}}
\csromanhangingbegin
\hangingheaditem{65}{\csromanfolio{7}โกธ สนฺตุสิโต โลเก, (อิจฺจายสฺมา ติสฺสเมตฺเตยฺโย,)}
\hangingbody{กสฺส โน สนฺติ อิญฺชิตา.}
\hangingbody{โก อุภนฺตมภิญฺญาย, มชฺเฌ มนฺตา น ลิปฺปติ1.}
\hangingbody{กํ พฺรูสิ มหาปุริโสติ, โก อิธ สิพฺพินิมจฺจคาติ2.}
\csromanhangingend

\csromanhangingbegin
\hangingheaditem{66}{กาเมสุ พฺรหฺมจริยวา, (เมตฺเตยฺยาติ ภควา,)}
\hangingbody{วีตตโณฺห สทา สโต.}
\hangingbody{สงฺขาย นิพฺพุโต ภิกฺขุ, ตสฺส โน สนฺติ อิญฺชิตา.}
\csromanhangingend

\csromanhangingbegin
\hangingheaditem{67}{โส อุภนฺตมภิญฺญาย, มชฺเฌ มนฺตา น ลิปฺปติ.}
\hangingbody{ตํ พฺรูมิ มหาปุริโสติ, โส อิธ สิพฺพินิมจฺจคาติ.}
\csromanhangingend

\vspace{\tipitakaparskip}
\csromancenter{ติสฺสเมตฺเตยฺยมาณวปุจฺฉา ทุติยา.}
\csromanheader{2}{3. ปุณฺณกมาณวปุจฺฉา}
\proseitem{68}{อเนชํ มูลทสฺสาวิํ, (อิจฺจายสฺมา ปุณฺณโก,)}

\prose{อตฺถิ ปเญฺหน อาคมํ.}

\prose{กิํ นิสฺสิตา อิสโย มนุชา,}

\proseitem{69}{เย เกจิเม อิสโย มนุชา, (ปุณฺณกาติ ภควา,)}

\prose{อาสีสมานา ปุณฺณก อิตฺตตฺตํ.}

\prose{ชรํ สิตา ยญฺญมกปฺปยิํสุ.}

\proseitem{70}{\csromanfolio{8}เย เกจิเม อิสโย มนุชา, (อิจฺจายสฺมา ปุณฺณโก,)}

\prose{กจฺจิสุ เต ภควา ยญฺญปเถ อปฺปมตฺตา.}

\prose{อตารุํ ชาติญฺจ ชรญฺจ มาริส,}

\csromanhangingbegin
\hangingheaditem{71}{อาสีสนฺติ โถมยนฺติ อภิชปฺปนฺติ ชุหนฺติ, (ปุณฺณกาติ ภควา,)}
\hangingbody{กามาภิชปฺปนฺติ ปฏิจฺจ ลาภํ.}
\hangingbody{เต ยาชโยคา ภวราครตฺตา,}
\hangingbody{นาตริํสุ ชาติชรนฺติ พฺรูมิ.}
\csromanhangingend

\proseitem{72}{เต เจ นาตริํสุ ยาชโยคา, (อิจฺจายสฺมา ปุณฺณโก,)}

\prose{ยญฺเญหิ ชาติญฺจ ชรญฺจ มาริส.}

\csromangathameasure{อถ โก จรหิ เทวมนุสฺสโลเก, \\ อตาริ ชาติญฺจ ชรญฺจ มาริส.}
\csromangathagroup{\csromangathastanza{อถ โก จรหิ เทวมนุสฺสโลเก, \\ อตาริ ชาติญฺจ ชรญฺจ มาริส.}}

\csromanhangingbegin
\hangingheaditem{73}{สงฺขาย โลกสฺมิ ปโรปรานิ, (ปุณฺณกาติ ภควา,)}
\hangingbody{ยสฺสิญฺชิตํ นตฺถิ กุหิญฺจิ โลเก.}
\hangingbody{สนฺโต วิธูโม อนีโฆ นิราโส,}
\hangingbody{อตาริ โส ชาติชรนฺติ พฺรูมีติ.}
\csromanhangingend

\vspace{\tipitakaparskip}
\csromancenter{ปุณฺณกมาณวปุจฺฉา ตติยา.}
\csromanmatikaanchor{mtk.6}
\csromantocmarkref{subsection}{4. เมตฺตคูมาณวปุจฺฉา}{mtk.6}
\bookmark[dest={mtk.6},level=2]{4. เมตฺตคูมาณวปุจฺฉา}
\csromanheader{2}{4. เมตฺตคูมาณวปุจฺฉา}
\csromanhangingbegin
\hangingheaditem{74}{ปุจฺฉามิ ตํ ภควา พฺรูหิ เม’ตํ, (อิจฺจายสฺมา เมตฺตคู,)}
\hangingbody{มญฺญามิ ตํ เวทคุํ ภาวิตตฺตํ.}
\hangingbody{กุโต นุ ทุกฺขา สมุทาคตา อิเม,}
\hangingbody{เย เกจิ โลกสฺมิ’มเนกรูปา.}
\csromanhangingend

\csromanheader{2}{4. เมตฺตคูมาณวปุจฺฉา}
\csromanhangingbegin
\hangingheaditem{75}{\csromanfolio{9}ทุกฺขสฺส เว มํ ปภวํ อปุจฺฉสิ, (เมตฺตคูติ ภควา,)}
\hangingbody{ตํ เต ปวกฺขามิ ยถา ปชานํ.}
\hangingbody{อุปธินิทานา ปภวนฺติ ทุกฺขา,}
\hangingbody{เย เกจิ โลกสฺมิ’มเนกรูปา.}
\csromanhangingend

\csromanhangingbegin
\hangingheaditem{76}{โย เว อวิทฺวา อุปธิํ กโรติ,}
\hangingbody{ปุนปฺปุนํ ทุกฺข’มุเปติ มนฺโท.}
\hangingbody{ตสฺมา ปชานํ อุปธิํ น กยิรา,}
\hangingbody{ทุกฺขสฺส ชาติปฺปภวานุปสฺสี.}
\csromanhangingend

\proseitem{77}{ยํ ตํ อปุจฺฉิมฺห อกิตฺตยี โน,}

\prose{อญฺญํ ตํ ปุจฺฉาม ตทิงฺฆ พฺรูหิ.}

\csromangathameasure{กถํ นุ ธีรา วิตรนฺติ โอฆํ, \\ ชาติํ ชรํ โสกปริทฺทวญฺจ. \\ ตํ เม มุนิ สาธุ วิยากโรหิ,}
\csromangathagroup{\csromangathastanza{กถํ นุ ธีรา วิตรนฺติ โอฆํ, \\ ชาติํ ชรํ โสกปริทฺทวญฺจ. \\ ตํ เม มุนิ สาธุ วิยากโรหิ,}}

\csromanhangingbegin
\hangingheaditem{78}{กิตฺตยิสฺสามิ เต ธมฺมํ, (เมตฺตคูติ ภควา,)}
\hangingbody{ทิฏฺเฐ ธมฺเม อนีติหํ.}
\hangingbody{ยํ วิทิตฺวา สโต จรํ, ตเร โลเก วิสตฺติกํ.}
\csromanhangingend

\csromanhangingbegin
\hangingheaditem{79}{ตญฺจาหํ อภินนฺทามิ, มเหสิ ธมฺมมุตฺตมํ.}
\hangingbody{ยํ วิทิตฺวา สโต จรํ, ตเร โลเก วิสตฺติกํ.}
\csromanhangingend

\csromanhangingbegin
\hangingheaditem{80}{ยํ กิญฺจิ สมฺปชานาสิ, (เมตฺตคูติ ภควา,)}
\hangingbody{อุทฺธํ อโธ ติริยญฺจาปิ มชฺเฌ.}
\hangingbody{เอเตสุ นนฺทิญฺจ นิเวสนญฺจ,}
\hangingbody{ปนุชฺช วิญฺญาณํ ภเว น ติฏฺเฐ.}
\csromanhangingend

\csromanhangingbegin
\hangingheaditem{81}{เอวํวิหารี สโต อปฺปมตฺโต,}
\hangingbody{ภิกฺขุ จรํ หิตฺวา มมายิตานิ.}
\hangingbody{ชาติํ ชรํ โสกปริทฺทวญฺจ,}
\hangingbody{อิเธว วิทฺวา ปชเหยฺย ทุกฺขํ.}
\csromanhangingend

\proseitem{82}{\csromanfolio{10}เอตา’ภินนฺทามิ วโจ มเหสิโน,}

\prose{สุกิตฺติตํ โคตม’นูปธีกํ.}

\prose{อทฺธา หิ ภควา ปหาสิ ทุกฺขํ,}

\csromanheader{2}{(9)}
\proseitem{83}{เต จาปิ นูนปฺปชเหยฺยุ ทุกฺขํ,}

\prose{เย ตฺวํ มุนิ อฏฺฐิตํ โอวเทยฺย.}

\prose{ตํ ตํ นมสฺสามิ สเมจฺจ นาค,}

\csromanheader{2}{อปฺเปว มํ ภควา อฏฺฐิตํ โอวเทยฺย. (10)}
\proseitem{84}{ยํ พฺราหฺมณํ เวทคุ’มาภิชญฺญา,}

\prose{อกิญฺจนํ กามภเว อสตฺตํ.}

\prose{อทฺธา หิ โส โอฆมิมํ อตาริ,}

\csromanheader{2}{ติณฺโณ จ ปารํ อขิโล อกงฺโข. (11)}
\proseitem{85}{วิทฺวา จ โย เวทคู นโร อิธ,}

\prose{ภวาภเว สงฺคมิมํ วิสชฺช.}

\prose{โส วีตตโณฺห อนีโฆ นิราโส,}

\csromanheader{2}{อตาริ โส ชาติชรนฺติ พฺรูมีติ. (12)}
\vspace{\tipitakaparskip}
\csromancenter{เมตฺตคูมาณวปุจฺฉา จตุตฺถี.}
\csromanmatikaanchor{mtk.7}
\csromantocmarkref{subsection}{5. โธตกมาณวปุจฺฉา}{mtk.7}
\bookmark[dest={mtk.7},level=2]{5. โธตกมาณวปุจฺฉา}
\csromanheader{2}{5. โธตกมาณวปุจฺฉา}
\csromanhangingbegin
\hangingheaditem{86}{ปุจฺฉามิ ตํ ภควา พฺรูหิ เม’ตํ, (อิจฺจายสฺมา โธตโก,)}
\hangingbody{วาจา’ภิกงฺขามิ มเหสิ ตุยฺหํ.}
\hangingbody{ตว สุตฺวาน นิคฺโฆสํ, สิกฺเข นิพฺพานมตฺตโน.}
\csromanhangingend

\csromanhangingbegin
\hangingheaditem{87}{เตนหา’ตปฺปํ กโรหิ, (โธตกาติ ภควา,)}
\hangingbody{อิเธว นิปโก สโต.}
\hangingbody{อิโต สุตฺวาน นิคฺโฆสํ, สิกฺเข นิพฺพานมตฺตโน.}
\csromanhangingend

\csromanhangingbegin
\hangingheaditem{88}{ปสฺสามหํ เทวมนุสฺสโลเก,}
\hangingbody{อกิญฺจนํ พฺราหฺมณ’มิริยมานํ.}
\hangingbody{ตํ ตํ นมสฺสามิ สมนฺตจกฺขุ,}
\hangingbody{ปมุญฺจ มํ สกฺก กถํกถาหิ.}
\csromanhangingend

\csromanmatikaanchor{mtk.8}
\csromantocmarkref{subsection}{6. อุปสีวมาณวปุจฺฉา}{mtk.8}
\bookmark[dest={mtk.8},level=2]{6. อุปสีวมาณวปุจฺฉา}
\csromanheader{2}{6. อุปสีวมาณวปุจฺฉา}
\csromanhangingbegin
\hangingheaditem{89}{\csromanfolio{11}นาหํ สหิสฺสามิ ปโมจนาย,}
\hangingbody{กถํกถิํ โธตก กญฺจิ โลเก.}
\hangingbody{ธมฺมญฺจ เสฏฺฐํ อภิชานมาโน1,}
\hangingbody{เอวํ ตุวํ โอฆมิมํ ตเรสิ.}
\csromanhangingend

\csromanhangingbegin
\hangingheaditem{90}{อนุสาส พฺรหฺเม กรุณายมาโน,}
\hangingbody{วิเวกธมฺมํ ยมหํ วิชญฺญํ.}
\hangingbody{ยถาหํ อากาโสว อพฺยาปชฺชมาโน,}
\hangingbody{อิเธว สนฺโต อสิโต จเรยฺยํ.}
\csromanhangingend

\csromanhangingbegin
\hangingheaditem{91}{กิตฺตยิสฺสามิ เต สนฺติํ, (โธตกาติ ภควา,)}
\hangingbody{ทิฏฺเฐ ธมฺเม อนีติหํ.}
\hangingbody{ยํ วิทิตฺวา สโต จรํ,}
\hangingbody{ตเร โลเก วิสตฺติกํ.}
\csromanhangingend

\csromanhangingbegin
\hangingheaditem{92}{ตญฺจาหํ อภินนฺทามิ, มเหสิ สนฺติ’มุตฺตมํ.}
\hangingbody{ยํ วิทิตฺวา สโต จรํ, ตเร โลเก วิสตฺติกํ.}
\csromanhangingend

\csromanhangingbegin
\hangingheaditem{93}{ยํ กิญฺจิ สมฺปชานาสิ, (โธตกาติ ภควา,)}
\hangingbody{อุทฺธํ อโธ ติริยญฺจาปิ มชฺเฌ,}
\hangingbody{เอตํ วิทิตฺวา สงฺโคติ โลเก.}
\hangingbody{ภวาภวาย มากาสิ ตณฺหนฺติ.}
\csromanhangingend

\vspace{\tipitakaparskip}
\csromancenter{โธตกมาณวปุจฺฉา ปญฺจมี.}
\csromanheader{2}{6. อุปสีวมาณวปุจฺฉา}
\csromanhangingbegin
\hangingheaditem{94}{เอโก อหํ สกฺก มหนฺตโมฆํ, (อิจฺจายสฺมา อุปสีโว,)}
\hangingbody{อนิสฺสิโต โน วิสหามิ ตาริตุํ.}
\hangingbody{อารมฺมณํ พฺรูหิ สมนฺตจกฺขุ,}
\hangingbody{ยํ นิสฺสิโต โอฆมิมํ ตเรยฺยํ.}
\csromanhangingend

\csromanhangingbegin
\hangingheaditem{95}{\csromanfolio{12}อากิญฺจญฺญํ เปกฺขมาโน สติมา, (อุปสีวาติ ภควา,)}
\hangingbody{นตฺถีติ นิสฺสาย ตรสฺสุ โอฆํ.}
\hangingbody{กาเม ปหาย วิรโต กถาหิ,}
\hangingbody{ตณฺหกฺขยํ นตฺตมหา’ภิปสฺส.}
\csromanhangingend

\csromanhangingbegin
\hangingheaditem{96}{สพฺเพสุ กาเมสุ โย วีตราโค, (อิจฺจายสฺมา อุปสีโว,)}
\hangingbody{อากิญฺจญฺญํ นิสฺสิโต หิตฺวา มญฺญํ.}
\hangingbody{สญฺญาวิโมกฺเข ปรเม วิมุตฺโต1,}
\hangingbody{ติฏฺเฐ นุ โส ตตฺถ อนานุยายี2.}
\csromanhangingend

\csromanhangingbegin
\hangingheaditem{97}{สพฺเพสุ กาเมสุ โย วีตราโค, (อุปสีวาติ ภควา,)}
\hangingbody{อากิญฺจญฺญํ นิสฺสิโต หิตฺวา มญฺญํ.}
\hangingbody{สญฺญาวิโมกฺเข ปรเม วิมุตฺโต,}
\hangingbody{ติฏฺเฐยฺย โส ตตฺถ อนานุยายี.}
\csromanhangingend

\csromanhangingbegin
\hangingheaditem{98}{ติฏฺเฐ เจ โส ตตฺถ อนานุยายี,}
\hangingbody{ปูคมฺปิ วสฺสานํ สมนฺตจกฺขุ.}
\hangingbody{ตตฺเถว โส สีติสิยา วิมุตฺโต,}
\hangingbody{จเวถ วิญฺญาณํ ตถาวิธสฺส.}
\csromanhangingend

\csromanhangingbegin
\hangingheaditem{99}{อจฺจิ ยถา วาตเวเคน ขิตฺตา, (อุปสีวาติ ภควา,)}
\hangingbody{อตฺถํ ปเลติ น อุเปติ สงฺขํ.}
\hangingbody{เอวํ มุนี นามกายา วิมุตฺโต,}
\hangingbody{อตฺถํ ปเลติ น อุเปติ สงฺขํ.}
\csromanhangingend

\proseitem{100}{อตฺถงฺคโต โส อุท วา โส นตฺถิ,}

\prose{อุทาหุ เว สสฺสติยา อโรโค.}

\prose{ตํ เม มุนี สาธุ วิยากโรหิ,}

\csromanhangingbegin
\hangingheaditem{101}{อตฺถงฺคตสฺส น ปมาณ’มตฺถิ, (อุปสีวาติ ภควา,)}
\hangingbody{เยน นํ วชฺชุํ ตํ ตสฺส นตฺถิ.}
\hangingbody{สพฺเพสุ ธมฺเมสุ สมูหเตสุ,}
\hangingbody{สมูหตา วาทปถาปิ สพฺเพติ.}
\csromanhangingend

\vspace{\tipitakaparskip}
\csromancenter{อุปสีวมาณวปุจฺฉา ฉฏฺฐี.}
\csromanmatikaanchor{mtk.9}
\csromantocmarkref{subsection}{7. นนฺทมาณวปุจฺฉา}{mtk.9}
\bookmark[dest={mtk.9},level=2]{7. นนฺทมาณวปุจฺฉา}
\csromanheader{2}{7. นนฺทมาณวปุจฺฉา}
\csromanheader{2}{7. นนฺทมาณวปุจฺฉา}
\csromanhangingbegin
\hangingheaditem{102}{\csromanfolio{13}“สนฺติ โลเก มุนโย”, (อิจฺจายสฺมา นนฺโท,)}
\hangingbody{ชนา วทนฺติ ตยิทํ กถํสุ.}
\hangingbody{ญาณูปปนฺนํ มุนิ โน วทนฺติ,}
\hangingbody{อุทาหุ เว ชีวิเตนูปปนฺนํ.}
\csromanhangingend

\csromanhangingbegin
\hangingheaditem{103}{น ทิฏฺฐิยา น สุติยา น ญาเณน,}
\hangingbody{มุนีธ นนฺท กุสลา วทนฺติ.}
\hangingbody{วิเสนิกตฺวา อนีฆา นิราสา.}
\hangingbody{จรนฺติ เย เต มุนโยติ พฺรูมิ.}
\csromanhangingend

\proseitem{104}{เย เกจิเม สมณพฺราหฺมณาเส, (อิจฺจายสฺมา นนฺโท,)}

\prose{ทิฏฺฐสฺสุเตนาปิ วทนฺติ สุทฺธิํ.}

\csromangathameasure{สีลพฺพเตนาปิ วทนฺติ สุทฺธิํ, \\ อเนกรูเปน วทนฺติ สุทฺธิํ. \\ กจฺจิสฺสุ เต ภควา ตตฺถ ยตา จรนฺตา, \\ อตารุ ชาติญฺจ ชรญฺจ มาริส.}
\csromangathagroup{\csromangathastanza{สีลพฺพเตนาปิ วทนฺติ สุทฺธิํ, \\ อเนกรูเปน วทนฺติ สุทฺธิํ. \\ กจฺจิสฺสุ เต ภควา ตตฺถ ยตา จรนฺตา, \\ อตารุ ชาติญฺจ ชรญฺจ มาริส.}}

\csromanhangingbegin
\hangingheaditem{105}{เย เกจิเม สมณพฺราหฺมณาเส, (นนฺทาติ ภควา,)}
\hangingbody{ทิฏฺฐสฺสุเตนาปิ วทนฺติ สุทฺธิํ.}
\hangingbody{สีลพฺพเตนาปิ วทนฺติ สุทฺธิํ,}
\hangingbody{อเนกรูเปน วทนฺติ สุทฺธิํ.}
\hangingbody{กิญฺจาปิ เต ตตฺถ ยตา จรนฺติ,}
\hangingbody{นาตริํสุ ชาติชรนฺติ พฺรูมิ.}
\csromanhangingend

\proseitem{106}{เย เกจิเม สมณพฺราหฺมณาเส, (อิจฺจายสฺมา นนฺโท,)}

\prose{ทิฏฺฐสฺสุเตนาปิ วทนฺติ สุทฺธิํ.}

\csromangathameasure{สีลพฺพเตนาปิ วทนฺติ สุทฺธิํ, \\ อเนกรูเปน วทนฺติ สุทฺธิํ. \\ เต เจ มุนิ พฺรูสิ อโนฆติณฺเณ, \\ อถ โก จรหิ เทวมนุสฺสโลเก. \\ อตาริ ชาติญฺจ ชรญฺจ มาริส,}
\csromangathagroup{\csromangathastanza{สีลพฺพเตนาปิ วทนฺติ สุทฺธิํ, \\ อเนกรูเปน วทนฺติ สุทฺธิํ. \\ เต เจ มุนิ พฺรูสิ อโนฆติณฺเณ, \\ อถ โก จรหิ เทวมนุสฺสโลเก.}\csromangathastanza{อตาริ ชาติญฺจ ชรญฺจ มาริส,}}

\csromanhangingbegin
\hangingheaditem{107}{\csromanfolio{14}นาหํ สพฺเพ สมณพฺราหฺมณาเส, (นนฺทาติ ภควา,)}
\hangingbody{ชาติชราย นิวุตาติ พฺรูมิ.}
\hangingbody{เย สีธ ทิฏฺฐํ ว สุตํ มุตํ วา,}
\hangingbody{สีลพฺพตํ วาปิ ปหาย สพฺพํ.}
\hangingbody{อเนกรูปมฺปิ ปหาย สพฺพํ,}
\hangingbody{ตณฺหํ ปริญฺญาย อนาสวาเส.}
\hangingbody{เต เว นรา โอฆติณฺณาติ พฺรูมิ.}
\csromanhangingend

\csromanhangingbegin
\hangingheaditem{108}{เอตา’ภินนฺทามิ วโจ มเหสิโน,}
\hangingbody{สุกิตฺติตํ โคตม’นูปธีกํ.}
\hangingbody{เย สีธ ทิฏฺฐํ ว สุตํ มุตํ วา,}
\hangingbody{สีลพฺพตํ วาปิ ปหาย สพฺพํ.}
\hangingbody{อเนกรูปมฺปิ ปหาย สพฺพํ,}
\hangingbody{ตณฺหํ ปริญฺญาย อนาสวาเส.}
\hangingbody{อหมฺปิ เต โอฆติณฺณาติ พฺรูมีติ.}
\csromanhangingend

\vspace{\tipitakaparskip}
\csromancenter{นนฺทมาณวปุจฺฉา สตฺตมา.}
\csromanmatikaanchor{mtk.10}
\csromantocmarkref{subsection}{8. เหมกมาณวปุจฺฉา}{mtk.10}
\bookmark[dest={mtk.10},level=2]{8. เหมกมาณวปุจฺฉา}
\csromanheader{2}{8. เหมกมาณวปุจฺฉา}
\csromanhangingbegin
\hangingheaditem{109}{เย เม ปุพฺเพ วิยากํสุ, (อิจฺจายสฺมา เหมโก,)}
\hangingbody{หุรํ โคตมสาสนา.}
\hangingbody{“อิจฺจาสิ อิติ ภวิสฺสติ”, สพฺพํ ตํ อิติหีติหํ.}
\hangingbody{สพฺพํ ตํ ตกฺกวฑฺฒนํ, นาหํ ตตฺถ อภิรมิํ. (1)}
\csromanhangingend

\csromanhangingbegin
\hangingheaditem{110}{ตฺวญฺจ เม ธมฺมมกฺขาหิ, ตณฺหานิคฺฆาตนํ มุนิ.}
\hangingbody{ยํ วิทิตฺวา สโต จรํ, ตเร โลเก วิสตฺติกํ.}
\csromanhangingend

\csromanhangingbegin
\hangingheaditem{111}{อิธ ทิฏฺฐสุตมุตวิญฺญาเตสุ, ปิยรูเปสุ เหมก.}
\hangingbody{ฉนฺทราควิโนทนํ, นิพฺพานปท’มจฺจุตํ.}
\csromanhangingend

\csromanhangingbegin
\hangingheaditem{112}{เอตทญฺญาย เย สตา, ทิฏฺฐธมฺมา’ภินิพฺพุตา.}
\hangingbody{อุปสนฺตา จ เต สทา, ติณฺณา โลเก วิสตฺติกนฺติ.}
\csromanhangingend

\vspace{\tipitakaparskip}
\csromancenter{เหมกมาณวปุจฺฉา อฏฺฐมา.}
\csromanmatikaanchor{mtk.11}
\csromanmatikaanchor{mtk.12}
\csromantocmarkref{subsection}{9. โตเทยฺยมาณวปุจฺฉา}{mtk.11}
\bookmark[dest={mtk.11},level=2]{9. โตเทยฺยมาณวปุจฺฉา}
\csromantocmarkref{subsection}{10. กปฺปมาณวปุจฺฉา}{mtk.12}
\bookmark[dest={mtk.12},level=2]{10. กปฺปมาณวปุจฺฉา}
\csromanheader{2}{10. กปฺปมาณวปุจฺฉา \textbf{9. โตเทยฺยมาณวปุจฺฉา}}
\csromanhangingbegin
\hangingheaditem{113}{\csromanfolio{15}ยสฺมิํ กามา น วสนฺติ, (อิจฺจายสฺมา โตเทยฺโย,)}
\hangingbody{ตณฺหา ยสฺส น วิชฺชติ.}
\hangingbody{กถํกถา จ โย ติณฺโณ,}
\hangingbody{วิโมกฺโข ตสฺส กีทิโส.}
\csromanhangingend

\csromanhangingbegin
\hangingheaditem{114}{ยสฺมิํ กามา น วสนฺติ, (โตเทยฺยาติ ภควา,)}
\hangingbody{ตณฺหา ยสฺส น วิชฺชติ.}
\hangingbody{กถํกถา จ โย ติณฺโณ,}
\hangingbody{วิโมกฺโข ตสฺส นา’ปโร.}
\csromanhangingend

\csromanhangingbegin
\hangingheaditem{115}{นิราสโส โส อุท อาสสาโน\footnote{อาสยาโน (ก)},}
\hangingbody{ปญฺญาณวา โส อุท ปญฺญกปฺปี.}
\hangingbody{มุนิํ อหํ สกฺก ยถา วิชญฺญํ,}
\hangingbody{ตํ เม วิยาจิกฺข สมนฺตจกฺขุ.}
\csromanhangingend

\csromanhangingbegin
\hangingheaditem{116}{นิราสโส โส น จ อาสสาโน,}
\hangingbody{ปญฺญาณวา โส น จ ปญฺญกปฺปี.}
\hangingbody{เอวมฺปิ โตเทยฺย มุนิํ วิชาน,}
\hangingbody{อกิญฺจนํ กามภเว อสตฺตนฺติ.}
\csromanhangingend

\vspace{\tipitakaparskip}
\csromancenter{โตเทยฺยมาณวปุจฺฉา นวมา.}
\csromanheader{2}{10. กปฺปมาณวปุจฺฉา}
\csromanhangingbegin
\hangingheaditem{117}{มชฺเฌ สรสฺมิํ ติฏฺฐตํ, (อิจฺจายสฺมา กปฺโป,)}
\hangingbody{โอเฆ ชาเต มหพฺภเย.}
\hangingbody{ชรามจฺจุปเรตานํ, ทีปํ ปพฺรูหิ มาริส.}
\hangingbody{ตฺวญฺจ เม ทีปมกฺขาหิ, ยถายิทํ นาปรํ สิยา.}
\csromanhangingend

\csromanhangingbegin
\hangingheaditem{118}{มชฺเฌ สรสฺมิํ ติฏฺฐตํ, (กปฺปาติ ภควา,)}
\hangingbody{โอเฆ ชาเต มหพฺภเย.}
\hangingbody{ชรามจฺจุปเรตานํ, ทีปํ ปพฺรูมิ กปฺป เต.}
\csromanhangingend

\csromanhangingbegin
\hangingheaditem{119}{\csromanfolio{16}อกิญฺจนํ อนาทานํ, เอตํ ทีปํ อนาปรํ.}
\hangingbody{นิพฺพานํ อิติ นํ พฺรูมิ, ชรามจฺจุปริกฺขยํ.}
\csromanhangingend

\csromanhangingbegin
\hangingheaditem{120}{เอตทญฺญาย เย สตา, ทิฏฺฐธมฺมาภินิพฺพุตา.}
\hangingbody{น เต มารวสานุคา, น เต มารสฺส ปฏฺฐคูติ1.}
\csromanhangingend

\vspace{\tipitakaparskip}
\csromancenter{กปฺปมาณวปุจฺฉา ทสมา.}
\csromanmatikaanchor{mtk.13}
\csromantocmarkref{subsection}{11. ชตุกณฺณิมาณวปุจฺฉา}{mtk.13}
\bookmark[dest={mtk.13},level=2]{11. ชตุกณฺณิมาณวปุจฺฉา}
\csromanheader{2}{11. ชตุกณฺณิมาณวปุจฺฉา}
\csromanhangingbegin
\hangingheaditem{121}{สุตฺวานหํ วีรมกามกามิํ, (อิจฺจายสฺมา ชตุกณฺณิ,)}
\hangingbody{โอฆาติคํ ปุฏฺฐุมกามมาคมํ.}
\hangingbody{สนฺติปทํ พฺรูหิ สหชเนตฺต,}
\hangingbody{ยถาตจฺฉํ ภควา พฺรูหิ เม’ตํ.}
\csromanhangingend

\csromanhangingbegin
\hangingheaditem{122}{ภควา หิ กาเม อภิภุยฺย อิริยติ,}
\hangingbody{อาทิจฺโจว ปถวิํ เตชี เตชสา.}
\hangingbody{ปริตฺตปญฺญสฺส เม ภูริปญฺญ,}
\hangingbody{อาจิกฺข ธมฺมํ ยมหํ วิชญฺญํ.}
\hangingbody{ชาติชราย อิธ วิปฺปหานํ.}
\csromanhangingend

\csromanhangingbegin
\hangingheaditem{123}{กาเมสุ วินย เคธํ, (ชตุกณฺณีติ ภควา,)}
\hangingbody{เนกฺขมฺมํ ทฏฺฐุ เขมโต.}
\hangingbody{อุคฺคหิตํ นิรตฺตํ วา, มา เต วิชฺชิตฺถ กิญฺจนํ.}
\csromanhangingend

\csromanhangingbegin
\hangingheaditem{124}{ยํ ปุพฺเพ ตํ วิโสเสหิ, ปจฺฉา เต มาหุ กิญฺจนํ.}
\hangingbody{มชฺเฌ เจ โน คเหสฺสสิ, อุปสนฺโต จริสฺสสิ.}
\csromanhangingend

\csromanhangingbegin
\hangingheaditem{125}{สพฺพโส นามรูปสฺมิํ, วีตเคธสฺส พฺราหฺมณ.}
\hangingbody{อาสวาสฺส น วิชฺชนฺติ, เยหิ มจฺจุวสํ วเชติ.}
\csromanhangingend

\vspace{\tipitakaparskip}
\csromancenter{ชตุกณฺณิมาณวปุจฺฉา เอกาทสมา.}
\csromanmatikaanchor{mtk.14}
\csromanmatikaanchor{mtk.15}
\csromantocmarkref{subsection}{12. ภทฺราวุธมาณวปุจฺฉา}{mtk.14}
\bookmark[dest={mtk.14},level=2]{12. ภทฺราวุธมาณวปุจฺฉา}
\csromantocmarkref{subsection}{13. อุทยมาณวปุจฺฉา}{mtk.15}
\bookmark[dest={mtk.15},level=2]{13. อุทยมาณวปุจฺฉา}
\csromanheader{2}{13. อุทยมาณวปุจฺฉา \textbf{12. ภทฺราวุธมาณวปุจฺฉา}}
\csromanhangingbegin
\hangingheaditem{126}{\csromanfolio{17}โอกญฺชหํ ตณฺหจฺฉิทํ อเนชํ, (อิจฺจายสฺมา ภทฺราวุโธ,)}
\hangingbody{นนฺทิญฺชหํ โอฆติณฺณํ วิมุตฺตํ.}
\hangingbody{กปฺปญฺชหํ อภิยาเจ สุเมธํ,}
\hangingbody{สุตฺวาน นาคสฺส อปนมิสฺสนฺติ อิโต.}
\csromanhangingend

\proseitem{127}{นานาชนา ชนปเทหิ สงฺคตา,}

\prose{ตว วีร วากฺยํ อภิกงฺขมานา.}

\prose{เตสํ ตุวํ สาธุ วิยากโรหิ,}

\csromanhangingbegin
\hangingheaditem{128}{อาทานตณฺหํ วินเยถ สพฺพํ, (ภทฺราวุธาติ ภควา,)}
\hangingbody{อุทฺธํ อโธ ติริยญฺจาปิ มชฺเฌ.}
\hangingbody{ยํ ยญฺหิ โลกสฺมิมุปาทิยนฺติ,}
\hangingbody{เตเนว มาโร อเนฺวติ ชนฺตุํ.}
\csromanhangingend

\csromanhangingbegin
\hangingheaditem{129}{ตสฺมา ปชานํ น อุปาทิเยถ,}
\hangingbody{ภิกฺขุ สโต กิญฺจนํ สพฺพโลเก.}
\hangingbody{อาทานสตฺเต อิติ เปกฺขมาโน,}
\hangingbody{ปชํ อิมํ มจฺจุเธยฺเย วิสตฺตนฺติ.}
\csromanhangingend

\vspace{\tipitakaparskip}
\csromancenter{ภทฺราวุธมาณวปุจฺฉา ทฺวาทสมา.}
\csromanheader{2}{13. อุทยมาณวปุจฺฉา}
\csromanhangingbegin
\hangingheaditem{130}{ฌายิํ วิรชมาสีนํ, (อิจฺจายสฺมา อุทโย)}
\hangingbody{กตกิจฺจํ อนาสวํ.}
\hangingbody{ปารคุํ สพฺพธมฺมานํ, อตฺถิ ปเญฺหน อาคมํ.}
\hangingbody{อญฺญาวิโมกฺขํ ปพฺรูหิ, อวิชฺชาย ปเภทนํ.}
\csromanhangingend

\csromanhangingbegin
\hangingheaditem{131}{ปหานํ กามจฺฉนฺทานํ, (อุทยาติ ภควา)}
\hangingbody{โทมนสฺสาน จูภยํ.}
\hangingbody{ถินสฺส จ ปนูทนํ, กุกฺกุจฺจานํ นิวารณํ.}
\csromanhangingend

\csromanhangingbegin
\hangingheaditem{132}{\csromanfolio{18}อุเปกฺขาสติสํสุทฺธํ, ธมฺมตกฺกปุเรชวํ.}
\hangingbody{อญฺญาวิโมกฺขํ ปพฺรูมิ, อวิชฺชาย ปเภทนํ.}
\csromanhangingend

\csromanhangingbegin
\hangingheaditem{133}{กิํสุ สํโยชโน โลโก, กิํสุ ตสฺส วิจารณํ.}
\hangingbody{กิสฺสสฺส วิปฺปหาเนน, นิพฺพานํ อิติ วุจฺจติ.}
\csromanhangingend

\csromanhangingbegin
\hangingheaditem{134}{นนฺทิสํโยชโน โลโก, วิตกฺกสฺส วิจารณํ.}
\hangingbody{ตณฺหาย วิปฺปหาเนน, นิพฺพานํ อิติ วุจฺจติ.}
\csromanhangingend

\csromanhangingbegin
\hangingheaditem{135}{กถํ สตสฺส จรโต, วิญฺญาณํ อุปรุชฺฌติ.}
\hangingbody{ภควนฺตํ ปุฏฺฐุมาคมฺม, ตํ สุโณม วโจ ตว. (6)}
\csromanhangingend

\csromanhangingbegin
\hangingheaditem{136}{อชฺฌตฺตญฺจ พหิทฺธา จ, เวทนํ นาภินนฺทโต.}
\hangingbody{เอวํ สตสฺส จรโต, วิญฺญาณํ อุปรุชฺฌตีติ.}
\csromanhangingend

\vspace{\tipitakaparskip}
\csromancenter{อุทยมาณวปุจฺฉา เตรสมา.}
\csromanmatikaanchor{mtk.16}
\csromantocmarkref{subsection}{14. โปสาลมาณวปุจฺฉา}{mtk.16}
\bookmark[dest={mtk.16},level=2]{14. โปสาลมาณวปุจฺฉา}
\csromanheader{2}{14. โปสาลมาณวปุจฺฉา}
\csromanhangingbegin
\hangingheaditem{137}{โย อตีตํ อาทิสติ, (อิจฺจายสฺมา โปสาโล,)}
\hangingbody{อเนโช ฉินฺนสํสโย.}
\hangingbody{ปารคุํ สพฺพธมฺมานํ, อตฺถิ ปเญฺหน อาคมํ.}
\csromanhangingend

\csromanhangingbegin
\hangingheaditem{138}{วิภูตรูปสญฺญิสฺส, สพฺพกายปฺปหายิโน.}
\hangingbody{อชฺฌตฺตญฺจ พหิทฺธา จ, นตฺถิ กิญฺจีติ ปสฺสโต.}
\hangingbody{ญาณํ สกฺกา’นุปุจฺฉามิ, กถํ เนยฺโย ตถาวิโธ.}
\csromanhangingend

\csromanhangingbegin
\hangingheaditem{139}{วิญฺญาณฏฺฐิติโย สพฺพา, (โปสาลาติ ภควา,)}
\hangingbody{อภิชานํ ตถาคโต.}
\hangingbody{ติฏฺฐนฺตเมนํ ชานาติ, วิมุตฺตํ ตปฺปรายณํ.}
\csromanhangingend

\csromanhangingbegin
\hangingheaditem{140}{อากิญฺจญฺญสมฺภวํ ญ ตฺวา, นนฺที สํโยชนํ อิติ.}
\hangingbody{เอวเมตํ อภิญฺญาย, ตโต ตตฺถ วิปสฺสติ.}
\hangingbody{เอตํ1 ญาณํ ตถํ ตสฺส, พฺราหฺมณสฺส วุสีมโตติ.}
\csromanhangingend

\vspace{\tipitakaparskip}
\csromancenter{โปสาลมาณวปุจฺฉา จุทฺทสมา.}
\csromanmatikaanchor{mtk.17}
\csromanmatikaanchor{mtk.18}
\csromantocmarkref{subsection}{15. โมฆราชมาณวปุจฺฉา}{mtk.17}
\bookmark[dest={mtk.17},level=2]{15. โมฆราชมาณวปุจฺฉา}
\csromantocmarkref{subsection}{16. ปิงฺคิยมาณวปุจฺฉา}{mtk.18}
\bookmark[dest={mtk.18},level=2]{16. ปิงฺคิยมาณวปุจฺฉา}
\csromanheader{2}{16. ปิงฺคิยมาณวปุจฺฉา \textbf{15. โมฆราชมาณวปุจฺฉา}}
\csromanhangingbegin
\hangingheaditem{141}{\csromanfolio{19}ทฺวาหํ สกฺกํ อปุจฺฉิสฺสํ, (อิจฺจายสฺมา โมฆราชา,)}
\hangingbody{น เม พฺยากาสิ จกฺขุมา.}
\hangingbody{ยาวตติยญฺจ เทวีสิ,}
\hangingbody{พฺยากโรตีติ เม สุตํ.}
\csromanhangingend

\csromanhangingbegin
\hangingheaditem{142}{อยํ โลโก ปโร โลโก,}
\hangingbody{พฺรหฺมโลโก สเทวโก.}
\hangingbody{ทิฏฺฐิํ เต นาภิชานาติ,}
\hangingbody{โคตมสฺส ยสสฺสิโน.}
\csromanhangingend

\csromanhangingbegin
\hangingheaditem{143}{เอวํ อภิกฺกนฺตทสฺสาวิํ, อตฺถิ ปเญฺหน อาคมํ.}
\hangingbody{กถํ โลกํ อเวกฺขนฺตํ, มจฺจุราชา น ปสฺสติ.}
\csromanhangingend

\csromanhangingbegin
\hangingheaditem{144}{สุญฺญโต โลกํ อเวกฺขสฺสุ, โมฆราช สทา สโต.}
\hangingbody{อตฺตานุทิฏฺฐิํ อูหจฺจ, เอวํ มจฺจุตโร สิยา.}
\hangingbody{เอวํ โลกํ อเวกฺขนฺตํ, มจฺจุราชา น ปสฺสตีติ.}
\csromanhangingend

\vspace{\tipitakaparskip}
\csromancenter{โมฆราชมาณวปุจฺฉา ปนฺนรสมา.}
\csromanheader{2}{16. ปิงฺคิยมาณวปุจฺฉา}
\csromanhangingbegin
\hangingheaditem{145}{ชิณฺโณหมสฺมิ อพโล วีตวณฺโณ, (อิจฺจายสฺมา ปิงฺคิโย,)}
\hangingbody{เนตฺตา น สุทฺธา สวนํ น ผาสุ.}
\hangingbody{มา’หํ นสฺสํ โมมุโห อนฺตราว,}
\hangingbody{อาจิกฺข ธมฺมํ ยมหํ วิชญฺญํ.}
\hangingbody{ชาติชราย อิธ วิปฺปหานํ.}
\csromanhangingend

\csromanhangingbegin
\hangingheaditem{146}{ทิสฺวาน รูเปสุ วิหญฺญมาเน, (ปิงฺคิยาติ ภควา,)}
\hangingbody{รุปฺปนฺติ รูเปสุ ชนา ปมตฺตา.}
\hangingbody{ตสฺมา ตุวํ ปิงฺคิย อปฺปมตฺโต,}
\hangingbody{ชหสฺสุ รูปํ อปุนพฺภวาย.}
\csromanhangingend

\csromanhangingbegin
\hangingheaditem{147}{\csromanfolio{20}ทิสา จตสฺโส วิทิสา จตสฺโส,}
\hangingbody{อุทฺธํ อโธ ทส ทิสา อิมาโย.}
\hangingbody{น ตุยฺหํ อทิฏฺฐํ อสุตํ อมุตํ1,}
\hangingbody{อโถ อวิญฺญาตํ กิญฺจน’มตฺถิ2 โลเก.}
\hangingbody{อาจิกฺข ธมฺมํ ยมหํ วิชญฺญํ,}
\hangingbody{ชาติชราย อิธ วิปฺปหานํ.}
\csromanhangingend

\csromanhangingbegin
\hangingheaditem{148}{ตณฺหาธิปนฺเน มนุเช เปกฺขมาโน, (ปิงฺคิยาติ ภควา,)}
\hangingbody{สนฺตาปชาเต ชรสา ปเรเต,}
\hangingbody{ตสฺมา ตุวํ ปิงฺคิย อปฺปมตฺโต.}
\hangingbody{ชหสฺสุ ตณฺหํ อปุนพฺภวายาติ.}
\csromanhangingend

\vspace{\tipitakaparskip}
\csromancenter{ปิงฺคิยมาณวปุจฺฉา โสฬสมา.}
\csromanmatikaanchor{mtk.19}
\csromantocmarkref{subsubsection}{ปารายนตฺถุติคาถา}{mtk.19}
\bookmark[dest={mtk.19},level=3]{ปารายนตฺถุติคาถา}
\csromanheader{3}{\textbf{ปารายนานุคีติคาถา} \textbf{ปารายนตฺถุติคาถา}}
\prose{\csromanfolio{21}อิทมโวจ ภควา มคเธสุ วิหรนฺโต ปาสาณเก เจติเย, ปริจารกโสฬสานํ\footnote{ปริจารกโสฬสนฺนํ (สฺยา, ก)} พฺราหฺมณานํ อชฺฌิฏฺโฐ ปุฏฺโฐ ปุฏฺโฐ ปญฺหํ\footnote{ปเญฺห (สี, อิ)} พฺยากาสิ.\csromanspacer{} เอกเมกสฺส เจปิ ปญฺหสฺส อตฺถมญฺญาย ธมฺมมญฺญาย ธมฺมานุธมฺมํ ปฏิปชฺเชยฺย, คจฺเฉยฺเยว ชรามรณสฺส ปารํ, ปารงฺคมนียา อิเม ธมฺมาติ, ตสฺมา อิมสฺส ธมฺมปริยายสฺส ปารายนนฺเตว\footnote{ปารายณํเตฺวว (สี-ฏฺฐ)} อธิวจนํ.}

\csromanhangingbegin
\hangingheaditem{149}{อชิโต ติสฺสเมตฺเตยฺโย, ปุณฺณโก อถ เมตฺตคู.}
\hangingbody{โธตโก อุปสีโว จ, นนฺโท จ อถ เหมโก.}
\csromanhangingend

\csromanhangingbegin
\hangingheaditem{150}{โตเทยฺยกปฺปา ทุภโย, ชตุกณฺณี จ ปณฺฑิโต.}
\hangingbody{ภทฺราวุโธ อุทโย จ, โปสาโล จาปิ พฺราหฺมโณ.}
\hangingbody{โมฆราชา จ เมธาวี, ปิงฺคิโย จ มหา-อิสิ.}
\csromanhangingend

\csromanhangingbegin
\hangingheaditem{151}{เอเต พุทฺธํ อุปาคจฺฉุํ, สมฺปนฺนจรณํ อิสิํ.}
\hangingbody{ปุจฺฉนฺตา นิปุเณ ปเญฺห, พุทฺธเสฏฺฐํ อุปาคมุํ.}
\csromanhangingend

\csromanhangingbegin
\hangingheaditem{152}{เตสํ พุทฺโธ ปพฺยากาสิ, ปเญฺห ปุฏฺโฐ ยถาตถํ.}
\hangingbody{ปญฺหานํ เวยฺยากรเณน, โตเสสิ พฺราหฺมเณ มุนิ.}
\csromanhangingend

\csromanhangingbegin
\hangingheaditem{153}{เต โตสิตา จกฺขุมตา, พุทฺเธนาทิจฺจพนฺธุนา.}
\hangingbody{พฺรหฺมจริยมจริํสุ, วรปญฺญสฺส สนฺติเก.}
\csromanhangingend

\csromanhangingbegin
\hangingheaditem{154}{เอกเมกสฺส ปญฺหสฺส, ยถา พุทฺเธน เทสิตํ.}
\hangingbody{ตถา โย ปฏิปชฺเชยฺย, คจฺเฉ ปารํ อปารโต.}
\csromanhangingend

\csromanhangingbegin
\hangingheaditem{155}{อปารา ปารํ คจฺเฉยฺย, ภาเวนฺโต มคฺคมุตฺตมํ.}
\hangingbody{มคฺโค โส ปารํ คมนาย, ตสฺมา ปารายนํ อิติ.}
\csromanhangingend

\csromanmatikaanchor{mtk.20}
\csromantocmarkref{subsubsection}{ปารายนานุคีติคาถา}{mtk.20}
\bookmark[dest={mtk.20},level=3]{ปารายนานุคีติคาถา}
\csromanheader{3}{\textbf{ปารายนานุคีติคาถา}}
\csromanhangingbegin
\hangingheaditem{156}{ปารายนมนุคายิสฺสํ, (อิจฺจายสฺมา ปิงฺคิโย,)}
\hangingbody{ยถา’ทฺทกฺขิ ตถา’กฺขาสิ, วิมโล ภูริเมธโส.}
\hangingbody{นิกฺกาโม นิพฺพโน4 นาโค, กิสฺส เหตุ มุสา ภเณ.(1)}
\csromanhangingend

\csromanhangingbegin
\hangingheaditem{157}{ปหีนมลโมหสฺส, มานมกฺขปฺปหายิโน.}
\hangingbody{หนฺทาหํ กิตฺตยิสฺสามิ, คิรํ วณฺณูปสญฺหิตํ.}
\csromanhangingend

\csromanhangingbegin
\hangingheaditem{158}{\csromanfolio{22}ตโมนุโท พุทฺโธ สมนฺตจกฺขุ,}
\hangingbody{โลกนฺตคู สพฺพภวาติวตฺโต.}
\hangingbody{อนาสโว สพฺพทุกฺขปฺปหีโน,}
\hangingbody{สจฺจวฺหโย พฺรหฺเม อุปาสิโต เม.}
\csromanhangingend

\csromanhangingbegin
\hangingheaditem{159}{ทิโช ยถา กุพฺพนกํ ปหาย,}
\hangingbody{พหุปฺผลํ กานนมาวเสยฺย.}
\hangingbody{เอวมฺปหํ อปฺปทสฺเส ปหาย,}
\hangingbody{มโหทธิํ หํโสริว อชฺฌปตฺโต.}
\csromanhangingend

\csromanhangingbegin
\hangingheaditem{160}{เยเม ปุพฺเพ วิยากํสุ,}
\hangingbody{หุรํ โคตมสาสนา. “อิจฺจาสิ อิติ ภวิสฺสติ”,}
\hangingbody{สพฺพํ ตํ อิติหีติหํ, สพฺพํ ตํ ตกฺกวฑฺฒนํ.}
\csromanhangingend

\csromanhangingbegin
\hangingheaditem{161}{เอโก ตมนุทาสิโน, ชุติมา โส ปภงฺกโร.}
\hangingbody{โคตโม ภูริปญฺญาโณ, โคตโม ภูริเมธโส.}
\csromanhangingend

\csromanhangingbegin
\hangingheaditem{162}{โย เม ธมฺมมเทเสสิ, สนฺทิฏฺฐิกมกาลิกํ.}
\hangingbody{ตณฺหกฺขยมนีติกํ, ยสฺส นตฺถิ อุปมา กฺวจิ.}
\csromanhangingend

\csromanhangingbegin
\hangingheaditem{163}{กิํ นุ ตมฺหา วิปฺปวสสิ, มุหุตฺตมปิ ปิงฺคิย.}
\hangingbody{โคตมา ภูริปญฺญาณา, โคตมา ภูริเมธสา.}
\csromanhangingend

\proseitem{164}{โย เต ธมฺมมเทเสสิ, สนฺทิฏฺฐิกมกาลิกํ.}

\prose{ตณฺหกฺขยมนีติกํ, ยสฺส นตฺถิ อุปมา กฺวจิ. (9)}

\proseitem{165}{นาหํ ตมฺหา วิปฺปวสามิ, มุหุตฺตมปิ พฺราหฺมณ.}

\prose{โคตมา ภูริปญฺญาณา, โคตมา ภูริเมธสา. (10)}

\proseitem{166}{โย เม ธมฺมมเทเสสิ, สนฺทิฏฺฐิกมกาลิกํ.}

\prose{ตณฺหกฺขยมนีติกํ, ยสฺส นตฺถิ อุปมา กฺวจิ. (11)}

\proseitem{167}{ปสฺสามิ นํ มนสา จกฺขุนาว,}

\prose{รตฺตินฺทิวํ พฺราหฺมณ อปฺปมตฺโต.}

\prose{นมสฺสมาโน วิวเสมิ รตฺติํ,}

\csromanheader{2}{เตเนว มญฺญามิ อวิปฺปวาสํ. (12)}
\csromanheader{2}{ปารายนานุคีติคาถา}
\proseitem{168}{\csromanfolio{23}สทฺธา จ ปีติ จ มโน สติ จ,}

\prose{นา’เปนฺติ’เม โคตมสาสนมฺหา.}

\prose{ยํ ยํ ทิสํ วชติ ภูริปญฺโญ,}

\csromanheader{2}{ส เตน เตเนว นโตหมสฺมิ. (13)}
\proseitem{169}{ชิณฺณสฺส เม ทุพฺพลถามกสฺส,}

\prose{เตเนว กาโย น ปเลติ ตตฺถ.}

\prose{สํกปฺปยนฺตาย\footnote{สํกปฺปยตฺตาย (สี)} วชามิ นิจฺจํ,}

\csromanheader{2}{มโน หิ เม พฺราหฺมณ เตน ยุตฺโต. (14)}
\proseitem{170}{ปงฺเก สยาโน ปริผนฺทมาโน,}

\prose{ทีปา ทีปํ อุปลฺลวิํ.}

\prose{อถทฺทสาสิํ สมฺพุทฺธํ, โอฆติณฺณมนาสวํ. (15)}

\proseitem{171}{ยถา อหู วกฺกลิ มุตฺตสทฺโธ,}

\prose{ภทฺราวุโธ อาฬวิ โคตโม จ.}

\prose{เอวเมว ตฺวมฺปิ ปมุญฺจสฺสุ สทฺธํ,}

\csromanheader{2}{คมิสฺสสิ ตฺวํ ปิงฺคิย มจฺจุเธยฺยสฺส ปารํ\footnote{มจฺจุเธยฺยปารํ (สี)}. (16)}
\proseitem{172}{เอส ภิยฺโย ปสีทามิ, สุตฺวาน มุนิโน วโจ.}

\prose{วิวฏฺฏจฺฉโท สมฺพุทฺโธ, อขิโล ปฏิภานวา. (17)}

\proseitem{173}{อธิเทเว อภิญฺญาย, สพฺพํ เวทิ ปโรปรํ.}

\prose{ปญฺหานนฺตกโร สตฺถา, กงฺขีนํ ปฏิชานตํ. (18)}

\proseitem{174}{อสํหีรํ อสํกุปฺปํ,}

\prose{ยสฺส นตฺถิ อุปมา กฺวจิ.}

\prose{อทฺธา คมิสฺสามิ น เมตฺถ กงฺขา,}

\csromanheader{2}{เอวํ มํ ธาเรหิ อธิมุตฺตจิตฺตนฺติ\footnote{อชิตมาณวปุจฺฉาย ปฏฺฐาย ยาวปารายนานุคีติคาถาปริโยสานา สฺยา...โปตฺถเก นตฺถิ.}. (19)}
\nitthitam{\textbf{ปารายนานุคีติคาถา นิฏฺฐิตา.}}
\chapterheadpageread{\textbf{ปารายนวคฺค}}
\csromanmatikaanchor{mtk.21}
\csromanmatikaanchor{mtk.22}
\csromantocmarknopage{section}{ปารายนวคฺค}{mtk.21}
\bookmark[dest={mtk.21},level=1]{ปารายนวคฺค}
\csromantocmarkref{subsection}{1. อชิตมาณวปุจฺฉานิทฺเทส}{mtk.22}
\bookmark[dest={mtk.22},level=2]{1. อชิตมาณวปุจฺฉานิทฺเทส}
\csromanheader{2}{1. อชิตมาณวปุจฺฉานิทฺเทส}
\proseitem{1}{\csromanfolio{24}\textbf{เกนสฺสุ นิวุโต โลโก,} (\textbf{อิจฺจา}ยสฺมา \textbf{อชิโต,})}

\prose{\textbf{เกนสฺสุ นปฺปกาสติ.}}

\csromangathameasure{\textbf{กิสฺสาภิเลปนํ พฺรูสิ}\textbf{,} \\ \textbf{กิํสุ ตสฺส มหพฺภยํ.}}
\csromangathagroup{\csromangathastanza{\textbf{กิสฺสาภิเลปนํ พฺรูสิ}\footnote{พฺรูหิ (สฺยา)}\textbf{,} \\ \textbf{กิํสุ ตสฺส มหพฺภยํ.}}}

\prose{\textbf{เกนสฺสุ นิวุโต โลโก}ติ \textbf{โลโก}ติ นิรย\textbf{โลโก} ติรจฺฉาน\textbf{โลโก} เปตฺติวิสย\textbf{โลโก} มนุสฺส\textbf{โลโก} เทว\textbf{โลโก} ขนฺธ\textbf{โลโก} ธาตุ\textbf{โลโก} อายตน\textbf{โลโก} อยํ \textbf{โลโก} ปโร \textbf{โลโก} พฺรหฺม\textbf{โลโก} เทว\textbf{โลโก}.\csromanspacer{} อยํ วุจฺจติ \textbf{โลโก}.\csromanspacer{} อยํ \textbf{โลโก} เกน อาวุโต นิวุโต โอวุโต\footnote{โอผุโต (สฺยา)} ปิหิโต ปฏิจฺฉนฺโน ปฏิกุชฺชิโตติ เกนสฺสุ นิวุโต \textbf{โลโก}.}

\prose{\textbf{อิจฺจายสฺมา อชิโต}ติ \textbf{อิจฺจา}ติ ปทสนฺธิ ปทสํสคฺโค ปทปาริปูรี อกฺขรสมวาโย พฺยญฺชนสิลิฏฺฐตา ปทานุปุพฺพตาเปตํ\footnote{ปทานุปุพฺพตาเมตํ (พหูสุ)} \textbf{อิจฺจา}ติ.\csromanspacer{} \textbf{อายสฺมา}ติ ปิยวจนํ ครุวจนํ สคารวสปฺปติสฺสาธิวจนเมตํ อายสฺมาติ.\csromanspacer{} \textbf{อชิโต}ติ ตสฺส พฺราหฺมณสฺส นามํ สงฺขา สมญฺญา ปญฺญตฺติ โวหาโร นามํ นามกมฺมํ นามเธยฺยํ นิรุตฺติ พฺยญฺชนํ อภิลาโปติ \textbf{อิจฺจา}ยสฺมา \textbf{อชิโต}.}

\prose{\textbf{เกนสฺสุ นปฺปกาสตี}ติ เกน \textbf{โลโก} นปฺปกาสติ น ภาสติ น ตปติ น วิโรจติ น ญายติ น ปญฺญายตีติ เกนสฺสุ นปฺปกาสติ.}

\prose{\textbf{กิสฺสาภิเลปนํ พฺรูสี}ติ กิํ โลกสฺส เลปนํ ลคฺคนํ พนฺธนํ อุปกฺกิเลโส.\csromanspacer{} เกน \textbf{โลโก} ลิตฺโต สํลิตฺโต อุปลิตฺโต กิลิฏฺโฐ สํกิลิฏฺโฐ มกฺขิโต สํสฏฺโฐ ลคฺโค ลคฺคิโต ปลิพุทฺโธ พฺรูสิ อาจิกฺขสิ เทเสสิ ปญฺญเปสิ\footnote{ปญฺญาเปสิ (ก)} ปฏฺฐเปสิ วิวรสิ วิภชสิ อุตฺตานีกโรสิ\footnote{อุตฺตานิํ กโรสิ (ก)} ปกาเสสีติ กิสฺสาภิเลปนํ พฺรูสิ.}

\csromanheader{2}{1. อชิตมาณวปุจฺฉานิทฺเทส}
\prose{\csromanfolio{25}\textbf{กิํสุ ตสฺส มหพฺภยนฺ}ติ กิํ โลกสฺส ภยํ มหพฺภยํ ปีฬนํ ฆฏฺฏนํ อุปทฺทโว อุปสคฺโคติ กิํสุ ตสฺส มหพฺภยํ.\csromanspacer{} เตนาห โส พฺราหฺมโณ\csromandash{}}

\prose{เกนสฺสุ นิวุโต โลโก, (อิจฺจายสฺมา อชิโต.)}

\prose{เกนสฺสุ นปฺปกาสติ.}

\csromangathameasure{กิสฺสาภิเลปนํ พฺรูสิ, \\ \textbf{กิํสุ ตสฺส มหพฺภยนฺ}ติ.}
\csromangathagroup{\csromangathastanza{กิสฺสาภิเลปนํ พฺรูสิ, \\ \textbf{กิํสุ ตสฺส มหพฺภยนฺ}ติ.}}

\proseitem{2}{\textbf{อวิชฺชาย นิวุโต โลโก}, (\textbf{อชิตา}ติ \textbf{ภควา},)}

\prose{\textbf{เววิจฺฉา ปมาทา นปฺปกาสติ.}}

\csromangathameasure{\textbf{ชปฺปาภิเลปนํ พฺรูมิ,} \\ \textbf{ทุกฺข’มสฺส มหพฺภยํ.}}
\csromangathagroup{\csromangathastanza{\textbf{ชปฺปาภิเลปนํ พฺรูมิ,} \\ \textbf{ทุกฺข’มสฺส มหพฺภยํ.}}}

\prose{\textbf{อวิชฺชาย นิวุโต โลโก}ติ \textbf{อวิชฺชา}ติ ทุกฺเข อญฺญาณํ ทุกฺขสมุทเย อญฺญาณํ ทุกฺขนิโรเธ อญฺญาณํ ทุกฺขนิโรธคามินิยา ปฏิปทาย อญฺญาณํ ปุพฺพนฺเต อญฺญาณํ อปรนฺเต อญฺญาณํ ปุพฺพนฺตาปรนฺเต อญฺญาณํ อิทปฺปจฺจยตาปฏิจฺจสมุปฺปนฺเนสุ ธมฺเมสุ อญฺญาณํ, ยํ เอวรูปํ อญฺญาณํ อทสฺสนํ อนภิสมโย อนนุโพโธ อสมฺโพโธ อปฺปฏิเวโธ อสํคาหนา อปริโยคาหนา อสมเปกฺขนา อปจฺจเวกฺขณา\footnote{อปจฺจเวกฺขนา (สฺยา)} อปจฺจเวกฺขณกมฺมํ ทุมฺเมชฺฌํ พาลฺยํ อสมฺปชญฺญํ โมโห ปโมโห สมฺโมโห \textbf{อวิชฺชา} อวิชฺโชโฆ \textbf{อวิชฺชา}โยโค \textbf{อวิชฺชา}นุสโย \textbf{อวิชฺชา}ปริยุฏฺฐานํ \textbf{อวิชฺชา}ลงฺคี โมโห อกุสลมูลํ.\csromanspacer{} อยํ วุจฺจติ \textbf{อวิชฺชา}.}

\prose{\textbf{โลโก}ติ นิรยโลโก ติรจฺฉานโลโก เปตฺติวิสยโลโก มนุสฺสโลโก เทวโลโก ขนฺธโลโก ธาตุโลโก อายตนโลโก อยํ โลโก ปโร โลโก พฺรหฺมโลโก เทวโลโก.\csromanspacer{} อยํ วุจฺจติ โลโก.\csromanspacer{} อยํ โลโก อิมาย \textbf{อวิชฺชา}ย อาวุโต นิวุโต โอวุโต ปิหิโต ปฏิจฺฉนฺโน ปฏิกุชฺชิโตติ \textbf{อวิชฺชา}ย นิวุโต โลโก.}

\prose{\textbf{อชิตา}ติ \textbf{ภควา} ตํ พฺราหฺมณํ นาเมน อาลปติ.\csromanspacer{} \textbf{ภควา}ติ คารวาธิวจนํ.\csromanspacer{} อปิ จ ภคฺคราโคติ \textbf{ภควา}, ภคฺคโทโสติ \textbf{ภควา}, \csromanfolio{26}ภคฺคโมโหติ ภควา, ภคฺคมาโนติ ภควา, ภคฺคทิฏฺฐีติ ภควา, ภคฺคกณฺฏโกติ ภควา, ภคฺคกิเลโสติ ภควา, ภชิ วิภชิ ปวิภชิ ธมฺมรตนนฺติ ภควา, ภวานํ อนฺตกโรติ ภควา, ภวิตกาโย, ภาวิตสีโล, ภาวิตจิตฺโต\footnote{ภาวิตกาโยติ ภควา, ภาวิตสีโลติ ภาวิตจิตฺโตติ (สฺยา)}, ภาวิตปญฺโญติ ภควา, ภชิ วา ภควา อรญฺญวนปตฺถานิ ปนฺตานิ เสนาสนานิ อปฺปสทฺทานิ อปฺปนิคฺโฆสานิ วิชนวาตานิ มนุสฺสราหสฺเสยฺยกานิ\footnote{มนุสฺสราหเสยฺยกานิ (สฺยา)} ปฏิสลฺลานสารุปฺปานีติ ภควา, ภาคี วา ภควา จีวรปิณฺฑปาตเสนาสนคิลานปจฺจยเภสชฺชปริกฺขารานนฺติ ภควา, ภาคี วา ภควา อตฺถรสสฺส ธมฺมรสสฺส วิมุตฺติรสสฺส อธิสีลสฺส อธิจิตฺตสฺส อธิปญฺญายาติ ภควา, ภาคี วา ภควา จตุนฺนํ ฌานานํ จตุนฺนํ อปฺปมญฺญานํ จตุนฺนํ อรูปสมาปตฺตีนนฺติ ภควา, ภาคี วา ภควา อฏฺฐนฺนํ วิโมกฺขานํ อฏฺฐนฺนํ อภิภายตนานํ นวนฺนํ อนุปุพฺพสมาปตฺตีนนฺติ ภควา, ภาคี วา ภควา ทสนฺนํ สญฺญาภาวนานํ กสิณสมาปตฺตีนํ อานาปานสฺสติสมาธิสฺส อสุภสมาปตฺติยาติ ภควา, ภาคี วา ภควา จตุนฺนํ สติปฏฺฐานานํ จตุนฺนํ สมฺมปฺปธานานํ จตุนฺนํ อิทฺธิปาทานํ ปญฺจนฺนํ อินฺทฺริยานํ ปญฺจนฺนํ พลานํ สตฺตนฺนํ โพชฺฌงฺคานํ อริยสฺส อฏฺฐงฺคิกสฺส มคฺคสฺสาติ ภควา, ภาคี วา ภควา ทสนฺนํ ตถาคตพลานํ จตุนฺนํ เวสารชฺชานํ จตุนฺนํ ปฏิสมฺภิทานํ ฉนฺนํ อภิญฺญานํ ฉนฺนํ พุทฺธธมฺมานนฺติ ภควา.\csromanspacer{} ภควาติ เนตํ นามํ มาตรา กตํ, น ปิตรา กตํ, น ภาตรา กตํ, น ภคินิยา กตํ, น มิตฺตามจฺเจหิ กตํ, น ญาติสาโลหิเตหิ กตํ, น สมณพฺราหฺมเณหิ กตํ, น เทวตาหิ กตํ, วิโมกฺขนฺติกเมตํ พุทฺธานํ ภควนฺตานํ โพธิยา มูเล สห สพฺพญฺญุตญาณสฺส ปฏิลาภา สจฺฉิกา ปญฺญตฺติ, ยทิทํ ภควาติ อชิตาติ ภควา.}

\prose{\textbf{เววิจฺฉา ปมาทา นปฺปกาสตี}ติ เววิจฺฉํ วุจฺจติ ปญฺจ มจฺฉริยานิ อาวาสมจฺฉริยํ กุลมจฺฉริยํ ลาภมจฺฉริยํ, วณฺณมจฺฉริยํ ธมฺมมจฺฉริยํ, ยํ เอวรูปํ มจฺเฉรํ มจฺฉรายนา มจฺฉรายิตตฺตํ เววิจฺฉํ กทริยํ กฏุกญฺจุกตา อคฺคหิตตฺตํ จิตฺตสฺส, อิทํ วุจฺจติ มจฺฉริยํ.\csromanspacer{} อปิ จ ขนฺธมจฺฉริยมฺปิ มจฺฉริยํ, ธาตุมจฺฉริยมฺปิ มจฺฉริยํ, อายตนมจฺฉริยมฺปิ มจฺฉริยํ, คาโห วุจฺจติ มจฺฉริยํ.\csromanspacer{} ปมาโท วตฺตพฺโพ, กายทุจฺจริเต วา วจีทุจฺจริเต วา}

\vspace{\tipitakaparskip}
\csromancenter{อชิตมาณวปุจฺฉานิทฺเทส มโนทุจฺจริเต วา ปญฺจสุ กามคุเณสุ วา จิตฺตสฺส โวสคฺโค\footnote{โวสฺสคฺโค (พหูสุ)} โวสคฺคานุปฺปทานํ กุสลานํ ธมฺมานํ ภาวนาย อสกฺกจฺจกิริยตา อสาตจฺจกิริยตา อนฏฺฐิตกิริยตา\footnote{อนิฏฺฐิตกิริยตา (ก) อภิ 2. 363 ปิฏฺเฐ.} โอลีนวุตฺติตา นิกฺขิตฺตจฺฉนฺทตา นิกฺขิตฺตธุรตา อนาเสวนา อภาวนา อพหุลีกมฺมํ อนธิฏฺฐานํ อนนุโยโค ปมาโท, โย เอวรูโป ปมาโท ปมชฺชนา ปมชฺชิตตฺตํ.\csromanspacer{} อยํ วุจฺจติ ปมาโท.\csromanspacer{}\csromanspacer{}\textbf{เววิจฺฉา ปมาทา นปฺปกาสตี}ติ อิมินา จ มจฺฉริเยน อิมินา จ ปมาเทน โลโก นปฺปกาสติ น ภาสติ น ตปติ น วิโรจติ น ญายติ น ปญฺญายตีติ เววิจฺฉา ปมาทา นปฺปกาสติ.}
\prose{\csromanfolio{27}\textbf{ชปฺปาภิเลปนํ พฺรูมี}ติ ชปฺปา วุจฺจติ ตณฺหา, โย ราโค สาราโค อนุนโย อนุโรโธ นนฺที\footnote{นนฺทิ (สฺยา)} นนฺทิราโค จิตฺตสฺส สาราโค อิจฺฉา มุจฺฉา อชฺโฌสานํ เคโธ ปลิเคโธ สงฺโค ปงฺโก เอชา มายา ชนิกา สญฺชนนี สิพฺพินี ชาลินี สริตา วิสตฺติกา สุตฺตํ วิสฏา\footnote{โสตฺตํ วิสตา (สฺยา)} อายูหนี ทุติยา ปณิธิ ภวเนตฺติ วนํ วนโถ สนฺถโว\footnote{สนฺธโว (ก) อภิ 2. 376 ปิฏฺเฐ.} สิเนโห อเปกฺขา ปฏิพนฺธุ อาสา อาสีสนา\footnote{อาสิํสนา (สฺยา)} อาสีสิตตฺตํ รูปาสา สทฺทาสา คนฺธาสา รสาสา โผฏฺฐพฺพาสา ลาภาสา ธนาสา ปุตฺตาสา ชีวิตาสา ชปฺปา ปชปฺปา อภิชปฺปา ชปฺปนา ชปฺปิตตฺตํ โลลุปฺปํ โลลุปฺปายนา โลลุปฺปายิตตฺตํ ปุจฺฉญฺชิกตา สาธุกมฺยตา อธมฺมราโค วิสมโลโภ นิกนฺติ นิกามนา ปตฺถนา ปิหนา สมฺปตฺถนา กามตณฺหา ภวตณฺหา วิภวตณฺหา รูปตณฺหา อรูปตณฺหา นิโรธตณฺหา รูปตณฺหา สทฺทตณฺหา คนฺธตณฺหา รสตณฺหา โผฏฺฐพฺพตณฺหา ธมฺมตณฺหา โอโฆ โยโค คนฺโถ อุปาทานํ อาวรณํ นีวรณํ ฉทนํ พนฺธนํ อุปกฺกิเลโส อนุสโย ปริยุฏฺฐานํ ลตา เววิจฺฉํ ทุกฺขมูลํ ทุกฺขนิทานํ ทุกฺขปฺปภโว มารปาโส มารพฬิสํ มารามิสํ มารวิสโย มารนิวาโส มารโคจโร มารพนฺธนํ ตณฺหานที ตณฺหาชาลํ ตณฺหาคทฺทุลํ ตณฺหาสมุทฺโท อภิชฺฌา โลโภ อกุสลมูลํ, อยํ วุจฺจติ ชปฺปา.\csromanspacer{} โลกสฺส เลปนํ ลคฺคนํ พนฺธนํ อุปกฺกิเลโส อิมาย ชปฺปาย โลโก ลิตฺโต สํลิตฺโต อุปลิตฺโต กิลิฏฺโฐ สํกิลิฏฺโฐ มกฺขิโต สํสฏฺโฐ ลคฺโค ลคฺคิโต ปลิพุทฺโธติ พฺรูมิ อาจิกฺขามิ เทเสมิ ปญฺญเปมิ ปฏฺฐเปมิ วิวรามิ วิภชามิ อุตฺตานีกโรมิ ปกาเสมีติ ชปฺปาภิเลปนํ พฺรูมิ.}

\prose{\csromanfolio{28}\textbf{ทุกฺข’มสฺส มหพฺภยนฺ}ติ \textbf{ทุกฺขนฺ}ติ ชาติทุกฺขํ ชราทุกฺขํ พฺยาธิทุกฺขํ มรณทุกฺขํ โสกปริเทวทุกฺข- โทมนสฺสุปายาสทุกฺขํ เนรยิกํ ทุกฺขํ ติรจฺฉานโยนิกํ ทุกฺขํ เปตฺติวิสยิกํ ทุกฺขํ มานุสิกํ ทุกฺขํ คพฺโภกฺกนฺติมูลกํ ทุกฺขํ คพฺภฏฺฐิติมูลกํ\footnote{คพฺเภฐิติมูลกํ (สฺยา, ก)} ทุกฺขํ คพฺภวุฏฺฐานมูลกํ ทุกฺขํ ชาตสฺสูปนิพนฺธกํ ทุกฺขํ ชาตสฺส ปราเธยฺยกํ ทุกฺขํ อตฺตูปกฺกมทุกฺขํ ปรูปกฺกมทุกฺขํ สงฺขารทุกฺขํ วิปริณามทุกฺขํ, จกฺขุโรโค โสตโรโค ฆานโรโค ชิวฺหาโรโค กายโรโค สีสโรโค กณฺณโรโค มุขโรโค ทนฺตโรโค กาโส สาโส ปินาโส ฑาโห\footnote{qอโห (สฺยา)} ชโร กุจฺฉิโรโค มุจฺฉา ปกฺขนฺทิกา สูลา วิสูจิกา กุฏฺฐํ คณฺโฑ กิลาโส โสโส อปมาโร ททฺทุ กณฺฑุ กจฺฉุ รขสา\footnote{รกฺขสา (ก)} วิตจฺฉิกา โลหิตปิตฺตํ\footnote{โลหิตํ ปิตฺตํ (พหูสุ)} มธุเมโห อํสา ปิฬกา ภคนฺทลา ปิตฺตสมุฏฺฐานา อาพาธา เสมฺหสมุฏฺฐานา อาพาธา วาตสมุฏฺฐานา อาพาธา สนฺนิปาติกา อาพาธา อุตุปริณามชา อาพาธา วิสมปริหารชา อาพาธา โอปกฺกมิกา อาพาธา กมฺมวิปากชา อาพาธา สีตํ อุณฺหํ ชิฆจฺฉา ปิปาสา อุจฺจาโร ปสฺสาโว ฑํสมกสวาตาตปสรีสปสมฺผสฺสํ ทุกฺขํ มาตุมรณํ ทุกฺขํ ปิตุมรณํ ทุกฺขํ ภาตุมรณํ ทุกฺขํ ภคินิมรณํ ทุกฺขํ ปุตฺตมรณํ ทุกฺขํ ธีตุมรณํ ทุกฺขํ ญาติพฺยสนํ ทุกฺขํ โรคพฺยสนํ ทุกฺขํ โภคพฺยสนํ ทุกฺขํ สีลพฺยสนํ ทุกฺขํ ทิฏฺฐิพฺยสนํ ทุกฺขํ เยสํ ธมฺมานํ อาทิโต สมุทาคมนํ ปญฺญายติ.\csromanspacer{} อตฺถงฺคมโต นิโรโธ ปญฺญายติ.\csromanspacer{} กมฺมสนฺนิสฺสิโต วิปาโก, วิปากสนฺนิสฺสิตํ กมฺมํ, นามสนฺนิสฺสิตํ รูปํ รูปสนฺนิสฺสิตํ นามํ, ชาติยา อนุคตํ ชราย อนุสฏํ พฺยาธินา อภิภูตํ มรเณน อพฺภาหตํ ทุกฺเข ปติฏฺฐิตํ อตาณํ อเลณํ อสรณํ อสรณีภูตํ.\csromanspacer{} อิทํ วุจฺจติ ทุกฺขํ, อิทํ ทุกฺขํ โลกสฺส ภยํ มหาภยํ ปีฬนํ ฆฏฺฏนํ อุปทฺทโว อุปสคฺโคติ ทุกฺขมสฺส มหพฺภยํ.\csromanspacer{} เตนาห ภควา\csromandash{}}

\prose{อวิชฺชาย นิวุโต โลโก, (อชิตาติ ภควา,)}

\prose{เววิจฺฉา ปมาทา นปฺปกาสติ.}

\csromangathameasure{ชปฺปาภิเลปนํ พฺรูมิ, ทุกฺข’มสฺส มหพฺภยนฺติ.}
\csromangathagroup{\csromangathastanza{ชปฺปาภิเลปนํ พฺรูมิ, ทุกฺข’มสฺส มหพฺภยนฺติ.}}

\csromanheader{2}{1. อชิตมาณวปุจฺฉานิทฺเทส}
\proseitem{3}{\csromanfolio{29}\textbf{สวนฺติ สพฺพธิ โสตา}, (\textbf{อิจฺจา}ยสฺมา อชิโต,)}

\prose{\textbf{โสตานํ กิํ นิวารณํ.}}

\csromangathameasure{โสตานํ สํวรํ พฺรูหิ, เกน โสตา ปิธียเร.}
\csromangathagroup{\csromangathastanza{โสตานํ สํวรํ พฺรูหิ, เกน โสตา ปิธียเร\footnote{ปิถิยฺยเร (สฺยา), ปิถียเร (สี-ฏฺฐ)}.}}

\prose{\textbf{สวนฺติ สพฺพธิ โสตา}ติ \textbf{โสตา}ติ ตณฺหาโสโต ทิฏฺฐิโสโต กิเลสโสโต ทุจฺจริตโสโต อวิชฺชาโสโต\textbf{.}\csromanspacer{} \textbf{สพฺพธี}ติ สพฺเพสุ อายตเนสุ\textbf{.}\csromanspacer{} \textbf{สวนฺตี}ติ สวนฺติ อาสวนฺติ สนฺทนฺติ ปวตฺตนฺติ\textbf{.}\csromanspacer{} จกฺขุโต รูเป สวนฺติ อาสวนฺติ สนฺทนฺติ ปวตฺตนฺติ\textbf{.}\csromanspacer{} โสตโต สทฺเท สวนฺติ\textbf{.}\csromanspacer{} ฆานโต คนฺเธ สวนฺติ\textbf{.}\csromanspacer{} ชิวฺหาโต รเส สวนฺติ\textbf{.}\csromanspacer{} กายโต โผฏฺฐพฺเพ สวนฺติ\textbf{.}\csromanspacer{} มนโต ธมฺเม สวนฺติ อาสวนฺติ สนฺทนฺติ ปวตฺตนฺติ\textbf{.}\csromanspacer{} จกฺขุโต รูปตณฺหา สวนฺติ อาสวนฺติ สนฺทนฺติ ปวตฺตนฺติ\textbf{.}\csromanspacer{} โสตโต สทฺทตณฺหา สวนฺติ อาสวนฺติ สนฺทนฺติ ปวตฺตนฺติ\textbf{.}\csromanspacer{} ฆานโต คนฺธตณฺหา สวนฺติ\textbf{.}\csromanspacer{} ชิวฺหาโต รสตณฺหา สวนฺติ\textbf{.}\csromanspacer{} กายโต โผฏฺฐพฺพตณฺหา สวนฺติ\textbf{.}\csromanspacer{} มนโต ธมฺมตณฺหา สวนฺติ อาสวนฺติ สนฺทนฺติ ปวตฺตนฺตีติ สวนฺติ สพฺพธิ \textbf{โสตา.}}

\prose{\textbf{อิจฺจายสฺมา อชิโต}ติ \textbf{อิจฺจา}ติ ปทสนฺธิ ฯเปฯ ปทานุปุพฺพตาเปตํ \textbf{อิจฺจา}ติ ฯเปฯ \textbf{อิจฺจา}ยสฺมา อชิโต\textbf{.}}

\prose{\textbf{โสตานํ กิํ นิวารณนฺ}ติ \textbf{โสตา}นํ กิํ อาวรณํ นีวรณํ สํวรณํ รกฺขนํ โคปนนฺติ \textbf{โสตา}นํ กิํ นิวารณํ\textbf{.}}

\prose{\textbf{โสตานํ สํวรํ พฺรูหี}ติ \textbf{โสตา}นํ อาวรณํ นีวรณํ สํวรณํ รกฺขนํ โคปนํ พฺรูหิ อาจิกฺข เทเสหิ ปญฺญเปหิ ปฏฺฐเปหิ วิวราหิ วิภชาหิ อุตฺตานีกโรหิ ปกาเสหีติ \textbf{โสตา}นํ สํวรํ พฺรูหิ\textbf{.}}

\prose{\textbf{เกน โสตา ปิธียเร}ติ เกน \textbf{โสตา} ปิธียนฺติ ปจฺฉิชฺชนฺติ น สวนฺติ น อาสวนฺติ น สนฺทนฺติ นปฺปวตฺตนฺตีติ เกน \textbf{โสตา} ปิธียเร\textbf{.}\csromanspacer{} เตนาห โส พฺราหฺมโณ\csromandash{}}

\prose{\textbf{สวนฺติ สพฺพธิ โสตา}, (\textbf{อิจฺจา}ยสฺมา อชิโต,)}

\prose{\textbf{โสตานํ กิํ นิวารณํ.}}

\csromangathameasure{โสตานํ สํวรํ พฺรูหิ, เกน โสตา ปิธียเร.}
\csromangathagroup{\csromangathastanza{โสตานํ สํวรํ พฺรูหิ, เกน โสตา ปิธียเร.}}

\proseitem{4}{\csromanfolio{30}ยานิ โสตานิ โลกสฺมิํ, (\textbf{อชิตา}ติ ภควา,)}

\prose{\textbf{สติ เตสํ นิวารณํ.}}

\csromangathameasure{โสตานํ สํวรํ พฺรูมิ, ปญฺญาเยเต ปิธียเร.}
\csromangathagroup{\csromangathastanza{โสตานํ สํวรํ พฺรูมิ, ปญฺญาเยเต ปิธียเร.}}

\prose{\textbf{ยานิ โสตานิ โลกสฺมินฺ}ติ ยานิ เอตานิ โสตานิ มยา กิตฺติตานิ ปกิตฺติตานิ อาจิกฺขิตานิ เทสิตานิ ปญฺญปิตานิ ปฏฺฐปิตานิ วิวริตานิ วิภชิตานิ\footnote{วิภตฺตานิ (ก)} อุตฺตานีกตานิ ปกาสิตานิ.\csromanspacer{} เสยฺยถิทํ\footnote{เสยฺยถีทํ (สฺยา)}, ตณฺหาโสโต ทิฏฺฐิโสโต กิเลสโสโต ทุจฺจริตโสโต อวิชฺชาโสโต.\csromanspacer{} \textbf{โลกสฺมินฺ}ติ อปายโลเก มนุสฺสโลเก เทวโลเก ขนฺธโลเก ธาตุโลเก อายตนโลเกติ ยานิ โสตานิ โลกสฺมิํ.\csromanspacer{} \textbf{อชิตา}ติ ภควา ตํ พฺราหฺมณํ นาเมน อาลปติ.}

\prose{\textbf{สติ เตสํ นิวารณนฺ}ติ \textbf{สตี}ติ ยา สติ อนุสฺสติ ปฏิสฺสติ สติ สรณตา ธารณตา อปิลาปนตา อสมฺมุสฺสนตา, สติ สตินฺทฺริยํ สติพลํ สมฺมาสติ สติสมฺโพชฺฌงฺโค เอกายนมคฺโค.\csromanspacer{} อยํ วุจฺจติ สติ.\csromanspacer{} \textbf{นิวารณนฺ}ติ อาวรณํ นีวรณํ สํวรณํ รกฺขนํ โคปนนฺติ สติ เตสํ นิวารณํ.}

\prose{\textbf{โสตานํ สํวรํ พฺรูมี}ติ โสตานํ อาวรณํ นีวรณํ สํวรณํ รกฺขนํ โคปนํ พฺรูมิ อาจิกฺขามิ ฯเปฯ อุตฺตานีกโรมิ ปกาเสมีติ โสตานํ สํวรํ พฺรูมิ.}

\prose{\textbf{ปญฺญาเยเต ปิธียเร}ติ \textbf{ปญฺญา}ติ ยา \textbf{ปญฺญา} ปชานนา ฯเปฯ อโมโห ธมฺมวิจโย สมฺมาทิฏฺฐิ.\csromanspacer{} \textbf{ปญฺญาเยเต ปิธียเร}ติ \textbf{ปญฺญา}เยเต โสตา ปิธียนฺติ ปจฺฉิชฺชนฺติ น สวนฺติ น อาสวนฺติ น สนฺทนฺติ นปฺปวตฺตนฺติ. “สพฺเพ สงฺขารา อนิจฺจา”ติ ชานโต ปสฺสโต \textbf{ปญฺญา}เยเต โสตา ปิธียนฺติ ปจฺฉิชฺชนฺติ น สวนฺติ น อาสวนฺติ น สนฺทนฺติ นปฺปวตฺตนฺติ. “สพฺเพ สงฺขารา ทุกฺขา”ติ ชานโต ปสฺสโต \textbf{ปญฺญา}เยเต โสตา ปิธียนฺติ ปจฺฉิชฺชนฺติ น สวนฺติ น อาสวนฺติ น สนฺทนฺติ นปฺปวตฺตนฺติ. “สพฺเพ สงฺขารา อนตฺตา”ติ ชานโต ปสฺสโต \textbf{ปญฺญา}เยเต โสตา ปิธียนฺติ ปจฺฉิชฺชนฺติ น สวนฺติ น อาสวนฺติ น สนฺทนฺติ นปฺปวตฺตนฺติ. “อวิชฺชาปจฺจยา สงฺขารา”ติ ชานโต ปสฺสโต \textbf{ปญฺญา}เยเต โสตา ปิธียนฺติ}

\vspace{\tipitakaparskip}
\csromancenter{อชิตมาณวปุจฺฉานิทฺเทส ปจฺฉิชฺชนฺติ น สวนฺติ น อาสวนฺติ น สนฺทนฺติ นปฺปวตฺตนฺติ. “สงฺขารปจฺจยา วิญฺญาณนฺ”ติ. “วิญฺญาณปจฺจยา นามรูปนฺ”ติ. “นามรูปปจฺจยา สฬายตนนฺ”ติ. “สฬายตนปจฺจยา ผสฺโส”ติ. “ผสฺสปจฺจยา เวทนา”ติ. “เวทนาปจฺจยา ตณฺหา”ติ. “ตณฺหาปจฺจยา อุปาทานนฺ”ติ. “อุปาทานปจฺจยา ภโว”ติ. “ภวปจฺจยา ชาตี”ติ. “ชาติปจฺจยา ชรามรณนฺ”ติ ชานโต ปสฺสโต ปญฺญาเยเต โสตา ปิธียนฺติ ปจฺฉิชฺชนฺติ น สวนฺติ น อาสวนฺติ น สนฺทนฺติ นปฺปวตฺตนฺติ. “อวิชฺชานิโรธา สงฺขารนิโรโธ”ติ. “สงฺขารนิโรธา วิญฺญาณนิโรโธ”ติ. “วิญฺญาณนิโรธา นามรูปนิโรโธ”ติ. “นามรูปนิโรธา สฬายตนนิโรโธ”ติ “สฬายตนนิโรธา ผสฺสนิโรโธ”ติ. “ผสฺสนิโรธา เวทนานิโรโธ”ติ. “เวทนานิโรธา ตณฺหานิโรโธ”ติ. “ตณฺหานิโรธา อุปาทานนิโรโธ”ติ. “อุปาทานนิโรธา ภวนิโรโธ”ติ. “ภวนิโรธา ชาตินิโรโธ”ติ. “ชาตินิโรธา ชรามรณนิโรโธ”ติ ชานโต ปสฺสโต ปญฺญาเยเต โสตา ปิธียนฺติ ปจฺฉิชฺชนฺติ น สวนฺติ น อาสวนฺติ น สนฺทนฺติ นปฺปวตฺตนฺติ. “อิทํ ทุกฺขนฺ”ติ. “อยํ ทุกฺขสมุทโย”ติ. “อยํ ทุกฺขนิโรโธ”ติ. “อยํ ทุกฺขนิโรธคามินี ปฏิปทา”ติ ชานโต ปสฺสโต ปญฺญาเยเต โสตา ปิธียนฺติ ปจฺฉิชฺชนฺติ น สวนฺติ น อาสวนฺติ น สนฺทนฺติ นปฺปวตฺตนฺติ. “อิเม ธมฺมา อาสวา”ติ. “อยํ อาสวสมุทโย”ติ. “อยํ อาสวนิโรโธ”ติ. “อยํ อาสวนิโรธคามินี ปฏิปทา”ติ ชานโต ปสฺสโต ปญฺญาเยเต โสตา ปิธียนฺติ ปจฺฉิชฺชนฺติ น สวนฺติ น อาสวนฺติ น สนฺทนฺติ นปฺปวตฺตนฺติ. “อิเม ธมฺมา อภิญฺเญยฺยา”ติ. “อิเม ธมฺมา ปริญฺเญยฺยา”ติ. “อิเม ธมฺมา ปหาตพฺพา”ติ. “อิเม ธมฺมา ภาเวตพฺพา”ติ. “อิเม ธมฺมา สจฺฉิกาตพฺพา”ติ ชานโต ปสฺสโต ปญฺญาเยเต โสตา ปิธิยนฺติ ปจฺฉิชฺชนฺติ น สวนฺติ น อาสวนฺติ น สนฺทนฺติ นปฺปวตฺตนฺติ.\csromanspacer{} ฉนฺนํ ผสฺสายตนานํ สมุทยญฺจ อตฺถงฺคมญฺจ อสฺสาทญฺจ อาทีนวญฺจ นิสฺสรณญฺจ ชานโต ปสฺสโต ปญฺญาเยเต โสตา ปิธียนฺติ ปจฺฉิชฺชนฺติ น สวนฺติ น อาสวนฺติ น สนฺทนฺติ นปฺปวตฺตนฺติ.\csromanspacer{} ปญฺจนฺนํ อุปาทานกฺขนฺธานํ สมุทยญฺจ อตฺถงฺคมญฺจ อสฺสาทญฺจ อาทีนวญฺจ นิสฺสรณญฺจ ชานโต ปสฺสโต.\csromanspacer{} จตุนฺนํ มหาภูตานํ สมุทยญฺจ อตฺถงฺคมญฺจ อสฺสาทญฺจ อาทีนวญฺจ นิสฺสรณญฺจ ชานโต ปสฺสโต. “ยํ กิญฺจิ สมุทยธมฺมํ, สพฺพํ ตํ นิโรธธมฺมนฺ”ติ ชานโต ปสฺสโต ปญฺญาเยเต โสตา ปิธียนฺติ ปจฺฉิชฺชนฺติ น สวนฺติ น}
\prosecont{\csromanfolio{32}อาสวนฺติ น สนฺทนฺติ นปฺปวตฺตนฺตีติ \textbf{ปญฺญา}เยเต ปิธียเร.\csromanspacer{} เตนาห ภควา\csromandash{}}

\prose{ยานิ โสตานิ โลกสฺมิํ, (อชิตาติ ภควา,)}

\prose{สติ เตสํ นิวารณํ.}

\csromangathameasure{โสตานํ สํวรํ พฺรูมิ, ปญฺญาเยเต ปิธียเรติ.}
\csromangathagroup{\csromangathastanza{โสตานํ สํวรํ พฺรูมิ, ปญฺญาเยเต ปิธียเรติ.}}

\proseitem{5}{ปญฺญา เจว สติ จาปิ, (อิจฺจายสฺมา อชิโต,)}

\prose{\textbf{นามรูปญฺจ มาริส.}}

\csromangathameasure{เอตํ เม ปุฏฺโฐ ปพฺรูหิ, กตฺเถตํ อุปรุชฺฌติ.}
\csromangathagroup{\csromangathastanza{เอตํ เม ปุฏฺโฐ ปพฺรูหิ, กตฺเถตํ อุปรุชฺฌติ.}}

\prose{\textbf{ปญฺญา เจว สติ จาปี}ติ \textbf{ปญฺญา}ติ ยา \textbf{ปญฺญา} ปชานนา วิจโย ปวิจโย ธมฺมวิจโย สลฺลกฺขณา อุปลกฺขณา ปจฺจุปลกฺขณา ปณฺฑิจฺจํ โกสลฺลํ เนปุญฺญํ เวภพฺยา จินฺตา อุปปริกฺขา ภูรี\footnote{ภูริ (ก)} เมธา ปริณายิกา วิปสฺสนา สมฺปชญฺญํ ปโตโท \textbf{ปญฺญา} ปญฺญินฺทฺริยํ \textbf{ปญฺญา}พลํ \textbf{ปญฺญา}สตฺถํ \textbf{ปญฺญา}ปาสาโท \textbf{ปญฺญา}-อาโลโก \textbf{ปญฺญา}-โอภาโส \textbf{ปญฺญา}ปชฺโชโต \textbf{ปญฺญา}รตนํ อโมโห ธมฺมวิจโย สมฺมาทิฏฺฐิ, \textbf{สตี}ติ ยา สติ อนุสฺสติ ฯเปฯ สมฺมา\textbf{สตี}ติ \textbf{ปญฺญา} เจว สติจาปิ, อิจฺจายสฺมา อชิโต.}

\prose{\textbf{นามรูปญฺจ มาริสา}ติ \textbf{นามนฺ}ติ จตฺตาโร อรูปิโน ขนฺธา.\csromanspacer{} \textbf{รูปนฺ}ติ จตฺตาโร จ มหาภูตา จตุนฺนญฺจ มหาภูตานํ อุปาทายรูปํ.\csromanspacer{} \textbf{มาริสา}ติ ปิยวจนํ ครุวจนํ สคารวสปฺปติสฺสาธิวจนเมตํ มาริสาติ นามรูปญฺจ มาริส.}

\prose{\textbf{เอตํ เม ปุฏฺโฐ ปพฺรูหี}ติ \textbf{เอตํ เม}ติ ยํ ปุจฺฉามิ ยํ ยาจามิ ยํ อชฺเฌสามิ ยํ ปสาเทมิ.\csromanspacer{} \textbf{ปุฏฺโฐ}ติ ปุจฺฉิโต ยาจิโต อชฺเฌสิโต ปสาทิโต.\csromanspacer{} \textbf{ปพฺรูหี}ติ พฺรูหิ อาจิกฺขาหิ เทเสหิ ปญฺญเปหิ ปฏฺฐเปหิ วิวราหิ วิภชาหิ\footnote{วิวเรหิ วิภเชหิ (ก)} อุตฺตานีกโรหิ ปกาเสหีติ \textbf{เอตํ เม} ปุฏฺโฐ ปพฺรูหิ.}

\prose{\textbf{กตฺเถตํ อุปรุชฺฌตี}ติ กตฺเถตํ นิรุชฺฌติ วูปสมฺมติ อตฺถํ คจฺฉติ ปฏิปฺปสฺสมฺภตีติ กตฺเถตํ อุปรุชฺฌติ.\csromanspacer{} เตนาห โส พฺราหฺมโณ\csromandash{}}

\csromanheader{2}{1. อชิตมาณวปุจฺฉานิทฺเทส}
\prose{\csromanfolio{33}ปญฺญา เจว สติ จาปิ, (อิจฺจายสฺมา อชิโต,)}

\prose{นามรูปญฺจ มาริส.}

\csromangathameasure{เอตํ เม ปุฏฺโฐ ปพฺรูหิ, กตฺเถตํ อุปรุชฺฌตีติ.}
\csromangathagroup{\csromangathastanza{เอตํ เม ปุฏฺโฐ ปพฺรูหิ, กตฺเถตํ อุปรุชฺฌตีติ.}}

\proseitem{6}{ยเมตํ ปญฺหํ อปุจฺฉิ, \textbf{อชิต ตํ วทามิ เต}.}

\csromangathameasure{ยตฺถ นามญฺจ รูปญฺจ, อเสสํ อุปรุชฺฌติ. \\ วิญฺญาณสฺส นิโรเธน, เอตฺเถตํ อุปรุชฺฌติ.}
\csromangathagroup{\csromangathastanza{ยตฺถ นามญฺจ รูปญฺจ, อเสสํ อุปรุชฺฌติ. \\ วิญฺญาณสฺส นิโรเธน, เอตฺเถตํ อุปรุชฺฌติ.}}

\prose{\textbf{ยเมตํ ปญฺหํ อปุจฺฉี}ติ \textbf{ยเมตนฺ}ติ ปญฺญญฺจ สติญฺจ นามรูปญฺจ.\csromanspacer{} \textbf{อปุจฺฉี}ติ อปุจฺฉสิ ยาจสิ อชฺเฌสสิ\footnote{อชฺเฌสิ (ก)} ปสาเทสีติ ยเมตํ ปญฺหํ อปุจฺฉิ.}

\prose{\textbf{อชิต ตํ วทามิ เต}ติ \textbf{อชิตา}ติ ภควา ตํ พฺราหฺมณํ นาเมน อาลปติ.\csromanspacer{} \textbf{ตนฺ}ติ ปญฺญญฺจ สติญฺจ นามรูปญฺจ.\csromanspacer{} \textbf{วทามี}ติ วทามิ อาจิกฺขามิ เทเสมิ ปญฺญเปมิ ปฏฺฐเปมิ วิวรามิ วิภชามิ อุตฺตานีกโรมิ ปกาเสมีติ \textbf{อชิต ตํ วทามิ เต}.}

\prose{\textbf{ยตฺถ นามญฺจ รูปญฺจ, อเสสํ อุปรุชฺฌตี}ติ \textbf{นามนฺ}ติ จตฺตาโร อรูปิโน ขนฺธา.\csromanspacer{} \textbf{รูปนฺ}ติ จตฺตาโร จ มหาภูตา จตุนฺนญฺจ มหาภูตานํ อุปาทายรูปํ.\csromanspacer{} \textbf{อเสสนฺ}ติ สพฺเพน สพฺพํ สพฺพถา สพฺพํ อเสสํ นิสฺเสสํ ปริยาทิยนวจนเมตํ\footnote{ปริยาทายวจนเมตํ (สฺยา, ก)} อเสสนฺติ.\csromanspacer{} \textbf{อุปรุชฺฌตี}ติ นิรุชฺฌติ วูปสมฺมติ อตฺถํ คจฺฉติ ปฏิปฺปสฺสมฺภตีติ ยตฺถ นามญฺจ รูปญฺจ, อเสสํ อุปรุชฺฌติ.}

\prose{\textbf{วิญฺญาณสฺส นิโรเธน, เอตฺเถตํ อุปรุชฺฌตี}ติ โสตาปตฺติมคฺคญาเณน อภิสงฺขารวิญฺญาณสฺส นิโรเธน สตฺต ภเว ฐเปตฺวา อนมตคฺเค สํสาเร เย อุปฺปชฺเชยฺยุํ นามญฺจ รูปญฺจ, เอตฺเถเต นิรุชฺฌนฺติ วูปสมฺมนฺติ อตฺถํ คจฺฉนฺติ ปฏิปฺปสฺสมฺภนฺติ.\csromanspacer{} สกทาคามิมคฺคญาเณน อภิสงฺขารวิญฺญาณสฺส นิโรเธน เทฺว ภเว ฐเปตฺวา ปญฺจสุ ภเวสุ เย อุปฺปชฺเชยฺยุํ นามญฺจ รูปญฺจ.\csromanspacer{} เอตฺเถเต นิรุชฺฌนฺติ วูปสมฺมนฺติ อตฺถํ คจฺฉนฺติ ปฏิปฺปสฺสมฺภนฺติ.\csromanspacer{} อนาคามิมคฺคญาเณน อภิสงฺขารวิญฺญาณสฺส นิโรเธน เอกํ ภวํ ฐเปตฺวา รูปธาตุยา วา อรูปธาตุยา วา เย อุปฺปชฺเชยฺยุํ นามญฺจ รูปญฺจ, เอตฺเถเต นิรุชฺฌนฺติ วูปสมฺมนฺติ อตฺถํ คจฺฉนฺติ ปฏิปฺปสฺสมฺภนฺติ.\csromanspacer{} อรหตฺตมคฺคญาเณน อภิสงฺขารวิญฺญาณสฺส นิโรเธน เย อุปฺปชฺเชยฺยุํ นามญฺจ รูปญฺจ, เอตฺเถเต นิรุชฺฌนฺติ วูปสมฺมนฺติ อตฺถํ คจฺฉนฺติ \csromanfolio{34}ปฏิปฺปสฺสมฺภนฺติ.\csromanspacer{} อรหโต อนุปาทิเสสาย นิพฺพานธาตุยา ปรินิพฺพายนฺตสฺส จริมวิญฺญาณสฺส นิโรเธน ปญฺญา จ สติ จ นามญฺจ รูปญฺจ, เอตฺเถเต นิรุชฺฌนฺติ วูปสมฺมนฺติ อตฺถํ คจฺฉนฺติ ปฏิปฺปสฺสมฺภนฺตีติ วิญฺญาณสฺส นิโรเธน, เอตฺเถตํ อุปรุชฺฌติ.\csromanspacer{} เตนาห ภควา\csromandash{}}

\csromangathameasure{ยเมตํ ปญฺหํ อปุจฺฉิ, อชิต ตํ วทามิ เต. \\ ยตฺถ นามญฺจ รูปญฺจ, อเสสํ อุปรุชฺฌติ. \\ วิญฺญาณสฺส นิโรเธน, เอตฺเถตํ อุปรุชฺฌตีติ.}
\csromangathagroup{\csromangathastanza{ยเมตํ ปญฺหํ อปุจฺฉิ, อชิต ตํ วทามิ เต. \\ ยตฺถ นามญฺจ รูปญฺจ, อเสสํ อุปรุชฺฌติ. \\ วิญฺญาณสฺส นิโรเธน, เอตฺเถตํ อุปรุชฺฌตีติ.}}

\proseitem{7}{\textbf{เย จ สงฺขาตธมฺมาเส}, เย จ \textbf{เสขา}\footnote{เสกฺขา (สฺยา, ก)} \textbf{ปุถู อิธ.}}

\csromangathameasure{เตสํ เม นิปโก อิริยํ, ปุฏฺโฐ ปพฺรูหิ มาริส.}
\csromangathagroup{\csromangathastanza{เตสํ เม นิปโก อิริยํ, ปุฏฺโฐ ปพฺรูหิ มาริส.}}

\prose{\textbf{เย จ สงฺขาตธมฺมาเส}ติ สงฺขาตธมฺมา วุจฺจนฺติ อรหนฺโต ขีณาสวา, กิํ การณา สงฺขาตธมฺมา วุจฺจนฺติ อรหนฺโต ขีณาสวา.\csromanspacer{} เต สงฺขาตธมฺมา ญาตธมฺมา ตุลิตธมฺมา ตีริตธมฺมา วิภูตธมฺมา วิภาวิตธมฺมา. “สพฺเพ สงฺขารา อนิจฺจา”ติ สงฺขาตธมฺมา ญาตธมฺมา ตุลิตธมฺมา ตีริตธมฺมา วิภูตธมฺมา วิภาวิตธมฺมา. “สพฺเพ สงฺขารา ทุกฺขา”ติ สงฺขาตธมฺมา ฯเปฯ. “สพฺเพ ธมฺมา อนตฺตา”ติ สงฺขาตธมฺมา. “อวิชฺชาปจฺจยา สงฺขารา”ติ สงฺขาตธมฺมา ฯเปฯ. “ยํ กิญฺจิ สมุทยธมฺมํ, สพฺพํ ตํ นิโรธธมฺมนฺ”ติ สงฺขาตธมฺมา ญาตธมฺมา ตุลิตธมฺมา ตีริตธมฺมา วิภูตธมฺมา วิภาวิตธมฺมา.\csromanspacer{} อถ วา เตสํ ขนฺธา สงฺขาตา ธาตุโย สงฺขาตา อายตนานิ สงฺขาตา คติโย สงฺขาตา อุปปตฺติโย สงฺขาตา ปฏิสนฺธิ สงฺขาตา ภวา สงฺขาตา สํสารา สงฺขาตา วฏฺฏา สงฺขาตา.\csromanspacer{} อถ วา เต ขนฺธปริยนฺเต ฐิตา ธาตุปริยนฺเต ฐิตา อายตนปริยนฺเต ฐิตา คติปริยนฺเต ฐิตา อุปปตฺติปริยนฺเต ฐิตา ปฏิสนฺธิปริยนฺเต ฐิตา ภวปริยนฺเต ฐิตา สํสารปริยนฺเต ฐิตา วฏฺฏปริยนฺเต ฐิตา อนฺติเม ภเว ฐิตา อนฺติเม สมุสฺสเย ฐิตา อนฺติมเทหธรา อรหนฺโต.}

\csromangathameasure{เตสํ จายํ ปจฺฉิมโก, จริโมยํ สมุสฺสโย. \\ ชาติมรณสํสาโร, นตฺถิ เนสํ ปุนพฺภโวติ.}
\csromangathagroup{\csromangathastanza{เตสํ จายํ\footnote{ยายํ (ก)} ปจฺฉิมโก, จริโมยํ สมุสฺสโย. \\ ชาติมรณสํสาโร, นตฺถิ เนสํ ปุนพฺภโวติ.}}

\prose{ตํการณา สงฺขาตธมฺมา วุจฺจนฺติ อรหนฺโต ขีณาสวาติ.\csromanspacer{} \textbf{เย} จ \textbf{สงฺขาตธมฺมาเส, เย จ เสขา ปุถู อิธา}ติ \textbf{เสขา}ติ กิํ การณา}

\vspace{\tipitakaparskip}
\csromancenter{อชิตมาณวปุจฺฉานิทฺเทส วุจฺจนฺติ เสขา, สิกฺขนฺตีติ เสขา, กิญฺจ สิกฺขนฺติ, อธิสีลมฺปิ สิกฺขนฺติ.\csromanspacer{} อธิจิตฺตมฺปิ สิกฺขนฺติ.\csromanspacer{} อธิปญฺญมฺปิ สิกฺขนฺติ.\csromanspacer{} กตมา อธิสีลสิกฺขา.\csromanspacer{} อิธ ภิกฺขุ สีลวา โหติ ปาติโมกฺขสํวรสํวุโต วิหรติ อาจารโคจรสมฺปนฺโน อณุมตฺเตสุ วชฺเชสุ ภยทสฺสาวี, สมาทาย สิกฺขติ สิกฺขาปเทสุ ขุทฺทโก สีลกฺขนฺโธ มหนฺโต สีลกฺขนฺโธ สีลํ ปติฏฺฐา อาทิ จรณํ สํยโม สํวโร มุขํ ปมุขํ กุสลานํ ธมฺมานํ สมาปตฺติยา.\csromanspacer{} อยํ อธิสีลสิกฺขา.}
\prose{\csromanfolio{35}กตมา อธิจิตฺตสิกฺขา.\csromanspacer{} อิธ ภิกฺขุ วิวิจฺเจว กาเมหิ ฯเปฯ ปฐมํ ฌานํ.\csromanspacer{} ทุติยํ ฌานํ.\csromanspacer{} ตติยํ ฌานํ.\csromanspacer{} จตุตฺถํ ฌานํ อุปสมฺปชฺช วิหรติ.\csromanspacer{} อยํ อธิจิตฺตสิกฺขา.}

\prose{กตมา อธิปญฺญาสิกฺขา.\csromanspacer{} อิธ ภิกฺขุ ปญฺญวา โหติ อุทยตฺถคามินิยา ปญฺญาย สมนฺนาคโต อริยาย นิพฺเพธิกาย สมฺมา ทุกฺขกฺขยคามินิยา.\csromanspacer{} โส “อิทํ ทุกฺขนฺ”ติ ยถาภูตํ ปชานาติ. “อยํ ทุกฺขสมุทโย”ติ. “อยํ ทุกฺขนิโรโธ”ติ. “อยํ ทุกฺขนิโรธคามินี ปฏิปทา”ติ ยถาภูตํ ปชานาติ. “อิเม อาสวา”ติ. “อยํ อาสวสมุทโย”ติ. “อยํ อาสวนิโรโธ”ติ. “อยํ อาสวนิโรธคามินี ปฏิปทา”ติ ยถาภูตํ ปชานาติ.\csromanspacer{} อยํ อธิปญฺญาสิกฺขา.\csromanspacer{} อิมา ติสฺโส สิกฺขาโย อาวชฺชนฺตา สิกฺขนฺติ.\csromanspacer{} ชานนฺตา สิกฺขนฺติ.\csromanspacer{} ปสฺสนฺตา สิกฺขนฺติ.\csromanspacer{} จิตฺตํ อธิฏฺฐหนฺตา สิกฺขนฺติ.\csromanspacer{} สทฺธาย อธิมุจฺจนฺตา สิกฺขนฺติ.\csromanspacer{} วีริยํ\footnote{วิริยํ (สฺยา)} ปคฺคณฺหนฺตา สิกฺขนฺติ.\csromanspacer{} สติํ อุปฏฺฐเปนฺตา สิกฺขนฺติ.\csromanspacer{} จิตฺตํ สมาทหนฺตา สิกฺขนฺติ.\csromanspacer{} ปญฺญาย ปชานนฺตา สิกฺขนฺติ.\csromanspacer{} อภิญฺเญยฺยํ อภิชานนฺตา สิกฺขนฺติ.\csromanspacer{} ปริญฺเญยฺยํ ปริชานนฺตา สิกฺขนฺติ.\csromanspacer{} ปหาตพฺพํ ปชหนฺตา สิกฺขนฺติ.\csromanspacer{} ภาเวตพฺพํ ภาเวนฺตา สิกฺขนฺติ.\csromanspacer{} สจฺฉิกาตพฺพํ สจฺฉิกโรนฺตา สิกฺขนฺติ อาจรนฺติ สมาจรนฺติ สมาทาย วตฺตนฺติ.\csromanspacer{} ตํการณา วุจฺจนฺติ เสขา.\csromanspacer{} \textbf{ปุถู}ติ พหุกา.\csromanspacer{} เอเต เสขา โสตาปนฺนา จ ปฏิปนฺนา จ สกทาคามิโน จ ปฏิปนฺนา จ อนาคามิโน จ ปฏิปนฺนา จ อรหนฺโต จ ปฏิปนฺนา จ.\csromanspacer{} \textbf{อิธา}ติ อิมิสฺสา ทิฏฺฐิยา อิมิสฺสา ขนฺติยา อิมิสฺสา รุจิยา อิมสฺมิํ อาทาเย อิมสฺมิํ ธมฺเม อิมสฺมิํ วินเย อิมสฺมิํ ธมฺมวินเย อิมสฺมิํ ปาวจเน อิมสฺมิํ พฺรหฺมจริเย อิมสฺมิํ สตฺถุสาสเน อิมสฺมิํ อตฺตภาเว อิมสฺมิํ มนุสฺสโลเกติ เย จ เสขา ปุถู อิธ.}

\prose{\csromanfolio{36}\textbf{เตสํ เม นิปโก อิริยํ, ปุฏฺโฐ ปพฺรูหิ มาริสา}ติ ตฺวมฺปิ นิปโก ปณฺฑิโต ปญฺญวา พุทฺธิมา ญาณี เมธาวี.\csromanspacer{} เตสํ สงฺขาตธมฺมานญฺจ เสกฺขานญฺจ อิริยํ จริยํ วุตฺติ ปวตฺติ อาจรํ โคจรํ วิหารํ ปฏิปทํ.\csromanspacer{} \textbf{ปุฏฺโฐ}ติ ปุจฺฉิโต ยาจิโต อชฺเฌสิโต ปสาทิโต.\csromanspacer{} \textbf{ปพฺรูหี}ติ พฺรูหิ อาจิกฺขาหิ เทเสหิ ปญฺญเปหิ ปฏฺฐเปหิ วิวราหิ วิภชาหิ อุตฺตานีกโรหิ ปกาเสหิ.\csromanspacer{} \textbf{มาริสา}ติ ปิยวจนํ ครุวจนํ สคารวสปฺปติสฺสาธิวจนเมตํ มาริสาติ เตสํ เม นิปโก อิริยํ, ปุฏฺโฐ ปพฺรูหิ มาริส.\csromanspacer{} เตนาห โส พฺราหฺมโณ\csromandash{}}

\csromangathameasure{เย จ สงฺขาตธมฺมาเส, เย จ เสขา ปุถู อิธ.}
\csromangathagroup{\csromangathastanza{เย จ สงฺขาตธมฺมาเส, เย จ เสขา ปุถู อิธ.}}

\prose{\textbf{เตสํ เม นิปโก อิริยํ, ปุฏฺโฐ ปพฺรูหิ มาริสา}ติ.}

\proseitem{8}{กาเมสุ นาภิคิชฺเฌยฺย, มนสา’นาวิโล สิยา.}

\csromangathameasure{กุสโล สพฺพธมฺมานํ, สโต ภิกฺขุ ปริพฺพเช.}
\csromangathagroup{\csromangathastanza{กุสโล สพฺพธมฺมานํ, สโต ภิกฺขุ ปริพฺพเช.}}

\prose{\textbf{กาเมสุ นาภิคิชฺเฌยฺยา}ติ \textbf{กามา}ติ อุทฺทานโต เทฺว \textbf{กามา} วตฺถุ\textbf{กามา} จ กิเลส\textbf{กามา} จ, กตเม วตฺถุ\textbf{กามา}.\csromanspacer{} มนาปิกา รูปา มนาปิกา สทฺทา มนาปิกา คนฺธา มนาปิกา รสา มนาปิกา โผฏฺฐพฺพา, อตฺถรณา ปาวุรณา\footnote{ปาปุรณา (สฺยา) 2. ราชฐานิโย (ก) 3. นตฺถิ สฺยา-โปตฺถํเก, ขุ 7. 2 ปิฏฺเฐปิ.} ทาสิทาสา อเชฬกา กุกฺกุฏสูกรา หตฺถิควาสฺสวฬวา เขตฺตํ วตฺถุ หิรญฺญํ สุวณฺณํ คามนิคมราชธานิโย2 รฏฺฐญฺจ ชนปโท จ โกโส จ โกฏฺฐาคารญฺจ, ยํ กิญฺจิ รชนียวตฺถุ วตฺถุ\textbf{กามา}.\csromanspacer{} อปิ จ อตีตา \textbf{กามา} อนาคตา \textbf{กามา} ปจฺจุปฺปนฺนา \textbf{กามา}, อชฺฌตฺตา \textbf{กามา} พหิทฺธา \textbf{กามา} อชฺฌตฺตพหิทฺธา \textbf{กามา}, หีนา \textbf{กามา} มชฺฌิมา \textbf{กามา} ปณีตา \textbf{กามา}, อาปายิกา \textbf{กามา} มานุสิกา \textbf{กามา} ทิพฺพา \textbf{กามา}, ปจฺจุปฏฺฐิตา \textbf{กามา}, นิมฺมิตา \textbf{กามา} ปรนิมฺมิตา \textbf{กามา}, ปริคฺคหิตา \textbf{กามา} อปริคฺคหิตา \textbf{กามา} มมายิตา \textbf{กามา} อมมายิตา \textbf{กามา}, สพฺเพปิ \textbf{กามา}วจรา ธมฺมา, สพฺเพปิ รูปาวจรา ธมฺมา, สพฺเพปิ อรูปาวจรา ธมฺมา, ตณฺหาวตฺถุกา ตณฺหารมฺมณา, กามนียฏฺเฐน รชนียฏฺเฐน มทนียฏฺเฐน รมณียฏฺเฐน3 \textbf{กามา}.\csromanspacer{} อิเม วุจฺจนฺติ วตฺถุ\textbf{กามา}.}

\prose{กตเม กิเลส\textbf{กามา}.\csromanspacer{} ฉนฺโท กาโม ราโค กาโม ฉนฺทราโค กาโม สงฺกปฺโป กาโม ราโค กาโม สงฺกปฺปราโค กาโม.\csromanspacer{} โย กาเมสุ กามจฺฉนฺโท กามราโค กามนนฺที กามตณฺหา}

\vspace{\tipitakaparskip}
\csromancenter{อชิตมาณวปุจฺฉานิทฺเทส กามสิเนโห กามปิปาสา กามปริฬาโห กามเคโธ กามมุจฺฉา กามชฺโฌสานํ กาโมโฆ กามโยโค กามุปาทานํ กามจฺฉนฺทนีวรณํ.}
\vspace{0.5\baselineskip}
\csromangathameasure{อทฺทสํ กาม เต มูลํ, สงฺกปฺปา กาม ชายสิ. \\ น ตํ สงฺกปฺปยิสฺสามิ, เอวํ กาม น เหหิสีติ.}
\csromangathagroup{\csromangathastanza{\csromanfolio{37}อทฺทสํ กาม เต มูลํ, สงฺกปฺปา กาม ชายสิ. \\ น ตํ สงฺกปฺปยิสฺสามิ, เอวํ กาม น เหหิสีติ.}}

\prose{อิเม วุจฺจนฺติ กิเลสกามา.\csromanspacer{} \textbf{เคโธ} วุจฺจติ ตณฺหา, โย ราโค สาราโค ฯเปฯ อภิชฺฌา โลโภ อกุสลมูลํ.\csromanspacer{} \textbf{กาเมสุ นาภิคิชฺเฌยฺยา}ติ กิเลสกาเมน วตฺถุกาเมสุ นาภิคิชฺเฌยฺย น ปลิคิชฺเฌยฺย น ปลิพุนฺเธยฺย\footnote{ปลิพุชฺเฌยฺย (สฺยา)} อคิทฺโธ อสฺส อคธิโต อมุจฺฉิโต อนชฺฌาปนฺโน\footnote{อนชฺโฌปนฺโน (สฺยา)} วีตเคโธ วิคตเคโธ จตฺตเคโธ วนฺตเคโธ มุตฺตเคโธ ปหีนเคโธ ปฏินิสฺสฏฺฐเคโธ วีตราโค วิคตราโค จตฺตราโค วนฺตราโค มุตฺตราโค ปหีนราโค ปฏินิสฺสฏฺฐราโค นิจฺฉาโต นิพฺพุโต สีติภูโต สุขปฺปฏิสํเวที พฺรหฺมภูเตน อตฺตนา วิหเรยฺยาติ กาเมสุ นาภิคิชฺเฌยฺย.}

\prose{\textbf{มนสา’นาวิโล สิยา}ติ \textbf{มโน}ติ ยํ จิตฺตํ \textbf{มโน} มานสํ หทยํ ปณฺฑรํ \textbf{มโน} มนายตนํ มนินฺทฺริยํ วิญฺญาณํ วิญฺญาณกฺขนฺโธ ตชฺชา \textbf{มโน}วิญฺญาณธาตุ.\csromanspacer{} กายทุจฺจริเตน จิตฺตํ อาวิลํ โหติ ลุฬิตํ เอริตํ ฆฏฺฏิตํ จลิตํ ภนฺตํ อวูปสนฺตํ.\csromanspacer{} วจีทุจฺจริเตน.\csromanspacer{} มโนทุจฺจริเตน.\csromanspacer{} ราเคน.\csromanspacer{} โทเสน.\csromanspacer{} โมเหน.\csromanspacer{} โกเธน.\csromanspacer{} อุปนาเหน.\csromanspacer{} มกฺเขน.\csromanspacer{} ปฬาเสน.\csromanspacer{} อิสฺสาย.\csromanspacer{} มจฺฉริเยน.\csromanspacer{} มายาย.\csromanspacer{} สาเฐยฺเยน.\csromanspacer{} ถมฺเภน.\csromanspacer{} สารมฺเภน.\csromanspacer{} มาเนน.\csromanspacer{} อติมาเนน.\csromanspacer{} มเทน.\csromanspacer{} ปมาเทน.\csromanspacer{} สพฺพกิเลเสหิ.\csromanspacer{} สพฺพทุจฺจริเตหิ.\csromanspacer{} สพฺพฑาเหหิ.\csromanspacer{} สพฺพปริฬาเหหิ.\csromanspacer{} สพฺพสนฺตาเปหิ.\csromanspacer{} สพฺพากุสลาภิสงฺขาเรหิ จิตฺตํ อาวิลํ โหติ ลุฬิตํ เอริตํ ฆฏฺฏิตํ จลิตํ ภนฺตํ อวูปสนฺตํ.\csromanspacer{} \textbf{มนสา’นาวิโล สิยา}ติ จิตฺเตน อนาวิโล สิยา อลุฬิโต อเนริโต อฆฏฺฏิโต อจลิโต อภนฺโต วูปสนฺโต อาวิลกเร กิเลเส ชเหยฺย ปชเหยฺย วิโนเทยฺย พฺยนฺตีกเรยฺย\footnote{พฺยนฺติํ กเรยฺย (ก)} อนภาวํ คเมยฺย, อาวิลกเรหิ กิเลเสหิ จ อารโต\footnote{อารโต อสฺส (ก) ขุ 7. 54 ปิฏฺเฐ ปสฺส.} วิรโต ปฏิวิรโต นิกฺขนฺโต นิสฺสโฏ วิปฺปมุตฺโต วิสญฺญุตฺโต วิมริยาทิกเตน เจตสา วิหเรยฺยาติ มนสฺสา’นาวิโล สิยา.}

\prose{\csromanfolio{38}\textbf{กุสโล สพฺพธมฺมานนฺ}ติ “สพฺเพ สงฺขารา อนิจฺจา”ติ กุสโล \textbf{สพฺพธมฺมา}นํ, “สพฺเพ สงฺขารา ทุกฺขา”ติ กุสโล \textbf{สพฺพธมฺมา}นํ, “สพฺเพ ธมฺมา อนตฺตา”ติ กุสโล \textbf{สพฺพธมฺมา}นํ ฯเปฯ “อวิชฺชาปจฺจยา สงฺขารา”ติ กุสโล \textbf{สพฺพธมฺมา}นํ. “ยํ กิญฺจิ สมุทยธมฺมํ, สพฺพํ ตํ นิโรธธมฺมนฺ”ติ กุสโล \textbf{สพฺพธมฺมา}นํ.\csromanspacer{} เอวมฺปิ กุสโล \textbf{สพฺพธมฺมา}นํ.}

\prose{อถ วา อนิจฺจโต กุสโล \textbf{สพฺพธมฺมา}นํ.\csromanspacer{} ทุกฺขโต.\csromanspacer{} โรคโต.\csromanspacer{} คณฺฑโต.\csromanspacer{} สลฺลโต.\csromanspacer{} อฆโต.\csromanspacer{} อาพาธโต.\csromanspacer{} ปรโต.\csromanspacer{} ปโลกโต. ¢ติโต.\csromanspacer{} อุปทฺทวโต.\csromanspacer{} ภยโต.\csromanspacer{} อุปสคฺคโต.\csromanspacer{} จลโต.\csromanspacer{} ปภงฺคุโต.\csromanspacer{} อทฺธุวโต\footnote{อธุวโต (ก) ขุ 7. 40 ปิฏฺเฐปิ.}.\csromanspacer{} อตาณโต.\csromanspacer{} อเลณโต.\csromanspacer{} อสรณโต.\csromanspacer{} อสรณีภูตโต.\csromanspacer{} ริตฺตโต.\csromanspacer{} ตุจฺฉโต.\csromanspacer{} สุญฺญโต.\csromanspacer{} อนตฺตโต.\csromanspacer{} อาทีนวโต.\csromanspacer{} วิปริณามธมฺมโต.\csromanspacer{} อสารกโต.\csromanspacer{} อฆมูลโต.\csromanspacer{} วธกโต.\csromanspacer{} วิภวโต.\csromanspacer{} สาสวโต.\csromanspacer{} สงฺขตโต.\csromanspacer{} มารามิ\textbf{สโต}.\csromanspacer{} ชาติธมฺมโต.\csromanspacer{} ชราธมฺมโต.\csromanspacer{} พฺยาธิธมฺมโต.\csromanspacer{} มรณธมฺมโต.\csromanspacer{} โสกปริเทวทุกฺขโทมนสฺสุปายาสธมฺมโต.\csromanspacer{} สํกิเลสิกธมฺมโต.\csromanspacer{} สมุทยโต.\csromanspacer{} อตฺถงฺคมโต.\csromanspacer{} อสฺสาทโต.\csromanspacer{} อาทีนวโต.\csromanspacer{} นิสฺสรณโต กุสโล \textbf{สพฺพธมฺมา}นํ.\csromanspacer{} เอวมฺปิ กุสโล \textbf{สพฺพธมฺมา}นํ.}

\prose{อถ วา ขนฺธกุสโล ธาตุกุสโล อายตนกุสโล ปฏิจฺจสมุปฺปาทกุสโล สติปฏฺฐานกุสโล สมฺมปฺปธานกุสโล อิทฺธิปาทกุสโล อินฺทฺริยกุสโล พลกุสโล โพชฺฌงฺคกุสโล มคฺคกุสโล ผลกุสโล นิพฺพานกุสโล.\csromanspacer{} เอวมฺปิ กุสโล \textbf{สพฺพธมฺมา}นํ.}

\prose{อถ วา \textbf{สพฺพธมฺมา} วุจฺจนฺติ ทฺวาทสายตนานิ.\csromanspacer{} จกฺขุ เจว\footnote{จกฺขุญฺเจว (ก)} รูปา จ, โสตญฺจ สทฺทา จ, ฆานญฺจ คนฺธา จ, ชิวฺหา จ รสา จ, กาโย จ โผฏฺฐพฺพา จ, มโน จ ธมฺมา จ.\csromanspacer{} ยโต จ อชฺฌตฺติกพาหิเรสุ อายตเนสุ ฉนฺทราโค ปหีโน โหติ อุจฺฉินฺนมูโล ตาลาวตฺถุกโต อนภาวํกโต\footnote{อนภาวงฺคโต (สฺยา)} อายติํ อนุปฺปาทธมฺโม, เอตฺตาวตาปิ กุสโล \textbf{สพฺพธมฺมา}นนฺติ กุสโล \textbf{สพฺพธมฺมา}นํ.}

\prose{\textbf{สโต ภิกฺขุ ปริพฺพเช}ติ \textbf{สโต}ติ จตูหิ การเณหิ \textbf{สโต}, กาเย กายานุปสฺสนาสติปฏฺฐานํ ภาเวนฺโต \textbf{สโต}, เวทนาสุ}

\vspace{\tipitakaparskip}
\csromancenter{อชิตมาณวปุจฺฉานิทฺเทส เวทนานุปสฺสนาสติปฏฺฐานํ ภเวนฺโต สโต, จิตฺเต จิตฺตานุปสฺสนาสติปฏฺฐานํ ภาเวนฺโต สโต, ธมฺเมสุ ธมฺมานุปสฺสนาสติปฏฺฐานํ ภาเวนฺโต สโต.}
\prose{\csromanfolio{39}อปเรหิปิ จตูหิ การเณหิ สโต, อสติปริวชฺชนาย สโต, สติกรณียานํ ธมฺมานํ กตตฺตา สโต, สติปริพนฺธานํ\footnote{สติปฏิปกฺขานํ (สฺยา) ขุ 7. 7 ปิฏฺเฐ.} ธมฺมานํ หตตฺตา สโต, สตินิมิตฺตานํ ธมฺมานํ อสมฺมุฏฺฐตฺตา\footnote{อปฺปมุฏฺฐตฺตา (สฺยา)} สโต.}

\prose{อปเรหิปิ จตูหิ การเณหิ สโต, สติยา สมนฺนาคตตฺตา สโต, สติยา วสิตตฺตา สโต, สติยา ปาคุญฺเญน สมนฺนาคตตฺตา สโต, สติยา อปจฺโจโรหณตาย สโต.}

\prose{อปเรหิปิ จตูหิ การเณหิ สโต, สติยา สมนฺนาคตตฺตา สโต, สนฺตตฺตา สโต, สมิตตฺตา สโต, สนฺตธมฺมสมนฺนาคตตฺตา สโต.\csromanspacer{} พุทฺธานุสฺสติยา สโต.\csromanspacer{} ธมฺมานุสฺสติยา สโต.\csromanspacer{} สํฆานุสฺสติยา สโต.\csromanspacer{} สีลานุสฺสติยา สโต.\csromanspacer{} จาคานุสฺสติยา สโต.\csromanspacer{} เทวตานุสฺสติยา สโต.\csromanspacer{} อานาปานสฺสติยา สโต.\csromanspacer{} มรณสฺสติยา สโต.\csromanspacer{} กายคตาสติยา สโต.\csromanspacer{} อุปสมานุสฺสติยา สโต.\csromanspacer{} ยา สติ อนุสฺสติ ฯเปฯ สมฺมาสติ สติสมฺโพชฺฌงฺโค เอกายนมคฺโค.\csromanspacer{} อยํ วุจฺจติ สติ.\csromanspacer{} อิมาย สติยา อุเปโต โหติ สมุเปโต อุปาคโต สมุปาคโต อุปปนฺโน สมุปปนฺโน\footnote{สมฺปนฺโน (ก)} สมนฺนาคโต, โส วุจฺจติ สโต.\csromanspacer{}\csromanspacer{}\textbf{ภิกฺขู}ติ สตฺตนฺนํ ธมฺมานํ ภินฺนตฺตา ภิกฺขุ.\csromanspacer{} สกฺกายทิฏฺฐิ ภินฺนา โหติ.\csromanspacer{} วิจิกิจฺฉา ภินฺนา โหติ.\csromanspacer{} สีลพฺพตปรามาโส ภินฺโน โหติ.\csromanspacer{} ราโค ภินฺโน โหติ.\csromanspacer{} โทโส ภินฺโน โหติ.\csromanspacer{} โมโห ภินฺโน โหติ.\csromanspacer{} มาโน ภินฺโน โหติ.\csromanspacer{} ภินฺนา โหนฺติ ปาปกา อกุสลา ธมฺมา สํกิเลสิกา โปโนภวิกา\footnote{โปโนพฺภวิกา (สฺยา, ก)} สทรา ทุกฺขวิปากา อายติํ ชาติชรามรณิยา.}

\prose{ปชฺเชน กเตน\footnote{ปชฺโชตกเตน (ก) ขุ 1. 357 ปิฏฺเฐ.} อตฺตนา, (สภิยาติ ภควา,)}

\prose{ปรินิพฺพานคโต วิติณฺณกงฺโข.}

\csromangathameasure{วิภวญฺจ ภวญฺจ วิปฺปหาย, \\ วุสิตวา ขีณปุนพฺภโว ส ภิกฺขูติ.}
\csromangathagroup{\csromangathastanza{วิภวญฺจ ภวญฺจ วิปฺปหาย, \\ วุสิตวา ขีณปุนพฺภโว ส ภิกฺขูติ.}}

\prose{\csromanfolio{40}\textbf{สโต ภิกฺขุ ปริพฺพเช}ติ สโต ภิกฺขุ ปริพฺพเช, สโต คจฺเฉยฺย.\csromanspacer{} สโต ติฏฺเฐยฺย.\csromanspacer{} สโต นิสีเทยฺย.\csromanspacer{} สโต เสยฺยํ กปฺเปยฺย.\csromanspacer{} สโต อภิกฺกเมยฺย.\csromanspacer{} สโต ปฏิกฺกเมยฺย.\csromanspacer{} สโต อาโลเกยฺย.\csromanspacer{} สโต วิโลเกยฺย.\csromanspacer{} สโต สมิญฺเชยฺย.\csromanspacer{} สโต ปสาเรยฺย.\csromanspacer{} สโต สงฺฆาฏิปตฺตจีวรํ ธาเรยฺย.\csromanspacer{} สโต จเรยฺย วิหเรยฺย อิริเยยฺย วตฺเตยฺย ปาเลยฺย ยเปยฺย ยาเปยฺยาติ สโต ภิกฺขุ ปริพฺพเช.\csromanspacer{} เตนาห ภควา\csromandash{}}

\csromangathameasure{กาเมสุ นาภิคิชฺเฌยฺย, มนสา’นาวิโล สิยา. \\ กุสโล สพฺพธมฺมานํ, สโต ภิกฺขุ ปริพฺพเชติ.}
\csromangathagroup{\csromangathastanza{กาเมสุ นาภิคิชฺเฌยฺย, มนสา’นาวิโล สิยา. \\ กุสโล สพฺพธมฺมานํ, สโต ภิกฺขุ ปริพฺพเชติ.}}

\prose{สห คาถาปริโยสานา เย เต พฺราหฺมเณน สทฺธิํ เอกจฺฉนฺทา เอกปโยคา เอกาธิปฺปายา เอกวาสนวาสิตา.\csromanspacer{} เตสํ อเนกปาณสหสฺสานํ วิรชํ วีตมลํ ธมฺมจกฺขุํ อุทปาทิ “ยํ กิญฺจิ สมุทยธมฺมํ, สพฺพํ ตํ นิโรธธมฺมนฺ”ติ.\csromanspacer{} ตสฺส พฺราหฺมณสฺส อนุปาทาย อาสเวหิ จิตฺตํ วิมุจฺจิ.\csromanspacer{} สห อรหตฺตปฺปตฺตา อชินชฏาวากจีรติทณฺฑกมณฺฑลุเกสา จ มสฺสู จ อนฺตรหิตา, ภณฺฑุกาสายวตฺถวสโน สงฺฆาฏิปตฺตจีวรธโร อนฺวตฺถปฏิปตฺติยา ปญฺชลิโก ภควนฺตํ นมสฺสมาโน นิสินฺโน โหติ “สตฺถา เม ภนฺเต ภควา, สาวโกหมสฺมี”ติ.}

\vspace{\tipitakaparskip}
\csromancenter{อชิตมาณวปุจฺฉานิทฺเทโส ปฐโม.}
\csromanmatikaanchor{mtk.23}
\csromantocmarkref{subsection}{2. ติสฺสเมตฺเตยฺยมาณวปุจฺฉานิทฺเทส}{mtk.23}
\bookmark[dest={mtk.23},level=2]{2. ติสฺสเมตฺเตยฺยมาณวปุจฺฉานิทฺเทส}
\csromanheader{2}{2. ติสฺสเมตฺเตยฺยมาณวปุจฺฉานิทฺเทส}
\proseitem{9}{\textbf{โกธ สนฺตุสิโต โลเก}, (อิจฺจายสฺมา ติสฺสเมตฺเตยฺโย,)}

\prose{\textbf{กสฺส โน สนฺติ อิญฺชิตา.}}

\csromangathameasure{\textbf{โก อุภนฺตมภิญฺญาย,} \\ \textbf{มชฺเฌ มนฺตา น ลิปฺปติ.} \\ \textbf{กํ พฺรูสิ มหาปุริโสติ,} \\ โก อิธ สิพฺพินิมจฺจคา.}
\csromangathagroup{\csromangathastanza{\textbf{โก อุภนฺตมภิญฺญาย,} \\ \textbf{มชฺเฌ มนฺตา น ลิปฺปติ.} \\ \textbf{กํ พฺรูสิ มหาปุริโสติ,} \\ โก อิธ สิพฺพินิมจฺจคา.}}

\prose{\textbf{โกธ สนฺตุสิโต โลเก}ติ โก โลเก ตุฏฺโฐ สนฺตุฏฺโฐ อตฺตมโน ปริปุณฺณสงฺกปฺโปติ โกธ สนฺตุสิโต โลเก.}

\csromanheader{2}{2. ติสฺสเมตฺเตยฺยมาณวปุจฺฉานิทฺเทส}
\prose{\csromanfolio{41}\textbf{อิจฺจายสฺมา ติสฺสเมตฺเตยฺโย}ติ \textbf{อิจฺจา}ติ ปทสนฺธิ ปทสํสคฺโค ปทปาริปูรี อกฺขรสมวาโย พฺยญฺชนสิลิฏฺฐตา ปทานุปุพฺพตาเปตํ \textbf{อิจฺจา}ติ.\csromanspacer{} \textbf{อายสฺมา}ติ ปิยจนํ ครุวจนํ สคารวสปฺปติสฺสาธิวจนเมตํ อายสฺมาติ.\csromanspacer{} \textbf{ติสฺโส}ติ ตสฺส พฺราหฺมณสฺส นามํ สงฺขา สมญฺญา ปญฺญตฺติ โวหาโร นามํ นามกมฺมํ นามเธยฺยํ นิรุตฺติ พฺยญฺชนํ อภิลาโป.\csromanspacer{} \textbf{เมตฺเตยฺโย}ติ ตสฺส พฺราหฺมณสฺส โคตฺตํ สงฺขา สมญฺญา ปญฺญตฺติ โวหาโรติ \textbf{อิจฺจา}ยสฺมา ติสฺสเมตฺเตยฺโย.}

\prose{\textbf{กสฺส โน สนฺติ อิญฺชิตา}ติ ตณฺหิญฺชิตํ ทิฏฺฐิญฺชิตํ มานิญฺชิตํ กิเลสิญฺชิตํ กามิญฺชิตํ.\csromanspacer{} กสฺสิเม อิญฺชิตา นตฺถิ น สนฺติ น สํวิชฺชนฺติ นุปลพฺภนฺติ, ปหีนา สมุจฺฉินฺนา วูปสนฺตา ปฏิปสฺสทฺธา อภพฺพุปฺปตฺติกา ญาณคฺคินา ทฑฺฒาติ กสฺส โน สนฺติ อิญฺชิตา.}

\prose{\textbf{โก อุภนฺตมภิญฺญายา}ติ โก อุโภ อนฺเต อภิญฺญาย ชานิตฺวา ตุลยิตฺวา ตีรยิตฺวา วิภาวยิตฺวา วิภูตํ กตฺวาติ โก อุภนฺตมภิญฺญาย.}

\prose{\textbf{มชฺเฌ มนฺตา น ลิปฺปตี}ติ มชฺเฌ มนฺตาย น ลิปฺปติ, อลิตฺโต อนุปลิตฺโต นิกฺขนฺโต นิสฺสโฏ วิปฺปมุตฺโต วิสญฺญุตฺโต วิมริยาทิกเตน เจตสา วิหรตีติ มชฺเฌ มนฺตา น ลิปฺปติ.}

\prose{\textbf{กํ พฺรูสิ มหาปุริโส}ติ มหาปุริโส อคฺคปุริโส เสฏฺฐปุริโส วิเสฏฺฐปุริโส ปาโมกฺขปุริโส อุตฺตมปุริโส ปธานปุริโส ปวรปุริโสติ กํ พฺรูสิ กํ กเถสิ กํ มญฺญสิ กํ ภณสิ กํ ปสฺสสิ กํ โวหรสีติ กํ พฺรูสิ มหาปุริโสติ.}

\prose{\textbf{โก อิธ สิพฺพินิมจฺจคา}ติ โก อิธ สิพฺพินิํ ตณฺหํ อชฺฌคา อุปจฺจคา อติกฺกนฺโต สมติกฺกนฺโต วีติวตฺโตติ โก อิธ สิพฺพินิมจฺจคา.\csromanspacer{} เตนาห โส พฺราหฺมโณ\csromandash{}}

\vspace{\tipitakaparskip}
\csromancenter{โกธ สนฺตุสิโต โลเก, (\textbf{อิจฺจา}ยสฺมา ติสฺสเมตฺเตยฺโย,)}
\prose{\textbf{กสฺส โน สนฺติ อิญฺชิตา}.}

\csromangathameasure{โก อุภนฺตมภิญฺญาย, \\ มชฺเฌ มนฺตา น ลิปฺปติ. \\ \textbf{กํ พฺรูสิ มหาปุริโส}ติ, \\ \textbf{โก อิธ สิพฺพินิมจฺจคา}ติ.}
\csromangathagroup{\csromangathastanza{โก อุภนฺตมภิญฺญาย, \\ มชฺเฌ มนฺตา น ลิปฺปติ. \\ \textbf{กํ พฺรูสิ มหาปุริโส}ติ, \\ \textbf{โก อิธ สิพฺพินิมจฺจคา}ติ.}}

\csromanhangingbegin
\hangingheaditem{10}{\csromanfolio{42}\textbf{กาเมสุ พฺรหฺมจริยวา}, (\textbf{เมตฺเตยฺยา}ติ \textbf{ภควา},)}
\hangingbody{\textbf{วีตตโณฺห สทา สโต.}}
\hangingbody{\textbf{สงฺขาย นิพฺพุโต ภิกฺขุ,}}
\hangingbody{ตสฺส โน สนฺติ อิญฺชิตา.}
\csromanhangingend

\prose{\textbf{กาเมสุ พฺรหฺมจริยวา}ติ \textbf{กามา}ติ อุทฺทานโต เทฺว \textbf{กามา} วตฺถุ\textbf{กามา} จ กิเลส\textbf{กามา} จ ฯเปฯ.\csromanspacer{} อิเม วุจฺจนฺติ วตฺถุ\textbf{กามา} ฯเปฯ.\csromanspacer{} อิเม วุจฺจนฺติ กิเลส\textbf{กามา}.\csromanspacer{} \textbf{พฺรหฺมจริยํ} วุจฺจติ อสทฺธมฺมสมาปตฺติยา อารติ วิรติ ปฏิวิรติ เวรมณี อกิริยา อกรณํ อนชฺฌาปตฺติ เวลา-อนติกฺกโม.\csromanspacer{} อปิ จ นิปฺปริยาเยน พฺรหฺมจริยํ วุจฺจติ อริโย อฏฺฐงฺคิโก มคฺโค.\csromanspacer{} เสยฺยถิทํ, สมฺมาทิฏฺฐิ สมฺมาสงฺกปฺโป สมฺมาวาจา สมฺมากมฺมนฺโต สมฺมา-อาชีโว สมฺมาวายาโม สมฺมาสติ สมฺมาสมาธิ.\csromanspacer{} โย อิมินา อริเยน อฏฺฐงฺคิเกน มคฺเคน อุเปโต สมุเปโต อุปาคโต สมุปาคโต อุปปนฺโน สมุปปนฺโน สมนฺนาคโต.\csromanspacer{} โส วุจฺจติ พฺรหฺมจริยวา, ยถา จ ธเนน ธนวาติ วุจฺจติ.\csromanspacer{} โภเคน โภควาติ วุจฺจติ.\csromanspacer{} ยเสน ยสวาติ วุจฺจติ.\csromanspacer{} สิปฺเปน สิปฺปวาติ วุจฺจติ.\csromanspacer{} สีเลน สีลวาติ วุจฺจติ.\csromanspacer{} วีริเยน วีริยวาติ วุจฺจติ.\csromanspacer{} ปญฺญาย ปญฺญวาติ วุจฺจติ.\csromanspacer{} วิชฺชาย วิชฺชวาติ วุจฺจติ.\csromanspacer{} เอวเมว โย อิมินา อริเยน อฏฺฐงฺคิเกน มคฺเคน อุเปโต สมุเปโต อุปาคโต สมุปาคโต อุปปนฺโน สมุปปนฺโน สมนฺนาคโต.\csromanspacer{} โส วุจฺจติ พฺรหฺมจริยวาติ กาเมสุ พฺรหฺมจริยวา.}

\prose{\textbf{เมตฺเตยฺยา}ติ \textbf{ภควา} ตํ พฺราหฺมณํ โคตฺเตน อาลปติ.\csromanspacer{} \textbf{ภควา}ติ คารวาธิวจนเมตํ ฯเปฯ สจฺฉิกา ปญฺญตฺติ, ยทิทํ \textbf{ภควา}ติ \textbf{เมตฺเตยฺยา}ติ \textbf{ภควา}.}

\prose{\textbf{วีตตโณฺห สทา สโต}ติ \textbf{ตณฺหา}ติ รูป\textbf{ตณฺหา} ฯเปฯ ธมฺม\textbf{ตณฺหา}, ยสฺเสสา \textbf{ตณฺหา} ปหีนา สมุจฺฉินฺนา วูปสนฺตา ปฏิปสฺสทฺธา อภพฺพุปฺปตฺติกา ญาณคฺคินา ทฑฺฒา.\csromanspacer{} โส วุจฺจติ วีตตโณฺห จตฺตตโณฺห วนฺตตโณฺห มุตฺตตโณฺห ปหีนตโณฺห ปฏินิสฺสฏฺฐตโณฺห วีตราโค จตฺตราโค วนฺตราโค มุตฺตราโค ปหีนราโค ปฏินิสฺสฏฺฐราโค นิจฺฉาโต นิพฺพุโต สีติภูโต สุขปฺปฏิสํเวที พฺรหฺมภูเตน อตฺตนา วิหรติ.\csromanspacer{} \textbf{สทา}ติ สทา สพฺพทา สพฺพกาลํ นิจฺจกาลํ ธุวกาลํ สตตํ สมิตํ อพฺโพกิณฺณํ}

\vspace{\tipitakaparskip}
\csromancenter{ติสฺสเมตฺเตยฺยมาณวปุจฺฉานิทฺเทส โปงฺขานุโปงฺขํ\footnote{โปขานุโปขํ (สฺยา)} อุทกูมิกชาตํ อวีจิ สนฺตติ สหิตํ\footnote{อวีจิ สมงฺคิสหิตํ (สฺยา) ขุ 7. 14 ปิฏฺเฐปิ.} ผสฺสิตํ\footnote{ผุสิตํ (สฺยา)} ปุเรภตฺตํ ปจฺฉาภตฺตํ ปุริมยามํ มชฺฌิมยามํ ปจฺฉิมยามํ กาเฬ ชุเณฺห วสฺเส เหมนฺเต คิเมฺห ปุริเม วโยขนฺเธ มชฺฌิเม วโยขนฺเธ ปจฺฉิเม วโยขนฺเธ.\csromanspacer{} \textbf{สโต}ติ จตูหิ การเณหิ สโต, กาเย กายานุปสฺสนาสติปฏฺฐานํ ภาเวนฺโต สโต, เวทนาสุ เวทนานุปสฺสนาสติปฏฺฐานํ ภาเวนฺโต สโต, จิตฺเต จิตฺตานุปสฺสนาสติปฏฺฐานํ ภาเวนฺโต สโต, ธมฺเมสุ ธมฺมานุปสฺสนาสติปฏฺฐานํ ภาเวนฺโต สโต ฯเปฯ.\csromanspacer{} โส วุจฺจติ สโตติ วีตตโณฺห สทา สโต.}
\prose{\csromanfolio{43}\textbf{สงฺขาย นิพฺพุโต ภิกฺขู}ติ สงฺขา วุจฺจติ ญาณํ.\csromanspacer{} ยา ปญฺญา ปชานนา วิจโย ปวิจโย ฯเปฯ อโมโห ธมฺมวิจโย สมฺมาทิฏฺฐิ.\csromanspacer{} \textbf{สงฺขายา}ติ สงฺขาย ชานิตฺวา ตุลยิตฺวา ตีรยิตฺวา วิภาวยิตฺวา วิภูตํ กตฺวา, “สพฺเพ สงฺขารา อนิจฺจา”ติ สงฺขาย ชานิตฺวา ตุ ลยิตฺวา ตีรยิตฺวา วิภาวยิตฺวา วิภูตํ กตฺวา, “สพฺเพ สงฺขารา ทุกฺขา”ติ. “สพฺเพ ธมฺมา อนตฺตา”ติ ฯเปฯ “อวิชฺชาปจฺจยา สงฺขารา”ติ. “ยํ กิญฺจิ สมุทยธมฺมํ, สพฺพํ ตํ นิโรธธมฺมนฺ”ติ สงฺขาย ชานิตฺวา ตุลยิตฺวา ตีรยิตฺวา วิภาวยิตฺวา วิภูตํ กตฺวา.}

\prose{อถ วา อนิจฺจโต สงฺขาย ชานิตฺวา.\csromanspacer{} ทุกฺขโต.\csromanspacer{} โรคโต.\csromanspacer{} คณฺฑโต.\csromanspacer{} สลฺลโต ฯเปฯ.\csromanspacer{} นิสฺสรณโต สงฺขาย ชานิตฺวา ตุลยิตฺวา ตีรยิตฺวา วิภาวยิตฺวา วิภูตํ กตฺวา.\csromanspacer{} \textbf{นิพฺพุโต}ติ ราคสฺส นิพฺพาปิตตฺตา นิพฺพุโต.\csromanspacer{} โทสสฺส นิพฺพาปิตตฺตา นิพฺพุโต.\csromanspacer{} โมหสฺส นิพฺพาปิตตฺตา นิพฺพุโต.\csromanspacer{} โกธสฺส.\csromanspacer{} อุปนาหสฺส.\csromanspacer{} มกฺขสฺส.\csromanspacer{} ปฬาสสฺส.\csromanspacer{} อิสฺสาย.\csromanspacer{} มจฺฉริยสฺส.\csromanspacer{} มายาย.\csromanspacer{} สาเฐยฺยสฺส.\csromanspacer{} ถมฺภสฺส.\csromanspacer{} สารมฺภสฺส.\csromanspacer{} มานสฺส.\csromanspacer{} อติมานสฺส.\csromanspacer{} มทสฺส.\csromanspacer{} ปมาทสฺส.\csromanspacer{} สพฺพกิเลสานํ.\csromanspacer{} สพฺพทุจฺจริตานํ.\csromanspacer{} สพฺพทรถานํ.\csromanspacer{} สพฺพปริฬาหานํ.\csromanspacer{} สพฺพสนฺตาปานํ.\csromanspacer{} สพฺพากุสลาภิสงฺขารานํ นิพฺพาปิตตฺตา นิพฺพุโต.\csromanspacer{} \textbf{ภิกฺขู}ติ สตฺตนฺนํ ธมฺมานํ ภินฺนตฺตา ภิกฺขุ ฯเปฯ วุสิตวา ขีณปุนพฺภโว ส ภิกฺขูติ สงฺขาย นิพฺพุโต ภิกฺขุ.}

\prose{\textbf{ตสฺส โน สนฺติ อิญฺชิตา}ติ \textbf{ตสฺสา}ติ อรหโต ขีณาสวสฺส.\csromanspacer{} \textbf{อิญฺชิตา}ติ ตณฺหิญฺชิตํ ทิฏฺฐิญฺชิตํ มานิญฺชิตฺตํ กิเลสิญฺชิตํ กามิญฺชิตํ.\csromanspacer{} ตสฺสิเม \csromanfolio{44}อิญฺชิตา นตฺถิ น สนฺติ น สํวิชฺชนฺติ นุปลพฺภนฺติ, ปหีนา สมุจฺฉินฺนา วูปสฺ\textbf{อนฺตา} ปฏิปสฺสทฺธา อภพฺพุปฺปตฺติกา ญาณคฺคินา ทฑฺฒาติ ตสฺส โน สนฺติ อิญฺชิตา.\csromanspacer{} เตนาห ภควา\csromandash{}}

\prose{กาเมสุ พฺรหฺมจริยวา, (เมตฺเตยฺยาติ ภควา,)}

\prose{วีตตโณฺห สทา สโต.}

\csromangathameasure{สงฺขาย นิพฺพุโต ภิกฺขุ, \\ ตสฺส โน สนฺติ อิญฺชิตาติ.}
\csromangathagroup{\csromangathastanza{สงฺขาย นิพฺพุโต ภิกฺขุ, \\ ตสฺส โน สนฺติ อิญฺชิตาติ.}}

\csromanhangingbegin
\hangingheaditem{11}{โส อุภนฺตมภิญฺญาย, มชฺเฌ มฺ\textbf{อนฺตา} น ลิปฺปติ.}
\hangingbody{ตํ พฺรูมิ มหาปุริโสติ, โส อิธ สิพฺพินิมจฺจคา.}
\csromanhangingend

\prose{\textbf{โส อุภนฺตมภิญฺญาย, มชฺเฌ มนฺตา น ลิปฺปตี}ติ \textbf{อนฺตา}ติ ผสฺโส เอโก อนฺโต, ผสฺสสมุทโย ทุติโย อนฺโต, ผสฺสนิโรโธ มชฺเฌ.\csromanspacer{} อตีตํ เอโก อนฺโต, อนาคตํ ทุติโย อนฺโต, ปจฺจุปฺปนฺนํ มชฺเฌ.\csromanspacer{} สุขา เวทนา เอโก อนฺโต, ทุกฺขา เวทนา ทุติโย อนฺโต, อทุกฺขมสุขา เวทนา มชฺเฌ.\csromanspacer{} นามํ เอโก อนฺโต, รูปํ ทุติโย อนฺโต, วิญฺญาณํ มชฺเฌ.\csromanspacer{} ฉ อชฺฌตฺติกานิ อายตนานิ เอโก อนฺโต, ฉ พาหิรานิ อายตนานิ ทุติโย อนฺโต, วิญฺญาณํ มชฺเฌ.\csromanspacer{} สกฺกาโย เอโก อนฺโต, สกฺกายสมุทโย ทุติโย อนฺโต, สกฺกายนิโรโธ มชฺเฌ.\csromanspacer{} \textbf{มนฺตา} วุจฺจติ ปญฺญา, ยา ปญฺญา ปชานนา ฯเปฯ อโมโห ธมฺมวิจโย สมฺมาทิฏฺฐิ.}

\prose{\textbf{เลปา}ติ เทฺว เลปา ตณฺหาเลโป จ ทิฏฺฐิเลโป จ.\csromanspacer{} กตโม ตณฺหาเลโป.\csromanspacer{} ยาวตา ตณฺหาสงฺขาเตน สีมกตํ โอธิกตํ\footnote{มริยาทิกตํ โอธิกตํ (สฺยา)} ปริยนฺตกตํ ปริคฺคหิตํ มมายิตํ อิทํ มม, เอตํ มม, เอตฺตกํ มม, เอตฺตาวตา มม รูปา สทฺทา คนฺธา รสา โผฏฺฐพฺพา อตฺถรณา ปาวุรณา ทาสิทาสา อเชฬกา กุกฺกุฏสูกรา หตฺถิควาสฺสวฬวา เขตฺตํ วตฺถุ หิรญฺญํ สุวณฺณํ คามนิคมราชธานิโย รฏฺฐญฺจ ชนปโท จ โกโส จ โกฏฺฐาคารญฺจ, เกวลมฺปิ มหาปถวิํ ตณฺหาวเสน มมายติ, ยาวตา อฏฺฐสตตณฺหาวิจริตํ.\csromanspacer{} อยํ ตณฺหาเลโป.}

\csromanheader{2}{2. ติสฺสเมตฺเตยฺยมาณวปุจฺฉานิทฺเทส}
\prose{\csromanfolio{45}กตโม ทิฏฺฐิเลโป.\csromanspacer{} วีสติวตฺถุกา สกฺกายทิฏฺฐิ, ทสวตฺถุกา มิจฺฉาทิฏฺฐิ, ทสวตฺถุกา อนฺตคฺคาหิกา ทิฏฺฐิ.\csromanspacer{} ยา เอวรูปา ทิฏฺฐิ ทิฏฺฐิคตํ ทิฏฺฐิคหนํ ทิฏฺฐิกนฺตาโร ทิฏฺฐิวิสูกายิกํ ทิฏฺฐิวิปฺผนฺทิตํ ทิฏฺฐิสํโยชนํ คาโห ปฏิคฺคาโห อภินิเวโส ปรามาโส กุมฺมคฺโค มิจฺฉาปโถ มิจฺฉตฺตํ ติตฺถายตนํ วิปริยาสคฺคาโห\footnote{วิปริเยสคฺคาโห (พหูสุ)} วิปรีตคฺคาโห วิปลฺลาสคฺคาโห มิจฺฉาคาโห, อยาถาวกสฺมิํ ยาถาวกนฺติ คาโห, ยาวตา ทฺวาสฏฺฐิ ทิฏฺฐิคตานิ.\csromanspacer{} อยํ ทิฏฺฐิเลโป.}

\prose{\textbf{โส อุภนฺตมภิญฺญาย, มชฺเฌ มนฺตา น ลิปฺปตี}ติ โส อุโภ จ อนฺเต มชฺฌญฺจ มนฺตาย อภิญฺญาย ชานิตฺวา ตุลยิตฺวา ตีรยิตฺวา วิภาวยิตฺวา วิภูตํ กตฺวา น ลิปฺปติ น ปลิปฺปติ น อุปลิปฺปติ, อลิตฺโต อสํลิตฺโต อนุปลิตฺโต นิกฺขนฺโต นิสฺสโฏ วิปฺปมุตฺโต วิสญฺญุตฺโต วิมริยาทิกเตน เจตสา วิหรตีติ โส อุภนฺตมภิญฺญาย, มชฺเฌ มนฺตา น ลิปฺปติ.}

\prose{\textbf{ตํ พฺรูมิ มหาปุริโสติ} มหาปุริโส อคฺคปุริโส เสฏฺฐปุริโส วิเสฏฺฐปุริโส ปาโมกฺขปุริโส อุตฺตมปุริโส ปวรปุริโส ตํ พฺรูมิ ตํ กเถมิ ตํ ภณามิ ตํ ทีเปมิ ตํ โวหรามิ.}

\prose{อายสฺมา สาริปุตฺโต\footnote{ปสฺส สํ 3. 137 ปิฏฺเฐ.} ภควนฺตํ เอตทโวจ “มหาปุริโส มหาปุริโส”ติ ภนฺเต วุจฺจติ.\csromanspacer{} กิตฺตาวตา นุ โข ภนฺเต มหาปุริโส โหตีติ.\csromanspacer{} วิมุตฺตจิตฺตตฺตา ขฺวาหํ สาริปุตฺต “มหาปุริโส”ติ วทามิ, อวิมุตฺตจิตฺตตฺตา โน “มหาปุริโส”ติ วทามิ.}

\prose{กถญฺจ สาริปุตฺต วิมุตฺตจิตฺโต โหติ.\csromanspacer{} อิธ สาริปุตฺต ภิกฺขุ อชฺฌตฺตํ กาเย กายานุปสฺสี วิหรติ อาตาปี สมฺปชาโน สติมา วิเนยฺย โลเก อภิชฺฌาโทมนสฺสํ.\csromanspacer{} ตสฺส กาเย กายานุปสฺสิโน วิหรโต จิตฺตํ วิรชฺชติ วิมุจฺจติ อนุปาทาย อาสเวหิ.\csromanspacer{} เวทนาสุ.\csromanspacer{} จิตฺเต.\csromanspacer{} ธมฺเมสุ ธมฺมานุปสฺสี วิหรติ อาตาปี สมฺปชาโน สติมา วิเนยฺย โลเก อภิชฺฌาโทมนสฺสํ.\csromanspacer{} ตสฺส ธมฺเมสุ ธมฺมานุปสฺสิโน วิหรโต จิตฺตํ วิรชฺชติ วิมุจฺจติ อนุปาทาย อาสเวหิ.\csromanspacer{} เอวํ โข สาริปุตฺต ภิกฺขุ วิมุตฺตจิตฺโต โหติ, วิมุตฺตจิตฺตตฺตา ขฺวาหํ สาริปุตฺต}

\prose{\csromanfolio{46}“มหาปุริโส”ติ วทามิ.\csromanspacer{} อวิมุตฺตจิตฺตตฺตา โน “มหาปุริโส”ติ วทามีติ ตํ พฺรูมิ มหาปุริโสติ.}

\prose{\textbf{โส อิธ สิพฺพินิมจฺจคา}ติ สิพฺพินี วุจฺจติ ตณฺหา, โย ราโค สาราโค ฯเปฯ อภิชฺฌา โลโภ อกุสลมูลํ.\csromanspacer{} ยสฺเสสา สิพฺพินี ตณฺหา ปหีนา สมุจฺฉินฺนา วูปสนฺตา ปฏิปสฺสทฺธา อภพฺพุปฺปตฺติกา ญาณคฺคินา ทฑฺฒา.\csromanspacer{} โส สิพฺพินิํ ตณฺหํ อจฺจคา อุปจฺจคา อติกฺกนฺโต สมติกฺกนฺโต วีติวตฺโตติ โส อิธ สิพฺพินิมจฺจคา.\csromanspacer{} เตนาห ภควา\csromandash{}}

\csromangathameasure{โส อุภนฺตมภิญฺญาย, มชฺเฌ มนฺตา น ลิปฺปติ. \\ ตํ พฺรูมิ มหาปุริโสติ, โส อิธ สิพฺพินิมจฺจคาติ.}
\csromangathagroup{\csromangathastanza{โส อุภนฺตมภิญฺญาย, มชฺเฌ มนฺตา น ลิปฺปติ. \\ ตํ พฺรูมิ มหาปุริโสติ, โส อิธ สิพฺพินิมจฺจคาติ.}}

\prose{สห คาถาปริโยสานา เย เต พฺราหฺมเณน สทฺธิํ เอกจฺฉนฺทา เอกปโยคา เอกาธิปฺปายา เอกวาสนวาสิตา, เตสํ อเนกปาณสหสฺสานํ วิรชํ วีตมลํ ธมฺมจกฺขุํ อุทปาทิ “ยํ กิญฺจิ สมุทยธมฺมํ, สพฺพํ ตํ นิโรธธมฺมนฺ”ติ.\csromanspacer{} ตสฺส พฺราหฺมณสฺส อนุปาทาย อาสเวหิ จิตฺตํ วิมุจฺจิ.\csromanspacer{} สห อรหตฺตปฺปตฺตา อชินชฏาวากจีรติทณฺฑกมณฺฑลุเกสา จ มสฺสู จ อนฺตรหิตา.\csromanspacer{} ภณฺฑุกาสายวตฺถวสโน สงฺฆาฏิปตฺตจีวรธโร อนฺวตฺถปฏิปตฺติยา ปญฺชลิโก ภควนฺตํ นมสฺสมาโน นิสินฺโน โหติ “สตฺถา เม ภนฺเต ภควา, สาวโกหมสฺมี”ติ.}

\vspace{\tipitakaparskip}
\csromancenter{ติสฺสเมตฺเตยฺยมาณวปุจฺฉานิทฺเทโส ทุติโย.}
\csromanmatikaanchor{mtk.24}
\csromantocmarkref{subsection}{3. ปุณฺณกมาณวปุจฺฉานิทฺเทส}{mtk.24}
\bookmark[dest={mtk.24},level=2]{3. ปุณฺณกมาณวปุจฺฉานิทฺเทส}
\csromanheader{2}{3. ปุณฺณกมาณวปุจฺฉานิทฺเทส}
\proseitem{12}{อเนชํ มูลทสฺสาวิํ, (อิจฺจายสฺมา ปุณฺณโก,)}

\prose{\textbf{อตฺถิ ปเญฺหน อาคมํ.}}

\csromangathameasure{\textbf{กิํ นิสฺสิตา อิสโย มนุชา,} \\ \textbf{ยญฺญมกปฺปยิํสุ ปุถุ’ธ โลเก,} \\ ปุจฺฉามิ ตํ ภควา พฺรูหิ เมตํ.}
\csromangathagroup{\csromangathastanza{\textbf{กิํ นิสฺสิตา อิสโย มนุชา,} \\ \textbf{ยญฺญมกปฺปยิํสุ ปุถุ’ธ โลเก,} \\ ปุจฺฉามิ ตํ ภควา พฺรูหิ เมตํ.}}

\prose{\textbf{อเนชํ มูลทสฺสาวินฺ}ติ \textbf{เอชา} วุจฺจติ ตณฺหา, โย ราโค สาราโค ฯเปฯ อภิชฺฌา โลโภ อกุสลมูลํ.\csromanspacer{} สา \textbf{เอชา} ตณฺหา พุทฺธสฺส ภควโต}

\vspace{\tipitakaparskip}
\csromancenter{ปุณฺณกมาณวปุจฺฉานิทฺเทส ปหีนา อุจฺฉินฺนมูลา ตาลาวตฺถุกตา อนภาวํกตา อายติํ อนุปฺปาทธมฺมา.\csromanspacer{} ตสฺมา พุทฺโธ อเนโช, เอชาย ปหีนตฺตา อเนโช.\csromanspacer{} ภควา ลาเภปิ น อิญฺชติ, อลาเภปิ น อิญฺชติ, ยเสปิ น อิญฺชติ, อยเสปิ น อิญฺชติ, ปสํสายปิ น อิญฺชติ, นินฺทายปิ น อิญฺชติ, สุเขปิ น อิญฺชติ, ทุกฺเขปิ น อิญฺชติ น จลติ น เวธติ นปฺปเวธตีติ อเนชํ.\csromanspacer{}\csromanspacer{}\textbf{มูลทสฺสาวินฺ}ติ ภควา มูลทสฺสาวี เหตุทสฺสาวี นิทานทสฺสาวี สมฺภวทสฺสาวี ปภวทสฺสาวี สมุฏฺฐานทสฺสาวี อาหารทสฺสาวี อารมฺมณทสฺสาวี ปจฺจยทสฺสาวี สมุทยทสฺสาวี.}
\prose{\csromanfolio{47}ตีณิ อกุสลมูลานิ, โลโภ อกุสลมูลํ, โทโส อกุสลมูลํ, โมโห อกุสลมูลํ.}

\prose{วุตฺตํ เหตํ ภควตา\csromandash{}\footnote{ปสฺส อํ 1. 266 ปิฏฺเฐ.} ตีณิมานิ ภิกฺขเว นิทานานิ กมฺมานํ สมุทยาย.\csromanspacer{} กตมานิ ตีณิ, โลโภ นิทานํ กมฺมานํ สมุทยาย, โทโส นิทานํ กมฺมานํ สมุทยาย, โมโห นิทานํ กมฺมานํ สมุทยาย.\csromanspacer{} น ภิกฺขเว โลภเชน กมฺเมน โทสเชน กมฺเมน โมหเชน กมฺเมน เทวา ปญฺญายนฺติ, มนุสฺสา ปญฺญายนฺติ, ยา วา ปนญฺญาปิ กาจิ สุคติโย.\csromanspacer{} อถ โข ภิกฺขเว โลภเชน กมฺเมน โทสเชน กมฺเมน โมหเชน กมฺเมน นิรโย ปญฺญายติ, ติรจฺฉานโยนิ ปญฺญายติ, เปตฺติวิสโย ปญฺญายติ.\csromanspacer{} ยา วา ปนญฺญาปิ กาจิ ทุคฺคติโย นิรเย ติรจฺฉานโยนิยา เปตฺติวิสเย อตฺตภาวาภินิพฺพตฺติยา.\csromanspacer{} อิมานิ ตีณิ อกุสลมูลานีติ ภควา ชานาติ ปสฺสติ.\csromanspacer{} เอวมฺปิ ภควา มูลทสฺสาวี ฯเปฯ สมุทยทสฺสาวี.\csromanspacer{} ตีณิ กุสลมูลานิ, อโลโภ กุสลมูลํ, อโทโส กุสลมูลํ, อโมโห กุสลมูลํ.}

\prose{วุตฺตํ เหตํ ภควตา\csromandash{}ตีณิมานิ ฯเปฯ.\csromanspacer{} น ภิกฺขเว อโลภเชน กมฺเมน อโทสเชน กมฺเมน อโมหเชน กมฺเมน นิรโย ปญฺญายติ, ติรจฺฉานโยนิ ปญฺญายติ, เปตฺติวิสโย ปญฺญายติ, ยา วา ปนญฺญาปิ กาจิ ทุคฺคติโย.\csromanspacer{} อถ โข ภิกฺขเว อโลภเชน กมฺเมน อโทสเชน กมฺเมน อโมหเชน กมฺเมน เทวา ปญฺญายนฺติ, มนุสฺสา ปญฺญายนฺติ, ยา วา ปนญฺญาปิ กาจิ สุคติโย เทเว จ มนุสฺเส จ อตฺตภาวาภินิพฺพตฺติยา.\csromanspacer{} อิมานิ ตีณิ กุสลมูลานีติ ภควา ชานาติ ปสฺสติ.\csromanspacer{} เอวมฺปิ ภควา มูลทสฺสาวี ฯเปฯ สมุทยทสฺสาวี.}

\prose{\csromanfolio{48}วุตฺตํ เหตํ ภควตา\csromandash{}เย เกจิ ภิกฺขเว ธมฺมา อกุสลา อกุสลภาคิยา อกุสลปกฺขิกา, สพฺเพ เต อวิชฺชามูลกา อวิชฺชาสโมสรณา.\csromanspacer{} อวิชฺชา สมุคฺฆาตา สพฺเพ เต สมุคฺฆาตํ คจฺฉนฺตีติ ภควา ชานาติ ปสฺสติ.\csromanspacer{} เอวมฺปิ ภควา มูลทสฺสาวี \textbf{ฯเปฯ} สมุทยทสฺสาวี.}

\prose{วุตฺตํ เหตํ ภควตา\csromandash{}เย เกจิ ภิกฺขเว ธมฺมา กุสลา กุสลภาคิยา กุสลปกฺขิกา, สพฺเพ เต อปฺปมาทมูลกา อปฺปมาทสโมสรณา.\csromanspacer{} อปฺปมาโท เตสํ ธมฺมานํ อคฺคมกฺขายตีติ ภควา ชานาติ ปสฺสติ.\csromanspacer{} เอวมฺปิ ภควา มูลทสฺสาวี \textbf{ฯเปฯ} สมุทยทสฺสาวี.}

\prose{อถ วา ภควา ชานาติ ปสฺสติ “อวิชฺชา มูลํ สงฺขารานํ.\csromanspacer{} สงฺขารา มูลํ วิญฺญาณสฺส.\csromanspacer{} วิญฺญาณํ มูลํ นามรูปสฺส.\csromanspacer{} นามรูปํ มูลํ สฬายตนสฺส.\csromanspacer{} สฬายตนํ มูลํ ผสฺสสฺส.\csromanspacer{} ผสฺโส มูลํ เวทนาย.\csromanspacer{} เวทนา มูลํ ตณฺหาย.\csromanspacer{} ตณฺหา มูลํ อุปาทานสฺส.\csromanspacer{} อุปาทานํ มูลํ ภวสฺส.\csromanspacer{} ภโว มูลํ ชาติยา.\csromanspacer{} ชาติ มูลํ ชรามรณสฺสา”ติ ภควา ชานาติ ปสฺสติ.\csromanspacer{} เอวมฺปิ ภควา มูลทสฺสาวี \textbf{ฯเปฯ} สมุทยทสฺสาวี.}

\prose{อถ วา ภควา ชานาติ ปสฺสติ “จกฺขุ มูลํ จกฺขุโรคานํ.\csromanspacer{} โสตํ มูลํ โสตโรคานํ.\csromanspacer{} ฆานํ มูลํ ฆานโรคานํ.\csromanspacer{} ชิวฺหา มูลํ ชิวฺหาโรคานํ.\csromanspacer{} กาโย มูลํ กายโรคานํ.\csromanspacer{} มโน มูลํ เจตสิกานํ ทุกฺขานนฺ”ติ ภควา ชานาติ ปสฺสติ.\csromanspacer{} เอวมฺปิ ภควา มูลทสฺสาวี เหตุทสฺสาวี นิทานทสฺสาวี สมฺภวทสฺสาวี ปภวทสฺสาวี สมุฏฺฐานทสฺสาวี อาหารทสฺสาวี อารมฺมณทสฺสาวี ปจฺจยทสฺสาวี สมุทยทสฺสาวีติ อเนชํ มูลทสฺสาวี.}

\prose{\textbf{อิจฺจายสฺมา ปุณฺณโก}ติ \textbf{อิจฺจา}ติ ปทสนฺธิ \textbf{ฯเปฯ} อายสฺมา ปุณฺณโก.}

\prose{\textbf{อตฺถิ ปเญฺหน อาคมนฺ}ติ ปเญฺหน อตฺถิโก อาคโตมฺหิ\footnote{ปญฺหตฺถิกามฺห อาคตา (พหูสุ) ปสฺส ขุ 7. 369 ปิฏฺเฐ.}, ปญฺหํ ปุจฺฉิตุกาโม อาคโตมฺหิ, ปญฺหํ โสตุกาโม อาคโตมฺหีติ เอวมฺปิ อตฺถิ ปเญฺหน อาคมํ.\csromanspacer{}\csromanspacer{}อถ วา ปญฺหตฺถิกานํ ปญฺหํ ปุจฺฉิตุกามานํ ปญฺหํ โสตุกามานํ อาคมนํ อภิกฺกมนํ อุปสงฺกมนํ ปยิรุปาสนํ อตฺถีติ เอวมฺปิ อตฺถิ ปเญฺหน อาคมํ.\csromanspacer{}\csromanspacer{}อถ วา ปญฺหาคโม ตุยฺหํ อตฺถิ, ตฺวมฺปิ ปหุ, วิสวี อลมตฺโต, มยา ปุจฺฉิตํ กเถตุํ วิสชฺเชตุํ, วหสฺเสตํ ภารนฺติ\footnote{วิสชฺเชตุํ สนฺทสฺเสตุํ ภณิตุนฺติ (สฺยา) วหสฺสุ + เอตํ.} เอวมฺปิ อตฺถิ ปเญฺหน อาคมํ.}

\csromanheader{2}{3. ปุณฺณกมาณวปุจฺฉานิทฺเทส}
\prose{\csromanfolio{49}\textbf{กิํ นิสฺสิตา อิสโย มนุชา}ติ กิํ นิสฺสิตา อาสิตา อลฺลีนา อุปคตา อชฺโฌสิตา อธิมุตฺตา.\csromanspacer{} \textbf{อิสโย}ติ อิสินามกา เย เกจิ อิสิปพฺพชฺชํ ปพฺพชิตา อาชีวกา นิคณฺฐา ชฏิลา ตาปสา.\csromanspacer{} \textbf{มนุชา}ติ มนุสฺสา วุจฺจนฺตีติ กิํ นิสฺสิตา อิสโย มนุชา.}

\prose{\textbf{ขตฺติยา พฺราหฺมณา เทวตานนฺ}ติ \textbf{ขตฺติยา}ติ เย เกจิ ขตฺติยชาติกา.\csromanspacer{} \textbf{พฺราหฺมณา}ติ เย เกจิ โภวาทิกา.\csromanspacer{} \textbf{เทวตานนฺ}ติ อาชีวกสาวกานํ อาชีวกา เทวตา.\csromanspacer{} นิคณฺฐสาวกานํ นิคณฺฐา เทวตา.\csromanspacer{} ชฏิลสาวกานํ ชฏิลา เทวตา.\csromanspacer{} ปริพฺพาชกสาวกานํ ปริพฺพาชกา เทวตา.\csromanspacer{} อวิรุทฺธกสาวกานํ อวิรุทฺธกา\footnote{อวรุทฺธกสาวกานํ อวรุทฺธกา (สฺยา)} เทวตา.\csromanspacer{} หตฺถิวติกานํ หตฺถี เทวตา.\csromanspacer{} อสฺสวติกานํ อสฺสา เทวตา.\csromanspacer{} โควติกานํ คาโว เทวตา.\csromanspacer{} กุกฺกุรวติกานํ กุกฺกุรา เทวตา.\csromanspacer{} กากวติกานํ กากา เทวตา.\csromanspacer{} วาสุเทววติกานํ วาสุเทโว เทวตา.\csromanspacer{} พลเทววติกานํ พลเทโว เทวตา.\csromanspacer{} ปุณฺณภทฺทวติกานํ ปุณฺณภทฺโท เทวตา.\csromanspacer{} มณิภทฺทวติกานํ มณิภทฺโท เทวตา.\csromanspacer{} อคฺคิวติกานํ อคฺคิ เทวตา.\csromanspacer{} นาควติกานํ นาคา เทวตา.\csromanspacer{} สุปณฺณวติกานํ สุปณฺณา เทวตา.\csromanspacer{} ยกฺขวติกานํ ยกฺขา เทวตา.\csromanspacer{} อสุรวติกานํ อสุรา เทวตา.\csromanspacer{} คนฺธพฺพวติกานํ คนฺธพฺพา เทวตา.\csromanspacer{} มหาราชวติกานํ มหาราชาโน เทวตา.\csromanspacer{} จนฺทวติกานํ จนฺโท เทวตา.\csromanspacer{} สูริยวติกานํ สูริโย เทวตา.\csromanspacer{} อินฺทวติกานํ อินฺโท เทวตา.\csromanspacer{} พฺรหฺมวติกานํ พฺรหฺมา เทวตา.\csromanspacer{} เทววติกานํ เทโว เทวตา.\csromanspacer{} ทิสาวติกานํ ทิสา เทวตา.\csromanspacer{} เย เยสํ ทกฺขิเณยฺยา เต เตสํ เทวตาติ ขตฺติยพฺราหฺมณา เทวตานํ.}

\prose{\textbf{ยญฺญมกปฺปยิํสุ ปุถูธ โลเก}ติ \textbf{ยญฺญํ} วุจฺจติ เทยฺยธมฺโม จีวรปิณฺฑปาตเสนาสนคิลานปจฺจยเภสชฺชปริกฺขารํ อนฺนํ ปานํ วตฺถํ ยานํ มาลาคนฺธาวิเลปนํ\footnote{มาลาคนฺธํ วิเลปนํ (สฺยา) ขุ 1. 240 ปิฏฺเฐ.} เสยฺยาวสถปทีเปยฺยํ.\csromanspacer{} \textbf{ยญฺญมกปฺปยิํสู}ติ เยปิ \textbf{ยญฺญํ} เอสนฺติ คเวสนฺติ ปริเยสนฺติ จีวรปิณฺฑปาตเสนาสนคิลานปจฺจยเภสชฺชปริกฺขารํ อนฺนํ ปานํ วตฺถํ ยานํ มาลาคนฺธวิเลปนํ เสยฺยาวสถปทีเปยฺยํ.\csromanspacer{} เตปิ \textbf{ยญฺญํ} กปฺเปนฺติ.\csromanspacer{} เยปิ \textbf{ยญฺญํ} อภิสงฺขาโรนฺติ จีวรปิณฺฑปาตเสนาสนคิลานปจฺจยเภสชฺชปริกฺขารํ อนฺนํ ปานํ ฯเปฯ เสยฺยาวสถปทีเปยฺยํ.\csromanspacer{} เตปิ \textbf{ยญฺญํ} กปฺเปนฺติ.\csromanspacer{} เยปิ \textbf{ยญฺญํ} เทนฺติ ยชนฺติ ปริจฺจชนฺติ จีวร ปิณฺฑปาต เสนาสน คิลาน ปจฺจย เภสชฺช ปริกฺขารํ อนฺนํ ปานํ ฯเปฯ}

\prose{\csromanfolio{50}เสยฺยาวสถปทีเปยฺยํ.\csromanspacer{} เตปิ ยญฺญํ กปฺเปนฺติ.\csromanspacer{}\csromanspacer{}\textbf{ปุถู}ติ ยญฺญา วา เอเต ปุถู.\csromanspacer{} ยญฺญยาชกา\footnote{ยญฺญยชกา (สฺยา)} วา เอเต ปุถู.\csromanspacer{} ทกฺขิเณยฺยา วา เอเต ปุถู.\csromanspacer{} กถํ ยญฺญา วา เอเต ปุถู.\csromanspacer{} พหุกานํ เอเต ยญฺญา จีวรปิณฺฑปาตเสนาสนคิลานปจฺจยเภสชฺชปริกฺขารา อนฺนํ ปานํ วตฺถํ ยานํ มาลํ คนฺธํ วิเลปนํ เสยฺยาวสถปทีเปยฺยํ.\csromanspacer{} เอวํ ยญฺญา วา เอเต ปุถู.}

\prose{กถํ ยญฺญยาชกา วา เอเต ปุถู.\csromanspacer{} พหุกา เอเต ยญฺญยาชกา ขตฺติยา จ พฺราหฺมณา จ เวสฺสา จ สุทฺทา จ คหฏฺฐา จ ปพฺพชิตา จ เทวา จ มนุสฺสา จ.\csromanspacer{} เอวํ ยญฺญยาชกา วา เอเต ปุถู.}

\prose{กถํ ทกฺขิเณยฺยา วา เอเต ปุถู.\csromanspacer{} พหุกา เอเต ทกฺขิเณยฺยา ปุถู สมณพฺราหฺมณา กปณทฺธิกวนิพฺพกยาจกา\footnote{...วณิพฺพกสาวกา (สฺยา) ขุ 1. 240 ปิฏฺเฐ.}.\csromanspacer{} เอวํ ทกฺขิเณยฺยา วา เอเต ปุถู.\csromanspacer{}\csromanspacer{}\textbf{อิธ โลเก}ติ มนุสฺสโลเกติ ยญฺญมกปฺปยิํสุ ปุถูธ โลเก.}

\prose{\textbf{ปุจฺฉามิ ตํ ภควา พฺรูหิ เมตนฺ}ติ \textbf{ปุจฺฉา}ติ ติสฺโส \textbf{ปุจฺฉา} อทิฏฺฐโชตนา \textbf{ปุจฺฉา} ทิฏฺฐสํสนฺทนา \textbf{ปุจฺฉา} วิมติจฺเฉทนา \textbf{ปุจฺฉา}.\csromanspacer{} กตมา อทิฏฺฐโชตนา \textbf{ปุจฺฉา}.\csromanspacer{} ปกติยา ลกฺขณํ อญฺญาตํ โหติ อทิฏฺฐํ อตุลิตํ อตีริตํ อวิภูตํ อวิภาวิตํ, ตสฺส ญาณาย ทสฺสนาย ตุลนาย ตีรณาย วิภูตตฺถาย วิภาวนตฺถาย ปญฺหํ ปุจฺฉติ.\csromanspacer{} อยํ อทิฏฺฐโชตนา \textbf{ปุจฺฉา}.}

\prose{กตมา ทิฏฺฐสํสนฺทนา \textbf{ปุจฺฉา}.\csromanspacer{} ปกติยา ลกฺขณํ ญาตํ โหติ ทิฏฺฐํ ตุลิตํ ตีริตํ วิภูตํ วิภาวิตํ.\csromanspacer{} อญฺเญหิ ปณฺฑิเตหิ สทฺธิํ สํสนฺทนตฺถาย ปญฺหํ ปุจฺฉติ.\csromanspacer{} อยํ ทิฏฺฐสํสนฺทนา \textbf{ปุจฺฉา}.}

\prose{กตมา วิมติจฺเฉทนา \textbf{ปุจฺฉา}.\csromanspacer{} ปกติยา สํสยปกฺขนฺโท\footnote{สํสยปกฺขนฺโน (สฺยา)} โหติ วิมติปกฺขนฺโท เทฺวฬฺหกชาโต “เอวํ นุ โข, น นุ โข, กิํ นุ โข, กถํ นุ โข”ติ.\csromanspacer{} โส วิมติจฺเฉทนตฺถาย ปญฺหํ ปุจฺฉติ.\csromanspacer{} อยํ วิมติจฺเฉทนา \textbf{ปุจฺฉา}.\csromanspacer{} อิมา ติสฺโส \textbf{ปุจฺฉา}.}

\prose{อปราปิ ติสฺโส \textbf{ปุจฺฉา} มนุสฺส\textbf{ปุจฺฉา} อมนุสฺส\textbf{ปุจฺฉา} นิมฺมิต\textbf{ปุจฺฉา}.\csromanspacer{} กตมา มนุสฺส\textbf{ปุจฺฉา}.\csromanspacer{} มนุสฺสา พุทฺธํ ภควนฺตํ อุปสงฺกมิตฺวา ปุจฺฉนฺติ, ภิกฺขู ปุจฺฉนฺติ, ภิกฺขุนิโย ปุจฺฉนฺติ, อุปาสกา ปุจฺฉนฺติ, อุปาสิกาโย ปุจฺฉนฺติ,}

\vspace{\tipitakaparskip}
\csromancenter{ปุณฺณกมาณวปุจฺฉานิทฺเทส ราชาโน ปุจฺฉนฺติ, ขตฺติยา ปุจฺฉนฺติ, พฺราหฺมณา ปุจฺฉนฺติ, เวสฺสา ปุจฺฉนฺติ, สุทฺทา ปุจฺฉนฺติ, คหฏฺฐา ปุจฺฉนฺติ, ปพฺพชิตา ปุจฺฉนฺติ, อยํ มนุสฺสปุจฺฉา.}
\prose{\csromanfolio{51}กตมา อมนุสฺสปุจฺฉา.\csromanspacer{} อมนุสฺสา พุทฺธํ ภควนฺตํ อุปสงฺกมิตฺวา ปญฺหํ ปุจฺฉนฺติ, นาคา ปุจฺฉนฺติ, สุปณฺณา ปุจฺฉนฺติ, ยกฺขา ปุจฺฉนฺติ, อสุรา ปุจฺฉนฺติ, คนฺธพฺพา ปุจฺฉนฺติ, มหาราชาโน ปุจฺฉนฺติ, อินฺทา ปุจฺฉนฺติ, พฺรหฺมาโน ปุจฺฉนฺติ, เทวตาโย ปุจฺฉนฺติ.\csromanspacer{} อยํ อมนุสฺสปุจฺฉา.}

\prose{กตมา นิมฺมิตปุจฺฉา.\csromanspacer{} ยํ \textbf{ภควา} รูปํ อภินิมฺมินาติ มโนมยํ สพฺพงฺคปจฺจงฺคํ อหีนินฺทฺริยํ.\csromanspacer{} โส นิมฺมิโต พุทฺธํ ภควนฺตํ อุปสงฺกมิตฺวา ปญฺหํ ปุจฺฉติ.\csromanspacer{} \textbf{ภควา} วิสชฺเชติ\footnote{วิสฺสชฺเชติ (ก)}.\csromanspacer{} อยํ นิมฺมิตปุจฺฉา.\csromanspacer{} อิมา ติสฺโส ปุจฺฉา.}

\prose{อปราปิ ติสฺโส ปุจฺฉา อตฺตตฺถปุจฺฉา ปรตฺถปุจฺฉา อุภยตฺถปุจฺฉา.\csromanspacer{} อปราปิ ติสฺโส ปุจฺฉา ทิฏฺฐธมฺมิกตฺถปุจฺฉา สมฺปรายิกตฺถปุจฺฉา ปรมตฺถปุจฺฉา.\csromanspacer{} อปราปิ ติสฺโส ปุจฺฉา อนวชฺชตฺถปุจฺฉา นิกฺกิเลสตฺถปุจฺฉา โวทานตฺถปุจฺฉา.\csromanspacer{} อปราปิ ติสฺโส ปุจฺฉา อตีตปุจฺฉา อนาคตปุจฺฉา ปจฺจุปฺปนฺนปุจฺฉา.\csromanspacer{} อปราปิ ติสฺโส ปุจฺฉา อชฺฌตฺตปุจฺฉา พหิทฺธาปุจฺฉา อชฺฌตฺตพหิทฺธาปุจฺฉา.\csromanspacer{} อปราปิ ติสฺโส ปุจฺฉา กุสลปุจฺฉา อกุสลปุจฺฉา อพฺยากตปุจฺฉา.\csromanspacer{} อปราปิ ติสฺโส ปุจฺฉา ขนฺธปุจฺฉา ธาตุปุจฺฉา อายตนปุจฺฉา.\csromanspacer{} อปราปิ ติสฺโส ปุจฺฉา สติปฏฺฐานปุจฺฉา สมฺมปฺปธานปุจฺฉา อิทฺธิปาทปุจฺฉา.\csromanspacer{} อปราปิ ติสฺโส ปุจฺฉา อินฺทฺริยปุจฺฉา พลปุจฺฉา โพชฺฌงฺคปุจฺฉา.\csromanspacer{} อปราปิ ติสฺโส ปุจฺฉา มคฺคปุจฺฉา ผลปุจฺฉา นิพฺพานปุจฺฉา.\csromanspacer{} \textbf{ปุจฺฉามิ ต}นฺติ ปุจฺฉามิ ตํ ยาจามิ ตํ อชฺเฌสามิ ตํ ปสาเทมิ ตํ “กถยสฺสุ เม”ติ ปุจฺฉามิ ตํ.\csromanspacer{} \textbf{ภควา}ติ คารวาธิวจนเมตํ ฯเปฯ สจฺฉิกา ปญฺญตฺติ ยทิทํ \textbf{ภควา}ติ.\csromanspacer{} \textbf{พฺรูหิเมตนฺ}ติ พฺรูหิ อาจิกฺขาหิ เทเสหิ ปญฺญเปหิ ปฏฺฐเปหิ วิวราหิ วิภชาหิ อุตฺตานีกโรหิ ปกาเสหีติ ปุจฺฉามิ ตํ \textbf{ภควา} พฺรูหิ เมตํ.\csromanspacer{} เตนาห โส พฺราหฺมโณ\csromandash{}}

\vspace{\tipitakaparskip}
\csromancenter{อเนชํ มูลทสฺสาวิํ, (อิจฺจายสฺมา ปุณฺณโก,)}
\prose{อตฺถิ ปเญฺหน อาคมํ.}

\csromangathameasure{กิํ นิสฺสิตา อิสโย มนุชา, \\ ยญฺญมกปฺปยิํสุ ปุถูธ โลเก, \\ \textbf{ปุจฺฉามิ ต}ํ \textbf{ภควา} พฺรูหิ เมตนฺติ.}
\csromangathagroup{\csromangathastanza{กิํ นิสฺสิตา อิสโย มนุชา, \\ ยญฺญมกปฺปยิํสุ ปุถูธ โลเก, \\ \textbf{ปุจฺฉามิ ต}ํ \textbf{ภควา} พฺรูหิ เมตนฺติ.}}

\proseitem{13}{\csromanfolio{52}\textbf{เย เกจิเม อิสโย มนุชา}, (ปุณฺณกาติ ภควา,)}

\csromangathameasure{\textbf{ยญฺญมกปฺปยิํสุ ปุถูธ โลเก,} \\ \textbf{อาสีสมานา ปุณฺณก อิตฺถตฺตํ}\textbf{.} \\ ชรํ สิตา ยญฺญมกปฺปยิํสุ\textbf{.}}
\csromangathagroup{\csromangathastanza{\textbf{ยญฺญมกปฺปยิํสุ ปุถูธ โลเก,} \\ \textbf{อาสีสมานา ปุณฺณก อิตฺถตฺตํ}\footnote{อิตฺถตํ (สฺยา), อิตฺถภาวํ (ก)}\textbf{.} \\ ชรํ สิตา ยญฺญมกปฺปยิํสุ\textbf{.}}}

\prose{\textbf{เย เกจิเม อิสโย มนุชา}ติ \textbf{เย เกจี}ติ สพฺเพน สพฺพํ สพฺพถา สพฺพํ อเสสํ นิสฺเสสํ ปริยาทิยนวจนเมตํ \textbf{เย เกจี}ติ\textbf{.}\csromanspacer{} \textbf{อิสโย}ติ อิสินามกา เย เกจิ อิสิปพฺพชฺชํ ปพฺพชิตา อาชีวกา นิคณฺฐา ชฏิลา ตาปสา\textbf{.}\csromanspacer{} \textbf{มนุชา}ติ มนุสฺสา วุจฺจนฺตีติ เย เกจิเม อิสโย มนุชา, ปุณฺณกาติ ภควา\textbf{.}}

\prose{\textbf{ขตฺติยา พฺราหฺมณา เทวตานนฺ}ติ \textbf{ขตฺติยา}ติ เย เกจิ ขตฺติยชาติกา\textbf{.}\csromanspacer{} \textbf{พฺราหฺมณา}ติ เย เกจิ โภวาทิกา\textbf{.}\csromanspacer{} \textbf{เทวตานนฺ}ติ อาชีวกสาวกานํ อาชีวกา เทวตา ฯเปฯ\textbf{.}\csromanspacer{} ทิสาวติกานํ ทิสา เทวตา\textbf{.}\csromanspacer{} เย เยสํ ทกฺขิเณยฺยา, เต เตสํ เทวตาติ \textbf{ขตฺติยา} พฺราหฺมณา เทวตานํ\textbf{.}}

\prose{\textbf{ยญฺญมกปฺปยิํสุ ปุถูธ โลเก}ติ \textbf{ยญฺญํ} วุจฺจติ เทยฺยธมฺโม, จีวรปิณฺฑปาตเสนาสนคิลานปจฺจยเภสชฺชปริกฺขารํ อนฺนํ ปานํ ฯเปฯ เสยฺยาวสถปทีเปยฺยํ\textbf{.}\csromanspacer{} \textbf{ยญฺญมกปฺปยิํสู}ติ เยปิ \textbf{ยญฺญํ} เอสนฺติ คเวสนฺติ ปริเยสนฺติ ฯเปฯ เสยฺยาวสถปทีเปยฺยํ, เตปิ \textbf{ยญฺญํ} กปฺเปนฺติ\textbf{.}\csromanspacer{} \textbf{ปุถู}ติ ยญฺญา วา เอเต ปุถู, ยญฺญยาชกา วา เอเต ปุถู, ทกฺขิเณยฺยา วา เอเต ปุถู ฯเปฯ\textbf{.}\csromanspacer{} เอวํ ทกฺขิเณยฺยา วา เอเต ปุถู\textbf{.}\csromanspacer{} \textbf{อิธ โลเก}ติ มนุสฺสโลเกติ ยญฺญมกปฺปยิํสุ ปุถูธ โลเก\textbf{.}}

\prose{\textbf{อาสีสมานา ปุณฺณก อิตฺถตฺตนฺ}ติ \textbf{อาสีสมานา}ติ รูปปฏิลาภํ \textbf{อาสีสมานา} สทฺทปฏิลาภํ \textbf{อาสีสมานา} คนฺธปฏิลาภํ \textbf{อาสีสมานา} รสปฏิลาภํ \textbf{อาสีสมานา} โผฏฺฐพฺพปฏิลาภํ \textbf{อาสีสมานา} ปุตฺตปฏิลาภํ \textbf{อาสีสมานา} ทารปฏิลาภํ \textbf{อาสีสมานา} ธนปฏิลาภํ \textbf{อาสีสมานา} ยสปฏิลาภํ \textbf{อาสีสมานา} อิสฺสริยปฏิลาภํ \textbf{อาสีสมานา} ขตฺติยมหาสาลกุเล อตฺตภาวปฏิลาภํ \textbf{อาสีสมานา}}

\vspace{\tipitakaparskip}
\csromancenter{\textbf{ปุณฺณก}มาณวปุจฺฉานิทฺเทส พฺราหฺมณมหาสาลกุเล อตฺตภาวปฏิลาภํ อาสีสมานา คหปติมหาสาลกุเล อตฺตภาวปฏิลาภํ อาสีสมานา จาตุมหาราชิเกสุ\footnote{จาตุมฺมหาราชิเกสุ (สฺยา)} เทเวสุ อตฺตภาวปฏิลาภํ อาสีสมานา ตาวติํเสสุ เทเวสุ.\csromanspacer{} ยาเมสุ เทเวสุ.\csromanspacer{} ตุสิเตสุ เทเวสุ.\csromanspacer{} นิมฺมานรตีสุ เทเวสุ.\csromanspacer{} ปรนิมฺมิตวสวตฺตีสุ เทเวสุ.\csromanspacer{} พฺรหฺมกายิเกสุ เทเวสุ อตฺตภาวปฏิลาภํ อาสีสมานา อิจฺฉมานา สาทิยมานา ปตฺถยมานา ปิหยมานา อภิชปฺปมานาติ อาสีสมานา.}
\prose{\csromanfolio{53}\textbf{ปุณฺณก} \textbf{อิตฺถตฺตนฺ}ติ เอตฺถ อตฺตภาวาภินิพฺพตฺติํ อาสีสมานา เอตฺถ ขตฺติยมหาสาลกุเล อตฺตภาวาภินิพฺพตฺติํ อาสีสมานา ฯเปฯ.\csromanspacer{} เอตฺถ พฺรหฺมกายิเกสุ เทเวสุ อตฺตภาวาภินิพฺพตฺติํ อาสีสมานา อิจฺฉยมานา สาทิยมานา ปตฺถยมานา ปิหยมานา อภิชปฺปมานาติ อาสีสมานา \textbf{ปุณฺณก} อิตฺถตฺตํ.}

\prose{\textbf{ชรํ สิตา ยญฺญมกปฺปยิํสู}ติ ชรานิสฺสิตา พฺยาธินิสฺสิตา มรณนิสฺสิตา โสกปริเทวทุกฺขโทมนสฺสุปายาสนิสฺสิตา.\csromanspacer{} ยเทว เต ชาตินิสฺสิตา, ตเทว เต ชรานิสฺสิตา.\csromanspacer{} ยเทว เต ชรานิสฺสิตา, ตเทว เต พฺยาธินิสฺสิตา.\csromanspacer{} ยเทว เต พฺยาธินิสฺสิตา, ตเทว เต มรณนิสฺสิตา.\csromanspacer{} ยเทว เต มรณนิสฺสิตา, ตเทว เต โสกปริเทวทุกฺขโทมนสฺสุปายาสนิสฺสิตา.\csromanspacer{} ยเทว เต โสกปริเทวทุกฺขโทมนสฺสุปายาสนิสฺสิตา, ตเทว เต คตินิสฺสิตา.\csromanspacer{} ยเทว เต คตินิสฺสิตา, ตเทว เต อุปปตฺตินิสฺสิตา.\csromanspacer{} ยเทว เต อุปปตฺตินิสฺสิตา, ตเทว เต ปฏิสนฺธินิสฺสิตา.\csromanspacer{} ยเทว เต ปฏิสนฺธินิสฺสิตา, ตเทว เต ภวนิสฺสิตา.\csromanspacer{} ยเทว เต ภวนิสฺสิตา, ตเทว เต สํสารนิสฺสิตา.\csromanspacer{} ยเทว เต สํสารนิสฺสิตา, ตเทว เต วฏฺฏนิสฺสิตา อลฺลีนา อุปคตา อชฺโฌสิตา อธิมุตฺตาติ ชรํ สิตา ยญฺญมกปฺปยิํสุ.\csromanspacer{} เตนาห ภควา\csromandash{}}

\vspace{\tipitakaparskip}
\csromancenter{เย เกจิเม อิสโย มนุชา, (ปุณฺณกาติ ภควา,)}
\vspace{0.5\baselineskip}
\csromangathameasure{ยญฺญมกปฺปยิํสุ ปุถูธ โลเก, \\ อาสีสมานา \textbf{ปุณฺณก} อิตฺถตฺตํ.}
\csromangathagroup{\csromangathastanza{ยญฺญมกปฺปยิํสุ ปุถูธ โลเก, \\ อาสีสมานา \textbf{ปุณฺณก} อิตฺถตฺตํ.}}

\prose{\textbf{ชรํ สิตา ยญฺญมกปฺปยิํสู}ติ.}

\csromanhangingbegin
\hangingheaditem{14}{\csromanfolio{54}\textbf{เย เกจิเม อิสโย มนุชา}, (อิจฺจายสฺมา ปุณฺณาโก,)}
\hangingbody{\textbf{ขตฺติยา พฺรหฺมณา เทวตานํ,}}
\hangingbody{\textbf{ยญฺญมกปฺปยิํสุ ปุถูธ โลเก,}}
\hangingbody{\textbf{กจฺจิสุ เต ภควา ยญฺญปเถ อปฺปมตฺตา.}}
\hangingbody{\textbf{อตารุ1} \textbf{ชาติญฺจ ชรญฺจ มาริส,}}
\hangingbody{ปุจฺฉามิ ตํ \textbf{ภควา} พฺรูหิ เมตํ.}
\csromanhangingend

\prose{\textbf{เย เกจิเม อิสโย มนุชา}ติ \textbf{เย เกจี}ติ ฯเปฯ.}

\prose{\textbf{กจฺจิสุ เต ภควา ยญฺญปเถ อปฺปมตฺตา}ติ \textbf{กจฺจิสู}ติ สํสยปุจฺฉา วิมติปุจฺฉา เทฺวฬฺหกปุจฺฉา อเนกํสปุจฺฉา “เอวํ นุ โข, น นุ โข กิํ นุ โข, กถํ นุ โข”ติ กจฺจิสุ.\csromanspacer{} \textbf{เต}ติ ยญฺญยาชกา วุจฺจนฺติ.\csromanspacer{} \textbf{ภควา}ติ คารวาธิวจนํ ฯเปฯ สจฺฉิกา ปญฺญตฺติ, ยทิทํ \textbf{ภควา}ติ กจฺจิสุ เต \textbf{ภควา}.\csromanspacer{}\csromanspacer{}\textbf{ยญฺญปเถ อปฺปมตฺตา}ติ ยญฺโญเยว วุจฺจติ ยญฺญปโถ.\csromanspacer{} ยถา อริยมคฺโค อริยปโถ เทวมคฺโค เทวปโถ พฺรหฺมมคฺโค พฺรหฺมปโถ.\csromanspacer{} เอวเมว ยญฺโญเยว วุจฺจติ ยญฺญปโถ.\csromanspacer{} \textbf{อปฺปมตฺตา}ติ ยญฺญปเถ อปฺปมตฺตา สกฺกจฺจการิโน สาตจฺจการิโน อฏฺฐิตการิโน อโนลีนวุตฺติโน อนิกฺขิตฺตจฺฉนฺทา อนิกฺขิตฺตธุรา ตจฺจริตา ตพฺพหุลา ตคฺครุกา ตนฺนินฺนา ตปฺโปณา ตปฺปพฺภารา ตทธิมุตฺตา ตทธิปเตยฺยาติ เตปิ ยญฺญปเถ อปฺปมตฺตา.\csromanspacer{} เยปิ ยญฺญํ เอสนฺติ คเวสนฺติ ปริเยสนฺติ จีวรปิณฺฑปาตเสนาสนคิลานปจฺจยเภสชฺชปริกฺขารํ อนฺนํ ปานํ ฯเปฯ เสยฺยาวสถปทีเปยฺยํ สกฺกจฺจการิโน ฯเปฯ ตทธิปเตยฺยา.\csromanspacer{} \textbf{เต}ปิ ยญฺญปเถ อปฺปมตฺตา.\csromanspacer{} เยปิ ยญฺญํ อภิสงฺขโรนฺติ จีวรปิณฺฑปาตเสนาสนคิลานปจฺจยเภสชฺชปริกฺขารํ อนฺนํ ปานํ ฯเปฯ เสยฺยาวสถปทีเปยฺยํ สกฺกจฺจการิโน ฯเปฯ ตทธิปเตยฺยา.\csromanspacer{} \textbf{เต}ปิ ยญฺญปเถ อปฺปมตฺตา.\csromanspacer{} เยปิ ยญฺญํ เทนฺติ ยชนฺติ ปริจฺจชนฺติ จีวรปิณฺฑปาตเสนาสนคิลานปจฺจยเภสชฺชปริกฺขารํ อนฺนํ ปานํ ฯเปฯ เสยฺยาวสถปทีเปยฺยํ สกฺกจฺจการิโน ฯเปฯ ตทธิปเตยฺยา.\csromanspacer{} \textbf{เต}ปิ ยญฺญปเถ อปฺปมตฺตาติ กจฺจิสุ เต \textbf{ภควา} ยญฺญปเถ อปฺปมตฺตา.}

\prose{\textbf{อตารุ ชาติญฺจ ชรญฺจ มาริสา}ติ ชรามรณํ อตริํสุ อุตฺตริํสุ ปตริํสุ สมติกฺกมิํสุ วีติวตฺติํสุ.\csromanspacer{} \textbf{มาริสา}ติ ปิยวจนํ ครุวจนํ สคารวสปฺปติสฺสาธิวจนเมตํ มาริสาติ อตารุ ชาติญฺจ ชรญฺจ มาริส.}

\csromanheader{2}{3. ปุณฺณกมาณวปุจฺฉานิทฺเทส}
\prose{\csromanfolio{55}\textbf{ปุจฺฉามิ ตํ ภควา พฺรูหิ เมตนฺ}ติ \textbf{ปุจฺฉามิ} \textbf{ตนฺ}ติ \textbf{ปุจฺฉามิ} ตํ ยาจามิ ตํ อชฺเฌสามิ ตํ ปสาเทมิ ตํ กถยสฺสุ เมติ \textbf{ปุจฺฉามิ} ตํ.\csromanspacer{} \textbf{ภควา}ติ คารวาธิวจนํ ฯเปฯ สจฺฉิกา ปญฺญตฺติ, ยทิทํ \textbf{ภควา}ติ.\csromanspacer{} \textbf{พฺรูหิ เมตนฺ}ติ พฺรูหิ อาจิกฺขาหิ เทเสหิ ปญฺญเปหิ ปฏฺฐเปหิ วิวราหิ วิภชาหิ อุตฺตานีกโรหิ ปกาเสหีติ \textbf{ปุจฺฉามิ} ตํ \textbf{ภควา} พฺรูหิ เมตํ.\csromanspacer{} เตนาห โส พฺราหฺมโณ\csromandash{}}

\prose{เย เกจิเม อิสโย มนุชา, (อิจฺจายสฺมา ปุณฺณโก,)}

\prose{ขตฺติยา พฺรหฺมณา เทวตานํ.}

\csromangathameasure{ยญฺญมกปฺปยิํสุ ปุถูธ โลเก, \\ กจฺจิสุ เต \textbf{ภควา} ยญฺญปเถ อปฺปมตฺตา. \\ อตารุ ชาติญฺจ ชรญฺจ มาริส,}
\csromangathagroup{\csromangathastanza{ยญฺญมกปฺปยิํสุ ปุถูธ โลเก, \\ กจฺจิสุ เต \textbf{ภควา} ยญฺญปเถ อปฺปมตฺตา. \\ อตารุ ชาติญฺจ ชรญฺจ มาริส,}}

\prose{\textbf{ปุจฺฉามิ ตํ ภควา พฺรูหิ เม ตนฺ}ติ.}

\csromanhangingbegin
\hangingheaditem{15}{\textbf{อาสีสนฺติ}\footnote{อาสิํสนฺติ (สฺยา)} \textbf{โถมยนฺติ อภิชปฺปนฺติ ชุหนฺติ, (ปุณฺณกาติ ภควา,)}}
\hangingbody{\textbf{กามาภิชปฺปนฺติ ปฏิจฺจ ลาภํ.}}
\hangingbody{\textbf{เต ยาชโยคา ภวราครตฺตา,}}
\hangingbody{นาตริํสุ ชาติชรนฺติ พฺรูมิ.}
\csromanhangingend

\prose{\textbf{อาสีสนฺติ โถมยนฺติ อภิชปฺปนฺติ ชุหนฺติ}ติ \textbf{อาสีสนฺตี}ติ รูปปฏิลาภํ อาสีสนฺติ.\csromanspacer{} สทฺทปฏิลาภํ อาสีสนฺติ.\csromanspacer{} คนฺธปฏิลาภํ อาสีสนฺติ.\csromanspacer{} รสปฏิลาภํ อาสีสนฺติ.\csromanspacer{} โผฏฺฐพฺพปฏิลาภํ อาสีสนฺติ.\csromanspacer{} ปุตฺตปฏิลาภํ อาสีสนฺติ.\csromanspacer{} ทารปฏิลาภํ อาสีสนฺติ.\csromanspacer{} ธนปฏิลาภํ อาสีสนฺติ.\csromanspacer{} ยสปฏิลาภํ อาสีสนฺติ.\csromanspacer{} อิสฺสริยปฏิลาภํ อาสีสนฺติ ขตฺติยมหาสาลกุเล อตฺตภาวปฏิลาภํ อาสีสนฺติ.\csromanspacer{} พฺราหฺมณมหาสาลกุเล.\csromanspacer{} คหปติมหาสาลกุเล อตฺตภาวปฏิลาภํ อาสีสนฺติ.\csromanspacer{} จาตุมหาราชิเกสุ เทเวสุ ฯเปฯ.\csromanspacer{} พฺรหฺมกายิเกสุ เทเวสุ อตฺตภาวปฏิลาภํ อาสีสนฺติ อิจฺฉนฺติ สาทิยนฺติ ปตฺถยนฺติ ปิหยนฺตีติ อาสีสนฺติ.}

\prose{\textbf{โถมยนฺตี}ติ ยญฺญํ วา โถเมนฺติ ผลํ วา โถเมนฺติ ทกฺขิเณยฺเย วา โถเมนฺติ.\csromanspacer{} กถํ ยญฺญํ โถเมนฺติ.\csromanspacer{} สุจิํ ทินฺนํ\footnote{ปิยํ ทินฺนํ (สฺยา)} มนาปํ ทินฺนํ ปณีตํ ทินฺนํ \csromanfolio{56}กาเลน ทินฺนํ กปฺปิยํ ทินฺนํ วิเจยฺย ทินฺนํ อนวชฺชํ ทินฺนํ อภิณฺหํ ทินฺนํ ททํ จิตฺตํ ปสาทิตนฺติ โถเมนฺติ กิตฺ\textbf{เต}นฺติ วณฺเณนฺติ ปสํสนฺติ.\csromanspacer{} เอวํ ยญฺญํ โถเมนฺติ.}

\prose{กถํ ผลํ โถเมนฺติ.\csromanspacer{} อิโต นิทานํ รูปปฏิลาโภ ภวิสฺสติ ฯเปฯ.\csromanspacer{} พฺรหฺมกายิเกสุ เทเวสุ อตฺตภาวปฏิลาโภ ภวิสฺสตีติ โถเมนฺติ กิตฺ\textbf{เต}นฺติ วณฺเณนฺติ ปสํสนฺติ เอวํ ผลํ โถเมนฺติ.}

\prose{กถํ ทกฺขิเณยฺเย โถเมนฺติ.\csromanspacer{} ทกฺขิเณยฺยา ชาติสมฺปนฺนา โคตฺตสมฺปนฺนา อชฺฌายกา มนฺตธรา ติณฺณํ เวทานํ ปารคู สนิฆณฺฑุเกฏุภานํ สากฺขรปฺปเภทานํ อิติหาสปญฺจมานํ ปทกา เวยฺยากรณา โลกายตมหาปุริสลกฺขเณสุ อนวยาติ, วีตราคา วา ราควินยาย วา ปฏิปนฺนา, วีตโทสา วา โทสวินยาย วา ปฏิปนฺนา, วีตโมหา วา โมหวินยาย วา ปฏิปนฺนา, สทฺธาสมฺปนฺนา สีลสมฺปนฺนา สมาธิสมฺปนฺนา ปญฺญาสมฺปนฺนา วิมุตฺติสมฺปนฺนา วิมุตฺติญาณทสฺสนสมฺปนฺนาติ โถเมนฺติ กิตฺ\textbf{เต}นฺติ วณฺเณนฺติ ปสํสนฺติ, เอวํ ทกฺขิเณยฺเย โถเมนฺตีติ อาสีสนฺติ โถมยนฺติ.}

\prose{\textbf{อภิชปฺปนฺตี}ติ รูปปฏิลาภํ อภิชปฺปนฺติ, สทฺทปฏิลาภํ อภิชปฺปนฺติ.\csromanspacer{} คนฺธปฏิลาภํ อภิชปฺปนฺติ.\csromanspacer{} รสปฏิลาภํ อภิชปฺปนฺติ ฯเปฯ.\csromanspacer{} พฺรหฺมกายิเกสุ เทเวสุ อตฺตภาวปฏิลาภํ อภิชปฺปนฺตีติ อาสีสนฺติ โถมยนฺติ อภิชปฺปนฺติ.\csromanspacer{} \textbf{ชุหนฺตี}ติ ชุหนฺติ เทนฺติ ยชนฺติ ปริจฺจชนฺติ จีวรปิณฺฑปาตเสนาสนคิลานปจฺจยเภสชฺชปริกฺขารํ อนฺนํ ปานํ วตฺถํ ยานํ มาลาคนฺธวิเลปนํ เสยฺยาวสถปทีเปยฺยนฺติ อาสีสนฺติ โถมยนฺติ อภิชปฺปนฺติ ชุหนฺติ.\csromanspacer{} ปุณฺณกาติ ภควา.}

\prose{\textbf{กามาภิชปฺปนฺติ ปฏิจฺจ ลาภนฺ}ติ รูปปฏิลาภํ ปฏิจฺจ กาเม อภิชปฺปนฺติ สทฺทปฏิลาภํ ปฏิจฺจ กาเม อภิชปฺปนฺติ ฯเปฯ.\csromanspacer{} พฺรหฺมกายิเกสุ เทเวสุ อตฺตภาวปฏิลาภํ ปฏิจฺจ กาเม อภิชปฺปนฺติ ปชปฺปนฺตีติ กามาภิชปฺปนฺติ ปฏิจฺจ ลาภํ.}

\prose{\textbf{เต ยาชโยคา ภวราครตฺตา, นาตริํสุ ชาติชรนฺติ พฺรูมี}ติ \textbf{เต}ติ ยญฺญยาชกา วุจฺจนฺติ.\csromanspacer{} \textbf{ยาชโยคา}ติ ยาชโยเคสุ ยุตฺตา ปยุตฺตา อายุตฺตา สมายุตฺตา ตจฺจริตา ตพฺพหุลา ตคฺครุกา ตนฺนินฺนา ตปฺโปณา ตปฺปพฺภารา ตทธิมุตฺตา ตทธิป\textbf{เต}ยฺยาติ \textbf{เต} ยาชโยคา.\csromanspacer{} \textbf{ภวราครตฺตา}ติ ภวราโค วุจฺจติ โย ภเวสุ ภวจฺฉนฺโท ภวราโค}

\vspace{\tipitakaparskip}
\csromancenter{ปุณฺณกมาณวปุจฺฉานิทฺเทส ภวนนฺที ภวตณฺหา ภวสิเนโห ภวปริฬาโห ภวมุจฺฉา ภวชฺโฌสานํ.\csromanspacer{} ภวราเคน ภเวสุ รตฺตา คิทฺธา คธิตา มุจฺฉิตา อชฺโฌสนฺนา ลคฺคา ลคฺคิตา ปลิพุทฺธาติ เต ยาชโยคา ภวราครตฺตา.}
\prose{\csromanfolio{57}\textbf{นาตริํสุ ชาติชรนฺติ พฺรูมี}ติ เต ยาชโยคา ภวราครตฺตา ชาติชรามรณํ นาตริํสุ น อุตฺตริํสุ น ปตริํสุ น สมติกฺกมิํสุ น วีติวตฺติํสุ ชาติชรามรณา อนิกฺขนฺตา อนิสฺสฏา อนติกฺกนฺตา อสมติกฺกนฺตา อวีติวตฺตา อนฺโตชาติชรามรเณ ปริวตฺตนฺติ.\csromanspacer{} อนฺโตสํสารปเถ ปริวตฺตนฺติ.\csromanspacer{} ชาติยา อนุคตา ชราย อนุสฏา พฺยาธินา อภิภูตา มรเณน อพฺภาหตา อตาณา อเลณา อสรณา อสรณีภูตาติ พฺรูมิ อาจิกฺขามิ เทเสมิ ปญฺญเปมิ ปฏฺฐเปมิ วิวรามิ วิภชามิ อุตฺตานีกโรมิ ปกาเสมีติ เต ยาชโยคา ภวราครตฺตา, นาตริํสุ ชาติชรนฺติ พฺรูมิ.\csromanspacer{} เตนาห ภควา\csromandash{}}

\prose{อาสีสนฺติ โถมยนฺติ อภิชปฺปนฺติ ชุหนฺติ, (ปุณฺณกาติ ภควา,)}

\prose{กามา’ภิชปฺปนฺติ ปฏิจฺจ ลาภํ.}

\prose{เต ยาชโยคา ภวราครตฺตา,}

\prose{\textbf{นาตริํสุ ชาติชรนฺติ พฺรูมี}ติ.}

\csromanhangingbegin
\hangingheaditem{16}{\textbf{เต เจ นาตริํสุ ยาชโยคา}, (อิจฺจายสฺมา ปุณฺณโก,)}
\hangingbody{\textbf{ยญฺเญหิ ชาติญฺจ ชรญฺจ มาริส.}}
\hangingbody{\textbf{อถ โก จรหิ เทวมนุสฺสโลเก,}}
\hangingbody{\textbf{อตาริ ชาติญฺจ ชรญฺจ มาริส.}}
\hangingbody{ปุจฺฉามิ ตํ ภควา พฺรูหิ เมตํ.}
\csromanhangingend

\prose{\textbf{เต เจ นาตริํสุ ยาชโยคา}ติ เต ยญฺญยาชกา ยาชโยคา ภวราครตฺตา ชาติชรามรณํ นาตริํสุ น อุตฺตริํสุ น ปตริํสุ น สมติกฺกมิํสุ น วีติวตฺติํสุ, ชาติชรามรณา อนิกฺขนฺตา อนิสฺสฏา อนติกฺกนฺตา อสมติกฺกนฺตา อวีติวตฺตา อนฺโตชาติชรามรเณ ปริวตฺตนฺติ.\csromanspacer{} อนฺโตสํสารปเถ ปริวตฺตนฺติ.\csromanspacer{} ชาติยา อนุคตา ชราย อนุสฏา พฺยาธินา อภิภูตา มรเณน อพฺภาหตา อตาณา อเลณา อสรณา อสรณีภูตาติ เต เจ นาตริํสุ ยาชโยคา.}

\prose{\textbf{อิจฺจายสฺมา ปุณฺณโก}ติ อิจฺจาติ ปทสนฺธิ ฯเปฯ อายสฺมา ปุณฺณโก.}

\prose{\csromanfolio{58}\textbf{ยญฺเญหิ ชาติญฺจ ชรญฺจ มาริสา}ติ \textbf{ยญฺเญหี}ติ ยญฺเญหิ ปหูเตหิ ยญฺเญหิ วิวิเธหิ ยญฺเญหิ ปุถูหิ.\csromanspacer{} \textbf{มาริสา}ติ ปิยวจนํ ครุวจนํ สคารวสปฺปติสฺสาธิวจนเมตํ มาริสาติ ยญฺเญหิ ชาติญฺจ ชรญฺจ มาริส.}

\prose{\textbf{อถ โก จรหิ เทวมนุสฺสโลเก, อตาริ ชาติญฺจ ชรญฺจ มาริสา}ติ อถ โก เอโส สเทวเก โลเก สมารเก สพฺรหฺมเก สสฺสมณพฺราหฺมณิยา ปชาย สเทวมนุสฺสาย ชาติชรามรณํ อตริ อุตฺตริ ปตริ สมติกฺกมิ วีติวตฺตยิ\footnote{วีติวตฺติ (ก)}.\csromanspacer{} \textbf{มาริสา}ติ ปิยวจนํ ครุวจนํ สคารวสปฺปติสฺสาธิวจนเมตํ มาริสาติ อถ โก จรหิ เทวมนุสฺสโลเก, อตาริ ชาติญฺจ ชรญฺจ มาริส.}

\prose{\textbf{ปุจฺฉามิ ตํ ภควา พฺรูหิ เมตนฺ}ติ \textbf{ปุจฺฉามิ ตนฺ}ติ ปุจฺฉามิ ตํ ยาจามิ ตํ อชฺเฌสามิ ตํ ปสาเทมิ ตํ, กถยสฺสุ เมตนฺติ ปุจฺฉามิ ตํ.\csromanspacer{} \textbf{ภควา}ติ คารวาธิวจนํ ฯเปฯ สจฺฉิกา ปญฺญตฺติ, ยทิทํ \textbf{ภควา}ติ.\csromanspacer{} \textbf{พฺรูหิ เมตนฺ}ติ พฺรูหิ อาจิกฺขาหิ เทเสหิ ปญฺญเปหิ ปฏฺฐเปหิ วิวราหิ วิภชาหิ อุตฺตานีกโรหิ ปกาเสหีติ ปุจฺฉามิ ตํ \textbf{ภควา} พฺรูหิ เมตํ.\csromanspacer{} เตนาห โส พฺราหฺมโณ\csromandash{}}

\prose{เต เจ นาตริํสุ ยาชโยคา, (อิจฺจายสฺมา ปุณฺณโก,)}

\prose{ยญฺเญหิ ชาติญฺจ ชรญฺจ มาริส.}

\prose{อถ โก จรหิ เทวมนุสฺสโลเก,}

\prose{อตาริ ชาติญฺจ ชรญฺจ มาริส.}

\prose{\textbf{ปุจฺฉามิ ตํ ภควา พฺรูหิ เมตนฺ}ติ.}

\csromanhangingbegin
\hangingheaditem{17}{สงฺขาย โลกสฺมิ\footnote{โลกสฺมิํ (สฺยา, ก)} \textbf{ปโรปรานิ, (ปุณฺณกาติ ภควา,)}}
\hangingbody{\textbf{ยสฺสิญฺชิตํ นตฺถิ กุหิญฺจิ โลเก.}}
\hangingbody{\textbf{สนฺโต วิธูโม อนีโฆ นิราโส,}}
\hangingbody{อตาริ โส ชาติชรนฺติ พฺรูมิ.}
\csromanhangingend

\prose{\textbf{สงฺขาย โลกสฺมิ ปโรปรานี}ติ สงฺขา วุจฺจติ ญาณํ.\csromanspacer{} ยา ปญฺญา ปชานนา ฯเปฯ อโมโห ธมฺมวิจโย สมฺมาทิฏฺฐิ.\csromanspacer{} \textbf{ปโรปรานี}ติ \textbf{โอรํ} วุจฺจติ สกตฺตภาโว.\csromanspacer{} \textbf{ปรํ} วุจฺจติ ปรตฺตภาโว.\csromanspacer{} โอรํ วุจฺจติ สกรูปเวทนาสญฺญาสงฺขารวิญฺญาณํ.\csromanspacer{} \textbf{ปรํ} วุจฺจติ ปรรูปเวทนาสญฺญาสงฺขาร-}

\vspace{\tipitakaparskip}
\csromancenter{ปุณฺณกมาณวปุจฺฉานิทฺเทส วิญฺญาณํ\textbf{.}\csromanspacer{} โอรํ วุจฺจติ ฉ อชฺฌตฺติกานิ อายตนานิ\textbf{.}\csromanspacer{} ปรํ วุจฺจติ ฉ พาหิรานิ อายตนานิ\textbf{.}\csromanspacer{} โอรํ วุจฺจติ มนุสฺสโลโก\textbf{.}\csromanspacer{} ปรํ วุจฺจติ เทวโลโก\textbf{.}\csromanspacer{} โอรํ วุจฺจติ กามธาตุ\textbf{.}\csromanspacer{} ปรํ วุจฺจติ รูปธาตุ อรูปธาตุ\textbf{.}\csromanspacer{} โอรํ วุจฺจติ กามธาตุ รูปธาตุ\textbf{.}\csromanspacer{} ปรํ วุจฺจติ อรูปธาตุ\textbf{.}\csromanspacer{} สงฺขาย โลกสฺมิ ปโรปรานีติ ปโรปรานิ อนิจฺจโต สงฺขาย, ทุกฺขโต\textbf{.}\csromanspacer{} โรคโต\textbf{.}\csromanspacer{} คณฺฑโต ฯเปฯ\textbf{.}\csromanspacer{} นิสฺสรณโต สงฺขาย ชานิตฺวา ตุลยิตฺวา ตีรยิตฺวา วิภาวยิตฺวา วิภูตํ กตฺวาติ สงฺขาย โลกสฺมิ ปโรปรานิ\textbf{.}\csromanspacer{} \textbf{ปุณฺณกาติ ภควา}ติ \textbf{ปุณฺณกาติ ภควา} ตํ พฺราหฺมณํ นาเมน อาลปติ\textbf{.}\csromanspacer{} \textbf{ภควา}ติ คารวาธิวจนเมตํ ฯเปฯ ยทิทํ \textbf{ภควา}ติ \textbf{ปุณฺณกาติ ภควา.}}
\prose{\csromanfolio{59}\textbf{ยสฺสิญฺชิตํ นตฺถิ กุหิญฺจิ โลเก}ติ \textbf{ยสฺสา}ติ อรหโต ขีณาสวสฺส\textbf{.}\csromanspacer{} \textbf{อิญฺชิตนฺ}ติ ตณฺหิญฺชิตํ ทิฏฺฐิญฺชิตํ มานิญฺชิตํ กิเลสิญฺชิตํ กามิญฺชิตํ\textbf{.}\csromanspacer{} ยสฺสิเม อิญฺชิตา นตฺถิ น สนฺติ น สํวิชฺชนฺติ นุปลพฺภนฺติ, ปหีนา สมุจฺฉินฺนา วูปสนฺตา ปฏิปสฺสทฺธา อภพฺพุปฺปตฺติกา ญาณคฺคินา ทฑฺฒา\textbf{.}\csromanspacer{} \textbf{กุหิญฺจี}ติ กุหิญฺจิ กิสฺมิญฺจิ กตฺถจิ อชฺฌตฺตํ วา พหิทฺธา วา อชฺฌตฺตพหิทฺธา วา\textbf{.}\csromanspacer{} \textbf{โลเก}ติ อปายโลเก ฯเปฯ อายตนโลเกติ ยสฺสิญฺชิตํ นตฺถิ กุหิญฺจิ โลเก\textbf{.}}

\prose{\textbf{สนฺโต วิธูโม อนีโฆ นิราโส, อตาริ โส ชาติชรนฺติ พฺรูมี}ติ \textbf{สนฺโต}ติ ราคสฺส สนฺตตฺตา \textbf{สนฺโต.}\csromanspacer{} โทสสฺส\textbf{.}\csromanspacer{} โมหสฺส\textbf{.}\csromanspacer{} โกธสฺส\textbf{.}\csromanspacer{} อุปนาหสฺส\textbf{.}\csromanspacer{} มกฺขสฺส ฯเปฯ\textbf{.}\csromanspacer{} สพฺพากุสลาภิสงฺขารานํ สนฺตตฺตา สมิตตฺตา วูปสมิตตฺตา วิชฺฌาตตฺตา\footnote{นิชฺฌาตตฺตา (ก) ขุ 7. 53 ปิฏฺเฐ.} นิพฺพุตตฺตา วิคตตฺตา ปฏิปสฺสทฺธตฺตา \textbf{สนฺโต} อุป\textbf{สนฺโต} วูป\textbf{สนฺโต} นิพฺพุโต ปฏิปสฺสทฺโธติ \textbf{สนฺโต.}\csromanspacer{} \textbf{วิธูโม}ติ กายทุจฺจริตํ วิธูมิตํ วิธมิตํ โสสิตํ วิโสสิตํ พฺยนฺตีกตํ\footnote{พฺยนฺติกตํ (ก)}\textbf{.}\csromanspacer{} วจีทุจฺจริตํ\textbf{.}\csromanspacer{} มโนทุจฺจริตํ วิธูมิตํ วิธมิตํ โสสิตํ วิโสสิตํ พฺยนฺตีกตํ ราโค, โทโส, โมโห วิธูมิโต วิธมิโต โสสิโต วิโสสิโต พฺยนฺตีกโต\textbf{.}\csromanspacer{} โกโธ\textbf{.}\csromanspacer{} อุปนาโห\textbf{.}\csromanspacer{} มกฺโข\textbf{.}\csromanspacer{} ปฬาโส\textbf{.}\csromanspacer{} อิสฺสา\textbf{.}\csromanspacer{} มจฺฉริยํ\textbf{.}\csromanspacer{} มายา\textbf{.}\csromanspacer{} สาเฐยฺยํ\textbf{.}\csromanspacer{} ถมฺโภ\textbf{.}\csromanspacer{} สารมฺโภ\textbf{.}\csromanspacer{} มาโน\textbf{.}\csromanspacer{} อติมาโน\textbf{.}\csromanspacer{} มโท\textbf{.}\csromanspacer{} ปมาโท\textbf{.}\csromanspacer{} สพฺเพ กิเลสา\textbf{.}\csromanspacer{} สพฺเพ ทุจฺจริตา\textbf{.}\csromanspacer{} สพฺเพ ทรถา\textbf{.}\csromanspacer{} สพฺเพ ปริฬาหา\textbf{.}\csromanspacer{} สพฺเพ สนฺตาปา\textbf{.}\csromanspacer{} สพฺพากุสลาภิสงฺขารา วิธูมิตา วิธมิตา โสสิตา วิโสสิตา พฺยนฺตีกตา\textbf{.}\csromanspacer{} อถ วา โกโธ วุจฺจติ ธูโม\textbf{.}}

\csromangathameasure{มาโน หิ เต พฺราหฺมณ ขาริภาโร, \\ โกโธ ธูโม ภสฺมนิ โมสวชฺชํ. \\ ชิวฺหา สุชา หทยํ โชติฐานํ, \\ อตฺตา สุทนฺโต ปุริสสฺส โชติ.}
\csromangathagroup{\csromangathastanza{\csromanfolio{60}มาโน หิ เต พฺราหฺมณ ขาริภาโร, \\ โกโธ ธูโม ภสฺมนิ\footnote{คมฺมนิ (สฺยา)} โมสวชฺชํ. \\ ชิวฺหา สุชา หทยํ\footnote{คมฺมนิ (สฺยา)} โชติฐานํ, \\ อตฺตา สุทนฺโต ปุริสสฺส โชติ.}}

\prose{อปิ จ ทสหากาเรหิ โกโธ ชายติ, “อนตฺถํ เม อจรี”ติ โกโธ ชายติ, “อนตฺถํ เม จรตี”ติ โกโธ ชายติ, “อนตฺถํ เม จริสฺสตี”ติ โกโธ ชายติ. “ปิยสฺส เม มนาปสฺส อนตฺถํ อจริ.\csromanspacer{} อนตฺถํ จรติ.\csromanspacer{} อนตฺถํ จริสฺสตี”ติ โกโธ ชายติ. “อปฺปิยสฺส เม อมนาปสฺส อตฺถํ อจริ.\csromanspacer{} อตฺถํ จรติ.\csromanspacer{} อตฺถํ จริสฺสตี”ติ โกโธ ชายติ.\csromanspacer{} อฏฺฐาเน วา ปน โกโธ ชายติ.\csromanspacer{} โย เอวรูโป จิตฺตสฺส อาฆาโต ปฏิฆาโต ปฏิฆํ ปฏิวิโรโธ โกโป ปโกโป สมฺปโกโป โทโส ปโทโส สมฺปโทโส จิตฺตสฺส พฺยาปตฺติ มโนปโทโส โกโธ กุชฺฌนา กุชฺฌิตตฺตํ โทโส ทุสฺสนา ทุสฺสิตตฺตํ พฺยาปตฺติ พฺยาปชฺชนา พฺยาปชฺชิตตฺตํ วิโรโธ ปฏิวิโรโธ จณฺฑิกฺกํ อสุโรโป\footnote{อสฺสุโรโป (สฺยา)} อนตฺตมนตา จิตฺตสฺส.\csromanspacer{} อยํ วุจฺจติ โกโธ.}

\prose{อปิ จ โกธสฺส อธิมตฺตปริตฺตตา เวทิตพฺพา, อตฺถิ กญฺจิ\footnote{กิญฺจิ (ก) ขุ 7. 165 ปิฏฺเฐ.} กาลํ โกโธ จิตฺตาวิลกรณมตฺโต โหติ, น จ ตาว มุขกุลานวิกุลาโน โหติ.\csromanspacer{} อตฺถิ กญฺจิ กาลํ โกโธ มุขกุลานวิกุลานมตฺโต โหติ.\csromanspacer{} น จ ตาว หนุสญฺโจปโน โหติ.\csromanspacer{} อตฺถิ กญฺจิ กาลํ โกโธ หนุสญฺโจปนมตฺโต โหติ, น จ ตาว ผรุสวาจํ นิจฺฉารโณ\footnote{ผรุสวาจนิจฺฉารโณ (สฺยา)} โหติ.\csromanspacer{} อตฺถิ กญฺจิ กาลํ โกโธ ผรุสวาจํ นิจฺฉารณมตฺโต โหติ, น จ ตาว ทิสาวิทิสานุวิโลกโน โหติ.\csromanspacer{} อตฺถิ กญฺจิ กาลํ โกโธ ทิสาวิทิสานุวิโลกนมตฺโต โหติ, น จ ตาว ทณฺฑสตฺถปรามสโน โหติ.\csromanspacer{} อตฺถิ กญฺจิ กาลํ โกโธ ทณฺฑสตฺถปรามสนมตฺโต โหติ, น จ ตาว ทณฺฑสตฺถ-อพฺภุกฺกิรโณ โหติ.\csromanspacer{} อตฺถิ กญฺจิ กาลํ โกโธ ทณฺฑสตฺถ-อพฺภุกฺกิรณมตฺโต โหติ, น จ ตาว ทณฺฑสตฺถ-อภินิปาตโน โหติ.\csromanspacer{} อตฺถิ กญฺจิ กาลํ โกโธ ทณฺฑสตฺถ- อภินิปาตนมตฺโต โหติ,}

\vspace{\tipitakaparskip}
\csromancenter{ปุณฺณกมาณวปุจฺฉานิทฺเทส น จ ตาว สมฺภญฺชนปลิภญฺชโน โหติ.\csromanspacer{} อตฺถิ กญฺจิ กาลํ โกโธ สมฺภญฺชนปลิภญฺชนมตฺโต โหติ, น จ ตาว องฺคมงฺค-อปกฑฺฒโน โหติ.\csromanspacer{} อตฺถิ กญฺจิ กาลํ โกโธ องฺคมงฺค-อปกฑฺฒนมตฺโต โหติ, น จ ตาว ชีวิตาโวโรปโน\footnote{ชีวิตปนาสโน (สฺยา) ขุ 7. 166 ปิฏฺเฐ.} โหติ.\csromanspacer{} อตฺถิ กญฺจิ กาลํ โกโธ ชีวิตาโวโรปนมตฺโต โหติ, น จ ตาว สพฺพจาคปริจฺจาคาย สณฺฐิโต โหติ.\csromanspacer{} ยโต โกโธ ปรํ ปุคฺคลํ ฆาเตตฺวา อตฺตานํ ฆาเตติ.\csromanspacer{} เอตฺตาวตา โกโธ ปรมุสฺสทคโต ปรมเวปุลฺลปตฺโต โหติ.\csromanspacer{} ยสฺส โส โหติ โกโธ ปหีโน สมุจฺฉินฺโน วูปสนฺโต ปฏิปสฺสทฺโธ อภพฺพุปฺปตฺติโก ญาณคฺคินา ทฑฺโฒ.\csromanspacer{} โส วุจฺจติ วิธูโม.}
\prose{\csromanfolio{61}โกธสฺส ปหีนตฺตา วิธูโม โกธวตฺถุสฺส ปริญฺญาตตฺตา วิธูโม โกธเหตุสฺส ปริญฺญาตตฺตา วิธูโม โกธเหตุสฺส อุปจฺฉินฺนตฺตา วิธูโม.\csromanspacer{} \textbf{อนีโฆ}ติ ราโค นีโฆ, โทโส นีโฆ, โมโห นีโฆ, โกโธ นีโฆ, อุปนาโห นีโฆ ฯเปฯ.\csromanspacer{} สพฺพากุสลาภิสงฺขารา นีฆา, ยสฺเสเต นีฆา ปหีนา สมุจฺฉินฺนา วูปสนฺตา ปฏิปสฺสทฺธา อภพฺพุปฺปตฺติกา ญาณคฺคินา ทฑฺฒา, โส วุจฺจติ อนีโฆ.}

\prose{\textbf{นิราโส}ติ \textbf{อาสา} วุจฺจติ ตณฺหา.\csromanspacer{} โย ราโค สาราโค ฯเปฯ อภิชฺฌา โลโภ อกุสลมูลํ, ยสฺเสสา \textbf{อาสา} ตณฺหา ปหีนา สมุจฺฉินฺนา วูปสนฺตา ปฏิปสฺสทฺธา อภพฺพุปฺปตฺติกา ญาณคฺคินา ทฑฺฒา, โส วุจฺจติ นิราโส.\csromanspacer{} \textbf{ชาตี}ติ ยา เตสํ เตสํ สตฺตานํ ตมฺหิ ตมฺหิ สตฺตนิกาเย ชาติ สญฺชาติ โอกฺกนฺติ อภินิพฺพตฺติ ขนฺธานํ ปาตุภาโว อายตนานํ ปฏิลาโภ.\csromanspacer{} \textbf{ชรา}ติ ยา เตสํ เตสํ สตฺตานํ ตมฺหิ ตมฺหิ สตฺตนิกาเย ชรา ชีรณตา ขณฺฑิจฺจํ ปาลิจฺจํ วลิตฺตจตา อายุโน สํหานิ อินฺทฺริยานํ ปริปาโก.\csromanspacer{} \textbf{สนฺโต วิธูโม อนีโฆ นิราโส, อตาริ โส ชาติชรนฺติ พฺรูมี}ติ โย สนฺโต จ วิธูโม จ อนีโฆ จ นิราโส จ, โส ชาติชรามรณํ อตริ อุตฺตริ ปตริ สมติกฺกมิ วีติวตฺตยีติ พฺรูมิ อาจิกฺขามิ เทเสมิ ปญฺญเปมิ ปฏฺฐเปมิ วิวรามิ วิภชามิ อุตฺตานีกโรมิ ปกาเสมีติ สนฺโต วิธูโม อนีโฆ นิราโส, อตาริ โส ชาติชรนฺติ พฺรูมิ.\csromanspacer{} เตนาห ภควา\csromandash{}}

\prose{\csromanfolio{62}สงฺขาย โลกสฺมิ ปโรปรานิ, (ปุณฺณกาติ ภควา,)}

\prose{ยสฺสิญฺชิตํ นตฺถิ กุหิญฺจิ โลเก.}

\csromangathameasure{สนฺโต วิธูโม อนีโฆ นิราโส, \\ อตาริ โส ชาติชรนฺติ พฺรูมีติ.}
\csromangathagroup{\csromangathastanza{สนฺโต วิธูโม อนีโฆ นิราโส, \\ อตาริ โส ชาติชรนฺติ พฺรูมีติ.}}

\prose{สหคาถาปริโยสานา ฯเปฯ ปญฺชลิโก ภควนฺตํ นมสฺสมาโน นิสินฺโน โหติ “สตฺถา เม ภนฺเต ภควา, สาวโกหมสฺมี”ติ.}

\vspace{\tipitakaparskip}
\csromancenter{ปุณฺณกมาณวปุจฺฉานิทฺเทโส ตติโย.}
\csromanmatikaanchor{mtk.25}
\csromantocmarkref{subsection}{4. เมตฺตคูมาณวปุจฺฉานิทฺเทส}{mtk.25}
\bookmark[dest={mtk.25},level=2]{4. เมตฺตคูมาณวปุจฺฉานิทฺเทส}
\csromanheader{2}{4. เมตฺตคูมาณวปุจฺฉานิทฺเทส}
\csromanhangingbegin
\hangingheaditem{18}{ปุจฺฉามิ ตํ ภควา พฺรูหิ เม’ตํ, (อิจฺจายสฺมา เมตฺตคู,)}
\hangingbody{\textbf{มญฺญามิ ตํ เวทคู ภาวิตตฺตํ.}}
\hangingbody{\textbf{กุโต นุ ทุกฺขา สมุทาคตา อิเม,}}
\hangingbody{เย เกจิ โลกสฺมิ’มเนกรูปา.}
\csromanhangingend

\prose{\textbf{ปุจฺฉามิตํ ภควา พฺรูหิ เม’ตนฺ}ติ \textbf{ปุจฺฉามี}ติ ติสฺโส ปุจฺฉา อทิฏฺฐโชตนา ปุจฺฉา ทิฏฺฐสํสนฺทนา ปุจฺฉา วิมติจฺเฉทนา ปุจฺฉา.\csromanspacer{} กตมา อทิฏฺฐโชตนา ปุจฺฉา, ปกติยา ลกฺขณํ อญฺญาตํ โหติ อทิฏฺฐํ อตุลิตํ อตีริตํ อวิภูตํ อวิภาวิตํ, ตสฺส ญาณาย ทสฺสนาย ตุลนาย ตีรณาย วิภูตตฺถาย วิภาวนตฺถาย ปญฺหํ ปุจฺฉติ, อยํ อทิฏฺฐโชตนา ปุจฺฉา.}

\prose{กตมา ทิฏฺฐสํสนฺทนา ปุจฺฉา, ปกติยา ลกฺขณํ ญาตํ โหติ ทิฏฺฐํ ตุลิตํ ตีริตํ วิภูตํ วิภาวิตํ, อญฺเญหิ ปณฺฑิเตหิ สทฺธิํ สํสนฺทนตฺถาย ปญฺหํ ปุจฺฉติ, อยํ ทิฏฺฐสํสนฺทนา ปุจฺฉา.}

\prose{กตมา วิมติจฺเฉทนา ปุจฺฉา, ปกติยา สํสยปกฺขนฺโท โหติ วิมติปกฺขนฺโท เทฺวฬฺหกชาโต “เอวํ นุ โข, น นุ โข กิํ นุ โข กถํ นุ โข”ติ, โส วิมติจฺเฉทนตฺถาย ปญฺหํ ปุจฺฉติ.\csromanspacer{} อยํ วิมติจฺเฉทนา ปุจฺฉา.\csromanspacer{} อิมา ติสฺโส ปุจฺฉา.}

\prose{อปราปิ ติสฺโส ปุจฺฉา มนุสฺสปุจฺฉา อมนุสฺสปุจฺฉา นิมฺมิตปุจฺฉา.\csromanspacer{} กตมา มนุสฺส ปุจฺฉา, มนุสฺสา พุทฺธํ ภควนฺตํ อุปสงฺกมิตฺวา ปญฺหํ ปุจฺฉนฺติ, ภิกฺขู ปุจฺฉนฺติ, ภิกฺขุนิโย ปุจฺฉนฺติ, อุปาสกา ปุจฺฉนฺติ, อุปาสิกาโย ปุจฺฉนฺติ, ราชาโน}

\vspace{\tipitakaparskip}
\csromancenter{เมตฺตาคูมาณวปุจฺฉานิทฺเทส ปุจฺฉนฺติ, ขตฺติยา ปุจฺฉนฺติ, พฺราหฺมณา ปุจฺฉนฺติ, เวสฺสา ปุจฺฉนฺติ, สุทฺทา ปุจฺฉนฺติ, คหฏฺฐา ปุจฺฉนฺติ, ปพฺพชิตา ปุจฺฉนฺติ, อยํ มนุสฺสปุจฺฉา.}
\prose{\csromanfolio{63}กตมา อมนุสฺสปุจฺฉา, อมนุสฺสา พุทฺธํ ภควนฺตํ อุปสงฺกมิตฺวา ปญฺหํ ปุจฺฉนฺติ, นาคา ปุจฺฉนฺติ, สุปณฺณา ปุจฺฉนฺติ, ยกฺขา ปุจฺฉนฺติ, อสุรา ปุจฺฉนฺติ, คนฺธพฺพา ปุจฺฉนฺติ, มหาราชาโน ปุจฺฉนฺติ, อินฺทา ปุจฺฉนฺติ, พฺรหฺมา ปุจฺฉนฺติ, เทวา ปุจฺฉนฺติ, อยํ อมนุสฺสปุจฺฉา.}

\prose{กตมา นิมฺมิตปุจฺฉา, \textbf{ภควา} รูปํ อภินิมฺมินาติ มโนมยํ สพฺพงฺคปจฺจงฺคํ อหีนินฺทฺริยํ, โส นิมฺมิโต พุทฺธํ ภควนฺตํ อุปสงฺกมิตฺวา ปญฺหํ ปุจฺฉติ, \textbf{ภควา} วิสชฺเชติ, อยํ นิมฺมิตปุจฺฉา.\csromanspacer{} อิมา ติสฺโส ปุจฺฉา.}

\prose{อปราปิ ติสฺโส ปุจฺฉา อตฺตตฺถปุจฺฉา ปรตฺถปุจฺฉา อุภยตฺถปุจฺฉา.\csromanspacer{}\csromanspacer{}อปราปิ ติสฺโส ปุจฺฉา ทิฏฺฐธมฺมิกตฺถปุจฺฉา สมฺปรายิกตฺถปุจฺฉา ปรมตฺถปุจฺฉา.\csromanspacer{}\csromanspacer{}อปราปิ ติสฺโส ปุจฺฉา อนวชฺชตฺถปุจฺฉา นิกฺกิเลสตฺถปุจฺฉา โวทานตฺถปุจฺฉา.\csromanspacer{}\csromanspacer{}อปราปิ ติสฺโส ปุจฺฉา อตีตปุจฺฉา อนาคตปุจฺฉา ปจฺจุปฺปนฺนปุจฺฉา.\csromanspacer{}\csromanspacer{}อปราปิ ติสฺโส ปุจฺฉา อชฺฌตฺตปุจฺฉา พหิทฺธาปุจฺฉา อชฺฌตฺตพหิทฺธาปุจฺฉา.\csromanspacer{}\csromanspacer{}อปราปิ ติสฺโส ปุจฺฉา กุสลปุจฺฉา อกุสลปุจฺฉา อพฺยากตปุจฺฉา.\csromanspacer{}\csromanspacer{}อราปิ ติสฺโส ปุจฺฉา ขนฺธปุจฺฉา ธาตุปุจฺฉา อายตนปุจฺฉา.\csromanspacer{}\csromanspacer{}อปราปิ ติสฺโส ปุจฺฉา สติปฏฺฐานปุจฺฉา สมฺมปฺปธานปุจฺฉา อิทฺธิปาทปุจฺฉา.\csromanspacer{}\csromanspacer{}อปราปิ ติสฺโส ปุจฺฉา อินฺทฺริยปุจฺฉา พลปุจฺฉา โพชฺฌงฺคปุจฺฉา.\csromanspacer{}\csromanspacer{}อปราปิ ติสฺโส ปุจฺฉา มคฺคปุจฺฉา ผลปุจฺฉา นิพฺพานปุจฺฉา.}

\prose{\textbf{ปุจฺฉามิ ตนฺ}ติ ปุจฺฉามิ ตํ ยาจามิ ตํ อชฺเฌสามิ ตํ ปสาเทมิ ตํ “กถยสฺสุ เม”ติ ปุจฺฉามิ ตํ.\csromanspacer{} \textbf{ภควา}ติ คารวาธิวจนเมตํ \textbf{ฯเปฯ} สจฺฉิกา ปญฺญตฺติ, ยทิทํ \textbf{ภควา}ติ.\csromanspacer{} \textbf{พฺรูหิ เม ตนฺ}ติ พฺรูหิ อาจิกฺขาหิ เทเสหิ ปญฺญเปหิ ปฏฺฐเปหิ วิวราหิ วิภชาหิ อุตฺตานีกโรหิ ปกาเสหีติ ปุจฺฉามิ ตํ \textbf{ภควา} พฺรูหิ เมตํ.}

\prose{\textbf{อิจฺจายสฺมา เมตฺตคู}ติ อิจฺจาติ ปทสนฺธิ \textbf{ฯเปฯ} อิจฺจายสฺมา เมตฺตคู.}

\prose{\textbf{มญฺญามิ ตํ เวทคู ภาวิตตฺตนฺ}ติ \textbf{เวทคู}ติ ตํ มญฺญามิ, ภาวิตตฺโตติ ตํ มญฺญามิ เอวํ ชานามิ เอวํ อาชานามิ เอวํ ปฏิชานามิ เอวํ ปฏิวิชฺฌามิ.\csromanspacer{} \textbf{เวทคู ภาวิตตฺโต}ติ กถญฺจ \textbf{ภควา} \textbf{เวทคู}, \textbf{เวทา} วุจฺจนฺติ จตูสุ มคฺเคสุ ญาณํ ปญฺญา ปญฺญินฺทฺริยํ ปญฺญาพลํ \textbf{ฯเปฯ} ธมฺมวิจยสมฺโพชฺฌงฺโค \csromanfolio{64}วีมสา วิปสฺสนา สมฺมาทิฏฺฐิ.\csromanspacer{} ภควา เตหิ เวเทหิ ชาติชรามรณสฺส อนฺตคโต อนฺตปฺปตฺโต โกฏิคโต โกฏิปฺปตฺโต ปริยนฺตคโต ปริยนฺตปฺปตฺโต โวสานคโต โวสานปฺปตฺโต ตาณคโต ตาณปฺปตฺโต เลณคโต เลณปฺปตฺโต สรณคโต สรณปฺปตฺโต อภยคโต อภยปฺปตฺโต อจฺจุตคโต อจฺจุตปฺปตฺโต อมตคโต อมตปฺปตฺโต นิพฺพานคโต นิพฺพานปฺปตฺโต.\csromanspacer{} เวทานํ วา อนฺตคโตติ เวทคู.\csromanspacer{} เวเทหิ วา อนฺตคโตติ เวทคู.\csromanspacer{} สตฺตนฺนํ วา ธมฺมานํ วิทิตตฺตา เวทคู.\csromanspacer{} สกฺกายทิฏฺฐิ วิทิตา โหติ.\csromanspacer{} วิจิกิจฺฉา วิทิตา โหติ.\csromanspacer{} สีลพฺพตปรามาโส วิทิโต โหติ.\csromanspacer{} ราโค โทโส โมโห มาโน วิทิโต โหติ.\csromanspacer{} วิทิตาสฺส โหนฺติ ปาปกา อกุสลา ธมฺมา สํกิเลสิกา โปโนภวิกา สทรา ทุกฺขวิปากา อายติํ ชาติชรามรณิยา.}

\vspace{\tipitakaparskip}
\csromancenter{เวทานิ วิเจยฺย เกวลานิ, (สภิยาติ ภควา,)}
\prose{สมณานํ ยานีธตฺถิ\footnote{ยานิ ปตฺถิ (สฺยา), ยานิ อตฺถิ (ก) ขุ 1. 360 ปิฏฺเฐสุ.} พฺราหฺมณานํ.}

\csromangathameasure{สพฺพเวทนาสุ วีตราโค, \\ สพฺพํ เวทมติจฺจ เวทคู โสติ.}
\csromangathagroup{\csromangathastanza{สพฺพเวทนาสุ วีตราโค, \\ สพฺพํ เวทมติจฺจ เวทคู โสติ.}}

\prose{เอวํ ภควา เวทคู.}

\prose{กถํ ภควา ภาวิตตฺโต.\csromanspacer{} ภควา ภาวิตกาโย ภาวิตสีโล ภาวิตจิตฺโต ภาวิตปญฺโญ ภาวิตสติปฏฺฐาโน ภาวิตสมฺมปฺปธาโน ภาวิต-อิทฺธิปาโท ภาวิต-อินฺทฺริโย ภาวิตพโล ภาวิตโพชฺฌงฺโค ภาวิตมคฺโค ปหีนกิเลโส ปฏิวิทฺธากุปฺโป สจฺฉิกตนิโรโธ.\csromanspacer{} ทุกฺขํ ตสฺส ปริญฺญาตํ, สมุทโย ปหีโน, มคฺโค ภาวิโต, นิโรโธ สจฺฉิกโต.\csromanspacer{} อภิญฺเญยฺยํ อภิญฺญาตํ, ปริญฺเญยฺยํ ปริญฺญาตํ, ปหาตพฺพํ ปหีนํ, ภาเวตพฺพํ ภาวิตํ.\csromanspacer{} สจฺฉิกาตพฺพํ สจฺฉิกตํ.\csromanspacer{} อปริตฺโต มหนฺโต คมฺภีโร อปฺปเมยฺโย ทุปฺปริโยคาโฬฺห พหุรตโน สาครูปโม\footnote{สาครสโม (ก)} ฉฬงฺคุเปกฺขาย สมนฺนาคโต โหติ.}

\prose{จกฺขุนา รูปํ ทิสฺวา เนว สุมโน โหติ น ทุมฺมโน.\csromanspacer{} อุเปกฺขโก วิหรติ สโต สมฺปชาโน.\csromanspacer{} โสเตน สทฺทํ สุตฺวา.\csromanspacer{} ฆาเนน คนฺธํ ฆายิตฺวา.\csromanspacer{} ชิวฺหาย รสํ สายิตฺวา.\csromanspacer{} กาเยน โผฏฺฐพฺพํ ผุสิตฺวา.\csromanspacer{} มนสา}

\vspace{\tipitakaparskip}
\csromancenter{เมตฺตาคูมาณวปุจฺฉานิทฺเทส ธมฺมํ วิญฺญาย เนว สุมโน โหติ น ทุมฺมโน.\csromanspacer{} อุเปกฺขโก วิหรติ สโต สมฺปชาโน.}
\prose{\csromanfolio{65}จกฺขุนา รูปํ ทิสฺวา มนาปํ รูปํ นาภิคิชฺฌติ นาภิหํสติ\footnote{นาภิปิหยติ (สฺยา) ขุ 7. 186 ปิฏฺเฐ.} น ราคํ ชเนติ.\csromanspacer{} ตสฺส ฐิโตว กาโย โหติ, ฐิตํ จิตฺตํ อชฺฌตฺตํ สุสณฺฐิตํ สุวิมุตฺตํ.\csromanspacer{} จกฺขุนา โข ปเนว รูปํ ทิสฺวา อมนาปํ น มงฺกุ โหติ อปฺปติฏฺฐิตจิตฺโต\footnote{อปฺปติฏฺฐีนจิตฺโต (สฺยา)} อลีนมนโส\footnote{อาทินมนโส (สฺยา) ขุ 7. 186 ปิฏฺเฐปิ.} อพฺยาปนฺนเจตโส.\csromanspacer{} ตสฺส ฐิโตว กาโย โหติ, ฐิตํ จิตฺตํ อชฺฌตฺตํ สุสณฺฐิตํ สุวิมุตฺตํ.\csromanspacer{} โสเตน สทฺทํ สุตฺวา.\csromanspacer{} ฆาเนน คนฺธํ ฆายิตฺวา.\csromanspacer{} ชิวฺหาย รสํ สายิตฺวา.\csromanspacer{} กาเยน โผฏฺฐพฺพํ ผุสิตฺวา.\csromanspacer{} มนสา ธมฺมํ วิญฺญาย มนาปํ นาภิคิชฺฌติ นาภิหํสติ น ราคํ ชเนติ.\csromanspacer{} ตสฺส ฐิโตว กาโย โหติ, ฐิตํ จิตฺตํ อชฺฌตฺตํ สุสณฺฐิตํ สุวิมุตฺตํ.\csromanspacer{} มนสา เยว โข ปน ธมฺมํ วิญฺญาย อมนาปํ น มงฺกุ โหติ อปฺปติฏฺฐิตจิตฺโต อลีนมนโส อพฺยาปนฺนเจตโส.\csromanspacer{} ตสฺส ฐิโตว กาโย โหติ, ฐิตํ จิตฺตํ อชฺฌตฺตํ สุสณฺฐิตํ สุวิมุตฺตํ.}

\prose{จกฺขุนา รูปํ ทิสฺวา มนาปามนาเปสุ รูเปสุ ฐิโตว กาโย โหติ, ฐิตํ จิตฺตํ อชฺฌตฺตํ สุสณฺฐิตํ สุวิมุตฺตํ.\csromanspacer{} โสเตน สทฺทํ สุตฺวา ฯเปฯ.\csromanspacer{} มนสา ธมฺมํ วิญฺญาย มนาปามนาเปสุ ธมฺเมสุ ฐิโตว กาโย โหติ ฐิตํ จิตฺตํ อชฺฌตฺตํ สุสณฺฐิตํ สุวิมุตฺตํ.}

\prose{จกฺขุนา รูปํ ทิสฺวา รชนีเย น รชฺชติ.\csromanspacer{} ทุสฺสนีเย\footnote{โทสนีเย (สฺยา, ก) ขุ 7. 186 ปิฏฺเฐปิ.} น ทุสฺสติ.\csromanspacer{} โมหนีเย น มุยฺหติ.\csromanspacer{} โกปนีเย น กุปฺปติ.\csromanspacer{} มทนีเย น มชฺชติ.\csromanspacer{} กิเลสนีเย น กิลิสฺสติ.\csromanspacer{} โสเตน สทฺทํ สุตฺวา ฯเปฯ.\csromanspacer{} มนสา ธมฺมํ วิญฺญาย รชนีเย น รชฺชติ.\csromanspacer{} ทุสฺสนีเย น ทุสฺสติ.\csromanspacer{} โมหนีเย น มุยฺหติ.\csromanspacer{} โกปนีเย น กุปฺปติ.\csromanspacer{} มทนีเย น มชฺชติ.\csromanspacer{} กิเลสนีเย น กิลิสฺสติ.}

\prose{ทิฏฺเฐ ทิฏฺฐมตฺโต, สุเต สุตมตฺโต, มุเต มุตมตฺโต, วิญฺญาเต วิญฺญาตมตฺโต, ทิฏฺเฐ น ลิมฺปติ, สุเต น ลิมฺปติ, มุเต น ลิมฺปติ, วิญฺญาเต น ลิมฺปติ, ทิฏฺเฐ อนูปโย\footnote{อนุปโย (สฺยา), อนุสโย (ก) ขุ 7. 187 ปิฏฺเฐ.} อนปาโย อนิสฺสิโต \csromanfolio{66}อปฺปฏิพทฺโธ วิปฺปมุตฺโต วิสญฺญุตฺโต วิมริยาทิกเตน เจตสา วิหรติ.\csromanspacer{} สุเต.\csromanspacer{} มุเต.\csromanspacer{} วิญฺญาเต อนูปโย\footnote{อนุปโย (สฺยา), อนุสโย (ก) ขุ 7. 187 ปิฏฺเฐ.} อนปาโย อนิสฺสิโต อปฺปฏิพทฺโธ วิปฺปมุตฺโต วิสญฺญุตฺโต วิมริยาทิกเตน เจตสา วิหรติ.}

\prose{สํวิชฺชติ ภควโต จกฺขุ, ปสฺสติ ภควา จกฺขุนา รูปํ, ฉนฺทราโค ภควโต นตฺถิ, สุวิมุตฺตจิตฺโต ภควา.\csromanspacer{} สํวิชฺชติ ภควโต โสตํ, สุณาติ ภควา โสเตน สทฺทํ, ฉนฺทราโค ภควโต นตฺถิ, สุวิมุตฺตจิตฺโต ภควา.\csromanspacer{} สํวิชฺชติ ภควโต ฆานํ, ฆายติ ภควา ฆาเนน คนฺธํ, ฉนฺทราโค ภควโต นตฺถิ, สุวิมุตฺตจิตฺโต ภควา.\csromanspacer{} สํวิชฺชติ ภควโต ชิวฺหา, สายติ ภควา ชิวฺหาย รสํ, ฉนฺทราโค ภควโต นตฺถิ, สุวิมุตฺตจิตฺโต ภควา.\csromanspacer{} สํวิชฺชติ ภควโต กาโย, ผุสติ ภควา กาเยน โผฏฺฐพฺพํ, ฉนฺทราโค ภควโต นตฺถิ, สุวิมุตฺตจิตฺโต ภควา.\csromanspacer{} สํวิชฺชติ ภควโต มโน, วิชานาติ ภควา มนสา ธมฺมํ, ฉนฺทราโค ภควโต นตฺถิ, สุวิมุตฺตจิตฺโต ภควา.}

\prose{จกฺขุ รูปารามํ รูปรตํ รูปสมฺมุทิตํ, ตํ ภควโต\footnote{ภควตา (สฺยา) ขุ 7. 187 ปิฏฺเฐ.} ทนฺตํ คุตฺตํ รกฺขิตํ สํวุตํ, ตสฺส จ สํวราย ธมฺมํ เทเสติ.\csromanspacer{} โสตํ สทฺทารามํ สทฺทรตํ.\csromanspacer{} ฆานํ คนฺธารามํ คนฺธรตํ.\csromanspacer{} ชิวฺหา รสารามา รสรตา รสสมฺมุทิตา, สา ภควโต ทนฺตา คุตฺตา รกฺขิตา สํวุตา, ตสฺส จ สํวราย ธมฺมํ เทเสติ.\csromanspacer{} กาโย โผฏฺฐพฺพาราโม โผฏฺฐพฺพรโต โผฏฺฐพฺพสมฺมุทิโต.\csromanspacer{} มโน ธมฺมาราโม ธมฺมรโต ธมฺมสมฺมุทิโต, โส ภควโต ทนฺโต คุตฺโต รกฺขิโต สํวุโต, ตสฺส จ สํวราย ธมฺมํ เทเสติ.}

\csromangathameasure{“ทนฺตํ นยนฺติ สมิติํ, ทนฺตํ ราชาภิรูหติ. \\ ทนฺโต เสฏฺโฐ มนุสฺเสสุ, โย’ติวากฺยํ ติติกฺขติ. \\ วรมสฺสตรา ทนฺตา, อาชานียา จ สินฺธวา. \\ กุญฺชรา จ มหานาคา, อตฺตทนฺโต ตโต วรํ. \\ น หิ เอเตหิ ยาเนหิ, คจฺเฉยฺย อคตํ ทิสํ. \\ ยถา’ตฺตนา สุทนฺเตน, ทนฺโต ทนฺเตน คจฺฉติ.}
\csromangathagroup{\csromangathastanza{“ทนฺตํ นยนฺติ สมิติํ, ทนฺตํ ราชาภิรูหติ. \\ ทนฺโต เสฏฺโฐ มนุสฺเสสุ, โย’ติวากฺยํ ติติกฺขติ.}\csromangathastanza{วรมสฺสตรา ทนฺตา, อาชานียา จ\footnote{อาชานิยาว (สฺยา) ขุ 1. 59 ปิฏฺเฐ.} สินฺธวา. \\ กุญฺชรา จ\footnote{อาชานิยาว (สฺยา) ขุ 1. 59 ปิฏฺเฐ.} มหานาคา, อตฺตทนฺโต ตโต วรํ.}\csromangathastanza{น หิ เอเตหิ ยาเนหิ, คจฺเฉยฺย อคตํ ทิสํ. \\ ยถา’ตฺตนา สุทนฺเตน, ทนฺโต ทนฺเตน คจฺฉติ.}}

\csromanheader{2}{4. เมตฺตาคูมาณวปุจฺฉานิทฺเทส}
\csromangathameasure{วิธาสุ น วิกมฺปนฺติ, วิปฺปมุตฺตา ปุนพฺภวา. \\ ทนฺตภูมิํ อนุปฺปตฺตา, เต โลเก วิชิตาวิโน. \\ ยสฺสินฺทฺริยานิ ภาวิตานิ, \\ อชฺฌตฺตญฺจ พหิทฺธา จ สพฺพโลเก. \\ นิพฺพิชฺฌ อิมํ ปรญฺจ โลกํ, \\ กาลํ กงฺขติ ภาวิโต ส ทนฺโต”ติ.}
\csromangathagroup{\csromangathastanza{\csromanfolio{67}วิธาสุ น วิกมฺปนฺติ, วิปฺปมุตฺตา ปุนพฺภวา. \\ ทนฺตภูมิํ อนุปฺปตฺตา, เต โลเก วิชิตาวิโน.}\csromangathastanza{ยสฺสินฺทฺริยานิ ภาวิตานิ, \\ อชฺฌตฺตญฺจ พหิทฺธา จ สพฺพโลเก. \\ นิพฺพิชฺฌ อิมํ ปรญฺจ โลกํ, \\ กาลํ กงฺขติ ภาวิโต ส ทนฺโต”ติ\footnote{สุทนฺโตติ (สฺยา) ขุ 1. 358 ปิฏฺเฐ; ขุ 7. 188 ปิฏฺเฐปิ.}.}}

\prose{เอวํ ภควา ภาวิตตฺโตติ.}

\prose{\textbf{มญฺญามิ ตํ เวทคู ภาวิตตฺตํ.}\csromanspacer{}\textbf{ กุโต นุ ทุกฺขา สมุทาคตา อิเม}ติ \textbf{กุโต นู}ติ สํสยปุจฺฉา วิมติปุจฺฉา เทฺวฬฺหกปุจฺฉา อเนกํสปุจฺฉา “เอวํ นุ โข, นนุ โข, กิํ นุ โข, กถํ นุ โข”ติ กุโต นุ.\csromanspacer{} \textbf{ทุกฺขา}ติ ชาติทุกฺขํ ชราทุกฺขํ พฺยาธิทุกฺขํ มรณทุกฺขํ โสกปริเทวทุกฺขโทมนสฺสุปายาสทุกฺขํ,พฺยสนํ ทุกฺขํ เนรยิกํ ทุกฺขํ ติรจฺฉานโยนิกํ ทุกฺขํ เปตฺติวิสยิกํ ทุกฺขํ, มานุสิกํ ทุกฺขํ คพฺโภกฺกนฺติมูลกํ ทุกฺขํ คพฺภฏฺฐิติมูลกํ ทุกฺขํ คพฺภวุฏฺฐานมูลกํ ทุกฺขํ ชาตสฺสูปนิพนฺธกํ ทุกฺขํ ชาตสฺส ปราเธยฺยกํ ทุกฺขํ, อตฺตูปกฺกมํ ทุกฺขํ ปรูปกฺกมํ ทุกฺขํ ทุกฺขทุกฺขํ สงฺขารทุกฺขํ วิปริณามทุกฺขํ, จกฺขุโรโค โสตโรโค ฆานโรโค ชิวฺหาโรโค กายโรโค สีสโรโค กณฺณโรโค มุขโรโค ทนฺตโรโค กาโส สาโส ปินาโส ฑาโห ชโร กุจฺฉิโรโค มุจฺฉา ปกฺขนฺทิกา สูลา วิสูจิกา กุฏฺฐํ คณฺโฑ กิลาโส โสโส อปมาโร ททฺทุ กณฺฑุ กจฺฉุ รขสา วิตจฺฉิกา โลหิตปิตฺตํ มธุเมโห อํสา ปิฬกา ภคนฺทลา ปิตฺตสมุฏฺฐานา อาพาธา เสมฺหสมุฏฺฐานา อาพาธา วาตสมุฏฺฐานา อาพาธา สนฺนิปาติกา อาพาธา อุตุปริณามชา อาพาธา วิสมปริหารชา อาพาธา โอปกฺกมิกา อาพาธา กมฺมวิปากชา อาพาธา สีตํ อุณฺหํ ชิฆจฺฉา ปิปาสา อุจฺจาโร ปสฺสาโว ฑํสมกสวาตาตปสรีสป- สมฺผสฺสํ ทุกฺขํ, มาตุมรณํ ทุกฺขํ ปิตุมรณํ ทุกฺขํ ภาตุมรณํ ทุกฺขํ ภคินิมรณํ ทุกฺขํ ปุตฺตมรณํ ทุกฺขํ ธีตุมรณํ ทุกฺขํ ญาติพฺยสนํ ทุกฺขํ โรคพฺยสนํ ทุกฺขํ โภคพฺยสนํ ทุกฺขํ สีลพฺยสนํ ทุกฺขํ ทิฏฺฐิพฺยสนํ ทุกฺขํ, เยสํ ธมฺมานํ อาทิโต สมุทาคมนํ ปญฺญายติ, อตฺถงฺคมโต นิโรโธ ปญฺญายติ, กมฺมสนฺนิสฺสิโต วิปาโก, วิปากสนฺนิสฺสิตํ กมฺมํ, นามสนฺนิสฺสิตํ รูปํ,}

\prose{\csromanfolio{68}รูปสนฺนิสฺสิตํ นามํ, ชาติยา อนุคตํ, ชราย อนุสฏํ, พฺยาธินา อภิภูตํ, มรเณน อพฺภาหตํ, ทุกฺเข ปติฏฺฐิตํ, อตาณํ อเลณํ อสรณํ อสรณีภูตํ, อิเม วุจฺจนฺติ ทุกฺขา.\csromanspacer{} อิเม ทุกฺขา กุโต สมุทาคตา กุโต ชาตา กุโต สญฺชาตา กุโต นิพฺพตฺตา กุโต อภินิพฺพตฺตา กุโต ปาตุภูตา, กิํนิทานา กิํสมุทยา กิํชาติกา กิํปภวาติ อิเมสํ ทุกฺขานํ มูลํ ปุจฺฉติ, เหตุํ ปุจฺฉติ, นิทานํ ปุจฺฉติ, สมฺภวํ ปุจฺฉติ, ปภวํ ปุจฺฉติ, สมุฏฺฐานํ ปุจฺฉติ, อาหารํ ปุจฺฉติ, อารมฺมณํ ปุจฺฉติ, ปจฺจยํ ปุจฺฉติ, สมุทยํ ปุจฺฉติ ปปุจฺฉติ ยาจติ อชฺเฌสติ ปสาเทตีติ กุโต นุ ทุกฺขา สมุทาคตา อิเม.}

\prose{\textbf{เย เกจิ โลกสฺมิ’มเนกรูปา}ติ \textbf{เย เกจี}ติ สพฺเพน สพฺพํ สพฺพถา สพฺพํ อเสสํ นิสฺเสสํ ปริยาทิยนวจนเมตํ \textbf{เย เกจี}ติ.\csromanspacer{} \textbf{โลกสฺมินฺ}ติ อปายโลเก มนุสฺสโลเก เทวโลเก ขนฺธโลเก ธาตุโลเก อายตนโลเก.\csromanspacer{} \textbf{อเนกรูปา}ติ อเนกวิธา นานาปฺปการา ทุกฺขาติ เย เกจิ โลกสฺมิ’มเนกรูปา.\csromanspacer{} เตนาห โส พฺราหฺมโณ\csromandash{}}

\prose{ปุจฺฉามิ ตํ ภควา พฺรูหิ เม’ตํ, (อิจฺจายสฺมา \textbf{เมตฺตคู},)}

\prose{มญฺญามิ ตํ เวทคู ภาวิตตฺตํ.}

\prose{กุโต นุ ทุกฺขา สมุทาคตา อิเม,}

\prose{\textbf{เย เกจิ โลกสฺมิ’มเนกรูปา}ติ.}

\csromanhangingbegin
\hangingheaditem{19}{ทุกฺขสฺส เว มํ ปภวํ อปุจฺฉสิ, (\textbf{เมตฺตคู}ติ ภควา,)}
\hangingbody{\textbf{ตํ เต ปวกฺขามิ ยถา ปชานํ.}}
\hangingbody{\textbf{อุปธินิทานา ปภวนฺติ ทุกฺขา,}}
\hangingbody{\textbf{เย เกจิ โลกสฺมิ’มเนกรูปา}.}
\csromanhangingend

\prose{\textbf{ทุกฺขสฺส เว มํ ปภวํ อปุจฺฉสี}ติ \textbf{ทุกฺขสฺสา}ติ ชาติทุกฺขสฺส ชราทุกฺขสฺส พฺยาธิทุกฺขสฺส มรณทุกฺขสฺส โสกปริเทวทุกฺขโทมนสฺสุปายาสทุกฺขสฺส.\csromanspacer{} \textbf{ปภวํ อปุจฺฉสี}ติ ทุกฺขสฺส มูลํ ปุจฺฉสิ, เหตุํ ปุจฺฉสิ, นิทานํ ปุจฺฉสิ, สมฺภวํ ปุจฺฉสิ, ปภวํ ปุจฺฉสิ, สมุฏฺฐานํ ปุจฺฉสิ, อาหารํ ปุจฺฉสิ, อารมฺมณํ ปุจฺฉสิ, ปจฺจยํ ปุจฺฉสิ, สมุทยํ ปุจฺฉสิ ยาจสิ อชฺเฌสสิ ปสาเทสีติ ทุกฺขสฺส เว มํ ปภวํ อปุจฺฉสิ.\csromanspacer{} \textbf{เมตฺตคู}ติ ภควา ตํ พฺราหฺมณํ นาเมน อาลปติ.}

\vspace{\tipitakaparskip}
\csromancenter{เมตฺตาคูมาณวปุจฺฉานิทฺเทส \textbf{ภควา}ติ คารวาธิวจนเมตํ ฯเปฯ สจฺฉิกา ปญฺญตฺติ, ยทิทํ \textbf{ภควา}ติ เมตฺตคูติ \textbf{ภควา}.}
\prose{\csromanfolio{69}\textbf{ตํ เต ปวกฺขามิ ยถา ปชานนฺ}ติ \textbf{ตนฺ}ติ ทุกฺขสฺส มูลํ ปวกฺขามิ, เหตุํ ปวกฺขามิ, นิทานํ ปวกฺขามิ, สมฺภวํ ปวกฺขามิ, ปภวํ ปวกฺขามิ, สมุฏฺฐานํ ปวกฺขามิ, อาหารํ ปวกฺขามิ, อารมฺมณํ ปวกฺขามิ, ปจฺจยํ ปวกฺขามิ, สมุทยํ ปวกฺขามิ อาจิกฺขิสฺสามิ เทเสสฺสามิ ปญฺญเปสฺสามิ ปฏฺฐเปสฺสามิ วิวริสฺสามิ วิภชิสฺสามิ อุตฺตานีกริสฺสามิ ปกาเสสฺสามีติ ตํ เต ปวกฺขามิ.\csromanspacer{} \textbf{ยถา ปชานนฺ}ติ ยถา ปชานนฺโต อาชานนฺโต วิชานนฺโต ปฏิวิชานนฺโต ปฏิวิชฺฌนฺโต.\csromanspacer{} น อิติหีติหํ น อิติกิราย น ปรมฺปราย น ปิฏกสมฺปทาย\footnote{น ปิฏกสมฺปทาเนน (ก) ขุ 7. 280 ปิฏฺเฐ.} น ตกฺกเหตุ น นยเหตุ น อาการปริวิตกฺเกน น ทิฏฺฐินิชฺฌานกฺขนฺติยา สามํ สยมภิญฺญาตํ อตฺตปจฺจกฺขธมฺมํ, ตํ กถยิสฺสามีติ ตํ เต ปวกฺขามิ ยถา ปชานํ.}

\prose{\textbf{อุปธินิทานา ปภวนฺติ ทุกฺขา}ติ \textbf{อุปธี}ติ ทส \textbf{อุปธี}, ตณฺหูปธิ, ทิฏฺฐูปธิ, กิเลสูปธิ, กมฺมูปธิ, ทุจฺจริตูปธิ, อาหารูปธิ, ปฏิฆูปธิ, จตสฺโส อุปาทินฺนธาตุโย \textbf{อุปธี}, ฉ อชฺฌตฺติกานิ อาย\textbf{ตนฺ}อานิ \textbf{อุปธี}, ฉ วิญฺญาณกายา \textbf{อุปธี}, สพฺพมฺปิ ทุกฺขํ ทุกฺขมนฏฺเฐน\footnote{ทุกฺขฏฺเฐน (สฺยา)} อุปธิ, อิเม วุจฺจนฺติ ทส \textbf{อุปธี}.\csromanspacer{} \textbf{ทุกฺขา}ติ ชาติทุกฺขํ ชราทุกฺขํ พฺยาธิทุกฺขํ มรณทุกฺขํ โสกปริเทวทุกฺขโทมนสฺสุปายาส- ทุกฺขํ เนรยิกํ ทุกฺขํ ฯเปฯ ทิฏฺฐิพฺยสนํ ทุกฺขํ.\csromanspacer{} เยสํ ธมฺมานํ อาทิโต สมุทาคมนํ ปญฺญายติ, อตฺถงฺคมโต นิโรโธ ปญฺญายติ, กมฺมสนฺนิสฺสิโต วิปาโก, วิปากสนฺนิสฺสิตํ กมฺมํ, นามสนฺนิสฺสิตํ รูปํ.\csromanspacer{} รูปสนฺนิสฺสิตํ นามํ.\csromanspacer{} ชาติยา อนุคตํ ชราย อนุสฏํ, พฺยาธินา อภิภูตํ, มรเณน อพฺภาหตํ, ทุกฺเข ปติฏฺฐิตํ อตาณํ อเลณํ อสรณํ อสรณีภูตํ.\csromanspacer{} อิเม วุจฺจนฺติ ทุกฺขา.\csromanspacer{} อิเม ทุกฺขา อุปธินิทานา อุปธิเหตุกา อุปธิปจฺจยา อุปธิการณา โหนฺติ ปภวนฺติ สมฺภวนฺติ ชายนฺติ สญฺชายนฺติ นิพฺพตฺ\textbf{ตนฺ}ติ ปาตุภวนฺตีติ อุปธินิทานา ปภวนฺติ ทุกฺขา.}

\prose{\textbf{เย เกจิ โลกสฺมิ’มเนกรูปา}ติ \textbf{เย เกจี}ติ สพฺเพน สพฺพํ สพฺพถา สพฺพํ อเสสํ นิสฺเสสํ ปริยาทิยนวจนเมตํ \textbf{เย เกจี}ติ.\csromanspacer{} \textbf{โลกสฺมินฺ}ติ อปายโลเก มนุสฺสโลเก เทวโลเก ขนฺธโลเก \csromanfolio{70}ธาตุโลเก อายตนโลเก.\csromanspacer{} \textbf{อเนกรูปา}ติ อเนกวิธา นานปฺปการา ทุกฺขาติ เย เกจิ โลกสฺมิ’มเนกรูปา.\csromanspacer{} เตนาห ภควา\csromandash{}}

\prose{ทุกฺขสฺส เว มํ ปภวํ อปุจฺฉสิ, (เมตฺตคูติ ภควา,)}

\prose{ตํ เต ปวกฺขามิ ยถา ปชานํ.}

\prose{อุปธินิทานา ปภวนฺติ ทุกฺขา.}

\prose{เย เกจิ โลกสฺมิ’มเนกรูปาติ.}

\csromanhangingbegin
\hangingheaditem{20}{โย เว อวิทฺวา อุปธิํ กโรติ,}
\hangingbody{\textbf{ปุนปฺปุนํ ทุกฺข’มุเปติ มนฺโท.}}
\hangingbody{\textbf{ตสฺมา ปชานํ อุปธิํ น กยิรา,}}
\hangingbody{ทุกฺขสฺส ชาติปฺปภวานุปสฺสี.}
\csromanhangingend

\prose{\textbf{โย เว อวิทฺวา อุปธิํ กโรตี}ติ \textbf{โย}ติ \textbf{โย} ยาทิโส ยถายุตฺโต ยถาวิหิโต ยถาปกาโร ยํฐานปฺปตฺโต ยํธมฺมสมนฺนาคโต ขตฺติ\textbf{โย} วา พฺราหฺมโณ วา เวสฺโส วา สุทฺโท วา คหฏฺโฐ วา ปพฺพชิโต วา เทโว วา มนุสฺโส วา.\csromanspacer{} \textbf{อวิทฺวา}ติ อวิชฺชาคโต อญฺญาณี อวิภาวี ทุปฺปญฺโญ.\csromanspacer{} \textbf{อุปธิํ กโรตี}ติ ตณฺหูปธิํ กโรติ, ทิฏฺฐูปธิํ กโรติ, กิเลสูปธิํ กโรติ, กมฺมูปธิํ กโรติ, ทุจฺจริตูปธิํ กโรติ, อาหารูปธิํ กโรติ, ปฏิฆูปธิํ กโรติ.\csromanspacer{} จตสฺโส อุปาทินฺนธาตุ\textbf{โย} อุปธี กโรติ, ฉ อชฺฌตฺติกานิ อายตนานิ อุปธี กโรติ, ฉ วิญฺญาณกาเย อุปธี กโรติ ชเนติ สญฺชเนติ นิพฺพตฺเตติ อภินิพฺพตฺเตตีติ อวิทฺวา อุปธิํ กโรติ.}

\prose{\textbf{ปุนปฺปุนํ ทุกฺข’มุเปติ มนฺโท}ติ ปุนปฺปุนํ ชาติทุกฺขํ ชราทุกฺขํ พฺยาธิทุกฺขํ มรณทุกฺขํ โสกปริเทวทุกฺข- โทมนสฺสุปายาสทุกฺขํ เอติ สมุเปติ อุปคจฺฉติ คณฺหาติ ปรามสติ อภินิวิสตีติ ปุนปฺปุนํ ทุกฺข’มุเปติ.\csromanspacer{} \textbf{มนฺโท}ติ มนฺโท โมมุโห อวิทฺวา อวิชฺชาคโต อญฺญาณี อวิภาวี ทุปฺปญฺโญติ ปุนปฺปุนํ ทุกฺข’มุเปติ มนฺโท.}

\prose{\textbf{ตสฺมา ปชานํ อุปธิํ น กยิรา}ติ \textbf{ตสฺมา}ติ ตํการณา ตํเหตุ ตปฺปจฺจยา ตํนิทานา เอตํ อาทีนวํสมฺปสฺสมาโน อุปธีสูติ \textbf{ตสฺมา}.\csromanspacer{} \textbf{ปชานนฺ}ติ ปชานนฺโต อาชานนฺโต วิชานนฺโต ปฏิวิชานนฺโต ปฏิวิชฺฌนฺโต, “สพฺเพ สงฺขารา อนิจฺจา”ติ ปชานนฺโต อาชานนฺโต วิชานนฺโต ปฏิวิชานนฺโต}

\vspace{\tipitakaparskip}
\csromancenter{เมตฺตาคูมาณวปุจฺฉานิทฺเทส ปฏิวิชฺฌนฺโต, “สพฺเพ สงฺขารา ทุกฺขา”ติ ฯเปฯ “สพฺเพ ธมฺมา อนตฺตา”ติ ฯเปฯ. “ยํ กิญฺจิ สมุทยธมฺมํ, สพฺพํ ตํ นิโรธธมฺมนฺ”ติ ปชานนฺโต อาชานนฺโต วิชานนฺโต ปฏิวิชานนฺโต ปฏิวิชฺฌนฺโต.\csromanspacer{} \textbf{อุปธิํ น กยิรา}ติ ตณฺหูปธิํ น กเรยฺย, ทิฏฺฐูปธิํ น กเรยฺย, กิเลสูปธิํ น กเรยฺย, ทุจฺจริตูปธิํ น กเรยฺย, อาหารูปธิํ น กเรยฺย, ปฏิฆูปธิํ น กเรยฺย, จตสฺโส อุปาทินฺนธาตุโย อุปธี น กเรยฺย, ฉ อชฺฌตฺติกานิ อายตนานิ อุปธี น กเรยฺย, ฉ วิญฺญาณกาเย อุปธี น กเรยฺย น ชเนยฺย น สญฺชเนยฺย น นิพฺพตฺเตยฺย นาภินิพฺพตฺเตยฺยาติ ตสฺมา ปชานํ อุปธิํ น กยิรา.}
\prose{\csromanfolio{71}\textbf{ทุกฺขสฺสา}ติ ชาติทุกฺขสฺส ชราทุกฺขสฺส พฺยาธิทุกฺขสฺส มรณทุกฺขสฺส โสกปริเทวทุกฺขโทมนสฺสุปายาสทุกฺขสฺส.\csromanspacer{} \textbf{ปภวานุปสฺสี}ติ ทุกฺขสฺส มูลานุปสฺสี เหตานุปสฺสี นิทานานุปสฺสี สมฺภวานุปสฺสี ปภวานุปสฺสี สมุฏฺฐานานุปสฺสี อาหารานุปสฺสี อารมฺมณานุปสฺสี ปจฺจยานุปสฺสี สมุทยานุปสฺสี.\csromanspacer{} \textbf{อนุปสฺสนา} วุจฺจติ ญาณํ.\csromanspacer{} ยา ปญฺญา ปชานนา ฯเปฯ อโมโห ธมฺมวิจโย สมฺมาทิฏฺฐิ.\csromanspacer{} อิมาย อนุปสฺสนาย ปญฺญาย อุเปโต โหติ สมุเปโต อุปาคโต สมุปาคโต อุปปนฺโน สมุปปนฺโน สมนฺนาคโต, โส วุจฺจติ อนุปสฺสีติ ทุกฺขสฺส ชาติปฺปภวานุปสฺสี.\csromanspacer{} เตนาห ภควา\csromandash{}}

\csromangathameasure{โย เว อวิทฺวา อุปธิํ กโรติ, \\ ปุนปฺปุนํ ทุกฺข’มุเปติ มนฺโท. \\ ตสฺมา ปชานํ อุปธิํ น กยิรา, \\ ทุกฺขสฺส ชาติปฺปภวานุปสฺสีติ.}
\csromangathagroup{\csromangathastanza{โย เว อวิทฺวา อุปธิํ กโรติ, \\ ปุนปฺปุนํ ทุกฺข’มุเปติ มนฺโท. \\ ตสฺมา ปชานํ อุปธิํ น กยิรา, \\ ทุกฺขสฺส ชาติปฺปภวานุปสฺสีติ.}}

\proseitem{21}{\textbf{ยํ ตํ อปุจฺฉิมฺห อกิตฺตยี โน},}

\prose{\textbf{อญฺญํ ตํ ปุจฺฉาม ตทิงฺฆ พฺรูหิ.}}

\csromangathameasure{\textbf{กถํ นุ ธีรา วิตรนฺติ โอฆํ,} \\ \textbf{ชาติํ ชรํ โสกปริทฺทวญฺจ.} \\ \textbf{ตํ เม มุนี สาธุ วิยากโรหิ,}}
\csromangathagroup{\csromangathastanza{\textbf{กถํ นุ ธีรา วิตรนฺติ โอฆํ,} \\ \textbf{ชาติํ ชรํ โสกปริทฺทวญฺจ.} \\ \textbf{ตํ เม มุนี สาธุ วิยากโรหิ,}}}

\prose{\textbf{ยํ ตํ อปุจฺฉิมฺห อกิตฺตยี โน}ติ ยํ ตํ อปุจฺฉิมฺห อยาจิมฺห อชฺเฌสิมฺห ปสาทิมฺห.\csromanspacer{} \textbf{อกิตฺตยี โน}ติ กิตฺติตํ\footnote{อกิตฺติ ตํ (สฺยา) เอวมีทิเสสุ ปเทสุ อตีตวิภตฺติวเสน. ขุ 7. 217 ปิฏฺเฐ.} ปกิตฺติตํ อาจิกฺขิตํ เทสิตํ \csromanfolio{72}ปญฺญปิตํ\footnote{ปญฺญาปิตํ (ก)} ปฏฺฐปิตํ วิวริตํ วิภตฺตํ อุตฺตานีกตํ ปกาสิ\textbf{ตนฺ}ติ ยํ ตํ อปุจฺฉิมฺห อกิตฺตยี โน.}

\prose{\textbf{อญฺญํ ตํ ปุจฺฉาม ตทิงฺฆ พฺรูหี}ติ อญฺญํ ตํ ปุจฺฉาม, อญฺญํ ตํ ยาจาม, อญฺญํ ตํ อชฺเฌสาม, อญฺญํ ตํ ปสาเทม, อุตฺตริ ตํ ปุจฺฉาม.\csromanspacer{} \textbf{ตทิงฺฆ พฺรูหี}ติ อิงฺฆ พฺรูหิ อาจิกฺขาหิ เทเสหิ ปญฺญเปหิ ปฏฺฐเปหิ วิวราหิ วิภชาหิ อุตฺตานีกโรหิ ปกาเสหีติ อญฺญํ ตํ ปุจฺฉาม ตทิงฺฆ พฺรูหิ.}

\prose{\textbf{กถํ นุ ธีรา วิตรนฺติ โอฆํ, ชาติํ ชรํ โสกปริทฺทวญฺจา}ติ \textbf{กถํ นู}ติ สํสยปุจฺฉา วิมติปุจฺฉา เทฺวฬฺหกปุจฺฉา อเนกํสปุจฺฉา, “เอวํ นุ โข, นนุ โข, กิํ นุ โข, กถํ นุ โข”ติ กถํ นุ.\csromanspacer{} \textbf{ธีรา}ติ ธีรา ปณฺฑิตา ปญฺญวนฺโต พุทฺธิมนฺโต ญาณิโน วิภาวิโน เมธาวิโน.\csromanspacer{} \textbf{โอฆนฺ}ติ กาโมฆํ ภโวฆํ ทิฏฺโฐฆํ อวิชฺโชฆํ.\csromanspacer{} \textbf{ชาตี}ติ ยา เตสํ เตสํ สตฺตานํ ตมฺหิ ตมฺหิ สตฺตนิกาเย ชาติ สญฺชาติ โอกฺกนฺติ นิพฺพตฺติ อภินิพฺพตฺติ ขนฺธานํ ปาตุภาโว อายตนานํ ปฏิลาโภ.\csromanspacer{} \textbf{ชรา}ติ ยา เตสํ เตสํ สตฺตานํ ตมฺหิ ตมฺหิ สตฺตนิกาเย ชรา ชีรณตา ขณฺฑิจฺจํ ปาลิจฺจํ วลิตฺตจตา อายุโน สํหานิ อินฺทฺริยานํ ปริปาโก.\csromanspacer{} \textbf{โสโก}ติ ญาติพฺยสเนน วา ผุฏฺฐสฺส โภคพฺยสเนน วา ผุฏฺฐสฺส โรคพฺยสเนน วา ผุฏฺฐสฺส สีลพฺยสเนน วา ผุฏฺฐสฺส ทิฏฺฐิพฺยสเนน วา ผุฏฺฐสฺส อญฺญตรญฺญตเรน พฺยสเนน วา สมนฺนาคตสฺส อญฺญตรญฺญตเรน ทุกฺขธมฺเมน วา ผุฏฺฐสฺส โสโก โสจนา โสจิตตฺตํ อนฺโตโสโก อนฺโต ปริโสโก อนฺโตฑาโห อนฺโตปริฑาโห เจตโส ปริชฺฌายนา โทมนสฺสํ โสกสลฺลํ.\csromanspacer{} \textbf{ปริเทโว}ติ ญาติพฺยสเนน วา ผุฏฺฐสฺส โภคพฺยสเนน วา ผุฏฺฐสฺส โรคพฺยสเนน วา ผุฏฺฐสฺส สีลพฺยสเนน วา ผุฏฺฐสฺส ทิฏฺฐิพฺยสเนน วา ผุฏฺฐสฺส อญฺญตรญฺญตเรน พฺยสเนน วา สมนฺนาคตสฺส อญฺญตรญฺญตเรน ทุกฺขธมฺเมน วา ผุฏฺฐสฺส อาเทโว ปริเทโว อาเทวนา ปริเทวนา อาเทวิตตฺตํ ปริเทวิตตฺตํ วาจา ปลาโป\footnote{ลาโป ปลาโป (สฺยา) อภิ 1. 105 ปิฏฺเฐปิ.} วิปฺปลาโป ลาลปฺโป ลาลปฺปนา ลาลปฺปิตตฺตํ\footnote{ลาลปฺปายนา ลาลปฺปายิตตฺตํ (พหูสุ) ชราสุตฺตนิทฺเทสฏฺฐกถา โอโลเกตพฺพา.}.}

\prose{\textbf{กถํ นุ ธีรา วิตรนฺติ โอฆํ, ชาติํ ชรํ โสกปริทฺทวญฺจา}ติ ธีรา กถํ โอฆญฺจ ชาติญฺจ ชรญฺจ โสกญฺจ ปริเทวญฺจ ตรนฺติ อุตฺตรนฺติ ปตรนฺติ}

\vspace{\tipitakaparskip}
\csromancenter{เมตฺตาคูมาณวปุจฺฉานิทฺเทส สมติกฺกมนฺติ วีติวตฺ\textbf{ตนฺ}ตีติ กถํ นุ ธีรา วิตรนฺติ โอฆํ, ชาติํ ชรํ โสกปริทฺทวญฺจ.}
\prose{\csromanfolio{73}\textbf{ตํ เม มุนี สาธุ วิยากโรหี}ติ \textbf{ตนฺ}ติ ยํ ปุจฺฉามิ, ยํ ยาจามิ, ยํ อชฺเฌสามิ, ยํ ปสาเทมิ.\csromanspacer{} \textbf{มุนี}ติ โมนํ วุจฺจติ ญาณํ.\csromanspacer{} ยา ปญฺญา ปชานนา ฯเปฯ อโมโห ธมฺมวิจโย สมฺมาทิฏฺฐิ.\csromanspacer{} ภควา เตน ญาเณน สมนฺนาคโต มุนิ โมนปฺปตฺโต.\csromanspacer{} ตีณิ โมเนยฺยานิ กายโมเนยฺยํ วจีโมเนยฺยํ มโนโมเนยฺยํ.}

\prose{กตมํ กายโมเนยฺยํ, ติวิธานํ กายทุจฺจริตานํ ปหานํ กายโมเนยฺยํ, ติวิธํ กายสุจริตํ กายโมเนยฺยํ, กายารมฺมเณ ญาณํ กายโมเนยฺยํ, กายปริญฺญา กายโมเนยฺยํ, ปริญฺญาสหคโต มคฺโค กายโมเนยฺยํ, กาเย ฉนฺทราคสฺส ปหานํ กายโมเนยฺยํ, กายสงฺขารนิโรโธ จตุตฺถชฺฌานสมาปตฺติ กายโมเนยฺยํ.\csromanspacer{} อิทํ กายโมเนยฺยํ.}

\prose{กตมํ วจีโมเนยฺยํ, จตุพฺพิธานํ วจีทุจฺจริตานํ ปหานํ วจีโมเนยฺยํ, จตุพฺพิธํ วจีสุจริตํ วจีโมเนยฺยํ, วาจารมฺมเณ ญาณํ วจีโมเนยฺยํ, วาจาปริญฺญา วจีโมเนยฺยํ, ปริญฺญาสหคโต มคฺโค วจีโมเนยฺยํ, วาจาย ฉนฺทราคสฺส ปหานํ วจีโมเนยฺยํ, วจีสงฺขารนิโรโธ ทุติยชฺฌานสมาปตฺติ วจีโมเนยฺยํ.\csromanspacer{} อิทํ วจีโมเนยฺยํ.}

\prose{กตมํ มโนโมเนยฺยํ, ติวิธานํ มโนทุจฺจริตานํ ปหานํ มโนโมเนยฺยํ, ติวิธํ มโนสุจริตํ มโนโมเนยฺยํ, จิตฺตารมฺมเณ ญาณํ มโนโมเนยฺยํ, จิตฺตปริญฺญา มโนโมเนยฺยํ, ปริญฺญาสหคโต มคฺโค มโนโมเนยฺยํ, จิตฺเต ฉนฺทราคสฺส ปหานํ มโนโมเนยฺยํ, จิตฺตสงฺขารนิโรโธ สญฺญาเวทยิ\textbf{ตนฺ}อิโรธสมาปตฺติ มโนโมเนยฺยํ.\csromanspacer{} อิทํ มโนโมเนยฺยํ.}

\csromangathameasure{กายมุนิํ วจีมุนิํ, มโนมุนิมนาสวํ. \\ มุนิํ โมเนยฺยสมฺปนฺนํ, อาหุ สพฺพปฺปหายินํ. \\ กายมุนิํ วจีมุนิํ, มโนมุนิมนาสวํ. \\ มุนิํ โมเนยฺยสมฺปนฺนํ, อาหุ นินฺหาตปาปกนฺติ.}
\csromangathagroup{\csromangathastanza{กายมุนิํ วจีมุนิํ\footnote{วาจามุนิํ (พหูสุ) ขุ 1. 234 ปิฏฺเฐ.}, มโนมุนิมนาสวํ. \\ มุนิํ โมเนยฺยสมฺปนฺนํ, อาหุ สพฺพปฺปหายินํ.}\csromangathastanza{กายมุนิํ วจีมุนิํ, มโนมุนิมนาสวํ. \\ มุนิํ โมเนยฺยสมฺปนฺนํ, อาหุ นินฺหาตปาปกนฺติ.}}

\prose{\csromanfolio{74}อิเมหิ ตีหิ โมเนยฺเยหิ ธมฺเมหิ สมนฺนาคตา, ฉ มุนิโน\footnote{มุนโย (สฺยา) ขุ 7. 44 ปิฏฺเฐปิ.} อคารมุนิโน อนคารมุนิโน เสขมุนิโน\footnote{เสกฺขมุนิโน (สฺยา, ก)} อเสขมุนิโน ปจฺเจกมุนิโน มุนิมุนิโนติ.\csromanspacer{} กตเม อคารมุนิโน, เย เต อคาริกา ทิฏฺฐปทา วิญฺญาตสาสนา, อิเม อคารมุนิโน.\csromanspacer{}\csromanspacer{}กตเม อนคารมุนิโน, เย เต ปพฺพชิตา ทิฏฺฐปทา วิญฺญาตสาสนา, อิเม อนคารมุนิโน.\csromanspacer{} สตฺต เสขา เสขมุนิโน.\csromanspacer{} อรหนฺโต อเสขมุนิโน.\csromanspacer{} ปจฺเจกสมฺพุทฺธา ปจฺเจกมุนิโน.\csromanspacer{} ตถาคตา อรหนฺโต สมฺมาสมฺพุทฺธา มุนิมุนิโน.}

\csromangathameasure{น โมเนน มุนี โหติ, มูฬฺหรูโป อวิทฺทสุ. \\ โย จ ตุลํว ปคฺคยฺห, วรมาทาย ปณฺฑิโต. \\ ปาปานิ ปริวชฺเชติ, ส มุนี เตน โส มุนิ. \\ โย มุนาติ อุโภ โลเก, “มุนิ” เตน ปวุจฺจติ. \\ อสตญฺจ สตญฺจ ญตฺวา ธมฺมํ, \\ อชฺฌตฺตํ พหิทฺธา จ สพฺพโลเก. \\ เทวมนุสฺเสหิ ปูชนีโย, \\ สงฺคชาลมติจฺจ โส มุนีติ.}
\csromangathagroup{\csromangathastanza{น โมเนน มุนี\footnote{มุนิ (สฺยา, ก) ขุ 1. 51 ปิฏฺเฐ.} โหติ, มูฬฺหรูโป อวิทฺทสุ. \\ โย จ ตุลํว ปคฺคยฺห, วรมาทาย ปณฺฑิโต.}\csromangathastanza{ปาปานิ ปริวชฺเชติ, ส มุนี เตน โส มุนิ. \\ โย มุนาติ อุโภ โลเก, “มุนิ” เตน ปวุจฺจติ.}\csromangathastanza{อสตญฺจ สตญฺจ ญตฺวา ธมฺมํ, \\ อชฺฌตฺตํ พหิทฺธา จ สพฺพโลเก. \\ เทวมนุสฺเสหิ ปูชนีโย\footnote{ปูชิโต (สฺยา, ก) ขุ 7. 45 ปิฏฺเฐปิ.}, \\ สงฺคชาลมติจฺจ\footnote{ปูชิโต (สฺยา, ก) ขุ 7. 45 ปิฏฺเฐปิ.} โส มุนีติ.}}

\prose{\textbf{สาธุ วิยากโรหี}ติ ตํ สาธุ อาจิกฺขาหิ เทเสหิ ปญฺญเปหิ ปฏฺฐเปหิ วิวราหิ วิภชาหิ อุตฺตานีกโรหิ ปกาเสหีติ ตํ เม มุนี สาธุ วิยากโรหิ.\csromanspacer{} \textbf{ตถา หิ เต วิทิโต เอส ธมฺโม}ติ ตถา หิ เต วิทิโต ตุลิโต ตีริโต วิภูโต วิภาวิโต เอส ธมฺโมติ ตถา หิ เต วิทิโต เอส ธมฺโม.\csromanspacer{} เตนาห โส พฺราหฺมโณ\csromandash{}}

\csromangathameasure{ยํ ตํ อปุจฺฉิมฺห อกิตฺตยี โน, \\ อญฺญํ ตํ ปุจฺฉาม ตทิงฺฆ พฺรูหิ. \\ กถํ นุ ธีรา วิตรนฺติ โอฆํ, \\ ชาติํ ชรํ โสกปริทฺทวญฺจ. \\ ตํ เม มุนี สาธุ วิยากโรหิ,}
\csromangathagroup{\csromangathastanza{ยํ ตํ อปุจฺฉิมฺห อกิตฺตยี โน, \\ อญฺญํ ตํ ปุจฺฉาม ตทิงฺฆ พฺรูหิ. \\ กถํ นุ ธีรา วิตรนฺติ โอฆํ, \\ ชาติํ ชรํ โสกปริทฺทวญฺจ.}\csromangathastanza{ตํ เม มุนี สาธุ วิยากโรหิ,}}

\prosecont{\textbf{ตถา หิ เต วิทิโต เอส ธมฺโม}ติ.}

\csromanheader{2}{4. เมตฺตาคูมาณวปุจฺฉานิทฺเทส}
\csromanhangingbegin
\hangingheaditem{22}{\csromanfolio{75}กิตฺตยิสฺสามิ เต ธมฺมํ, (\textbf{เมตฺตคู}ติ ภควา,)}
\hangingbody{\textbf{ทิฏฺเฐ ธมฺเม อนีติหํ.}}
\hangingbody{ยํ วิทิตฺวา สโต จรํ, ตเร โลเก วิสตฺติกํ.}
\csromanhangingend

\prose{\textbf{กิตฺตยิสฺสามิ เต ธมฺมนฺ}ติ \textbf{ธมฺมนฺ}ติ อาทิกลฺยาณํ มชฺเฌกลฺยาณํ ปริโยสานกลฺยาณํ สาตฺถํ สพฺยญฺชนํ เกวลปริปุณฺณํ ปริสุทฺธํ พฺรหฺมจริยํ จตฺตาโร สติปฏฺฐาเน จตฺตาโร สมฺมปฺปธาเน จตฺตาโร อิทฺธิปาเท ปญฺจินฺทฺริยานิ ปญฺจ พลานิ สตฺต โพชฺฌงฺเค อริยํ อฏฺฐงฺคิกํ มคฺคํ นิพฺพานญฺจ นิพฺพานคามินิญฺจ ปฏิปทํ กิตฺตยิสฺสามิ อาจิกฺขิสฺสามิ เทเสสฺสามิ ปญฺญเปสฺสามิ ปฏฺฐเปสฺสามิ วิวริสฺสามิ วิภชิสฺสามิ อุตฺตานีกริสฺสามิ ปกาสิสฺสามีติ กิตฺตยิสฺสามิ เต ธมฺมํ.\csromanspacer{} \textbf{เมตฺตคู}ติ ภควา ตํ พฺราหฺมณํ นาเมน อาลปติ.}

\prose{\textbf{ทิฏฺเฐ ธมฺเม อนีติหนฺ}ติ \textbf{ทิฏฺเฐ ธมฺเม}ติ \textbf{ทิฏฺเฐ ธมฺเม} ญาเต ธมฺเม ตุลิเต ธมฺเม ตีริเต ธมฺเม วิภูเต ธมฺเม วิภาวิเต ธมฺเม “สพฺเพ สงฺขารา อนิจฺจา”ติ ฯเปฯ. “ยํ กิญฺจิ สมุทยธมฺมํ, สพฺพํ ตํ นิโรธ\textbf{ธมฺมนฺ}”ติ \textbf{ทิฏฺเฐ ธมฺเม} ญาเต ธมฺเม ตุลิเต ธมฺเม ตีริเต ธมฺเม วิภูเต ธมฺเม วิภาวิเต ธมฺเมติ เอวมฺปิ \textbf{ทิฏฺเฐ ธมฺเม} กถยิสฺสามิ.\csromanspacer{}\csromanspacer{}อถ วา ทุกฺเข ทิฏฺเฐ ทุกฺขํ กถยิสฺสามิ, สมุทเย ทิฏฺเฐ สมุทยํ กถยิสฺสามิ, มคฺเค ทิฏฺเฐ มคฺคํ กถยิสฺสามิ, นิโรเธ ทิฏฺเฐ นิโรธํ กถยิสฺสามีติ เอวมฺปิ \textbf{ทิฏฺเฐ ธมฺเม} กถยิสฺสามิ.\csromanspacer{}\csromanspacer{}อถ วา \textbf{ทิฏฺเฐ ธมฺเม} สนฺทิฏฺฐิกํ อกาลิกํ เอหิปสฺสิกํ โอปเนยฺยิกํ ปจฺจตฺตํ เวทิตพฺพํ วิญฺญูหีติ เอวมฺปิ \textbf{ทิฏฺเฐ ธมฺเม} กถยิสฺสามีติ \textbf{ทิฏฺเฐ ธมฺเม}.\csromanspacer{} \textbf{อนีติหนฺ}ติ น อิติหีติหํ น อิติกิราย น ปรมฺปราย น ปิฏกสมฺปทาย น ตกฺกเหตุ น นยเหตุ น อาการปริวิตกฺเกน น ทิฏฺฐินิชฺฌานกฺขนฺติยา สามํ สยมภิญฺญาตํ อตฺตปจฺจกฺขธมฺมํ, ตํ กถยิสฺสามีติ \textbf{ทิฏฺเฐ ธมฺเม} อนีติหํ.}

\prose{\textbf{ยํ วิทิตฺวา สโต จรนฺ}ติ ยํ วิทิตํ กตฺวา ตุลยิตฺวา ตีรยิตฺวา วิภาวยิตฺวา วิภูตํ กตฺวา, “สพฺเพ สงฺขารา อนิจฺจา”ติ วิทิตํ กตฺวา ตุลยิตฺวา ตีรยิตฺวา วิภาวยิตฺวา วิภูตํ กตฺวา, “สพฺเพ สงฺขารา ทุกฺขา”ติ ฯเปฯ. “สพฺเพ ธมฺมา อนตฺตา”ติ. “ยํ กิญฺจิ สมุทยธมฺมํ, สพฺพํ ตํ นิโรธ\textbf{ธมฺมนฺ}”ติ วิทิตํ กตฺวา ตุลยิตฺวา ตีรยิตฺวา วิภาวยิตฺวา วิภูตํ กตฺวา.\csromanspacer{} \textbf{สโต}ติ จตูหิ การเณหิ สโต, กาเย กายานุปสฺสนาสติปฏฺฐานํ ภาเวนฺโต สโต ฯเปฯ โส วุจฺจติ สโต. \csromanfolio{76}\textbf{จรนฺ}ติ จรนฺโต วิหรนฺโต อิริยนฺโต วตฺเตนฺโต ปาเลนฺโต ยเปนฺโต ยาเปนฺโตติ ยํ วิทิตฺวา สโต จรํ.}

\prose{\textbf{ตเร โลเก วิสตฺติกนฺ}ติ วิสตฺติกา วุจฺจติ ตณฺหา.\csromanspacer{} โย ราโค สาราโค ฯเปฯ อภิชฺฌา โลโภ อกุสลมูลํ.\csromanspacer{} \textbf{วิสตฺติกา}ติ เกนฏฺเฐน วิสตฺติกา, วิสตาติ วิสตฺติกา, วิสาลาติ วิสตฺติกา, วิสฏาติ วิสตฺติกา, วิสมาติ วิสตฺติกา, วิสกฺกตีติ วิสตฺติกา, วิสํหรตีติ วิสตฺติกา, วิสํวาทิกาติ วิสตฺติกา, วิสมูลาติ วิสตฺติกา, วิสผลาติ วิสตฺติกา, วิสปริโภคาติ วิสตฺติกา, วิสาลา วา ปน สา ตณฺหา รูเป สทฺเท คนฺเธ รเส โผฏฺฐพฺเพ กุเล คเณ อาวาเส ลาเภ ยเส ปสํสาย สุเข จีวเร ปิณฺฑปาเต เสนาสเน คิลานปจฺจยเภสชฺชปริกฺขาเร กามธาตุยา รูปธาตุยา อรูปธาตุยา กามภเว รูปภเว อรูปภเว สญฺญาภเว อสญฺญาภเว เนวสญฺญานาสญฺญาภเว เอกโวการภเว จตุโวการภเว ปญฺจโวการภเว อตีเต อนาคเต ปจฺจุปฺปนฺเน ทิฏฺฐสุตมุตวิญฺญาตพฺเพสุ ธมฺเมสุ วิสฏา วิตฺถตาติ วิสตฺติกา.\csromanspacer{} \textbf{โลเก}ติ อปายโลเก มนุสฺสโลเก เทวโลเก ขนฺธโลเก ธาตุโลเก อาย\textbf{ตนฺ}อโลเก.\csromanspacer{} \textbf{ตเร โลเก วิสตฺติกนฺ}ติ โลเก เวสา วิสตฺติกา\footnote{ยา สา โลเก วิสตฺติกา (สฺยา) กามสุตฺตนิทฺเทสฏฺฐกถา โอโลเกตพฺพา.}, โลเก เวตํ วิสตฺติกํ สโต ตเรยฺย อุตฺตเรยฺย ปตเรยฺย สมติกฺกเมยฺย วีติวตฺเตยฺยาติ ตเร โลเก วิสตฺติกํ.\csromanspacer{} เตนาห ภควา\csromandash{}}

\prose{กิตฺตยิสฺสามิ เต ธมฺมํ, (เมตฺตคูติ ภควา,)}

\prose{ทิฏฺเฐ ธมฺเม อนีติหํ.}

\csromangathameasure{ยํ วิทิตฺวา สโต จรํ, ตเร โลเก วิสตฺติกนฺติ.}
\csromangathagroup{\csromangathastanza{ยํ วิทิตฺวา สโต จรํ, ตเร โลเก วิสตฺติกนฺติ.}}

\csromanhangingbegin
\hangingheaditem{23}{ตญฺจาหํ อภินนฺทามิ, มเหสิ ธมฺม’มุตฺตมํ.}
\hangingbody{ยํ วิทิตฺวา สโต จรํ, ตเร โลเก วิสตฺติกํ.}
\csromanhangingend

\prose{\textbf{ตญฺจาหํ อภินนฺทามี}ติ \textbf{ตนฺ}ติ ตุยฺหํ วจนํ พฺยปฺปถํ\footnote{พฺยปถํ (สฺยา, ก)} เทสนํ อนุสาสนํ อนุสิฏฺฐํ\footnote{เทสนํ อนุสนฺธิํ (สฺยา)}.\csromanspacer{} \textbf{นนฺทามี}ติ อภินนฺทามิ โมทามิ อนุโมทามิ อิจฺฉามิ สาทิยามิ ยาจามิ ปตฺถยามิ ปิหยามิ อภิชปฺปามีติ ตญฺจาหํ อภินนฺทามิ.}

\csromanheader{2}{4. เมตฺตาคูมาณวปุจฺฉานิทฺเทส}
\prose{\csromanfolio{77}\textbf{มเหสิ ธมฺม’มุตฺตมนฺ}ติ \textbf{มเหสี}ติ กิํ มเหสิ ภควา, มหนฺตํ สีลกฺขนฺธํ เอสี คเวสี\footnote{เอสิ คเวสิ (สฺยา) ขุ 7. 266 ปิฏฺเฐ.} ปริเยสีติ มเหสิ, มหนฺตํ สมาธิกฺขนฺธํ.\csromanspacer{} มหนฺตํ ปญฺญากฺขนฺธํ.\csromanspacer{} มหนฺตํ วิมุตฺติกฺขนฺธํ.\csromanspacer{} มหนฺตํ วิมุตฺติญาณทสฺสนกฺขนฺธํ เอสี คเวสี ปริเยสีติ มเหสิ.\csromanspacer{} มหโต ตโมกายสฺส ปทาลนํ เอสี คเวสี ปริเยสีติ มเหสิ, มหโต วิปลฺลาสสฺส ปเภทนํ เอสี คเวสี ปริเยสีติ มเหสิ, มหโต ตณฺหาสลฺลสฺส อพฺพหนํ\footnote{อพฺพูหนํ (พหูสุ) อพฺพูหํ (สี-ฏฺฐ)} เอสี คเวสี ปริเยสีติ มเหสิ, มหโต ทิฏฺฐิสํฆาฏสฺส วินิเวฐนํ เอสี คเวสี ปริเยสีติ มเหสิ, มหโต มานธชสฺส ปาตนํ เอสี คเวสี ปริเยสีติ มเหสิ, มหโต อภิสงฺขารสฺส วูปสมํ เอสี คเวสี ปริเยสีติ มเหสิ, มหโต โอฆสฺส นิตฺถรณํ เอสี คเวสี ปริเยสีติ มเหสิ, มหโต ภารสฺส นิกฺเขปนํ เอสี คเวสี ปริเยสีติ มเหสิ, มหโต สํสารวฏฺฏสฺส อุปจฺเฉทํ เอสี คเวสี ปริเยสีติ มเหสิ, มหโต สนฺตาปสฺส นิพฺพาปนํ เอสี คเวสี ปริเยสีติ มเหสิ, มหโต ปริฬาหสฺส ปฏิปฺปสฺสทฺธํ เอสี คเวสี ปริเยสีติ มเหสิ, มหโต ธมฺมธชสฺส อุสฺสาปนํ เอสี คเวสี ปริเยสีติ มเหสิ, มหนฺเต สติปฏฺฐาเน.\csromanspacer{} มหนฺเต สมฺมปฺปธาเน.\csromanspacer{} มหนฺเต อิทฺธิปาเท.\csromanspacer{} มหนฺตานิ อินฺทฺริยานิ.\csromanspacer{} มหนฺตานิ พลานิ.\csromanspacer{} มหนฺเต โพชฺฌงฺเค.\csromanspacer{} มหนฺตํ อริยํ อฏฺฐงฺคิกํ มคฺคํ.\csromanspacer{} มหนฺตํ ปรมตฺถํ อมตํ นิพฺพานํ เอสี คเวสี ปริเยสีติ มเหสิ, มเหสกฺเขหิ สตฺเตหิ เอสิโต คเวสิโต ปริเยสิโต “กหํ พุทฺโธ, กหํ ภควา, กหํ เทวเทโว, กหํ นราสโภ”ติ มเหสิ.\csromanspacer{} \textbf{ธมฺม’มุตฺตมนฺ}ติ ธมฺมมุตฺตมํ วุจฺจติ อมตํ นิพฺพานํ.\csromanspacer{} โย โส สพฺพสงฺขารสมโถ สพฺพูปธิปฏินิสฺสคฺโค ตณฺหกฺขโย วิราโค นิโรโธ นิพฺพานํ.\csromanspacer{} \textbf{อุตฺตมนฺ}ติ อคฺคํ เสฏฺฐํ วิเสฏฺฐํ ปาโมกฺขํ อุตฺตมํ ปวรํ ธมฺมนฺติ มเหสิ ธมฺม’มุตฺตมํ.}

\prose{\textbf{ยํ วิทิตฺวา สโต จรนฺ}ติ วิทิตํ กตฺวา ตุลยิตฺวา ตีรยิตฺวา วิภาวยิตฺวา วิภูตํ กตฺวา, “สพฺเพ สงฺขารา อนิจฺจา”ติ วิทิตํ กตฺวา ตุลยิตฺวา ตีรยิตฺวา วิภาวยิตฺวา วิภูตํ กตฺวา, “สพฺเพ สงฺขารา ทุกฺขา”ติ ฯเปฯ. “สพฺเพ ธมฺมา อนตฺตา”ติ. “ยํ กิญฺจิ สมุทยธมฺมํ, สพฺพํ ตํ นิโรธธมฺมนฺ”ติ วิทิตํ กตฺวา ตุลยิตฺวา ตีรยิตฺวา วิภาวยิตฺวา วิภูตํ กตฺวา.\csromanspacer{} \textbf{สโต}ติ จตูหิ การเณหิ สโต, กาเย}

\prose{\csromanfolio{78}กายานุปสฺสนาสติปฏฺฐานํ ภาเวนฺโต สโต.\csromanspacer{} เวทนาสุ.\csromanspacer{} จิตฺเต.\csromanspacer{} ธมฺเมสุ ธมฺมานุปสฺสนาสติปฏฺฐานํ ภาเวนฺโต สโต ฯเปฯ โส วุจฺจติ สโต.\csromanspacer{} \textbf{จรนฺ}ติ จรนฺโต วิหรนฺโต อิริยนฺโต วตฺเตนฺโต ปาเลนฺโต ยเปนฺโต ยาเปนฺโตติ ยํ วิทิตฺวา สโต จรํ.}

\prose{\textbf{ตเร โลเก วิสตฺติกนฺ}ติ วิสตฺติกา วุจฺจติ ตณฺหา.\csromanspacer{} โย ราโค สาราโค ฯเปฯ อภิชฺฌา โลโภ อกุสลมูลํ.\csromanspacer{} \textbf{วิสตฺติกา}ติ เกนฏฺเฐน วิสตฺติกา ฯเปฯ วิสฏา วิตฺถตาติ วิสตฺติกา.\csromanspacer{} \textbf{โลเก}ติ อปายโลเก ฯเปฯ อายตนโลเก.\csromanspacer{} \textbf{ตเร โลเก วิสตฺติกนฺ}ติ โลเก เวสา วิสตฺติกา, โลเก เวตํ วิสตฺติกํ สโต ตเรยฺย อุตฺตเรยฺย ปตเรยฺย สมติกฺกเมยฺย วีติวตฺเตยฺยาติ ตเร โลเก วิสตฺติกํ.\csromanspacer{} เตนาห โส พฺราหฺมโณ\csromandash{}}

\csromangathameasure{ตญฺจาหํ อภินนฺทามิ, มเหสิ ธมฺม’มุตฺตมํ. \\ ยํ วิทิตฺวา สโต จรํ, ตเร โลเก วิสตฺติกนฺติ.}
\csromangathagroup{\csromangathastanza{ตญฺจาหํ อภินนฺทามิ, มเหสิ ธมฺม’มุตฺตมํ. \\ ยํ วิทิตฺวา สโต จรํ, ตเร โลเก วิสตฺติกนฺติ.}}

\csromanhangingbegin
\hangingheaditem{24}{ยํ กิญฺจิ สมฺปชานาสิ, (\textbf{เมตฺตคู}ติ \textbf{ภควา},)}
\hangingbody{\textbf{อุทฺธํ อโธ ติริยญฺจาปิ มชฺเฌ.}}
\hangingbody{\textbf{เอเตสุ นนฺทิญฺจ นิเวสนญฺจ,}}
\hangingbody{ปนุชฺช วิญฺญาณํ ภเว น ติฏฺเฐ.}
\csromanhangingend

\prose{\textbf{ยํ กิญฺจิ สมฺปชานาสี}ติ ยํ กิญฺจิ ปชานาสิ อาชานาสิ วิชานาสิ ปฏิวิชานาสิ ปฏิวิชฺฌสีติ ยํ กิญฺจิ สมฺปชานาสิ.\csromanspacer{} \textbf{เมตฺตคู}ติ \textbf{ภควา} ตํ พฺราหฺมณํ นาเมน อาลปติ.\csromanspacer{} \textbf{ภควา}ติ คารวาธิวจนเมตํ ฯเปฯ สจฺฉิกา ปญฺญตฺติ, ยทิทํ \textbf{ภควา}ติ \textbf{เมตฺตคู}ติ \textbf{ภควา}.}

\prose{\textbf{อุทฺธํ อโธ ติริยญฺจาปิ มชฺเฌ}ติ \textbf{อุทฺธนฺ}ติ อนาคตํ\footnote{อุทฺธํ วุจฺจติ อนาคตํ (สฺยา, ก)}.\csromanspacer{} \textbf{อโธ}ติ อตีตํ.\csromanspacer{} \textbf{ติริยญฺจาปิ มชฺเฌ}ติ ปจฺจุปฺปนฺนํ.\csromanspacer{} \textbf{อุทฺธนฺ}ติ เทวโลโก.\csromanspacer{} \textbf{อโธ}ติ นิรยโลโก.\csromanspacer{} \textbf{ติริยญฺจาปิ มชฺเฌ}ติ มนุสฺสโลโก.\csromanspacer{} อถ วา \textbf{อุทฺธนฺ}ติ กุสลา ธมฺมา.\csromanspacer{} \textbf{อโธ}ติ อกุสลา ธมฺมา.\csromanspacer{} \textbf{ติริยญฺจาปิ มชฺเฌ}ติ อพฺยากตา ธมฺมา.\csromanspacer{} \textbf{อุทฺธนฺ}ติ อรูปธาตุ.\csromanspacer{} \textbf{อโธ}ติ กามธาตุ.\csromanspacer{} \textbf{ติริยญฺจาปิ มชฺเฌ}ติ รูปธาตุ.\csromanspacer{} \textbf{อุทฺธนฺ}ติ สุขา เวทนา.\csromanspacer{} \textbf{อโธ}ติ ทุกฺขา เวทนา.\csromanspacer{} \textbf{ติริยญฺจาปิ มชฺเฌ}ติ อทุกฺขมสุขา เวทนา.\csromanspacer{} \textbf{อุทฺธนฺ}ติ อุทฺธํ ปาทตลา.}

\vspace{\tipitakaparskip}
\csromancenter{เมตฺตาคูมาณวปุจฺฉานิทฺเทส \textbf{อโธ}ติ อโธ เกสมตฺถกา.\csromanspacer{} \textbf{ติริยญฺจาปิ มชฺเฌ}ติ เวมชฺเฌติ อุทฺธํ อโธ ติริยญฺจาปิ มชฺเฌ.}
\prose{\csromanfolio{79}\textbf{เอเตสุ นนฺทิญฺจ นิเวสนญฺจ, ปนุชฺช วิญฺญาณํ ภเว น ติฏฺเฐ}ติ \textbf{เอเตสู}ติ อาจิกฺขิเตสุ เทสิเตสุ ปญฺญปิเตสุ ปฏฺฐปิเตสุ วิวริเตสุ วิภชิเตสุ อุตฺตานีกเตสุ ปกาสิเตสุ.\csromanspacer{} \textbf{นนฺที} วุจฺจติ ตณฺหา.\csromanspacer{} โย ราโค สาราโค ฯเปฯ อภิชฺฌา โลโภ อกุสลมูลํ.\csromanspacer{} \textbf{นิเวสนนฺ}ติ เทฺว นิเวสนา, ตณฺหานิเวสนา จ ทิฏฺฐินิเวสนา จ.\csromanspacer{} กตมา ตณฺหานิเวสนา, ยาวตา ตณฺหาสงฺขาเตน ฯเปฯ อยํ ตณฺหานิเวสนา.\csromanspacer{} กตมา ทิฏฺฐินิเวสนา, วีสติวตฺถุกา สกฺกายทิฏฺฐิ ฯเปฯ อยํ ทิฏฺฐินิเวสนา.}

\prose{\textbf{ปนุชฺช วิญฺญาณนฺ}ติ ปุญฺญาภิสงฺขารสหคตํ วิญฺญาณํ อปุญฺญาภิสงฺขารสหคตํ วิญฺญาณํ อาเนญฺชาภิสงฺขารสหคตํ วิญฺญาณํ, เอเตสุ นนฺทิญฺจ นิเวสนญฺจ อภิสงฺขารสหคตญฺจ วิญฺญาณํ นุชฺช ปนุชฺช นุท ปนุท ชห ปชห วิโนเทหิ พฺยนฺตีกโรหิ อนภาวํ คเมหีติ เอเตสุ นนฺทิญฺจ นิเวสนญฺจ, ปนุชฺช วิญฺญาณํ.}

\prose{\textbf{ภเว น ติฏฺเฐ}ติ \textbf{ภวา}ติ เทฺว \textbf{ภวา} กมฺมภโว จ ปฏิสนฺธิโก จ ปุนพฺภโว.\csromanspacer{} กตโม กมฺมภโว, ปุญฺญาภิสงฺขาโร อปุญฺญาภิสงฺขาโร อาเนญฺชาภิสงฺขาโร, อยํ กมฺมภโว.\csromanspacer{} กตโม ปฏิสนฺธิโก ปุนพฺภโว, ปฏิสนฺธิกา รูปํ เวทนา สญฺญา สงฺขารา วิญฺญาณํ, อยํ ปฏิสนฺธิโก ปุนพฺภโว.\csromanspacer{} \textbf{ภเว น ติฏฺเฐ}ติ นนฺทิญฺจ นิเวสนญฺจ อภิสงฺขารสหคตํ วิญฺญาณญฺจ กมฺมภวญฺจ ปฏิสนฺธิกญฺจ ปุนพฺภวํ ปชหนฺโต วิโนเทนฺโต พฺยนฺตีกโรนฺโต อนภาวํ คเมนฺโต กมฺมภเว น ติฏฺเฐยฺย, ปฏิสนฺธิเก ปุนพฺภเว น ติฏฺเฐยฺย น สนฺติฏฺเฐยฺยาติ ปนุชฺช วิญฺญาณํ ภเว น ติฏฺเฐ.\csromanspacer{} เตนาห ภควา\csromandash{}}

\prose{ยํ กิญฺจิ สมฺปชานาสิ, (เมตฺตคูติ ภควา,)}

\prose{อุทฺธํ อโธ ติริยญฺจาปิ มชฺเฌ.}

\prose{เอเตสุ นนฺทิญฺจ นิเวสนญฺจ,}

\prose{ปนุชฺช วิญฺญาณํ ภเว น ติฏฺเฐติ.}

\csromanhangingbegin
\hangingheaditem{25}{เอวํวิหารี สโต อปฺปมตฺโต,}
\hangingbody{\textbf{ภิกฺขุ จรํ หิตฺวา มมายิตานิ.}}
\hangingbody{\textbf{ชาติํ ชรํ โสกปริทฺทวญฺจ,}}
\hangingbody{อิเธว วิทฺวา ปชเหยฺย ทุกฺขํ.}
\csromanhangingend

\prose{\csromanfolio{80}\textbf{เอวํวิหารี สโต อปฺปมตฺโต}ติ \textbf{เอวํวิหารี}ติ นนฺทิญฺจ นิเวสนญฺจ อภิสงฺขารสหคตวิญฺญาณญฺจ กมฺมภวญฺจ ปฏิสนฺธิกญฺจ ปุนพฺภวํ ปชหนฺโต วิโนเทนฺโต พฺยนฺตีกโรนฺโต อนภาวํ คเมนฺโตติ \textbf{เอวํวิหารี}.\csromanspacer{} \textbf{สโต}ติ จตูหิ การเณหิ สโต, กาเย กายานุปสฺสนาสติปฏฺฐานํ ภาเวนฺโต ฯเปฯ โส วุจฺจติ สโต.\csromanspacer{} \textbf{อปฺปมตฺโต}ติ สกฺกจฺจการี สาตจฺจการี อฏฺฐิตการี อโนลีนวุตฺตีโก\footnote{อโนลีนวุตฺติโก (ก) ขุ 7. 45 ปิฏฺเฐ.} อนิกฺขิตฺตจฺฉนฺโท อนิกฺขิตฺตธุโร อปฺปมตฺโต กุสเลสุ ธมฺเมสุ “กถาหํ\footnote{กทาหํ (สฺยา)} อปริปูรํ วา สีลกฺขนฺธํ ปริปูเรยฺยํ, ปริปูรํ วา สีลกฺขนฺธํ ตตฺถ ตตฺถ ปญฺญาย อนุคฺคเณฺหยฺยนฺ”ติ โย ตตฺถ ฉนฺโท จ วายาโม จ อุสฺสาโห จ อุสฺโสฬฺหี จ อปฺปฏิวานี\footnote{อปฺปฏิวานิ (ก) ขุ 7. 45 ปิฏฺเฐ.} จ สติ จ สมฺปชญฺญญฺจ อาตปฺปํ ปธานํ อธิฏฺฐานํ อนุโยโค อปฺปมตฺโต อปฺปมาโท กุสเลสุ ธมฺเมสุ. “กถาหํ อปริปูรํ วา สมาธิกฺขนฺธํ.\csromanspacer{} ปญฺญากฺขนฺธํ.\csromanspacer{} วิมุตฺติกฺขนฺธํ.\csromanspacer{} วิมุตฺติญาณทสฺสนกฺขนฺธํ ปริปูเรยฺยํ, ปริปูรํ วา วิมุตฺติญาณทสฺสนกฺขนฺธํ ตตฺถ ตตฺถ ปญฺญาย อนุคฺคเณฺหยฺยนฺ”ติ โย ตตฺถ ฉนฺโท จ วายาโม จ อุสฺสาโห จ อุสฺโสฬฺหี จ อปฺปฏิวานี จ สติ จ สมฺปชญฺญญฺจ อาตปฺปํ ปธานํ อธิฏฺฐานํ อนุโยโค อปฺปมตฺโต อปฺปมาโท กุสเลสุ ธมฺเมสุ. “กถาหํ อปริญฺญาตํ วา ทุกฺขํ ปริชาเนยฺยํ, อปฺปหีเน วา กิเลเส ปชเหยฺยํ, อภาวิตํ วา มคฺคํ ภาเวยฺยํ, อสจฺฉิกตํ วา นิโรธํ สจฺฉิกเรยฺยนฺ”ติ โย ตตฺถ ฉนฺโท จ วายาโม จ อุสฺสาโห จ อุสฺโสฬฺหี จ อปฺปฏิวานี จ สติ จ สมฺปชญฺญญฺจ อาตปฺปํ ปธานํ อธิฏฺฐานํ อนุโยโค อปฺปมตฺโต อปฺปมาโท กุสเลสุ ธมฺเมสูติ \textbf{เอวํวิหารี} สโต อปฺปมตฺโต.}

\prose{\textbf{ภิกฺขุ จรํ หิตฺวา มมายิตานี}ติ \textbf{ภิกฺขู}ติ ปุถุชฺชนกลฺยาณโก\footnote{กลฺยาณปุถุชฺชโน (สฺยา), เอวมีทิเสสุ ฐาเนสุ.} วา ภิกฺขุ เสกฺโข วา ภิกฺขุ.\csromanspacer{} \textbf{จรนฺ}ติ จรนฺโต วิหรนฺโต อิริยนฺโต วตฺเตนฺโต ปาเลนฺโต ยเปนฺโต ยาเปนฺโต.\csromanspacer{} \textbf{มมตฺตา}ติ เทฺว มมตฺตา ตณฺหามมตฺตญฺจ ทิฏฺฐิมมตฺตญฺจ ฯเปฯ อิทํ ตณฺหามมตฺตํ ฯเปฯ อิทํ ทิฏฺฐิมมตฺตํ.\csromanspacer{} ตณฺหามมตฺตํ ปหาย ทิฏฺฐิมมตฺตํ ปฏินิสฺสชฺชิตฺวา มมตฺเต ชหิตฺวา จชิตฺวา ปชหิตฺวา วิโนเทตฺวา พฺยนฺตีกริตฺวา อนภาวํ คเมตฺวาติ ภิกฺขุ จรํ หิตฺวา มมายิตานิ.}

\csromanheader{2}{4. เมตฺตาคูมาณวปุจฺฉานิทฺเทส}
\prose{\csromanfolio{81}\textbf{ชาติํ ชรํ โสกปริทฺทวญฺจ, อิเธว วิทฺวา ปชเหยฺย ทุกฺขนฺ}ติ \textbf{ชาตี}ติ ยา เตสํ เตสํ สตฺตานํ ฯเปฯ.\csromanspacer{} \textbf{ชรนฺ}ติ ยา เตสํ เตสํ สตฺตานํ ฯเปฯ.\csromanspacer{} \textbf{โสโก}ติ ญาติพฺยสเนน วา ผุฏฺฐสฺส ฯเปฯ.\csromanspacer{} \textbf{ปริเทโว}ติ ญาติพฺยสเนน วา ผุฏฺฐสฺส ฯเปฯ.\csromanspacer{} \textbf{อิธา}ติ อิมิสฺสา ทิฏฺฐิยา ฯเปฯ.\csromanspacer{} อิมสฺมิํ มนุสฺสโลเก.\csromanspacer{} \textbf{วิทฺวา}ติ วิชฺชาคโต ญาณี วิภาวี เมธาวี.\csromanspacer{} \textbf{ทุกฺขนฺ}ติ ชาติทุกฺขํ ฯเปฯ โทมนสฺสุปายาสทุกฺขํ.\csromanspacer{}\csromanspacer{}\textbf{ชาติํ ชรํ โสกปริทฺทวญฺจ, อิเธว วิทฺวา ปชเหยฺย ทุกฺขนฺ}ติ วิชฺชาคโต ญาณี วิภาวี เมธาวี อิเธว ชาติญฺจ ชรญฺจ โสกปริทฺทวญฺจ ทุกฺขญฺจ ปชเหยฺย วิโนเทยฺย พฺยนฺตีกเรยฺย อนภาวํ คเมยฺยาติ ชาติํ ชรํ โสกปริทฺทวญฺจ, อิเธว วิทฺวา ปชเหยฺย ทุกฺขํ.\csromanspacer{} เตนาห ภควา\csromandash{}}

\csromangathameasure{เอวํวิหารี สโต อปฺปมตฺโต, \\ ภิกฺขุ จรํ หิตฺวา มมายิตานิ.}
\csromangathagroup{\csromangathastanza{เอวํวิหารี สโต อปฺปมตฺโต, \\ ภิกฺขุ จรํ หิตฺวา มมายิตานิ.}}

\prose{ชาติํ ชรํ โสกปริทฺทวญฺจ,}

\prose{อิเธว วิทฺวา ปชเหยฺย ทุกฺขนฺติ.}

\proseitem{26}{\textbf{เอตา’ภินนฺทามิ วโจ มเหสิโน},}

\prose{\textbf{สุกิตฺติตํ โคตม’นูปธีกํ.}}

\prose{\textbf{อทฺธา หิ ภควา ปหาสิ ทุกฺขํ,}}

\csromanheader{2}{(9)}
\prose{\textbf{เอตา’ภินนฺทามิ วโจ มเหสิโน}ติ \textbf{เอตนฺ}ติ ตุยฺหํ วจนํ พฺยปฺปถํ เทสนํ อนุสาสนํ อนุสิฏฺฐํ นนฺทามิ อภินนฺทามิ โมทามิ อนุโมทามิ อิจฺฉามิ สาทิยามิ ปตฺถยามิ ปิหยามิ อภิชปฺปามิ.\csromanspacer{} \textbf{มเหสิโน}ติ กิํ มเหสิ ภควา, มหนฺตํ สีลกฺขนฺธํ เอสี คเวสี ปริเยสีติ มเหสิ ฯเปฯ “กหํ นราสโภ”ติ มเหสีติ เอตาภินนฺทามิ วโจ \textbf{มเหสิโน}.}

\prose{\textbf{สุกิตฺติตํ โคตมนูปธีกนฺ}ติ \textbf{สุกิตฺติตนฺ}ติ สุกิตฺติตํ สุ-อาจิกฺขิตํ สุเทสิตํ สุปญฺญปิตํ สุปฏฺฐปิตํ สุวิวริตํ สุวิภชิตํ สุ- อุตฺตานีกตํ สุปกาสิตนฺติ สุกิตฺติตํ.\csromanspacer{} \textbf{โคตมนูปธีกนฺ}ติ อุปธี วุจฺจนฺติ กิเลสา จ ขนฺธา จ อภิสงฺขารา จ.\csromanspacer{} อุปธิปฺปหานํ อุปธิวูปสมํ อุปธิปฏินิสฺสคฺคํ อุปธิปฏิปสฺสทฺธํ อมตํ นิพฺพานนฺติ สุกิตฺติตํ โคตมนูปธีกํ.}

\prose{\csromanfolio{82}\textbf{อทฺธา หิ ภควา ปหาสิ ทุกฺขนฺ}ติ \textbf{อทฺธา}ติ เอกํสวจนํ นิสฺสํสยวจนํ นิกฺกงฺขาวจนํ อเทฺวชฺฌวจนํ อเทฺวฬฺหกวจนํ นิโรธวจนํ อปณฺณกวจนํ อวตฺถาปนวจนเมตํ \textbf{อทฺธา}ติ.\csromanspacer{} \textbf{ภควา}ติ คารวาธิวจนเมตํ ฯเปฯ สจฺฉิกา ปญฺญตฺติ, ยทิทํ \textbf{ภควา}ติ.\csromanspacer{} \textbf{ปหาสิ ทุกฺขนฺ}ติ ชาติทุกฺขํ ชราทุกฺขํ พฺยาธิทุกฺขํ มรณทุกฺขํ โสกปริเทวทุกฺขโทมนสฺสุปายาสทุกฺขํ ปหาสิ ปชหิ วิโนเทสิ พฺยนฺตีกโรสิ อนภาวํ คเมสีติ \textbf{อทฺธา} หิ \textbf{ภควา} ปหาสิ ทุกฺขํ.}

\prose{\textbf{ตถา หิ เต วิทิโต เอส ธมฺโม}ติ ตถา หิ เต วิทิโต ตุลิโต ตีริโต วิภูโต วิภาวิโต เอส ธมฺโมติ ตถา หิ เต วิทิโต เอส ธมฺโม.\csromanspacer{} เตนาห โส พฺราหฺมโณ\csromandash{}}

\csromangathameasure{เอตา’ภินนฺทามิ วโจ มเหสิโน, \\ สุกิตฺติตํ โคตม’นูปธีกํ.}
\csromangathagroup{\csromangathastanza{เอตา’ภินนฺทามิ วโจ มเหสิโน, \\ สุกิตฺติตํ โคตม’นูปธีกํ.}}

\prose{อทฺธา หิ \textbf{ภควา} ปหาสิ ทุกฺขํ,}

\prose{\textbf{ตถา หิ เต วิทิโต เอส ธมฺโม}ติ.}

\proseitem{27}{เต จาปิ นูนปฺปชเหยฺยุ ทุกฺขํ,}

\prose{\textbf{เย ตฺวํ มุนี อฏฺฐิตํ โอวเทยฺย.}}

\prose{\textbf{ตํ ตํ นมสฺสามิ สเมจฺจ นาคํ,}}

\csromanheader{2}{\textbf{อปฺเปว มํ ภควา อฏฺฐิตํ โอวเทยฺย. (10)}}
\prose{\textbf{เต จาปิ นูนปฺปชเหยฺยุ ทุกฺขนฺ}ติ \textbf{เต จาปี}ติ ขตฺติยา จ พฺราหฺมณา จ เวสฺสา จ สุทฺทา จ คหฏฺฐา จ ปพฺพชิตา จ เทวา จ มนุสฺสา จ.\csromanspacer{} \textbf{ปชเหยฺยุ ทุกฺขนฺ}ติ ชาติทุกฺขํ ชราทุกฺขํ พฺยาธิทุกฺขํ มรณทุกฺขํ โสกปริเทว ทุกฺขโทมนสฺสุปายาสทุกฺขํ ปชเหยฺยุํ วิโนเทยฺยุํ พฺยนฺตีกเรยฺยุํ อนภาวํ คเมยฺยุนฺติ เต จาปิ นูนปฺปชเหยฺยุ ทุกฺขํ.}

\prose{\textbf{เย ตฺวํ มุนี อฏฺฐิตํ โอวเทยฺยา}ติ \textbf{เย}ติ ขตฺติ\textbf{เย} จ พฺราหฺมเณ จ เวสฺเส จ สุทฺเท จ คหฏฺเฐ จ ปพฺพชิเต จ เทเว จ มนุสฺเส จ.\csromanspacer{} \textbf{ตฺวนฺ}ติ ภควนฺตํ ภณติ.\csromanspacer{} \textbf{มุนี}ติ โมนํ วุจฺจติ ญาณํ ฯเปฯ สงฺคชาลมติจฺจ โส มุนิ.\csromanspacer{} \textbf{อฏฺฐิตํ โอวเทยฺยา}ติ อฏฺฐิตํ โอวเทยฺย, สกฺกจฺจํ โอวเทยฺย, อภิณฺหํ โอวเทยฺย, ปุนปฺปุนํ โอวเทยฺย อนุสาเสยฺยาติ \textbf{เย} ตฺวํ มุนี อฏฺฐิตํ โอวเทยฺย.}

\prose{\textbf{ตํ ตํ นมสฺสามิ สเมจฺจ นาคนฺ}ติ \textbf{ตนฺ}ติ ภควนฺตํ ภณติ.\csromanspacer{} \textbf{นมสฺสามี}ติ กา\textbf{เย}น วา นมสฺสามิ, วาจาย วา นมสฺสามิ, จิตฺเตน วา นมสฺสามิ,}

\vspace{\tipitakaparskip}
\csromancenter{เมตฺตาคูมาณวปุจฺฉานิทฺเทส อนฺวตฺถปฏิปตฺติยา วา นมสฺสามิ, ธมฺมานุธมฺมปฏิปตฺติยา วา นมสฺสามิ สกฺกโรมิ ครุํ กโรมิ\footnote{ครุกโรมิ (สฺยา)} มาเนมิ ปูเชมิ.\csromanspacer{} \textbf{สเมจฺจา}ติ สเมจฺจ อภิสเมจฺจ สมาคนฺตฺวา อภิสมาคนฺตฺวา สมฺมุขา ตํ นมสฺสามิ.\csromanspacer{} \textbf{นาคนฺ}ติ นาโค จ ภควา, อาคุํ น กโรตีติ นาโค, น คจฺฉตีติ นาโค, น อาคจฺฉตีติ นาโค, กถํ ภควา อาคุํ น กโรตีติ นาโค.\csromanspacer{} อาคุ วุจฺจติ ปาปกา อกุสลา ธมฺมา สํกิเลสิกา โปโนภวิกา สทรา ทุกฺขวิปากา อายติํ ชาติชรามรณิยา.}
\csromancenter{อาคุํ น กโรติ กิญฺจิ โลเก, (สภิยาติ ภควา,)}
\prose{\csromanfolio{83}สพฺพสํโยเค\footnote{สพฺพโยเค (ก) ขุ 1. 359 ปิฏฺเฐ.} วิสชฺช พนฺธนานิ.}

\csromangathameasure{สพฺพตฺถ น สชฺชตี วิมุตฺโต, \\ นาโค ตาทิ ปวุจฺจเต ตถตฺตาติ.}
\csromangathagroup{\csromangathastanza{สพฺพตฺถ น สชฺชตี วิมุตฺโต, \\ นาโค ตาทิ ปวุจฺจเต ตถตฺตาติ.}}

\prose{เอวํ ภควา อาคุํ น กโรตีติ นาโค.}

\prose{กถํ ภควา น คจฺฉตีติ นาโค, ภควา น ฉนฺทาคติํ คจฺฉติ, น โทสาคติํ คจฺฉติ, น โมหาคติํ คจฺฉติ, น ภยาคติํ คจฺฉติ, น ราควเสน คจฺฉติ, น โทสวเสน คจฺฉติ, น โมหวเสน คจฺฉติ, น มานวเสน คจฺฉติ, น ทิฏฺฐิวเสน คจฺฉติ, น อุทฺธจฺจวเสน คจฺฉติ, น วิจิกิจฺฉาวเสน คจฺฉติ, น อนุสยวเสน คจฺฉติ, น วคฺเคหิ ธมฺเมหิ ยายติ นียติ\footnote{นิยฺยติ (สฺยา, ก)} วุยฺหติ สํหรียติ.\csromanspacer{} เอวํ ภควา น คจฺฉตีติ นาโค.}

\prose{กถํ ภควา น อาคจฺฉตีติ นาโค, โสตาปตฺติมคฺเคน เย กิเลสา ปหีนา, เต กิเลเส น ปุเนติ น ปจฺเจติ น ปจฺจาคจฺฉติ.\csromanspacer{} สกทาคามิมคฺเคน.\csromanspacer{} อนาคามิมคฺเคน.\csromanspacer{} อรหตฺตมคฺเคน เย กิเลสา ปหีนา, เต กิเลเส น ปุเนติ น ปจฺเจติ น ปจฺจาคจฺฉติ.\csromanspacer{} เอวํ ภควา น อาคจฺฉตีติ นาโคติ ตํ ตํ นมสฺสามิ สเมจฺจ นาค.}

\prose{\textbf{อปฺเปว มํ ภควา อฏฺฐิตํ โอวเทยฺยา}ติ อปฺเปว มํ ภควา อฏฺฐิตํ โอวเทยฺย, สกฺกจฺจํ โอวเทยฺย, อภิณฺหํ โอวเทยฺย, ปุนปฺปุนํ โอวเทยฺย, อนุสาเสยฺยาติ อปฺเปว มํ ภควา อฏฺฐิตํ โอวเทยฺย.\csromanspacer{} เตนาห โส พฺราหฺมโณ\csromandash{}}

\prose{\csromanfolio{84}เต จาปิ นูนปฺปชเหยฺยุ ทุกฺขํ.}

\prose{เย ตฺวํ มุนี อฏฺฐิตํ โอวเทยฺย.}

\csromangathameasure{ตํ ตํ นมสฺสามิ สเมจฺจ นาคํ, \\ อปฺเปว มํ ภควา อฏฺฐิตํ โอวเทยฺยาติ.}
\csromangathagroup{\csromangathastanza{ตํ ตํ นมสฺสามิ สเมจฺจ นาคํ, \\ อปฺเปว มํ ภควา อฏฺฐิตํ โอวเทยฺยาติ.}}

\proseitem{28}{ยํ พฺราหฺมณํ เวทคุ’มาภิชญฺญา,}

\prose{\textbf{อกิญฺจนํ กามภเว อสตฺตํ.}}

\prose{\textbf{อทฺธา หิ โส โอฆมิมํ อตาริ,}}

\csromanheader{2}{\textbf{ติณฺโณ จ ปารํ อขิโล อกงฺโข. (11)}}
\prose{\textbf{ยํ พฺราหฺมณํ เวทคุมาภิชญฺญา}ติ \textbf{พฺราหฺมโณ}ติ สตฺตนฺนํ ธมฺมานํ พาหิตตฺตา \textbf{พฺราหฺมโณ}, สกฺกายทิฏฺฐิ พาหิตา โหติ, วิจิกิจฺฉา พาหิตา โหติ, สีลพฺพตปรามาโส พาหิโต โหติ, ราโค พาหิโต โหติ, โทโส พาหิโต โหติ, โมโห พาหิโต โหติ, มาโน พาหิโต โหติ, พาหิตาสฺส โหนฺติ ปาปกา อกุสลา ธมฺมา สํกิเลสิกา โปโนภวิกา สทรา ทุกฺขวิปากา อายติํ ชาติชรามรณิยา.}

\prose{พาหิตฺวา สพฺพปาปกานิ, (สภิยาติ ภควา,)}

\prose{วิมโล สาธุสมาหิโต ฐิตตฺโต.}

\csromangathameasure{สํสารมติจฺจ เกวลี โส, \\ อสิโต ตาทิ ปวุจฺจเต ส พฺรหฺมา.}
\csromangathagroup{\csromangathastanza{สํสารมติจฺจ เกวลี โส, \\ อสิโต\footnote{อนิสฺสิโต (สฺยา) ขุ 1. 358 ปิฏฺเฐ} ตาทิ ปวุจฺจเต ส พฺรหฺมา.}}

\prose{\textbf{เวทคู}ติ เวโท วุจฺจติ จตูสุ มคฺเคสุ ญาณํ ฯเปฯ สพฺพํ เวทมติจฺจ เวทคู โสติ.\csromanspacer{} \textbf{อภิชญฺญา}ติ อภิชาเนยฺย อาชาเนยฺย วิชาเนยฺย ปฏิวิชาเนยฺย ปฏิวิชฺเฌยฺยาติ ยํ พฺราหฺมณํ เวทคุมาภิชญฺญา.}

\prose{\textbf{อกิญฺจนํ กามภเว อสตฺตนฺ}ติ \textbf{อกิญฺจนนฺ}ติ ราคกิญฺจนํ โทสกิญฺจนํ โมหกิญฺจนํ มานกิญฺจนํ ทิฏฺฐิกิญฺจนํ กิเลสกิญฺจนํ ทุจฺจริตกิญฺจนํ.\csromanspacer{} ยสฺเสเต กิญฺจนา ปหีนา สมุจฺฉินฺนา วูปสนฺตา ปฏิปสฺสทฺธา อภพฺพุปฺปตฺติกา ญาณคฺคินา ทฑฺฒา, โส วุจฺจติ อกิญฺจโน.\csromanspacer{} \textbf{กามา}ติ อุทฺทานโต เทฺว กามา วตฺถุกามา จ กิเลสกามา จ ฯเปฯ.\csromanspacer{} อิเม วุจฺจนฺติ วตฺถุกามา ฯเปฯ.\csromanspacer{} อิเม วุจฺจนฺติ กิเลสกามา.\csromanspacer{} \textbf{ภวา}ติ เทฺว ภวา กมฺมภโว จ ปฏิสนฺธิโก จ}

\vspace{\tipitakaparskip}
\csromancenter{เมตฺตาคูมาณวปุจฺฉานิทฺเทส ปุนพฺภโว ฯเปฯ อยํ ปฏิสนฺธิโก ปุนพฺภโว.\csromanspacer{}\csromanspacer{}\textbf{อกิญฺจนํ กามภเว อสตฺตนฺ}ติ อกิญฺจนํ ปุคฺคลํ กามภเว จ อสตฺตํ อลคฺคํ อลคฺคิตํ อปลิพุทฺธํ นิกฺขนฺตํ นิสฺสฏํ วิปฺปมุตฺตํ วิสญฺญุตฺตํ วิมริยาทิกเตน เจตสา วิหรนฺตนฺติ อกิญฺจนํ กามภเว อสตฺตํ.}
\prose{\csromanfolio{85}\textbf{อทฺธา หิ โส โอฆมิมํ อตารี}ติ \textbf{อทฺธา}ติ เอกํสวจนํ ฯเปฯ อวตฺถาปนวจนเมตํ \textbf{อทฺธา}ติ.\csromanspacer{} \textbf{โอฆนฺ}ติ กาโมฆํ ภโวฆํ ทิฏฺโฐฆํ อวิชฺโชฆํ.\csromanspacer{} \textbf{อตารี}ติ อุตฺตริ ปตริ สมติกฺกมิ วีติวตฺตยีติ \textbf{อทฺธา} หิ โส โอฆมิมํ อตาริ.}

\prose{\textbf{ติณฺโณ จ ปารํ อขิโล อกงฺโข}ติ \textbf{ติณฺโณ}ติ กาโมฆํ \textbf{ติณฺโณ}, ภโวฆํ \textbf{ติณฺโณ}, ทิฏฺโฐฆํ \textbf{ติณฺโณ}, อวิชฺโชฆํ \textbf{ติณฺโณ}, สํสารปถํ \textbf{ติณฺโณ} อุตฺ\textbf{ติณฺโณ} นิตฺถิณฺโณ\footnote{นิตฺติณฺโณ (สฺยา)} อติกฺกนฺโต สมติกฺกนฺโต วีติวตฺโต.\csromanspacer{} โส วุตฺถวาโส\footnote{วุฏฺฐวาโส (สฺยา) ขุ 7. 16 ปิฏฺเฐปิ.} จิณฺณจรโณ คตทฺโธ คตทิโส คตโกฏิโก ปาลิตพฺรหฺมจริโย อุตฺตมทิฏฺฐิปฺปตฺโต ภาวิตมคฺโค, ปหีนกิเลโส ปฏิวิทฺธากุปฺโป สจฺฉิกตนิโรโธ, ทุกฺขํ ตสฺส ปริญฺญาตํ, สมุทโย ปหีโน, มคฺโค ภาวิโต, นิโรโธ สจฺฉิกโต, อภิญฺเญยฺยํ อภิญฺญาตํ, ปริญฺเญยฺยํ ปริญฺญาตํ, ปหาตพฺพํ ปหีนํ, ภาเวตพฺพํ ภาวิตํ, สจฺฉิกาตพฺพํ สจฺฉิกตํ.\csromanspacer{} โส อุกฺขิตฺตปลิโฆ สํกิณฺณปริกฺโข อพฺพุเฬฺหสิโก นิรคฺคโฬ อริโย ปนฺนทฺธโช ปนฺนภาโร วิสญฺญุตฺโต ปญฺจงฺควิปฺปหีโน ฉฬงฺคสมนฺนาคโต เอการกฺโข จตุราปสฺเสโน ปนุณฺณปจฺเจกสจฺโจ\footnote{ปณุนฺนปจฺเจกสจฺโจ (ก)} สมวยสฏฺเฐสโน อนาวิลสงฺกปฺโป ปสฺสทฺธกายสงฺขาโร สุวิมุตฺตจิตฺโต สุวิมุตฺตปญฺโญ เกวลี วุสิตวา อุตฺตมปุริโส ปรมปุริโส ปรมปตฺติปฺปตฺโต, โส เนว อาจินาติ น อปจินาติ อปจินิตฺวา ฐิโต, เนว ปชหติ น อุปาทิยติ ปชหิตฺวา ฐิโต, เนว วิสิเนติ น อุสฺสิเนติ วิสิเนตฺวา ฐิโต, เนว วิธูเปติ น สนฺธูเปติ วิธูเปตฺวา ฐิโต, อเสกฺเขน สีลกฺขนฺเธน สมนฺนาคตตฺตา ฐิโต, อเสกฺเขน สมาธิกฺขนฺเธน.\csromanspacer{} ปญฺญากฺขนฺเธน.\csromanspacer{} วิมุตฺติกฺขนฺเธน.\csromanspacer{} วิมุตฺติญาณทสฺสนกฺขนฺเธน สมนฺนาคตตฺตา ฐิโต, สจฺจํ สมฺปฏิปาทยิตฺวา\footnote{ปฏิปาทยิตฺวา (สฺยา)} ฐิโต, เอชํ สมติกฺกมิตฺวา ฐิโต, กิเลสคฺคิํ ปริยาทิยิตฺวา ฐิโต, อปริคมนตาย ฐิโต, กถํ\footnote{กฏํ (สฺยา) กามสุตฺตนิทฺเทสฏฺฐกถา โอโลเกตพฺพา.} สมาทาย ฐิโต, วิมุตฺติปฏิเสวนตาย \csromanfolio{86}ฐิโต, เมตฺตาย ปาริสุทฺธิยา ฐิโต, กรุณาย.\csromanspacer{} มุทิตาย.\csromanspacer{} อุเปกฺขาย ปาริสุทฺธิยา ฐิโต, อจฺจนฺตปาริสุทฺธิยา ฐิโต, อตมฺมยตาย\footnote{อกมฺมญฺญตาย (สฺยา)} ปาริสุทฺธิยา ฐิโต, วิมุตฺตตฺตา ฐิโต, สนฺตุสฺสิตตฺตา ฐิโต, ขนฺธปริยนฺเต ฐิโต, ธาตุปริยนฺเต ฐิโต, อายตนปริยนฺเต ฐิโต, คติปริยนฺเต ฐิโต, อุปปตฺติปริยนฺเต ฐิโต, ปฏิสนฺธิปริยนฺเต ฐิโต, ภวปริยนฺเต ฐิโต, สํสารปริยนฺเต ฐิโต, วฏฺฏปริยนฺเต ฐิโต, อนฺติมภเว ฐิโต, อนฺติเม สมุสฺสเย ฐิโต, อนฺติมเทหธโร อรหา.}

\csromangathameasure{ตสฺสายํ ปจฺฉิมโก ภโว, จริโมยํ สมุสฺสโย. \\ ชาติมรณสํสาโร, นตฺถิ ตสฺส ปุนพฺภโวติ.}
\csromangathagroup{\csromangathastanza{ตสฺสายํ ปจฺฉิมโก ภโว, จริโมยํ สมุสฺสโย. \\ ชาติมรณสํสาโร, นตฺถิ ตสฺส ปุนพฺภโวติ.}}

\prose{\textbf{ติณฺโณ จ ปารนฺ}ติ ปารํ วุจฺจติ อมตํ นิพฺพานํ.\csromanspacer{} โย โส สพฺพสงฺขารสมโถ สพฺพูปธิปฏินิสฺสคฺโค ตณฺหกฺขโย วิราโค นิโรโธ นิพฺพานํ, โส ปารคโต ปารปฺปตฺโต อนฺตคโต อนฺตปฺปตฺโต โกฏิคโต โกฏิปฺปตฺโต ปริยนฺตคโต ปริยนฺตปฺปตฺโต โวสานคโต โวสานปฺปตฺโต ตาณคโต ตาณปฺปตฺโต เลณคโต เลณปฺปตฺโต สรณคโต สรณปฺปตฺโต อภยคโต อภยปฺปตฺโต อจฺจุตคโต อจฺจุตปฺปตฺโต อมตคโต อมตปฺปตฺโต นิพฺพานคโต นิพฺพานปฺปตฺโต.\csromanspacer{} โส วุตฺถวาโส จิณฺณจรโณ ฯเปฯ ชาติมรณสํสาโร นตฺถิ ตสฺส ปุนพฺภโวติ ติณฺโณ จ ปารํ.}

\prose{\textbf{อขิโล}ติ ราโค ขิโล, โทโส ขิโล, โมโห ขิโล, โกโธ ขิโล, อุปนาโห ขิโล ฯเปฯ สพฺพากุสลาภิสงฺขารา ขิลา.\csromanspacer{} ยสฺเสเต ขิลา ปหีนา สมุจฺฉินฺนา วูปสนฺตา ปฏิปสฺสทฺธา อภพฺพุปฺปตฺติกา ญาณคฺคินา ทฑฺฒา, โส วุจฺจติ อขิโล.\csromanspacer{} \textbf{อกงฺโข}ติ ทุกฺเข กงฺขา, ทุกฺขสมุทเย กงฺขา, ทุกฺขนิโรเธ กงฺขา, ทุกฺขนิโรธคามินิยา ปฏิปทาย กงฺขา, ปุพฺพนฺเต กงฺขา, อปรนฺเต กงฺขา, ปุพฺพนฺตาปรนฺเต กงฺขา, อิทปฺปจฺจยตาปฏิจฺจสมุปฺปนฺเนสุ ธมฺเมสุ กงฺขา, ยา เอวรูปา กงฺขา กงฺขายนา กงฺขายิตตฺตํ วิมติ วิจิกิจฺฉา เทฺวฬฺหกํ เทฺวธาปโถ สํสโย อเนกํสคฺคาโห อาสปฺปนา ปริสปฺปนา อปริโยคาหนา ฉมฺภิตตฺตํ จิตฺตสฺส มโนวิเลโข, ยสฺเสเต กงฺขา ปหีนา สมุจฺฉินฺนา วูปสนฺตา ปฏิปสฺสทฺธา อภพฺพุปฺปตฺติกา ญาณคฺคินา ทฑฺฒา,}

\vspace{\tipitakaparskip}
\csromancenter{เมตฺตาคูมาณวปุจฺฉานิทฺเทส โส วุจฺจติ อกงฺโขติ ติณฺโณ จ ปารํ อขิโล อกงฺโข.\csromanspacer{} เตนาห ภควา\csromandash{}}
\vspace{0.5\baselineskip}
\csromangathameasure{ยํ พฺราหฺมณํ เวทคุมาภิชญฺญา, \\ อกิญฺจนํ กามภเว อสตฺตํ. \\ อทฺธา หิ โส โอฆมิมํ อตาริ, \\ ติณฺโณ จ ปารํ อขิโล อกงฺโขติ.}
\csromangathagroup{\csromangathastanza{\csromanfolio{87}ยํ พฺราหฺมณํ เวทคุมาภิชญฺญา, \\ อกิญฺจนํ กามภเว อสตฺตํ. \\ อทฺธา หิ โส โอฆมิมํ อตาริ, \\ ติณฺโณ จ ปารํ อขิโล อกงฺโขติ.}}

\proseitem{29}{วิทฺวา จ โย เวทคู นโร อิธ,}

\prose{\textbf{ภวาภเว สงฺคมิมํ วิสชฺช.}}

\prose{\textbf{โส วีตตโณฺห อนีโฆ นิราโส,}}

\csromanheader{2}{\textbf{อตาริ โส ชาติชรนฺติ พฺรูมิ. (12)}}
\prose{\textbf{วิทฺวา จ โย เวทคู นโร อิธา}ติ \textbf{วิทฺวา}ติ วิชฺชาคโต ญาณี วิภาวี เมธาวี.\csromanspacer{} \textbf{โย}ติ โย ยาทิโส ฯเปฯ มนุสฺโส วา.\csromanspacer{} \textbf{เวทคู}ติ เวโท วุจฺจติ จตูสุ มคฺเคสุ ญาณํ.\csromanspacer{} ปญฺญา ปญฺญินฺทฺริยํ ปญฺญาพลํ ธมฺมวิจยสมฺโพชฺฌงฺโค วีมํสา วิปสฺสนา สมฺมาทิฏฺฐิ\footnote{ญาณํ ฯเปฯ สพฺพเวทมติจฺจ เวทคู โสติ (สฺยา) ปสฺส ขุ 7. 158 ปิฏฺเฐ.}.\csromanspacer{} เตหิ เวเทหิ ชาติชรามรณสฺส อนฺตคโต อนฺตปฺปตฺโต โกฏิคโต โกฏิปฺปตฺโต ปริยนฺตคโต ปริยนฺตปฺปตฺโต โวสานคโต โวสานปฺปตฺโต ตาณคโต ตาณปฺปตฺโต เลณคโต เลณปฺปตฺโต สรณคโต สรณปฺปตฺโต อภยคโต อภยปฺปตฺโต อจฺจุตคโต อจฺจุตปฺปตฺโต อมตคโต อมตปฺปตฺโต นิพฺพานคโต นิพฺพานปฺปตฺโต.\csromanspacer{} เวทานํ วา อนฺตคโตติ เวทคู, เวเทหิ วา อนฺตคโตติ เวทคู, สตฺตนฺนํ วา ธมฺมานํ วิทิตตฺตา เวทคู, สกฺกายทิฏฺฐิ วิทิตา โหติ.\csromanspacer{} วิจิกิจฺฉา.\csromanspacer{} สีลพฺพตปรามาโส.\csromanspacer{} ราโค.\csromanspacer{} โทโส.\csromanspacer{} โมโห.\csromanspacer{} มาโน วิทิโต โหติ, วิทิตาสฺส โหนฺติ ปาปกา อกุสลา ธมฺมา สํกิเลสิกา โปโนภวิกา สทรา ทุกฺขวิปากา อายติํ ชาติชรามรณิยา.}

\prose{เวทานิ วิเจยฺย เกวลานิ, (สภิยาติ ภควา,)}

\prose{สมณานํ ยานีธตฺถิ พฺราหฺมณานํ.}

\csromangathameasure{สพฺพเวทนาสุ วีตราโค, \\ สพฺพํ เวทมติจฺจ เวทคู โส.}
\csromangathagroup{\csromangathastanza{สพฺพเวทนาสุ วีตราโค, \\ สพฺพํ เวทมติจฺจ เวทคู โส.}}

\prose{\csromanfolio{88}\textbf{นโร}ติ สตฺโต นโร มานโว โปโส ปุคฺคโล ชีโว ชาคุ\footnote{ชาตุ (สฺยา)} ชนฺตุ อินฺทคุ\footnote{อินฺทคู (สฺยา)} มนุโช.\csromanspacer{} \textbf{อิธา}ติ อิมิสฺสา ทิฏฺฐิยา ฯเปฯ อิมสฺมิํ มนุสฺสโลเกติ วิทฺวา จ โย เวทคู นโร อิธ.}

\prose{\textbf{ภวาภเว สงฺคมิมํ วิสชฺชา}ติ \textbf{ภวาภเว}ติ \textbf{ภวาภเว} กมฺมภเว ปุนพฺภเว กามภเว, กมฺมภเว กามภเว ปุนพฺภเว รูปภเว, กมฺมภเว รูปภเว ปุนพฺภเว อรูปภเว, กมฺมภเว อรูปภเว ปุนพฺภเว ปุนปฺปุนภเว, ปุนปฺปุนคติยา ปุนปฺปุน-อุปปตฺติยา ปุนปฺปุนปฏิสนฺธิยา ปุนปฺปุน- อตฺตภาวาภินิพฺพตฺติยา.\csromanspacer{} \textbf{สงฺคา}ติ สตฺต สงฺคา ราคสงฺโค โทสสงฺโค โมหสงฺโค มานสงฺโค ทิฏฺฐิสงฺโค กิเลสสงฺโค ทุจฺจริตสงฺโค.\csromanspacer{} วฺ\textbf{อิสชฺชา}ติ สงฺเค โวสชฺเชตฺวา วา วิสชฺช.\csromanspacer{} อถ วา สงฺเค พนฺเธ วิพนฺเธ อาพนฺเธ ลคฺเค ลคฺคิเต ปลิพุทฺเธ พนฺธเน โผฏยิตฺวา\footnote{โมจยิตฺวา (สฺยา)} วา วิสชฺช.\csromanspacer{} ยถา ยานํ วา วยฺหํ วา รถํ วา สกฏํ วา สนฺทมานิกํ วา สชฺชํ วิสชฺชํ กโรนฺติ วิโกเปนฺติ, เอวเมว เต สงฺเค โวสชฺเชตฺวา วา วิสชฺช.\csromanspacer{} อถ วา สงฺเค พนฺเธ วิพนฺเธ อาพนฺเธ ลคฺเค ลคฺคิเต ปลิพุทฺเธ พนฺธเน โผฏยิตฺวา วา วฺ\textbf{อิสชฺชา}ติ \textbf{ภวาภเว} สงฺคมิมํ วิสชฺช.}

\prose{\textbf{โส วีตตโณฺห อนีโฆ นิราโส, อตาริ โส ชาติชรนฺติ พฺรูมี}ติ \textbf{ตณฺหา}ติ รูป\textbf{ตณฺหา} ฯเปฯ ธมฺม\textbf{ตณฺหา}.\csromanspacer{} ยสฺเสสา \textbf{ตณฺหา} ปหีนา สมุจฺฉินฺนา วูปสนฺตา ปฏิปสฺสทฺธา อภพฺพุปฺปตฺติกา ญาณคฺคินา ทฑฺฒา, โส วุจฺจติ วีตตโณฺห วิคตตโณฺห จตฺตตโณฺห วนฺตตโณฺห มุตฺตตโณฺห ปหีนตโณฺห ปฏินิสฺสฏฺฐตโณฺห, วีตราโค จตฺตราโค ปหีนราโค ปฏินิสฺสฏฺฐราโค นิจฺฉาโต นิพฺพุโต สีติภูโต สุขปฺปฏิสํเวที พฺรหฺมภูเตน อตฺตนา วิหรตีติ โส วีตตโณฺห.\csromanspacer{} \textbf{อนีโฆ}ติ ราโค นีโฆ, โทโส นีโฆ, โมโห นีโฆ, โกโธ นีโฆ, อุปนาโห นีโฆ ฯเปฯ สพฺพากุสลาภิสงฺขารา นีฆา.\csromanspacer{} ยสฺเสเต นีฆา ปหีนา สมุจฺฉินฺนา วูปสนฺตา ปฏิปสฺสทฺธา อภพฺพุปฺปตฺติกา ญาณคฺคินา ทฑฺฒา, โส วุจฺจติ อนีโฆ.\csromanspacer{} \textbf{นิราโส}ติ อาสา วุจฺจติ \textbf{ตณฺหา}.\csromanspacer{} โย ราโค สาราโค ฯเปฯ อภิชฺฌา โลโภ อกุสลมูลํ.\csromanspacer{} ยสฺเสสา อาสา \textbf{ตณฺหา} ปหีนา สมุจฺฉินฺนา วูปสนฺตา ปฏิปสฺสทฺธา อภพฺพุปฺปตฺติกา ญาณคฺคินา ทฑฺฒา, โส วุจฺจติ นิราโส.\csromanspacer{} \textbf{ชาตี}ติ ยา เตสํ เตสํ สตฺตานํ ฯเปฯ อายตนานํ ปฏิลาโภ.\csromanspacer{} \textbf{ชรา}ติ ยา เตสํ เตสํ}

\vspace{\tipitakaparskip}
\csromancenter{โธตกมาณวปุจฺฉานิทฺเทส สตฺตานํ ฯเปฯ อินฺทฺริยานํ ปริปาโก.\csromanspacer{} อยํ วุจฺจติ ชรา.\csromanspacer{} \textbf{โส วีตตโณฺห อนีโฆ นิราโส, อตาริ โส ชาติชรนฺติ พฺรูมี}ติ โย โส วีตตโณฺห อนีโฆ จ นิราโส จ, โส โข ชาติชรามรณํ อตริ อุตฺตริ ปตริ สมติกฺกมิ วีติวตฺตยีติ พฺรูมิ อาจิกฺขามิ เทเสมิ ปญฺญเปมิ ปฏฺฐเปมิ วิวรามิ วิภชามิ อุตฺตานีกโรมิ ปกาเสมีติ โส วีตตโณฺห อนีโฆ นิราโส, อตาริ โส ชาติชรนฺติ พฺรูมิ.\csromanspacer{} เตนาห \textbf{ภควา}\csromandash{}}
\vspace{0.5\baselineskip}
\csromangathameasure{วิทฺวา จ โย เวทคู นโร อิธ, \\ ภวาภเว สงฺคมิมํ วิสชฺช. \\ โส วีตตโณฺห อนีโฆ นิราโส, \\ อตาริ โส ชาติชรนฺติ พฺรูมีติ.}
\csromangathagroup{\csromangathastanza{\csromanfolio{89}วิทฺวา จ โย เวทคู นโร อิธ, \\ ภวาภเว สงฺคมิมํ วิสชฺช. \\ โส วีตตโณฺห อนีโฆ นิราโส, \\ อตาริ โส ชาติชรนฺติ พฺรูมีติ.}}

\prose{สห คาถาปริโยสานา ฯเปฯ “สตฺถา เม ภนฺเต \textbf{ภควา}, สาวโกหมสฺมี”ติ.}

\vspace{\tipitakaparskip}
\csromancenter{เมตฺตคูมาณวปุจฺฉานิทฺเทโส จตุตฺโถ.}
\csromanmatikaanchor{mtk.26}
\csromantocmarkref{subsection}{5. โธตกมาณวปุจฺฉานิทฺเทส}{mtk.26}
\bookmark[dest={mtk.26},level=2]{5. โธตกมาณวปุจฺฉานิทฺเทส}
\csromanheader{2}{5. โธตกมาณวปุจฺฉานิทฺเทส}
\csromanhangingbegin
\hangingheaditem{30}{ปุจฺฉามิ ตํ \textbf{ภควา} พฺรูหิ เม’ตํ, (\textbf{อิจฺจา}ยสฺมา \textbf{โธตโก},)}
\hangingbody{\textbf{วาจาภิกงฺขามิ มเหสิ ตุยฺหํ.}}
\hangingbody{ตว สุตฺวาน นิคฺโฆสํ, สิกฺเข นิพฺพานมตฺตโน.}
\csromanhangingend

\prose{\textbf{ปุจฺฉามิ ตํ ภควา พฺรูหิ เม’ตนฺ}ติ \textbf{ปุจฺฉามี}ติ ติสฺโส ปุจฺฉา อทิฏฺฐโชตนา ปุจฺฉา ทิฏฺฐสํสนฺทนา ปุจฺฉา วิมติจฺเฉทนา ปุจฺฉา ฯเปฯ.\csromanspacer{} อิมา ติสฺโส ปุจฺฉา ฯเปฯ นิพฺพานปุจฺฉา.\csromanspacer{} \textbf{ปุจฺฉามิ ตนฺ}ติ ปุจฺฉามิ ตํ, ยาจามิ ตํ, อชฺเฌสามิ ตํ, ปสาเทมิ ตํ, “กถยสฺสุ เม”ติ ปุจฺฉามิ ตํ.\csromanspacer{} \textbf{ภควา}ติ คารวาธิวจนเมตํ ฯเปฯ สจฺฉิกา ปญฺญตฺติ, ยทิทํ \textbf{ภควา}ติ.\csromanspacer{} \textbf{พฺรูหิ เมตนฺ}ติ พฺรูหิ อาจิกฺขาหิ เทเสหิ ปญฺญเปหิ ปฏฺฐเปหิ วิวราหิ วิภชาหิ อุตฺตานีกโรหิ ปกาเสหีติ ปุจฺฉามิ ตํ \textbf{ภควา} พฺรูหิ \textbf{เม’ตํ.}}

\prose{\textbf{อิจฺจายสฺมา โธตโก}ติ \textbf{อิจฺจา}ติ ปทสนฺธิ ฯเปฯ.\csromanspacer{} \textbf{อายสฺมา}ติ ปิยวจนํ ครุวจนํ สคารวสปฺปติสฺสาธิวจนเมตํ อายสฺมาติ.\csromanspacer{} \textbf{โธตโก}ติ \csromanfolio{90}ตสฺส พฺราหฺมณสฺส นามํ สงฺขา สมญฺญา ปญฺญตฺติ โวหาโร นามํ นามกมฺมํ นามเธยฺยํ นิรุตฺติ พฺยญฺชนํ อภิลาโปติ อิจฺจายสฺมา โธตโก.}

\prose{\textbf{วาจาภิกงฺขามิ มเหสิ ตุยฺหนฺ}ติ ตุยฺหํ วจนํ พฺยปฺปถํ เทสนํ อนุสาสนํ อนุสิฏฺฐํ กงฺขามิ อภิกงฺขามิ อิจฺฉามิ สาทิยามิ ปตฺถยามิ ปิหยามิ อภิชปฺปามิ.\csromanspacer{} \textbf{มเหสี}ติ กิํ มเหสิ ภควา, มหนฺตํ สีลกฺขนฺธํ เอสี คเวสี ปริเยสีติ มเหสิ ฯเปฯ. “กหํ นราสโภ”ติ \textbf{มเหสี}ติ วาจาภิกงฺขามิ มเหสิ ตุยฺหํ.}

\prose{\textbf{ตว สุตฺวาน นิคฺโฆสนฺ}ติ ตุยฺหํ วจนํ พฺยปฺปถํ เทสนํ อนุสาสนํ อนุสิฏฺฐํ สุตฺวา สุณิตฺวา อุคฺคเหตฺวา อุปธารยิตฺวา อุปลกฺขยิตฺวาติ ตว สุตฺวาน นิคฺโฆสํ.}

\prose{\textbf{สิกฺเข นิพฺพานมตฺตโน}ติ \textbf{สิกฺขา}ติ ติสฺโส \textbf{สิกฺขา} อธิสีล\textbf{สิกฺขา} อธิจิตฺต\textbf{สิกฺขา} อธิปญฺญา\textbf{สิกฺขา} ฯเปฯ อยํ อธิปญฺญา\textbf{สิกฺขา}.\csromanspacer{} นิพฺพาน\textbf{มตฺตโน}ติ อตฺตโน ราคสฺส นิพฺพาปนาย, โทสสฺส นิพฺพาปนาย, โมหสฺส นิพฺพาปนาย, โกธสฺส นิพฺพาปนาย, อุปนาหสฺส นิพฺพาปนาย ฯเปฯ สพฺพากุสลาภิสงฺขารานํ สมาย อุปสมาย วูปสมาย นิพฺพาปนาย ปฏินิสฺสคฺคาย ปฏิปสฺสทฺธิยา อธิสีลมฺปิ สิกฺเขยฺย, อธิจิตฺตมฺปิ สิกฺเขยฺย, อธิปญฺญมฺปิ สิกฺเขยฺย, อิมา ติสฺโส \textbf{สิกฺขา}โย อาวชฺชนฺโต สิกฺเขยฺย, ชานนฺโต สิกฺเขยฺย, ปสฺสนฺโต สิกฺเขยฺย, ปจฺจเวกฺขนฺโต สิกฺเขยฺย, จิตฺตํ ปทหนฺโต สิกฺเขยฺย, สทฺธาย อธิมุจฺจนฺโต สิกฺเขยฺย, วีริยํ ปคฺคณฺหนฺโต สิกฺเขยฺย, สติํ อุปฏฺฐเปนฺโต สิกฺเขยฺย, จิตฺตํ สมาทหนฺโต สิกฺเขยฺย, ปญฺญาย ปชานนฺโต สิกฺเขยฺย, อภิญฺเญยฺยํ อภิชานนฺโต สิกฺเขยฺย, ปริญฺเญยฺยํ ปริชานนฺโต สิกฺเขยฺย, ปหาตพฺพํ ปชหนฺโต สิกฺเขยฺย, ภาเวตพฺพํ ภาเวนฺโต สิกฺเขยฺย, สจฺฉิกาตพฺพํ สจฺฉิกโรนฺโต สิกฺเขยฺย อาจเรยฺย สมาจเรยฺย สมาทาย วตฺเตยฺยาติ สิกฺเข นิพฺพาน\textbf{มตฺตโน}.\csromanspacer{} เตนาห โส พฺราหฺมโณ\csromandash{}}

\prose{ปุจฺฉามิ ตํ ภควา พฺรูหิ เมตํ, (อิจฺจายสฺมา โธตโก,)}

\prose{วาจาภิกงฺขามิ มเหสิ ตุยฺหํ.}

\csromangathameasure{ตว สุตฺวาน นิคฺโฆสํ, สิกฺเข นิพฺพานมตฺตโนติ.}
\csromangathagroup{\csromangathastanza{ตว สุตฺวาน นิคฺโฆสํ, สิกฺเข นิพฺพานมตฺตโนติ.}}

\csromanheader{2}{5. โธตกมาณวปุจฺฉานิทฺเทส}
\csromanhangingbegin
\hangingheaditem{31}{\csromanfolio{91}เตนหาตปฺปํ กโรหิ, (\textbf{โธตกา}ติ \textbf{ภควา},)}
\hangingbody{\textbf{อิเธว นิปโก สโต.}}
\hangingbody{อิโต สุตฺวาน นิคฺโฆสํ, สิกฺเข นิพฺพาน\textbf{มตฺตโน}.}
\csromanhangingend

\prose{\textbf{เตนหาตปฺปํ กโรหี}ติ อาตปฺปํ กโรหิ, อุสฺสาหํ กโรหิ, อุสฺโสฬฺหิํ กโรหิ, ถามํ กโรหิ, ธิติํ กโรหิ, วีริยํ กโรหิ, ฉนฺทํ ชเนหิ สญฺชเนหิ อุปฏฺฐเปหิ สมุฏฺฐเปหิ นิพฺพตฺเตหิ อภินิพฺพตฺเตหีติ เตนหาตปฺปํ กโรหิ.}

\prose{\textbf{โธตกา}ติ \textbf{ภควา} ตํ พฺราหฺมณํ นาเมน อาลปติ.\csromanspacer{} \textbf{ภควา}ติ คารวาธิวจนเมตํ ฯเปฯ สจฺฉิกา ปญฺญตฺติ, ยทิทํ \textbf{ภควา}ติ \textbf{โธตกา}ติ \textbf{ภควา}.}

\prose{\textbf{อิเธว นิปโก สโต}ติ \textbf{อิธา}ติ อิมิสฺสา ทิฏฺฐิยา อิมิสฺสา ขนฺติยา อิมิสฺสา รุจิยา อิมสฺมิํ อาทาเย อิมสฺมิํ ธมฺเม อิมสฺมิํ วินเย อิมสฺมิํ ธมฺมวินเย อิมสฺมิํ ปาวจเน อิมสฺมิํ พฺรหฺมจริเย อิมสฺมิํ สตฺถุสาสเน อิมสฺมิํ อตฺตภาเว อิมสฺมิํ มนุสฺสโลเก.\csromanspacer{} \textbf{นิปโก}ติ นิปโก ปณฺฑิโต ปญฺญวา พุทฺธิมา ญาณี วิภาวี เมธาวี.\csromanspacer{} \textbf{สโต}ติ จตูหิ การเณหิ สโต, กาเย กายานุปสฺสนาสติปฏฺฐานํ ภาเวนฺโต สโต ฯเปฯ โส วุจฺจติ สโตติ อิเธว นิปโก สโต.}

\prose{\textbf{อิโต สุตฺวาน นิคฺโฆสนฺ}ติ อิโต มยฺหํ วจนํ พฺยปฺปถํ เทสนํ อนุสาสนํ อนุสิฏฺฐํ สุตฺวา สุณิตฺวา อุคฺคณฺหิตฺวา อุปธารยิตฺวา อุปลกฺขยิตฺวาติ อิโต สุตฺวาน นิคฺโฆสํ.}

\prose{\textbf{สิกฺเข นิพฺพานมตฺตโน}ติ \textbf{สิกฺขา}ติ ติสฺโส \textbf{สิกฺขา} อธิสีล\textbf{สิกฺขา} อธิจิตฺต\textbf{สิกฺขา} อธิปญฺญา\textbf{สิกฺขา} ฯเปฯ อยํ อธิปญฺญา\textbf{สิกฺขา}.\csromanspacer{} นิพฺพาน\textbf{มตฺตโน}ติ อตฺตโน ราคสฺส นิพฺพาปนาย, โทสสฺส นิพฺพาปนาย, โมหสฺส นิพฺพาปนาย, โกธสฺส นิพฺพาปนาย, อุปนาหสฺส นิพฺพาปนาย ฯเปฯ สพฺพากุสลาภิสงฺขารานํ สมาย อุปสมาย วูปสมาย นิพฺพาปนาย ปฏินิสฺสคฺคาย ปฏิปสฺสทฺธิยา อธิสีลมฺปิ สิกฺเขยฺย, อธิจิตฺตมฺปิ สิกฺเขยฺย, อธิปญฺญมฺปิ สิกฺเขยฺย, อิมา ติสฺโส \textbf{สิกฺขา}โย อาวชฺชนฺโต สิกฺเขยฺย, ชานนฺโต สิกฺเขยฺย ฯเปฯ สจฺฉิกาตพฺพํ สจฺฉิกโรนฺโต สิกฺเขยฺย \csromanfolio{92}อาจเรยฺย สมาจเรยฺย, สมาทาย วตฺเตยฺยาติ สิกฺเข นิพฺพานมตฺ\textbf{ตนฺ}โอ.\csromanspacer{} เตนาห ภควา\csromandash{}}

\prose{เตนหาตปฺปํ กโรหิ, (โธตกาติ ภควา,)}

\prose{อิเธว นิปโก สโต.}

\csromangathameasure{อิโต สุตฺวาน นิคฺโฆสํ, สิกฺเข นิพฺพานมตฺตโนติ.}
\csromangathagroup{\csromangathastanza{อิโต สุตฺวาน นิคฺโฆสํ, สิกฺเข นิพฺพานมตฺตโนติ.}}

\csromanhangingbegin
\hangingheaditem{32}{\textbf{ปสฺสามหํ เทวมนุสฺสโลเก},}
\hangingbody{\textbf{อกิญฺจนํ พฺราหฺมณมิริยมานํ.}}
\hangingbody{\textbf{ตํ ตํ นมสฺสามิ สมนฺตจกฺขุ,}}
\hangingbody{ปมุญฺจ มํ สกฺก กถํกถาหิ.}
\csromanhangingend

\prose{\textbf{ปสฺสามหํ เทวมนุสฺสโลเก}ติ \textbf{เทวา}ติ ตโย \textbf{เทวา} สมฺมุติ\textbf{เทวา} อุปปตฺติ\textbf{เทวา} วิสุทฺธิ\textbf{เทวา}.\csromanspacer{} กตเม สมฺมุติ\textbf{เทวา}.\csromanspacer{} สมฺมุติ\textbf{เทวา} วุจฺจนฺติ ราชาโน จ ราชกุมารา จ เทวิโย จ, อิเม วุจฺจนฺติ สมฺมุติ\textbf{เทวา}.\csromanspacer{} กตเม อุปปตฺติ\textbf{เทวา}, อุปปตฺติ\textbf{เทวา} วุจฺจนฺติ จาตุมหาราชิกา \textbf{เทวา} ตาวติํสา \textbf{เทวา} ยามา \textbf{เทวา} ตุสิตา \textbf{เทวา} นิมฺมานรตี \textbf{เทวา} ปรนิมฺมิตวสวตฺตี \textbf{เทวา} พฺรหฺมกายิกา \textbf{เทวา}, เย จ \textbf{เทวา} ตทุตฺตริ\footnote{ตตฺรุปริ (สฺยา)}, อิเม วุจฺจนฺติ อุปปตฺติ\textbf{เทวา}.\csromanspacer{} กตเม วิสุทฺธิ\textbf{เทวา}, วิสุทฺธิ\textbf{เทวา} วุจฺจนฺติ ตถาคตสาวกา อรหนฺโต ขีณาสวา, เย จ ปจฺเจกพุทฺธา, อิเม วุจฺจนฺติ วิสุทฺธิ\textbf{เทวา}.\csromanspacer{} ภควา สมฺมุติ\textbf{เทวา}นญฺจ อุปปตฺติ\textbf{เทวา}นญฺจ วิสุทฺธิ\textbf{เทวา}นญฺจ เทโว จ อติเทโว จ เทวาติเทโว จ สีหสีโห นาคนาโค คณิคณี มุนิมุนี ราชราชา.\csromanspacer{} \textbf{ปสฺสามหํ เทวมนุสฺสโลเก}ติ มนุสฺสโลเก เทวํ ปสฺสามิ, อติเทวํ ปสฺสามิ, \textbf{เทวา}ติเทวํ ปสฺสามิ ทกฺขามิ โอโลเกมิ นิชฺฌายามิ อุปปริกฺขามีติ ปสฺสามหํ เทวมนุสฺสโลเก.}

\prose{\textbf{อกิญฺจนํ พฺราหฺมณมิริยมานนฺ}ติ \textbf{อกิญฺจนนฺ}ติ ราคกิญฺจนํ โทสกิญฺจนํ โมหกิญฺจนํ มานกิญฺจนํ ทิฏฺฐิกิญฺจนํ กิเลสกิญฺจนํ ทุจฺจริตกิญฺจนํ, เต กิญฺจนา พุทฺธสฺส ภควโต ปหีนา อุจฺฉินฺนมูลา ตาลาวตฺถุกตา อนภาวํกตา อายติํ อนุปฺปาทธมฺมา, ตสฺมา พุทฺโธ อกิญฺจโน.\csromanspacer{} \textbf{พฺราหฺมโณ}ติ ภควา สตฺตนฺนํ ธมฺมานํ พาหิตตฺตา พฺราหฺมโณ, สกฺกายทิฏฺฐิ พาหิตา โหติ, วิจิกิจฺฉา พาหิตา โหติ, สีลพฺพตปรามาโส พาหิโต โหติ, ราโค พาหิโต โหติ, โทโส}

\vspace{\tipitakaparskip}
\csromancenter{โธตกมาณวปุจฺฉานิทฺเทส พาหิโต โหติ, โมโห พาหิโต โหติ, มาโน พาหิโต โหติ, พาหิตาสฺส โหนฺติ ปาปกา อกุสลา ธมฺมา สํกิเลสิกา โปโนภวิกา สทรา ทุกฺขวิปากา อายติํ ชาติชรามรณิยา.}
\csromancenter{พาหิตฺวา สพฺพปาปกานิ, (สภิยาติ ภควา,)}
\prose{\csromanfolio{93}วิมโล สาธุสมาหิโต ฐิตตฺโต.}

\csromangathameasure{สํสารมติจฺจ เกวลี โส, \\ อสิโต ตาทิ ปวุจฺจเต ส พฺรหฺมาติ.}
\csromangathagroup{\csromangathastanza{สํสารมติจฺจ เกวลี โส, \\ อสิโต ตาทิ ปวุจฺจเต ส พฺรหฺมาติ.}}

\prose{\textbf{อิริยมานนฺ}ติ จรนฺตํ วิหรนฺตํ อิริยนฺตํ วตฺเตนฺตํ ปาเลนฺตํ ยเปนฺตํ ยาเปนฺ\textbf{ตนฺ}ติ อกิญฺจนํ พฺราหฺมณมิริยมานํ.}

\prose{\textbf{ตํ ตํ นมสฺสามิ สมนฺตจกฺขู}ติ \textbf{ตนฺ}ติ ภควนฺตํ ภณติ.\csromanspacer{} \textbf{นมสฺสามี}ติ กาเยน วา นมสฺสามิ, วาจาย วา นมสฺสามิ, จิตฺเตน วา นมสฺสามิ, อนฺวตฺถปฏิปตฺติยา วา นมสฺสามิ, ธมฺมานุธมฺมปฏิปตฺติยา วา นมสฺสามิ สกฺกโรมิ ครุํ กโรมิ มาเนมิ ปูเชมิ.\csromanspacer{} \textbf{สมนฺตจกฺขู}ติ สมนฺตจกฺขุ วุจฺจติ สพฺพญฺญุตญาณํ.\csromanspacer{} ภควา สพฺพญฺญุตญาเณน อุเปโต สมุเปโต อุปาคโต สมุปาคโต อุปปนฺโน สมุปปนฺโน สมนฺนาคโต.}

\csromangathameasure{“น ตสฺส อทฺทิฏฺฐมิธตฺถิ กิญฺจิ, \\ อโถ อวิญฺญาตมชานิตพฺพํ. \\ สพฺพํ อภิญฺญาสิ ยทตฺถิ เนยฺยํ, \\ ตถาคโต เตน \textbf{สมนฺตจกฺขู}”ติ.}
\csromangathagroup{\csromangathastanza{“น ตสฺส อทฺทิฏฺฐมิธตฺถิ\footnote{อทิฏฺฐมิธตฺถิ (สฺยา, ก) ขุ 7. 280 ปิฏฺเฐปิ.} กิญฺจิ, \\ อโถ อวิญฺญาตมชานิตพฺพํ. \\ สพฺพํ อภิญฺญาสิ ยทตฺถิ เนยฺยํ, \\ ตถาคโต เตน \textbf{สมนฺตจกฺขู}”ติ.}}

\prose{ตํ ตํ นมสฺสามิ สมนฺตจกฺขุ.}

\prose{\textbf{ปมุญฺจ มํ สกฺก กถํกถาหี}ติ \textbf{สกฺกา}ติ สกฺโก, ภควา สกฺยกุลา ปพฺพชิโตติปิ สกฺโก.\csromanspacer{} อถ วา อฑฺโฒ\footnote{อทฺโธ (สฺยา, ก)} มหทฺธโน ธนวาติปิ สกฺโก.\csromanspacer{} ตสฺสิมานิ ธนานิ.\csromanspacer{} เสยฺยถิทํ, สทฺธาธนํ สีลธนํ หิริธนํ โอตฺตปฺปธนํ สุตธนํ จาคธนํ ปญฺญาธนํ สติปฏฺฐานธนํ สมฺมปฺปธานธนํ อิทฺธิปาทธนํ อินฺทฺริยธนํ พลธนํ โพชฺฌงฺคธนํ มคฺคธนํ ผลธนํ นิพฺพานธนํ, อิเมหิ อเนกวิเธหิ ธนร\textbf{ตนฺ}เอหิ อฑฺโฒ มหทฺธโน ธนวาติปิ สกฺโก.\csromanspacer{} อถ วา สกฺโก ปหุ วิสวี อลมตฺโต สูโร วีโร วิกฺกนฺโต อภีรู อจฺฉมฺภี อนุตฺราสี อปลายี ปหีนภยเภรโว วิคตโลมหํโสติปิ \csromanfolio{94}สกฺโก.\csromanspacer{} กถํกถา วุจฺจติ วิจิกิจฺฉา, ทุกฺเข กงฺขา ทุกฺขสมุทเย กงฺขา ทุกฺขนิโรเธ กงฺขา ทุกฺขนิโรธคามินิยา ปฏิปทาย กงฺขา ปุพฺพนฺเต กงฺขา อปรนฺเต กงฺขา ปุพฺพนฺตาปรนฺเต กงฺขา อิทปฺปจฺจยตาปฏิจฺจสมุปฺปนฺเนสุ ธมฺเมสุ กงฺขา, ยา เอวรูปา กงฺขา กงฺขายนา กงฺขายิตตฺตํ วิมติ วิจิกิจฺฉา เทฺวฬฺหกํ เทฺวธาปโถ สํสโย อเนกํสคฺคาโห อาสปฺปนา ปริสปฺปนา อปริโยคาหนา ฉมฺภิตตฺตํ จิตฺตสฺส มโนวิเลโข.\csromanspacer{} \textbf{ปมุญฺจ มํ สกฺก กถํกถาหี}ติ มุญฺจ มํ ปมุญฺจ มํ โมเจหิ มํ ปโมเจหิ มํ อุทฺธร มํ สมุทฺธร มํ วุฏฺฐาเปหิ มํ กถํกถาสลฺลโตติ ปมุญฺจ มํ สกฺก กถํกถาหิ.\csromanspacer{} เตนาห โส พฺราหฺมโณ\csromandash{}}

\csromangathameasure{ปสฺสามหํ เทวมนุสฺสโลเก, \\ อกิญฺจนํ พฺราหฺมณมิริยมานํ. \\ ตํ ตํ นมสฺสามิ สมนฺตจกฺขุ, \\ \textbf{ปมุญฺจ มํ สกฺก กถํกถาหี}ติ.}
\csromangathagroup{\csromangathastanza{ปสฺสามหํ เทวมนุสฺสโลเก, \\ อกิญฺจนํ พฺราหฺมณมิริยมานํ. \\ ตํ ตํ นมสฺสามิ สมนฺตจกฺขุ, \\ \textbf{ปมุญฺจ มํ สกฺก กถํกถาหี}ติ.}}

\csromanhangingbegin
\hangingheaditem{33}{\textbf{นาหํ สหิสฺสามิ} ปโมจนาย,}
\hangingbody{\textbf{กถํกถิํ โธตก กญฺจิ โลเก.}}
\hangingbody{\textbf{ธมฺมญฺจ เสฏฺฐํ อาชานมาโน,}}
\hangingbody{เอวํ ตุวํ โอฆมิมํ ตเรสิ.}
\csromanhangingend

\prose{\textbf{นาหํ สหิสฺสามิ}\footnote{สมีหามิ (ก)} \textbf{ปโมจนายา}ติ นาหํ ตํ สกฺโกมิ มุญฺจิตุํ ปมุญฺจิตุํ โมเจตุํ ปโมเจตุํ อุทฺธริตุํ สมุทฺธริตุํ อุฏฺฐาเปตุํ สมุฏฺฐาเปตุํ กถํกถาสลฺลโตติ, เอวมฺปิ นาหํ สหิสฺสามิ ปโมจนาย.\csromanspacer{}\csromanspacer{}อถ วา น อีหามิ น สมีหามิ น อุสฺสหามิ น วายมามิ น อุสฺสาหํ กโรมิ น อุสฺโสฬฺหิํ กโรมิ น ถามํ กโรมิ น ธิติํ กโรมิ น วีริยํ กโรมิ น ฉนฺทํ ชเนมิ น สญฺชเนมิ น นิพฺพตฺเตมิ น อภินิพฺพตฺเตมิ อสฺสทฺเธ ปุคฺคเล อจฺฉนฺทิเก กุสีเต หีนวีริเย อปฺปฏิปชฺชมาเน ธมฺมเทสนายาติ, เอวมฺปิ นาหํ สหิสฺสามิ ปโมจนาย.\csromanspacer{}\csromanspacer{}อถ วา นตฺถญฺโญ โกจิ โมเจตา, เต ยทิ โมเจยฺยุํ, สเกน ถาเมน สเกน พเลน สเกน วีริเยน สเกน ปรกฺกเมน สเกน ปุริสถาเมน สเกน ปุริสพเลน สเกน ปุริสวีริเยน สเกน ปุริสปรกฺกเมน อตฺตนา}

\vspace{\tipitakaparskip}
\csromancenter{โธตกมาณวปุจฺฉานิทฺเทส สมฺมาปฏิปทํ อนุโลมปฏิปทํ อปจฺจนีกปฏิปทํ อนฺวตฺถปฏิปทํ ธมฺมานุธมฺมปฏิปทํ ปฏิปชฺชมานา โมเจยฺยุนฺติ, เอวมฺปิ นาหํ สหิสฺสามิ ปโมจนาย.}
\prose{\csromanfolio{95}วุตฺตเญฺหตํ ภควตา\csromandash{}โส วต จุนฺท อตฺตนา ปลิปปลิปนฺโน ปรํ ปลิปปลิปนฺนํ อุทฺธริสฺสตีติ เนตํ ฐานํ วิชฺชติ, โส วต จุนฺท อตฺตนา อทนฺโต อวินีโต อปรินิพฺพุโต ปรํ ทเมสฺสติ วิเนสฺสติ ปรินิพฺพาเปสฺสตีติ เนตํ ฐานํ วิชฺชตีติ, เอวมฺปิ นาหํ สหิสฺสามิ ปโมจนาย.\csromanspacer{} วุตฺตเญฺหตํ ภควตา\csromandash{}}

\csromangathameasure{“อตฺตนา หิ กตํ ปาปํ, อตฺตนา สํกิลิสฺสติ. \\ อตฺตนา อกตํ ปาปํ, อตฺตนาว วิสุชฺฌติ. \\ สุทฺธี อสุทฺธิ ปจฺจตฺตํ, นาญฺโญ อญฺญํ วิโสธเย”ติ.}
\csromangathagroup{\csromangathastanza{“อตฺตนา หิ\footnote{อตฺตนาว (พหูสุ) ขุ 1. 38 ปิฏฺเฐ.} กตํ ปาปํ, อตฺตนา สํกิลิสฺสติ. \\ อตฺตนา อกตํ ปาปํ, อตฺตนาว วิสุชฺฌติ. \\ สุทฺธี อสุทฺธิ ปจฺจตฺตํ, นาญฺโญ อญฺญํ วิโสธเย”ติ.}}

\prose{เอวมฺปิ นาหํ สหิสฺสามิ ปโมจนาย.}

\prose{วุตฺตเญฺหตํ ภควตา\csromandash{}เอวเมว โข พฺราหฺมณ ติฏฺฐเตว นิพฺพานํ, ติฏฺฐติ นิพฺพานคามิมคฺโค, ติฏฺฐามหํ สมาทเปตา, อถ จ ปน มม สาวกา มยา เอวํ โอวทิยมานา เอวํ อนุสาสิยมานา อปฺเปกจฺเจ อจฺจนฺตนิฏฺฐํ นิพฺพานํ อาราเธนฺติ, เอกจฺเจ นาราเธนฺตีติ, เอตฺถ กฺยาหํ พฺราหฺมณ กโรมิ, มคฺคกฺขายี พฺราหฺมณ ตถาคโต มคฺคํ พุทฺโธ อาจิกฺขติ อตฺตนา ปฏิปชฺชมานา มุจฺเจยฺยุนฺติ\footnote{มุญฺเจยฺยุนฺติ (สฺยา)}, เอวมฺปิ นาหํ สหิสฺสามิ ปโมจนาย.}

\prose{\textbf{กถํกถิํ โธตก กญฺจิ โลเก}ติ กถํกถิํ ปุคฺคลํ สกงฺขํ สขิลํ สเทฺวฬฺหกํ สวิจิกิจฺฉํ.\csromanspacer{} \textbf{กญฺจี}ติ กญฺจิ ขตฺติยํ วา พฺราหฺมณํ วา เวสฺสํ วา สุทฺทํ วา คหฏฺฐํ วา ปพฺพชิตํ วา เทวํ วา มนุสฺสํ วา.\csromanspacer{} \textbf{โลเก}ติ อปายโลเก ฯเปฯ อายตนโลเกติ กถํกถิํ โธตก กญฺจิ โลเก.}

\prose{\textbf{ธมฺมญฺจ เสฏฺฐํ อาชานมาโน}ติ ธมฺมํ เสฏฺฐํ วุจฺจติ อมตํ นิพฺพานํ.\csromanspacer{} โย โส สพฺพสงฺขารสมโถ สพฺพูปธิปฏินิสฺสคฺโค ตณฺหกฺขโย วิราโค นิโรโธ นิพฺพานํ.\csromanspacer{} \textbf{เสฏฺฐนฺ}ติ อคฺคํ เสฏฺฐํ วิเสฏฺฐํ ปาโมกฺขํ อุตฺตมํ ปวรํ ธมฺมํ อาชานมาโน วิชานมาโน ปฏิวิชานมาโน ปฏิวิชฺฌมาโนติ ธมฺมญฺจ เสฏฺฐํ อาชานมาโน.}

\prose{\csromanfolio{96}\textbf{เอวํ ตุวํ โอฆมิมํ ตเรสี}ติ เอวํ กาโมฆํ ภโวฆํ ทิฏฺโฐฆํ อวิชฺโชฆํ ตเรยฺยาสิ อุตฺตเรยฺยาสิ ปตเรยฺยาสิ สมติกฺกเมยฺยาสิ วีติวตฺเตยฺยาสีติ เอวํ ตุวํ โอฆมิมํ ตเรสิ.\csromanspacer{} เตนาห ภควา\csromandash{}}

\csromangathameasure{นาหํ สหิสฺสามิ ปโมจนาย\textbf{,} \\ กถํกถิํ โธตก กญฺจิ โลเก. \\ ธมฺมญฺจ เสฏฺฐํ อาชานมาโน\textbf{,}}
\csromangathagroup{\csromangathastanza{นาหํ สหิสฺสามิ ปโมจนาย\textbf{,} \\ กถํกถิํ โธตก กญฺจิ โลเก. \\ ธมฺมญฺจ เสฏฺฐํ อาชานมาโน\textbf{,}}}

\prose{\textbf{เอวํ ตุวํ โอฆมิมํ ตเรสี}ติ.}

\csromanhangingbegin
\hangingheaditem{34}{\textbf{อนุสาส พฺรหฺเม กรุณายมาโน,}}
\hangingbody{\textbf{วิเวกธมฺมํ ยมหํ วิชญฺญํ.}}
\hangingbody{\textbf{ยถาหํ อากาโสว1} \textbf{อพฺยาปชฺชมาโน2,}}
\hangingbody{อิเธว สนฺโต อสิโต จเรยฺยํ.}
\csromanhangingend

\prose{\textbf{อนุสาส พฺรหฺเม กรุณายมาโน}ติ อนุสาส พฺรหฺเม อนุคฺคณฺห พฺรหฺเม อนุกมฺป พฺรหฺเมติ อนุสาส พฺรหฺเม.\csromanspacer{} \textbf{กรุณายมาโน}ติ กรุณายมาโน อนุทยมาโน\footnote{อนุทฺทยมาโน (พหูสุ) 4. นตฺถิ ... สฺยา...โปตฺถเก.} อนุรกฺขมาโน อนุคฺคณฺหมาโน อนุกมฺปมาโนติ อนุสาส พฺรหฺเม กรูณายมาโน.}

\prose{\textbf{วิเวกธมฺมํ ยมหํ วิชญฺญนฺ}ติ วิเวกธมฺมํ วุจฺจติ อมตํ นิพฺพานํ.\csromanspacer{} โย โส สพฺพสงฺขารสมโถ สพฺพูปธิปฏินิสฺสคฺโค ตณฺหกฺขโย วิราโค นิโรโธ นิพฺพานํ.\csromanspacer{} \textbf{ยมหํ วิชญฺญนฺ}ติ ยมหํ ชาเนยฺยํ อาชาเนยฺยํ วิชาเนยฺยํ ปฏิวิชาเนยฺยํ ปฏิวิชฺเฌยฺยํ อธิคจฺเฉยฺยํ ผสฺเสยฺยํ สจฺฉิกเรยฺยนฺติ วิเวกธมฺมํ ยมหํ วิชญฺญํ.}

\prose{\textbf{ยถาหํ อากาโสว อพฺยาปชฺชมาโน}ติ ยถา อากาโส อ ปชฺชติ น คณฺหติ4 น พชฺฌติ น ปลิพชฺฌติ\textbf{,} เอวํ อปชฺชมาโน อคณฺหมาโน4 อพชฺฌมาโน อปลิพชฺฌมาโนติ เอวมฺปิ อากาโสว \textbf{อพฺยาปชฺชมาโน}.\csromanspacer{} ยถา อากาโส น รชฺชติ ลาขาย วา หลิทฺทิยา\footnote{หลิทฺเทน (สฺยา)} วา นีลิยา\footnote{นีเลน (สฺยา)} วา มญฺเชฏฺฐาย วา\textbf{,} เอวํ อรชฺชมาโน อทุสฺสมาโน อมุยฺหมาโน อกิลิสฺสมาโนติ\footnote{อกิลิยมาโน (สฺยา)} เอวมฺปิ อากาโสว \textbf{อพฺยาปชฺชมาโน}.\csromanspacer{} ยถา อากาโส}

\vspace{\tipitakaparskip}
\csromancenter{โธตกมาณวปุจฺฉานิทฺเทส น กุปฺปติ น พฺยาปชฺชติ น ปติลียติ\footnote{ปติฏฺฐิยติ (ก)} น ปฏิหญฺญติ, เอวํ อกุปฺปมาโน อพฺยาปชฺชมาโน อปฺปติลียมาโน อปฺปฏิหญฺญมาโน อปฺปฏิหตมาโนติ เอวมฺปิ อากาโสว อพฺยาปชฺชมาโน.}
\prose{\csromanfolio{97}\textbf{อิเธว สนฺโต อสิโต จเรยฺยนฺ}ติ \textbf{อิเธว สนฺโต}ติ \textbf{อิเธว สนฺโต} อิเธว สมาโน อิเธว นิสินฺโน สมาโน อิมสฺมิํเยว อาสเน นิสินฺโน สมาโน อิมิสฺสาเยว ปริสาย นิสินฺโน สมาโนติ เอวมฺปิ \textbf{อิเธว สนฺโต}.\csromanspacer{}\csromanspacer{}อถ วา \textbf{อิเธว สนฺโต} อุปสนฺโต วูปสนฺโต นิพฺพุโต ปฏิปฺปสฺสทฺโธติ เอวมฺปิ \textbf{อิเธว สนฺโต}.\csromanspacer{} \textbf{อสิโต}ติ เทฺว นิสฺสยา ตณฺหานิสฺสโย จ ทิฏฺฐินิสฺสโย จ ฯเปฯ อยํ ตณฺหานิสฺสโย ฯเปฯ อยํ ทิฏฺฐินิสฺสโย.\csromanspacer{} ตณฺหานิสฺสยํ ปหาย, ทิฏฺฐินิสฺสยํ ปฏินิสฺสชฺชิตฺวา, จกฺขุํ อนิสฺสิโต, โสตํ อนิสฺสิโต, ฆานํ อนิสฺสิโต, ชิวฺหํ อนิสฺสิโต, กายํ อนิสฺสิโต, มนํ อนิสฺสิโต, รูเป.\csromanspacer{} สทฺเท.\csromanspacer{} คนฺเธ.\csromanspacer{} รเส.\csromanspacer{} โผฏฺฐพฺเพ.\csromanspacer{} ธมฺเม.\csromanspacer{} กุลํ.\csromanspacer{} คณํ.\csromanspacer{} อาวาสํ.\csromanspacer{} ลาภํ.\csromanspacer{} ยสํ.\csromanspacer{} ปสํสํ.\csromanspacer{} สุขํ.\csromanspacer{} จีวรํ.\csromanspacer{} ปิณฺฑปาตํ.\csromanspacer{} เสนาสนํ.\csromanspacer{} คิลานปจฺจยเภสชฺชปริกฺขารํ.\csromanspacer{} กามธาตุํ.\csromanspacer{} รูปธาตุํ.\csromanspacer{} อรูปธาตุํ.\csromanspacer{} กามภวํ.\csromanspacer{} รูปภวํ.\csromanspacer{} อรูปภวํ.\csromanspacer{} สญฺญาภวํ.\csromanspacer{} อสญฺญาภวํ.\csromanspacer{} เนวสญฺญานาสญฺญาภวํ.\csromanspacer{} เอกโวการภวํ.\csromanspacer{} จตุโวการภวํ.\csromanspacer{} ปญฺจโวการภวํ.\csromanspacer{} อตีตํ.\csromanspacer{} อนาคตํ.\csromanspacer{} ปจฺจุปฺปนฺนํ.\csromanspacer{} ทิฏฺฐสุตมุตวิญฺญาตพฺเพ\footnote{ทิฏฺฐํ, สุตํ, มุตํ , วิญฺญาตํ, สพฺเพ. ขุ 7. 103 ปิฏฺเฐ, 189 ปิฏฺเฐปิ ปสฺสิตพฺพํ.} ธมฺเม อสิโต อนิสฺสิโต อนลฺลีโน อนุปคโต อนชฺโฌสิโต อนธิมุตฺโต นิกฺขนฺโต นิสฺสโฏ\footnote{นิสฺสฏฺโฐ (สฺยา)} วิปฺปมุตฺโต วิสํยุตฺโต วิมริยาทิกเตน เจตสา.\csromanspacer{} \textbf{จเรยฺยนฺ}ติ จเรยฺยํ วิหเรยฺยํ อิริเยยฺยํ วตฺเตยฺยํ ยเปยฺยํ ยาเปยฺยนฺติ \textbf{อิเธว สนฺโต} อสิโต จเรยฺยํ.\csromanspacer{} เตนาห โส พฺราหฺมโณ\csromandash{}}

\csromangathameasure{อนุสาส พฺรหฺเม กรุณายมาโน, \\ วิเวกธมฺมํ ยมหํ วิชญฺญํ. \\ ยถาหํ อากาโสว อพฺยาปชฺชมาโน,}
\csromangathagroup{\csromangathastanza{อนุสาส พฺรหฺเม กรุณายมาโน, \\ วิเวกธมฺมํ ยมหํ วิชญฺญํ. \\ ยถาหํ อากาโสว อพฺยาปชฺชมาโน,}}

\prose{\textbf{อิเธว สนฺโต อสิโต จเรยฺยนฺ}ติ.}

\csromanhangingbegin
\hangingheaditem{35}{กิตฺตยิสฺสามิ เต สนฺติํ, (โธตกาติ ภควา,)}
\hangingbody{\textbf{ทิฏฺเฐ ธมฺเม อนีติหํ.}}
\hangingbody{ยํ วิทิตฺวา สโต จรํ, ตเร โลเก วิสตฺติกํ.}
\csromanhangingend

\prose{\csromanfolio{98}\textbf{กิตฺตยิสฺสามิ เต สนฺตินฺ}ติ รฺ\textbf{อา}คสฺส สนฺติํ, โทสสฺส สนฺติํ, โมหสฺส สนฺติํ, โกธสฺส สนฺติํ, อุปนฺ\textbf{อา}หสฺส.\csromanspacer{} มกฺขสฺส.\csromanspacer{} ปฬฺ\textbf{อา}สสฺส.\csromanspacer{} อิสฺสฺ\textbf{อา}ย.\csromanspacer{} มจฺฉริยสฺส.\csromanspacer{} มฺ\textbf{อา}ยฺ\textbf{อา}ย.\csromanspacer{} สฺ\textbf{อา}เฐยฺยสฺส.\csromanspacer{} ถมฺภสฺส.\csromanspacer{} สฺ\textbf{อา}รมฺภสฺส.\csromanspacer{} มฺ\textbf{อา}นสฺส.\csromanspacer{} อติมฺ\textbf{อา}นสฺส.\csromanspacer{} มทสฺส.\csromanspacer{} ปมฺ\textbf{อา}ทสฺส.\csromanspacer{} สพฺพกิเลสฺ\textbf{อา}นํ.\csromanspacer{} สพฺพทุจฺจริตฺ\textbf{อา}นํ.\csromanspacer{} สพฺพทรถฺ\textbf{อา}นํ.\csromanspacer{} สพฺพปริฬฺ\textbf{อา}หฺ\textbf{อา}นํ.\csromanspacer{} สพฺพสนฺตฺ\textbf{อา}ปฺ\textbf{อา}นํ.\csromanspacer{} สพฺพฺ\textbf{อา}กุสลฺ\textbf{อา}ภิสงฺขฺ\textbf{อา}รฺ\textbf{อา}นํ สนฺติํ อุปสนฺติํ วูปสนฺติํ นิพฺพุติํ ปฏิปฺปสฺสทฺธิํ กิตฺตยิสฺสฺ\textbf{อา}มิ ปกิตฺตยิสฺสฺ\textbf{อา}มิ \textbf{อา}จิกฺขิสฺสฺ\textbf{อา}มิ เทเสสฺสฺ\textbf{อา}มิ ปญฺญเปสฺสฺ\textbf{อา}มิ ปฏฺฐเปสฺสฺ\textbf{อา}มิ วิวริสฺสฺ\textbf{อา}มิ วิภชิสฺสฺ\textbf{อา}มิ อุตฺตฺ\textbf{อา}นีกริสฺสฺ\textbf{อา}มิ ปกฺ\textbf{อา}สิสฺสฺ\textbf{อา}มีติ กิตฺตยิสฺสฺ\textbf{อา}มิ เต สนฺติํ.}

\prose{\textbf{โธตกาติ ภควา}ติ \textbf{โธตกาติ ภควา} ตํ พฺรฺ\textbf{อา}หฺมณํ นฺ\textbf{อา}เมน \textbf{อา}ลปติ.\csromanspacer{} ภควฺ\textbf{อา}ติ คฺ\textbf{อา}รวฺ\textbf{อา}ธิวจนเมตํ ฯเปฯ สจฺฉิกฺ\textbf{อา} ปญฺญตฺติ, ยทิทํ ภควฺ\textbf{อา}ติ \textbf{โธตกาติ ภควา}.}

\prose{\textbf{ทิฏฺเฐ ธมฺเม อนีติหนฺ}ติ \textbf{ทิฏฺเฐ ธมฺเม}ติ \textbf{ทิฏฺเฐ ธมฺเม} ญฺ\textbf{อา}เต ธมฺเม ตุลิเต ธมฺเม ตีริเต ธมฺเม วิภูเต ธมฺเม วิภฺ\textbf{อา}วิเต ธมฺเม “สพฺเพ สงฺขฺ\textbf{อา}รฺ\textbf{อา} อนิจฺจฺ\textbf{อา}”ติ ฯเปฯ. “ยํ กิญฺจิ สมุทยธมฺมํ, สพฺพํ ตํ นิโรธธมฺมนฺ”ติ \textbf{ทิฏฺเฐ ธมฺเม} ญฺ\textbf{อา}เต ธมฺเม ตุลิเต ธมฺเม ตีริเต ธมฺเม วิภฺ\textbf{อา}วิเต ธมฺเม วิภูเต ธมฺเมติ เอวมฺปิ \textbf{ทิฏฺเฐ ธมฺเม}.\csromanspacer{}\csromanspacer{}อถ วฺ\textbf{อา} ทุกฺเข ทิฏฺเฐ ทุกฺขํ กถยิสฺสฺ\textbf{อา}มิ, สมุทเย ทิฏฺเฐ สมุทยํ กถยิสฺสฺ\textbf{อา}มิ, มคฺเค ทิฏฺเฐ มคฺคํ กถยิสฺสฺ\textbf{อา}มิ, นิโรเธ ทิฏฺเฐ นิโรธํ กถยิสฺสฺ\textbf{อา}มีติ เอวมฺปิ \textbf{ทิฏฺเฐ ธมฺเม}.\csromanspacer{}\csromanspacer{}อถ วฺ\textbf{อา} สนฺทิฏฺฐิกํ อกฺ\textbf{อา}ลิกํ เอหิปสฺสิกํ โอปเนยฺยิกํ\footnote{โอปนยิกํ (สฺยา, ก)} ปจฺจตฺตํ เวทิตพฺพํ วิญฺญูหีติ เอวมฺปิ \textbf{ทิฏฺเฐ ธมฺเม}.\csromanspacer{} \textbf{อนีติหนฺ}ติ น อิติหีติหํ น อิติกิรฺ\textbf{อา}ย น ปรมฺปรฺ\textbf{อา}ย น ปิฏกสมฺปทฺ\textbf{อา}ย น ตกฺกเหตุ น นยเหตุ น \textbf{อา}กฺ\textbf{อา}รปริวิตกฺเกน น ทิฏฺฐินิชฺฌฺ\textbf{อา}นกฺขนฺติยฺ\textbf{อา} สฺ\textbf{อา}มํ สยมภิญฺญฺ\textbf{อา}ตํ อตฺตปจฺจกฺขธมฺมํ, ตํ กถยิสฺสฺ\textbf{อา}มีติ \textbf{ทิฏฺเฐ ธมฺเม} อนีติหํ.}

\prose{\textbf{ยํ วิทิตฺวา สโต จรนฺ}ติ ยํ วิทิตํ กตฺวฺ\textbf{อา} ตุลยิตฺวฺ\textbf{อา} ตีรยิตฺวฺ\textbf{อา} วิภฺ\textbf{อา}วยิตฺวฺ\textbf{อา} วิภูตํ กตฺวฺ\textbf{อา}, “สพฺเพ สงฺขฺ\textbf{อา}รฺ\textbf{อา} อนิจฺจฺ\textbf{อา}”ติ วิทิตํ กตฺวฺ\textbf{อา} ตุลยิตฺวฺ\textbf{อา} ตีรยิตฺวฺ\textbf{อา} วิภฺ\textbf{อา}วยิตฺวฺ\textbf{อา} วิภูตํ กตฺวฺ\textbf{อา}, “สพฺเพ สงฺขฺ\textbf{อา}รฺ\textbf{อา} ทุกฺขฺ\textbf{อา}”ติ ฯเปฯ. “สพฺเพ ธมฺมฺ\textbf{อา} อนตฺตฺ\textbf{อา}”ติ ฯเปฯ. “ยํ กิญฺจิ สมุทยธมฺมํ, สพฺพํ ตํ นิโรธธมฺมนฺ”ติ วิทิตํ กตฺวฺ\textbf{อา} ตุลยิตฺวฺ\textbf{อา} ตีรยิตฺวฺ\textbf{อา} วิภฺ\textbf{อา}วยิตฺวฺ\textbf{อา} วิภูตํ}

\vspace{\tipitakaparskip}
\csromancenter{โธตกมาณวปุจฺฉานิทฺเทส กตฺวา.\csromanspacer{} \textbf{สโต}ติ จตูหิ การเณหิ สโต, กาเย กายานุปสฺสนาสติปฏฺฐานํ ภาเวนฺโต สโต ฯเปฯ โส วุจฺจติ สโต.\csromanspacer{} \textbf{จรนฺ}ติ จรนฺโต วิหรนฺโต อิริยนฺโต วตฺเตนฺโต ปาเลนฺโต ยเปนฺโต ยาเปนฺโตติ ยํ วิทิตฺวา สโต จรํ.}
\prose{\csromanfolio{99}\textbf{ตเร โลเก วิสตฺติกนฺ}ติ วิสตฺติกา วุจฺจติ ตณฺหา.\csromanspacer{} โย ราโค สาราโค ฯเปฯ อภิชฺฌา โลโภ อกุสลมูลํ.\csromanspacer{} \textbf{วิสตฺติกา}ติ เกนฏฺเฐน วิสตฺติกา ฯเปฯ วิสฏา วิตฺถตาติ วิสตฺติกา.\csromanspacer{} \textbf{โลเก}ติ อปายโลเก ฯเปฯ อาย\textbf{ตนฺ}อโลเก.\csromanspacer{} \textbf{ตเร โลเก วิสตฺติกนฺ}ติ โลเก เวสา วิสตฺติกา, โลเก เวตํ วิสตฺติกํ สโต ตเรยฺย อุตฺตเรยฺย ปตเรยฺย สมติกฺกเมยฺย วีติวตฺเตยฺยาติ ตเร โลเก วิสตฺติกํ.\csromanspacer{} เตนาห ภควา\csromandash{}}

\prose{กิตฺตยิสฺสามิ เต สนฺติํ, (โธตกาติ ภควา,)}

\prose{ทิฏฺเฐ ธมฺเม อนีติหํ.}

\csromangathameasure{ยํ วิทิตฺวา สโต จรํ, ตเร โลเก วิสตฺติกนฺติ.}
\csromangathagroup{\csromangathastanza{ยํ วิทิตฺวา สโต จรํ, ตเร โลเก วิสตฺติกนฺติ.}}

\csromanhangingbegin
\hangingheaditem{36}{ตญฺจาหํ อภินนฺทามิ, มเหสิ สนฺติ’มุตฺตมํ.}
\hangingbody{ยํ วิทิตฺวา สโต จรํ, ตเร โลเก วิสตฺติกํ.}
\csromanhangingend

\prose{\textbf{ตญฺจาหํ อภินนฺทามี}ติ \textbf{ตนฺ}ติ ตุยฺหํ วจนํ พฺยปฺปถํ เทสนํ อนุสาสนํ อนุสิฏฺฐํ นนฺทามิ อภินนฺทามิ โมทามิ อนุโมทามิ อิจฺฉามิ สาทิยามิ ปตฺถยามิ ปิหยามิ อภิชปฺปามีติ ตญฺจาหํ อภินนฺทามิ.}

\prose{\textbf{มเหสิ สนฺติ’มุตฺตมนฺ}ติ \textbf{มเหสี}ติ กิํ มเหสิ ภควา, มหนฺตํ สีลกฺขนฺธํ เอสี คเวสี ปริเยสีติ มเหสิ, มหนฺตํ สมาธิกฺขนฺธํ ฯเปฯ “กหํ นราสโภ”ติ มเหสิ.\csromanspacer{} \textbf{สนฺติ’มุตฺตมนฺ}ติ สนฺติ วุจฺจติ อมตํ นิพฺพานํ.\csromanspacer{} โย โส สพฺพสงฺขารสมโถ สพฺพูปธิปฏินิสฺสคฺโค ตณฺหกฺขโย วิราโค นิโรโธ นิพฺพานํ.\csromanspacer{} \textbf{อุตฺตมนฺ}ติ อคฺคํ เสฏฺฐํ วิเสฏฺฐํ ปาโมกฺขํ อุตฺตมํ ปวรนฺติ มเหสิ สนฺติ’มุตฺตมํ.}

\prose{\textbf{ยํ วิทิตฺวา สโต จรนฺ}ติ ยํ วิทิตํ กตฺวา ฯเปฯ “สพฺเพ สงฺขารา อนิจฺจา”ติ วิทิตํ กตฺวา ตุลยิตฺวา ตีรยิตฺวา วิภาวยิตฺวา วิภูตํ กตฺวา, “สพฺเพ สงฺขารา ทุกฺขา”ติ. “สพฺเพ ธมฺมา อนตฺตา”ติ ฯเปฯ. “ยํ กิญฺจิ สมุทยธมฺมํ, สพฺพํ ตํ นิโรธธมฺมนฺ”ติ วิทิตํ กตฺวา ตุลยิตฺวา ตีรยิตฺวา \csromanfolio{100}วิภาวยิตฺวา วิภูตํ กตฺวา.\csromanspacer{} \textbf{สโต}ติ จตูหิ การเณหิ สโต, กาเย กายานุปสฺสนาสติปฏฺฐานํ ภเวนฺโต สโต ฯเปฯ โส วุจฺจติ สโต.\csromanspacer{} \textbf{จรนฺ}ติ จรนฺโต ฯเปฯ ยาเปนฺโตติ ยํ วิทิตฺวา สโต จรํ.}

\prose{\textbf{ตเร โลเก วิสตฺติกนฺ}ติ วิสตฺติกา วุจฺจติ ตณฺหา.\csromanspacer{} โย ราโค สาราโค ฯเปฯ อภิชฺฌา โลโภ อกุสลมูลํ.\csromanspacer{} \textbf{วิสตฺติกา}ติ เกนฏฺเฐน วิสตฺติกา ฯเปฯ วิสฏา วิตฺถตาติ วิสตฺติกา.\csromanspacer{} \textbf{โลเก}ติ อปายโลเก ฯเปฯ อายตนโลเก.\csromanspacer{} \textbf{ตเร โลเก วิสตฺติกนฺ}ติ โลเก เวสา วิสตฺติกา, โลเก เวตํ วิสตฺติกํ สโต ตเรยฺยํ อุตฺตเรยฺยํ ฯเปฯ วีติวตฺเตยฺยนฺติ ตเร โลเก วิสตฺติกํ.\csromanspacer{} เตนาห โส พฺราหฺมโณ\csromandash{}}

\csromangathameasure{ตญฺจาหํ อภินนฺทามิ, มเหสิ สนฺติ’มุตฺตมํ. \\ ยํ วิทิตฺวา สโต จรํ, ตเร โลเก วิสตฺติกนฺติ.}
\csromangathagroup{\csromangathastanza{ตญฺจาหํ อภินนฺทามิ, มเหสิ สนฺติ’มุตฺตมํ. \\ ยํ วิทิตฺวา สโต จรํ, ตเร โลเก วิสตฺติกนฺติ.}}

\csromanhangingbegin
\hangingheaditem{37}{ยํ กิญฺจิ สมฺปชานาสิ, (\textbf{โธตกาติ ภควา},)}
\hangingbody{\textbf{อุทฺธํ อโธ ติริยญฺจาปิ มชฺเฌ.}}
\hangingbody{\textbf{เอตํ วิทิตฺวา สงฺโคติ โลเก,}}
\hangingbody{ภวาภวาย มากาสิ ตณฺหํ.}
\csromanhangingend

\prose{\textbf{ยํ กิญฺจิ สมฺปชานาสี}ติ ยํ กิญฺจิ สมฺปชานาสิ อาชานาสิ ปฏิวิชานาสิ ปฏิวิชฺฌสีติ ยํ กิญฺจิ สมฺปชานาสิ.\csromanspacer{} \textbf{โธตกาติ ภควา}ติ \textbf{โธตกาติ ภควา} ตํ พฺราหฺมณํ นาเมน อาลปติ.\csromanspacer{} \textbf{ภควา}ติ คารวาธิวจนเมตํ ฯเปฯ สจฺฉิกา ปญฺญตฺติ, ยทิทํ \textbf{ภควา}ติ \textbf{โธตกาติ ภควา}.}

\prose{\textbf{อุทฺธํ อโธ ติริยญฺจาปิ มชฺเฌ}ติ \textbf{อุทฺธนฺ}ติ อนาคตํ.\csromanspacer{} \textbf{อโธ}ติ อตีตํ.\csromanspacer{} \textbf{ติริยญฺจาปิ มชฺเฌ}ติ ปจฺจุปฺปนฺนํ.\csromanspacer{} \textbf{อุทฺธนฺ}ติ เทวโลโก.\csromanspacer{} \textbf{อโธ}ติ อปายโลโก.\csromanspacer{} \textbf{ติริยญฺจาปิ มชฺเฌ}ติ มนุสฺสโลโก.\csromanspacer{} อถ วา \textbf{อุทฺธนฺ}ติ กุสลา ธมฺมา.\csromanspacer{} \textbf{อโธ}ติ อกุสลา ธมฺมา.\csromanspacer{} \textbf{ติริยญฺจาปิ มชฺเฌ}ติ อพฺยากตา ธมฺมา.\csromanspacer{} \textbf{อุทฺธนฺ}ติ อรูปธาตุ.\csromanspacer{} \textbf{อโธ}ติ กามธาตุ.\csromanspacer{} \textbf{ติริยญฺจาปิ มชฺเฌ}ติ รูปธาตุ.\csromanspacer{} \textbf{อุทฺธนฺ}ติ สุขา เวทนา.\csromanspacer{} \textbf{อโธ}ติ ทุกฺขา เวทนา.\csromanspacer{} \textbf{ติริยญฺจาปิ มชฺเฌ}ติ อทุกฺขมสุขา เวทนา.\csromanspacer{} \textbf{อุทฺธนฺ}ติ อุทฺธํ ปาทตลา.}

\vspace{\tipitakaparskip}
\csromancenter{อุปสีวมาณวปุจฺฉานิทฺเทส \textbf{อโธ}ติ อโธ เกสมตฺถกา.\csromanspacer{} \textbf{ติริยญฺจาปิ มชฺเฌ}ติ เวมชฺเฌติ อุทฺธํ อโธ ติริยญฺจาปิ มชฺเฌ.}
\prose{\csromanfolio{101}\textbf{เอตํ วิทิตฺวา สงฺโคติ โลเก}ติ สงฺโค เอโส, ลคฺคนํ เอตํ, พนฺธนํ เอตํ, ปลิโพโธ เอโสติ ญ ตฺวา ชานิตฺวา ตุลยิตฺวา ตีรยิตฺวา วิภาวยิตฺวา วิภูตํ กตฺวาติ เอตํ วิทิตฺวา สงฺโคติ โลเก.}

\prose{\textbf{ภวาภวาย มากาสิ ตณฺหนฺ}ติ \textbf{ตณฺหา}ติ รูป\textbf{ตณฺหา} สทฺท\textbf{ตณฺหา} ฯเปฯ ธมฺม\textbf{ตณฺหา}.\csromanspacer{} \textbf{ภวาภวายา}ติ ภวาภวาย กมฺมภวาย ปุนพฺภวาย กามภวาย, กมฺมภวาย กามภวาย ปุนพฺภวาย รูปภวาย, กมฺมภวาย รูปภวาย ปุนพฺภวาย อรูปภวาย, กมฺมภวาย อรูปภวาย ปุนพฺภวาย ปุนปฺปุนพฺภวาย, ปุนปฺปุนคติยา ปุนปฺปุน-อุปปตฺติยา ปุนปฺปุนปฏิสนฺธิยา ปุนปฺปุน- อตฺตภาวาภินิพฺพตฺติยา ตณฺหํ มากาสิ มา ชเนสิ มา สญฺชเนสิ มา นิพฺพตฺเตสิ มาภินิพฺพตฺเตสิ, ปชห วิโนเทหิ พฺยนฺตีกโรหิ อนภาวํ คเมหีติ ภวาภวาย มากาสิ ตณฺหนฺติ.\csromanspacer{} เตนาห ภควา\csromandash{}}

\prose{ยํ กิญฺจิ สมฺปชานาสิ, (โธตกาติ ภควา,)}

\prose{อุทฺธํ อโธ ติริยญฺจาปิ มชฺเฌ.}

\prose{\textbf{เอตํ วิทิตฺวา สงฺโคติ โลเก},}

\prose{\textbf{ภวาภวาย มากาสิ ตณฺหนฺ}ติ.}

\prose{สห คาถาปริโยสานา ฯเปฯ “สตฺถา เม ภนฺเต ภควา, สาวโกหมสฺมี”ติ.}

\vspace{\tipitakaparskip}
\csromancenter{โธตกมาณวปุจฺฉานิทฺเทโส ปญฺจโม.}
\csromanmatikaanchor{mtk.27}
\csromantocmarkref{subsection}{6. อุปสีวมาณวปุจฺฉานิทฺเทส}{mtk.27}
\bookmark[dest={mtk.27},level=2]{6. อุปสีวมาณวปุจฺฉานิทฺเทส}
\csromanheader{2}{6. อุปสีวมาณวปุจฺฉานิทฺเทส}
\csromanhangingbegin
\hangingheaditem{38}{เอโก อหํ สกฺก มหนฺตโมฆํ, (อิจฺจายสฺมา อุปสีโว,)}
\hangingbody{\textbf{อนิสฺสิโต โน วิสหามิ ตาริตุํ.}}
\hangingbody{\textbf{อารมฺมณํ1} \textbf{พฺรูหิ สมนฺตจกฺขุ,}}
\hangingbody{ยํ นิสฺสิโต โอฆมิมํ ตเรยฺยํ.}
\csromanhangingend

\prose{\csromanfolio{102}\textbf{เอโก อหํ สกฺก มหนฺตโมฆนฺ}ติ \textbf{เอโก}ติ ปุคฺคโล วา เม ทุติโย นตฺถิ, ธมฺโม วา เม ทุติโย นตฺถิ, ยํ วา ปุคฺคลํ นิสฺสาย, ธมฺมํ วา นิสฺสาย มหนฺตํ กาโมฆํ ภโวฆํ ทิฏฺโฐฆํ อวิชฺโชฆํ ตเรยฺยํ อุตฺตเรยฺยํ ปตเรยฺยํ สมติกฺกเมยฺยํ วีติวตฺเตยฺยนฺติ.\csromanspacer{} \textbf{สกฺกา}ติ สกฺโก, ภควา สกฺยกุลา ปพฺพชิโตติปิ สกฺโก.\csromanspacer{} อถ วา อฑฺโฒ มหทฺธโน ธนวาติปิ สกฺโก.\csromanspacer{} ตสฺสิมานิ ธนานิ.\csromanspacer{} เสยฺยถิทํ, สทฺธาธนํ สีลธนํ หิริธนํ โอตฺตปฺปธนํ สุตธนํ จาคธนํ ปญฺญาธนํ สติปฏฺฐานธนํ ฯเปฯ นิพฺพานธนํ.\csromanspacer{} อิเมหิ อเนเกหิ ธนรตเนหิ อฑฺโฒ มหทฺธโน ธนวาติปิ สกฺโก.\csromanspacer{} อถ วา สกฺโก ปหุ วิสวี อลมตฺโต สูโร วีโร วิกฺกนฺโต อภีรู อฉมฺภี อนุตฺราสี อปลายี, ปหีนภยเภรโว วิคตโลมหํโสติปิ สกฺโกติ \textbf{เอโก} อหํ สกฺก มหนฺตโมฆํ.}

\prose{\textbf{อิจฺจายสฺมา อุปสีโว}ติ \textbf{อิจฺจา}ติ ปทสนฺธิ ฯเปฯ.\csromanspacer{} \textbf{อายสฺมา}ติ ปิยวจนํ ฯเปฯ.\csromanspacer{} \textbf{อุปสีโว}ติ ตสฺส พฺราหฺมณสฺสนามํ ฯเปฯ อภิลาโปติ \textbf{อิจฺจา}ยสฺมา \textbf{อุปสีโว}.}

\prose{\textbf{อนิสฺสิโต โน วิสหามิ ตาริตุนฺ}ติ \textbf{อนิสฺสิโต}ติ ปุคฺคลํ วา \textbf{อนิสฺสิโต} ธมฺมํ วา \textbf{อนิสฺสิโต} โน วิสหามิ น อุสฺสหามิ น สกฺโกมิ น ปฏิพโล.\csromanspacer{} มหนฺตํ กาโมฆํ ภโวฆํ ทิฏฺโฐฆํ อวิชฺโชฆํ ตริตุํ อุตฺตริตุํ ปตริตุํ สมติกฺกมิตุํ วีติวตฺติตุนฺติ \textbf{อนิสฺสิโต} โน วิสหามิ ตาริตุํ.}

\prose{\textbf{อารมฺมณํ พฺรูหิ สมนฺตจกฺขู}ติ อารมฺมณํ อาลมฺพณํ นิสฺสยํ อุปนิสฺสยํ พฺรูหิ อาจิกฺขาหิ เทเสหิ ปญฺญเปหิ ปฏฺฐเปหิ วิวราหิ วิภชาหิ อุตฺตานีกโรหิ ปกาเสหิ.\csromanspacer{} \textbf{สมนฺตจกฺขู}ติ สมนฺตจกฺขุ วุจฺจติ สพฺพญฺญุตญาณํ.\csromanspacer{} ภควา เตน สพฺพญฺญุตญาเณน อุเปโต สมุเปโต อุปาคโต สมุปาคโต อุปปนฺโน สมุปปนฺโน สมนฺนาคโต.}

\csromangathameasure{น ตสฺส อทิฏฺฐมิธตฺถิ กิญฺจิ, \\ อโถ อวิญฺญาตมชานิตพฺพํ. \\ สพฺพํ อภิญฺญาสิ ยทตฺถิ เนยฺยํ, \\ ตถาคโต เตน \textbf{สมนฺตจกฺขู}ติ.}
\csromangathagroup{\csromangathastanza{น ตสฺส อทิฏฺฐมิธตฺถิ กิญฺจิ, \\ อโถ อวิญฺญาตมชานิตพฺพํ. \\ สพฺพํ อภิญฺญาสิ ยทตฺถิ เนยฺยํ, \\ ตถาคโต เตน \textbf{สมนฺตจกฺขู}ติ.}}

\prose{อารมฺมณํ พฺรูหิ สมนฺตจกฺขุ.}

\csromanheader{2}{6. อุปสีวมาณวปุจฺฉานิทฺเทส}
\prose{\csromanfolio{103}\textbf{ยํ นิสฺสิโต โอฆมิมํ ตเรยฺยนฺ}ติ \textbf{ยํ นิสฺสิโต}ติ ยํ ปุคฺคลํ วา นิสฺสิโต ธมฺมํ วา นิสฺสิโต มหนฺตํ กาโมฆํ ภโวฆํ ทิฏฺโฐฆํ อวิชฺโชฆํ ตเรยฺยํ อุตฺตเรยฺยํ ปตเรยฺยํ สมติกฺกเมยฺยํ วีติวตฺเตยฺยนฺติ \textbf{ยํ นิสฺสิโต} โอฆมิมํ ตเรยฺยํ.\csromanspacer{} เตนาห โส พฺราหฺมโณ\csromandash{}}

\prose{เอโก อหํ สกฺก มหนฺตโมฆํ, (อิจฺจายสฺมา อุปสีโว,)}

\prose{อนิสฺสิโต โน วิสหามิ ตาริตุํ.}

\prose{อารมฺมณํ พฺรูหิ สมนฺตจกฺขุ,}

\prose{\textbf{ยํ นิสฺสิโต โอฆมิมํ ตเรยฺยนฺ}ติ.}

\csromanhangingbegin
\hangingheaditem{39}{\textbf{อากิญฺจญฺญํ เปกฺขมาโน สติมา}, (อุปสีวาติ ภควา,)}
\hangingbody{\textbf{นตฺถีติ นิสฺสาย ตรสฺสุ โอฆํ.}}
\hangingbody{\textbf{กาเม ปหาย วิรโต กถาหิ,}}
\hangingbody{\textbf{ตณฺหกฺขยํ นตฺตมหา’ภิปสฺส1}.}
\csromanhangingend

\prose{\textbf{อากิญฺจญฺญํ เปกฺขมาโน สติมา}ติ โส พฺราหฺมโณ ปกติยา อากิญฺจญฺญายตนสมาปตฺติํ ลาภีเยว นิสฺสยํ น ชานาติ “อยํ เม นิสฺสโย”ติ.\csromanspacer{} ตสฺส ภควา นิสฺสยญฺจ อาจิกฺขติ อุตฺตริญฺจ นิยฺยานปถํ.\csromanspacer{} อากิญฺจญฺญายตนสมาปตฺติํ สโต สมาปชฺชิตฺวา ตโต วุฏฺฐหิตฺวา ตตฺถ ชาเต จิตฺตเจตสิเก ธมฺเม อนิจฺจโต เปกฺขมาโน, ทุกฺขโต.\csromanspacer{} โรคโต.\csromanspacer{} คณฺฑโต.\csromanspacer{} สลฺลโต.\csromanspacer{} อฆโต.\csromanspacer{} อาพาธโต.\csromanspacer{} ปรโต.\csromanspacer{} ปโลกโต. ¢ติโต.\csromanspacer{} อุปทฺทวโต.\csromanspacer{} ภยโต.\csromanspacer{} อุปสคฺคโต.\csromanspacer{} จลโต.\csromanspacer{} ปภงฺคุโต.\csromanspacer{} อทฺธุวโต.\csromanspacer{} อตาณโต.\csromanspacer{} อเลณโต.\csromanspacer{} อสรณโต.\csromanspacer{} อสรณีภูตโต.\csromanspacer{} ริตฺตโต.\csromanspacer{} ตุจฺฉโต.\csromanspacer{} สุญฺญโต.\csromanspacer{} อนตฺตโต.\csromanspacer{} อาทีนวโต.\csromanspacer{} วิปริณามธมฺมโต.\csromanspacer{} อสารกโต.\csromanspacer{} อฆมูลโต.\csromanspacer{} ภวโต.\csromanspacer{} วิภวโต.\csromanspacer{} สาสวโต.\csromanspacer{} สงฺขตโต.\csromanspacer{} มารามิสโต.\csromanspacer{} ชาติธมฺมโต.\csromanspacer{} ชราธมฺมโต.\csromanspacer{} พฺยาธิธมฺมโต.\csromanspacer{} มรณธมฺมโต.\csromanspacer{} โสกปริเทวทุกฺขโทมนสฺสุปายาสธมฺมโต.\csromanspacer{} สมุทยธมฺมโต.\csromanspacer{} อตฺถงฺคมโต.\csromanspacer{} อสฺสาทโต.\csromanspacer{} อาทีนวโต.\csromanspacer{} นิสฺสรณโต เปกฺขมาโน ทกฺขมาโน โอโลกยมาโน นิชฺฌายมาโน อุปปริกฺขมาโน.}

\prose{\csromanfolio{104}\textbf{สติมา}ติ ยา สติ อนุสฺสติ ปฏิสฺสติ ฯเปฯ สมฺมาสติ, อยํ วุจฺจติ สติ.\csromanspacer{} อิมาย สติยา อุเปโต โหติ ฯเปฯ สมนฺนาคโต.\csromanspacer{} โส วุจฺจติ สติมาติ อากิญฺจญฺญํ เปกฺขมาโน สติมา.}

\prose{\textbf{อุปสีวาติ ภควา}ติ \textbf{อุปสีวาติ ภควา} ตํ พฺราหฺมณํ นาเมน อาลปติ.\csromanspacer{} \textbf{ภควา}ติ คารวาธิวจนเมตํ ฯเปฯ สจฺฉิกา ปญฺญตฺติ, ยทิทํ \textbf{ภควา}ติ \textbf{อุปสีวาติ ภควา}.}

\prose{\textbf{นตฺถีติ นิสฺสาย ตรสฺสุ โอฆนฺ}ติ “นตฺถิ กิญฺจี”ติ อากิญฺจญฺญายตนสมาปตฺติ.\csromanspacer{} กิํการณา “นตฺถิ กิญฺจี”ติ อากิญฺจญฺญายตนสมาปตฺติ, วิญฺญาณญฺจายตนสมาปตฺติํ สโต สมาปชฺชิตฺวา ตโต วุฏฺฐหิตฺวา ตญฺเญว วิญฺญาณํ อภาเวติ วิภาเวติ อนฺตรธาเปติ “นตฺถิ กิญฺจี”ติ ปสฺสติ, ตํการณา “นตฺถิ กิญฺจี”ติ อากิญฺจญฺญายตนสมาปตฺติํ นิสฺสาย อุปนิสฺสาย อาลมฺพณํ กริตฺวา กาโมฆํ ภโวฆํ ทิฏฺโฐฆํ อวิชฺโชฆํ ตรสฺสุ อุตฺตรสฺสุ ปตรสฺสุ สมติกฺกมสฺสุ วีติวตฺตสฺสูติ นตฺถีติ นิสฺสาย ตรสฺสุ โอฆํ.}

\prose{\textbf{กาเม ปหาย วิรโต กถาหี}ติ \textbf{กามา}ติ อุทฺทานโต เทฺว \textbf{กามา} วตฺถุ\textbf{กามา} จ กิเลส\textbf{กามา} จ ฯเปฯ.\csromanspacer{} อิเม วุจฺจนฺติ วตฺถุ\textbf{กามา} ฯเปฯ.\csromanspacer{} อิเม วุจฺจนฺติ กิเลส\textbf{กามา}.\csromanspacer{} \textbf{กาเม ปหายา}ติ วตฺถุกาเม ปริชานิตฺวา กิเลสกาเม ปหาย ปชหิตฺวา วิโนเทตฺวา พฺยนฺตีกริตฺวา อนภาวํ คเมตฺวาติ กาเม ปหาย.\csromanspacer{} \textbf{วิรโต กถาหี}ติ กถํกถา วุจฺจติ วิจิกิจฺฉา.\csromanspacer{} ทุกฺเข กงฺขา ฯเปฯ ฉมฺภิตตฺตํ จิตฺตสฺส มโนวิเลโข, กถํกถาย อารโต วิรโต ปฏิวิรโต นิกฺขนฺโต นิสฺสโฏ วิปฺปมุตฺโต วิสญฺญุตฺโต วิมริยาทิกเตน เจตสา วิหรตีติ เอวมฺปิ วิรโต กถาหิ.\csromanspacer{}\csromanspacer{}อถ วา ทฺวตฺติํสาย ติรจฺฉานกถาย อารโต วิรโต ปฏิวิรโต นิกฺขนฺโต นิสฺสโฏ วิปฺปมุตฺโต วิสญฺญุตฺโต วิมริยาทิกเตน เจตสา วิหรตีติ เอวมฺปิ วิรโต กถาหีติ กาเม ปหาย วิรโต กถาหิ.}

\prose{\textbf{ตณฺหกฺขยํ นตฺตมหา’ภิปสฺสา}ติ \textbf{ตณฺหา}ติ รูป\textbf{ตณฺหา} ฯเปฯ ธมฺม\textbf{ตณฺหา}.\csromanspacer{} \textbf{นตฺตํ} วุจฺจติ รตฺติ.\csromanspacer{} \textbf{อโห}ติ ทิวโส.\csromanspacer{} รตฺติญฺจ ทิวา จ ตณฺหกฺขยํ ราคกฺขยํ โทสกฺขยํ โมหกฺขยํ คติกฺขยํ อุปปตฺติกฺขยํ ปฏิสนฺธิกฺขยํ ภวกฺขยํ}

\vspace{\tipitakaparskip}
\csromancenter{อุปสีวมาณวปุจฺฉานิทฺเทส สํสารกฺขยํ วฏฺฏกฺขยํ ปสฺส อภิปสฺส ทกฺข โอโลกย นิชฺฌาย อุปปริกฺขาติ ตณฺหกฺขยํ นตฺตมหา’ภิปสฺส.\csromanspacer{} เตนาห ภควา\csromandash{}}
\prose{\csromanfolio{105}อากิญฺจญฺญํ เปกฺขมาโน สติมา, (อุปสีวาติ ภควา,)}

\prose{นตฺถีติ นิสฺสาย ตรสฺสุ โอฆํ.}

\csromangathameasure{กาเม ปหาย วิรโต กถาหิ, \\ ตณฺหกฺขยํ นตฺตมหา’ภิปสฺสาติ.}
\csromangathagroup{\csromangathastanza{กาเม ปหาย วิรโต กถาหิ, \\ ตณฺหกฺขยํ นตฺตมหา’ภิปสฺสาติ.}}

\proseitem{40}{\textbf{สพฺเพสุ กาเมสุ โย วีตราโค}, (\textbf{อิจฺจา}ยสฺมา \textbf{อุปสีโว},)}

\prose{\textbf{อากิญฺจญฺญํ นิสฺสิโต หิตฺวา มญฺญํ.}}

\prose{\textbf{สญฺญาวิโมกฺเข ปรเมธิมุตฺโต,}}

\csromanheader{2}{\textbf{ติฏฺเฐ นุ โส ตตฺถ อนานุยายี.(3)}}
\prose{\textbf{สพฺเพสุ กาเมสุ โย วีตราโค}ติ \textbf{สพฺเพสู}ติ สพฺเพน สพฺพํ สพฺพถา สพฺพํ อเสสํ นิสฺเสสํ ปริยาทิยนวจนเมตํ \textbf{สพฺเพสู}ติ.\csromanspacer{} \textbf{กาเมสู}ติ \textbf{กามา}ติ อุทฺทานโต เทฺว \textbf{กามา} วตฺถุ\textbf{กามา} จ กิเลส\textbf{กามา} จ ฯเปฯ.\csromanspacer{} อิเม วุจฺจนฺติ วตฺถุ\textbf{กามา} ฯเปฯ.\csromanspacer{} อิเม วุจฺจนฺติ กิเลส\textbf{กามา}.\csromanspacer{} \textbf{สพฺเพสุ กาเมสุ โย วีตราโค}ติ สพฺเพสุ กาเมสุ โย วีตราโค วิคตราโค จตฺตราโค วนฺตราโค มุตฺตราโค ปหีนราโค ปฏินิสฺสฏฺฐราโค วิกฺขมฺภนโตติ สพฺเพสุ กาเมสุ โย วีตราโค.}

\prose{\textbf{อิจฺจายสฺมา อุปสีโว}ติ \textbf{อิจฺจา}ติ ปทสนฺธิ ฯเปฯ.\csromanspacer{} \textbf{อายสฺมา}ติ ปิยวจนํ ฯเปฯ.\csromanspacer{} \textbf{อุปสีโว}ติ ตสฺส พฺราหฺมณสฺส นามํ ฯเปฯ อภิลาโปติ \textbf{อิจฺจา}ยสฺมา \textbf{อุปสีโว}.}

\prose{\textbf{อากิญฺจญฺญํ นิสฺสิโต หิตฺวา มญฺญนฺ}ติ เหฏฺฐิมา ฉ สมาปตฺติโย หิตฺวา จชิตฺวา ปริจฺจชิตฺวา อติกฺกมิตฺวา สมติกฺกมิตฺวา วีติวตฺติตฺวา อากิญฺจญฺญายตนสมาปตฺติํ นิสฺสิโต อลฺลีโน อุปคโต สมุปคโต อชฺโฌสิโต อธิมุตฺโตติ อากิญฺจญฺญํ นิสฺสิโต หิตฺวา มญฺญํ.}

\prose{\textbf{สญฺญาวิโมกฺเข ปรเมธิมุตฺโต}ติ สญฺญาวิโมกฺขา วุจฺจนฺติ สตฺต สญฺญาสมาปตฺติโย.\csromanspacer{} ตาสํ สญฺญาสมาปตฺตีนํ อากิญฺจญฺญายตนสมาปตฺติวิโมกฺโข\footnote{วิโมกฺขา (ก) เอวมญฺเญสุ ปเทสุ พหุวจเนน.} อคฺโค จ เสฏฺโฐ จ วิเสฏฺโฐ จ ปาโมกฺโข จ อุตฺตโม จ ปวโร จ, ปรเม อคฺเค เสฏฺเฐ วิเสฏฺเฐ ปาโมกฺเข อุตฺตเม}

\prose{\csromanfolio{106}ปวเร อธิมุตฺติวิโมกฺเขน อธิมุตฺโต ตตฺราธิมุตฺโต ตทธิมุตฺโต ตจฺจริโต ตพฺพหุโล ตคฺครุโก ตนฺนินฺโน ตปฺโปโณ ตปฺปพฺภาโร ตทธิมุตฺโต ตทธิปเตยฺโยติ สญฺญาวิโมกฺเข ปรเมธิมุตฺโต.}

\prose{\textbf{ติฏฺเฐ นุ โส ตตฺถ อนานุยายี}ติ \textbf{ติฏฺเฐ นู}ติ สํสยปุจฺฉา วิมติปุจฺฉา เทฺวฬฺหกปุจฺฉา อเนกํสปุจฺฉา “เอวํ นุ โข, นนุ โข, กิํ นุ โข, กถํ นุ โข”ติ ติฏฺเฐ นุ.\csromanspacer{} \textbf{ตตฺถา}ติ อากิญฺจญฺญายตเน.\csromanspacer{} \textbf{อนานุยายี}ติ อนานุยายี อวิจฺจมาโน\footnote{อเวธมาโน (สฺยา)} อวิคจฺฉมาโน อนนฺตรธายมาโน อปริหายมาโน.\csromanspacer{}\csromanspacer{}อถ วา อรชฺชมาโน อทุสฺสมาโน อมุยฺหมาโน อกิลิสฺสมาโนติ ติฏฺเฐ นุ โส ตตฺถ อนานุยายี.\csromanspacer{} เตนาห โส พฺราหฺมโณ\csromandash{}}

\prose{\textbf{สพฺเพสุ กาเมสุ โย วีตราโค}, (อิจฺจายสฺมา อุปสีโว,)}

\prose{\textbf{อากิญฺจญฺญํ นิสฺสิโต หิตฺวา มญฺญํ.}}

\prose{\textbf{สญฺญาวิโมกฺเข ปรเมธิมุตฺโต,}}

\prose{\textbf{ติฏฺเฐ นุ โส ตตฺถ อนานุยายี}ติ.}

\csromanhangingbegin
\hangingheaditem{41}{\textbf{สพฺเพสุ กาเมสุ โย วีตราโค}, (\textbf{อุปสีวาติ ภควา},)}
\hangingbody{\textbf{อากิญฺจญฺญํ นิสฺสิโต หิตฺวา มญฺญํ.}}
\hangingbody{\textbf{สญฺญาวิโมกฺเข ปรเมธิมุตฺโต,}}
\hangingbody{ติฏฺเฐยฺย โส ตตฺถ อนานุยายี.}
\csromanhangingend

\prose{\textbf{สพฺเพสุ กาเมสุ โย วีตราโค}ติ \textbf{สพฺเพสู}ติ สพฺเพน สพฺพํ สพฺพถา สพฺพํ อเสสํ นิสฺเสสํ ปริยาทิยนวจนเมตํ \textbf{สพฺเพสู}ติ.\csromanspacer{} \textbf{กาเมสู}ติ \textbf{กามา}ติ อุทฺทานโต เทฺว \textbf{กามา} วตฺถุ\textbf{กามา} จ กิเลส\textbf{กามา} จ ฯเปฯ.\csromanspacer{} อิเม วุจฺจนฺติ วตฺถุ\textbf{กามา} ฯเปฯ.\csromanspacer{} อิเม วุจฺจนฺติ กิเลส\textbf{กามา}.\csromanspacer{} \textbf{สพฺเพสุ กาเมสุ โย วีตราโค}ติ สพฺเพสุ กาเมสุ โย วีตราโค ฯเปฯ ปฏินิสฺสฏฺฐราโค วิกฺขมฺภนโตติ สพฺเพสุ กาเมสุ โย วีตราโค.}

\prose{\textbf{อุปสีวาติ ภควา}ติ \textbf{อุปสีวาติ ภควา} ตํ พฺราหฺมณํ นาเมน อาลปติ.\csromanspacer{} \textbf{ภควา}ติ คารวาธิวจนเมตํ ฯเปฯ สจฺฉิกา ปญฺญตฺติ, ยทิทํ \textbf{ภควา}ติ \textbf{อุปสีวาติ ภควา}.}

\csromanheader{2}{6. อุปสีวมาณวปุจฺฉานิทฺเทส}
\prose{\csromanfolio{107}\textbf{อากิญฺจญฺญํ นิสฺสิโต หิตฺวา มญฺญนฺ}ติ เหฏฺฐิมา ฉ สมาปตฺติโย หิตฺวา จชิตฺวา ปริจฺจชิตฺวา อติกฺกมิตฺวา สมติกฺกมิตฺวา วีติวตฺติตฺวา อากิญฺจญฺญายตนสมาปตฺติํ นิสฺสิโต อลฺลีโน อุปคโต สมุปคโต อชฺโฌสิโต อธิมุตฺโตติ อากิญฺจญฺญํ นิสฺสิโต หิตฺวา มญฺญํ.}

\prose{\textbf{สญฺญาวิโมกฺเข ปรเมธิมุตฺโต}ติ สญฺญาวิโมกฺขา วุจฺจนฺติ สตฺต สญฺญาสมาปตฺติโย.\csromanspacer{} ตาสํ สญฺญาสมาปตฺตีนํ อากิญฺจญฺญายตนสมาปตฺติวิโมกฺโข อคฺโค จ เสฏฺโฐ จ วิเสฏฺโฐ จ ปาโมกฺโข จ อุตฺตโม จ ปวโร จ, ปรเม อคฺเค เสฏฺเฐ วิเสฏฺเฐ ปาโมกฺเข อุตฺตเม ปวเร อธิมุตฺติวิโมกฺเขน อธิมุตฺโต ตตฺราธิมุตฺโต ตทธิมุตฺโต ฯเปฯ ตทธิปเตยฺโยติ สญฺญาวิโมกฺเข ปรเมธิมุตฺโต.}

\prose{\textbf{ติฏฺฐยฺย โส ตตฺถ อนานุยายี}ติ \textbf{ติฏฺเฐยฺยา}ติ ติฏฺเฐยฺย สฏฺฐิกปฺปสหสฺสานิ.\csromanspacer{} \textbf{ตตฺถา}ติ อากิญฺจญฺญายตเน.\csromanspacer{} \textbf{อนานุยายี}ติ อนานุยายี อวิจฺจมาโน อวิคจฺฉมาโน อนนฺตรธายมาโน อปริหายมาโน.\csromanspacer{}\csromanspacer{}อถ วา อรชฺชมาโน อทุสฺสมาโน อมุยฺหมาโน อกิลิสฺสมาโนติ ติฏฺเฐยฺย โส ตตฺถ อนานุยายี.\csromanspacer{} เตนาห ภควา\csromandash{}}

\prose{สพฺเพสุ กาเมสุ โย วีตราโค, (อุปสีวาติ ภควา,)}

\prose{อากิญฺจญฺญํ นิสฺสิโต หิตฺวา มญฺญํ.}

\prose{\textbf{สญฺญาวิโมกฺเข ปรเมธิมุตฺโต},}

\prose{ติฏฺเฐยฺย โส ตตฺถ อนานุยายีติ.}

\csromanhangingbegin
\hangingheaditem{42}{\textbf{ติฏฺเฐ เจ โส ตตฺถ อนานุยายี},}
\hangingbody{\textbf{ปูคมฺปิ วสฺสานิ1} \textbf{สมนฺตจกฺขุ.}}
\hangingbody{\textbf{ตตฺเถว โส สีติสิยา วิมุตฺโต,}}
\hangingbody{จเวถ วิญฺญาณํ ตถาวิธสฺส.}
\csromanhangingend

\prose{\textbf{ติฏฺเฐ เจ โส ตตฺถ อนานุยายี}ติ สเจ โส ติฏฺเฐยฺย สฏฺฐิกปฺปสหสฺสานิ.\csromanspacer{} \textbf{ตตฺถา}ติ อากิญฺจญฺญายตเน.\csromanspacer{} \textbf{อนานุยายี}ติ อนานุยายี อวิจฺจมาโน อวิคจฺฉมาโน อนนฺตรธายมาโน อปริหายมาโน.\csromanspacer{} อถ วา อรชฺชมาโน อทุสฺสมาโน อมุยฺหมาโน อกิลิสฺสมาโนติ ติฏฺเฐ เจ โส ตตฺถ อนานุยายี.}

\prose{\csromanfolio{108}\textbf{ปูคมฺปิ วสฺสานิ สมนฺตจกฺขู}ติ \textbf{ปูคมฺปิ วสฺสานี}ติ ปูคมฺปิ วสฺสานิ พหูนิ วสฺสานิ\footnote{พหุนฺนํ วสฺสานํ (สฺยา)} พหูนิ วสฺสสตานิ พหูนิ วสฺสสหสฺสานิ พหูนิ วสฺสสตสหสฺสานิ พหูนิ กปฺปานิ พหูนิ กปฺปสตานิ พหูนิ กปฺปสหสฺสานิ พหูนิ กปฺปสตสหสฺสานิ.\csromanspacer{} \textbf{สมนฺตจกฺขู}ติ สมนฺตจกฺขุ วุจฺจติ สพฺพญฺญุตญาณํ ฯเปฯ ตถาคโต เตน \textbf{สมนฺตจกฺขู}ติ ปูคมฺปิ วสฺสานิ สมนฺตจกฺขุ.}

\prose{\textbf{ตตฺเถว โส สีติสิยา วิมุตฺโต, จเวถ วิญฺญาณํ ตถาวิธสฺสา}ติ ตตฺเถว โส สีติภาวมนุปฺปตฺโต นิจฺโจ ธุโว สสฺสโต อวิปริณามธมฺโม สสฺสติสมํ ตเถว ติฏฺเฐยฺย.\csromanspacer{} อถ วา ตสฺส วิญฺญาณํ จเวยฺย อุจฺฉิชฺเชยฺย นสฺเสยฺย วินสฺเสยฺย น ภเวยฺยาติ ปุนพฺภวปฏิสนฺธิวิญฺญาณํ นิพฺพตฺเตยฺย กามธาตุยา วา รูปธาตุยา วา อรูปธาตุยา วาติ อากิญฺจญฺญายตนํ สมาปนฺนสฺส สสฺสตญฺจ อุจฺเฉทญฺจ ปุจฺฉติ.\csromanspacer{} อุทาหุ ตตฺเถว อนุปาทิเสสาย นิพฺพานธาตุยา ปรินิพฺพาเยยฺย.\csromanspacer{} อถ วา ตสฺส วิญฺญาณํ จเวยฺย ปุน ปฏิสนฺธิวิญฺญาณํ นิพฺพตฺเตยฺย กามธาตุยา วา รูปธาตุยา วา อรูปธาตุยา วาติ อากิญฺจญฺญายตนํ อุปปนฺนสฺส ปรินิพฺพานญฺจ ปฏิสนฺธิญฺจ ปุจฺฉติ.\csromanspacer{} \textbf{ตถาวิธสฺสา}ติ ตถาวิธสฺส ตาทิสสฺส ตสฺสณฺฐิตสฺส ตปฺปการสฺส ตปฺปฏิภาคสฺส อากิญฺจญฺญายตนํ อุปปนฺนสฺสาติ ตตฺเถว โส สีติสิยา วิมุตฺโต, จเวถ วิญฺญาณํ ตถาวิธสฺส.\csromanspacer{} เตนาห โส พฺราหฺมโณ\csromandash{}}

\prose{ติฏฺเฐ เจ โส ตตฺถ อนานุยายี,}

\prose{ปูคมฺปิ วสฺสานิ สมนฺตจกฺขุ.}

\prose{ตตฺเถว โส สีติสิยา วิมุตฺโต,}

\prose{จเวถ วิญฺญาณํ ตถาวิธสฺสาติ.}

\csromanhangingbegin
\hangingheaditem{43}{\textbf{อจฺจิ ยถา วาตเวเคน ขิตฺตา}, (อุปสีวาติ ภควา,)}
\hangingbody{\textbf{อตฺถํ ปเลติ น อุเปติ สงฺขํ.}}
\hangingbody{\textbf{เอวํ มุนี นามกายา วิมุตฺโต,}}
\hangingbody{\textbf{อตฺถํ ปเลติ น อุเปติ สงฺขํ.}}
\csromanhangingend

\prose{\textbf{อจฺจิ ยถา วาตเวเคน ขิตฺตา}ติ อจฺจิ วุจฺจติ ชาลสิขา.\csromanspacer{} \textbf{วาตา}ติ ปุรตฺถิมา วาตา ปจฺฉิมา วาตา อุตฺตรา วาตา ทกฺขิณา วาตา สรชา วาตา อรชา วาตา สีตา วาตา อุณฺหา วาตา ปริตฺตา วาตา}

\vspace{\tipitakaparskip}
\csromancenter{อุปสีวมาณวปุจฺฉานิทฺเทส อธิมตฺตา วาตา\footnote{กาฬวาตา (ก)} เวรมฺภวาตา ปกฺขวาตา สุปณฺณวาตา ตาลปณฺณวาตา วิธูปนวาตา.\csromanspacer{} \textbf{วาตเวเคน ขิตฺตา}ติ วาตเวเคน ขิตฺตา\footnote{ขิตฺตํ (สฺยา) เอวมญฺเญสุ ปเทสุ นิคฺคหีตนฺตวเสน.} อุกฺขิตฺตา นุนฺนา ปนุนฺนา ขมฺภิตา วิกฺขมฺภิตาติ อจฺจิ ยถา วาตเวเคน ขิตฺตา.\csromanspacer{} \textbf{อุปสีวาติ ภควา}ติ \textbf{อุปสีวาติ ภควา} ตํ พฺราหฺมณํ นาเมน อาลปติ.\csromanspacer{} \textbf{ภควา}ติ คารวาธิวจนเมตํ ฯเปฯ สจฺฉิกา ปญฺญตฺติ, ยทิทํ \textbf{ภควา}ติ \textbf{อุปสีวาติ ภควา}.}
\prose{\csromanfolio{109}\textbf{อตฺถํ ปเลติ น อุเปติ สงฺขนฺ}ติ \textbf{อตฺถํ ปเลตี}ติ อตฺถํ ปเลติ, อตฺถํ คเมติ, อตฺถํ คจฺฉติ นิรุชฺฌติ วูปสมติ ปฏิปฺปสฺสมฺภติ.\csromanspacer{} \textbf{น อุเปติ สงฺขนฺ}ติ สงฺขํ\footnote{อมุกํ นาม ทิสํ คโตติ สงฺขํ (สฺยา)} น อุเปติ, อุทฺเทสํ น อุเปติ, คณนํ น อุเปติ, ปณฺณตฺติํ น อุเปติ, “ปุรตฺถิมํ วา ทิสํ คตา, ปจฺฉิมํ วา ทิสํ คตา, อุตฺตรํ วา ทิสํ คตา, ทกฺขิณํ วา ทิสํ คตา, อุทฺธํ วา คตา, อโธ วา คตา, ติริยํ วา คตา, วิทิสํ วา คตา”ติ, โส เหตุ นตฺถิ, ปจฺจโย นตฺถิ, การณํ นตฺถิ, เยน สงฺขํ คจฺเฉยฺยาติ อตฺถํ ปเลติ น อุเปติ สงฺขํ.}

\prose{\textbf{เอวํ มุนี นามกายา วิมุตฺโต}ติ \textbf{เอวนฺ}ติ โอปมฺมสมฺปฏิปาทนํ.\csromanspacer{} \textbf{มุนี}ติ โมนํ วุจฺจติ ญาณํ ฯเปฯ สงฺคชาลมติจฺจ โส มุนิ.\csromanspacer{} \textbf{นามกายา วิมุตฺโต}ติ โส มุนิ ปกติยา ปุพฺเพว รูปกายา วิมุตฺโต.\csromanspacer{} ตทงฺคํ สมติกฺกมา\footnote{ตทงฺคํ สมติกฺกมฺม (ก)} วิกฺขมฺภนปฺปหาเนน ปหีโน.\csromanspacer{} ตสฺส มุนิโน ภวนฺตํ อาคมฺม จตฺตาโร อริยมคฺคา ปฏิลทฺธา โหนฺติ, จตุนฺนํ อริยมคฺคานํ ปฏิลทฺธตฺตา นามกาโย จ รูปกาโย จ ปริญฺญาตา โหนฺติ.\csromanspacer{} นามกายสฺส จ รูปกายสฺส จ ปริญฺญาตตฺตา นามกายา จ รูปกายา จ มุตฺโต วิมุตฺโต สุวิมุตฺโต อจฺจนฺต-อนุปาทาวิโมกฺเขนาติ เอวํ มุนี นามกายา วิมุตฺโต.}

\prose{\textbf{อตฺถํ ปเลติ น อุเปติ สงฺขนฺ}ติ \textbf{อตฺถํ ปเลตี}ติ อนุปาทิเสสาย นิพฺพานธาตุยา ปรินิพฺพายติ.\csromanspacer{} \textbf{น อุเปติ สงฺขนฺ}ติ อนุปาทิเสสาย นิพฺพานธาตุยา ปรินิพฺพุโต สงฺขํ น อุเปติ, อุทฺเทสํ น อุเปติ, คณนํ น อุเปติ, ปณฺณตฺติํ น อุเปติ “ขตฺติโย”ติ วา “พฺราหฺมโณ”ติ วา “เวสฺโส”ติ วา “สุทฺโท”ติ วา “คหฏฺโฐ”ติ วา “ปพฺพชิโต”ติ วา “เทโว”ติ วา “มนุสฺโส”ติ วา “รูปี”ติ วา “อรูปี”ติ วา “สญฺญี”ติ เวตฺ “อสญฺญี”ติ วา “เนวสญฺญีนาสญฺญี”ติ วา, โส เหตุ นตฺถิ, ปจฺจโย นตฺถิ, การณํ \csromanfolio{110}นตฺถิ, เยน สงฺขํ คจฺเฉยฺยาติ อตฺถํ ปเลติ น อุเปติ สงฺขํ.\csromanspacer{} เตนาห \textbf{ภควา}\csromandash{}}

\prose{อจฺจิ ยถา วาตเวเคน ขิตฺตา, (อุปสีวาติ ภควา,)}

\prose{อตฺถํ ปเลติ น อุเปติ สงฺขํ.}

\csromangathameasure{เอวํ \textbf{มุนี} นามกายา วิมุตฺโต, \\ อตฺถํ ปเลติ น อุเปติ สงฺขนฺติ.}
\csromangathagroup{\csromangathastanza{เอวํ \textbf{มุนี} นามกายา วิมุตฺโต, \\ อตฺถํ ปเลติ น อุเปติ สงฺขนฺติ.}}

\proseitem{44}{อตฺถงฺคโต โส อุท วา โส นตฺถิ,}

\prose{\textbf{อุทาหุ เว สสฺสติยา อโรโค.}}

\prose{\textbf{ตํ เม มุนี สาธุ วิยากโรหิ,}}

\prose{\textbf{อตฺถงฺคโต โส อุท วา โส นตฺถี}ติ โส อตฺถงฺคโต อุทาหุ นตฺถิ โส นิรุทฺโธ อุจฺฉินฺโน วินฏฺโฐติ อตฺถงฺคโต โส อุท วา โส นตฺถิ.}

\prose{\textbf{อุทาหุ เว สสฺสติยา อโรโค}ติ อุทาหุ นิจฺโจ ธุโว สสฺสโต อวิปริณามธมฺโม สสฺสติสมํ ตเถว ติฏฺเฐยฺยาติ อุทาหุ เว สสฺสติยา อโรโค.}

\prose{\textbf{ตํ เม มุนี สาธุ วิยากโรหี}ติ \textbf{ตนฺ}ติ ยํ ปุจฺฉามิ, ยํ ยาจามิ, ยํ อชฺเฌสามิ, ยํ ปสาเทมิ.\csromanspacer{} \textbf{มุนี}ติ โมนํ วุจฺจติ ญาณํ ฯเปฯ สงฺคชาลมติจฺจ โส มุนิ.\csromanspacer{} \textbf{สาธุ วิยากโรหี}ติ สาธุ อาจิกฺขาหิ เทเสหิ ปญฺญเปหิ ปฏฺฐเปหิ วิวราหิ วิภชาหิ อุตฺตานีกโรหิ ปกาเสหีติ ตํ เม มุนี สาธุ วิยากโรหิ.}

\prose{\textbf{ตถา หิ เต วิทิโต เอส ธมฺโม}ติ ตถา หิ เต วิทิโต ตุลิโต ตีริโต วิภูโต วิภาวิโต เอส ธมฺโมติ, ตถา หิ เต วิทิโต เอส ธมฺโม.\csromanspacer{} เตนาห โส พฺราหฺมโณ\csromandash{}}

\prose{อตฺถงฺคโต โส อุท วา โส นตฺถิ,}

\prose{\textbf{อุทาหุ เว สสฺสติยา อโรโค.}}

\csromangathameasure{\textbf{ตํ เม มุนี สาธุ วิยากโรหิ,} \\ \textbf{ตถา หิ เต วิทิโต เอส ธมฺโม}ติ.}
\csromangathagroup{\csromangathastanza{\textbf{ตํ เม มุนี สาธุ วิยากโรหิ,} \\ \textbf{ตถา หิ เต วิทิโต เอส ธมฺโม}ติ.}}

\csromanheader{2}{6. อุปสีวมาณวปุจฺฉานิทฺเทส}
\csromanhangingbegin
\hangingheaditem{45}{\csromanfolio{111}อตฺถงฺคตสฺส น ปมาณ’มตฺถิ, (\textbf{อุปสีวาติ ภควา},)}
\hangingbody{\textbf{เยน นํ วชฺชุํ ตํ ตสฺส นตฺถิ.}}
\hangingbody{\textbf{สพฺเพสุ ธมฺเมสุ สมูหเตสุ,}}
\hangingbody{สมูหตา วาทปถาปิ สพฺเพ.}
\csromanhangingend

\prose{\textbf{อตฺถงฺคตสฺส น ปมาณ’มตฺถี}ติ อตฺถงฺคตสฺส อนุปาทิเสสาย นิพฺพานธาตุยา ปรินิพฺพุตสฺส รูปปมาณํ นตฺถิ.\csromanspacer{} เวทนาปมาณํ นตฺถิ.\csromanspacer{} สญฺญาปมาณํ นตฺถิ.\csromanspacer{} สงฺขารปมาณํ นตฺถิ.\csromanspacer{} วิญฺญาณปมาณํ นตฺถิ น สติ น สํวิชฺชติ นุปลพฺภติ ปหีนํ สมุจฺฉินฺนํ วูปสนฺตํ ปฏิปฺปสฺสทฺธํ อภพฺพุปฺปตฺติกํ ญาณคฺคินา ทฑฺฒนฺติ อตฺถงฺคตสฺส น ปมาณ’มตฺถิ.\csromanspacer{} \textbf{อุปสีวาติ ภควา}ติ \textbf{อุปสีวาติ ภควา} ตํ พฺราหฺมณํ นาเมน อาลปติ.\csromanspacer{} \textbf{ภควา}ติ คารวาธิวจนเมตํ ฯเปฯ สจฺฉิกา ปญฺญตฺติ, ยทิทํ \textbf{ภควา}ติ \textbf{อุปสีวาติ ภควา}.}

\prose{\textbf{เยน นํ วชฺชุํ ตํ ตสฺส นตฺถี}ติ เยน ตํ ราเคน\footnote{เยน ราเคน (สฺยา, ก) ขุ 7. 192 ปิฏฺเฐ.} วเทยฺยุํ.\csromanspacer{} เยน โทเสน วเทยฺยุํ.\csromanspacer{} เยน โมเหน วเทยฺยุํ.\csromanspacer{} เยน มาเนน วเทยฺยุํ.\csromanspacer{} ยาย ทิฏฺฐิยา วเทยฺยุํ.\csromanspacer{} เยน อุทฺธจฺเจน วเทยฺยุํ.\csromanspacer{} ยาย วิจิกิจฺฉาย วเทยฺยุํ.\csromanspacer{} เยหิ อนุสเยหิ วเทยฺยุํ “รตฺโต”ติ วา “ทุฏฺโฐ”ติ วา “มูโฬฺห”ติ วา “วินิพทฺโธ”ติ วา “ปรามฏฺโฐ”ติ วา “วิกฺเขปคโต”ติ วา “อนิฏฺฐงฺคโต”ติ\footnote{อนิฏฺฐา คโตติ (ก)} วา “ถามคโต”ติ วา, เต อภิสงฺขารา ปหีนา, อภิสงฺขารานํ ปหีนตฺตา คติยา เยน ตํ วเทยฺยุํ “เนรยิโก”ติ วา “ติรจฺฉานโยนิโก”ติ วา “เปตฺติวิสยิโก”ติ วา “มนุสฺโส”ติ วา “เทโว”ติ วา “รูปี”ติ วา “อรูปี”ติ วา “สญฺญี”ติ วา “อสญฺญี”ติ วา “เนวสญฺญีนาสญฺญี”ติ วา, โส เหตุ นตฺถิ, ปจฺจโย นตฺถิ, การณํ นตฺถิ.\csromanspacer{} เยน วเทยฺยุํ กเถยฺยุํ ภเณยฺยุํ ทีเปยฺยุํ โวหเรยฺยุนฺติ เยน นํ วชฺชุํ ตํ ตสฺส นตฺถิ.}

\prose{\textbf{สพฺเพสุ ธมฺเมสุ สมูหเตสู}ติ สพฺเพสุ ธมฺเมสุ สพฺเพสุ ขนฺเธสุ สพฺเพสุ อายตเนสุ สพฺพาสุ ธาตูสุ สพฺพาสุ คตีสุ สพฺพาสุ อุปปตฺตีสุ สพฺพาสุ ปฏิสนฺธีสุ สพฺเพสุ ภเวสุ สพฺเพสุ สํสาเรสุ สพฺเพสุ วฏฺเฏสุ อูหเตสุ สมูหเตสุ อุทฺธเตสุ สมุทฺธเตสุ อุปฺปาฏิเตสุ สมุปฺปาฏิเตสุ ปหีเนสุ สมุจฺฉินฺเนสุ วูปสนฺเตสุ ปฏิปฺปสฺสทฺเธสุ อภพฺพุปฺปตฺติเกสุ ญาณคฺคินา ทฑฺเฒสูติ สพฺเพสุ ธมฺเมสุ สมูหเตสุ.}

\prose{\csromanfolio{112}\textbf{สมูหตา วาทปถาปิ สพฺเพ}ติ วาทปถา วุจฺจนฺติ กิเลสา จ ขนฺธา จ อภิสงฺขารา จ.\csromanspacer{} ตสฺส วาทา จ วาทปถา จ อธิวจนานิ จ อธิวจนปถา จ นิรุตฺติ จ นิรุตฺติปถา จ ปญฺญตฺติ จ ปญฺญตฺติปถา จ อูหตา สมูหตา อุทฺธตา สมุทฺธตา อุปฺปาฏิตา สมุปฺปาฏิตา ปหีนา สมุจฺฉินฺนา วูปสนฺตา ปฏิปฺปสฺสทฺธา อภพฺพุปฺปตฺติกา ญาณคฺคินา ทฑฺฑาติ สมูหตา วาทปถาปิ สพฺเพ.\csromanspacer{} เตนาห ภควา\csromandash{}}

\prose{อตฺถงฺคตสฺส น ปมาณมตฺถิ, (อุปสีวาติ ภควา,)}

\prose{เยน นํ วชฺชุํ ตํ ตสฺส นตฺถิ.}

\prose{สพฺเพสุ ธมฺเมสุ สมูหเตสุ,}

\prose{\textbf{สมูหตา วาทปถาปิ สพฺเพ}ติ.}

\prose{สห คาถาปริโยสานา เย เต พฺราหฺมเณน สทฺธิํ ฯเปฯ ปญฺชลิโก นมสฺสมาโน นิสินฺโน โหติ “สตฺถา เม ภนฺเต ภควา, สาวโกหมสฺมี”ติ.}

\vspace{\tipitakaparskip}
\csromancenter{อุปสีวมาณวปุจฺฉานิทฺเทโส ฉฏฺโฐ.}
\csromanmatikaanchor{mtk.28}
\csromantocmarkref{subsection}{7. นนฺทมาณวปุจฺฉานิทฺเทส}{mtk.28}
\bookmark[dest={mtk.28},level=2]{7. นนฺทมาณวปุจฺฉานิทฺเทส}
\csromanheader{2}{7. นนฺทมาณวปุจฺฉานิทฺเทส}
\csromanhangingbegin
\hangingheaditem{46}{“สนฺติ โลเก มุนโย”, (\textbf{อิจฺจา}ยสฺมา \textbf{นนฺโท},)}
\hangingbody{\textbf{ชนา วทนฺติ ตยิทํ กถํสุ.}}
\hangingbody{\textbf{ญาณูปปนฺนํ มุนิ โน วทนฺติ,}}
\hangingbody{\textbf{อุทาหุ เว ชีวิเตนูปปนฺนํ1}.}
\csromanhangingend

\prose{\textbf{สนฺติ โลเก มุนโย}ติ \textbf{สนฺตี}ติ สนฺติ สํวิชฺชนฺติ อตฺถิ อุปลพฺภนฺติ.\csromanspacer{} \textbf{โลเก}ติ อปายโลเก ฯเปฯ อายตนโลเก.\csromanspacer{} \textbf{มุนโย}ติ มุนินามกา อาชีวกา นิคณฺฐา ชฏิลา ตาปสา (เทวา โลเก มุนโยติ สญฺชานนฺติ, น จ เต มุนโย)2ติ สนฺติ โลเก มุนโย.\csromanspacer{} \textbf{อิจฺจายสฺมา นนฺโท}ติ \textbf{อิจฺจา}ติ ปทสนฺธิ ฯเปฯ.\csromanspacer{} \textbf{อายสฺมา}ติ ปิยวจนํ ฯเปฯ.\csromanspacer{} \textbf{นนฺโท}ติ ตสฺส พฺราหฺมณสฺส นามํ ฯเปฯ อภิลาโปติ \textbf{อิจฺจา}ยสฺมา \textbf{นนฺโท}.}

\csromanheader{2}{7. นนฺทมาณวปุจฺฉานิทฺเทส}
\prose{\csromanfolio{113}\textbf{ชนา วทนฺติ ตยิทํ กถํสู}ติ \textbf{ชนา}ติ ขตฺติยา จ พฺราหฺมณา จ เวสฺสา จ สุทฺทา จ คหฏฺฐา จ ปพฺพชิตา จ เทวา จ มนุสฺสา จ.\csromanspacer{} \textbf{วทนฺตี}ติ กเถนฺติ ภณนฺติ ทีปยนฺติ โวหรนฺติ.\csromanspacer{} \textbf{ตยิทํ กถํสู}ติ สํสยปุจฺฉา วิมติปุจฺฉา เทฺวฬฺหกปุจฺฉา อเนกํสปุจฺฉา “เอวํ นุ โข, น นุ โข, กิํ นุ โข, กถํ นุ โข”ติ \textbf{ชนา} วทนฺติ ตยิทํ กถํสุ.}

\prose{\textbf{ญาณูปปนฺนํ มุนิ โน วทนฺตี}ติ อฏฺฐ สมาปตฺติญาเณน วา ปญฺจาภิญฺญาญาเณน วา อุเปตํ สมุเปตํ อุปาคตํ สมุปาคตํ อุปปนฺนํ สมุปปนฺนํ สมนฺนาคตํ มุนิํ วทนฺติ กเถนฺติ ภณนฺติ ทีปยนฺติ โวหรนฺตีติ ญาณูปปนฺนํ มุนิ โน วทนฺติ.}

\prose{\textbf{อุทาหุ เว ชีวิเตนูปปนฺนนฺ}ติ อุทาหุ อเนกวิวิธ-อติปรมทุกฺกร- การิกลูขชีวิตานุโยเคน อุเปตํ สมุเปตํ อุปาคตํ สมุปาคตํ อุปปนฺนํ สมุปปนฺนํ สมนฺนาคตํ มุนิํ วทนฺติ กเถนฺติ ภณนฺติ ทีปยนฺติ โวหรนฺตีติ อุทาหุ เว ชีวิเตนูปปนฺนํ.\csromanspacer{} เตนาห โส พฺราหฺมโณ\csromandash{}}

\prose{สนฺติ โลเก มุนโย, (อิจฺจายสฺมา นนฺโท,)}

\prose{ชนา วทนฺติ ตยิทํ กถํสุ.}

\csromangathameasure{ญาณูปปนฺนํ มุนิ โน วทนฺติ, \\ อุทาหุ เว ชีวิเตนุปปนฺนนฺติ.}
\csromangathagroup{\csromangathastanza{ญาณูปปนฺนํ มุนิ โน วทนฺติ, \\ อุทาหุ เว ชีวิเตนุปปนฺนนฺติ.}}

\csromanhangingbegin
\hangingheaditem{47}{น ทิฏฺฐิยา น สุติยา น ญาเณน,}
\hangingbody{\textbf{มุนีธ นนฺท กุสลา วทนฺติ.}}
\hangingbody{\textbf{วิเสนิกตฺวา1} \textbf{อนีฆา นิราสา,}}
\hangingbody{จรนฺติ เย เต มุนโยติ พฺรูมิ.}
\csromanhangingend

\prose{\textbf{น ทิฏฺฐิยา น สุติยา น ญาเณนา}ติ \textbf{น ทิฏฺฐิยา}ติ น ทิฏฺฐสุทฺธิยา.\csromanspacer{} \textbf{น สุติยา}ติ น สุตสุทฺธิยา.\csromanspacer{} \textbf{น ญาเณนา}ติ นปิ อฏฺฐสมาปตฺติญาเณน นปิ ปญฺจาภิญฺญาญาเณน นปิ มิจฺฉาญาเณนาติ \textbf{น ทิฏฺฐิยา} น สุติยา น ญาเณน.}

\prose{\textbf{มุนีธ นนฺท กุสลา วทนฺตี}ติ \textbf{กุสลา}ติ เย เต ขนฺธ\textbf{กุสลา} ธาตุ\textbf{กุสลา} อายตน\textbf{กุสลา} ปฏิจฺจสมุปฺปาท\textbf{กุสลา} สติปฏฺฐาน\textbf{กุสลา} สมฺมปฺปธาน\textbf{กุสลา} อิทฺธิปาท\textbf{กุสลา} อินฺทฺริย\textbf{กุสลา} พล\textbf{กุสลา} \csromanfolio{114}โพชฺฌงฺคกุสลา มคฺคกุสลา ผลกุสลา นิพฺพานกุสลา ทิฏฺฐสุทฺธิยา วา สุตสุทฺธิยา วา อฏฺฐสมาปตฺติญาเณน วา ปญฺจาภิญฺญาญาเณน วา มิจฺฉาญาเณน วา ทิฏฺเฐน วา สุเตน วา อุเปตํ สมุเปตํ อุปาคตํ สมุปาคตํ อุปปนฺนํ สมุปปนฺนํ สมนฺนาคตํ มุนิํ น วทนฺติ น กเถนฺติ น ภณนฺติ น ทีปยนฺติ น โวหรนฺตีติ มุนีธ นนฺท กุสลา วทนฺติ.}

\prose{\textbf{วิเสนิกตฺวา อนีฆา นิราสา, จรนฺติ เย เต มุนโยติ พฺรูมี}ติ \textbf{เสนา} วุจฺจติ มาร\textbf{เสนา}, กายทุจฺจริตํ มาร\textbf{เสนา}, วจีทุจฺจริตํ มาร\textbf{เสนา}, มโนทุจฺจริตํ มาร\textbf{เสนา}, ราโค มาร\textbf{เสนา}, โทโส มาร\textbf{เสนา}, โมโห มาร\textbf{เสนา}, โกโธ.\csromanspacer{} อุปนาโห.\csromanspacer{} มกฺโข.\csromanspacer{} ปฬาโส.\csromanspacer{} อิสฺสา.\csromanspacer{} มจฺฉริยํ.\csromanspacer{} มายา.\csromanspacer{} สาเฐยฺยํ.\csromanspacer{} ถมฺโภ.\csromanspacer{} สารมฺโภ.\csromanspacer{} มาโน.\csromanspacer{} อติมาโน.\csromanspacer{} มโท.\csromanspacer{} ปมาโท.\csromanspacer{} สพฺเพ กิเลสา.\csromanspacer{} สพฺเพ ทุจฺจริตา.\csromanspacer{} สพฺเพ ทรถา.\csromanspacer{} สพฺเพ ปริฬาหา.\csromanspacer{} สพฺเพ สนฺตาปา.\csromanspacer{} สพฺพากุสลาภิสงฺขารา มาร\textbf{เสนา}.\csromanspacer{} วุตฺตเญฺหตํ ภควตา\csromandash{}}

\csromangathameasure{กามา เต ปฐมา เสนา, ทุติยา อรติ วุจฺจติ. \\ ตติยา ขุปฺปิปาสา เต, จตุตฺถี ตณฺหา ปวุจฺจติ. \\ ปญฺจมํ ถินมิทฺธํ เต, ฉฏฺฐา ภีรู ปวุจฺจติ. \\ สตฺตมี วิจิกิจฺฉา เต, มกฺโข ถมฺโภ เต อฏฺฐโม. \\ ลาโภ สิโลโก สกฺกาโร, \\ มิจฺฉาลทฺโธ จ โย ยโส. \\ โย จตฺตานํ สมุกฺกํเส, ปเร จ อวชานาติ. \\ เอสา นมุจิ เต เสนา, กณฺหสฺสาภิปฺปหารินี. \\ น นํ อสูโร ชินาติ, เชตฺวา จ ลภเต สุขนฺติ.}
\csromangathagroup{\csromangathastanza{กามา เต ปฐมา เสนา, ทุติยา อรติ วุจฺจติ. \\ ตติยา ขุปฺปิปาสา เต, จตุตฺถี ตณฺหา ปวุจฺจติ.}\csromangathastanza{ปญฺจมํ ถินมิทฺธํ เต, ฉฏฺฐา ภีรู ปวุจฺจติ. \\ สตฺตมี วิจิกิจฺฉา เต, มกฺโข ถมฺโภ เต อฏฺฐโม.}\csromangathastanza{ลาโภ สิโลโก สกฺกาโร, \\ มิจฺฉาลทฺโธ จ โย ยโส.}\csromangathastanza{โย จตฺตานํ สมุกฺกํเส, ปเร จ อวชานาติ. \\ เอสา นมุจิ เต เสนา\footnote{เอสา เต นมุจิ เสนา (สฺยา, ก) ขุ 1. 343 ปิฏฺเฐ.}, กณฺหสฺสาภิปฺปหารินี. \\ น นํ อสูโร ชินาติ, เชตฺวา จ ลภเต สุขนฺติ.}}

\prose{ยโต จตูหิ อริยมคฺเคหิ สพฺพา จ มาร\textbf{เสนา} สพฺเพ จ ปฏิเสนิกรา\footnote{วิเสนิํกตฺวา (ก) ขุ 7. 135 ปิฏฺเฐ.} กิเลสา ชิตา จ ปราชิตา จ ภคฺคา วิปฺปลุคฺคา\footnote{วิปฺปลุคฺคตา (สฺยา), ปสฺส ขุ 7. 74 ปิฏฺเฐ.} ปรมฺมุขา, เตน วุจฺจนฺติ วิเสนิกตฺวา.\csromanspacer{} \textbf{อนีฆา}ติ ราโค นีโฆ.\csromanspacer{} โทโส นีโฆ.\csromanspacer{} โมโห นีโฆ.\csromanspacer{} โกโธ นีโฆ.\csromanspacer{} อุปนาโห นีโฆ ฯเปฯ สพฺพากุสลาภิสงฺขารา นีฆา, เยสํ เอเต นีฆา ปหีนา สมุจฺฉินฺนา วูปสนฺตา ปฏิปฺปสฺสทฺธา อภพฺพุปฺปตฺติกา ญาณคฺคินา ทฑฺฒา.\csromanspacer{} เต วุจฺจนฺติ อนีฆา.}

\vspace{\tipitakaparskip}
\csromancenter{นนฺทมาณวปุจฺฉานิทฺเทส \textbf{นิราสา}ติ อาสา วุจฺจติ ตณฺหา, โย ราโค สาราโค ฯเปฯ อวิชฺชา โลโภ อกุสลมูลํ.\csromanspacer{} เยสํ เอสา อาสา ตณฺหา ปหีนา สมุจฺฉินฺนา วูปสนฺตา ปฏิปฺปสฺสทฺธา อภพฺพุปฺปตฺติกา ญาณคฺคินา ทฑฺฒา.\csromanspacer{} เต วุจฺจนฺติ นิราสา อรหนฺโต ขีณาสวา.\csromanspacer{} \textbf{วิเสนิกตฺวา อนีฆา นิราสา, จรนฺติ เย เต มุนโยติ พฺรูมี}ติ เย เต วิเสนิกตฺวาว อนีฆา จ นิราสา จ จรนฺติ วิหรนฺติ อิริยนฺติ วตฺเตนฺติ ปาเลนฺติ ยเปนฺติ ยาเปนฺติ.\csromanspacer{} เต โลเก มุนโยติ พฺรูมิ อาจิกฺขามิ เทเสมิ ปญฺญเปมิ ปฏฺฐเปมิ วิวรามิ วิภชามิ อุตฺตานีกโรมิ ปกาเสมีติ วิเสนิกตฺวา อนีฆา นิราสา, จรนฺติ เย เต มุนโยติ พฺรูมิ.\csromanspacer{} เตนาห ภควา\csromandash{}}
\vspace{0.5\baselineskip}
\csromangathameasure{น ทิฏฺฐิยา น สุติยา น ญาเณน, \\ มุนีธ นนฺท กุสลา วทนฺติ.}
\csromangathagroup{\csromangathastanza{\csromanfolio{115}น ทิฏฺฐิยา น สุติยา น ญาเณน, \\ มุนีธ นนฺท กุสลา วทนฺติ.}}

\prosecont{วิเสนิกตฺวา อนีฆา นิราสา,}

\prose{จรนฺติ เย เต มุนโยติ พฺรูมีติ.}

\csromanhangingbegin
\hangingheaditem{48}{\textbf{เย เกจิเม สมณพฺราหฺมณาเส}, (\textbf{อิจฺจา}ยสฺมา \textbf{นนฺโท},)}
\hangingbody{\textbf{ทิฏฺถสฺสุเตนาปิ วทนฺติ สุทฺธิํ.}}
\hangingbody{\textbf{สีลพฺพเตนาปิ วทนฺติ สุทฺธิํ,}}
\hangingbody{\textbf{อเนกรูเปน วทนฺติ สุทฺธิํ.}}
\hangingbody{\textbf{กจฺจิสฺสุ เต ภควา ตตฺถ ยตา จรนฺตา,}}
\hangingbody{\textbf{อตารุ ชาติญฺจ ชรญฺจ มาริส.}}
\hangingbody{ปุจฺฉามิ ตํ ภควา พฺรูหิ เมตํ.}
\csromanhangingend

\prose{\textbf{เย เกจิเม สมณพฺราหฺมณาเส}ติ \textbf{เย เกจี}ติ สพฺเพน สพฺพํ สพฺพถา สพฺพํ อเสสํ นิสฺเสสํ ปริยาทิยนวจนเมตํ \textbf{เย เกจี}ติ.\csromanspacer{} \textbf{สมณา}ติ เย เกจิ อิโต พหิทฺธา ปพฺพชฺชูปคตา ปริพฺพาชกสมาปนฺนา.\csromanspacer{} \textbf{พฺราหฺมณา}ติ เย เกจิ โภวาทิกาติ เย เกจิเม สมณพฺราหฺมณาเส.\csromanspacer{} \textbf{อิจฺจายสฺมา นนฺโท}ติ \textbf{อิจฺจา}ติ ปทสนฺธิ ฯเปฯ.\csromanspacer{} \textbf{อายสฺมา}ติ ปิยวจนํ ฯเปฯ.\csromanspacer{} \textbf{นนฺโท}ติ ตสฺส พฺราหฺมณสฺส นามํ ฯเปฯ อภิลาโปติ \textbf{อิจฺจา}ยสฺมา \textbf{นนฺโท}.}

\prose{\textbf{ทิฏฺฐสฺสุเตนาปิ วทนฺติ สุทฺธินฺ}ติ ทิฏฺเฐนปิ สุทฺธิํ วิสุทฺธิํ ปริสุทฺธิํ, มุตฺติํ วิมุตฺติํ ปริมุตฺติํ วทนฺติ กเถนฺติ ภณนฺติ ทีปยนฺติ โวหรนฺติ.\csromanspacer{} สุเตนปิ สุทฺธิํ วิสุทฺธิํ ปริสุทฺธิํ, มุตฺติํ วิมุตฺติํ ปริมุตฺติํ วทนฺติ กเถนฺติ ภณนฺติ \csromanfolio{116}ทีปยนฺติ โวหรนฺติ.\csromanspacer{} ทิฏฺฐสฺสุเตนปิ สุทฺธิํ วิสุทฺธิํ ปริสุทฺธิํ, มุตฺติํ วิมุตฺติํ ปริมุตฺติํ วทนฺติ กเถนฺติ ภณนฺติ ทีปยนฺติ โวหรนฺตีติ ทิฏฺฐสฺสุเตนาปิ วทนฺติ สุทฺธิํ.}

\prose{\textbf{สีลพฺพเตนาปิ วทนฺติ สุทฺธินฺ}ติ สีเลนปิ สุทฺธิํ วิสุทฺธิํ ปริสุทฺธิํ, มุตฺติํ วิมุตฺติํ ปริมุตฺติํ วทนฺติ กเถนฺติ ภณนฺติ ทีปยนฺติ โวหรนฺติ.\csromanspacer{} วเตนปิ สุทฺธิํ วิสุทฺธิํ ปริสุทฺธิํ, มุตฺติํ วิมุตฺติํ ปริมุตฺติํ วทนฺติ กเถนฺติ ภณนฺติ ทีปยนฺติ โวหรนฺติ.\csromanspacer{} สีลพฺพเตนาปิ สุทฺธิํ วิสุทฺธิํ ปริสุทฺธิํ, มุตฺติํ วิมุตฺติํ ปริมุตฺติํ วทนฺติ กเถนฺติ ภณนฺติ ทีปยนฺติ โวหรนฺตีติ สีลพฺพเตนาปิ วทนฺติ สุทฺธิํ.}

\prose{\textbf{อเนกรูเปน วทนฺติ สุทฺธินฺ}ติ อเนกวิธโกตูหลมงฺคเลน\footnote{อเนกวิธวตฺต กุตูหลมงฺคเลน (สฺยา)} สุทฺธิํ วิสุทฺธิํ ปริสุทฺธิํ, มุตฺติํ วิมุตฺติํ ปริมุตฺติํ วทนฺติ กเถนฺติ ภณนฺติ ทีปยนฺติ โวหรนฺตีติ อเนกรูเปน วทนฺติ สุทฺธิํ.}

\prose{\textbf{กจฺจิสฺสุ เต ภควา ตตฺถ ยตา จรนฺตา}ติ \textbf{กจฺจิสฺสู}ติ สํสยปุจฺฉา วิมติปุจฺฉา เทฺวฬฺหกปุจฺฉา อเนกํสปุจฺฉา “เอวํ นุ โข, น นุ โข, กิํ นุ โข, กถํ นุ โข”ติ กจฺจิสฺสุ.\csromanspacer{} \textbf{เต}ติ ทิฏฺฐิคติกา.\csromanspacer{} \textbf{ภควา}ติ คารวาธิวจนเมตํ ฯเปฯ สจฺฉิกา ปญฺญตฺติ, ยทิทํ \textbf{ภควา}ติ กจฺจิสฺสุ เต \textbf{ภควา}.\csromanspacer{} \textbf{ตตฺถ ยตา จรนฺตา}ติ \textbf{ตตฺถา}ติ สกาย ทิฏฺฐิยา สกาย ขนฺติยา สกาย รุจิยา สกาย ลทฺธิยา.\csromanspacer{} \textbf{ยตา}ติ ยตฺตา ปฏิยตฺตา\footnote{ยตา ปฏิยตา (สฺยา)} คุตฺตา โคปิตา รกฺขิตา สํวุตา.\csromanspacer{} \textbf{จรนฺตา}ติ จรนฺตา วิหรนฺตา อิริยนฺตา วตฺเตนฺตา ปาเลนฺตา ยเปนฺตา ยาเปนฺตาติ กจฺจิสฺสุ เต \textbf{ภควา} ตตฺถ ยตา จรนฺตา.}

\prose{\textbf{อตารุ ชาติญฺจ ชรญฺจ มาริสา}ติ ชาติชรามรณํ อตริํสุ อุตฺตริํสุ ปตริํสุ สมติกฺกมิํสุ วีติวตฺติํสุ.\csromanspacer{} \textbf{มาริสา}ติ ปิยวจนํ ครุวจนํ สคารวสปฺปติสฺสาธิวจนเมตํ มาริสาติ อตารุ ชาติญฺจ ชรญฺจ มาริส.}

\prose{\textbf{ปุจฺฉามิ ตํ ภควา พฺรูหิ เมตนฺ}ติ \textbf{ปุจฺฉามิ ตนฺ}ติ ปุจฺฉามิ ตํ, ยาจามิ ตํ, อชฺเฌสามิ ตํ, “กถยสฺสุ เม”ติ ปุจฺฉามิ ตํ.\csromanspacer{} \textbf{ภควา}ติ ฯเปฯ สจฺฉิกา ปญฺญตฺติ, ยทิทํ \textbf{ภควา}ติ.\csromanspacer{} \textbf{พฺรูหิ เมตนฺ}ติ พฺรูหิ อาจิกฺขาหิ เทเสหิ}

\vspace{\tipitakaparskip}
\csromancenter{นนฺทมาณวปุจฺฉานิทฺเทส ปญฺญเปหิ ปฏฺฐเปหิ วิวราหิ วิภชาหิ อุตฺตานีกโรหิ ปกาเสหีติ ปุจฺฉามิ ตํ \textbf{ภควา} พฺรูหิ เมตํ.\csromanspacer{} เตนาห โส พฺราหฺมโณ\csromandash{}}
\prose{\csromanfolio{117}\textbf{เย เกจิเม สมณพฺราหฺมณาเส}, (อิจฺจายสฺมา นนฺโท,)}

\prose{\textbf{ทิฏฺฐสฺสุเตนาปิ วทนฺติ สุทฺธิํ.}}

\csromangathameasure{\textbf{สีลพฺพเตนาปิ วทนฺติ สุทฺธิํ,} \\ \textbf{อเนกรูเปน วทนฺติ สุทฺธิํ.} \\ กจฺจิสฺสุ เต \textbf{ภควา} ตตฺถ ยตา จรนฺตา, \\ อตารุ ชาติญฺจ ชรญฺจ มาริส. \\ ปุจฺจามิ ตํ ภควา พฺรูหิ เมตนฺติ.}
\csromangathagroup{\csromangathastanza{\textbf{สีลพฺพเตนาปิ วทนฺติ สุทฺธิํ,} \\ \textbf{อเนกรูเปน วทนฺติ สุทฺธิํ.} \\ กจฺจิสฺสุ เต \textbf{ภควา} ตตฺถ ยตา จรนฺตา, \\ อตารุ ชาติญฺจ ชรญฺจ มาริส.}\csromangathastanza{ปุจฺจามิ ตํ ภควา พฺรูหิ เมตนฺติ.}}

\csromanhangingbegin
\hangingheaditem{49}{\textbf{เย เกจิเม สมณพฺราหฺมณาเส}, (\textbf{นนฺทาติ ภควา},)}
\hangingbody{\textbf{ทิฏฺฐสฺสุเตนาปิ วทนฺติ สุทฺธิํ.}}
\hangingbody{\textbf{สีลพฺพเตนาปิ วทนฺติ สุทฺธิํ,}}
\hangingbody{\textbf{อเนกรูเปน วทนฺติ สุทฺธิํ.}}
\hangingbody{\textbf{กิญฺจาปิ เต ตตฺถ ยตา จรนฺติ,}}
\hangingbody{นาตริํสุ ชาติชรนฺติ พฺรูมิ.}
\csromanhangingend

\prose{\textbf{เย เกจิเม สมณพฺราหฺมณาเส}ติ \textbf{เย เกจี}ติ สพฺเพน สพฺพํ สพฺพถา สพฺพํ อเสสํ นิสฺเสสํ ปริยาทิยนวจนเมตํ \textbf{เย เกจี}ติ.\csromanspacer{} \textbf{สมณา}ติ เย เกจิ อิโต พหิทฺธา ปพฺพชฺชูปคตา ปริพฺพาชกสมาปนฺนา.\csromanspacer{} \textbf{พฺราหฺมณา}ติ เย เกจิ โภวาทิกาติ เย เกจิเม สมณพฺราหฺมณาเส.\csromanspacer{} \textbf{นนฺทาติ ภควา}ติ \textbf{นนฺทาติ ภควา} ตํ พฺราหฺมณํ นาเมน อาลปติ.\csromanspacer{} \textbf{ภควา}ติ คารวาธิวจนเมตํ ฯเปฯ สจฺฉิกา ปญฺญตฺติ, ยทิทํ \textbf{ภควา}ติ \textbf{นนฺทาติ ภควา}.}

\prose{\textbf{ทิฏฺฐสฺสุเตนาปิ วทนฺติ สุทฺธินฺ}ติ ทิฏฺเฐนปิ สุทฺธิํ วิสุทฺธิํ ปริสุทฺธิํ, มุตฺติํ วิมุตฺติํ ปริมุตฺติํ วทนฺติ กเถนฺติ ภณนฺติ ทีปยนฺติ โวหรนฺติ.\csromanspacer{} สุเตนปิ สุทฺธิํ ฯเปฯ ทิฏฺฐสฺสุเตนปิ สุทฺธิํ วิสุทฺธิํ ปริสุทฺธิํ, มุตฺติํ วิมุตฺติํ ปริมุตฺติํ วทนฺติ กเถนฺติ ภณนฺติ ทีปยนฺติ โวหรนฺตีติ ทิฏฺฐสฺสุเตนาปิ วทนฺติ สุทฺธิํ.}

\prose{\textbf{สีลพฺพเตนาปิ วทนฺติ สุทฺธินฺ}ติ สีเลนปิ สุทฺธิํ วิสุทฺธิํ ปริสุทฺธิํ, มุตฺติํ วิมุตฺติํ ปริมุตฺติํ วทนฺติ กเถนฺติ ภณนฺติ ทีปยนฺติ โวหรนฺติ.\csromanspacer{} วเตนปิ สุทฺธิํ ฯเปฯ โวหรนฺติ.\csromanspacer{} สีลพฺพเตนาปิ สุทฺธิํ วิสุทฺธิํ ปริสุทฺธิํ, มุตฺติํ วิมุตฺติํ \csromanfolio{118}ปริมุตฺติํ วทนฺติ กเถนฺติ ภณนฺติ ทีปยนฺติ โวหรนฺตีติ สีลพฺพเตนาปิ วทนฺติ สุทฺธิํ.}

\prose{\textbf{อเนกรูเปน วทนฺติ สุทฺธินฺ}ติ อเนกวิธโกตูหลมงฺคเลน สุทฺธิํ วิสุทฺธิํ ปริสุทฺธิํ, มุตฺติํ วิมุตฺติํ ปริมุตฺติํ วทนฺติ กเถนฺติ ภณนฺติ ทีปยนฺติ โวหรนฺตีติ อเนกรูเปน วทนฺติ สุทฺธิํ.}

\prose{\textbf{กิญฺจาปิ เต ตตฺถ ยตา จรนฺตี}ติ \textbf{กิญฺจาปี}ติ ปทสนฺธิ ปทสํสคฺโค ปทปาริปูรี อกฺขรสมวาโย พฺยญฺชนสิลิฏฺฐตา ปทานุปุพฺพตาเปตํ \textbf{กิญฺจาปี}ติ.\csromanspacer{} \textbf{เต}ติ ทิฏฺฐิคติกา.\csromanspacer{} \textbf{ตตฺถา}ติ สกาย ทิฏฺฐิยา สกาย ขนฺติยา สกาย รุจิยา สกาย ลทฺธิยา.\csromanspacer{} \textbf{ยตา}ติ ยตฺตา ปฏิยตฺตา คุตฺตา โคปิตา รกฺขิตา สํวุตา.\csromanspacer{} \textbf{จรนฺตี}ติ จรนฺติ วิหรนฺติ อิริยนฺติ วตฺเตนฺติ ปาเลนฺติ ยเปนฺติ ยาเปนฺตีติ กิญฺจาปิ เต ตตฺถ ยตา จรนฺติ.}

\prose{\textbf{นาตริํสุ ชาติชรนฺติ พฺรูมี}ติ ชาติชรามรณํ น ตริํสุ น อุตฺตริํสุ น ปตริํสุ น สมติกฺกมิํสุ น วีติวตฺติํสุ, ชาติชรามรณา อนิกฺขนฺตา อนิสฺสฏา อนติกฺกนฺตา อสมติกฺกนฺตา อวีติวตฺตา, อนฺโตชาติชรามรเณ ปริวตฺเตนฺติ, อนฺโตสํสารปเถ ปริวตฺเตนฺติ, ชาติยา อนุคตา, ชราย อนุสฏา, พฺยาธินา อภิภูตา, มรเณน อพฺภาหตา อตาณา อเลณา อสรณา อสรณีภูตาติ พฺรูมิ อาจิกฺขามิ เทเสมิ ปญฺญเปมิ ปฏฺฐเปมิ วิวรามิ วิภชามิ อุตฺตานีกโรมิ ปกาเสมีติ นาตริํสุ ชาติชรนฺติ พฺรูมิ.\csromanspacer{} \textbf{เต}นาห ภควา\csromandash{}}

\prose{เย เกจิเม สมณพฺราหฺมณาเส, (นนฺทาติ ภควา,)}

\prose{\textbf{ทิฏฺฐสฺสุเตนาปิ วทนฺติ สุทฺธิํ.}}

\csromangathameasure{\textbf{สีลพฺพเตนาปิ วทนฺติ สุทฺธิํ,} \\ \textbf{อเนกรูเปน วทนฺติ สุทฺธิํ.}}
\csromangathagroup{\csromangathastanza{\textbf{สีลพฺพเตนาปิ วทนฺติ สุทฺธิํ,} \\ \textbf{อเนกรูเปน วทนฺติ สุทฺธิํ.}}}

\prose{กิญฺจาปิ เต ตตฺถ ยตา จรนฺติ,}

\prose{\textbf{นาตริํสุ ชาติชรนฺติ พฺรูมี}ติ.}

\csromanhangingbegin
\hangingheaditem{50}{เย เกจิเม สมณพฺราหฺมณาเส, (อิจฺจายสฺมา นนฺโท,)}
\hangingbody{\textbf{ทิฏฺฐสฺสุเตนาปิ วทนฺติ สุทฺธิํ.}}
\hangingbody{\textbf{สีลพฺพเตนาปิ วทนฺติ สุทฺธิํ,}}
\hangingbody{\textbf{อเนกรูเปน วทนฺติ สุทฺธิํ.}}
\csromanhangingend

\csromanheader{2}{7. นนฺทมาณวปุจฺฉานิทฺเทส}
\csromangathameasure{\textbf{เต เจ มุนี พฺรูสิ อโนฆติณฺเณ,} \\ \textbf{อถ โก จรหิ เทวมนุสฺสโลเก.} \\ \textbf{อตาริ ชาติญฺจ ชรญฺจ มาริส,} \\ ปุจฺฉามิ ตํ ภควา พฺรูหิ เมตํ.}
\csromangathagroup{\csromangathastanza{\csromanfolio{119}\textbf{เต เจ มุนี พฺรูสิ อโนฆติณฺเณ,} \\ \textbf{อถ โก จรหิ เทวมนุสฺสโลเก.} \\ \textbf{อตาริ ชาติญฺจ ชรญฺจ มาริส,} \\ ปุจฺฉามิ ตํ ภควา พฺรูหิ เมตํ.}}

\prose{\textbf{เย เกจิเม สมณพฺราหฺมณาเส}ติ \textbf{เย เกจี}ติ สพฺเพน สพฺพํ สพฺพถา สพฺพํ อเสสํ นิสฺเสสํ ปริยาทิยนวจนเมตํ \textbf{เย เกจี}ติ.\csromanspacer{} \textbf{สมณา}ติ เย เกจิ อิโต พหิทฺธา ปพฺพชฺชูปคตา ปริพฺพาชกสมาปนฺนา.\csromanspacer{} \textbf{พฺราหฺมณา}ติ เย เกจิ โภวาทิกาติ เย เกจิเม สมณพฺราหฺมณาเส.\csromanspacer{} \textbf{อิจฺจายสฺมา นนฺโท}ติ \textbf{อิจฺจา}ติ ปทสนฺธิ ฯเปฯ \textbf{อิจฺจา}ยสฺมา นนฺโท.}

\prose{\textbf{ทิฏฺฐสฺสุเตนาปิ วทนฺติ สุทฺธินฺ}ติ ทิฏฺเฐนปิ สุทฺธิํ วิสุทฺธิํ ปริสุทฺธิํ, มุตฺติํ วิมุตฺติํ ปริมุตฺติํ วทนฺติ กเถนฺติ ภณนฺติ ทีปยนฺติ โวหรนฺติ.\csromanspacer{} สุเตนาปิ สุทฺธิํ ฯเปฯ ทิฏฺฐสฺสุเตนาปิ สุทฺธิํ วิสุทฺธิํ ปริสุทฺธิํ, มุตฺติํ วิมุตฺติํ ปริมุตฺติํ วทนฺติ กเถนฺติ ภณนฺติ ทีปยนฺติ โวหรนฺตีติ ทิฏฺฐสฺสุเตนาปิ วทนฺติ สุทฺธิํ.}

\prose{\textbf{สีลพฺพเตนาปิ วทนฺติ สุทฺธินฺ}ติ สีเลนาปิ สุทฺธิํ วิสุทฺธิํ ปริสุทฺธิํ, มุตฺติํ วิมุตฺติํ ปริมุตฺติํ วทนฺติ กเถนฺติ ภณนฺติ ทีปยนฺติ โวหรนฺติ.\csromanspacer{} วเตนาปิ สุทฺธิํ ฯเปฯ โวหรนฺติ.\csromanspacer{} สีลพฺพเตนาปิ สุทฺธิํ วิสุทฺธิํ ปริสุทฺธิํ, มุตฺติํ วิมุตฺติํ ปริมุตฺติํ วทนฺติ กเถนฺติ ภณนฺติ ทีปยนฺติ โวหรนฺตีติ สีลพฺพเตนาปิ วทนฺติ สุทฺธิํ.}

\prose{\textbf{อเนกรูเปน วทนฺติ สุทฺธินฺ}ติ อเนกวิธโกตูหลมงฺคเลน สุทฺธิํ วิสุทฺธิํ ปริสุทฺธิํ, มุตฺติํ วิมุตฺติํ ปริมุตฺติํ วทนฺติ กเถนฺติ ภณนฺติ ทีปยนฺติ โวหรนฺตีติ อเนกรูเปน วทนฺติ สุทฺธิํ.}

\prose{\textbf{เต เจ มุนี พฺรูสิ อโนฆติณฺเณ}ติ \textbf{เต เจ}ติ ทิฏฺฐิคติเก.\csromanspacer{} \textbf{มุนี}ติ โมนํ วุจฺจติ ญาณํ ฯเปฯ สงฺคชาลมติจฺจ โส มุนิ.\csromanspacer{} \textbf{พฺรูสิ อโนฆติณฺเณ}ติ กาโมฆํ ภโวฆํ ทิฏฺโฐฆํ อวิชฺโชฆํ อติณฺเณ อนติกฺกนฺเต อสมติกฺกนฺเต อวีติวตฺเต, อนฺโตชาติชรามรเณ ปริวตฺเตนฺเต, อนฺโตสํสารปเถ ปริวตฺเตนฺเต, ชาติยา อนุคเต, ชราย อนุสเฏ, พฺยาธินา อภิภูเต, มรเณน อพฺภาหเต อตาเณ อเลเณ อสรเณ อสรณีภูเต.\csromanspacer{} \textbf{พฺรูสี}ติ พฺรูสิ อาจิกฺขสิ เทเสสิ ปญฺญเปสิ ปฏฺฐเปสิ}

\prose{\csromanfolio{120}วิวรสิ วิภชสิ อุตฺตานีกโรสิ ปกาเสสีติ เต เจ มุนี พฺรูสิ อโนฆติณฺเณ\textbf{.}}

\prose{\textbf{อถ โก จรหิ เทวมนุสฺสโลเก, อตาริ ชาติญฺจ ชรญฺจ มาริสา}ติ อถ โก เอโส สเทวเก โลเก สมารเก สพฺรหฺมเก สสฺสมณพฺราหฺมณิยา ปชาย สเทวมนุสฺสาย ชาติชรามรณํ อตริ อุตฺตริ ปตริ สมติกฺกมิ วีติวตฺตยิ\textbf{.}\csromanspacer{} \textbf{มาริสา}ติ ปิยวจนํ ครุวจนํ สคารวสปฺปติสฺสาธิวจนเมตํ มาริสาติ อถ โก จรหิ เทวมนุสฺสโลเก, อตาริ ชาติญฺจ ชรญฺจ มาริส\textbf{.}}

\prose{\textbf{ปุจฺฉามิ ตํ ภควา พฺรูหิ เมตนฺ}ติ \textbf{ปุจฺฉามิ ตนฺ}ติ ปุจฺฉามิ ตํ, ยาจามิ ตํ, อชฺเฌสามิ ตํ, ปสาเทมิ ตํ\textbf{.}\csromanspacer{} \textbf{ภควา}ติ คารวาธิวจนเมตํ ฯเปฯ สจฺฉิกา ปญฺญตฺติ, ยทิทํ \textbf{ภควา}ติ\textbf{.}\csromanspacer{} \textbf{พฺรูหิ เมตนฺ}ติ พฺรูหิ อาจิกฺขาหิ เทเสหิ ปญฺญเปหิ ปฏฺฐเปหิ วิวราหิ วิภชาหิ อุตฺตานีกโรหิ ปกาเสหีติ ปุจฺฉามิ ตํ \textbf{ภควา} พฺรูหิ เมตํ\textbf{.}\csromanspacer{} เตนาห โส พฺราหฺมโณ\csromandash{}}

\prose{เย เกจิเม สมณพฺราหฺมณาเส, (อิจฺจายสฺมา นนฺโท,)}

\prose{ทิฏฺฐสฺสุเตนาปิ วทนฺติ สุทฺธิํ\textbf{.}}

\csromangathameasure{สีลพฺพเตนาปิ วทนฺติ สุทฺธิํ, \\ อเนกรูเปน วทนฺติ สุทฺธิํ\textbf{.} \\ เต เจ มุนี พฺรูสิ อโนฆติณฺเณ,}
\csromangathagroup{\csromangathastanza{สีลพฺพเตนาปิ วทนฺติ สุทฺธิํ, \\ อเนกรูเปน วทนฺติ สุทฺธิํ\textbf{.} \\ เต เจ มุนี พฺรูสิ อโนฆติณฺเณ,}}

\prose{อถ โก จรหิ เทวมนุสฺสโลเก\textbf{.}}

\prose{อตาริ ชาติญฺจ ชรญฺจ มาริส,}

\prose{\textbf{ปุจฺฉามิ ตํ ภควา พฺรูหิ เมตนฺ}ติ\textbf{.}}

\csromanhangingbegin
\hangingheaditem{51}{นาหํ สพฺเพ สมณพฺราหฺมณาเส, (นนฺทาติ \textbf{ภควา},)}
\hangingbody{\textbf{ชาติชราย นิวุตาติ พฺรูมิ.}}
\hangingbody{\textbf{เย สีธ ทิฏฺฐํ ว สุตํ มุตํ วา,}}
\hangingbody{\textbf{สีลพฺพตํ วาปิ ปหาย สพฺพํ.}}
\hangingbody{\textbf{อเนกรูปมฺปิ ปหาย สพฺพํ,}}
\hangingbody{\textbf{ตณฺหํ ปริญฺญาย อนาสวาเส1.}}
\hangingbody{เต เว นรา โอฆติณฺณาติ พฺรูมิ\textbf{.}}
\csromanhangingend

\csromanheader{2}{7. นนฺทมาณวปุจฺฉานิทฺเทส}
\prose{\csromanfolio{121}\textbf{นาหํ สพฺเพ สมณพฺราหฺมณาเส, นนฺทาติ ภควา.}\csromanspacer{}\textbf{ ชาติชราย นิวุตาติ พฺรูมี}ติ นาหํ นนฺท สพฺเพ สมณพฺราหฺมณา ชาติชราย อาวุตา นิวุตา โอวุตา ปิหิตา ปฏิจฺฉนฺนา ปฏิกุชฺชิตาติ วทามิ.\csromanspacer{} อตฺถิ เต สมณพฺราหฺมณา.\csromanspacer{} เยสํ ชาติ จ ชรามรณญฺจ ปหีนา อุจฺฉินฺนมูลา ตาลาวตฺถุกตา อนภาวํกตา อายติํ อนุปฺปาทธมฺมาติ พฺรูมิ อาจิกฺขามิ เทเสมิ ปญฺญเปมิ ปฏฺฐเปมิ วิวรามิ วิภชามิ อุตฺตานีกโรมิ ปกาเสมีติ นาหํ สพฺเพ สมณพฺราหฺมณาเส.\csromanspacer{} นนฺทาติ ภควา ชาติชราย นิวุตาติ พฺรูมิ.}

\prose{\textbf{เย สีธ ทิฏฺฐํ ว สุตํ มุตํ วา, สีลพฺพตํ วาปิ ปหาย สพฺพนฺ}ติ เย สพฺพา ทิฏฺฐสุทฺธิโย ปหาย ชหิตฺวา ปชหิตฺวา วิโนเทตฺวา พฺยนฺตีกริตฺวา อนภาวํ คเมตฺวา.\csromanspacer{} เย สพฺพา สุตสุทฺธิโย ปหาย ฯเปฯ.\csromanspacer{} เย สพฺพา มุตสุทฺธิโย ปหาย.\csromanspacer{} เย สพฺพา ทิฏฺฐสุตมุตสุทฺธิโย ปหาย.\csromanspacer{} เย สพฺพา สีลสุทฺธิโย ปหาย.\csromanspacer{} เย สพฺพา วตสุทฺธิโย ปหาย.\csromanspacer{} เย สพฺพา สีลพฺพตสุทฺธิโย ปหาย ชหิตฺวา ปชหิตฺวา วิโนเทตฺวา พฺยนฺตีกริตฺวา อนภาวํ คเมตฺวาติ เย สีธ ทิฏฺฐํ ว สุตํ มุตํ วา, สีลพฺพตํ วาปิ ปหาย สพฺพํ.}

\prose{\textbf{อเนกรูปมฺปิ ปหาย สพฺพนฺ}ติ อเนกวิธโกตูหลมงฺคเลน สุทฺธิํ วิสุทฺธิํ ปริสุทฺธิํ, มุตฺติํ วิมุตฺติํ ปริมุตฺติํ ปหาย ชหิตฺวา ปชหิตฺวา วิโนเทตฺวา พฺยนฺตีกริตฺวา อนภาวํ คเมตฺวาติ อเนกรูปมฺปิ ปหาย สพฺพํ.}

\prose{\textbf{ตณฺหํ ปริญฺญาย อนาสวา เส, เต เว นรา โอฆติณฺณาติ พฺรูมี}ติ \textbf{ตณฺหา}ติ รูป\textbf{ตณฺหา} สทฺท\textbf{ตณฺหา} คนฺธ\textbf{ตณฺหา} รส\textbf{ตณฺหา} โผฏฺฐพฺพ\textbf{ตณฺหา} ธมฺม\textbf{ตณฺหา}.\csromanspacer{} \textbf{ตณฺหํ ปริญฺญายา}ติ ตณฺหํ ตีหิ ปริญฺญาหิ ปริชานิตฺวา ญาตปริญฺญาย ตีรณปริญฺญาย ปหานปริญฺญาย.\csromanspacer{} กตมา ญาตปริญฺญา, ตณฺหํ ชานาติ\footnote{ปชานาติ (สฺยา) ปริชานาติ (ก) ขุ 7. 40 ปิฏฺเฐปิ.} “อยํ รูป\textbf{ตณฺหา}, อยํ สทฺท\textbf{ตณฺหา}, อยํ คนฺธ\textbf{ตณฺหา}, อยํ รส\textbf{ตณฺหา}, อยํ โผฏฺฐพฺพ\textbf{ตณฺหา}, อยํ ธมฺม\textbf{ตณฺหา}”ติ ชานาติ ปสฺสติ.\csromanspacer{} อยํ ญาตปริญฺญา.}

\prose{กตมา ตีรณปริญฺญา, เอวํ ญาตํ กตฺวา ตณฺหํ ตีเรติ อนิจฺจโต.\csromanspacer{} ทุกฺขโต.\csromanspacer{} โรคโต.\csromanspacer{} คณฺฑโต ฯเปฯ นิสฺสรณโต ตีเรติ.\csromanspacer{} อยํ ตีรณปริญฺญา.}

\prose{\csromanfolio{122}กตมา ปหานปริญฺญา, เอวํ ตีรยิตฺวา ตณฺหํ ปชหติ วิโนเทติ พฺยนฺตีกโรติ อนภาวํ คเมติ.\csromanspacer{} วุตฺตเญฺหตํ ภควตา\csromandash{}โย ภิกฺขเว ตณฺหาย ฉนฺทราโค, ตํ ปชหถ.\csromanspacer{} เอวํ สา ตณฺหา ปหีนา ภวิสฺสติ อุจฺฉินฺนมูลา ตาลาวตฺถุกตา อนภาวํกตา อายติํ อนุปฺปาทธมฺมา.\csromanspacer{} อยํ ปหานปริญฺญา.\csromanspacer{}\csromanspacer{}\textbf{ตณฺหํ ปริญฺญายา}ติ ตณฺหํ อิมาหิ ตีหิ ปริญฺญาหิ ปริชานิตฺวา.\csromanspacer{} \textbf{อนาสวา}ติ จตฺตาโร อาสวา กามาสโว ภวาสโว ทิฏฺฐาสโว อวิชฺชาสโว.\csromanspacer{} เยสํ อิเม อาสวา ปหีนา อุจฺฉินฺนมูลา ตาลาวตฺถุกตา อนภาวํกตา อายติํ อนุปฺปาทธมฺมา.\csromanspacer{} เต วุจฺจนฺติ อนาสวา อรหนฺโต ขีณาสวา.\csromanspacer{} ตณฺหํ ปริญฺญาย อนาสวา.}

\prose{\textbf{เต เว นรา โอฆติณฺณาติ พฺรูมี}ติ เย ตณฺหํ ปริญฺญาย อนาสวา.\csromanspacer{} เต กาโมฆํ ติณฺณา, ภโวฆํ ติณฺณา, ทิฏฺโฐฆํ ติณฺณา, อวิชฺโชฆํ ติณฺณา, สพฺพสํสารปถํ ติณฺณา อุตฺติณฺณา นิตฺถิณฺณา อติกฺกนฺตา สมติกฺกนฺตา วีติวตฺตาติ พฺรูมิ อาจิกฺขามิ เทเสมิ ปญฺญเปมิ ปฏฺฐเปมิ วิวรามิ วิภชามิ อุตฺตานีกโรมิ ปกาเสมีติ ตณฺหํ ปริญฺญาย อนาสวาเส, เต เว นรา โอฆติณฺณาติ พฺรูมิ.\csromanspacer{} เตนาห ภควา\csromandash{}}

\prose{นาหํ สพฺเพ สมณพฺราหฺมณาเส, (นนฺทาติ ภควา,)}

\prose{ชาติชราย นิวุตาติ พฺรูมิ.}

\csromangathameasure{\textbf{เย สีธ ทิฏฺฐํ ว สุตํ มุตํ วา,} \\ \textbf{สีลพฺพตํ วาปิ ปหาย สพฺพํ.} \\ \textbf{อเนกรูปมฺปิ ปหาย สพฺพํ,} \\ \textbf{ตณฺหํ ปริญฺญาย อนาสวาเส.}}
\csromangathagroup{\csromangathastanza{\textbf{เย สีธ ทิฏฺฐํ ว สุตํ มุตํ วา,} \\ \textbf{สีลพฺพตํ วาปิ ปหาย สพฺพํ.} \\ \textbf{อเนกรูปมฺปิ ปหาย สพฺพํ,} \\ \textbf{ตณฺหํ ปริญฺญาย อนาสวาเส.}}}

\prose{\textbf{เต เว นรา โอฆติณฺณาติ พฺรูมี}ติ.}

\csromanhangingbegin
\hangingheaditem{52}{เอตา’ภินนฺทามิ วโจ มเหสิโน,}
\hangingbody{\textbf{สุกิตฺติตํ โคตม’นูปธีกํ.}}
\hangingbody{\textbf{เย สีธ ทิฏฺฐํ ว สุตํ มุตํ วา,}}
\hangingbody{\textbf{สีลพฺพตํ วาปิ ปหาย สพฺพํ.}}
\hangingbody{\textbf{อเนกรูปมฺปิ ปหาย สพฺพํ,}}
\hangingbody{\textbf{ตณฺหํ ปริญฺญาย อนาสวาเส.}}
\hangingbody{อหมฺปิ เต โอฆติณฺณาติ พฺรูมิ.}
\csromanhangingend

\csromanheader{2}{7. นนฺทมาณวปุจฺฉานิทฺเทส}
\prose{\csromanfolio{123}\textbf{เอตา’ภินนฺทามิ วโจ มเหสิโน}ติ \textbf{เอตนฺ}ติ ตุยฺหํ วจนํ พฺยปฺปถํ เทสนํ อนุสาสนํ อนุสิฏฺฐํ นนฺทามิ อภินนฺทามิ โมทามิ อนุโมทามิ อิจฺฉามิ สาทิยามิ ปตฺถยามิ ปิหยามิ อภิชปฺปามิ.\csromanspacer{} \textbf{มเหสิโน}ติ กิํ มเหสิ ภควา.\csromanspacer{} มหนฺตํ สีลกฺขนฺธํ เอสี คเวสี ปริเยสีติ มเหสิ ฯเปฯ “กหํ นราสโภ”ติ มเหสีติ เอตาภินนฺทามิ วโจ \textbf{มเหสิโน}.}

\prose{\textbf{สุกิตฺติตํ โคตม’นูปธีกนฺ}ติ \textbf{สุกิตฺติตนฺ}ติ สุกิตฺติตํ สุ-อาจิกฺขิตํ\footnote{สฺวาจิกฺขิตํ (ก)} สุเทสิตํ สุปญฺญปิตํ สุปฏฺฐปิตํ สุวิวฏํ สุวิภตฺตํ สุ- อุตฺตานีกตํ สุปกาสิตํ.\csromanspacer{} \textbf{โคตม’นูปธีกนฺ}ติ อุปธี วุจฺจนฺติ กิเลสา จ ขนฺธา จ อภิสงฺขารา จ.\csromanspacer{} อุปธิปฺปหานํ อุปธิวูปสมํ อุปธินิสฺสคฺคํ อุปธิปฏิปฺปสฺสทฺธํ อมตํ นิพฺพานนฺติ สุกิตฺติตํ โคตม’นูปธีกํ.}

\prose{\textbf{เย สีธ ทิฏฺฐํ ว สุตํ มุตํ วา, สีลพฺพตํ วาปิ ปหาย สพฺพนฺ}ติ เย สพฺพา ทิฏฺฐสุทฺธิโย ปหาย ชหิตฺวา ปชหิตฺวา วิโนเทตฺวา พฺยนฺตีกริตฺวา อนภาวํ คเมตฺวา.\csromanspacer{} เย สพฺพา สุตสุทฺธิโย ฯเปฯ.\csromanspacer{} เย สพฺพา มุตสุทฺธิโย.\csromanspacer{} เย สพฺพา ทิฏฺฐสุตมุตสุทฺธิโย.\csromanspacer{} เย สพฺพา สีลสุทฺธิโย.\csromanspacer{} เย สพฺพา วตสุทฺธิโย.\csromanspacer{} เย สพฺพา สีลพฺพตสุทฺธิโย ปหาย ชหิตฺวา ปชหิตฺวา วิโนเทตฺวา พฺยนฺตีกริตฺวา อนภาวํ คเมตฺวาติ เย สีธ ทิฏฺฐํ ว สุตํ มุตํ วา, สีลพฺพตํ วาปิ ปหาย สพฺพํ.}

\prose{\textbf{อเนกรูปมฺปิ ปหาย สพฺพนฺ}ติ อเนกวิธโกตูหลมงฺคเลน สุทฺธิํ วิสุทฺธิํ ปริสุทฺธิํ, มุตฺติํ วิมุตฺติํ ปริมุตฺติํ ปหาย ชหิตฺวา ปชหิตฺวา วิโนเทตฺวา พฺยนฺตีกริตฺวา อนภาวํ คเมตฺวาติ อเนกรูปมฺปิ ปหาย สพฺพํ.}

\prose{\textbf{ตณฺหํ ปริญฺญาย อนาสวาเส, อหมฺปิ เต โอฆติณฺณาติ พฺรูมี}ติ \textbf{ตณฺหา}ติ รูป\textbf{ตณฺหา}, สทฺท\textbf{ตณฺหา}, คนฺธ\textbf{ตณฺหา}, รส\textbf{ตณฺหา}, โผฏฺฐพฺพ\textbf{ตณฺหา}, ธมฺม\textbf{ตณฺหา}.\csromanspacer{} \textbf{ตณฺหํ ปริญฺญายา}ติ ตณฺหํ ตีหิ ปริญฺญาหิ ปริชานิตฺวา ญาตปริญฺญาย ตีรณปริญฺญาย\footnote{ติรณปริญฺญาย (สฺยา)} ปหานปริญฺญาย.\csromanspacer{} กตมา ญาตปริญฺญา, ตณฺหํ ชานาติ “อยํ รูป\textbf{ตณฺหา}, อยํ สทฺท\textbf{ตณฺหา}, อยํ คนฺธ\textbf{ตณฺหา}, อยํ รส\textbf{ตณฺหา}, อยํ โผฏฺฐพฺพ\textbf{ตณฺหา}, อยํ ธมฺม\textbf{ตณฺหา}”ติ ชานาติ ปสฺสติ.\csromanspacer{} อยํ ญาตปริญฺญา.}

\prose{กตมา ตีรณปริญฺญา, เอวํ ญาตํ กตฺวา ตณฺหํ ตีเรติ\footnote{ติเรติ (สฺยา)} อนิจฺจโต ทุกฺขโต โรคโต คณฺฑโต สลฺลโต อฆโต อาพาธโต}

\prose{\csromanfolio{124}ปรโต ปโลกโต อีติโต อุปทฺทวโต ภยโต อุปสคฺคโต จลโต ปภงฺคุโต อทฺธุวโต อตาณโต อเลณโต อสรณโต อสรณีภูตโต ริตฺตโต ตุจฺฉโต สุญฺญโต อนตฺตโต อาทีนวโต วิปริณามธมฺมโต อสารกโต อฆมูลโต วธกโต วิภวโต สาสวโต สงฺขตโต มารามิสโต ชาติธมฺมโต ชราธมฺมโต พฺยาธิธมฺมโต มรณธมฺมโต โสกปริเทวทุกฺขโทมนสฺสุปายาส- ธมฺมโต สํกิเลสธมฺมโต สมุทยโต อตฺถงฺคมโต อสฺสาทโต อาทีนวโต นิสฺสรณโต ตีเรติ.\csromanspacer{} อยํ ตีรณปริญฺญา.}

\prose{กตมา ปหานปริญฺญา, เอวํ ตีรยิตฺวา ตณฺหํ ปชหติ วิโนเทติ พฺยนฺตีกโรติ อนภาวํ คเมติ.\csromanspacer{} อยํ ปหานปริญฺญา.}

\prose{\textbf{ตณฺหํ ปริญฺญายา}ติ ตณฺหํ อิมาหิ ตีหิ ปริญฺญาหิ ปริชานิตฺวา.\csromanspacer{} อ\textbf{นาสวา}ติ จตฺตาโร อาสวา กามาสโว ภวาสโว ทิฏฺฐาสโว อวิชฺชาสโว.\csromanspacer{} เยสํ อิเม อาสวา ปหีนา อุจฺฉินฺนมูลา ตาลาวตฺถุกตา อนภาวํกตา อายติํ อนุปฺปาทธมฺมา.\csromanspacer{} เต วุจฺจนฺติ อ\textbf{นาสวา}.\csromanspacer{} อรหนฺโต ขีณาสวา.\csromanspacer{} \textbf{ตณฺหํ ปริญฺญาย อนาสวาเส, อหมฺปิ เต โอฆติณฺณาติ พฺรูมี}ติ เย ตณฺหํ ปริญฺญาย อ\textbf{นาสวา}, อหมฺปิ เต กาโมฆํ ติณฺณา, ภโวฆํ ติณฺณา, ทิฏฺโฐฆํ ติณฺณา, อวิชฺโชฆํ ติณฺณา, สพฺพสํสารปถํ ติณฺณา อุตฺติณฺณา นิตฺถิณฺณา อติกฺกนฺตา สมติกฺกนฺตา วีติวตฺตาติ พฺรูมิ วทามีติ ตณฺหํ ปริญฺญาย อ\textbf{นาสวา}เส, อหมฺปิ เต โอฆติณฺณาติ พฺรูมิ.\csromanspacer{} เตนาห โส พฺราหฺมโณ\csromandash{}}

\csromangathameasure{เอตา’ภินนฺทามิ วโจ มเหสิโน, \\ สุกิตฺติตํ โคตม’นูปธีกํ. \\ เย สีธ ทิฏฺฐํ ว สุตํ มุตํ วา, \\ สีลพฺพตํ วาปิ ปหาย สพฺพํ. \\ อเนกรูปมฺปิ ปหาย สพฺพํ. ตณฺหํ ปริญฺญาย อนาสวาเส. \\ อหมฺปิ เต โอฆติณฺณาติ พฺรูมีติ.}
\csromangathagroup{\csromangathastanza{เอตา’ภินนฺทามิ วโจ มเหสิโน, \\ สุกิตฺติตํ โคตม’นูปธีกํ. \\ เย สีธ ทิฏฺฐํ ว สุตํ มุตํ วา, \\ สีลพฺพตํ วาปิ ปหาย สพฺพํ. \\ อเนกรูปมฺปิ ปหาย สพฺพํ. ตณฺหํ ปริญฺญาย อนาสวาเส. \\ อหมฺปิ เต โอฆติณฺณาติ พฺรูมีติ.}}

\vspace{\tipitakaparskip}
\csromancenter{นนฺทมาณวปุจฺฉานิทฺเทโส สตฺตโม.}
\csromanmatikaanchor{mtk.29}
\csromantocmarkref{subsection}{8. เหมกมาณวปุจฺฉานิทฺเทส}{mtk.29}
\bookmark[dest={mtk.29},level=2]{8. เหมกมาณวปุจฺฉานิทฺเทส}
\csromanheader{2}{8. เหมกมาณวปุจฺฉานิทฺเทส \textbf{8. เหมกมาณวปุจฺฉานิทฺเทส}}
\csromanhangingbegin
\hangingheaditem{53}{\csromanfolio{125}เย เม ปุพฺเพ วิยากํสุ, (\textbf{อิจฺจา}ยสฺมา \textbf{เหมโก},)}
\hangingbody{\textbf{หุรํ โคตมสาสนา}.}
\hangingbody{อิจฺจาสิ อิติ ภวิสฺสติ, สพฺพํ ตํ อิติหีติหํ.}
\hangingbody{\textbf{สพฺพํ ตํ ตกฺกวฑฺฒนํ, นาหํ ตตฺถ อภิรมิํ. (1)}}
\csromanhangingend

\prose{\textbf{เย เม ปุพฺเพ วิยากํสู}ติ โย จ พาวรี พฺราหฺมโณ เย จญฺเญ ตสฺส อาจริยา, เต สกํ ทิฏฺฐิํ สกํ ขนฺติํ สกํ รุจิํ สกํ ลทฺธิํ สกํ อชฺฌาสยํ สกํ อธิปฺปายํ พฺยากํสุ อาจิกฺขิํสุ เทสยิํสุ ปญฺญปิํสุ ปฏฺฐปิํสุ วิวริํสุ วิภชิํสุ อุตฺตานี-อกํสุ ปกาเสสุนฺติ เย เม ปุพฺเพ วิยากํสุ.\csromanspacer{} \textbf{อิจฺจายสฺมา เหมโก}ติ \textbf{อิจฺจา}ติ ปทสนฺธิ ฯเปฯ ปทานุปุพฺพตาเปตํ \textbf{อิจฺจา}ติ.\csromanspacer{} \textbf{อายสฺมา}ติ ปิยวจนํ ฯเปฯ.\csromanspacer{} \textbf{เหมโก}ติ ตสฺส พฺราหฺมณสฺส นามํ ฯเปฯ อภิลาโปติ \textbf{อิจฺจา}ยสฺมา \textbf{เหมโก}.}

\prose{\textbf{หุรํ โคตมสาสนา}ติ หุรํ โคตมสาสนา ปรํ โคตมสาสนา ปุเร โคตมสาสนา ปฐมตรํ โคตมสาสนา พุทฺธสาสนา ชินสาสนา ตถาคตสาสนา\footnote{ตถาคตสาสนา เทวสาสนา (ก)} อรหนฺตสาสนาติ หุรํ โคตมสาสนา.}

\prose{\textbf{อิจฺจาสิ อิติ ภวิสฺสตี}ติ เอวํ กิร อาสิ, เอวํ กิร ภวิสฺสตีติ \textbf{อิจฺจา}สิ อิติ ภวิสฺสติ.}

\prose{\textbf{สพฺพํ ตํ อิติหีติหนฺ}ติ สพฺพํ ตํ อิติหีติหํ อิติกิราย ปรํปราย ปิฏกสมฺปทาย ตกฺกเหตุ นยเหตุ อาการปริวิตกฺเกน ทิฏฺฐินิชฺฌานกฺขนฺติยา น สามํ สยมภิญฺญาตํ น อตฺตปจฺจกฺขธมฺมํ กถยิํสูติ สพฺพํ ตํ อิติหีติหํ.}

\prose{\textbf{สพฺพํ ตํ ตกฺกวฑฺฒนนฺ}ติ สพฺพํ ตํ ตกฺกวฑฺฒนํ วิตกฺกวฑฺฒนํ สงฺกปฺปวฑฺฒนํ กามวิตกฺกวฑฺฑฺธนํ พฺยาปาทวิตกฺกวฑฺฒนํ วิหิํสาวิตกฺกวฑฺฒนํ ญาติวิตกฺกวฑฺฒนํ ชนปทวิตกฺกวฑฺฒนํ อมราวิตกฺกวฑฺฒนํ ปรานุทยตาปฏิสํยุตฺตวิตกฺกวฑฺฒนํ ลาภสกฺการสิโลก- ปฏิสํยุตฺตวิตกฺกวฑฺฒนํ อนวญฺญตฺติปฏิสํยุตฺตวิตกฺกวฑฺฒนนฺติ สพฺพํ ตํ ตกฺกวฑฺฒนํ.}

\prose{\csromanfolio{126}\textbf{นาหํ ตตฺถ อภิรมินฺ}ติ นาหํ ตตฺถ อภิรมิํ น วินฺทิํ นาธิคจฺฉิํ น ปฏิลภินฺติ นาหํ ตตฺถ อภิรมิํ.\csromanspacer{} เตนาห โส พฺราหฺมโณ\csromandash{}}

\prose{เย เม ปุพฺเพ วิยากํสุ, (อิจฺจายสฺมา เหมโก,)}

\prose{หุรํ โคตมสาสนา.}

\csromangathameasure{อิจฺจาสิ อิติ ภวิสฺสติ, สพฺพํ ตํ อิติหีติหํ. \\ สพฺพํ ตํ ตกฺกวฑฺฒนํ, นาหํ ตตฺถ อภิรมินฺติ.}
\csromangathagroup{\csromangathastanza{อิจฺจาสิ อิติ ภวิสฺสติ, สพฺพํ ตํ อิติหีติหํ. \\ สพฺพํ ตํ ตกฺกวฑฺฒนํ, นาหํ ตตฺถ อภิรมินฺติ.}}

\csromanhangingbegin
\hangingheaditem{54}{ตฺวญฺจ เม ธมฺมมกฺขาหิ, \textbf{ตณฺหา}นิคฺฆาตนํ มุนิ.}
\hangingbody{ยํ วิทิตฺวา สโต จรํ, ตเร โลเก วิสตฺติกํ.}
\csromanhangingend

\prose{\textbf{ตฺวญฺจ เม ธมฺมมกฺขาหี}ติ \textbf{ตฺวนฺ}ติ ภควนฺตํ ภณติ.\csromanspacer{} \textbf{ธมฺมมกฺขาหี}ติ \textbf{ธมฺมนฺ}ติ อาทิกลฺยาณํ มชฺเฌกลฺยาณํ ปริโยสานกลฺยาณํ สาตฺถํ สพฺยญฺชนํ เกวลปริปุณฺณํ ปริสุทฺธํ พฺรหฺมจริยํ จตฺตาโร สติปฏฺฐาเน จตฺตาโร สมฺมปฺปธาเน จตฺตาโร อิทฺธิปาเท ปญฺจินฺทฺริยานิ ปญฺจ พลานิ สตฺต โพชฺฌงฺเค อริยํ อฏฺฐงฺคิกํ มคฺคํ นิพฺพานญฺจ นิพฺพานคามินิญฺจ ปฏิปทํ อกฺขาหิ อาจิกฺขาหิ เทเสหิ ปญฺญเปหิ ปฏฺฐเปหิ วิวราหิ วิภชาหิ อุตฺตานีกโรหิ ปกาเสหีติ ตฺวญฺจ เม ธมฺมมกฺขาหิ.}

\prose{\textbf{ตณฺหานิคฺฆาตนํ มุนี}ติ \textbf{ตณฺหา}ติ รูป\textbf{ตณฺหา} ฯเปฯ ธมฺม\textbf{ตณฺหา}.\csromanspacer{} ตณฺหานิคฺฆาตนํ \textbf{ตณฺหา}ปหานํ \textbf{ตณฺหา}วูปสมํ \textbf{ตณฺหา}ปฏินิสฺสคฺคํ \textbf{ตณฺหา}ปฏิปฺปสฺสทฺธิํ อมตํ นิพฺพานํ.\csromanspacer{} \textbf{มุนี}ติ โมนํ วุจฺจติ ญาณํ ฯเปฯ สงฺคชาลมติจฺจ โส \textbf{มุนี}ติ \textbf{ตณฺหา}นิคฺฆาตนํ มุนิ.}

\prose{\textbf{ยํ วิทิตฺวา สโต จรนฺ}ติ ยํ วิทิตํ กตฺวา ตุลยิตฺวา ตีรยิตฺวา วิภาวยิตฺวา วิภูตํ กตฺวา. “สพฺเพ สงฺขารา อนิจฺจา”ติ วิทิตํ กตฺวา ตุลยิตฺวา ตีรยิตฺวา วิภาวยิตฺวา วิภูตํ กตฺวา. “สพฺเพ สงฺขารา ทุกฺขา”ติ ฯเปฯ. “สพฺเพ ธมฺมา อนตฺตา”ติ. “ยํ กิญฺจิ สมุทยธมฺมํ, สพฺพํ ตํ นิโรธ\textbf{ธมฺมนฺ}”ติ วิทิตํ กตฺวา ตุลยิตฺวา ตีรยิตฺวา วิภาวยิตฺวา วิภูตํ กตฺวา.\csromanspacer{} \textbf{สโต}ติ จตูหิ การเณหิ สโต, กาเย กายานุปสฺสนาสติปฏฺฐานํ ภาเวนฺโต สโต ฯเปฯ โส วุจฺจติ สโต.\csromanspacer{} \textbf{จรนฺ}ติ จรนฺโต วิหรนฺโต อิริยนฺโต วตฺเตนฺโต ปาเลนฺโต ยเปนฺโต ยาเปนฺโตติ ยํ วิทิตฺวา สโต จรํ.}

\prose{\textbf{ตเร โลเก วิสตฺติกนฺ}ติ วิสตฺติกา วุจฺจติ \textbf{ตณฺหา}.\csromanspacer{} โย ราโค สาราโค ฯเปฯ อภิชฺฌา โลโภ อกุสลมูลํ, \textbf{วิสตฺติกา}ติ}

\vspace{\tipitakaparskip}
\csromancenter{เหมกมาณวปุจฺฉานิทฺเทส เกนฏฺเฐน วิสตฺ\textbf{ติ}กา ฯเปฯ วิสฏา วิตฺถตา\textbf{ติ} วิสตฺ\textbf{ติ}กา.\csromanspacer{} \textbf{โลเกติ} อปายโลเก มนุสฺสโลเก เทวโลเก, ขนฺธโลเก ธาตุโลเก อายตนโลเก.\csromanspacer{} \textbf{ตเร โลเก วิสตฺติกนฺติ} โลเก เวสา วิสตฺ\textbf{ติ}กา, โลเก เวตํ วิสตฺ\textbf{ติ}กํ สโต ตเรยฺยํ อุตฺตเรยฺยํ ปตเรยฺยํ สม\textbf{ติ}กฺกเมยฺยํ วี\textbf{ติ}วตฺเตยฺยนฺ\textbf{ติ} ตเร โลเก วิสตฺ\textbf{ติ}กํ.\csromanspacer{} เตนาห โส พฺราหฺมโณ\csromandash{}}
\vspace{0.5\baselineskip}
\csromangathameasure{ตฺวญฺจ เม ธมฺมมกฺขาหิ, ตณฺหานิคฺฆาตนํ มุนิ. \\ ยํ วิทิตฺวา สโต จรํ, ตเร โลเก วิสตฺติกนฺติ.}
\csromangathagroup{\csromangathastanza{\csromanfolio{127}ตฺวญฺจ เม ธมฺมมกฺขาหิ, ตณฺหานิคฺฆาตนํ มุนิ. \\ ยํ วิทิตฺวา สโต จรํ, ตเร โลเก วิสตฺติกนฺติ.}}

\csromanhangingbegin
\hangingheaditem{55}{อิธ ทิฏฺฐสุตมุตวิญฺญาเตสุ, ปิยรูเปสุ เหมก.}
\hangingbody{ฉนฺทราควิโนทนํ, นิพฺพานปทมจฺจุตํ.}
\csromanhangingend

\prose{อิธ ทิฏฺฐสุตมุตวิญฺญาเตสู\textbf{ติ} ทิฏฺฐนฺ\textbf{ติ} \textbf{จกฺขุนา ทิฏฺฐํ.}\csromanspacer{} สุตนฺ\textbf{ติ} \textbf{โสเตน สุตํ.}\csromanspacer{} มุตนฺ\textbf{ติ} \textbf{ฆาเนน ฆายิตํ, ชิวฺหาย สายิตํ, กาเยน ผุฏฺฐํ.}\csromanspacer{} วิญฺญาตนฺ\textbf{ติ} \textbf{มนสา วิญฺญาตนฺติ อิธ ทิฏฺฐสุตมุตวิญฺญาเตสุ.}}

\prose{\textbf{ปิยรูเปสุ เหมกาติ} กิญฺจ โลเก ปิยรูปํ สาตรูปํ, จกฺขุ\footnote{จกฺขุํ (สฺยา, ก)} โลเก ปิยรูปํ สาตรูปํ.\csromanspacer{} โสตํ โลเก.\csromanspacer{} ฆานํ โลเก.\csromanspacer{} ชิวฺหา โลเก.\csromanspacer{} กาโย โลเก.\csromanspacer{} มโน โลเก ปิยรูปํ สาตรูปํ.\csromanspacer{} รูปา โลเก ปิยรูปํ สาตรูปํ.\csromanspacer{} สทฺทา โลเก.\csromanspacer{} คนฺธา โลเก.\csromanspacer{} รสา โลเก.\csromanspacer{} โผฏฺฐพฺพา โลเก.\csromanspacer{} ธมฺมา โลเก ปิยรูปํ สาตรูปํ.\csromanspacer{} จกฺขุวิญฺญาณํ โลเก ปิยรูปํ สาตรูปํ.\csromanspacer{} โสตวิญฺญาณํ โลเก ปิยรูปํ สาตรูปํ.\csromanspacer{} ฆานวิญฺญาณํ โลเก.\csromanspacer{} ชิวฺหาวิญฺญาณํ โลเก.\csromanspacer{} กายวิญฺญาณํ โลเก.\csromanspacer{} มโนวิญฺญาณํ โลเก ปิยรูปํ สาตรูปํ.\csromanspacer{} จกฺขุสมฺผสฺโส โลเก.\csromanspacer{} โสตสมฺผสฺโส โลเก.\csromanspacer{} ฆานสมฺผสฺโส โลเก.\csromanspacer{} ชิวฺหาสมฺผสฺโส โลเก.\csromanspacer{} กายสมฺผสฺโส โลเก.\csromanspacer{} มโนสมฺผสฺโส โลเก ปิยรูปํ สาตรูปํ.\csromanspacer{} จกฺขุสมฺผสฺสชา เวทนา โลเก ปิยรูปํ สาตรูปํ.\csromanspacer{} โสตสมฺผสฺสชา เวทนา.\csromanspacer{} ฆานสมฺผสฺสชา เวทนา.\csromanspacer{} ชิวฺหาสมฺผสฺสชา เวทนา.\csromanspacer{} กายสมฺผสฺสชา เวทนา.\csromanspacer{} มโนสมฺผสฺสชา เวทนา โลเก ปิยรูปํ สาตรูปํ.\csromanspacer{} รูปสญฺญา โลเก.\csromanspacer{} สทฺทสญฺญา โลเก.\csromanspacer{} คนฺธสญฺญา โลเก.\csromanspacer{} รสสญฺญา โลเก.\csromanspacer{} โผฏฺฐพฺพสญฺญา โลเก.\csromanspacer{} ธมฺมสญฺญา โลเก ปิยรูปํ สาตรูปํ. \csromanfolio{128}รูปสญฺจฺ\textbf{เอตนฺ}อา โลเก.\csromanspacer{} สทฺทสญฺจฺ\textbf{เอตนฺ}อา โลเก.\csromanspacer{} คนฺธสญฺจฺ\textbf{เอตนฺ}อา โลเก.\csromanspacer{} รสสญฺจฺ\textbf{เอตนฺ}อา โลเก.\csromanspacer{} โผฏฺฐพฺพสญฺจฺ\textbf{เอตนฺ}อา โลเก.\csromanspacer{} ธมฺมสญฺจฺ\textbf{เอตนฺ}อา โลเก ปิยรูปํ สาตรูปํ.\csromanspacer{} รูปตณฺหา โลเก.\csromanspacer{} สทฺทตณฺหา โลเก.\csromanspacer{} คนฺธตณฺหา โลเก.\csromanspacer{} รสตณฺหา โลเก.\csromanspacer{} โผฏฺฐพฺพตณฺหา โลเก.\csromanspacer{} ธมฺมตณฺหา โลเก ปิยรูปํ สาตรูปํ.\csromanspacer{} รูปวิตกฺโก โลเก.\csromanspacer{} สทฺทวิตกฺโก โลเก.\csromanspacer{} คนฺธวิตกฺโก โลเก.\csromanspacer{} รสวิตกฺโก โลเก.\csromanspacer{} โผฏฺฐพฺพวิตกฺโก โลเก.\csromanspacer{} ธมฺมวิตกฺโก โลเก ปิยรูปํ สาตรูปํ.\csromanspacer{} รูปวิจาโร โลเก ปิยรูปํ สาตรูปํ.\csromanspacer{} สทฺทวิจาโร โลเก.\csromanspacer{} คนฺธวิจาโร โลเก.\csromanspacer{} รสวิจาโร โลเก.\csromanspacer{} โผฏฺฐพฺพวิจาโร โลเก.\csromanspacer{} ธมฺมวิจาโร โลเก ปิยรูปํ สาตรูปนฺติ ปิยรูเปสุ เหมก.}

\prose{\textbf{ฉนฺทราควิโนทนนฺ}ติ \textbf{ฉนฺทราโค}ติ โย กาเมสุ กามจฺฉนฺโท กามราโค กามนนฺที กามตณฺหา กามสิเนโห กามปริฬาโห กามมุจฺฉา กามชฺโฌสานํ กาโมโฆ กามโยโค กามุปาทานํ กามจฺฉนฺทนีวรณํ.\csromanspacer{} \textbf{ฉนฺทราควิโนทนนฺ}ติ ฉนฺทราคปฺปหานํ ฉนฺทราควูปสมํ ฉนฺทราคปฏินิสฺสคฺคํ ฉนฺทราคปฏิปฺปสฺสทฺธํ อมตํ นิพฺพานนฺติ ฉนฺทราควิโนทนํ.}

\prose{\textbf{นิพฺพานปทมจฺจุตนฺ}ติ นิพฺพานปทํ ตาณปทํ เลณปทํ สรณปทํ อภยปทํ.\csromanspacer{} \textbf{อจฺจุตนฺ}ติ นิจฺจํ ธุวํ สสฺสตํ อวิปริณามธมฺมนฺติ นิพฺพานปทมจฺจุตํ.\csromanspacer{} เตนาห ภควา\csromandash{}}

\csromangathameasure{อิธ ทิฏฺฐสุตมุตวิญฺญาเตสุ, ปิยรูเปสุ เหมก. \\ ฉนฺทราควิโนทนํ, นิพฺพานปทมจฺจุตนฺติ.}
\csromangathagroup{\csromangathastanza{อิธ ทิฏฺฐสุตมุตวิญฺญาเตสุ, ปิยรูเปสุ เหมก. \\ ฉนฺทราควิโนทนํ, นิพฺพานปทมจฺจุตนฺติ.}}

\proseitem{56}{\textbf{เอตทญฺญาย เย สตา}, ทิฏฺฐธมฺมาภินิพฺพุตา.}

\csromangathameasure{อุปสนฺตา จ เต สทา, ติณฺณา โลเก วิสตฺติกํ.}
\csromangathagroup{\csromangathastanza{อุปสนฺตา จ เต สทา, ติณฺณา โลเก วิสตฺติกํ.}}

\prose{\textbf{เอตทญฺญาย เย สตา}ติ \textbf{เอตนฺ}ติ อมตํ นิพฺพานํ.\csromanspacer{} โย โส สพฺพสงฺขารสมโถ สพฺพูปธิปฏินิสฺสคฺโค ตณฺหกฺขโย วิราโค นิโรโธ นิพฺพานํ.\csromanspacer{} \textbf{อญฺญายา}ติ อญฺญาย ชานิตฺวา ตุลยิตฺวา ตีรยิตฺวา วิภาวยิตฺวา วิภูตํ กตฺวา. “สพฺเพ สงฺขารา อนิจฺจา”ติ อญฺญาย ชานิตฺวา ตุลยิตฺวา ตีรยิตฺวา วิภาวยิตฺวา วิภูตํ กตฺวา. “สพฺเพ สงฺขารา ทุกฺขา”ติ. “สพฺเพ ธมฺมา อนตฺตา”ติ ฯเปฯ. “ยํ กิญฺจิ สมุทยธมฺมํ, สพฺพํ ตํ นิโรธธมฺมนฺ”ติ อญฺญาย ชานิตฺวา ตุลยิตฺวา ตีรยิตฺวา}

\vspace{\tipitakaparskip}
\csromancenter{เหมกมาณวปุจฺฉานิทฺเทส วิภาวยิตฺวา วิภูตํ กตฺวา.\csromanspacer{} \textbf{เย}ติ อรหนฺโต ขีณาสวา.\csromanspacer{} \textbf{สตา}ติ จตูหิ การเณหิ สตา.\csromanspacer{} กาเย กายานุปสฺสนาสติปฏฺฐานํ ภาวิตตฺตา สตา ฯเปฯ เต วุจฺจนฺติ สตาติ เอตทญฺญาย เย สตา.}
\prose{\csromanfolio{129}\textbf{ทิฏฺฐธมฺมาภินิพฺพุตา}ติ \textbf{ทิฏฺฐธมฺมา}ติ \textbf{ทิฏฺฐธมฺมา} ญาตธมฺมา ตุลิตธมฺมา ตีริตธมฺมา วิภูตธมฺมา วิภาวิตธมฺมา. “สพฺเพ สงฺขารา อนิจฺจา”ติ \textbf{ทิฏฺฐธมฺมา} ฯเปฯ. “ยํ กิญฺจิ สมุทยธมฺมํ, สพฺพํ ตํ นิโรธธมฺมนฺ”ติ \textbf{ทิฏฺฐธมฺมา} ญาตธมฺมา ตุลิตธมฺมา ตีริตธมฺมา วิภูตธมฺมา วิภาวิตธมฺมา.\csromanspacer{} \textbf{อภินิพฺพุตา}ติ ราคสฺส นิพฺพาปิตตฺตา นิพฺพุตา.\csromanspacer{} โทสสฺส นิพฺพาปิตตฺตา นิพฺพุตา.\csromanspacer{} โมหสฺส นิพฺพาปิตตฺตา นิพฺพุตา.\csromanspacer{} โกธสฺส.\csromanspacer{} อุปานาหสฺส ฯเปฯ สพฺพากุสลาภิสงฺขารานํ สนฺตตฺตา สมิตตฺตา วูปสมิตตฺตา นิชฺฌาตตฺตา นิพฺพุตตฺตา วิคตตฺตา ปฏิปฺปสฺสทฺธตฺตา สนฺตา \textbf{อุปสนฺตา} วูปสนฺตา นิพฺพุตา ปฏิปฺปสฺสทฺธาติ \textbf{ทิฏฺฐธมฺมา}ภินิพฺพุตา.}

\prose{\textbf{อุปสนฺตา จ เต สทา}ติ \textbf{อุปสนฺตา}ติ ราคสฺส อุปสมิตตฺตา นิพฺพาปิตตฺตา \textbf{อุปสนฺตา}.\csromanspacer{} โทสสฺส.\csromanspacer{} โมหสฺส.\csromanspacer{} โกธสฺส.\csromanspacer{} อุปนาหสฺส ฯเปฯ สพฺพากุสลาภิสงฺขารานํ สนฺตตฺตา สมิตตฺตา วูปสมิตตฺตา นิชฺฌาตตฺตา นิพฺพุตตฺตา วิคตตฺตา ปฏิปฺปสฺสทฺธตฺตา สนฺตา \textbf{อุปสนฺตา} วูปสนฺตา นิพฺพุตา ปฏิปฺปสฺสทฺธาติ \textbf{อุปสนฺตา}.\csromanspacer{} \textbf{เต}ติ อรหนฺโต กฺกฺ.\csromanspacer{} \textbf{สทา}ติ สทา สพฺพกาลํ นิจฺจกาลํ ธุวกาลํ สตตํ สมิตํ อพฺโพกิณฺณํ โปงฺขานุโปงฺขํ อุทกูมิกชาตํ อวีจิ สนฺตติ สหิตํ ผสฺสิตํ ปุเรภตฺตํ ปจฺฉาภตฺตํ ปุริมยามํ มชฺฌิมยามํ ปจฺฉิมยามํ กาเฬ ชุเณฺห วสฺเส เหมนฺเต คิเมฺห ปุริเม วโยขนฺเธ มชฺฌิเม วโยขนฺเธ ปจฺฉิเม วโยขนฺเธติ \textbf{อุปสนฺตา} จ เต สทา.}

\prose{\textbf{ติณฺณา โลเก วิสตฺติกนฺ}ติ วิสตฺติกา วุจฺจติ ตณฺหา.\csromanspacer{} โย ราโค สาราโค ฯเปฯ อภิชฺฌา โลโภ อกุสลมูลํ.\csromanspacer{} \textbf{วิสตฺติกา}ติ เกนฏฺเฐน วิสตฺติกา ฯเปฯ วิสฏา วิตฺถตาติ วิสตฺติกา.\csromanspacer{} \textbf{โลเก}ติ อปายโลเก ฯเปฯ อายตนโลเก.\csromanspacer{} \textbf{ติณฺณา โลเก วิสตฺติกนฺ}ติ โลเก เวสา วิสตฺติกา, โลเก เวตํ วิสตฺติกํ ติณฺณา อุตฺติณฺณา นิตฺถิณฺณา อติกฺกนฺตา สมติกฺกนฺตา วีติวตฺตาติ ติณฺณา โลเก วิสตฺติกํ.\csromanspacer{} \textbf{เต}นาห ภควา\csromandash{}}

\csromangathameasure{เอตทญฺญาย เย สตา, ทิฏฺฐธมฺมาภินิพฺพุตา. \\ อุปสนฺตา จ เต สทา, ติณฺณา โลเก วิสตฺติกนฺติ.}
\csromangathagroup{\csromangathastanza{\csromanfolio{130}เอตทญฺญาย เย สตา, ทิฏฺฐธมฺมาภินิพฺพุตา. \\ อุปสนฺตา จ เต สทา, ติณฺณา โลเก วิสตฺติกนฺติ.}}

\prose{สห คาถาปริโยสานา ฯเปฯ “สตฺถา เม ภนฺเต ภควา, สาวโกหมสฺมี”ติ.}

\vspace{\tipitakaparskip}
\csromancenter{เหมกมาณวปุจฺฉานิทฺเทโส อฏฺฐโม.}
\csromanmatikaanchor{mtk.30}
\csromantocmarkref{subsection}{9. โตเทยฺยมาณวปุจฺฉานิทฺเทส}{mtk.30}
\bookmark[dest={mtk.30},level=2]{9. โตเทยฺยมาณวปุจฺฉานิทฺเทส}
\csromanheader{2}{9. โตเทยฺยมาณวปุจฺฉานิทฺเทส}
\csromanhangingbegin
\hangingheaditem{57}{ยสฺมิํ กามา น วสนฺติ, (\textbf{อิจฺจา}ยสฺมา \textbf{โตเทยฺโย},)}
\hangingbody{\textbf{ตณฺหา ยสฺส น วิชฺชติ.}}
\hangingbody{\textbf{กถํกถา จ โย ติณฺโณ,}}
\hangingbody{\textbf{วิโมกฺโข ตสฺส กีทิโส}.}
\csromanhangingend

\prose{\textbf{ยสฺมิํ กามา น วสนฺตี}ติ ยสฺมิํ กามา น วสนฺติ น สํวสนฺติ น อาวสนฺติ น ปริวสนฺตีติ ยสฺมิํ กามา น วสนฺติ.\csromanspacer{} \textbf{อิจฺจายสฺมา โตเทยฺโย}ติ \textbf{อิจฺจา}ติ’ปทสนฺธิ ฯเปฯ ปทานุปุพฺพตาเปตํ \textbf{อิจฺจา}ติ.\csromanspacer{} \textbf{อายสฺมา}ติ ปิยวจนํ ฯเปฯ.\csromanspacer{} \textbf{โตเทยฺโย}ติ ตสฺส พฺราหฺมณสฺส นามํ ฯเปฯ อภิลาโปติ \textbf{อิจฺจา}ยสฺมา \textbf{โตเทยฺโย}.}

\prose{\textbf{ตณฺหา ยสฺส น วิชฺชตี}ติ ตณฺหา ยสฺส นตฺถิ น สติ น สํวิชฺชติ นุปลพฺภติ ญาณคฺคินา ทฑฺฒาติ ตณฺหา ยสฺส น วิชฺชติ.}

\prose{\textbf{กถํกถา จ โย ติณฺโณ}ติ กถํกถา จ โย ติณฺโณ อุตฺติณฺโณ นิตฺถิณฺโณ อติกฺกนฺโต สมติกฺกนฺโต วีติวตฺโตติ กถํกถา จ โย ติณฺโณ.}

\prose{\textbf{วิโมกฺโข ตสฺส กีทิโส}ติ วิโมกฺโข ตสฺส กีทิโส กิํสณฺฐิโต กิํปกาโร กิํปฏิภาโค อิจฺฉิตพฺโพติ วิโมกฺขํ ปุจฺฉตีติ วิโมกฺโข ตสฺส กีทิโส.\csromanspacer{} เตนาห โส พฺราหฺมโณ\csromandash{}}

\prose{ยสฺมิํ กามา น วสนฺติ, (\textbf{อิจฺจา}ยสฺมา \textbf{โตเทยฺโย},)}

\prose{\textbf{ตณฺหา ยสฺส น วิชฺชติ.}}

\csromangathameasure{\textbf{กถํกถา จ โย ติณฺโณ,} \\ \textbf{วิโมกฺโข ตสฺส กีทิโส}ติ.}
\csromangathagroup{\csromangathastanza{\textbf{กถํกถา จ โย ติณฺโณ,} \\ \textbf{วิโมกฺโข ตสฺส กีทิโส}ติ.}}

\csromanheader{2}{9. โตเทยฺยมาณวปุจฺฉานิทฺเทส}
\csromanhangingbegin
\hangingheaditem{58}{\csromanfolio{131}ยสฺมิํ กามา น วสนฺติ, (\textbf{โตเทยฺยาติ ภควา},)}
\hangingbody{\textbf{ตณฺหา ยสฺส น วิชฺชติ.}}
\hangingbody{\textbf{กถํกถา จ โย ติณฺโณ,}}
\hangingbody{\textbf{วิโมกฺโข ตสฺส นา’ปโร}.}
\csromanhangingend

\prose{\textbf{ยสฺมิํ กามา น วสนฺตี}ติ \textbf{ยสฺมินฺ}ติ ยสฺมิํ ปุคฺคเล อรหนฺเต ขีณาสเว.\csromanspacer{} \textbf{กามา}ติ อุทฺทานโต เทฺว กามา วตฺถุกามา จ กิเลสกามา จ ฯเปฯ.\csromanspacer{} อิเม วุจฺจนฺติ วตฺถุกามา ฯเปฯ.\csromanspacer{} อิเม วุจฺจนฺติ กิเลสกามา.\csromanspacer{} \textbf{ยสฺมิํ กามา น วสนฺตี}ติ ยสฺมิํ กามา น วสนฺติ น สํวสนฺติ น อาวสนฺติ น ปริวสนฺตีติ ยสฺมิํ กามา น วสนฺติ.}

\prose{\textbf{โตเทยฺยาติ ภควา}ติ \textbf{โตเทยฺยาติ ภควา} ตํ พฺราหฺมณํ นาเมน อาลปติ.\csromanspacer{} \textbf{ภควา}ติ คารวาธิวจนเมตํ ฯเปฯ สจฺฉิกา ปญฺญตฺติ, ยทิทํ \textbf{ภควา}ติ \textbf{โตเทยฺยาติ ภควา}.}

\prose{\textbf{ตณฺหา ยสฺส น วิชฺชตี}ติ \textbf{ตณฺหา}ติ รูป\textbf{ตณฺหา} สทฺท\textbf{ตณฺหา} คนฺธ\textbf{ตณฺหา} รส\textbf{ตณฺหา} โผฏฺฐพฺพ\textbf{ตณฺหา} ธมฺม\textbf{ตณฺหา}.\csromanspacer{} \textbf{ยสฺสา}ติ อรหโต ขีณาสวสฺส.\csromanspacer{} \textbf{ตณฺหา ยสฺส น วิชฺชตี}ติ \textbf{ตณฺหา} ยสฺส นตฺถิ น อตฺถิ น สํวิชฺชติ นุปลพฺภติ, ปหีนา สมุจฺฉินฺนา วูปสนฺตา ปฏิปฺปสฺสทฺธา อภพฺพุปฺปตฺติกา ญาณคฺคินา ทฑฺฒาติ \textbf{ตณฺหา} ยสฺส น วิชฺชติ.}

\prose{\textbf{กถํกถา จ โย ติณฺโณ}ติ กถํกถา วุจฺจติ วิจิกิจฺฉา.\csromanspacer{} ทุกฺเข กงฺขา ฯเปฯ ฉมฺภิตตฺตํ จิตฺตสฺส มโนวิเลโข.\csromanspacer{} \textbf{โย}ติ โย โส อรหํ ขีณาสโว.\csromanspacer{} \textbf{กถํกถา จ โย ติณฺโณ}ติ กถํกถา จ โย ติณฺโณ อุตฺติณฺโณ นิตฺถิณฺโณ อติกฺกนฺโต สมติกฺกนฺโต วีติวตฺโตติ กถํกถา จ โย ติณฺโณ.}

\prose{\textbf{วิโมกฺโข ตสฺส นา’ปโร}ติ นตฺถิ ตสฺส อปโร วิโมกฺโข.\csromanspacer{} เยน วิโมกฺเขน วิมุจฺเจยฺย วิมุตฺโต โส.\csromanspacer{} กตํ ตสฺส วิโมกฺเขน กรณียนฺติ วิโมกฺโข ตสฺส นา’ปโร.\csromanspacer{} เตนาห \textbf{ภควา}\csromandash{}}

\prose{ยสฺมิํ กามา น วสนฺติ, (\textbf{โตเทยฺยาติ ภควา},)}

\prose{\textbf{ตณฺหา ยสฺส น วิชฺชติ.}}

\csromangathameasure{\textbf{กถํกถา จ โย ติณฺโณ,} \\ \textbf{วิโมกฺโข ตสฺส นา’ปโร}ติ.}
\csromangathagroup{\csromangathastanza{\textbf{กถํกถา จ โย ติณฺโณ,} \\ \textbf{วิโมกฺโข ตสฺส นา’ปโร}ติ.}}

\csromanhangingbegin
\hangingheaditem{59}{\csromanfolio{132}\textbf{นิราสโส โส อุท อาสสาโน},}
\hangingbody{\textbf{ปญฺญาณวา โส อุท ปญฺญกปฺปี.}}
\hangingbody{\textbf{มุนิํ อหํ สกฺก ยถา วิชญฺญํ,}}
\hangingbody{ตํ เม วิยาจิกฺข สมนฺตจกฺขุ.}
\csromanhangingend

\prose{\textbf{นิราสโส โส อุท อาสสาโน}ติ นิตฺตโณฺห โส, อุทาหุ สตโณฺห รูเป อาสีสติ\footnote{อาสิํสติ (สฺยา)}.\csromanspacer{} สทฺเท.\csromanspacer{} คนฺเธ.\csromanspacer{} รเส.\csromanspacer{} โผฏฺฐพฺเพ.\csromanspacer{} กุลํ.\csromanspacer{} คณํ.\csromanspacer{} อาวาสํ.\csromanspacer{} ลาภํ.\csromanspacer{} ยสํ.\csromanspacer{} ปสํสํ.\csromanspacer{} สุขํ.\csromanspacer{} จีวรํ.\csromanspacer{} ปิณฺฑปาตํ.\csromanspacer{} เสนาสนํ.\csromanspacer{} คิลานปจฺจยเภสชฺชปริกฺขารํ.\csromanspacer{} กามธาตุํ.\csromanspacer{} รูปธาตุํ.\csromanspacer{} อรูปธาตุํ.\csromanspacer{} กามภวํ.\csromanspacer{} รูปภวํ.\csromanspacer{} อรูปภวํ.\csromanspacer{} สญฺญาภวํ.\csromanspacer{} อสญฺญาภวํ.\csromanspacer{} เนวสญฺญานาสญฺญาภวํ.\csromanspacer{} เอกโวการภวํ.\csromanspacer{} จตุโวการภวํ.\csromanspacer{} ปญฺจโวการภวํ.\csromanspacer{} อตีตํ.\csromanspacer{} อนาคตํ.\csromanspacer{} ปจฺจุปฺปนฺนํ.\csromanspacer{} ทิฏฺฐสุตมุตวิญฺญาตพฺเพ ธมฺเม อาสีสติ สาทิยติ ปตฺเถติ ปิเหติ อภิชปฺปตีติ นิราสโส โส อุท อาสสาโน.}

\prose{\textbf{ปญฺญาณวา โส อุท ปญฺญกปฺปี}ติ \textbf{ปญฺญาณวา โส}ติ ปณฺฑิโต ปญฺญวา พุทฺธิมา ญาณี วิภาวี เมธาวี.\csromanspacer{} \textbf{อุท ปญฺญกปฺปี}ติ อุทาหุ อฏฺฐสมาปตฺติญาเณน วา ปญฺจาภิญฺญาญาเณน วา มิจฺฉาญาเณน วา ตณฺหากปฺปํ วา ทิฏฺฐิกปฺปํ วา กปฺเปติ ชเนติ สญฺชเนติ นิพฺพตฺเตติ อภินิพฺพตฺเตตีติ ปญฺญณวา โส อุท ปญฺญกปฺปี.}

\prose{\textbf{มุนิํ อหํ สกฺก ยถา วิชญฺญนฺ}ติ \textbf{สกฺกา}ติ สกฺโก, ภควา สกฺยกุลา ปพฺพชิโตติปิ สกฺโก.\csromanspacer{} อถ วา อฑฺโฒ มหทฺธโน ธนวาติปิ สกฺโก.\csromanspacer{} ตสฺสิมานิ ธนานิ.\csromanspacer{} เสยฺยถิทํ, สทฺธาธนํ สีลธนํ หิริธนํ โอตฺตปฺปธนํ สุตธนํ จาคธนํ ปญฺญาธนํ สติปฏฺฐานธนํ สมฺมปฺปธานธนํ อิทฺธิปาทธนํ อินฺทฺริยธนํ พลธนํ โพชฺฌงฺคธนํ มคฺคธนํ ผลธนํ นิพฺพานธนนฺติ, เตหิ อเนกวิเธหิ ธนรตเนหิ อฑฺโฒ มหทฺธโน ธนวาติปิ สกฺโก.\csromanspacer{} อถ วา ปหุ วิสวี อลมตฺโต สูโร วีโร วิกฺกนฺโต อภีรู อจฺฉมฺภี อนุตฺราสี อปลายี ปหีนภยเภรโว วิคตโลมหํโสติปิ สกฺโก.\csromanspacer{} \textbf{มุนิํ อหํ สกฺก ยถา วิชญฺญนฺ}ติ สกฺก ยถาหํ มุนิํ ชาเนยฺยํ อาชาเนยฺยํ วิชาเนยฺยํ ปฏิวิชาเนยฺยํ ปฏิวิชฺเฌยฺยนฺติ มุนิํ อหํ สกฺก ยถา วิชญฺญํ.}

\csromanheader{2}{9. โตเทยฺยมาณวปุจฺฉานิทฺเทส}
\prose{\csromanfolio{133}\textbf{ตํ เม วิยาจิกฺข สมนฺตจกฺขู}ติ \textbf{ตนฺ}ติ ยํ ปุจฺฉามิ, ยํ ยาจามิ, ยํ อชฺเฌสามิ, ยํ ปสาเทมิ.\csromanspacer{} \textbf{วิยาจิกฺขา}ติ อาจิกฺขาหิ เทเสหิ ปญฺญเปหิ ปฏฺฐเปหิ วิวราหิ วิภชาหิ อุตฺตานีกโรหิ ปกาเสหิ.\csromanspacer{} \textbf{สมนฺตจกฺขู}ติ สมนฺตจกฺขุ วุจฺจติ สพฺพญฺญุตญาณํ ฯเปฯ ตถาคโต เตน \textbf{สมนฺตจกฺขู}ติ ตํ เม วิยาจิกฺข สมนฺตจกฺขุ.\csromanspacer{} เตนาห โส พฺราหฺมโณ\csromandash{}}

\csromangathameasure{นิราสโส โส อุท อาสสาโน, \\ ปญฺญาณวา โส อุท ปญฺญกปฺปี. \\ มุนิํ อหํ สกฺก ยถา วิชญฺญํ,}
\csromangathagroup{\csromangathastanza{นิราสโส โส อุท อาสสาโน, \\ ปญฺญาณวา โส อุท ปญฺญกปฺปี. \\ มุนิํ อหํ สกฺก ยถา วิชญฺญํ,}}

\prose{\textbf{ตํ เม วิยาจิกฺข สมนฺตจกฺขู}ติ.}

\csromanhangingbegin
\hangingheaditem{60}{\textbf{นิราสโส โส น จ อาสสาโน},}
\hangingbody{\textbf{ปญฺญาณวา โส น จ ปญฺญกปฺปี.}}
\hangingbody{\textbf{เอวมฺปิ โตเทยฺย มุนิํ วิชาน,}}
\hangingbody{อกิญฺจนํ กามภเว อสตฺตํ.}
\csromanhangingend

\prose{\textbf{นิราสโส โส น จ อาสสาโน}ติ นิตฺตโณฺห โส, น โส สตโณฺห รูเป นาสีสติ.\csromanspacer{} สทฺเท.\csromanspacer{} คนฺเธ ฯเปฯ ทิฏฺฐสุตมุตวิญฺญาตพฺเพ ธมฺเม นาสีสติ น อิจฺฉติ น สาทิยติ น ปตฺเถติ น ปิเหติ นาภิชปฺปตีติ นิราสโส โส น จ อาสสาโน.}

\prose{\textbf{ปญฺญาณวา โส น จ ปญฺญกปฺปี}ติ \textbf{ปญฺญาณวา}ติ ปณฺฑิโต ปญฺญวา พุทฺธิมา ญาณี วิภาวี เมธาวี.\csromanspacer{} \textbf{น จ ปญฺญกปฺปี}ติ อฏฺฐสมาปตฺติญาเณน วา ปญฺจาภิญฺญาญาเณน วา มิจฺฉาญาเณน วา ตณฺหากปฺปํ วา น กปฺเปติ ทิฏฺฐิกปฺปํ วา น กปฺเปติ น ชเนติ น สญฺชเนติ น นิพฺพตฺเตติ นาภินิพฺพตฺเตตีติ \textbf{ปญฺญาณวา} โส น จ ปญฺญกปฺปี.}

\prose{\textbf{เอวมฺปิ โตเทยฺย มุนิํ วิชานา}ติ \textbf{มุนี}ติ โมนํ วุจฺจติ ญาณํ ฯเปฯ สงฺคชาลมติจฺจ โส มุนิ.\csromanspacer{} \textbf{เอวมฺปิ โตเทยฺย มุนิํ วิชานา}ติ โตเทยฺย เอวํ มุนิํ ชาน ปฏิชาน ปฏิวิชาน ปฏิวิชฺฌาติ เอวมฺปิ โตเทยฺย มุนิํ วิชาน.}

\prose{\textbf{อกิญฺจนํ กามภเว อสตฺตนฺ}ติ \textbf{อกิญฺจนนฺ}ติ ราคกิญฺจนํ โทสกิญฺจนํ โมหกิญฺจนํ มานกิญฺจนํ ทิฏฺฐิกิญฺจนํ กิเลสกิญฺจนํ ทุจฺจริตกิญฺจนํ.\csromanspacer{} ยสฺเสตานิ\footnote{ยสฺเสเต (สฺยา)} \csromanfolio{134}กิญฺจนานิ\footnote{กิญฺจนา (สฺยา)} ปหีนานิ สมุจฺฉินฺนานิ วูปสนฺตานิ ปฏิปฺปสฺสทฺธานิ อภพฺพุปฺปตฺติกานิ ญาณคฺคินา ทฑฺฒานิ.\csromanspacer{} โส วุจฺจติ อกิญฺจโน.\csromanspacer{} \textbf{กามา}ติ อุทฺทานโต เทฺว กามา วตฺถุกามา จ กิเลสกามา จ ฯเปฯ.\csromanspacer{} อิเม วุจฺจนฺติ วตฺถุกามา ฯเปฯ.\csromanspacer{} อิเม วุจฺจนฺติ กิเลสกามา.\csromanspacer{} \textbf{ภวา}ติ เทฺว ภวา กมฺมภโว จ ปฏิสนฺธิโก จ ปุนพฺภโว ฯเปฯ อยํ ปฏิสนฺธิโก ปุนพฺภโว.}

\prose{\textbf{อกิญฺจนํ กามภเว อสตฺตนฺ}ติ อกิญฺจนํ ปุคฺคลํ กาเม จ ภเว จ อสตฺตํ อลคฺคํ อลคฺคิตํ อปลิพุทฺธํ นิกฺขนฺตํ นิสฺสฏํ วิปฺปมุตฺตํ วิสญฺญุตฺตํ วิมริยาทิกเตน เจตสา วิหรนฺตนฺติ อกิญฺจนํ กามภเว อสตฺตํ.\csromanspacer{} เตนาห ภควา\csromandash{}}

\csromangathameasure{นิราสโส โส นา จ อาสสาโน, \\ ปญฺญาณวา โส น จ ปญฺญกปฺปี. \\ เอวมฺปิ โตเทยฺย มุนิํ วิชาน,}
\csromangathagroup{\csromangathastanza{นิราสโส โส นา จ อาสสาโน, \\ ปญฺญาณวา โส น จ ปญฺญกปฺปี. \\ เอวมฺปิ โตเทยฺย มุนิํ วิชาน,}}

\prose{\textbf{อกิญฺจนํ กามภเว อสตฺตนฺ}ติ.}

\prose{สห คาถาปริโยสานา ฯเปฯ “สตฺถา เม ภนฺเต ภควา, สาวโก หมสฺมี”ติ.}

\vspace{\tipitakaparskip}
\csromancenter{โตเทยฺยมาณวปุจฺฉานิทฺเทโส นวโม.}
\csromanmatikaanchor{mtk.31}
\csromantocmarkref{subsection}{10. กปฺปมาณวปุจฺฉานิทฺเทส}{mtk.31}
\bookmark[dest={mtk.31},level=2]{10. กปฺปมาณวปุจฺฉานิทฺเทส}
\csromanheader{2}{10. กปฺปมาณวปุจฺฉานิทฺเทส}
\csromanhangingbegin
\hangingheaditem{61}{มชฺเฌ สรสฺมิํ ติฏฺฐตํ, (อิจฺจายสฺมา กปฺโป,)}
\hangingbody{\textbf{โอเฆ ชาเต มหพฺภเย.}}
\hangingbody{\textbf{ชรามจฺจุปเรตานํ, ทีปํ ปพฺรูหิ มาริส.}}
\hangingbody{ตฺวญฺจ เม ทีปมกฺขาหิ, ยถายิทํ นาปรํ สิยา.}
\csromanhangingend

\prose{\textbf{มชฺเฌ สรสฺมิํ ติฏฺฐตนฺ}ติ สโร วุจฺจติ สํสาโร อาคมนํ คมนํ คมนาคมนํ กาลํ คติ ภวาภโว จุติ จ อุปปตฺติ จ นิพฺพตฺติ จ เภโท จ ชาติ จ ชรา จ มรณญฺจ.\csromanspacer{} สํสารสฺส ปุริมาปิ โกฏิ น ปญฺญายติ, ปจฺฉิมาปิ โกฏิ น ปญฺญายติ.\csromanspacer{} มชฺเฌว สํสาเร สตฺตา ฐิตา ปติฏฺฐิตา อลฺลีนา อุปคตา อชฺโฌสิตา อธิมุตฺตา.}

\csromanheader{2}{10. กปฺปมาณวปุจฺฉานิทฺเทส}
\prose{\csromanfolio{135}กถํ สํสารสฺส ปุริมา โกฏิ น ปญฺญายติ, เอตฺตกา ชาติโย วฏฺฏํ วตฺติ, ตโต ปรํ น วตฺตตีติ เหวํ นตฺถิ, เอวมฺปิ สํสารสฺส ปุริมา โกฏิ น ปญฺญายติ.\csromanspacer{} เอตฺตกานิ ชาติสตานิ วฏฺฏํ วตฺติ, ตโต ปรํ น วตฺตตีติ เหวํ นตฺถิ, เอวมฺปิ สํสารสฺส ปุริมา โกฏิ น ปญฺญายติ.\csromanspacer{} เอตฺตกานิ ชาติสหสฺสานิ วฏฺฏํ วตฺติ, ตโต ปรํ น วตฺตตีติ เหวํ นตฺถิ, เอวมฺปิ สํสารสฺส ปุริมา โกฏิ น ปญฺญายติ.\csromanspacer{} เอตฺตกานิ ชาติสตสหสฺสานิ วฏฺฏํ วตฺติ, ตโต ปรํ น วตฺตตีติ เหวํ นตฺถิ, เอวมฺปิ สํสารสฺส ปุริมา โกติ น ปญฺญายติ.\csromanspacer{} เอตฺตกา ชาติโกฏิโย วฏฺฏํ วตฺติ, ตโต ปรํ น วตฺตตีติ เหวํ นตฺถิ, เอวมฺปิ สํสารสฺส ปุริมา โกฏิ น ปญฺญายติ.\csromanspacer{} เอตฺตกานิ ชาติโกฏิสตานิ วฏฺฏํ วตฺติ, ตโต ปรํ น วตฺตตีติ เหวํ นตฺถิ, เอวมฺปิ สํสารสฺส ปุริมา โกฏิ น ปญฺญายติ.\csromanspacer{} เอตฺตกานิ ชาติโกฏิสหสฺสานิ วฏฺฏํ วตฺติ, ตโต ปรํ น วตฺตตีติ เหวํ นตฺถิ, เอวมฺปิ สํสารสฺส ปุริมา โกฏิ น ปญฺญายติ.\csromanspacer{} เอตฺตกานิ ชาติโกฏิสตสหสฺสานิ วฏฺฏํ วตฺติ, ตโต ปรํ น วตฺตตีติ เหวํ นตฺถิ, เอวมฺปิ สํสารสฺส ปุริมา โกฏิ น ปญฺญายติ.}

\prose{เอตฺตกานิ วสฺสานิ วฏฺฏํ วตฺติ, ตโต ปรํ น วตฺตตีติ เหวํ นตฺถิ, เอวมฺปิ สํสารสฺส ปุริมา โกฏิ น ปญฺญายติ.\csromanspacer{} เอตฺตกานิ วสฺสสตานิ วฏฺฏํ วตฺติ, ตโต ปรํ น วตฺตตีติ เหวํ นตฺถิ, เอวมฺปิ สํสารสฺส ปุริมา โกฏิ น ปญฺญายติ.\csromanspacer{} เอตฺตกานิ วสฺสสหสฺสานิ วฏฺฏํ วตฺติ, ตโต ปรํ น วตฺตตีติ เหวํ นตฺถิ, เอวมฺปิ สํสารสฺส ปุริมา โกฏิ น ปญฺญายติ.\csromanspacer{} เอตฺตกานิ วสฺสสตสหสฺสานิ วฏฺฏํ วตฺติ, ตโต ปรํ น วตฺตตีติ เหวํ นตฺถิ, เอวมฺปิ สํสารสฺส ปุริมา โกฏิ น ปญฺญายติ.\csromanspacer{} เอตฺตกา วสฺสโกฏิโย วฏฺฏํ วตฺติ, ตโต ปรํ น วตฺตตีติ เหวํ นตฺถิ, เอวมฺปิ สํสารสฺส ปุริมา โกฏิ น ปญฺญายติ.\csromanspacer{} เอตฺตกานิ วสฺสโกฏิสตานิ วฏฺฏํ วตฺติ, ตโต ปรํ น วตฺตตีติ เหวํ นตฺถิ, เอวมฺปิ สํสารสฺส ปุริมา โกฏิ น ปญฺญายติ.\csromanspacer{} เอตฺตกานิ วสฺสโกฏิสหสฺสานิ วฏฺฏํ วตฺติ, ตโต ปรํ น วตฺตตีติ เหวํ นตฺถิ, เอวมฺปิ สํสารสฺส ปุริมา โกฏิ น ปญฺญายติ.\csromanspacer{} เอตฺตกานิ วสฺสโกฏิสตสหสฺสานิ วฏฺฏํ วตฺติ, ตโต ปรํ น วตฺตตีติ เหวํ นตฺถิ, เอวมฺปิ สํสารสฺส ปุริมา โกฏิ น ปญฺญายติ.}

\prose{เอตฺตกานิ กปฺปานิ วฏฺฏํ วตฺติ, ตโต ปรํ น วตฺตตีติ เหวํ นตฺถิ, เอวมฺปิ สํสารสฺส ปุริมา โกฏิ น ปญฺญายติ.\csromanspacer{} เอตฺตกานิ กปฺปสตานิ}

\prose{\csromanfolio{136}วฏฺฏํ วตฺติ, ตโต ปรํ น วตฺตตีติ เหวํ นตฺถิ, เอวมฺปิ สํสารสฺส ปุริมา โกฏิ น ปญฺญายติ.\csromanspacer{} เอตฺตกานิ กปฺปสหสฺสานิ วฏฺฏํ วตฺติ, ตโต ปรํ น วตฺตตีติ เหวํ นตฺถิ, เอวมฺปิ สํสารสฺส ปุริมา โกฏิ น ปญฺญายติ.\csromanspacer{} เอตฺตกานิ กปฺปสตสหสฺสานิ วฏฺฏํ วตฺติ, ตโต ปรํ น วตฺตตีติ เหวํ นตฺถิ, เอวมฺปิ สํสารสฺส ปุริมา โกฏิ น ปญฺญายติ.\csromanspacer{} เอตฺตกา กปฺปโกฏิโย วฏฺฏํ วตฺติ, ตโต ปรํ น วตฺตตีติ เหวํ นตฺถิ, เอวมฺปิ สํสารสฺส ปุริมา โกฏิ น ปญฺญายติ.\csromanspacer{} เอตฺตกานิ กปฺปโกฏิสตานิ วฏฺฏํ วตฺติ, ตโต ปรํ น วตฺตตีติ เหวํ นตฺถิ, เอวมฺปิ สํสารสฺส ปุริมา โกฏิ น ปญฺญายติ.\csromanspacer{} เอตฺตกานิ กปฺปโกฏิสหสฺสานิ วฏฺฏํ วตฺติ, ตโต ปรํ น วตฺตตีติ เหวํ นตฺถิ.\csromanspacer{} เอวมฺปิ สํสารสฺส ปุริมา โกฏิ น ปญฺญายติ.\csromanspacer{} เอตฺตกานิ กปฺปโกฏิสตสหสฺสานิ วฏฺฏํ วตฺติ, ตโต ปรํ น วตฺตตีติ เหวํ นตฺถิ, เอวมฺปิ สํสารสฺส ปุริมา โกฏิ น ปญฺญายติ.}

\prose{วุตฺตํ เหตํ ภควตา\csromandash{}อนมตคฺโคยํ ภิกฺขเว สํสาโร ปุพฺพา โกฏิ น ปญฺญายติ อวิชฺชานีวรณานํ สตฺตานํ ตณฺหาสํโยชนานํ สนฺธาวตํ สํสรตํ.\csromanspacer{} เอวํ ทีฆรตฺตํ โข ภิกฺขเว ทุกฺขํ ปจฺจนุภูตํ ติพฺพํ ปจฺจนุภูตํ พฺยสนํ ปจฺจนุภูตํ กฏสี วฑฺฒิตา\footnote{กฏสีววฑฺฒิตํ (สฺยา) ปสฺส สํ 1. 387 ปิฏฺเฐ.}, ยาวญฺจิทํ ภิกฺขเว อลเมว สพฺพสงฺขาเรสุ นิพฺพินฺทิตุํ, อลํ วิรชฺชิตุํ, อลํ วิมุจฺจิตุนฺติ เอวมฺปิ สํสารสฺส ปุริมา โกฏิ น ปญฺญายติ.}

\prose{กถํ สํสารสฺส ปจฺฉิมา โกฏิ น ปญฺญายติ.\csromanspacer{} เอตฺตกา ชาติโย วฏฺฏํ วตฺติสฺสติ, ตโต ปรํ น วตฺติสฺสตีติ เหวํ นตฺถิ, เอวมฺปิ สํสารสฺส ปจฺฉิมา โกฏิ น ปญฺญายติ.\csromanspacer{} เอตฺตกานิ ชาติสตานิ.\csromanspacer{} เอตฺตกานิ ชาติสหสฺสานิ.\csromanspacer{} เอตฺตกานิ ชาติสตสหสฺสานิ ฯเปฯ.\csromanspacer{} เอตฺตกา ชาติโกฏิโย.\csromanspacer{} เอตฺตกานิ ชาติโกฏิสตานิ.\csromanspacer{} เอตฺตกานิ ชาติโกฏิสหสฺสานิ.\csromanspacer{} เอตฺตกานิ ชาติโกฏิสตสหสฺสานิ.\csromanspacer{} เอตฺตกานิ วสฺสานิ.\csromanspacer{} เอตฺตกานิ วสฺสาสตานิ.\csromanspacer{} เอตฺตกานิ วสฺสสหสฺสานิ.\csromanspacer{} เอตฺตกานิ วสฺสสตสหสฺสานิ.\csromanspacer{} เอตฺตกา วสฺสโกฏิโย.\csromanspacer{} เอตฺตกานิ วสฺสโกฏิสตานิ.\csromanspacer{} เอตฺตกานิ วสฺสโกฏิสหสฺสานิ.\csromanspacer{} เอตฺตกานิ วสฺสโกฏิสตสหสฺสานิ.\csromanspacer{} เอตฺตกานิ กปฺปานิ.\csromanspacer{} เอตฺตกานิ กปฺปสตานิ.\csromanspacer{} เอตฺตกานิ กปฺปสหสฺสานิ.\csromanspacer{} เอตฺตกานิ กปฺปสตสหสฺสานิ.\csromanspacer{} เอตฺตกา กปฺปโกฏิโย.\csromanspacer{} เอตฺตกานิ}

\vspace{\tipitakaparskip}
\csromancenter{กปฺปมาณวปุจฺฉานิทฺเทส กปฺปโกฏิสตานิ.\csromanspacer{} เอตฺตกานิ กปฺปโกฏิสหสฺสานิ.\csromanspacer{} เอตฺตกานิ กปฺปโกฏิสตสหสฺสานิ วฏฺฏํ วตฺติสฺสติ.\csromanspacer{} ตโต ปรํ น วตฺติสฺสตีติ เหวํ นตฺถิ, เอวมฺปิ สํสารสฺส ปจฺฉิมา โกฏิ น ปญฺญายติ.\csromanspacer{} เอวมฺปิ สํสารสฺส ปุริมาปิ โกฏิ น ปญฺญายติ, ปจฺฉิมาปิ โกฏิ น ปญฺญายติ.\csromanspacer{} มชฺเฌว สํสาเร สตฺตา ฐิตา ปติฏฺฐิตา อลฺลีนา อุปคตา อชฺโฌสิตา อธิมุตฺตาติ มชฺเฌ สรสฺมิํ ติฏฺฐตํ.\csromanspacer{} \textbf{อิจฺจายสฺมา กปฺโป}ติ \textbf{อิจฺจา}ติ ปทสนฺธิ ฯเปฯ.\csromanspacer{} \textbf{อายสฺมา}ติ ปิยวจนํ ฯเปฯ.\csromanspacer{} \textbf{กปฺโป}ติ ตสฺส พฺราหฺมณสฺส นามํ ฯเปฯ อภิลาโปติ \textbf{อิจฺจา}ยสฺมา \textbf{กปฺโป}.}
\prose{\csromanfolio{137}\textbf{โอเฆ ชาเต มหพฺภเย}ติ กาโมเฆ ภโวเฆ ทิฏฺโฐเฆ อวิชฺโชเฆ ชาเต สญฺชาเต นิพฺพตฺเต อภินิพฺพตฺเต ปาตุภูเต.\csromanspacer{} \textbf{มหพฺภเย}ติ ชาติภเย ชราภเย พฺยาธิภเย มรณภเยติ โอเฆ ชาเต มหพฺภเย.}

\prose{\textbf{ชรามจฺจุปเรตานนฺ}ติ ชราย ผุฏฺฐานํ ปเรตานํ สโมหิตานํ สมนฺนาคตานํ.\csromanspacer{} มจฺจุนา ผุฏฺฐานํ ปเรตานํ สโมหิตานํ สมนฺนาคตานํ, ชาติยา อนุคตานํ ชราย อนุสฏานํ พฺยาธินา อภิภูตานํ มรเณน อพฺภาหตานํ อตาณานํ อเลณานํ อสรณานํ อสรณีภูตานนฺติ ชรามจฺจุปเรตานํ.}

\prose{\textbf{ทีปํ ปพฺรูหิ มาริสา}ติ ทีปํ ตาณํ เลณํ สรณํ คติํ ปรายนํ\footnote{คติปรายนํ (สฺยา) เอวมุปริปิ.} พฺรูหิ อาจิกฺขาหิ เทเสหิ ปญฺญเปหิ ปฏฺฐเปหิ วิวราหิ วิภชาหิ อุตฺตานีกโรหิ ปกาเสหิ.\csromanspacer{} \textbf{มาริสา}ติ ปิยวจนํ ครุวจนํ สคารวสปฺปติสฺสาธิวจนเมตํ มาริสาติ ทีปํ ปพฺรูหิ มาริส.}

\prose{\textbf{ตฺวญฺจ เม ทีปมกฺขาหี}ติ \textbf{ตฺวนฺ}ติ ภควนฺตํ ภณติ.\csromanspacer{} \textbf{ทีปมกฺขาหี}ติ ทีปํ ตาณํ เลณํ สรณํ คติํ ปรายนํ อกฺขาหิ อาจิกฺขาหิ เทเสหิ ปญฺญเปหิ ปฏฺฐเปหิ วิวราหิ วิภชาหิ อุตฺตานีกโรหิ ปกาเสหีติ ตฺวญฺจ เม ทีปมกฺขาหิ.}

\prose{\textbf{ยถายิทํ นาปรํ สิยา}ติ ยถยิทํ ทุกฺขํ อิเธว นิรุชฺเฌยฺย วูปสเมยฺย อตฺถํ คจฺเฉยฺย ปฏิปฺปสฺสมฺเภยฺย, ปุนปฏิสนฺธิกํ ทุกฺขํ น นิพฺพตฺเตยฺย, กามธาตุยา วา รูปธาตุยา วา อรูปธาตุยา วา กามภเว วา รูปภเว \csromanfolio{138}วา อรูปภเว วา สญฺญาภเว วา อสญฺญาภเว วา เนวสญฺญานาสญฺญาภเว วา เอกโวการภเว วา จตุโวการภเว วา ปญฺจโวการภเว วา ปุนคติยา วา อุปปตฺติยา วา ปฏิสนฺธิยา วา ภเว วา สํสาเร วา วฏฺเฏ วา น ชเนยฺย น สญฺชเนยฺย น นิพฺพตฺเตยฺย นาภินิพฺพตฺเตยฺย อิเธว นิรุชฺเฌยฺย วูปสเมยฺย อตฺถํ คจฺเฉยฺย ปฏิปฺปสฺสมฺเภยฺยาติ ยถายิทํ นาปรํ สิยา.\csromanspacer{} เตนาห โส พฺราหฺมโณ\csromandash{}}

\prose{มชฺเฌ สรสฺมิํ ติฏฺฐตํ, (อิจฺจายสฺมา กปฺโป,)}

\prose{\textbf{โอเฆ ชาเต มหพฺภเย.}}

\csromangathameasure{ชรามจฺจุปเรตานํ, ทีปํ ปพฺรูหิ มาริส. \\ ตฺวญฺจ เม ทีปมกฺขาหิ, ยถายิทํ นาปรํ สิยาติ.}
\csromangathagroup{\csromangathastanza{ชรามจฺจุปเรตานํ, ทีปํ ปพฺรูหิ มาริส. \\ ตฺวญฺจ เม ทีปมกฺขาหิ, ยถายิทํ นาปรํ สิยาติ.}}

\csromanhangingbegin
\hangingheaditem{62}{มชฺเฌ สรสฺมิํ ติฏฺฐตํ, (\textbf{กปฺปาติ ภควา},)}
\hangingbody{\textbf{โอเฆ ชาเต มหพฺภเย.}}
\hangingbody{ชรามจฺจุปเรตานํ, ทีปํ ปพฺรูมิ กปฺป เต.}
\csromanhangingend

\prose{\textbf{มชฺเฌ สรสฺมิํ ติฏฺฐตนฺ}ติ สโร วุจฺจติ สํสาโร อาคมนํ คมนํ คมนาคมนํ กาลํ คติ ภวาภโว จุติ จ อุปปตฺติ จ นิพฺพตฺติ จ เภโท จ ชาติ จ ชรา จ มรณญฺจ.\csromanspacer{} สํสารสฺส ปุริมาปิ โกฏิ น ปญฺญายติ, ปจฺฉิมาปิ โกฏิ น ปญฺญายติ.\csromanspacer{} มชฺเฌว สํสาเร สตฺตา ฐิตา ปติฏฺฐิตา อลฺลีนา อุปคตา อชฺโฌสิตา อธิมุตฺตา.}

\prose{กถํ สํสารสฺส ปุริมา โกฏิ น ปญฺญายติ ฯเปฯ.\csromanspacer{} เอวํ สํสารสฺส ปุริมา โกฏิ น ปญฺญายติ.\csromanspacer{} กถํ สํสารสฺส ปจฺฉิมา โกฏิ น ปญฺญายติ ฯเปฯ.\csromanspacer{} เอวํ สํสารสฺส ปจฺฉิมา โกฏิ น ปญฺญายติ.\csromanspacer{} เอวํ สํสารสฺส ปุริมาปิ โกฏิ น ปญฺญายติ, ปจฺฉิมาปิ โกฏิ น ปญฺญายติ.\csromanspacer{} มชฺเฌว สํสาเร สตฺตา ฐิตา ปติฏฺฐิตา อลฺลีนา อุปคตา อชฺโฌสิตา อธิมุตฺตาติ มชฺเฌ สรสฺมิํ ติฏฺฐตํ.\csromanspacer{} \textbf{กปฺปาติ ภควา}ติ \textbf{กปฺปาติ ภควา} ตํ พฺราหฺมณํ นาเมน อาลปติ.\csromanspacer{} \textbf{ภควา}ติ คารวาธิวจนเมตํ ฯเปฯ สจฺฉิกา ปญฺญตฺติ, ยทิทํ \textbf{ภควา}ติ \textbf{กปฺปาติ ภควา}.}

\prose{\textbf{โอเฆ ชาเต มหพฺภเย}ติ กาโมเฆ ภโวเฆ ทิฏฺโฐเฆ อวิชฺโชเฆ ชาเต สญฺชาเต นิพฺพตฺเต อภินิพฺพตฺเต ปาตุภูเต.\csromanspacer{} \textbf{มหพฺภเย}ติ ชาติภเย ชราภเย พฺยาธิภเย มรณภเยติ โอเฆ ชาเต มหพฺภเย.}

\csromanheader{2}{10. กปฺปมาณวปุจฺฉานิทฺเทส}
\prose{\csromanfolio{139}\textbf{ชรามจฺจุปเรตานนฺ}ติ ชราย ผุฏฺฐานํ ปเรตานํ สโมหิตานํ สมนฺนาคตานํ.\csromanspacer{} มจฺจุนา ผุฏฺฐานํ ปเรตานํ สโมหิตานํ สมนฺนาคตานํ ชาติยา อนุคตานํ ชราย อนุสฏานํ พฺยาธินา อภิภูตานํ มรเณน อพฺภาหตานํ อตาณานํ อเลณานํ อสรณานํ อสรณีภูตานนฺติ ชรามจฺจุปเรตานํ.}

\prose{\textbf{ทีปํ ปพฺรูมิ กปฺป เต}ติ ทีปํ ตาณํ เลณํ สรณํ คติํ ปรายนํ พฺรูมิ อาจิกฺขามิ เทเสมิ ปญฺญเปมิ ปฏฺฐเปมิ วิวรามิ วิภชามิ อุตฺตานีกโรมิ ปกาเสมีติ ทีปํ ปพฺรูมิ กปฺป เต.\csromanspacer{} เตนาห ภควา\csromandash{}}

\prose{มชฺเฌ สรสฺมิํ ติฏฺฐตํ, (กปฺปาติ ภควา,)}

\prose{โอเฆ ชาเต มหพฺภเย.}

\csromangathameasure{ชรามจฺจุปเรตานํ, ทีปํ ปพฺรูมิ กปฺป เตติ.}
\csromangathagroup{\csromangathastanza{ชรามจฺจุปเรตานํ, ทีปํ ปพฺรูมิ กปฺป เตติ.}}

\csromanhangingbegin
\hangingheaditem{63}{อกิญฺจนํ อนฺ\textbf{อาทานํ}, เอตํ ทีปํ อนาปรํ.}
\hangingbody{นิพฺพานํ อิติ นํ พฺรูมิ, ชรามจฺจุปริกฺขยํ.}
\csromanhangingend

\prose{\textbf{อกิญฺจนํ อนาทานนฺ}ติ \textbf{กิญฺจนนฺ}ติ ราคกิญฺจนํ โทสกิญฺจนํ โมหกิญฺจนํ มานกิญฺจนํ ทิฏฺฐิกิญฺจนํ กิเลสกิญฺจนํ ทุจฺจริตกิญฺจนํ, กิญฺจนปฺปหานํ กิญฺจนวูปสมํ\footnote{กิญฺจนวูปสโม (สฺยา) เอวมีทิเสสุ ฐาเนสุ.} กิญฺจนปฏินิสฺสคฺคํ\footnote{กิญฺจนปฏินิสฺสคฺโค (สฺยา)} กิญฺจนปฏิปฺปสฺสทฺธิํ อมตํ นิพฺพานนฺติ อกิญฺจนํ.\csromanspacer{} \textbf{อนาทานนฺ}ติ \textbf{อาทานํ} วุจฺจติ ตณฺหา.\csromanspacer{} โย ราโค สาราโค ฯเปฯ อภิชฺฌา โลโภ อกุสลมูลํ, อาทานปฺปหานํ อาทานวูปสมํ อาทานปฏินิสฺสคฺคํ อาทานปฏิปฺปสฺสทฺธิํ อมตํ นิพฺพานนฺติ อกิญฺจนํ อนฺ\textbf{อาทานํ}.}

\prose{\textbf{เอตํ ทีปํ อนาปรนฺ}ติ เอตํ ทีปํ ตาณํ เลณํ สรณํ คติ ปรายนํ.\csromanspacer{} \textbf{อนาปรนฺ}ติ ตมฺหา ปโร อญฺโญ ทีโป นตฺถิ, อถ โข โส เอวํ ทีโป อคฺโค จ เสฏฺโฐ จ วิเสฏฺโฐ จ ปาโมกฺโข จ อุตฺตโม จ ปวโร จาติ เอตํ ทีปํ อนาปรํ.}

\prose{\textbf{นิพฺพานํ อิติ นํ พฺรูมี}ติ \textbf{วานํ} วุจฺจติ ตณฺหา.\csromanspacer{} โย ราโค สาราโค ฯเปฯ อภิชฺฌา โลโภ อกุสลมูลํ, วานปฺปหานํ วานวูปสมํ วานปฏินิสฺสคฺคํ วานปฏิปฺปสฺสทฺธิํ อมตํ นิพฺพานํ.\csromanspacer{} \textbf{อิตี}ติ ปทสนฺธิ ปทสํสคฺโค ปทปาริปูรี อกฺขรสมวาโย พฺยญฺชนสิลิฏฺฐตา ปทานุปุพฺพตาเปตํ อิตีติ. \csromanfolio{140}\textbf{พฺรูมี}ติ พฺรูมิ อาจิกฺขามิ เทเสมิ ปญฺญเปมิ ปฏฺฐเปมิ วิวรามิ วิภชามิ อุตฺตานีกโรมิ ปกาเสมีติ นิพฺพานํ อิติ นํ พฺรูมิ.}

\prose{\textbf{ชรามจฺจุปริกฺขยนฺ}ติ ชรามรณสฺส ปหานํ วูปสมํ ปฏินิสฺสคฺคํ ปฏิปฺปสฺสทฺธิํ อมตํ นิพฺพานนฺติ ชรามจฺจุปริกฺขยํ.\csromanspacer{} เตนาห ภควา\csromandash{}}

\prose{อกิญฺจนํ อนาทานํ, เอตํ ทีปํ อนาปรํ,}

\csromangathameasure{นิพฺพานํ อิติ นํ พฺรูมิ, ชรามจฺจุปริกฺขยนฺติ.}
\csromangathagroup{\csromangathastanza{นิพฺพานํ อิติ นํ พฺรูมิ, ชรามจฺจุปริกฺขยนฺติ.}}

\csromanhangingbegin
\hangingheaditem{64}{\textbf{เอตทญฺญาย เย สตา}, \textbf{ทิฏฺฐธมฺมา}ภินิพฺพุตา.}
\hangingbody{\textbf{น เต มารวสานุคา, น เต มารสฺส ปทฺธคู1}.}
\csromanhangingend

\prose{\textbf{เอตทญฺญาย เย สตา}ติ \textbf{เอตนฺ}ติ อมตํ นิพฺพานํ.\csromanspacer{} โย โส สพฺพสงฺขารสมโถ สพฺพูปธิปฏินิสฺสคฺโค ตณฺหกฺขโย วิราโค นิโรโธ นิพฺพานํ.\csromanspacer{} \textbf{อญฺญายา}ติ อญฺญาย ชานิตฺวา ตุลยิตฺวา ตีรยิตฺวา วิภาวยิตฺวา วิภูตํ กตฺวา, “สพฺเพ สงฺขารา อนิจฺจา”ติ ฯเปฯ. “ยํ กิญฺจิ สมุทยธมฺมํ, สพฺพํ ตํ นิโรธธมฺมนฺ”ติ อญฺญาย ชานิตฺวา ตุลยิตฺวา ตีรยิตฺวา วิภาวยิตฺวา วิภูตํ กตฺวา.\csromanspacer{} \textbf{เย}ติ อรหนฺโต ขีณาสวา.\csromanspacer{} \textbf{สตา}ติ จตูหิ การเณหิ สตา, กาเย กายานุปสฺสานาสติปฏฺฐานํ ภาเวนฺตา\footnote{ภาวิตตฺตา (ก)} สตา ฯเปฯ เต วุจฺจนฺติ สตาติ เอตทญฺญาย เย สตา.}

\prose{\textbf{ทิฏฺฐธมฺมาภินิพฺพุตา}ติ \textbf{ทิฏฺฐธมฺมา}ติ \textbf{ทิฏฺฐธมฺมา} ญาตธมฺมา ตุลิตธมฺมา ตีริตธมฺมา วิภูตธมฺมา วิภาวิตธมฺมา.\csromanspacer{} \textbf{อภินิพฺพุตา}ติ ราคสฺส นิพฺพาปิตตฺตา นิพฺพุตา, โทสสฺส ฯเปฯ สพฺพากุสลาภิสงฺขารานํ สนฺตตฺตา สมิตตฺตา วูปสมิตตฺตา นิชฺฌาตตฺตา นิพฺพุตตฺตา ปฏิปฺปสฺสทฺธตฺตา สนฺตา อุปสนฺตา วูปสนฺตา นิพฺพุตา ปฏิปฺปสฺสทฺธาติ \textbf{ทิฏฺฐธมฺมา}ภินิพฺพุตา.}

\prose{\textbf{น เต มารวสานุคา}ติ \textbf{มาโร}ติ โย โส \textbf{มาโร} กโณฺห อธิปติ อนฺตคู นมุจิ ปมตฺตพนฺธุ.\csromanspacer{} \textbf{น เต มารวสานุคา}ติ น เต มารสฺส วเส วตฺตนฺติ, นาปิ \textbf{มาโร} เตสุ วสํ วตฺเตติ, เต มารญฺจ มารปกฺขญฺจ มารปาสญฺจ มารพฬิสญฺจ\footnote{มารพลิสญฺจ (ก)} มารามิสญฺจ มารวิสยญฺจ มารนิวาสญฺจ มารโคจรญฺจ มารพนฺธนญฺจ อภิภุยฺย อภิภวิตฺวา อชฺโฌตฺถริตฺวา ปริยาทิยิตฺวา มทฺทิตฺวา จรนฺติ วิหรนฺติ อิริยนฺติ วตฺเตนฺติ ปาเลนฺติ ยเปนฺติ ยาเปนฺตีติ น เต มารวสานุคา.}

\csromanmatikaanchor{mtk.32}
\csromantocmarkref{subsection}{11. ชตุกณฺณิมาณวปุจฺฉานิทฺเทส}{mtk.32}
\bookmark[dest={mtk.32},level=2]{11. ชตุกณฺณิมาณวปุจฺฉานิทฺเทส}
\csromanheader{2}{11. ชตุกณฺณิมาณวปุจฺฉานิทฺเทส}
\prose{\csromanfolio{141}\textbf{น เต มารสฺส ปทฺธคู}ติ น เต มารสฺส ปทฺธา ปทฺธจรา\footnote{ปฏฺฐา ปฏฺฐจรา (สฺยา, ก)} ปริจาริกา สิยา, พุทฺธสฺส เต ภควโต ปทฺธา ปทฺธจรา ปริจาริกา สิยาติ น เต มารสฺส ปทฺธคู.\csromanspacer{} เตนาห ภควา\csromandash{}}

\csromangathameasure{เอตทญฺญาย เย สตา, ทิฏฺฐธมฺมาภินิพฺพุตา. \\ น เต มารวสานุคา, น เต มารสฺส ปทฺธคูติ.}
\csromangathagroup{\csromangathastanza{เอตทญฺญาย เย สตา, ทิฏฺฐธมฺมาภินิพฺพุตา. \\ น เต มารวสานุคา, น เต มารสฺส ปทฺธคูติ.}}

\prose{สห คาถาปริโยสานา ฯเปฯ “สตฺถา เม ภนฺเต ภควา, สาวโกหมสฺมี”ติ.}

\vspace{\tipitakaparskip}
\csromancenter{กปฺปมาณวปุจฺฉานิทฺเทโส ทสโม.}
\csromanheader{2}{11. ชตุกณฺณิมาณวปุจฺฉานิทฺเทส}
\csromanhangingbegin
\hangingheaditem{65}{สุตฺวานหํ วีร อกามกามิํ, (อิจฺจายสฺมา ชตุกณฺณิ,)}
\hangingbody{\textbf{โอฆาติคํ ปุฏฺฐุมกามมาคมํ.}}
\hangingbody{\textbf{สนฺติปทํ พฺรูหิ สหชเนตฺต,}}
\hangingbody{ยถาตจฺฉํ ภควา พฺรูหิ เม’ตํ.}
\csromanhangingend

\prose{\textbf{สุตฺวานหํ วีร อกามกามินฺ}ติ สุตฺวา สุณิตฺวา อุคฺคเหตฺวา อุปธาเรตฺวา อุปลกฺขยิตฺวา “อิติปิ โส ภควา อรหํ ฯเปฯ พุทฺโธ ภควา”ติ สุตฺวานหํ.\csromanspacer{} \textbf{วีรนฺ}ติ วีโร, ภควา วีริยวาติ วีโร, ปหูติ วีโร, วิสวีติ วีโร, อลมตฺโตติ วีโร, สูโรติ วีโร, วิกฺกนฺโต อภีรู อจฺฉมฺภี อนุตฺราสี อปลายี ปหีนภยเภรโว วิคตโลมหํโสติ วีโร.}

\csromangathameasure{วิรโต อิธ สพฺพปาปเกหิ, \\ นิรยทุกฺขํ อติจฺจ วีริยวาโส. \\ โส วีริยวา ปธานวา, \\ วีโร ตาทิ ปวุจฺจเต ตถตฺตาติ.}
\csromangathagroup{\csromangathastanza{วิรโต อิธ สพฺพปาปเกหิ, \\ นิรยทุกฺขํ อติจฺจ วีริยวาโส\footnote{วิริยวา (สฺยา) ขุ 1. 360 ปิฏฺเฐปิ.}. \\ โส วีริยวา ปธานวา, \\ วีโร ตาทิ ปวุจฺจเต ตถตฺตาติ.}}

\prose{สุตฺวานหํ วีร.\csromanspacer{}\csromanspacer{}\textbf{อกามกามินฺ}ติ \textbf{กามา}ติ อุทฺทานโต เทฺว \textbf{กามา} วตฺถุ\textbf{กามา} จ กิเลส\textbf{กามา} จ ฯเปฯ.\csromanspacer{} อิเม วุจฺจนฺติ วตฺถุ\textbf{กามา} ฯเปฯ.\csromanspacer{} อิเม วุจฺจนฺติ กิเลส\textbf{กามา}.\csromanspacer{} พุทฺธสฺส ภควโต วตฺถุ\textbf{กามา} ปริญฺญาตา, กิเลส\textbf{กามา} ปหีนา, วตฺถุ\textbf{กามา}นํ ปริญฺญาตตฺตา กิเลส\textbf{กามา}นํ \csromanfolio{142}ปหีนตฺตา \textbf{ภควา} น กาเม กาเมติ, น กาเม ปตฺเถติ, น กาเม ปิเหติ, น กาเม อภิชปฺปติ.\csromanspacer{} เย กาเม กาเมนฺติ, กาเม ปตฺเถนฺติ, กาเม ปิเหนฺติ, กาเม อภิชปฺปนฺติ, เต กามกามิโน ราคราคิโน สญฺญาสญฺญิโน, \textbf{ภควา} น กาเม กาเมติ, น กาเม ปตฺเถติ, น กาเม ปิเหติ, น กาเม อภิชปฺปติ, ตสฺมา พุทฺโธ อกาโม นิกฺกาโม จตฺตกาโม วนฺตกาโม มุตฺตกาโม ปหีนกาโม ปฏินิสฺสฏฺฐกาโม วีตราโค วิคตราโค จตฺตราโค วนฺตราโค มุตฺตราโค ปหีนราโค ปฏินิสฺสฏฺฐราโค นิจฺฉาโต นิพฺพุโต สีติภูโต สุขปฺปฏิสํเวที พฺรหฺมภูเตน อตฺตนา วิหรตีติ สุตฺวานหํ วีร อกามกามิํ.}

\prose{\textbf{อิจฺจายสฺมา ชตุกณฺณี}ติ \textbf{อิจฺจา}ติ ปทสนฺธิ ฯเปฯ ปทานุปุพฺพตาเปตํ \textbf{อิจฺจา}ติ.\csromanspacer{} \textbf{อายสฺมา}ติ ปิยวจนํ สคารวสปฺปติสฺสาธิวจนเมตํ อายสฺมาติ.\csromanspacer{} \textbf{ชตุกณฺณี}ติ ตสฺส พฺราหฺมณสฺส โคตฺตํ สงฺขา สมญฺญา ปญฺญตฺติ โวหาโรติ \textbf{อิจฺจา}ยสฺมา ชตุกณฺณิ.}

\prose{\textbf{โอฆาติคํ ปุฏฺฐุมกามมาคมนฺ}ติ \textbf{โอฆาติคนฺ}ติ โอฆาติคํ โอฆํ อติกฺกนฺตํ สมติกฺกนฺตํ วีติวตฺตนฺติ โอฆาติคํ.\csromanspacer{} \textbf{ปุฏฺฐุ}นฺติ ปุฏฺฐุํ ปุจฺฉิตุํ ยาจิตุํ อชฺเฌสิตุํ ปสาเทตุํ.\csromanspacer{} \textbf{อกามมาคมนฺ}ติ อกามํ ปุฏฺฐุํ นิกฺกามํ จตฺตกามํ วนฺตกามํ มุตฺตกามํ ปหีนกามํ ปฏินิสฺสฏฺฐกามํ วีตราคํ วิคตราคํ จตฺตราคํ วนฺตราคํ มุตฺตราคํ ปหีนราคํ ปฏินิสฺสฏฺฐราคํ อาคมฺหา อาคตมฺหา อุปาคตมฺหา สมฺปตฺตมฺหา ตยา สทฺธิํ สมาคตมฺหาติ โอฆาติคํ ปุฏฺฐุมกามมาคมํ.}

\prose{\textbf{สนฺติปทํ พฺรูหิ สหชเนตฺตา}ติ \textbf{สนฺตี}ติ เอเกน อากาเรน สนฺติปิ สนฺติปทมฺปิ\footnote{สนฺติปทนฺติ (ก)} ตํเยว อมตํ นิพฺพานํ.\csromanspacer{} โย โส สพฺพสงฺขารสมโถ สพฺพูปธิปฏินิสฺสคฺโค ตณฺหกฺขโย วิราโค นิโรโธ นิพฺพานํ.\csromanspacer{} วุตฺตํ เหตํ ภควตา\csromandash{}สนฺตเมตํ ปทํ, ปณีตเมตํ ปทํ, ยทิทํ สพฺพสงฺขารสมโถ สพฺพูปธิปฏินิสฺสคฺโค ตณฺหกฺขโย วิราโค นิโรโธ นิพฺพานนฺติ.\csromanspacer{} อถาปเรนากาเรน เย ธมฺมา สนฺตาธิคมาย สนฺติผุสนาย สนฺติสจฺฉิกิริยาย สํวตฺตนฺติ.\csromanspacer{} เสยฺยถิทํ, จตฺตาโร สติปฏฺฐานา จตฺตาโร สมฺมปฺปธานา จตฺตาโร อิทฺธิปาทา ปญฺจินฺทฺริยานิ ปญฺจ พลานิ สตฺต โพชฺฌงฺคา}

\vspace{\tipitakaparskip}
\csromancenter{ชตุกณฺณิมาณวปุจฺฉานิทฺเทส อริโย อฏฺฐงฺคิโก มคฺโค, อิเม วุจฺจนฺติ สนฺติปทา.\csromanspacer{} สนฺติปทํ ตาณปทํ เลณปทํ สรณปทํ อภยปทํ อจฺจุตปทํ อมตปทํ นิพฺพานปทํ พฺรูหิ อาจิกฺขาหิ เทเสหิ ปญฺญเปหิ ปฏฺฐเปหิ วิวราหิ วิภชาหิ อุตฺตานีกโรหิ ปกาเสหิ.\csromanspacer{} \textbf{สหชเนตฺตา}ติ \textbf{เนตฺตํ} วุจฺจติ สพฺพญฺญุตญฺญาณํ.\csromanspacer{} พุทฺธสฺส ภควโต เนตฺตญฺจ ชินภาโว จ โพธิยา มูเล อปุพฺพํ อจริมํ เอกสฺมิํ ขเณ อุปฺปนฺโน ตสฺมา พุทฺโธ สหชเนตฺโตติ สนฺติปทํ พฺรูหิ สหชเนตฺต.}
\prose{\csromanfolio{143}\textbf{ยถาตจฺฉํ ภควา พฺรูหิ เม’ตนฺ}ติ ยถาตจฺฉํ วุจฺจติ อมตํ นิพฺพานํ ฯเปฯ นิโรโธ นิพฺพานํ.\csromanspacer{} \textbf{ภควา}ติ คารวาธิวจนํ ฯเปฯ สจฺฉิกา ปญฺญตฺติ, ยทิทํ \textbf{ภควา}ติ.\csromanspacer{} \textbf{พฺรูหิ เมตนฺ}ติ พฺรูหิ อาจิกฺขาหิ ฯเปฯ ปกาเสหีติ ยถาตจฺฉํ \textbf{ภควา} พฺรูหิ เม’ตํ.\csromanspacer{} เตนาห โส พฺราหฺมโณ\csromandash{}}

\prose{สุตฺวานหํ วีร อกามกามิํ, (อิจฺจายสฺมา ชตุกณฺณิ)}

\prose{โอฆาติคํ ปุฏฺฐุมกามมาคมํ.}

\prose{สนฺติปทํ พฺรูหิ สหชเนตฺต,}

\prose{\textbf{ยถาตจฺฉํ ภควา พฺรูหิ เม’ตนฺ}ติ.}

\csromanhangingbegin
\hangingheaditem{66}{\textbf{ภควา} หิ กาเม อภิภุยฺย อิริยติ,}
\hangingbody{\textbf{อาทิจฺโจว ปถวิํ เตชี เตชสา.}}
\hangingbody{\textbf{ปริตฺตปญฺญสฺส เม ภูริปญฺโญ,}}
\hangingbody{\textbf{อาจิกฺข ธมฺมํ ยมหํ วิชญฺญํ.}}
\hangingbody{ชาติชราย อิธ วิปฺปหานํ.}
\csromanhangingend

\prose{\textbf{ภควา หิ กาเม อภิภุยฺย อิริยตี}ติ \textbf{ภควา}ติ คารวาธิวจนํ ฯเปฯ สจฺฉิกา ปญฺญตฺติ, ยทิทํ \textbf{ภควา}ติ.\csromanspacer{} \textbf{กามา}ติ อุทฺทานโต เทฺว กามา วตฺถุกามา จ กิเลสกามา จ ฯเปฯ.\csromanspacer{} อิเม วุจฺจนฺติ วตฺถุกามา ฯเปฯ.\csromanspacer{} อิเม วุจฺจนฺติ กิเลสกามา.\csromanspacer{} \textbf{ภควา} วตฺถุกาเม ปริชานิตฺวา กิเลสกาเม ปหาย อภิภุยฺย อภิภวิตฺวา อชฺโฌตฺถริตฺวา ปริยาทิยิตฺวา จรติ วิหรติ อิริยติ วตฺเตติ ปาเลติ ยเปติ ยาเปตีติ \textbf{ภควา} หิ กาเม อภิภุยฺย อิริยติ.}

\prose{\textbf{อาทิจฺโจว ปถวิํ เตชี เตชสา}ติ อาทิจฺโจ วุจฺจติ สูริโย\footnote{สุริโย (สฺยา)}.\csromanspacer{} ปถวี วุจฺจติ ชคตี\footnote{ชรา (สฺยา)}.\csromanspacer{} ยถา สูริโย เตชี เตเชน สมนฺนาคโต \csromanfolio{144}ปถวิํ อภิภุยฺย อภิภวิตฺวา อชฺโฌตฺถริตฺวา ปริยาทิยิตฺวา สนฺตาปยิตฺวา สพฺพํ อากาสคตํ ตมคตํ อภิวิหจฺจ อนฺธการํ วิธมิตฺวา อาโลกํ ทสฺสยิตฺวา อากาเส อนฺตลิกฺเข คคนปเถ\footnote{คมนปเถ (สฺยา) อฏฺฐกถา โอโลเกตพฺพา.} คจฺฉติ.\csromanspacer{} เอวเมว \textbf{ภควา} ญาณเตชี ญาณเตเชน สมนฺนาคโต สพฺพํ อภิสงฺขารสมุทยํ ฯเปฯ กิเลสตมํ อวิชฺชนฺธการํ วิธมิตฺวา ญาณาโลกํ ทสฺเสตฺวา วตฺถุกาเม ปริชานิตฺวา กิเลสกาเม ปหาย อภิภุยฺย อภิภวิตฺวา อชฺโฌตฺถริตฺวา ปริยาทิยิตฺวา มทฺทิตฺวา จรติ วิหรติ อิริยติ วตฺเตติ ปาเลติ ยเปติ ยาเปตีติ อาทิจฺโจว ปถวิํ เตชี เตชสา.}

\prose{\textbf{ปริตฺตปญฺญสฺส เม ภูริปญฺโญ}ติ อหมสฺมิ ปริตฺตปญฺโญ โอมกปญฺโญ ลามกปญฺโญ ฉตุกฺกปญฺโญ.\csromanspacer{} ตฺวมฺปิ มหาปญฺโญ ปุถุปญฺโญ หาสปญฺโญ ชวนปญฺโญ ติกฺขปญฺโญ นิพฺเพธิกปญฺโญ.\csromanspacer{} ภูริ วุจฺจติ ปถวี.\csromanspacer{} ภควา ตาย ปถวิสมาย ปญฺญาย วิปุลาย วิตฺถตาย สมนฺนาคโตติ ปริตฺตปญฺญสฺส เม ภูริปญฺโญ.}

\prose{\textbf{อาจิกฺข ธมฺมํ ยมหํ วิชญฺญนฺ}ติ \textbf{ธมฺมนฺ}ติ อาทิกลฺยาณํ มชฺเฌกลฺยาณํ ปริโยสานกลฺยาณํ สาตฺถํ สพฺยญฺชนํ เกวลปริปุณฺณํ ปริสุทฺธํ พฺรหฺมจริยํ จตฺตาโร สติปฏฺฐาเน ฯเปฯ นิพฺพานญฺจ นิพฺพานคามินิญฺจ ปฏิปทํ อาจิกฺขาหิ เทเสหิ ปญฺญเปหิ ปฏฺฐเปหิ วิวราหิ วิภชาหิ อุตฺตานีกโรหิ ปกาเสหิ.\csromanspacer{} \textbf{ยมหํ วิชญฺญนฺ}ติ ยมหํ ชาเนยฺยํ อาชาเนยฺยํ วิชาเนยฺยํ ปฏิชาเนยฺยํ ปฏิวิชฺเฌยฺยํ อธิคจฺเฉยฺยํ ผสฺเสยฺยํ สจฺฉิกเรยฺยนฺติ อาจิกฺข ธมฺมํ ยมหํ วิชญฺญํ.}

\prose{\textbf{ชาติชราย อิธ วิปฺปหานนฺ}ติ อิเธว ชาติชราย มรณสฺส ปหานํ วูปสมํ ปฏินิสฺสคฺคํ ปฏิปฺปสฺสทฺธิํ อมตํ นิพฺพานนฺติ ชาติชราย อิธ วิปฺปหานํ.\csromanspacer{} เตนาห โส พฺราหฺมณฺ\csromandash{}}

\csromangathameasure{ภควา หิ กาเม อภิภุยฺย อิริยติ, \\ อาทิจฺโจว ปถวิํ เตชี เตชสา.}
\csromangathagroup{\csromangathastanza{ภควา หิ กาเม อภิภุยฺย อิริยติ, \\ อาทิจฺโจว ปถวิํ เตชี เตชสา.}}

\prose{\textbf{ปริตฺตปญฺญสฺส เม ภูริปญฺโญ},}

\prose{อาจิกฺข ธมฺมํ ยมหํ วิชญฺญํ.}

\prose{\textbf{ชาติชราย อิธ วิปฺปหานนฺ}ติ.}

\csromanheader{2}{11. ชตุกณฺณิมาณวปุจฺฉานิทฺเทส}
\csromanhangingbegin
\hangingheaditem{67}{\csromanfolio{145}กาเมสุ วินย เคธํ (\textbf{ชตุกณฺณี}ติ \textbf{ภควา},)}
\hangingbody{\textbf{เนกฺขมฺมํ ทฏฺฐุ เขมโต.}}
\hangingbody{\textbf{อุคฺคหิตํ นิรตฺตํ วา}, มา เต วิชฺชิตฺถ กิญฺจนํ.}
\csromanhangingend

\prose{\textbf{กาเมสุ วินย เคธนฺ}ติ \textbf{กามา}ติ อุทฺทานโต เทฺว \textbf{กามา} วตฺถุ\textbf{กามา} จ กิเลส\textbf{กามา} จ ฯเปฯ.\csromanspacer{} อิเม วุจฺจนฺติ วตฺถุ\textbf{กามา} ฯเปฯ.\csromanspacer{} อิเม วุจฺจนฺติ กิเลส\textbf{กามา}.\csromanspacer{} \textbf{เคธนฺ}ติ เคโธ วุจฺจติ ตณฺหา.\csromanspacer{} โย ราโค สาราโค ฯเปฯ อภิชฺฌา โลโภ อกุสลมูลํ.\csromanspacer{} \textbf{กาเมสุ วินย เคธนฺ}ติ กาเมสุ เคธํ วินย ปฏิวินย ปชห วิโนเทหิ พฺยนฺตีกโรหิ อนภาวํ คเมหีติ กาเมสุ วินย เคธํ.\csromanspacer{} \textbf{ชตุกณฺณี}ติ \textbf{ภควา} ตํ พฺราหฺมณํ โคตฺเตน อาลปติ.\csromanspacer{} \textbf{ภควา}ติ คารวาธิวจนเมตํ ฯเปฯ สจฺฉิกา ปญฺญตฺติ, ยทิทํ \textbf{ภควา}ติ \textbf{ชตุกณฺณี}ติ \textbf{ภควา}.}

\prose{\textbf{เนกฺขมฺมํ ทฏฺฐุ เขมโต}ติ \textbf{เนกฺขมฺมนฺ}ติ สมฺมาปฏิปทํ อนุโลมปฏิปทํ อปจฺจนีกปฏิปทํ อนฺวตฺถปฏิปทํ ธมฺมานุธมฺมปฏิปทํ สีเลสุ ปริปูรการิตํ อินฺทฺริเยสุ คุตฺตทฺวารตํ โภชเน มตฺตญฺญุตํ ชาคริยานุโยคํ สติสมฺปชญฺญํ จตฺตาโร สติปฏฺฐาเน จตฺตาโร สมฺมปฺปธาเน จตฺตาโร อิทฺธิปาเท ปญฺจินฺทฺริยานิ ปญฺจ พลานิ สตฺต โพชฺฌงฺเค อริยํ อฏฺฐงฺคิกํ มคฺคํ นิพฺพานญฺจ นิพฺพานคามินิญฺจ ปฏิปทํ เขมโต ตาณโต เลณโต สรณโต สรณีภูตโต อภยโต อจฺจุตโต อมตโต นิพฺพานโต ทฏฺฐุํ ปสฺสิตฺวา ตุลยิตฺวา ตีรยิตฺวา วิภาวยิตฺวา วิภูตํ กตฺวาติ เนกฺขมฺมํ ทฏฺฐุ เขมโต.}

\prose{\textbf{อุคฺคหิตํ นิรตฺตํ วา}ติ \textbf{อุคฺคหิตนฺ}ติ ตณฺหาวเสน ทิฏฺฐิวเสน คหิตํ ปรามฏฺฐํ อภินิวิฏฺฐํ อชฺโฌสิตํ อธิมุตฺตํ.\csromanspacer{} \textbf{นิรตฺตํ วา}ติ นิรตฺตํ วา มุญฺจิตพฺพํ วิชหิตพฺพํ วิโนทิตพฺพํ พฺยนฺตีกาตพฺพํ อนภาวํ คเมตพฺพนฺติ อุคฺคหิตํ นิรตฺตํ วา.}

\prose{\textbf{มา เต วิชฺชิตฺถ กิญฺจนนฺ}ติ ราคกิญฺจนํ โทสกิญฺจนํ โมหกิญฺจนํ มานกิญฺจนํ ทิฏฺฐิกิญฺจนํ กิเลสกิญฺจน ทุจฺจริตกิญฺจนํ.\csromanspacer{} อิทํ กิญฺจนํ\footnote{อิเม กิญฺจนา (ก)} ตุยฺหํ มา วิชฺชิตฺถ มา ปวิชฺชิตฺถ มา สํวิชฺชิตฺถ ปชห วิโนเทหิ พฺยนฺตีกโรหิ อนภาวํ คเมหีติ มา เต วิชฺชิตฺถ กิญฺจนํ.\csromanspacer{} เตนาห \textbf{ภควา}\csromandash{}}

\prose{\csromanfolio{146}กาเมสุ วินย เคธํ, (ชตุกณฺณีติ ภควา,)}

\prose{เนกฺขมฺมํ ทฏฺฐุ เขมโต.}

\csromangathameasure{อุคฺคหิตํ นิรตฺตํ วา, มา เต วิชฺชิตฺถ กิญฺจนนฺติ.}
\csromangathagroup{\csromangathastanza{อุคฺคหิตํ นิรตฺตํ วา, มา เต วิชฺชิตฺถ กิญฺจนนฺติ.}}

\csromanhangingbegin
\hangingheaditem{68}{ยํ ปุพฺเพ ตํ วิโสเสหิ, ปจฺฉา เต มาหุ กิญฺจนํ.}
\hangingbody{มชฺเฌ เจ โน คเหสฺสสิ, อุปสนฺโต จริสฺสสิ.}
\csromanhangingend

\prose{\textbf{ยํ ปุพฺเพ ตํ วิโสเสหี}ติ อตีเต สงฺขาเร อารพฺภ เย กิเลสา อุปฺปชฺเชยฺยุํ, เต กิเลเส โสเสหิ วิโสเสหิ สุกฺขาเปหิ วิสุกฺขาเปหิ อพีชํ กโรหิ ปชห วิโนเทหิ พฺยนฺตีกโรหิ อนภาวํ คเมหีติ เอวมฺปิ ยํ ปุพฺเพ ตํ วิโสเสหิ.\csromanspacer{}\csromanspacer{}อถ วา เย อตีตา กมฺมาภิสงฺขารา อวิปกฺกวิปากา, เต กมฺมาภิสงฺขาเร โสเสหิ วิโสเสหิ สุกฺขาเปหิ วิสุกฺขาเปหิ อพีชํ\footnote{อวีชํ (สฺยา)} กโรหิ ปชห วิโนเทหิ พฺยนฺตีกโรหิ อนภาวํ คเมหีติ เอวมฺปิ ยํ ปุพฺเพ ตํ วิโสเสหิ.}

\prose{\textbf{ปจฺฉา เต มาหุ กิญฺจนนฺ}ติ ปจฺฉา วุจฺจติ อนาคเต สงฺขาเร อารพฺภ ราคกิญฺจนํ โทสกิญฺจนํ โมหกิญฺจนํ มานกิญฺจนํ ทิฏฺฐิกิญฺจนํ กิเลสกิญฺจนํ ทุจฺจริตกิญฺจนํ, อิทํ กิญฺจนํ ตุยฺหํ มา อหุ มา อโหสิ มา ชเนสิ\footnote{มา ชเนหิ (สฺยา) ตถาวเสเสสุ ทฺวีสุ ปเทสุปิ.} มา สญฺชเนสิ มาภินิพฺพตฺเตสิ ปชห วิโนเทหิ พฺยนฺตีกโรหิ อนภาวํ คเมหีติ ปจฺฉา เต มาหุ กิญฺจนํ.}

\prose{\textbf{มชฺเฌ เจ โน คเหสฺสสี}ติ มชฺเฌ วุจฺจติ ปจฺจุปฺปนฺนํ รูปํ เวทนา สญฺญา สงฺขารา วิญฺญาณํ, ปจฺจุปฺปนฺเน สงฺขาเร ตณฺหาวเสน ทิฏฺฐิวเสน น คเหสฺสสิ น คณฺหิสฺสสิ น ปรามสิสฺสสิ น นนฺทิสฺสสิ นาภินนฺทิสฺสสิ น อชฺโฌสิสฺสสิ, อภินนฺทนํ อภิวทนํ อชฺโฌสานํ คาหํ ปรามาสํ อภินิเวสํ ปชหิสฺสสิ วิโนเทสฺสสิ พฺยนฺตีกริสฺสสิ อนภาวํ คเมสฺสสีติ มชฺเฌ เจ โน คเหสฺสสิ.}

\prose{\textbf{อุปสนฺโต จริสฺสสี}ติ ราคสฺส อุปสมิตตฺตา อุปสนฺโต จริสฺสสิ, โทสสฺส ฯเปฯ สพฺพากุสลาภิสงฺขารานํ สนฺตตฺตา สมิตตฺตา อุปสมิตตฺตา วูปสมิตตฺตา นิชฺฌาตตฺตา นิพฺพุตตฺตา วิคตตฺตา ปฏิปฺปสฺสทฺธตฺตา สนฺโต อุปสนฺโต วูปสนฺโต นิพฺพุโต ปฏิปฺปสฺสทฺโธ จริสฺสสิ วิหริสฺสสิ}

\vspace{\tipitakaparskip}
\csromancenter{ชตุกณฺณิมาณวปุจฺฉานิทฺเทส อิริยิสฺสสิ วตฺติสฺสสิ ปาเลสฺสสิ ยเปสฺสสิ ยาเปสฺสสีติ อุปสนฺโต จริสฺสสิ.\csromanspacer{} เตนาห ภควา\csromandash{}}
\vspace{0.5\baselineskip}
\csromangathameasure{ยํ ปุพฺเพ ตํ วิโสเสหิ, ปจฺฉา เต มาหุ กิญฺจนํ. \\ มชฺเฌ เจ โน คเหสฺสสิ, อุปสนฺโต จริสฺสสีติ.}
\csromangathagroup{\csromangathastanza{\csromanfolio{147}ยํ ปุพฺเพ ตํ วิโสเสหิ, ปจฺฉา เต มาหุ กิญฺจนํ. \\ มชฺเฌ เจ โน คเหสฺสสิ, อุปสนฺโต จริสฺสสีติ.}}

\csromanhangingbegin
\hangingheaditem{69}{สพฺพโส นามรูปสฺมิํ, วีตเคธสฺส พฺราหฺมณ.}
\hangingbody{อาสวาสฺส น วิชฺชนฺติ, เยหิ มจฺจุวสํ วเช.}
\csromanhangingend

\prose{\textbf{สพฺพโส นามรูปสฺมิํ, วีตเคธสฺส พฺราหฺมณา}ติ \textbf{สพฺพโส}ติ สพฺเพน สพฺพํ สพฺพถา สพฺพํ อเสสํ นิสฺเสสํ ปริยาทิยนวจนเมตํ \textbf{สพฺพโส}ติ.\csromanspacer{} \textbf{นามนฺ}ติ จตฺตาโร อรูปิโน ขนฺธา.\csromanspacer{} \textbf{รูปนฺ}ติ จตฺตาโร จ มหาภูตา จตุนฺนญฺจ มหาภูตานํ อุปาทาย รูปํ.\csromanspacer{} เคโธ วุจฺจติ ตณฺหา.\csromanspacer{} โย ราโค สาราโค ฯเปฯ อภิชฺฌา โลโภ อกุสลมูลํ.\csromanspacer{} \textbf{สพฺพโส นามรูปสฺมิํ, วีตเคธสฺส พฺราหฺมณา}ติ \textbf{สพฺพโส} นามรูปสฺมิํ วีตเคธสฺส วิคตเคธสฺส จตฺตเคธสฺส วนฺตเคธสฺส มุตฺตเคธสฺส ปหีนเคธสฺส ปฏินิสฺสฏฺฐเคธสฺส, วีตราคสฺส วิคตราคสฺส จตฺตราคสฺส วนฺตราคสฺส มุตฺตราคสฺส ปหีนราคสฺส ปฏินิสฺสฏฺฐราคสฺสาติ \textbf{สพฺพโส} นามรูปสฺมิํ, วีตเคธสฺส พฺราหฺมณ.}

\prose{\textbf{อาสวาสฺส น วิชฺชนฺตี}ติ \textbf{อาสวา}ติ จตฺตาโร \textbf{อาสวา} กามาสโว ภวาสโว ทิฏฺฐาสโว อวิชฺชาสโว.\csromanspacer{} \textbf{อสฺสา}ติ อรหโต ขีณาสวสฺส.\csromanspacer{} \textbf{น วิชฺชนฺตี}ติ อิเม \textbf{อาสวา} ตสฺส นตฺถิ น สนฺติ น สํวิชฺชนฺติ นุปลพฺภนฺติ, ปหีนา สมุจฺฉินฺนา วูปสนฺตา ปฏิปฺปสฺสทฺธา อภพฺพุปฺปตฺติกา ญาณคฺคินา ทฑฺฒาติ \textbf{อาสวา}สฺส น วิชฺชนฺติ.}

\prose{\textbf{เยหิ มจฺจุวสํ วเช}ติ เยหิ อาสเวหิ มจฺจุโน วา วสํ คจฺเฉยฺย, มรณสฺส วา วสํ คจฺเฉยฺย, มารปกฺขสฺส วา วสํ คจฺเฉยฺย, เต \textbf{อาสวา} ตสฺส นตฺถิ น สนฺติ น สํวิชฺชนฺติ นุปลพฺภนฺติ, ปหีนา สมุจฺฉินฺนา วูปสนฺตา ปฏิปฺปสฺสทฺธา อภพฺพุปฺปตฺติกา ญาณคฺคินา ทฑฺฒาติ เยหิ มจฺจุวสํ วเช.\csromanspacer{} เตนาห ภควา\csromandash{}}

\prose{สพฺพโส นามรูปสฺมิํ, วีตเคธสฺส พฺราหฺมณ.}

\csromangathameasure{อาสวาสฺส น วิชฺชนฺติ, เยหิ มจฺจุวสํ วเชติ.}
\csromangathagroup{\csromangathastanza{อาสวาสฺส น วิชฺชนฺติ, เยหิ มจฺจุวสํ วเชติ.}}

\prose{\csromanfolio{148}สห คาถาปริโยสานา ฯเปฯ “สตฺถา เม ภนฺเต ภควา, สาวโกหมสฺมี”ติ.}

\vspace{\tipitakaparskip}
\csromancenter{ชาตุกณฺณิมาณวปุจฺฉานิทฺเทโส เอกาทสโม.}
\csromanmatikaanchor{mtk.33}
\csromantocmarkref{subsection}{12. ภทฺราวุธมาณวปุจฺฉานิทฺเทส}{mtk.33}
\bookmark[dest={mtk.33},level=2]{12. ภทฺราวุธมาณวปุจฺฉานิทฺเทส}
\csromanheader{2}{12. ภทฺราวุธมาณวปุจฺฉานิทฺเทส}
\csromanhangingbegin
\hangingheaditem{70}{โอกญฺชหํ ตณฺหจฺฉิทํ อเนชํ, (อิจฺจายสฺมา ภทฺราวุโธ,)}
\hangingbody{\textbf{นนฺทิญฺชหํ โอฆติณฺณํ วิมุตฺตํ.}}
\hangingbody{\textbf{กปฺปญฺชหํ อภิยาเจ สุเมธํ,}}
\hangingbody{\textbf{สุตฺวาน นาคสฺส อปนมิสฺสนฺติ1} อิโต.}
\csromanhangingend

\prose{\textbf{โอกญฺชหํ ตณฺหจฺฉิทํ อเนชนฺ}ติ \textbf{โอกญฺชหนฺ}ติ รูปธาตุยา โย ฉนฺโท โย ราโค ยา นนฺที ยา \textbf{ตณฺหา} เย อุปยุปาทานา\footnote{อุปายุปาทานา (ก)} เจตโส อธิฏฺฐานาภินิเวสานุสยา, เต พุทฺธสฺส ภควโต ปหีนา อุจฺฉินฺนมูลา ตาลาวตฺถุกตา อนภาวํกตา อายติํ อนุปฺปาทธมฺมา.\csromanspacer{} ตสฺมา พุทฺโธ โอกญฺชโห.\csromanspacer{} เวทนาธาตุยา.\csromanspacer{} สญฺญาธาตุยา.\csromanspacer{} สงฺขารธาตุยา.\csromanspacer{} วิญฺญาณธาตุยา โย ฉนฺโท โย ราโค ยา นนฺที ยา \textbf{ตณฺหา} เย อุปยุปาทานา เจตโส อธิฏฺฐานาภินิเวสานุสยา, เต พุทฺธสฺส ภควโต ปหีนา อุจฺฉินฺนมูลา ตาลาวตฺถุกตา อนภาวํกตา อายติํ อนุปฺปาทธมฺมา, ตสฺมา พุทฺโธ โอกญฺชโห.}

\prose{\textbf{ตณฺหจฺฉิทนฺ}ติ \textbf{ตณฺหา}ติ รูป\textbf{ตณฺหา} ฯเปฯ ธมฺม\textbf{ตณฺหา}, สา \textbf{ตณฺหา} พุทฺธสฺส ภควโต ฉินฺนา อุจฺฉินฺนา สมุจฺฉินฺนา วูปสนฺตา ปฏิปฺปสฺสทฺธา อภพฺพุปฺปตฺติกา ญาณคฺคินา ทฑฺฒา, ตสฺมา พุทฺโธ ตณฺหจฺฉิโท.\csromanspacer{} \textbf{อเนโช}ติ \textbf{เอชา} วุจฺจติ \textbf{ตณฺหา}, โย ราโค สาราโค ฯเปฯ อภิชฺฌา โลโภ อกุสลมูลํ.\csromanspacer{} สา \textbf{เอชา} \textbf{ตณฺหา} พุทฺธสฺส ภควโต ปหีนา อุจฺฉินฺนมูลา ตาลาวตฺถุกตา อนภาวํกตา อายติํ อนุปฺปาทธมฺมา, ตสฺมา พุทฺโธ อเนโช.\csromanspacer{} เอชาย ปหีนตฺตา อเนโช.\csromanspacer{} ภควา ลาเภปิ น อิญฺชติ, อลาเภปิ น อิญฺชติ, ยเสปิ น อิญฺชติ, อยเสปิ น อิญฺชติ, ปสํสายปิ น อิญฺชติ, นินฺทายปิ น อิญฺชติ, สุเขปิ น อิญฺชติ, ทุกฺเขปิ น อิญฺชติ น จลติ น เวธติ น ปเวธติ น สมฺปเวธตีติ ตสฺมา พุทฺโธ}

\vspace{\tipitakaparskip}
\csromancenter{ภทฺราวุธมาณวปุจฺฉานิทฺเทส อเนโชติ โอกญฺชหํ ตณฺหจฺฉิทํ อเนชํ.\csromanspacer{} \textbf{อิจฺจายสฺมา ภทฺราวุโธ}ติ \textbf{อิจฺจา}ติ ปทสนฺธิ ฯเปฯ.\csromanspacer{} \textbf{อายสฺมา}ติ ปิยวจนํ ฯเปฯ.\csromanspacer{} \textbf{ภทฺราวุโธ}ติ ตสฺส พฺราหฺมณสฺส นามํ ฯเปฯ อภิลาโปติ \textbf{อิจฺจา}ยสฺมา \textbf{ภทฺราวุโธ}.}
\prose{\csromanfolio{149}\textbf{นนฺทิญฺชหํ โอฆติณฺณํ วิมุตฺตนฺ}ติ นนฺที วุจฺจติ ตณฺหา.\csromanspacer{} โย ราโค สาราโค ฯเปฯ อภิชฺฌา โลโภ อกุสลมูลํ, สา นนฺที สา ตณฺหา พุทฺธสฺส ภควโต ปหีนา อุจฺฉินฺนมูลา ตาลาวตฺถุกตา อนภาวํกตา อายติํ อนุปฺปาทธมฺมา, ตสฺมา พุทฺโธ นนฺทิญฺชิโห.\csromanspacer{} \textbf{โอฆติณฺณนฺ}ติ ภควา กาโมฆํ ติณฺโณ, ภโวฆํ ติณฺโณ, ทิฏฺโฐฆํ ติณฺโณ, อวิชฺโชฆํ ติณฺโณ, สพฺพสํสารปถํ ติณฺโณ อุตฺติณฺโณ นิตฺถิณฺโณ อติกฺกนฺโต สมติกฺกนฺโต วีติวตฺโต, โส วุตฺถวาโส จิณฺณจรโณ ฯเปฯ ชาติมรณสํสาโร, นตฺถิ ตสฺส ปุนพฺภโวติ นนฺทิญฺชหํ โอฆติณฺณํ.\csromanspacer{} \textbf{วิมุตฺตนฺ}ติ ภควโต ราคา จิตฺตํ มุตฺตํ วิมุตฺตํ สุวิมุตฺตํ.\csromanspacer{} โทสา จิตฺตํ.\csromanspacer{} โมหา จิตฺตํ ฯเปฯ สพฺพากุสลาภิสงฺขาเรหิ จิตฺตํ มุตฺตํ วิมุตฺตํ สุวิมุตฺตนฺติ นนฺทิญฺชหํ โอฆติณฺณํ วิมุตฺตํ.}

\prose{\textbf{กปฺปญฺชหํ อภิยาเจ สุเมธนฺ}ติ \textbf{กปฺปา}ติ เทฺว \textbf{กปฺปา} ตณฺหากปฺโป จ ทิฏฺฐิกปฺโป จ ฯเปฯ อยํ ตณฺหากปฺโป ฯเปฯ อยํ ทิฏฺฐิกปฺโป.\csromanspacer{} พุทฺธสฺส ภควโต ตณฺหากปฺโป ปหีโน, ทิฏฺฐิกปฺโป ปฏินิสฺสฏฺโฐ, ตณฺหากปฺปสฺส ปหีนตฺตา, ทิฏฺฐิกปฺปสฺส ปฏินิสฺสฏฺฐตฺตา, ตสฺมา พุทฺโธ กปฺปญฺชโห.\csromanspacer{} \textbf{อภิยาเจ}ติ ยาจามิ อภิยาจามิ อชฺเฌสามิ สาทิยามิ ปตฺถยามิ ปิหยามิ ชปฺปามิ อภิชปฺปามิ.\csromanspacer{} สุเมธา วุจฺจติ ปญฺญา, ยา ปญฺญา ปชานนา ฯเปฯ อโมโห ธมฺมวิจโย สมฺมาทิฏฺฐิ.\csromanspacer{} ภควา อิมาย เมธาย ปญฺญาย อุเปโต สมุเปโต อุปาคโต สมุปาคโต อุปปนฺโน สมุปปนฺโน สมนฺนาคโต, ตสฺมา พุทฺโธ สุเมโธติ กปฺปญฺชหํ อภิยาเจ สุเมธํ.}

\prose{\textbf{สุตฺวาน นาคสฺส อปนมิสฺสนฺติ อิโต}ติ \textbf{นาคสฺสา}ติ นาโค, ภควา อาคุํ น กโรตีติ นาโค, น คจฺฉตีติ นาโค, น อาคจฺฉตีติ นาโค ฯเปฯ.\csromanspacer{} เอวํ ภควา น คจฺฉตีติ นาโค.\csromanspacer{} \textbf{สุตฺวาน นาคสฺส อปนมิสฺสนฺติ อิโต}ติ ตุยฺหํ วจนํ พฺยปฺปถํ เทสนํ อนุสาสนํ อนุสิฏฺฐํ สุตฺวา สุณิตฺวา อุคฺคเหตฺวา อุปธารยิตฺวา อุปลกฺขยิตฺวา อิโต อปนมิสฺสนฺติ วชิสฺสนฺติ ปกฺกมิสฺสนฺติ ทิสาวิทิสํ คมิสฺสนฺตีติ สุตฺวาน นาคสฺส อปนมิสฺสนฺติ อิโต.\csromanspacer{} เตนาห โส พฺราหฺมโณ\csromandash{}}

\prose{\csromanfolio{150}โอกญฺชหํ ตณฺหจฺฉิทํ อเนชํ, (อิจฺจายสฺมา ภทฺราวุโธ,)}

\prose{นนฺทิญฺชหํ โอฆติณฺณํ วิมุตฺตํ.}

\csromangathameasure{กปฺปญฺชหํ อภิยาเจ สุเมธํ, \\ สุตฺวาน นาคสฺส อปนมิสฺสนฺติ อิโตติ.}
\csromangathagroup{\csromangathastanza{กปฺปญฺชหํ อภิยาเจ สุเมธํ, \\ สุตฺวาน นาคสฺส อปนมิสฺสนฺติ อิโตติ.}}

\proseitem{71}{\textbf{นานาชนา ชนปเทหิ สงฺคตา},}

\prose{\textbf{ตว วีร วากฺยํ อภิกงฺขมานา.}}

\prose{\textbf{เตสํ ตุวํ สาธุ วิยากโรหิ,}}

\prose{\textbf{นานาชนา ชนปเทหิ สงฺคตา}ติ \textbf{นานาชนา}ติ ขตฺติยา จ พฺราหฺมณา จ เวสฺสา จ สุทฺทา จ คหฏฺฐา จ ปพฺพชิตา จ เทวา จ มนุสฺสา จ.\csromanspacer{} \textbf{ชนปเทหิ สงฺคตา}ติ องฺคา จ มคธา จ กลิงฺคา จ กาสิยา จ โกสลา จ วชฺชิยา จ มลฺลา จ เจติยมฺหา จ\footnote{เจติยมฺหา จ สาครมฺหา จ (สฺยา)} วํสา จ กุรุมฺหา จ ปญฺจาลา จ มจฺฉา จ สุรเสนา จ อสฺสกา จ อวนฺติยา จ โยนา\footnote{โยนกา (ก) ขุ 7. 120 ปิฏฺเฐปิ.} จ กมฺโพชา จ.\csromanspacer{} \textbf{สงฺคตา}ติ สงฺคตา สมาคตา สโมหิตา สนฺนิปติตาติ \textbf{นานาชนา} ชนปเทหิ สงฺคตา.}

\csromangathameasure{ตว วีร วากฺยํ อภิกงฺขมานาติ วีราติ วีโร, ภควา วีริยวาติ วีโร, ปหูติ วีโร, วิสวีติ วีโร, อลมตฺโตติ วีโร, วิคตโลมหํโสติปิ วีโร. \\ วิรโต อิธ สพฺพปาปเกหิ, \\ นิรยทุกฺขํ อติจฺจ วีริยวาโส. \\ โส วีริยวา ปธานวา, \\ วีโร ตาทิ ปวุจฺจเต ตถตฺตาติ.}
\csromangathagroup{\csromangathastanza{ตว วีร วากฺยํ อภิกงฺขมานาติ วีราติ วีโร, ภควา วีริยวาติ วีโร, ปหูติ วีโร, วิสวีติ วีโร, อลมตฺโตติ วีโร, วิคตโลมหํโสติปิ วีโร.}\csromangathastanza{วิรโต อิธ สพฺพปาปเกหิ, \\ นิรยทุกฺขํ อติจฺจ วีริยวาโส. \\ โส วีริยวา ปธานวา, \\ วีโร ตาทิ ปวุจฺจเต ตถตฺตาติ.}}

\prose{\textbf{ตว วีร วากฺยํ อภิกงฺขมานา}ติ ตุยฺหํ วจนํ พฺยปฺปถํ เทสนํ อนุสาสนํ อนุสิฏฺฐํ.\csromanspacer{} \textbf{อภิกงฺขมานา}ติ อภิกงฺขมานา อิจฺฉมานา สาทิยมานา ปตฺถยมานา ปิหยมานา อภิชปฺปมานาติ ตว วีร วากฺยํ อภิกงฺขมานา.}

\prose{\textbf{เตสํ ตุวํ สาธุ วิยากโรหี}ติ \textbf{เตสนฺ}ติ เตสํ ขตฺติยานํ พฺราหฺมณานํ เวสฺสานํ สุทฺทานํ คหฏฺฐานํ ปพฺพชิตานํ เทวานํ มนุสฺสานํ.\csromanspacer{} \textbf{ตุวนฺ}ติ ภควนฺตํ ภณติ.\csromanspacer{} \textbf{สาธุ วิยากโรหี}ติ สาธุ อาจิกฺขาหิ}

\vspace{\tipitakaparskip}
\csromancenter{ภทฺราวุธมาณวปุจฺฉานิทฺเทส เทเสหิ ปญฺญเปหิ ปฏฺฐเปหิ วิวราหิ วิภชาหิ อุตฺตานีกโรหิ ปกาเสหีติ เตสํ ตุวํ สาธุ วิยากโรหิ.}
\prose{\csromanfolio{151}\textbf{ตถา หิ เต วิทิโต เอส ธมฺโม}ติ ตถา หิ เต วิทิโต ตุลิโต ตีริโต วิภูโต วิภาวิโต เอส ธมฺโมติ ตถา หิ เต วิทิโต เอส ธมฺโม.\csromanspacer{} เตนาห โส พฺราหฺมโณ\csromandash{}}

\csromangathameasure{นานาชนา ชนปเทหิ สงฺคตา, \\ ตว วีร วากฺยํ อภิกงฺขมานา. \\ เตสํ ตุวํ สาธุ วิยากโรหิ,}
\csromangathagroup{\csromangathastanza{นานาชนา ชนปเทหิ สงฺคตา, \\ ตว วีร วากฺยํ อภิกงฺขมานา. \\ เตสํ ตุวํ สาธุ วิยากโรหิ,}}

\prose{\textbf{ตถา หิ เต วิทิโต เอส ธมฺโม}ติ.}

\csromanhangingbegin
\hangingheaditem{72}{อาทานตณฺหํ วินเยถ สพฺพํ, (\textbf{ภทฺราวุธาติ ภควา},)}
\hangingbody{\textbf{อุทฺธํ อโธ ติริยญฺจาปิ มชฺเฌ.}}
\hangingbody{\textbf{ยํ ยญฺหิ โลกสฺมิมุปาทิยนฺติ,}}
\hangingbody{เตเนว มาโร อเนฺวติ ชนฺตุํ.}
\csromanhangingend

\prose{\textbf{อาทานตณฺหํ วินเยถ สพฺพนฺ}ติ อาทานตณฺหํ วุจฺจติ รูปตณฺหา ฯเปฯ อาทานตณฺหาติ กิํการณา วุจฺจติ อาทานตณฺหา, ตาย ตณฺหาย รูปํ อาทิยนฺติ อุปาทิยนฺติ คณฺหนฺติ ปรามสนฺติ อภินิวิสนฺติ.\csromanspacer{} เวทนํ.\csromanspacer{} สญฺญํ.\csromanspacer{} สงฺขาเร.\csromanspacer{} วิญฺญาณํ.\csromanspacer{} คติํ.\csromanspacer{} อุปปตฺติํ.\csromanspacer{} ปฏิสนฺธิํ.\csromanspacer{} ภวํ.\csromanspacer{} สํสารํ.\csromanspacer{} วฏฺฏํ อาทิยนฺติ อุปาทิยนฺติ คณฺหนฺติ ปรามสนฺติ อภินิวิสนฺติ, ตํการณา วุจฺจติ อาทานตณฺหา.\csromanspacer{} \textbf{อาทานตณฺหํ วินเยถ สพฺพนฺ}ติ สพฺพํ อาทานตณฺหํ วินเยยฺย ปฏิวินเยยฺย ปชเหยฺย วิโนเทยฺย พฺยนฺตีกเรยฺย อนภาวํ คเมยฺยาติ อาทานตณฺหํ วินเยถ สพฺพํ, \textbf{ภทฺราวุธาติ ภควา}ติ \textbf{ภทฺราวุธาติ ภควา} ตํ พฺราหฺมณํ นาเมน อาลปติ.\csromanspacer{} \textbf{ภควา}ติ คารวาธิวจนเมตํ ฯเปฯ สจฺฉิกา ปญฺญตฺติ, ยทิทํ \textbf{ภควา}ติ \textbf{ภทฺราวุธาติ ภควา}.}

\prose{\textbf{อุทฺธํ อโธ ติริยญฺจาปิ มชฺเฌ}ติ \textbf{อุทฺธนฺ}ติ อนาคตํ.\csromanspacer{} \textbf{อโธ}ติ อตีตํ.\csromanspacer{} \textbf{ติริยญฺจาปิ มชฺเฌ}ติ ปจฺจุปฺปนฺนํ.\csromanspacer{} \textbf{อุทฺธนฺ}ติ เทวโลโก.\csromanspacer{} \textbf{อโธ}ติ นิรยโลโก.\csromanspacer{} \textbf{ติริยญฺจาปิ มชฺเฌ}ติ มนุสฺสโลโก.\csromanspacer{} อถ วา \textbf{อุทฺธนฺ}ติ กุสลา ธมฺมา.\csromanspacer{} \textbf{อโธ}ติ อกุสลา ธมฺมา.\csromanspacer{} \textbf{ติริยญฺจาปิ มชฺเฌ}ติ อพฺยากตา ธมฺมา.\csromanspacer{} \textbf{อุทฺธนฺ}ติ อรูปธาตุ.\csromanspacer{} \textbf{อโธ}ติ กามธาตุ.\csromanspacer{} \textbf{ติริยญฺจาปิ มชฺเฌ}ติ รูปธาตุ.\csromanspacer{} \textbf{อุทฺธนฺ}ติ สุขา เวทนา.\csromanspacer{} \textbf{อโธ}ติ ทุกฺขา เวทนา.\csromanspacer{} \textbf{ติริยญฺจาปิ มชฺเฌ}ติ อทุกฺขมสุขา เวทนา.\csromanspacer{} \textbf{อุทฺธนฺ}ติ อุทฺธํ ปาทตลา. \csromanfolio{152}\textbf{อโธ}ติ อโธ เกสมตฺถกา.\csromanspacer{} \textbf{ติริยญฺจาปิ มชฺเฌ}ติ เวมชฺเฌติ อุทฺธํ อโธ ติริยญฺจาปิ มชฺเฌ.}

\prose{\textbf{ยํ ยญฺหิ โลกสฺมิมุปาทิยนฺตี}ติ ยํ ยํ รูปคตํ เวทนาคตํ สญฺญาคตํ สงฺขารคตํ วิญฺญาณคตํ อาทิยนฺติ อุปาทิยนฺติ คณฺหนฺติ ปรามสนฺติ อภินิวิสนฺติ.\csromanspacer{} \textbf{โลกสฺมินฺ}ติ อปายโลเก ฯเปฯ อายตนโลเกติ ยํ ยญฺหิ โลกสฺมิมุปาทิยนฺติ.}

\prose{\textbf{เตเนว มาโร อเนฺวติ ชนฺตุนฺ}ติ เตเนว กมฺมาภิสงฺขารวเสน ปฏิสนฺธิโก ขนฺธมาโร ธาตุมาโร อายตนมาโร คติมาโร อุปปตฺติมาโร ปฏิสนฺธิมาโร ภวมาโร สํสารมาโร วฏฺฏมาโร อเนฺวติ อนุคจฺฉติ อนฺวายิโก โหติ.\csromanspacer{} \textbf{ชนฺตุนฺ}ติ สตฺตํ ชนํ นรํ มาณวํ\footnote{มาณวํ (สฺยา)} โปสํ ปุคฺคลํ ชีวํ ชาคุํ ชนฺตุํ อินฺทคุํ มนุชนฺติ เตเนว มาโร อเนฺวติ ชนฺตุํ.\csromanspacer{} เตนาห ภควา\csromandash{}}

\prose{อาทานตณฺหํ วินเยถ สพฺพํ, (ภทฺราวุธาติ ภควา,)}

\prose{อุทฺธํ อโธ ติริยญฺจาปิ มชฺเฌ.}

\csromangathameasure{ยํ ยญฺหิ โลกสฺมิมุปาทิยนฺติ, \\ \textbf{เตเนว มาโร อเนฺวติ ชนฺตุนฺ}ติ.}
\csromangathagroup{\csromangathastanza{ยํ ยญฺหิ โลกสฺมิมุปาทิยนฺติ, \\ \textbf{เตเนว มาโร อเนฺวติ ชนฺตุนฺ}ติ.}}

\csromanhangingbegin
\hangingheaditem{73}{ตสฺมา ปชานํ น อุปาทิเยถ,}
\hangingbody{\textbf{ภิกฺขุ สโต กิญฺจนํ สพฺพโลเก.}}
\hangingbody{\textbf{อาทานสตฺเต อิติ เปกฺขมาโน,}}
\hangingbody{ปชํ อิมํ มจฺจุเธยฺเย วิสตฺตํ.}
\csromanhangingend

\prose{\textbf{ตสฺมา ปชานํ น อุปาทิเยถา}ติ \textbf{ตสฺมา}ติ \textbf{ตสฺมา} ตํการณา ตํเหตุ ตปฺปจฺจยา ตํนิทานา เอตํ อาทีนวํ สมฺปสฺสมาโน อาทานตณฺหายาติ \textbf{ตสฺมา}.\csromanspacer{} \textbf{ปชานนฺ}ติ ชานนฺโต ปชานนฺโต อาชานนฺโต วิชานนฺโต ปฏิวิชานนฺโต ปฏิวิชฺฌนฺโต “สพฺเพ สงฺขารา อนิจฺจา”ติ ฯเปฯ. “ยํ กิญฺจิ สมุทยธมฺมํ, สพฺพํ ตํ นิโรธธมฺมนฺ”ติ ชานนฺโต ปชานนฺโต อาชานนฺโต วิชานนฺโต ปฏิวิชานนฺโต ปฏิวิชฺฌนฺโต.\csromanspacer{} \textbf{น อุปาทิเยถา}ติ รูปํ นาทิเยยฺย น อุปาทิเยยฺย น คเณฺหยฺย น ปรามเสยฺย นาภินิวิเสยฺย.\csromanspacer{} เวทนํ.\csromanspacer{} สญฺญํ.\csromanspacer{} สงฺขาเร.\csromanspacer{} วิญฺญาณํ, คติํ, อุปปตฺติํ, ปฏิสนฺธิํ, ภวํ.\csromanspacer{} สํสารํ.\csromanspacer{} วฏฺฏํ นาทิเยยฺย น อุปาทิเยยฺย น คเณฺหยฺย น ปรามเสยฺย นาภินิวิเสยฺยาติ \textbf{ตสฺมา} ปชานํ น อุปาทิเยถ.}

\csromanheader{2}{12. ภทฺราวุธมาณวปุจฺฉานิทฺเทส}
\prose{\csromanfolio{153}\textbf{ภิกฺขุ สโต กิญฺจนํ สพฺพโลเก}ติ \textbf{ภิกฺขู}ติ ปุถุชฺชนกลฺยาณโก วา ภิกฺขุ, เสกฺโข วา ภิกฺขุ.\csromanspacer{} \textbf{สโต}ติ จตูหิ การเณหิ สโต, กาเย กายานุปสฺสนาสติปฏฺฐานํ ภาเวนฺโต สโต ฯเปฯ โส วุจฺจติ สโตติ ภิกฺขุ สโต.\csromanspacer{} \textbf{กิญฺจนนฺ}ติ กิญฺจิ รูปคตํ เวทนาคตํ สญฺญาคตํ สงฺขารคตํ วิญฺญาณคตํ.\csromanspacer{} \textbf{สพฺพโลเก}ติ สพฺพ-อปายโลเก สพฺพมนุสฺสโลเก สพฺพเทวโลเก สพฺพขนฺธโลเก สพฺพธาตุโลเก สพฺพ-อายตนโลเกติ ภิกฺขุ สโต กิญฺจนํ สพฺพโลเก.}

\prose{\textbf{อาทานสตฺเต อิติ เปกฺขมาโน}ติ อาทานสตฺตา วุจฺจนฺติ เย รูปํ อาทิยนฺติ อุปาทิยนฺติ คณฺหนฺติ ปรามสนฺติ อภินิวิสนฺติ.\csromanspacer{} เวทนํ.\csromanspacer{} สญฺญํ.\csromanspacer{} สงฺขาเร.\csromanspacer{} วิญฺญาณํ.\csromanspacer{} คติํ.\csromanspacer{} อุปปตฺติํ.\csromanspacer{} ปฏิสนฺธิํ.\csromanspacer{} ภวํ.\csromanspacer{} สํสารํ.\csromanspacer{} วฏฺฏํ อาทิยนฺติ อุปาทิยนฺติ คณฺหนฺติ ปรามสนฺติ อภินิวิสนฺติ.\csromanspacer{} \textbf{อิตี}ติ ปทสนฺธิ ฯเปฯ ปทานุปุพฺพตาเปตํ อิตีติ.\csromanspacer{} \textbf{เปกฺขมาโน}ติ เปกฺขมาโน ทกฺขมาโน ทิสฺสมาโน ปสฺสมาโน โอโลกยมาโน นิชฺฌายมาโน อุปปริกฺขมาโนติ อาทานสตฺเต อิติ เปกฺขมาโน.}

\prose{\textbf{ปชํ อิมํ มจฺจุเธยฺเย วิสตฺตนฺ}ติ \textbf{ปชา}ติ สตฺตาธิวจนํ.\csromanspacer{} มจฺจุเธยฺยา วุจฺจนฺติ กิเลสา จ ขนฺธา จ อภิสงฺขารา จ, \textbf{ปชา} มจฺจุเธยฺเย มารเธยฺเย มรณเธยฺเย สตฺตา วิสตฺตา อาสตฺตา ลคฺคา ลคฺคิตา ปลิพุทฺธา, ยถา ภิตฺติขิเล วา นาคทนฺเต วา ภณฺฑํ สตฺตํ วิสตฺตํ อาสตฺตํ ลคฺคํ ลคฺคิตํ ปลิพุทฺธํ, เอวเมว \textbf{ปชา} มจฺจุเธยฺเย มารเธยฺเย มรณเธยฺเย สตฺตา วิสตฺตา อาสตฺตา ลคฺคา ลคฺคิตา ปลิพุทฺธาติ ปชํ อิมํ มจฺจุเธยฺเย วิสตฺตํ.\csromanspacer{} เตนาห ภควา\csromandash{}}

\prose{ตสฺมา \textbf{ปชา}นํ น อุปาทิเยถ,}

\prose{\textbf{ภิกฺขุ สโต กิญฺจนํ สพฺพโลเก}.}

\prose{\textbf{อาทานสตฺเต อิติ เปกฺขมาโน},}

\prose{\textbf{ปชํ อิมํ มจฺจุเธยฺเย วิสตฺตนฺ}ติ.}

\prose{สห คาถาปริโยสานา ฯเปฯ “สตฺถา เม ภนฺเต ภควา, สาวโกหมสฺมี”ติ.}

\vspace{\tipitakaparskip}
\csromancenter{ภทฺราวุธมาณวปุจฺฉานิทฺเทโส ทฺวาทสโม.}
\csromanmatikaanchor{mtk.34}
\csromantocmarkref{subsection}{13. อุทยมาณวปุจฺฉานิทฺเทส}{mtk.34}
\bookmark[dest={mtk.34},level=2]{13. อุทยมาณวปุจฺฉานิทฺเทส}
\csromanheader{2}{13. อุทยมาณวปุจฺฉานิทฺเทส}
\csromanhangingbegin
\hangingheaditem{74}{\csromanfolio{154}ฌายิํ วิรชมาสีนํ, (อิจฺจายสฺมา อุทโย,)}
\hangingbody{\textbf{กตกิจฺจํ อนาสวํ.}}
\hangingbody{\textbf{ปารคุํ สพฺพธมฺมานํ, อตฺถิ ปเญฺหน อาคมํ.}}
\hangingbody{\textbf{อญฺญาวิโมกฺขํ ปพฺรูหิ1}, อวิชฺชาย ปเภทนํ.}
\csromanhangingend

\prose{\textbf{ฌายิํ วิรชมาสีนนฺ}ติ \textbf{ฌายินฺ}ติ ฌายี, ภควา ปฐเมนปิ ฌาเนน ฌายี, ทุติเยนปิ ฌาเนน ฌายี, ตติเยนปิ ฌาเนน ฌายี, จตุตฺเถนปิ ฌาเนน ฌายี, สวิตกฺกสวิจาเรนปิ ฌาเนน ฌายี, อวิตกฺกวิจารมตฺเตนปิ ฌาเนน ฌายี, อวิตกฺก-อวิจาเรนปิ ฌาเนน ฌายี, สปฺปีติเกนปิ ฌาเนน ฌายี, นิปฺปีติเกนปิ ฌาเนน ฌายี, สาตสหคเตนปิ ฌาเนน ฌายี, อุเปกฺขาสหคเตนปิ ฌาเนน ฌายี, สุญฺญเตนปิ ฌาเนน ฌายี, อนิมิตฺเตนปิ ฌาเนน ฌายี, อปฺปณิหิเตนปิ ฌาเนน ฌายี, โลกิเยนปิ ฌาเนน ฌายี, โลกุตฺตเรนปิ ฌาเนน ฌายี, ฌานรโต เอกตฺตมนุยุตฺโต สทตฺถครุโกติ ฌายิํ.\csromanspacer{} \textbf{วิรชนฺ}ติ ราโค รโช, โทโส รโช, โมโห รโช, โกโธ รโช, อุปนาโห รโช ฯเปฯ สพฺพากุสลาภิสงฺขารา รชา, เต รชา พุทฺธสฺส ภควโต ปหีนา อุจฺฉินฺนมูลา ตาลาวตฺถุกตา อนภาวํกตา อายติํ อนุปฺปาทธมฺมา, ตสฺมา พุทฺโธ อรโช วิรโช นิรโช รชาปคโต รชวิปฺปหีโน รชวิปฺปยุตฺโต สพฺพรชวีติวตฺโต.}

\csromangathameasure{ราโค รโช น จ ปน เรณุ วุจฺจติ, \\ ราคสฺเสตํ อธิวจนํ รโชติ. \\ เอตํ รชํ วิปฺปชหิตฺวา จกฺขุมา, \\ ตสฺมา ชิโน วิคตรโชติ วุจฺจติ. \\ โทโส รโช น จ ปน เรณุ วุจฺจติ, \\ โทสสฺเสตํ อธิวจนํ รโชติ. \\ เอตํ รชํ วิปฺปชหิตฺวา จกฺขุมา, \\ ตสฺมา ชิโน วิคตรโชติ วุจฺจติ.}
\csromangathagroup{\csromangathastanza{ราโค รโช น จ ปน เรณุ วุจฺจติ, \\ ราคสฺเสตํ อธิวจนํ รโชติ. \\ เอตํ รชํ วิปฺปชหิตฺวา\footnote{ปฏิวิโนทิตฺวา (ก) ขุ 7. 406 ปิฏฺเฐปิ.} จกฺขุมา, \\ ตสฺมา ชิโน วิคตรโชติ วุจฺจติ.}\csromangathastanza{โทโส รโช น จ ปน เรณุ วุจฺจติ, \\ โทสสฺเสตํ อธิวจนํ รโชติ. \\ เอตํ รชํ วิปฺปชหิตฺวา จกฺขุมา, \\ ตสฺมา ชิโน วิคตรโชติ วุจฺจติ.}}

\csromanheader{2}{13. อุทยมาณวปุจฺฉานิทฺเทส}
\csromangathameasure{โมโห รโช น จ ปน เรณุ วุจฺจติ, \\ โมหสฺเสตํ อธิวจนํ รโชติ. \\ เอตํ รชํ วิปฺปชหิตฺวา จกฺขุมา, \\ ตสฺมา ชิโน วิคตรโชติ วุจฺจตีติ.}
\csromangathagroup{\csromangathastanza{\csromanfolio{155}โมโห รโช น จ ปน เรณุ วุจฺจติ, \\ โมหสฺเสตํ อธิวจนํ รโชติ. \\ เอตํ รชํ วิปฺปชหิตฺวา จกฺขุมา, \\ ตสฺมา ชิโน วิคตรโชติ วุจฺจตีติ.}}

\prose{วิรชํ.\csromanspacer{}\csromanspacer{}\textbf{อาสีนนฺ}ติ นิสินฺโน ภควา ปาสาณเก เจติเยติ อาสีโน.}

\csromangathameasure{นคสฺส ปสฺเส อาสีนํ, มุนิํ ทุกฺขสฺส ปารคุํ.}
\csromangathagroup{\csromangathastanza{นคสฺส\footnote{นครสฺส (ก)} ปสฺเส อาสีนํ, มุนิํ ทุกฺขสฺส ปารคุํ.}}

\vspace{\tipitakaparskip}
\csromancenter{สาวกา ปยิรุปาสนฺติ, เตวิชฺชา มจฺจุหายิโนติ\csromandash{}}
\prose{เอวมฺปิ ภควา อาสีโน.\csromanspacer{}\csromanspacer{}อถ วา ภควา สพฺโพสฺสุกฺกปฏิปฺปสฺสทฺธตฺตา อาสีโน วุตฺถวาโส จิณฺณจรโณ ฯเปฯ ชาติมรณสํสาโร, นตฺถิ ตสฺส ปุนพฺภโวติ เอวมฺปิ ภควา อาสีโนติ ฌายิํ วิรชมาสีนํ.}

\prose{\textbf{อิจฺจายสฺมา อุทโย}ติ \textbf{อิจฺจา}ติ ปทสนฺธิ ฯเปฯ.\csromanspacer{} \textbf{อายสฺมา}ติ ปิยวจนํ ฯเปฯ.\csromanspacer{} \textbf{อุทโย}ติ ตสฺส พฺราหฺมณสฺส นามํ ฯเปฯ อภิลาโปติ \textbf{อิจฺจา}ยสฺมา \textbf{อุทโย}.}

\prose{\textbf{กตกิจฺจํ อนาสวนฺ}ติ พุทฺธสฺส ภควโต กฺ\textbf{อิจฺจา}กิจฺจํ กรณียากรณียํ ปหีนํ อุจฺฉินฺนมูลํ ตาลาวตฺถุกตํ อนภาวํกตํ อายติํ อนุปฺปาทธมฺมํ, ตสฺมา พุทฺโธ กตกิจฺโจ.}

\csromangathameasure{ยสฺส จ วิสตา นตฺถิ, ฉินฺนโสตสฺส ภิกฺขุโน. \\ กิจฺจากิจฺจปฺปหีนสฺส, ปริฬาโห น วิชฺชตีติ.}
\csromangathagroup{\csromangathastanza{ยสฺส จ วิสตา\footnote{ยสฺส ปริปตา (สฺยา) ปสฺส ขุ 1. 390 ปิฏฺเฐ.} นตฺถิ, ฉินฺนโสตสฺส ภิกฺขุโน. \\ กิจฺจากิจฺจปฺปหีนสฺส, ปริฬาโห น วิชฺชตีติ.}}

\prose{\textbf{กตกิจฺจํ อนาสวนฺ}ติ \textbf{อาสวา}ติ จตฺตาโร \textbf{อาสวา} กามาสโว ภวาสโว ทิฏฺฐาสโว อวิชฺชาสโว, เต \textbf{อาสวา} พุทฺธสฺส ภควโต ปหีนา อุจฺฉินฺนมูลา ตาลาวตฺถุกตา อนภาวํกตา อายติํ อนุปฺปาทธมฺมา, ตสฺมา พุทฺโธ อนาสโวติ กตกิจฺจํ อนาสวํ.}

\prose{\textbf{ปารคุํ สพฺพธมฺมานนฺ}ติ ภควา สพฺพธมฺมานํ อภิญฺญาปารคู ปริญฺญาปารคู ปหานปารคู ภาวนาปารคู สจฺฉิกิริยาปารคู สมาปตฺติปารคู.\csromanspacer{} อภิญฺญาปารคู สพฺพธมฺมานํ, ปริญฺญาปารคู สพฺพทุกฺขานํ, ปหานปารคู สพฺพกิเลสานํ, ภาวนาปารคู จตุนฺนํ มคฺคานํ, สจฺฉิกิริยาปารคู นิโรธสฺส, สมาปตฺติปารคู สพฺพสมาปตฺตีนํ.\csromanspacer{} โส วสิปฺปตฺโต ปารมิปฺปตฺโต อริยสฺมิํ สีลสฺมิํ,}

\prose{\csromanfolio{156}วสิปฺปตฺโต ปารมิปฺปตฺโต อริยสฺมิํ สมาธิสฺมิํ, วสิปฺปตฺโต ปารมิปฺปตฺโต อริยาย ปญฺญาย, วสิปฺปตฺโต ปารมิปฺปตฺโต อริยาย วิมุตฺติยา.\csromanspacer{} โส ปารคโต ปารปฺปตฺโต อนฺตคโต อนฺตปฺปตฺโต โกฏิคโต โกฏิปฺปตฺโต ปริยนฺตคโต ปริยนฺตปฺปตฺโต โวสานคโต โวสานปฺปตฺโต ตาณคโต ตาณปฺปตฺโต เลณคโต เลณปฺปตฺโต สรณคโต สรณปฺปตฺโต อภยคโต อภยปฺปตฺโต อจฺจุตคโต อจฺจุตปฺปตฺโต อมตคโต อมตปฺปตฺโต นิพฺพานคโต นิพฺพานปฺปตฺโต, โส วุตฺถวาโส จิณฺณจรโณ ฯเปฯ ชาติมรณสํสาโร, นตฺถิ ตสฺส ปุนพฺภโวติ ปารคุํ สพฺพธมฺมานํ.}

\prose{\textbf{อตฺถิ ปเญฺหน อาคมนฺ}ติ ปเญฺหน อตฺถิโก อาคโตมฺหิ, ปญฺหํ ปุจฺฉิตุกาโม อาคโตมฺหิ, ปญฺหํ โสตุกาโม อาคโตมฺหีติ เอวมฺปิ อตฺถิ ปเญฺหน อาคมํ.\csromanspacer{} อถ วา ปญฺหตฺถิกานํ ปญฺหํ ปุจฺฉิตุกามานํ ปญฺหํ โสตุกามานํ อาคมนํ อภิกฺกมนํ อุปสงฺกมนํ ปยิรุปาสนํ อตฺถีติ เอวมฺปิ อตฺถิ ปเญฺหน อาคมํ.\csromanspacer{} อถ วา ปญฺหาคโม ตุยฺหํ อตฺถิ, ตฺวมฺปิ ปหุ, วิสวี อลมตฺโต มยา ปุจฺฉิตํ กเถตุํ วิสชฺเชตุํ วหสฺเสตํ ภารนฺติ เอวมฺปิ อตฺถิ ปเญฺหน อาคมํ.}

\prose{\textbf{อญฺญาวิโมกฺขํ ปพฺรูหี}ติ อญฺญาวิโมกฺโข วุจฺจติ อรหตฺตวิโมกฺโข, อรหตฺตวิโมกฺขํ ปพฺรูหิ อาจิกฺขาหิ เทเสหิ ปญฺญเปหิ ปฏฺฐเปหิ วิวราหิ วิภชาหิ อุตฺตานีกโรหิ ปกาเสหีติ อญฺญาวิโมกฺขํ ปพฺรูหิ.}

\prose{\textbf{อวิชฺชาย ปเภทนนฺ}ติ อวิชฺชาย เภทนํ ปเภทนํ ปหานํ วูปสมํ ปฏินิสฺสคฺคํ ปฏิปฺปสฺสทฺธํ อมตํ นิพฺพานนฺติ อวิชฺชาย ปเภทนํ.\csromanspacer{} เตนาห โส พฺราหฺมโณ\csromandash{}}

\prose{ฌายิํ วิรชมาสีนํ, (อิจฺจายสฺมา อุทโย,)}

\prose{กตกิจฺจํ อนาสวํ.}

\csromangathameasure{ปารคุํ สพฺพธมฺมานํ, อตฺถิ ปเญฺหน อาคมํ. \\ อญฺญาวิโมกฺขํ ปพฺรูหิ, อวิชฺชาย ปเภทนนฺติ.}
\csromangathagroup{\csromangathastanza{ปารคุํ สพฺพธมฺมานํ, อตฺถิ ปเญฺหน อาคมํ. \\ อญฺญาวิโมกฺขํ ปพฺรูหิ, อวิชฺชาย ปเภทนนฺติ.}}

\csromanheader{2}{13. อุทยมาณวปุจฺฉานิทฺเทส}
\csromanhangingbegin
\hangingheaditem{75}{\csromanfolio{157}ปหานํ กามจฺฉนฺทานํ, (\textbf{อุทยาติ ภควา},)}
\hangingbody{\textbf{โทมนสฺสาน จูภยํ.}}
\hangingbody{\textbf{ถินสฺส1} จ ปนูทนํ, กุกฺกุจฺจานํ นิวารณํ.}
\csromanhangingend

\prose{\textbf{ปหานํ กามจฺฉนฺทานนฺ}ติ \textbf{ฉนฺโท}ติ โย กาเมสุ กามจฺ\textbf{ฉนฺโท} กามราโค กามนนฺที กามตณฺหา กามสิเนโห กามปิปาสา กามปริฬาโห กามมุจฺฉา กามชฺโฌสานํ กาโมโฆ กามโยโค กามุปาทานํ กามจฺฉนฺทนีวรณํ.\csromanspacer{} \textbf{ปหานํ กามจฺฉนฺทานนฺ}ติ กามจฺฉนฺทานํ ปหานํ วูปสมํ ปฏินิสฺสคฺคํ ปฏิปฺปสฺสทฺธิํ อมตํ นิพฺพานนฺติ ปหานํ กามจฺฉนฺทานํ.\csromanspacer{} \textbf{อุทยาติ ภควา}ติ \textbf{อุทยาติ ภควา} ตํ พฺราหฺมณํ นาเมน อาลปติ.\csromanspacer{} \textbf{ภควา}ติ คารวาธิวจนเมตํ ฯเปฯ สจฺฉิกา ปญฺญตฺติ, ยทิทํ \textbf{ภควา}ติ \textbf{อุทยาติ ภควา}.}

\prose{\textbf{โทมนสฺสาน จูภยนฺ}ติ \textbf{โทมนสฺสา}ติ ยํ เจตสิกํ อสาตํ, เจตสิกํ ทุกฺขํ เจโตสมฺผสฺสชํ อสาตํ ทุกฺขํ เวทยิตํ, เจโตสมฺผสฺสชา อสาตา ทุกฺขา เวทนา.\csromanspacer{} \textbf{โทมนสฺสาน จูภยนฺ}ติ กามจฺฉนฺทสฺส จ โทมนสฺสสฺส จ อุภินฺนํ ปหานํ วูปสมํ ปฏินิสฺสคฺคํ ปฏิปฺปสฺสทฺธิํ อมตํ นิพฺพานนฺติ \textbf{โทมนสฺสา}น จูภยํ.}

\prose{\textbf{ถินสฺส จ ปนูทนนฺ}ติ \textbf{ถินนฺ}ติ ยา จิตฺตสฺส อกลฺยตา อกมฺมญฺญตา โอลียนา สลฺลียนา ลีนา ลียนา ลียิตตฺตํ ถินํ ถิยนา\footnote{ถีนํ ถียนา (สฺยา)} ถิยิตตฺตํ จิตฺตสฺส.\csromanspacer{} \textbf{ปนูทนนฺ}ติ ถินสฺส จ ปนูทนํ ปหานํ วูปสมํ ปฏินิสฺสคฺคํ ปฏิปฺปสฺสทฺธิํ อมตํ นิพฺพานนฺติ ถินสฺส จ ปนูทนํ.}

\prose{\textbf{กุกฺกุจฺจานํ นิวารณนฺ}ติ \textbf{กุกฺกุจฺจนฺ}ติ หตฺถกุกฺกุจฺจมฺปิ กุกฺกุจฺจํ, ปาทกุกฺกุจฺจมฺปิ กุกฺกุจฺจํ, หตฺถปาทกุกฺกุจฺจมฺปิ กุกฺกุจฺจํ, อกปฺปิเย กปฺปิยสญฺญิตา, กปฺปิเย อกปฺปิยสญฺญิตา, อวชฺเช วชฺชสญฺญิตา, วชฺเช อวชฺชสญฺญิตา, ยํ เอวรูปํ กุกฺกุจฺจํ กุกฺกุจฺจายนา กุกฺกุจฺจายิตตฺตํ เจตโส วิปฺปฏิสาโร มโนวิเลโข, อิทํ วุจฺจติ กุกฺกุจฺจํ.\csromanspacer{} อปิ จ ทฺวีหิ การเณหิ อุปฺปชฺชติ กุกฺกุจฺจํ เจตโส วิปฺปฏิสาโร มโนวิเลโข กตตฺตา จ อกตตฺตา จ.\csromanspacer{} กถํ กตตฺตา จ อกตตฺตา จ อุปฺปชฺชติ กุกฺกุจฺจํ เจตโส วิปฺปฏิสาโร มโนวิเลโข, “กตํ เม กายทุจฺจริตํ, อกตํ เม กายสุจริตนฺ”ติ อุปฺปชฺชติ กุกฺกุจฺจํ เจตโส วิปฺปฏิสาโร \csromanfolio{158}มโนวิเลโข, “กตํ เม วจีทุจฺจริตํ, อกตํ เม วจีสุจริตนฺ”ติ ฯเปฯ “กตํ เม มโนทุจฺจริตํ, อกตํ เม มโนสุจริตนฺ”ติ ฯเปฯ “กโต เม ปาณาติปาโต, อกตา เม ปาณาติปาตา เวรมณี”ติ ฯเปฯ “กตํ เม อทินฺนาทานํ, อกตา เม อทินฺนาทานา เวรมณี”ติ ฯเปฯ “กโต เม กาเมสุมิจฺฉาจาโร, อกตา เม กาเมสุมิจฺฉาจารา เวรมณี”ติ ฯเปฯ “กโต เม มุสาวาโท, อกตา เม มุสาวาทา เวรมณี”ติ ฯเปฯ “กตา เม ปิสุณา วาจา\footnote{ปิสุณวาจา (ก)}, อกตา เม ปิสุณาย วาจาย เวรมณี”ติ ฯเปฯ “กตา เม ผรุสา วาจา, อกตา เม ผรุสาย วาจาย เวรมณี”ติ ฯเปฯ “กโต เม สมฺผปฺปลาโป, อกตา เม สมฺผปฺปลาปา เวรมณี”ติ ฯเปฯ “กตา เม อภิชฺฌา, อกตา เม อนภิชฺฌา”ติ ฯเปฯ “กโต เม พฺยาปาโท, อกโต เม อพฺยาปาโท”ติ ฯเปฯ “กตา เม มิจฺฉาทิฏฺฐิ, อกตา เม สมฺมาทิฏฺถิ”ติ อุปฺปชฺชติ กุกฺกุจฺจํ เจตโส วิปฺปฏิสาโร มโนวิเลโข, เอวํ กตตฺตา จ อกตตฺตา จ อุปฺปชฺชติ กุกฺกุจฺจํ เจตโส วิปฺปฏิสาโร มโนวิเลโข.}

\prose{อถ วา “สีเลสุ’มฺหิ อปริปูรการี”ติ อุปฺปชฺชติ กุกฺกุจฺจํ เจตโส วิปฺปฏิสาโร มโนวิเลโข, “อินฺทฺริเยสุ’มฺหิ อคุตฺตทฺวาโร”ติ. “โภชเน อมตฺตญฺญุมฺหี”ติ. “ชาคริยํ อนนุยุตฺโตมฺหี”ติ. “น สติสมฺปชญฺเญน สมนฺนาคโตมฺหี”ติ. “อภาวิตา เม จตฺตาโร สติปฏฺฐานา”ติ. “จตฺตาโร สมฺมปฺปธานา”ติ. “จตฺตาโร อิทฺธิปาทา”ติ. “ปญฺจินฺทฺริยานี”ติ. “ปญฺจ พลานี”ติ. “สตฺต โพชฺฌงฺคา”ติ. “อริโย อฏฺฐงฺคิโก มคฺโค”ติ. “ทุกฺขํ เม อปริญฺญาตํ.\csromanspacer{} สมุทโย เม อปฺปหีโน.\csromanspacer{} มคฺโค เม อภาวิโต.\csromanspacer{} นิโรโธ เม อสจฺฉิกโต”ติ อุปฺปชฺชติ กุกฺกุจฺจํ เจตโส วิปฺปฏิสาโร มโนวิเลโข.}

\prose{\textbf{กุกฺกุจฺจานํ นิวารณนฺ}ติ กุกฺกุจฺจานํ อาวรณํ นีวรณํ ปหานํ อุปสมํ วูปสมํ ปฏินิสฺสคฺคํ ปฏิปฺปสฺสทฺธิํ อมตํ นิพฺพานนฺติ กุกฺกุจฺจานํ นิวารณํ.\csromanspacer{} เตนาห ภควา\csromandash{}}

\vspace{\tipitakaparskip}
\csromancenter{ปหานํ กามจฺฉนฺทานํ, (อุทยาติ ภควา,)}
\prose{โทมนสฺสาน จูภยํ.}

\csromangathameasure{ถินสฺส จ ปนูทนํ, กุกฺกุจฺจานํ นิวารณนฺติ.}
\csromangathagroup{\csromangathastanza{ถินสฺส จ ปนูทนํ, กุกฺกุจฺจานํ นิวารณนฺติ.}}

\csromanheader{2}{13. อุทยมาณวปุจฺฉานิทฺเทส}
\csromanhangingbegin
\hangingheaditem{76}{\csromanfolio{159}อุเปกฺขาสติสํสุทฺธํ, ธมฺมตกฺกปุเรชวํ.}
\hangingbody{อญฺญาวิโมกฺขํ ปพฺรูมิ, \textbf{อวิชฺชา}ย ปเภทนํ.}
\csromanhangingend

\prose{\textbf{อุเปกฺขาสติสํสุทฺธนฺ}ติ \textbf{อุเปกฺขา}ติ ยา จตุตฺเถ ฌาเน \textbf{อุเปกฺขา} อุเปกฺขนา อชฺฌุเปกฺขนา จิตฺตสมตา\footnote{จิตฺตสมโถ (สฺยา) ขุ 7. 402 ปิฏฺเฐปิ.} จิตฺตปฺปสฺสทฺธตา มชฺฌตฺตตา จิตฺตสฺส.\csromanspacer{} \textbf{สตี}ติ ยา จตุตฺเถ ฌาเน อุเปกฺขํ อารพฺภสติ อนุสฺสติ ฯเปฯ สมฺมาสติ.\csromanspacer{} \textbf{อุเปกฺขาสติสํสุทฺธนฺ}ติ จตุตฺเถ ฌาเน \textbf{อุเปกฺขา} จ สติ จ สุทฺธา โหนฺติ วิสุทฺธา สํสุทฺธา ปริสุทฺธา ปริโยทาตา อนงฺคณา วิคตูปกฺกิเลสา มุทุภูตา กมฺมนิยา ฐิตา อาเนญฺชปฺปตฺตาติ \textbf{อุเปกฺขา}สติสํสุทฺธํ.}

\prose{\textbf{ธมฺมตกฺกปุเรชวนฺ}ติ ธมฺมตกฺโก วุจฺจติ สมฺมาสงฺกปฺโป, โส อาทิโต โหติ, ปุรโต โหติ, ปุพฺพงฺคโม โหติ อญฺญาวิโมกฺขสฺสาติ เอวมฺปิ ธมฺมตกฺกปุเรชวํ.\csromanspacer{} อถ วา ธมฺมตกฺโก วุจฺจติ สมฺมาทิฏฺฐิ, สา อาทิโต โหติ, ปุรโต โหติ, ปุพฺพงฺคโม โหติ อญฺญาวิโมกฺขสฺสาติ เอวมฺปิ ธมฺมตกฺกปุเรชวํ.\csromanspacer{} อถ วา ธมฺมตกฺโก วุจฺจติ จตุนฺนํ มคฺคานํ ปุพฺพภาควิปสฺสนา, สา อาทิโต โหติ, ปุรโต โหติ, ปุพฺพงฺคโม โหติ อญฺญาวิโมกฺขสฺสาติ เอวมฺปิ ธมฺมตกฺกปุเรชวํ.}

\prose{\textbf{อญฺญาวิโมกฺขํ ปพฺรูมี}ติ อญฺญาวิโมกฺโข วุจฺจติ อรหตฺตวิโมกฺโข, อรหตฺตวิโมกฺขํ ปพฺรูมิ อาจิกฺขามิ เทเสมิ ปญฺญเปมิ ปฏฺฐเปมิ วิวรามิ วิภชามิ อุตฺตานีกโรมิ ปกาเสมีติ อญฺญาวิโมกฺขํ ปพฺรูมิ.}

\prose{\textbf{อวิชฺชาย ปเภทนนฺ}ติ \textbf{อวิชฺชา}ติ ทุกฺเข อญฺญาณํ ฯเปฯ \textbf{อวิชฺชา} โมโห อกุสลมูลํ.\csromanspacer{} \textbf{ปเภทนนฺ}ติ \textbf{อวิชฺชา}ย ปเภทนํ ปหานํ วูปสมํ ปฏินิสฺสคฺคํ ปฏิปฺปสฺสทฺธิํ อมตํ นิพฺพานนฺติ \textbf{อวิชฺชา}ย ปเภทนํ.\csromanspacer{} เตนาห ภควา\csromandash{}}

\prose{อุเปกฺขาสติสํสุทฺธํ, ธมฺมตกฺกปุเรชวํ.}

\csromangathameasure{อญฺญาวิโมกฺขํ ปพฺรูมิ, อวิชฺชาย ปเภทนนฺติ.}
\csromangathagroup{\csromangathastanza{อญฺญาวิโมกฺขํ ปพฺรูมิ, อวิชฺชาย ปเภทนนฺติ.}}

\csromanhangingbegin
\hangingheaditem{77}{กิํสุ สํโยชโน โลโก, กิํสุ ตสฺส วิจารณํ.}
\hangingbody{กิสฺสสฺส วิปฺปหาเนน, นิพฺพานํ อิติ วุจฺจติ.}
\csromanhangingend

\prose{\csromanfolio{160}\textbf{กิํสุ สํโยชโน โลโก}ติ โลกสฺส สํโยชนํ ลคฺคนํ พนฺธนํ อุปกฺกิเลโส, เกน โลโก ยุตฺโต ปยุตฺโต อายุตฺโต สมายุตฺโต ลคฺโค ลคฺคิโต ปลิพุทฺโธติ กิํสุ สํโยชโน โลโก.}

\prose{\textbf{กิํสุ ตสฺส วิจารณนฺ}ติ กิํสุ ตสฺส จารณํ วิจารณํ ปฏิวิจารณํ, เกน โลโก จรติ วิจรติ ปฏิวิจรตีติ กิํสุ ตสฺส วิจารณํ.\csromanspacer{} \textbf{กิสฺสสฺส วิปฺปหาเนน, นิพฺพานํ อิติ วุจฺจตี}ติ กิสฺสสฺส วิปฺปหาเนน วูปสเมน ปฏินิสฺสคฺเคน ปฏิปฺปสฺสทฺธิยา นิพฺพานํ อิติ วุจฺจติ ปวุจฺจติ กถียติ ภณียติ ทีปียติ โวหรียตีติ กิสฺสสฺส วิปฺปหาเนน, นิพฺพานํ อิติ วุจฺจติ.\csromanspacer{} เตนาห โส พฺราหฺมโณ\csromandash{}}

\csromangathameasure{กิํสุ สํโยชโน โลโก, กิํสุ ตสฺส วิจารณํ. \\ กิสฺสสฺส วิปฺปหาเนน, นิพฺพานํ อิติ วุจฺจตีติ.}
\csromangathagroup{\csromangathastanza{กิํสุ สํโยชโน โลโก, กิํสุ ตสฺส วิจารณํ. \\ กิสฺสสฺส วิปฺปหาเนน, นิพฺพานํ อิติ วุจฺจตีติ.}}

\csromanhangingbegin
\hangingheaditem{78}{\textbf{นนฺทิสํโยชโน โลโก}, วิตกฺกสฺส วิจารณา.}
\hangingbody{ตณฺหาย วิปฺปหาเนน, นิพฺพานํ อิติ วุจฺจติ.}
\csromanhangingend

\prose{\textbf{นนฺทิสํโยชโน โลโก}ติ นนฺที วุจฺจติ \textbf{ตณฺหา}.\csromanspacer{} โย ราโค สาราโค ฯเปฯ อภิชฺฌา โลโภ อกุสลมูลํ, อยํ วุจฺจติ นนฺที.\csromanspacer{} ยา นนฺที โลกสฺส สํโยชนํ ลคฺคนํ พนฺธนํ อุปกฺกิเลโส, อิมาย นนฺทิยา โลโก ยุตฺโต ปยุตฺโต อายุตฺโต สมายุตฺโต ลคฺโค ลคฺคิโต ปลิพุทฺโธติ นนฺทิสํโยชโน โลโก.}

\prose{\textbf{วิตกฺกสฺส วิจารณา}ติ \textbf{วิตกฺกา}ติ นว \textbf{วิตกฺกา} กามวิตกฺโก พฺยาปาทวิตกฺโก วิหิํสาวิตกฺโก ญาติวิตกฺโก ชนปทวิตกฺโก อมราวิตกฺโก ปรานุทยตาปฏิสํยุตฺโต วิตกฺโก ลาภสกฺการสิโลกปฏิสํยุตฺโต วิตกฺโก อนวญฺญตฺติปฏิสํยุตฺโต วิตกฺโก, อิเม วุจฺจนฺติ นว \textbf{วิตกฺกา}.\csromanspacer{} อิเม นว \textbf{วิตกฺกา} โลกสฺส จารณา วิจารณา ปฏิวิจารณา, อิเมหิ นวหิ วิตกฺเกหิ โลโก จรติ วิจรติ ปฏิวิจรตีติ วิตกฺกสฺส วิจารณา}

\prose{\textbf{ตณฺหาย วิปฺปหาเนน, นิพฺพานํ อิติ วุจฺจตี}ติ \textbf{ตณฺหา}ติ รูป\textbf{ตณฺหา} ฯเปฯ ธมฺม\textbf{ตณฺหา}.\csromanspacer{} \textbf{ตณฺหาย วิปฺปหาเนน, นิพฺพานํ อิติ วุจฺจตี}ติ \textbf{ตณฺหา}ย วิปฺปหาเนน วูปสเมน ปฏินิสฺสคฺเคน ปฏิปฺปสฺสทฺธิยา นิพฺพานํ อิติ วุจฺจติ ปวุจฺจติ กถียติ ภณียติ ทีปียติ โวหรียตีติ \textbf{ตณฺหา}ย วิปฺปหาเนน, นิพฺพานํ อิติ วุจฺจติ.\csromanspacer{} เตนาห ภควา\csromandash{}}

\csromanheader{2}{13. อุทยมาณวปุจฺฉานิทฺเทส}
\csromangathameasure{นนฺทิสํโยชโน โลโก, วิตกฺกสฺส วิจารณา. \\ ตณฺหาย วิปฺปหาเนน, นิพฺพานํ อิติ วุจฺจตีติ.}
\csromangathagroup{\csromangathastanza{\csromanfolio{161}นนฺทิสํโยชโน โลโก, วิตกฺกสฺส วิจารณา. \\ ตณฺหาย วิปฺปหาเนน, นิพฺพานํ อิติ วุจฺจตีติ.}}

\csromanhangingbegin
\hangingheaditem{79}{\textbf{กถํ สตสฺส จรโต}, วิญฺญาณํ อุปรุชฺฌติ.}
\hangingbody{\textbf{ภควนฺตํ ปุฏฺฐุมาคมา}, \textbf{ต}ํ สุโณม วโจ \textbf{ต}ว.}
\csromanhangingend

\prose{\textbf{กถํ สตสฺส จรโต}ติ กถํ ส\textbf{ต}สฺส สมฺปชานสฺส จรโต วิหรโต อิริยโต วตฺ\textbf{ต}ยโต ปาลยโต ยปยโต ยาปยโตติ กถํ ส\textbf{ต}สฺส จรโต.}

\prose{\textbf{วิญฺญาณํ อุปรุชฺฌตี}ติ วิญฺญาณํ นิรุชฺฌติ วูปสมฺมติ อตฺถํ คจฺฉติ ปฏิปฺปสฺสมฺภตีติ วิญฺญาณํ อุปรุชฺฌติ.}

\prose{\textbf{ภควนฺตํ ปุฏฺฐุมาคมา}ติ พุทฺธํ ภควนฺ\textbf{ต}ํ ปุฏฺฐุํ ปุจฺฉิตุํ ยาจิตุํ อชฺเฌสิตุํ ปสาเทตุํ อาคมฺหา อาค\textbf{ต}มฺหา อุปาค\textbf{ต}มฺหา สมฺปตฺ\textbf{ต}มฺหา \textbf{ต}ยา สทฺธิํ สมาค\textbf{ต}มฺหาติ \textbf{ภควนฺตํ ปุฏฺฐุมาคมา}.}

\prose{\textbf{ตํ สุโณม วโจ ตวา}ติ \textbf{ต}นฺติ ตุยฺหํ วจนํ พฺยปฺปถํ เทสนํ อนุสาสนํ อนุสิฏฺฐํ สุโณม อุคฺคณฺหาม ธาเรม อุปธาเรม อุปลกฺเขมาติ \textbf{ต}ํ สุโณม วโจ \textbf{ต}ว.\csromanspacer{} เตนาห โส พฺราหฺมโณ\csromandash{}}

\prose{\textbf{กถํ สตสฺส จรโต}, วิญฺญาณํ อุปรุชฺฌติ.}

\csromangathameasure{\textbf{ภควนฺตํ ปุฏฺฐุมาคมา}, \textbf{ต}ํ สุโณม วโจ \textbf{ต}วาติ.}
\csromangathagroup{\csromangathastanza{\textbf{ภควนฺตํ ปุฏฺฐุมาคมา}, \textbf{ต}ํ สุโณม วโจ \textbf{ต}วาติ.}}

\csromanhangingbegin
\hangingheaditem{80}{\textbf{อชฺฌตฺตญฺจ พหิทฺธา จ, เวทนํ นาภินนฺทโต}.}
\hangingbody{เอวํ ส\textbf{ต}สฺส จรโต, วิญฺญาณํ อุปรุชฺฌติ.}
\csromanhangingend

\prose{\textbf{อชฺฌตฺตญฺจ พหิทฺธา จ, เวทนํ นาภินนฺทโต}ติ อชฺฌตฺ\textbf{ต}ํ เวทนาสุ เวทนานุปสฺสี วิหรนฺโต เวทนํ นาภินนฺทติ นาภิวทติ น อชฺโฌเสติ\footnote{น อชฺโฌสาย ติฏฺฐติ (สฺยา)}, อภินนฺทนํ อภิวทนํ อชฺโฌสานํ คาหํ ปรามาสํ อภินิเวสํ ปชหติ วิโนเทติ พฺยนฺตีกโรติ อนภาวํ คเมติ.\csromanspacer{} พหิทฺธา เวทนาสุ เวทนานุปสฺสี วิหรนฺโต เวทนํ นาภินนฺทติ นาภิวทติ น อชฺโฌเสติ, อภินนฺทนํ อภิวทนํ อชฺโฌสานํ คาหํ ปรามาสํ อภินิเวสํ ปชหติ วิโนเทติ พฺยนฺตีกโรติ อนภาวํ คเมติ.\csromanspacer{} อชฺฌตฺ\textbf{ต}พหิทฺธา เวทนาสุ เวทนานุปสฺสี วิหรนฺโต เวทนํ นาภินนฺทติ นาภิวทติ น อชฺโฌเสติ, อภินนฺทนํ \csromanfolio{162}อภิวทนํ อชฺโฌสานํ คาหํ ปรามาสํ อภินิเวสํ ปชหติ วิโนเทติ พฺยนฺตีกโรติ อนภาวํ คเมติ.\csromanspacer{} อชฺฌตฺตํ สมุทยธมฺมานุปสฺสี เวทนาสุ เวทนานุปสฺสี\footnote{อิทํ ปทํ นตฺถิ สฺยา-โปตฺถเก.} วิหรนฺโต เวทนํ นาภินนฺทติ นาภิวทติ น อชฺโฌเสติ, อภินนฺทนํ อภิวทนํ อชฺโฌสานํ คาหํ ปรามฺสํ อภินิเวสํ ปชหติ วิโนเทติ พฺยนฺตีกโรติ อนภาวํ คเมติ.\csromanspacer{} อชฺฌตฺตํ วยธมฺมานุปสฺสี เวทนาสุ เวทนานุปสฺสี วิหรนฺโต ฯเปฯ.\csromanspacer{} อชฺฌตฺตํ สมุทยวยธมฺมานุปสฺสี เวทนาสุ เวทนานุปสฺสี วิหรนฺโต ฯเปฯ.\csromanspacer{} พหิทฺธา สมุทยธมฺมานุปสฺสี เวทนาสุ เวทนานุปสฺสี วิหรนฺโต เวทนํ นาภินนฺทติ นาภิวทติ น อชฺโฌเสติ, อภินนฺทนํ อภิวทนํ อชฺโฌสนํ คาหํ ปรามาสํ อภินิเวสํ ปชหติ วิโนเทติ พฺยนฺตีกโรติ อนภาวํ คเมติ.\csromanspacer{} พหิทฺธา วยธมฺมานุปสฺสี เวทนาสุ เวทนานุปสฺสี วิหรนฺโต ฯเปฯ.\csromanspacer{} พหิทฺธา สมุทยวยธมฺมานุปสฺสี เวทนาสุ เวทนานุปสฺสี วิหรนฺโต ฯเปฯ.\csromanspacer{} อชฺฌตฺตพหิทฺธา สมุทยธมฺมานุปสฺสี เวทนาสุ เวทนานุปสฺสี วิหรนฺโต ฯเปฯ.\csromanspacer{} อชฺฌตฺตพหิทฺธา วยธมฺมานุปสฺสี เวทนาสุ เวทนานุปสฺสี วิหรนฺโต ฯเปฯ.\csromanspacer{} อชฺฌตฺตพหิทฺธา สมุทยวยธมฺมานุปสฺสี เวทนาสุ เวทนานุปสฺสี วิหรนฺโต เวทนํ นาภินนฺทติ นาภิวทติ น อชฺโฌเสติ, อภินนฺทนํ อภิวทนํ อชฺโฌสานํ คาหํ ปรามาสํ อภินิเวสํ ปชหติ วิโนเทติ พฺยนฺตีกโรติ อนภาวํ คเมติ.\csromanspacer{} อิเมหิ ทฺวาทสหิ อากาเรหิ เวทนาสุ เวทนานุปสฺสี วิหรนฺโต ฯเปฯ อนภาวํ คเมติ.}

\prose{อถ วา เวทนํ อนิจฺจโต ปสฺสนฺโต เวทนํ นาภินนฺทติ นาภิวทติ น อชฺโฌเสติ, อภินนฺทนํ อภิวทนํ อชฺโฌสานํ คาหํ ปรามาสํ อภินิเวสํ ปชหติ วิโนเทติ พฺยนฺตีกโรติ อนภาวํ คเมติ.\csromanspacer{} เวทนํ ทุกฺขโต โรคโต คณฺฑโต สลฺลโต อฆโต อาพาธโต ฯเปฯ นิสฺสรณโต ปสฺสนฺโต เวทนํ นาภินนฺทติ นาภิวทติ น อชฺโฌเสติ, อภินนฺทนํ อภิวทนํ อชฺโฌสานํ คาหํ ปรามาสํ อภินิเวสํ ปชหติ วิโนเทติ พฺยนฺตีกโรติ อนภาวํ คเมติ.\csromanspacer{} อิเมหิ จตฺตาลีสาย\footnote{ทฺวาจตฺตาฬีสาย (สฺยา)} อากาเรหิ เวทนาสุ เวทนานุปสฺสี วิหรนฺโต เวทนํ นาภินนฺทติ นาภิวทติ น อชฺโฌเสติ, อภินนฺทนํ อภิวทนํ อชฺโฌสานํ คาหํ ปรามาสํ อภินิเวสํ ปชหติ วิโนเทติ พฺยนฺตีกโรติ อนภาวํ คเมตีติ อชฺฌตฺตญฺจ พหิทฺธา จ, เวทนํ นาภินนฺทโต.}

\csromanmatikaanchor{mtk.35}
\csromantocmarkref{subsection}{14. โปสาลมาณวปุจฺฉานิทฺเทส}{mtk.35}
\bookmark[dest={mtk.35},level=2]{14. โปสาลมาณวปุจฺฉานิทฺเทส}
\csromanheader{2}{14. โปสาลมาณวปุจฺฉานิทฺเทส}
\prose{\csromanfolio{163}\textbf{เอวํ สตสฺส จรโต}ติ เอวํ สตสฺส สมฺปชานสฺส จรโต วิหรโต อิริยโต วตฺตยโต ปาลยโต ยปยโต ยาปยโตติ เอวํ สตสฺส จรโต.}

\prose{\textbf{วิญฺญาณํ อุปรุชฺฌตี}ติ ปุญฺญาภิสงฺขารสหคตํ วิญฺญาณํ อปุญฺญาภิสงฺขารสหคตํ วิญฺญาณํ อาเนญฺชาภิสงฺขารสหคตํ วิญฺญาณํ นิรุชฺฌติ วูปสมฺมติ อตฺถํ คจฺฉติ ปฏิปฺปสฺสมฺภตีติ วิญฺญาณํ อุปรุชฺฌติ.\csromanspacer{} เตนาห ภควา\csromandash{}}

\csromangathameasure{อชฺฌตฺตญฺจ พหิทฺธา จ, เวทนํ นาภินนฺทโต. \\ เอวํ สตสฺส จรโต, วิญฺญาณํ อุปรุชฺฌตีติ.}
\csromangathagroup{\csromangathastanza{อชฺฌตฺตญฺจ พหิทฺธา จ, เวทนํ นาภินนฺทโต. \\ เอวํ สตสฺส จรโต, วิญฺญาณํ อุปรุชฺฌตีติ.}}

\prose{สห คาถาปริ\textbf{โย}สานา ฯเปฯ “สตฺถา เม ภนฺเต ภควา, สาวโกหมสฺมี”ติ.}

\vspace{\tipitakaparskip}
\csromancenter{อุทยมาณวปุจฺฉานิทฺเทโส เตรสโม.}
\csromanheader{2}{14. โปสาลมาณวปุจฺฉานิทฺเทส}
\csromanhangingbegin
\hangingheaditem{81}{โย อตีตํ อาทิสติ, (อิจฺจายสฺมา โปสาโล,)}
\hangingbody{\textbf{อเนโช ฉินฺนสํสโย.}}
\hangingbody{\textbf{ปารคุํ1} สพฺพธมฺมานํ, อตฺถิ ปเญฺหน อาคมํ.}
\csromanhangingend

\prose{\textbf{โย อตีตํ อาทิสตี}ติ \textbf{โย}ติ \textbf{โย} โส ภควา สยมฺภู อนาจริยโก ปุพฺเพ อนนุสฺสุเตสุ ธมฺเมสุ สามํ สจฺจานิ อภิสมฺพุชฺฌิ, ตตฺถ จ สพฺพญฺญุตํ ปตฺโต, พเลสุ จ วสีภาวํ.\csromanspacer{} \textbf{อตีตํ อาทิสตี}ติ ภควา อตฺตโน จ ปเรสญฺจ อตีตมฺปิ อาทิสติ, อนาคตมฺปิ อาทิสติ, ปจฺจุปฺปนฺนมฺปิ อาทิสติ.}

\prose{กถํ ภควา อตฺตโน อตีตํ อาทิสติ, ภควา อตฺตโน อตีตํ เอกมฺปิ ชาติํ อาทิสติ, เทฺวปิ ชาติ\textbf{โย} อาทิสติ, ติสฺโสปิ ชาติ\textbf{โย} อาทิสติ, จตสฺโสปิ ชาติ\textbf{โย} อาทิสติ, ปญฺจปิ ชาติ\textbf{โย} อาทิสติ, ทสปิ ชาติ\textbf{โย} อาทิสติ, วีสมฺปิ ชาติ\textbf{โย} อาทิสติ, ติํสมฺปิ ชาติ\textbf{โย} อาทิสติ, จตฺตาลีสมฺปิ ชาติ\textbf{โย} \csromanfolio{164}อาทิสติ, ปญฺญาสมฺปิ ชาติโย อาทิสติ, ชาติสตมฺปิ.\csromanspacer{} ชาติสหสฺสมฺปิ.\csromanspacer{} ชาติสตสหสฺสมฺปิ.\csromanspacer{} อเนเกปิ สํวฏฺฏกปฺเป.\csromanspacer{} อเนเกปิ วิวฏฺฏกปฺเป.\csromanspacer{} อเนเกปิ สํวฏฺฏวิวฏฺฏกปฺเป อาทิสติ, อมุตฺราสิํ เอวํนาโม เอวํโคตฺโต เอวํวณฺโณ เอวมาหาโร เอวํสุขทุกฺขปฺปฏิสํเวที เอวมายุปริยนฺโต โส ตโต จุโต อมุตฺร อุทปาทิํ, ตตฺราปาสิํ “เอวํนาโม เอวํโคตฺโต เอวํวณฺโณ เอวมาหาโร เอวํสุขทุกฺขปฺปฏิสํเวที เอวมายุปริยนฺโต, โส ตโต จุโต อิธูปปนฺโน”ติ อิติ สาการํ ส-อุทฺเทสํ อเนกวิหิตํ ปุพฺเพนิวาสํ อาทิสติ.\csromanspacer{} เอวํ ภควา อตฺตโน อตีตํ อาทิสติ.}

\prose{กถํ ภควา ปเรสํ อตีตํ อาทิสติ, ภควา ปเรสํ อตีตํ เอกมฺปิ ชาติํ อาทิสติ, เทฺวปิ ชาติโย อาทิสติ ฯเปฯ.\csromanspacer{} อเนเกปิ สํวฏฺฏวิวฏฺฏกปฺเป อาทิสติ, อมุตฺราสิ เอวํนาโม เอวํโคตฺโต เอวํวณฺโณ เอวมาหาโร เอวํสุขทุกฺขปฺปฏิสํเวที เอวมายุปริยนฺโต, โส ตโต จุโต อมุตฺร อุทปาทิ, ตตฺราปาสิ “เอวํนาโม เอวํโคตฺโต เอวํวณฺโณ เอวมาหาโร เอวํสุขทุกฺขปฺปฏิสํเวที เอวมายุปริยนฺโต, โส ตโต จุโต อิธูปปนฺโน”ติ อิติ สาการํ ส-อุทฺเทสํ อเนกวิหิตํ ปุพฺเพนิวาสํ อาทิสติ.\csromanspacer{} เอวํ ภควา ปเรสํ อตีตํ อาทิสติ.}

\prose{ภควา ปญฺจ ชาตกสตานิ ภาสนฺโต อตฺตโน จ ปเรสญฺจ อตีตํ อาทิสติ, มหาปทานิยสุตฺตนฺตํ\footnote{มหาธนิยสุตฺตํ (สฺยา)} ภาสนฺโต อตฺตโน จ ปเรสญฺจ อตีตํ อาทิสติ, มหาสุทสฺสนิยสุตฺตนฺตํ ภาสนฺโต อตฺตโน จ ปเรสญฺจ อตีตํ อาทิสติ, มหาโควินฺทิยสุตฺตนฺตํ ภาสนฺโต อตฺตโน จ ปเรสญฺจ อตีตํ อาทิสติ, มาฆเทวิยสุตฺตนฺตํ ภาสนฺโต อตฺตโน จ ปเรสญฺจ อาทิสติ.}

\prose{วุตฺตํ เหตํ ภควตา\csromandash{}อตีตํ โข จุนฺท อทฺธานํ อารพฺภ ตถาคตสฺส สตานุสาริญาณํ\footnote{สตานุสฺสริยญาณํ (ก) ปสฺส ที 3. 111 ปิฏฺเฐ.} โหติ, โส ยาวตกํ อากงฺขติ, ตาวตกํ อนุสฺสรติ.\csromanspacer{} อนาคตญฺจ โข จุนฺท.\csromanspacer{} ปจฺจุปฺปนฺนญฺจ โข จุนฺท อทฺธานํ อารพฺภ ตถาคตสฺส โพธิชํ ญาณํ อุปฺปชฺชติ “อยมนฺติมา ชาติ, นตฺถิทานิ ปุนพฺภโว”ติ.}

\csromanheader{2}{14. โปสาลมาณวปุจฺฉานิทฺเทส}
\prose{\csromanfolio{165}อินฺทฺริยปโรปริยตฺตญาณํ\footnote{อินฺทฺริยปโรปริยตฺติญาณํ (ก) อฏฺฐกถา โอโลเกตพฺพา.} ตถาคตสฺส ตถาคตพลํ, สตฺตานํ อาสยานุสยญาณํ ตถาคตสฺส ตถาคตพลํ, ยมกปาฏิหีเร ญาณํ\footnote{ยมกปาฏิหิริยญาณํ (สฺยา)} ตถาคตสฺส ตถาคตพลํ, มหากรุณาสมาปตฺติยา ญาณํ ตถาคตสฺส ตถาคตพลํ, สพฺพญฺญุตญาณํ ตถาคตสฺส ตถาคตพลํ, อนาวรณญาณํ ตถาคตสฺส ตถาคตพลํ, สพฺพตฺถ อสงฺคมปฺปฏิหตมนาวรณญาณํ ตถาคตสฺส ตหาคตพลํ.\csromanspacer{} เอวํ ภควา อตฺตโน จ ปเรสญฺจ อตีตมฺปิ อาทิสติ, อนาคตมฺปิ อาทิสติ, ปจฺจุปฺปนฺนมฺปิ อาทิสติ อาจิกฺขติ เทเสติ ปญฺญเปติ ปฏฺฐเปติ วิวรติ วิภชติ อุตฺตานีกโรติ ปกาเสตีติ โย อตีตํ อาทิสติ.}

\prose{\textbf{อิจฺจายสฺมา โปสาโล}ติ \textbf{อิจฺจา}ติ ปทสนฺธิ ฯเปฯ.\csromanspacer{} อา\textbf{ยสฺมา}ติ ปิยวจนํ ฯเปฯ.\csromanspacer{} \textbf{โปสาโล}ติ ตสฺส พฺราหฺมณสฺส นามํ ฯเปฯ อภิลาโปติ \textbf{อิจฺจายสฺมา} \textbf{โปสาโล}.}

\prose{\textbf{อเนโช ฉินฺนสํสโย}ติ เอชา วุจฺจติ ตณฺหา.\csromanspacer{} โย ราโค สาราโค ฯเปฯ อภิชฺฌา โลโภ อกุสลมูลํ.\csromanspacer{} สา เอชา ตณฺหา พุทฺธสฺส ภควโต ปหีนา อุจฺฉินฺนมูลา ตาลาวตฺถุกตา อนภาวํกตา อายติํ อนุปฺปาทธมฺมา, ตสฺมา พุทฺโธ อเนโช.\csromanspacer{} เอชาย ปหีนตฺตา อเนโช, ภควา ลาเภปิ น อิญฺชติ ฯเปฯ ทุกฺเขปิ น อิญฺชติ น จลติ น เวธติ นปฺปเวธติ น สมฺปเวธตีติ อเนโช.\csromanspacer{} \textbf{ฉินฺนสํสโย}ติ สํสโย วุจฺจติ วิจิกิจฺฉา.\csromanspacer{} ทุกฺเข กงฺขา ฯเปฯ ฉมฺภิตตฺตํ จิตฺตสฺส มโนวิเลโข, โส สํสโย พุทฺธสฺส ภควโต ปหีโน ฉินฺโน อุจฺฉินฺโน สมุจฺฉินฺโน วูปสนฺโต ปฏินิสฺสคฺโค ปฏิปฺปสฺสทฺโธ อภพฺพุปฺปตฺติโก ญาณคฺคินา ทฑฺโฒ, ตสฺมา พุทฺโธ ฉินฺนสํสโยติ อเนโช ฉินฺนสํสโย.}

\prose{\textbf{ปารคุํ สพฺพธมฺมานนฺ}ติ ภควา สพฺพธมฺมานํ อภิญฺญาปารคู ปริญฺญาปารคู ปหานปารคู ภาวนาปารคู สจฺฉิกิริยาปารคู สมาปตฺติปารคู.\csromanspacer{} อภิญฺญาปารคู สพฺพธมฺมานํ ฯเปฯ ชาติมรณสํสาโร, นตฺถิ ตสฺส ปุนพฺภโวติ ปารคู สพฺพธมฺมานํ.}

\prose{\textbf{อตฺถิ ปเญฺหน อาคมนฺ}ติ ปเญฺหน อตฺถิโก อาคโตมฺหิ ฯเปฯ วหสฺเสตํ ภารนฺติ อตฺถิ ปเญฺหน อาคมํ.\csromanspacer{} เตนาห โส พฺราหฺมโณ\csromandash{}}

\prose{\csromanfolio{166}โย อตีตํ อาทิสติ, (อิจฺจายสฺมา โปสาโล,)}

\prose{อเนโช ฉินฺนสํสโย.}

\csromangathameasure{ปารคุํ สพฺพธมฺมานํ, อตฺถิ ปเญฺหน อาคมนฺติ.}
\csromangathagroup{\csromangathastanza{ปารคุํ สพฺพธมฺมานํ, อตฺถิ ปเญฺหน อาคมนฺติ.}}

\csromanhangingbegin
\hangingheaditem{82}{วิภูตรูปสญฺญิสฺส, สพฺพกายปฺปหายิโน.}
\hangingbody{\textbf{อชฺฌตฺตญฺจ พหิทฺธา จ, นตฺถิ กิญฺจีติ ปสฺสโต.}}
\hangingbody{ญาณํ \textbf{สกฺกา}’นุปุจฺฉามิ, กถํ เนยฺโย ตถาวิโธ.}
\csromanhangingend

\prose{\textbf{วิภูตรูปสญฺญิสฺสา}ติ กตมา รูปสญฺญา, รูปาวจรสมาปตฺติํ สมาปนฺนสฺส วา อุปปนฺนสฺส วา ทิฏฺฐธมฺมสุขวิหาริสฺส วา สญฺญา สญฺชานนา สญฺชานิตตฺตํ, อยํ รูปสญฺญา.\csromanspacer{} \textbf{วิภูตรูปสญฺญิสฺสา}ติ จตสฺโส อรูปสมาปตฺติโย ปฏิลทฺธสฺส\footnote{ลาภิสฺส (สฺยา)} รูปสญฺญา วิภูตา โหนฺติ วิคตา อติกฺกนฺตา สมติกฺกนฺตา วีติวตฺตาติ วิภูตรูปสญฺญิสฺส.}

\prose{\textbf{สพฺพกายปฺปหายิโน}ติ สพฺโพ ตสฺส ปฏิสนฺธิโก รูปกาโย ปหีโน, ตทงฺคสมติกฺกมา วิกฺขมฺภนปฺปหาเนน ปหีโน ตสฺส รูปกาโยติ สพฺพกายปฺปหายิโน.}

\prose{\textbf{อชฺฌตฺตญฺจ พหิทฺธา จ, นตฺถิ กิญฺจีติ ปสฺสโตติ นตฺถิ กิญฺจี}ติ อากิญฺจญฺญายตนสมาปตฺติ.\csromanspacer{} กิํ การณา “นตฺถิ กิญฺจี”ติ อากิญฺจญฺญายตนสมาปตฺติ.\csromanspacer{} ยํ วิญฺญาณญฺจายตนสมาปตฺติํ สโต สมาปชฺชิตฺวา ตโต วุฏฺฐหิตฺวา ตญฺเญว วิญฺญาณํ อภาเวติ วิภาเวติ อนฺตรธาเปติ, “นตฺถิ กิญฺจี”ติ ปสฺสติ, ตํการณา “นตฺถิ กิญฺจี”ติ อากิญฺจญฺญายตนสมาปตฺตีติ อชฺฌตฺตญฺจ พหิทฺธา จ, นตฺถิ กิญฺจีติ ปสฺสโต.}

\prose{\textbf{ญาณํ สกฺกา’นุปุจฺฉามี}ติ \textbf{สกฺกา}ติ สกฺโก, ภควา สกฺยกุลา ปพฺพชิโตติปิ สกฺโก ฯเปฯ ปหีนภยเภรโว วิคตโลมหํโสติปิ สกฺโก.\csromanspacer{} \textbf{ญาณํ สกฺกา’นุปุจฺฉามี}ติ ตสฺส ญาณํ ปุจฺฉามิ, ปญฺญํ ปุจฺฉามิ, สมฺพุทฺธํ ปุจฺฉามิ.\csromanspacer{} กีทิสํ กิํ สณฺฐิตํ กิํปการํ กิํปฏิภาคํ ญาณํ อิจฺฉิตพฺพนฺติ ญาณํ \textbf{สกฺกา}’นุปุจฺฉามิ.}

\prose{\textbf{กถํ เนยฺโย ตถาวิโธ}ติ กถํ โส เนตพฺโพ วิเนตพฺโพ อนุเนตพฺโพ ปญฺญาเปตพฺโพ นิชฺฌาเปตพฺโพ เปกฺเขตพฺโพ ปสาเทตพฺโพ, กถํ เตน\footnote{กถมสฺส (สฺยา)} อุตฺตริ ญาณํ อุปฺปาเทตพฺพํ.\csromanspacer{} \textbf{ตถาวิโธ}ติ ตถาวิโธ}

\vspace{\tipitakaparskip}
\csromancenter{โปสาลมาณวปุจฺฉานิทฺเทส ตาทิโส ตสฺสณฺฐิโต ตปฺปกาโร ตปฺปฏิภาโค, โย โส อากิญฺจญฺญายตนสมาปตฺติลาภีติ กถํ เนยฺโย ตถาวิโธ.\csromanspacer{} เตนาห โส พฺราหฺมโณ\csromandash{}}
\vspace{0.5\baselineskip}
\csromangathameasure{วิภูตรูปสญฺญิสฺส, สพฺพกายปฺปหายิโน. \\ อชฺฌตฺตญฺจ พหิทฺธา จ, นตฺถิ กิญฺจีติ ปสฺสโต. \\ ญาณํ สกฺกา’นุปุจฺฉามิ, กถํ เนยฺโย ตถาวิโธติ.}
\csromangathagroup{\csromangathastanza{\csromanfolio{167}วิภูตรูปสญฺญิสฺส, สพฺพกายปฺปหายิโน. \\ อชฺฌตฺตญฺจ พหิทฺธา จ, นตฺถิ กิญฺจีติ ปสฺสโต. \\ ญาณํ สกฺกา’นุปุจฺฉามิ, กถํ เนยฺโย ตถาวิโธติ.}}

\csromanhangingbegin
\hangingheaditem{83}{\textbf{วิญฺญาณฏฺฐิติโย สพฺพา}, (โปสาลาติ ภควา,)}
\hangingbody{\textbf{อภิชานํ ตถาคโต.}}
\hangingbody{ติฏฺฐนฺตเมนํ ชานาติ, ธิมุตฺตํ ตปฺปรายณํ.}
\csromanhangingend

\prose{\textbf{วิญฺญาณฏฺฐิติโย สพฺพา}ติ ภควา อภิสงฺขารวเสน จตสฺโส วิญฺญาณฏฺฐิติโย ชานาติ, ปฏิสนฺธิวเสน สตฺต วิญฺญาณฏฺฐิติโย ชานาติ.\csromanspacer{} กถํ ภควา อภิสงฺขารวเสน จตสฺโส วิญฺญาณฏฺฐิติโย ชานาติ, วุตฺตํ เหตํ ภควตา\csromandash{}รูปุปยํ\footnote{รูปูปายํ (สฺยา, ก) ปสฺส สํ 2. 43 ปิฏฺเฐ.} วา ภิกฺขเว วิญฺญาณํ ติฏฺฐมานํ ติฏฺเฐยฺย\footnote{ติฏฺฐติ (สฺยา, ก)}, รูปารมฺมณํ รูปปฺปติฏฺฐํ นนฺทูปเสจนํ วุทฺธิํ วิรูฬฺหิํ เวปุลฺลํ อาปชฺเชยฺย.\csromanspacer{} เวทนุปยํ วา ภิกฺขเว ฯเปฯ.\csromanspacer{} สญฺญุปยํ วา ภิกฺขเว ฯเปฯ.\csromanspacer{} สงฺขารุปยํ วา ภิกฺขเว วิญฺญาณํ ติฏฺฐมานํ ติฏฺเฐยฺย, สงฺขารารมฺมณํ สงฺขารปฺปติฏฺฐํ นนฺทูปเสจนํ วุทฺธิํ วิรูฬฺหิํ เวปุลฺลํ อาปชฺเชยฺยา”ติ, เอวํ ภควา อภิสงฺขารวเสน จตสฺโส วิญฺญาณฏฺฐิติโย ชานาติ.}

\prose{กถํ ภควา ปฏิสนฺธิวเสน สตฺต วิญฺญาณฏฺฐิติโย ชานาติ, วุตฺตํ เหตํ ภควตา\csromandash{}สนฺติ ภิกฺขเว สตฺตา นานตฺตกายา นานตฺตสญฺญิโน, เสยฺยถาปิ มนุสฺสา เอกจฺเจ จ เทวา เอกจฺเจ จ วินิปาติกา, อยํ ปฐมา วิญฺญาณฏฺฐิติ.}

\prose{สนฺติ ภิกฺขเว สตฺตา นานตฺตกายา เอกตฺตสญฺญิโน, เสยฺยถาปิ เทวา พฺรหฺมกายิกา ปฐมาภินิพฺพตฺตา, อยํ ทุติยา วิญฺญาณฏฺฐิติ.}

\prose{สนฺติ ภิกฺขเว สตฺตา เอกตฺตกายา นานตฺตสญฺญิโน, เสยฺยถาปิ เทวา อาภสฺสรา, อยํ ตติยา วิญฺญาณฏฺฐิติ.}

\prose{\csromanfolio{168}สนฺติ ภิกฺขเว สตฺตา เอกตฺตกายา เอกตฺตสญฺญิโน, เสยฺยถาปิ เทวา สุภกิณฺหา, อยํ จตุตฺถี\footnote{จตุตฺถา (สฺยา) ปสฺส อํ 2. 428 ปิฏฺเฐ.} วิญฺญาณฏฺฐิติ.}

\prose{สนฺติ ภิกฺขเว สตฺตา สพฺพโส รูปสญฺญานํ สมติกฺกมา ปฏิฆสญฺญานํ อตฺถงฺคมา นานตฺตสญฺญานํ อมนสิการา “อนนฺโต อากาโส”ติ อากาสานญฺจายตนูปคา, อยํ ปญฺจมี\footnote{ปญฺจมา (สฺยา)} วิญฺญาณฏฺฐิติ.}

\prose{สนฺติ ภิกฺขเว สตฺตา สพฺพโส อากาสานญฺจายตนํ สมติกฺกมฺม “อนนฺตํ วิญฺญาณนฺ”ติ วิญฺญาณญฺจายตนูปคา, อยํ ฉฏฺฐี\footnote{ฉฏฺฐา (สฺยา)} วิญฺญาณฏฺฐิติ.}

\prose{สนฺติ ภิกฺขเว สตฺตา สพฺพโส วิญฺญาณญฺจายตนํ สมติกฺกมฺม “นตฺถิ กิญฺจี”ติ อากิญฺจญฺญายตนูปคา, อยํ สตฺตมี\footnote{สตฺตมา (สฺยา)} วิญฺญาณฏฺฐิติ.\csromanspacer{} เอวํ \textbf{ภควา} ปฏิสนฺธิวเสน สตฺต วิญฺญาณฏฺฐิติโย ชานาตีติ วิญฺญาณฏฺฐิติโย สพฺพา.}

\prose{\textbf{โปสาลาติ ภควา}ติ \textbf{โปสาลาติ ภควา} ตํ พฺราหฺมณํ นาเมน อาลปติ.\csromanspacer{} \textbf{ภควา}ติ คารวาธิวจนเมตํ ฯเปฯ สจฺฉิกา ปญฺญตฺติ, ยทิทํ \textbf{ภควา}ติ \textbf{โปสาลาติ ภควา}.}

\prose{\textbf{อภิชานํ ตถาคโต}ติ \textbf{อภิชานนฺ}ติ \textbf{อภิชานนฺ}โต วิชานนฺโต ปฏิวิชานนฺโต ปฏิวิชฺฌนฺโต ตถาคโต.\csromanspacer{} วุตฺตํ เหตํ ภควตา\footnote{ปสฺส ที 3. 111 ปิฏฺเฐ.} อตีตญฺเจปิ โข จุนฺท โหติ อภูตํ อตจฺฉํ อนตฺถสญฺหิตํ, น ตํ ตถาคโต พฺยากโรติ.\csromanspacer{} อตีตญฺเจปิ จุนฺท โหติ ภูตํ ตจฺฉํ อนตฺถสญฺหิตํ, ตมฺปิ ตถาคโต น พฺยากโรติ.\csromanspacer{} อตีตญฺเจปิ โข จุนฺท โหติ ภูตํ ตจฺฉํ อตฺถสญฺหิตํ, ตตฺร กาลญฺญู ตถาคโต โหติ ตสฺเสว ปญฺหสฺส เวยฺยากรณาย.\csromanspacer{} อนาคตญฺเจปิ จุนฺท โหติ ฯเปฯ.\csromanspacer{} ปจฺจุปฺปนฺนญฺเจปิ จุนฺท โหติ อภูตํ อตจฺฉํ อนตฺถสญฺหิตํ, น ตํ ตถาคโต พฺยากโรติ.\csromanspacer{} ปจฺจุปฺปนฺนํ เจปิ จุนฺท โหติ ภูตํ ตจฺฉํ อนตฺถสญฺหิตํ, ตมฺปิ ตถาคโต น พฺยากโรติ.\csromanspacer{} ปจฺจุปฺปนฺนญฺเจปิ จุนฺท โหติ ภูตํ ตจฺฉํ อตฺถสญฺหิตํ, ตตฺร กาลญฺญู ตถาคโต โหติ ตสฺส ปญฺหสฺส เวยฺยากรณาย.\csromanspacer{} อิติ โข จุนฺท อตีตานาคตปจฺจุปฺปนฺเนสุ ธมฺเมสุ ตถาคโต กาลวาที ภูตวาที อตฺถวาที ธมฺมวาที วินยวาที, ตสฺมา “ตถาคโต”ติ วุจฺจติ.}

\prose{ยํ โข จุนฺท สเทวกสฺส โลกสฺส สมารกสฺส สพฺรหฺมกสฺส สสฺสมณพฺราหฺมณิยา ปชาย สเทวมนุสฺสาย ทิฏฺฐํ สุตํ มุตํ วิญฺญาตํ ปตฺตํ}

\vspace{\tipitakaparskip}
\csromancenter{โปสาลมาณวปุจฺฉานิทฺเทส ปริเยสิตํ อนุวิจริตํ มนสา, สพฺพํ ตํ ตถาคเตน อภิสมฺพุทฺธํ, ตสฺมา “ตถาคโต”ติ วุจฺจติ.\csromanspacer{} ยญฺจ จุนฺท รตฺติํ ตถาคโต อนุตฺตรํ สมฺมาสมฺโพธิํ อภิสมฺพุชฺฌติ, ยญฺจ รตฺติํ อนุปาทิเสสาย นิพฺพานธาตุยา ปรินิพฺพายติ, ยํ เอตสฺมิํ อนฺตเร ภาสติ ลปติ นิทฺทิสติ.\csromanspacer{} สพฺพํ ตํ ตเถว โหติ โน อญฺญถา, ตสฺมา “ตถาคโต”ติ วุจฺจติ.\csromanspacer{} ยถาวาที จุนฺท ตถาคโต ตถาการี, ยถาการี ตถาวาที, อิติ ยถาวาที ตถาการี, ยถาการี ตถาวาที, ตสฺมา “ตถาคโต”ติ วุจฺจติ.\csromanspacer{} สเทวเก จุนฺท โลเก สมารเก สพฺรหฺมเก สสฺสมณพฺราหฺมณิยา ปชาย สเทวมนุสฺสาย ตถาคโต อภิภู อนภิภูโต อญฺญทตฺถุทโส วสวตฺถี, ตสฺมา “ตถาคโต”ติ วุจฺจตีติ อภิชานํ ตถาคโต.}
\prose{\csromanfolio{169}\textbf{ติฏฺฐนฺตเมนํ ชานาตี}ติ ภควา อิธตฺถญฺเญว\footnote{อิธฏฺฐญฺเญว (สฺยา)} ชานาติ กมฺมาภิสงฺขารวเสน “อยํ ปุคฺคโล กายสฺส เภทา ปรํ มรณา อปายํ ทุคฺคติํ วินิปาตํ นิรยํ อุปปชฺชิสฺสตี”ติ.\csromanspacer{} ภควา อิธตฺถญฺเญว ชานาติ กมฺมาภิสงฺขารวเสน “อยํ ปุคฺคโล กายสฺส เภทา ปรํ มรณา ติรจฺฉานโยนิํ อุปปชฺชิสฺสตี”ติ.\csromanspacer{} ภควา- อิธตฺถญฺเญว ชานาติ กมฺมาภิสงฺขารวเสน “อยํ ปุคฺคโล กายสฺส เภทา ปรํ มรณา เปตฺติวิสยํ อุปปชฺชิสฺสตี”ติ.\csromanspacer{} ภควา อิธตฺถญฺเญว ชานาติ กมฺมาภิสงฺขารวเสน “อยํ ปุคฺคโล กายสฺส เภทา ปรํ มรณา มนุสฺเสสุ อุปฺปชฺชิสฺสตี”ติ.\csromanspacer{} ภควา อิธตฺถญฺเญว ชานาติ กมฺมาภิสงฺขารวเสน “อยํ ปุคฺคโล สุปฺปฏิปนฺโน กายสฺส เภทา ปรํ มรณา สุคติํ สคฺคํ โลกํ อุปปชฺชิสฺสตี”ติ.}

\prose{วุตฺตํ เหตํ ภควตา\csromandash{}อิธ ปนาหํ สาริปุตฺต เอกจฺจํ ปุคฺคลํ เอวํ เจตสา เจโต ปริจฺจ ปชานามิ “ตถายํ ปุคฺคโล ปฏิปนฺโน, ตถา จ อิริยติ, ตญฺจ มคฺคํ สมารูโฬฺห, ยถา กายสฺส เภทา ปรํ มรณา อปายํ ทุคฺคติํ วินิปาตํ นิรยํ อุปปชฺชิสฺสตี”ติ.}

\prose{อิธ ปนาหํ สาริปุตฺต เอกจฺจํ ปุคฺคลํ เอวํ เจตสา เจโต ปริจฺจ ปชานามิ “ตถายํ ปุคฺคโล ปฏิปนฺโน, ตถา จ อิริยติ, ตญฺจ มคฺคํ สมารูโฬฺห, ยถา กายสฺส เภทา ปรํ มรณา ติรจฺฉานโยนิํ อุปปชฺชิสฺสตี”ติ.}

\prose{\csromanfolio{170}อิธ ปนาหํ สาริปุตฺต เอกจฺจํ ปุคฺคลํ เอวํ เจตสา เจโต ปริจฺจ ปชานามิ “ตถายํ ปุคฺคโล ปฏิปนฺโน, ตถา จ อิริยติ, ตญฺจ มคฺคํ สมารูโฬฺห, ยถา กายสฺส เภทา ปรํ มรณา เปตฺติวิสยํ อุปปชฺชิสฺสตี”ติ.}

\prose{อิธ ปนาหํ สาริปุตฺต เอกจฺจํ ปุคฺคลํ เอวํ เจตสา เจโต ปริจฺจ ปชานามิ “ตถายํ ปุคฺคโล ปฏิปนฺโน, ตถา จ อิริยติ, ตญฺจ มคฺคํ สมารูโฬฺห, ยถา กายสฺส เภทา ปรํ มรณา มนุสฺเสสุ อุปฺปชฺชิสฺสตี”ติ.}

\prose{อิธ ปนาหํ สาริปุตฺต เอกจฺจํ ปุคฺคลํ เอวํ เจตสา เจโต ปริจฺจ ปชานามิ “ตถายํ ปุคฺคโล ปฏิปนฺโน, ตถา จ อิริยติ, ตญฺจ มคฺคํ สมารูโฬฺห, ยถา กายสฺส เภทา ปรํ มรณา สุคติํ สคฺคํ โลกํ อุปปชฺชิสฺสตี”ติ.}

\prose{อิธ ปนาหํ สาริปุตฺต เอกจฺจํ ปุคฺคลํ เอวํ เจตสา เจโต ปริจฺจ ปชานามิ “ตถายํ ปุคฺคโล ปฏิปนฺโน, ตถา จ อิริยติ, ตญฺจ มคฺคํ สมารูโฬฺห, ยถา อาสวานํ ขยา อนาสวํ เจโตวิมุตฺติํ ปญฺญาวิมุตฺติํ ทิฏฺเฐว ธมฺเม สยํ อภิญฺญา สจฺฉิกตฺวา อุปสมฺปชฺช วิหรตี”ติ ติฏฺฐนฺตเมนํ ชานาติ.}

\prose{\textbf{ธิมุตฺตํ ตปฺปรายณนฺ}ติ ธิมุตฺตนฺติ อากิญฺจญฺญายตนํ ธิมุตฺตนฺติ วิโมกฺเขน ธิมุตฺตํ ตตฺราธิมุตฺตํ ตทธิมุตฺตํ ตทาธิปเตยฺยํ.\csromanspacer{} อถ วา ภควา ชานาติ “อยํ ปุคฺคโล รูปาธิมุตฺโต สทฺทาธิมุตฺโต คนฺธาธิมุตฺโต รสาธิมุตฺโต โผฏฺฐพฺพาธิมุตฺโต กุลาธิมุตฺโต คณาธิมุตฺโต อาวาสาธิมุตฺโต ลาภาธิมุตฺโต ยสาธิมุตฺโต ปสํสาธิมุตฺโต สุขาธิมุตฺโต จีวราธิมุตฺโต ปิณฺฑปาตาธิมุตฺโต เสนาสนาธิมุตฺโต คิลานปจฺจยเภสชฺชปริกฺขาราธิมุตฺโต สุตฺตนฺตาธิมุตฺโต วินยาธิมุตฺโต อภิธมฺมาธิมุตฺโต อารญฺญกงฺคาธิมุตฺโต ปิณฺฑปาติกงฺคาธิมุตฺโต ปํสุกูลิกงฺคาธิมุตฺโต เตจีวริกงฺคาธิมุตฺโต สปทานจาริกงฺคาธิมุตฺโต ขลุปจฺฉาภตฺติกงฺคาธิมุตฺโต เนสชฺชิกงฺคาธิมุตฺโต ยถาสนฺถติกงฺคาธิมุตฺโต ปฐมชฺฌานาธิมุตฺโต ทุติยชฺฌานาธิมุตฺโต ตติยชฺฌานาธิมุตฺโต จตุตฺถชฺฌานาธิมุตฺโต อากาสานญฺจายตนสมาปตฺตาธิมุตฺโต วิญฺญาณญฺจายตนสมาปตฺตาธิมุตฺโต อากิญฺจญฺญายตนสมาปตฺตาธิมุตฺโต เนวสญฺญานาสญฺญายตนสมาปตฺตาธิมุตฺโต”ติ ธิมุตฺตํ.}

\csromanheader{2}{14. โปสาลมาณวปุจฺฉานิทฺเทส}
\prose{\csromanfolio{171}\textbf{ตปฺปรายณนฺ}ติ อากิญฺจญฺญายตนมยํ ตปฺปรายณํ กมฺมปรายณํ วิปากปรายณํ กมฺมครุกํ ปฏิสนฺธิครุกํ.\csromanspacer{} อถ วา ภควา ชานาติ “อยํ ปุคฺคโล รูปปรายโณ ฯเปฯ เนวสญฺญานาสญฺญายตนสมาปตฺติปรายโณ”ติ ธิมุตฺตํ ตปฺปรายณํ.\csromanspacer{} เตนาห ภควา\csromandash{}}

\prose{วิญฺญาณฏฺฐิติโย สพฺพา, (โปสาลาติ ภควา,)}

\prose{อภิชานํ ตถาคโต.}

\csromangathameasure{ติฏฺฐนฺตเมนํ ชานาติ, ธิมุตฺตํ ตปฺปรายณนฺติ.}
\csromangathagroup{\csromangathastanza{ติฏฺฐนฺตเมนํ ชานาติ, ธิมุตฺตํ ตปฺปรายณนฺติ.}}

\csromanhangingbegin
\hangingheaditem{84}{\textbf{อากิญฺจญฺญาสมฺภวํ ญตฺวา}, นนฺทิสํโยชนํ อิติ.}
\hangingbody{\textbf{เอวเมตํ อภิญฺญาย, ตโต ตตฺถ วิปสฺสติ.}}
\hangingbody{เอตํ ญาณํ ตถํ ตสฺส, พฺราหฺมณสฺส วุสีมโต.}
\csromanhangingend

\prose{\textbf{อากิญฺจญฺญาสมฺภวํ ญตฺวา}ติ อากิญฺจญฺญาสมฺภโวติ วุจฺจติ อากิญฺจญฺญายตนสํวตฺตนิโก กมฺมาภิสงฺขาโร, อากิญฺจญฺญายตนสํวตฺตนิกํ กมฺมาภิสงฺขารํ อากิญฺจญฺญาสมฺภโวติ ญตฺวา, ลคฺคนนฺติ ญตฺวา, พนฺธนนฺติ ญตฺวา, ปลิโพโธติ ญตฺวา ชานิตฺวา ตุลยิตฺวา ตีรยิตฺวา วิภาวยิตฺวา วิภูตํ กตฺวาติ อากิญฺจญฺญาสมฺภวํ ญตฺวา.}

\prose{\textbf{นนฺทิสํโยชนํ อิตี}ติ นนฺทิสํโยชนํ วุจฺจติ อรูปราโค, อรูปราเคน ตํ กมฺมํ ลคฺคํ ลคฺคิตํ ปลิพุทฺธํ อรูปราคํ นนฺทิสํโยชนนฺติ ญ ตฺวา, ลคฺคนนฺติ ญ ตฺวา, พนฺธนนฺติ ญ ตฺวา, ปลิโพโธติ ญ ตฺวา ชานิตฺวา ตุลยิตฺวา ตีรยิตฺวา วิภาวยิตฺวา วิภูตํ กตฺวา.\csromanspacer{} \textbf{อิตี}ติ ปทสนฺธิ ปทสํสคฺโค ปทปาริปูรี อกฺขรสมวาโย พฺยญฺชนสิลิฏฺฐตา ปทานุปุพฺพตาเปตํ อิตีติ นนฺทิสํโยชนํ อิติ.}

\prose{\textbf{เอวเมตํ อภิญฺญายา}ติ เอวํ เอตํ อภิญฺญาย ชานิตฺวา ตุลยิตฺวา ตีรยิตฺวา วิภาวยิตฺวา วิภูตํ กตฺวาติ เอวเมตํ อภิญฺญาย.}

\prose{\textbf{ตโต ตตฺถ วิปสฺสตี}ติ \textbf{ตตฺถา}ติ อากิญฺจญฺญายตนํ สมาปชฺชิตฺวา ตโต วุฏฺฐหิตฺวา ตตฺถ ชาเต จิตฺตเจตสิเก ธมฺเม อนิจฺจโต วิปสฺสติ, ทุกฺขโต วิปสฺสติ, โรคโต ฯเปฯ นิสฺสรณโต วิปสฺสติ ทกฺขติ โอโลเกติ นิชฺฌายติ อุปปริกฺขตีติ ตโต ตตฺถ วิปสฺสติ.}

\prose{\textbf{เอตํ ญาณํ ตถํ ตสฺสา}ติ เอตํ ญาณํ ตจฺฉํ ภูตํ ยาถาวํ อวิปรีตํ ตสฺสาติ เอตํ ญาณํ ตถํ ตสฺส.}

\prose{\csromanfolio{172}\textbf{พฺราหฺมณสฺส วุสีมโต}ติ \textbf{พฺราหฺมโณ}ติ สตฺตนฺนํ ธมฺมานํ พาหิตตฺตา \textbf{พฺราหฺมโณ} ฯเปฯ อสิโต ตาทิ ปวุจฺจเต ส พฺรหฺมาติ.\csromanspacer{} \textbf{พฺราหฺมณสฺส วุสีมโต}ติ ปุถุชฺชนกลฺยาณํ อุปาทาย สตฺต เสกฺขา อปฺปตฺตสฺส ปตฺติยา อนธิคตสฺส อธิคมาย อสจฺฉิกตสฺส สจฺฉิกิริยาย วสนฺติ สํวสนฺติ อาวสนฺติ ปริวสนฺติ\textbf{,} อรหา วุสิตวา กตกรณีโย โอหิตภาโร อนุปฺปตฺตสทตฺโถ ปริกฺขีณภวสํโยชโน สมฺมทญฺญา วิมุตฺโต\textbf{,} โส วุตฺถวาโส จิณฺณจรโณ ฯเปฯ ชาติมรณสํสาโร\textbf{,} นตฺถิ ตสฺส ปุนพฺภโวติ พฺราหฺมณสฺส วุสีมโต.\csromanspacer{} เตนาห ภควา\csromandash{}}

\csromangathameasure{อากิญฺจญฺญาสมฺภวํ ญ ตฺวา, นนฺทิสํโยชนํ อิติ. \\ เอวเมตํ อภิญฺญาย, ตโต ตตฺถ วิปสฺสติ. \\ เอตํ ญาณํ ตถํ ตสฺส, พฺราหฺมณสฺส วุสีมโตติ.}
\csromangathagroup{\csromangathastanza{อากิญฺจญฺญาสมฺภวํ ญ ตฺวา, นนฺทิสํโยชนํ อิติ. \\ เอวเมตํ อภิญฺญาย, ตโต ตตฺถ วิปสฺสติ. \\ เอตํ ญาณํ ตถํ ตสฺส, พฺราหฺมณสฺส วุสีมโตติ.}}

\prose{สห คาถาปริโยสานา ฯเปฯ “สตฺถา เม ภนฺเต ภควา\textbf{,} สาวโกหมสฺมี”ติ.}

\vspace{\tipitakaparskip}
\csromancenter{โปสาลมาณวปุจฺฉานิทฺเทโส จุทฺทสโม.}
\csromanmatikaanchor{mtk.36}
\csromantocmarkref{subsection}{15. โมฆราชมาณวปุจฺฉานิทฺเทส}{mtk.36}
\bookmark[dest={mtk.36},level=2]{15. โมฆราชมาณวปุจฺฉานิทฺเทส}
\csromanheader{2}{15. โมฆราชมาณวปุจฺฉานิทฺเทส}
\csromanhangingbegin
\hangingheaditem{85}{ทฺวาหํ สกฺกํ อปุจฺฉิสฺสํ\textbf{,} (อิจฺจายสฺมา โมฆราชา\textbf{,})}
\hangingbody{\textbf{น เม พฺยากาสิ จกฺขุมา.}}
\hangingbody{\textbf{ยาวตติยญฺจ เทวีสิ1,}}
\hangingbody{พฺยากโรตีติ เม สุตํ.}
\csromanhangingend

\prose{\textbf{ทฺวาหํ สกฺกํ อปุจฺฉิสฺสนฺ}ติ โส \textbf{พฺราหฺมโณ} ทฺวิกฺขตฺตุํ พุทฺธํ ภควนฺตํ ปญฺหํ อปุจฺฉิ\textbf{,} ตสฺส ภควา ปญฺหํ ปุฏฺโฐ น พฺยากาสิ “ตทนฺตรา\footnote{จกฺขุสมนนฺตรา (สฺยา)} อิมสฺส พฺราหฺมณสฺส อินฺทฺริยปริปาโก ภวิสฺสตี”ติ.\csromanspacer{} \textbf{สกฺกนฺ}ติ สกฺโก\textbf{,} ภควา สกฺยกุลา ปพฺพชิโตติปิ สกฺโก.\csromanspacer{} อถ วา อฑฺโฒ มหทฺธโน ธนวาติปิ สกฺโก.\csromanspacer{} ตสฺสิมานิ ธนานิ.\csromanspacer{} เสยฺยถิทํ\textbf{,} สทฺธาธนํ สีลธนํ หิริธนํ โอตฺตปฺปธนํ สุตธนํ จาคธนํ ปญฺญาธนํ สติปฏฺฐานธนํ สมฺมปฺปธานธนํ อิทฺธิปาทธนํ อินฺทฺริยธนํ พลธนํ โพชฺฌงฺคธนํ มคฺคธนํ ผลธนํ นิพฺพานธนํ.\csromanspacer{} อิเมหิ อเนกวิเธหิ ธนรตเนหิ อฑฺโฒ มหทฺธโน ธนวาติปิ สกฺโก.\csromanspacer{} อถ วา}

\vspace{\tipitakaparskip}
\csromancenter{โมฆราชมาณวปุจฺฉานิทฺเทส สกฺโก ปหุ วิสวี อลมตฺโต สูโร วีโร วิกฺกนฺโต อภีรู อจฺฉมฺภี อนุตฺราสี อปลายี ปหีนภยเภรโว วิคตโลมหํโสติปิ สกฺโก.\csromanspacer{} \textbf{ทฺวาหํ สกฺกํ อปุจฺฉิสฺสนฺ}ติ ทฺวาหํ สกฺกํ อปุจฺฉิสฺสํ อยาจิสฺสํ อชฺเฌสิสฺสํ ปสาทยิสฺสนฺติ ทฺวาหํ สกฺกํ อปุจฺฉิสฺสํ.}
\prose{\csromanfolio{173}\textbf{อิจฺจายสฺมา โมฆราชา}ติ \textbf{อิจฺจา}ติ ปทสนฺธิ ฯเปฯ.\csromanspacer{} \textbf{อายสฺมา}ติ ปิยวจนํ ฯเปฯ.\csromanspacer{} \textbf{โมฆราชา}ติ ตสฺส พฺราหฺมณสฺส นามํ ฯเปฯ อภิลาโปติ \textbf{อิจฺจา}ยสฺมา \textbf{โมฆราชา}.}

\prose{\textbf{น เม พฺยากาสิ จกฺขุมา}ติ \textbf{น เม พฺยากาสี}ติ น เม พฺยากาสิ น อาจิกฺขิ น เทเสสิ น ปญฺญเปสิ น ปฏฺฐเปสิ น วิวริ น วิภชิ น อุตฺตานี-อกาสิ น ปกาเสสิ.\csromanspacer{} \textbf{จกฺขุมา}ติ ภควา ปญฺจหิ จกฺขูหิ \textbf{จกฺขุมา}, มํสจกฺขุนาปิ \textbf{จกฺขุมา}, ทิพฺพจกฺขุนาปิ\footnote{ทิพฺเพน จกฺขุนาปิ (ก)} \textbf{จกฺขุมา}, ปญฺญาจกฺขุนาปิ \textbf{จกฺขุมา}, พุทฺธจกฺขุนาปิ \textbf{จกฺขุมา}, สมนฺตจกฺขุนาปิ \textbf{จกฺขุมา}.}

\prose{กถํ ภควา มํสจกฺขุนาปิ \textbf{จกฺขุมา}.\csromanspacer{} มํสจกฺขุมฺหิ ภควโต ปญฺจ วณฺณา สํวิชฺชนฺติ นีโล จ วณฺโณ ปีตโก จ วณฺโณ โลหิตโก จ วณฺโณ กโณฺห จ วณฺโณ โอทาโต จ วณฺโณ.\csromanspacer{} ยตฺถ จ อกฺขิโลมานิ ปติฏฺฐิตานิ, ตํ นีลํ โหติ สุนีลํ ปาสาทิกํ ทสฺสเนยฺยํ อุมาปุปฺผสมานํ\footnote{อุมฺมารปุปฺผสมานํ (สฺยา) ขุ 7. 276 ปิฏฺเฐปิ. 3. อฬาริฏฺฐกสมานํ (สฺยา)}.\csromanspacer{} ตสฺส ปรโต ปีตกํ โหติ สุปีตกํ สุวณฺณวณฺณํ ปาสาทิกํ ทสฺสเนยฺยํ กณิการปุปฺผสมานํ, อุภโต จ อกฺขิกูฏานิ ภควโต โลหิตกานิ โหนฺติ สุโลหิตกานิ ปาสาทิกานิ ทสฺสเนยฺยานิ อินฺทโคปกสมานานิ.\csromanspacer{} มชฺเฌ กณฺหํ โหติ สุกณฺหํ อลูขํ สินิทฺธํ ปาสาทิกํ ทสฺสเนยฺยํ อทฺทาริฏฺฐกสมานํ3.\csromanspacer{} ตสฺส ปรโต โอทาตํ โหติ สุ-โอทาตํ เสตํ ปณฺฑรํ ปาสาทิกํ ทสฺสเนยฺยํ โอสธิตารกสมานํ.\csromanspacer{} เตน ภควา ปากติเกน มํสจกฺขุนา อตฺตภาวปริยาปนฺเนน ปุริมสุจริตกมฺมาภินิพฺพตฺเตน สมนฺตา โยชนํ ปสฺสติ ทิวา เจว รตฺติํ จ.\csromanspacer{} ยทา หิ จตุรงฺคสมนฺนาคโต อนฺธกาโร โหติ, สูริโย จ อตฺถงฺคโต\footnote{อตฺถงฺคมิโต (สฺยา, ก) 5. อกาลเมโฆ (สฺยา, ก) ปสฺส ขุ 7. 276 ปิฏฺเฐ.} โหติ, กาฬปกฺโข จ อุโปสโถ โหติ, ติพฺโพ จ วนสณฺโฑ โหติ, มหา จ กาฬเมโฆ5 อพฺภุฏฺฐิโต โหติ.\csromanspacer{} เอวรูเป จตุรงฺคสมนฺนาคเต อนฺธกาเร \csromanfolio{174}สมนฺตา โยชนํ ปสฺสติ, นตฺถิ โส กุฏฺโฏ วา กวาโฏ วา ปากาโร วา ปพฺพโต วา คจฺโฉ วา ลตา วา อาวรณํ รูปานํ ทสฺสนาย.\csromanspacer{} เอกํ เจ ติลผลํ นิมิตฺตํ กตฺวา ติลวาเห ปกฺขิเปยฺย, ตํเยว ติลผลํ อุทฺธเรยฺย.\csromanspacer{} เอวํ ปริสุทฺธํ ภควโต ปากติกํ มํสจกฺขุ.\csromanspacer{} เอวํ ภควา มํสจกฺขุนาปิ จกฺขุมา.}

\prose{กถํ ภควา ทิพฺเพน จกฺขุนาปิ จกฺขุมา, ภควา ทิพฺเพน จกฺขุนา วิสุทฺเธน อติกฺกนฺตมานุสเกน สตฺเต ปสฺสติ จวมาเน อุปปชฺชมาเน หีเน ปณีเต สุวณฺเณ ทุพฺพณฺเณ สุคเต ทุคฺคเต, ยถากมฺมูปเค สตฺเต ปชานาติ “อิเม วต โภนฺโต สตฺตา กายทุจฺจริเตน สมนฺนาคตา วจีทุจฺจริเตน สมนฺนาคตา มโนทุจฺจริเตน สมนฺนาคตา อริยานํ อุปวาทกา มิจฺฉาทิฏฺฐิกา มิจฺฉาทิฏฺฐิกมฺมสมาทานา, เต กายสฺส เภทา ปรํ มรณา อปายํ ทุคฺคติํ วินิปาตํ นิรยํ อุปปนฺนา, อิเม วา ปน โภนฺโต สตฺตา กายสุจริเตน สมนฺนาคตา วจีสุจริเตน สมนฺนาคตา มโนสุจริเตน สมนฺนาคตา อริยานํ อนุปวาทกา สมฺมาทิฏฺฐิกา สมฺมาทิฏฺฐิกมฺมสมาทานา, เต กายสฺส เภทา ปรํ มรณา สุคติํ สคฺคํ โลกํ อุปปนฺนา”ติ อิติ ทิพฺเพน จกฺขุนา วิสุทฺเธน อติกฺกนฺตมานุสเกน สตฺเต ปสฺสติ จวมาเน อุปปชฺชมาเน หีเน ปณีเต สุวณฺเณ ทุพฺพณฺเณ สุคเต ทุคฺคเต, ยถากมฺมูปเค สตฺเต ปชานาติ.\csromanspacer{} อากงฺขมาโน จ ภควา เอกมฺปิ โลกธาตุํ ปสฺเสยฺย, เทฺวปิ โลกธาตุโย ปสฺเสยฺย, ติสฺโสปิ โลกธาตุโย ปสฺเสยฺย, จตสฺโสปิ โลกธาตุโย ปสฺเสยฺย, ปญฺจปิ โลกธาตุโย ปสฺเสยฺย, ทสปิ โลกธาตุโย ปสฺเสยฺย, วีสมฺปิ โลกธาตุโย ปสฺเสยฺย, ติํสมฺปิ โลกธาตุโย ปสฺเสยฺย, จตฺตาลีสมฺปิ โลกธาตุโย ปสฺเสยฺย, ปญฺญาสมฺปิ โลกธาตุโย ปสฺเสยฺย, สตมฺปิ โลกธาตุโย ปสฺเสยฺย, สหสฺสิมฺปิ จูฬนิกํ โลกธาตุํ ปสฺเสยฺย, ทฺวิสหสฺสิมฺปิ มชฺฌิมิกํ โลกธาตุํ ปสฺเสยฺย, ติสหสฺสิมฺปิ โลกธาตุํ ปสฺเสยฺย, มหาสหสฺสิมฺปิ\footnote{ติสหสฺสิํ มหาสหสฺสมฺปิ (ก)} โลกธาตุํ ปสฺเสยฺย, ยาวตกํ วา\footnote{ยาวตา (สี, ก)} ปน อากงฺเขยฺย, ตาวตกํ ปสฺเสยฺย.\csromanspacer{} เอวํ ปริสุทฺธํ ภควโต ทิพฺพจกฺขุ.\csromanspacer{} เอวํ ภควา ทิพฺเพน จกฺขุนาปิ จกฺขุมา.}

\csromanheader{2}{15. โมฆราชมาณวปุจฺฉานิทฺเทส}
\prose{\csromanfolio{175}กถํ ภควา ปญฺญาจกฺขุนาปิ จกฺขุมา, ภควา มหาปญฺโญ ปุถุปญฺโญ ชวนปญฺโญ หาสปญฺโญ ติกฺขปญฺโญ นิพฺเพธิกปญฺโญ ปญฺญาปเภทกุสโล ปภินฺนญาโณ อธิคตปฏิสมฺภิทปฺปตฺโต จตุเวสารชฺชปฺปตฺโต ทสพลธารี ปุริสาสโภ ปุริสสีโห ปุริสนาโค ปุริสาชญฺโญ ปุริสโธรโยฺห อนนฺตญาโณ อนนฺตเตโช อนนฺตยโส อฑฺโฒ มหทฺธโน ธนวา เนตา วิเนตา อนุเนตา ปญฺญาเปตา นิชฺฌาเปตา เปกฺเขตา ปสาเทตา.\csromanspacer{} โส หิ ภควา อนุปฺปนฺนสฺส มคฺคสฺส อุปฺปาเทตา, อสญฺชาตสฺส มคฺคสฺส สญฺชเนตา, อนกฺขาตสฺส มคฺคสฺส อกฺขาตา, มคฺคญฺญู มคฺควิทู มคฺคโกวิโท มคฺคานุคา จ ปน เอตรหิ สาวกา วิหรนฺติ ปจฺฉา สมนฺนาคตา.}

\prose{โส หิ ภควา ชานํ ชานาติ, ปสฺสํ ปสฺสติ, จกฺกุภูโต ญาณภูโต ธมฺมภูโต พฺรหฺมภูโต วตฺตา ปวตฺตา อตฺถสฺส นินฺเนตา อมตสฺส ทาตา ธมฺมสฺสามี\footnote{ธมฺมสามิ (สฺยา, ก)} ตถาคโต.\csromanspacer{} นตฺถิ ตสฺส ภควโต อญฺญาตํ อทิฏฺฐํ อวิทิตํ อสจฺฉิกตํ อผสฺสิตํ\footnote{อผุสิตํ (สฺยา, ก)} ปญฺญาย, อตีตํ อนาคตํ ปจฺจุปฺปนฺนํ\footnote{อตีตานาคตปจฺจุปฺปนฺนํ (สฺยา)} อุปาทาย สพฺเพ ธมฺมา สพฺพากาเรน พุทฺธสฺส ภควโต ญาณมุเข อาปาถํ อาคจฺฉนฺติ.\csromanspacer{} ยํ กิญฺจิ เนยฺยํ นาม อตฺถิ ชานิตพฺพํ\footnote{อตฺถิ ธมฺมํ ชานิตพฺพํ (ก)} อตฺตตฺโถ วา ปรตฺโถ วา อุภยตฺโถ วา ทิฏฺฐิธมฺมฺมิโก วา อตฺโถ สมฺปรายิโก วา อตฺโถ อุตฺตาโน วา อตฺโถ คมฺภีโร วา อตฺโถ คุโฬฺห วา อตฺโถ ปฏิจฺฉนฺโน วา อตฺโถ เนยฺโย วา อตฺโถ นีโต วา อตฺโถ อนวชฺโช วา อตฺโถ นิกฺกิเลโส วา อตฺโถ โวทาโน วา อตฺโถ ปรมตฺโถ วา อตฺโถ\footnote{ปรมตฺโถ วา อตฺโถ (ก)}, สพฺพํ ตํ อนฺโต พุทฺธญาเณ ปริวตฺตติ, สพฺพํ กายกมฺมํ พุทฺธสฺส ภควโต ญาณานุปริวตฺติ, สพฺพํ วจีกมฺมํ ญาณานุปริวตฺติ, สพฺพํ มโนกมฺมํ ญาณานุปริวตฺติ.\csromanspacer{} อตีเต พุทฺธสฺส ภควโต อปฺปฏิหตํ ญาณํ.\csromanspacer{} อนาคเต อปฺปฏิหตํ ญาณํ, ปจฺจุปฺปนฺเน อปฺปฏิหตํ ญาณํ.\csromanspacer{} ยาวตกํ เนยฺยํ ตาวตกํ ญาณํ.\csromanspacer{} ยาวตกํ ญาณํ ตาวตกํ เนยฺยํ.\csromanspacer{} เนยฺยปริยนฺติกํ ญาณํ, ญาณปริยนฺติกํ เนยฺยํ, เนยฺยํ อติกฺกมิตฺวา ญาณํ นปฺปวตฺตติ, ญาณํ อติกฺกมิตฺวา เนยฺยปโถ นตฺถิ.\csromanspacer{} อญฺญมญฺญปริยนฺตฏฺฐายิโน เต ธมฺมา, ยถา ทฺวินฺนํ สมุคฺคปฏลานํ สมฺมาผุสิตานํ เหฏฺฐิมํ สมุคฺคปฏลํ}

\prose{\csromanfolio{176}อุปริมํ นาติวตฺตติ.\csromanspacer{} อุปริมํ สมุคฺคปฏลํ เหฏฺฐิมํ นาติวตฺตติ.\csromanspacer{} อญฺญมญฺญปริยนฺตฏฺฐายิโน.\csromanspacer{} เอวเมว พุทฺธสฺส ภควโต เนยฺยญฺจ ญาณญฺจ อญฺญมญฺญปริยนฺตฏฺฐายิโน.\csromanspacer{} ยาวตกํ เนยฺยํ ตาวตกํ ญาณํ.\csromanspacer{} ยาวตกํ ญาณํ ตาวตกํ เนยฺยํ.\csromanspacer{} เนยฺยปริยนฺติกํ ญาณํ, ญาณปริยนฺติกํ เนยฺยํ.\csromanspacer{} เนยฺยํ อติกฺกมิตฺวา ญาณํ นปฺปวตฺตติ, ญาณํ อติกฺกมิตฺวา เนยฺยปโถ นตฺถิ.\csromanspacer{} อญฺญมญฺญปริยนฺตฏฺฐายิโน เต ธมฺมา.\csromanspacer{} สพฺพธมฺเมสุ พุทฺธสฺส ภควโต ญาณํ ปวตฺตติ.\csromanspacer{} สพฺเพ ธมฺมา พุทฺธสฺส ภควโต อาวชฺชนปฏิพทฺธา อากงฺขปฏิพทฺธา มนสิการปฏิพทฺธา จิตฺตุปฺปาทปฏิพทฺธา.\csromanspacer{} สพฺพสตฺเตสุ พุทฺธสฺส ภควโต ญาณํ ปวตฺตติ.\csromanspacer{} สพฺเพสํ จ สตฺตานํ ภควา อาสยํ ชานาติ, อนุสยํ ชานาติ, จริตํ ชานาติ, อธิมุตฺติํ ชานาติ, อปฺปรชกฺเข มหารชกฺเข ติกฺขินฺทฺริเย มุทินฺทฺริเย สฺวากาเร ทฺวากาเร สุวิญฺญาปเย ทุวิญฺญาปเย ภพฺพาภพฺเพ สตฺเต ชานาติ.\csromanspacer{} สเทวโก โลโก สมารโก สพฺรหฺมโก สสฺสมณพฺราหฺมณี ปชา สเทวมนุสฺสา อนฺโตพุทฺธญาเณ ปริวตฺตติ.}

\prose{ยถา เย เกจิ มจฺฉกจฺฉปา อนฺตมโส ติมิติมิงฺคลํ\footnote{ติมิติปิงฺคลํ (ก)} อุปาทาย อนฺโตมหาสมุทฺเท ปริวตฺตนฺติ.\csromanspacer{} เอวเมว สเทวโก โลโก สมารโก โลโก สพฺรหฺมโก โลโก สสฺสมณพฺราหฺมณี ปชา สเทวมนุสฺสา อนฺโตพุทฺธญาเณ ปริวตฺตติ.\csromanspacer{}\csromanspacer{}ยถา เย เกจิ ปกฺขี อนฺตมโส ครุฬํ เวนเตยฺยํ อุปาทาย อากาสสฺส ปเทเส ปริวตฺตนฺติ.\csromanspacer{} เอวเมว เยปิ เต สาริปุตฺตสมา ปญฺญาย สมนฺนาคตา, เตปิ พุทฺธญาณสฺส ปเทเส ปริวตฺตนฺติ.\csromanspacer{} พุทฺธญาณํ เทวมนุสฺสานํ ปญฺญํ ผริตฺวา อภิภวิตฺวา ติฏฺฐติ.}

\prose{เยปิ เต ขตฺติยปณฺฑิตา พฺราหฺมณปณฺฑิตา คหปติปณฺฑิตา สมณปณฺฑิตา นิปุณา กตปรปฺปวาทา วาลเวธิรูปา โวภินฺทนฺตา\footnote{เต ภินฺทนฺตา (สฺยา, ก)} มญฺเญ จรนฺติ ปญฺญาคเตน ทิฏฺฐิคตานิ, เต ปเญฺห อภิสงฺขริตฺวา อภิสงฺขริตฺวา ตถาคตํ อุปสงฺกมิตฺวา ปุจฺฉนฺติ คูฬฺหานิ จ ปฏิจฺฉนฺนานิ.\csromanspacer{} กถิตา วิสชฺชิตา จ เต ปญฺหา ภควตา\footnote{ภควโต (ก)} โหนฺติ นิทฺทิฏฺฐการณา.\csromanspacer{} อุปกฺขิตฺตกา จ เต ภควโต สมฺปชฺชนฺติ.\csromanspacer{} อถ โข ภควาว ตตฺถ อติโรจติ.\csromanspacer{} ยทิทํ ปญฺญายาติ, เอวํ ภควา ปญฺญาจกฺขุนาปิ จกฺขุมา.}

\csromanheader{2}{15. โมฆราชมาณวปุจฺฉานิทฺเทส}
\prose{\csromanfolio{177}กถํ ภควา พุทฺธจกฺขุนาปิ จกฺขุมา, ภควา พุทฺธจกฺขุนา โลกํ โวโลเกนฺโต\footnote{โอโลเกนฺโต (ก)} อทฺทส สตฺเต อปฺปรชกฺเข มหารชกฺเข ติกฺขินฺทฺริเย มุทินฺทฺริเย สฺวากาเร ทฺวากาเร สุวิญฺญาปเย ทุวิญฺญาปเย อปฺเปกจฺเจ ปรโลกวชฺชภยทสฺสาวิโน\footnote{ปร...ทสฺสาวิเน (ก)} วิหรนฺเต.\csromanspacer{} เสยฺยถาปิ นาม อุปฺปลินิยํ วา ปทุมินิยํ วา ปุณฺฑรีกินิยํ วา อปฺเปกจฺจานิ อุปฺปลานิ วา ปทุมานิ วา ปุณฺฑรีกานิ วา อุทเก ชาตานิ อุทเก สํวฑฺฒานิ อุทกานุคฺคตานิ อนฺโตนิมุคฺคโปสีนิ\footnote{อนฺโตนิมฺมุคฺคโปสีนิ (ก)}, อปฺเปกจฺจานิ อุปฺปลานิ วา ปทุมานิ วา ปุณฺฑรีกานิ วา อุทเก ชาตานิ อุทเก สํวฑฺฒานิ สโมทกํ ฐิตานิ, อปฺเปกจฺจานิ อุปฺปลานิ วา ปทุมานิ วา ปุณฺฑรีกานิ วา อุทเก ชาตานิ อุทเก สํวฑฺฒานิ อุทกา อจฺจุคฺคมฺม ติฏฺฐนฺติ อนุปลิตฺตานิ อุทเกน.\csromanspacer{} เอวเมวํ ภควา พุทฺธจกฺขุนา โลกํ โวโลเกนฺโต อทฺทส สตฺเต อปฺปรชกฺเข มหารชกฺเข ติกฺขินฺทฺริเย มุทินฺทฺริเย สฺวากาเร ทฺวากาเร สุวิญฺญาปเย ทุวิญฺญาปเย อปฺเปกจฺเจ ปรโลกวชฺชภยทสฺสาวิโน วิหรนฺเต.\csromanspacer{} ชานาติ ภควา “อยํ ปุคฺคโล ราคจริโต, อยํ โทสจริโต, อยํ โมหจริโต, อยํ วิตกฺกจริโต, อยํ สทฺธาจริโต, อยํ ญาณจริโต”ติ.\csromanspacer{} ราคจริตสฺส ภควา ปุคฺคลสฺส อสุภกถํ กเถติ.\csromanspacer{} โทสจริตสฺส ภควา ปุคฺคลสฺส เมตฺตาภาวนํ อาจิกฺขติ.\csromanspacer{} โมหจริตสฺส ภควา ปุคฺคลสฺส อุทฺเทเส ปริปุจฺฉาย กาเลน ธมฺมสฺสวเน กาเลน ธมฺมสากจฺฉาย ครุสํวาเส นิเวเสติ.\csromanspacer{} วิตกฺกจริตสฺส ภควา ปุคฺคลสฺส อานาปานสฺสติํ อาจิกฺขติ.\csromanspacer{} สทฺธาจริตสฺส ภควา ปุคฺคลสฺส ปสาทนียํ นิมิตฺตํ อาจิกฺขติ พุทฺธสุโพธิํ\footnote{พุทฺธสุพุทฺธตํ (ก)} ธมฺมสุธมฺมตํ สํฆสุปฺปฏิปตฺติํ สีลานิ จ อตฺตโน.\csromanspacer{} ญาณจริตสฺส ภควา ปุคฺคลสฺส วิปสฺสนานิมิตฺตํ อาจิกฺขติ อนิจฺจาการํ ทุกฺขาการํ อนตฺตาการํ.}

\nitthitam{“เสเล ยถา ปพฺพตมุทฺธนิฏฺฐิโต,}
\prose{ยถาปิ ปสฺเส ชนตํ สมนฺตโต.}

\csromangathameasure{ตถูปมํ ธมฺมมยํ สุเมธ, \\ ปาสาทมารุยฺห สมนฺตจกฺขุ. \\ โสกาวติณฺณํ ชนตมเปตโสโก, \\ อเวกฺขสฺสุ ชาติชราภิภูตนฺ”ติ.}
\csromangathagroup{\csromangathastanza{ตถูปมํ ธมฺมมยํ สุเมธ, \\ ปาสาทมารุยฺห สมนฺตจกฺขุ. \\ โสกาวติณฺณํ\footnote{โสกาวกิณฺณํ (สฺยา)} ชนตมเปตโสโก, \\ อเวกฺขสฺสุ ชาติชราภิภูตนฺ”ติ.}}

\prose{\csromanfolio{178}เอวํ ภควา พุทฺธจกฺขุนาปิ จกฺขุมา.}

\prose{กถํ ภควา \textbf{สมนฺตจกฺขุ}นาปิ จกฺขุมา, \textbf{สมนฺตจกฺขุ} วุจฺจติ สพฺพญฺญุตญาณํ, ภควา สพฺพญฺญุตญาเณน อุเปโต สมุเปโต อุปาคโต สมุปาคโต อุปปนฺโน สมุปปนฺโน สมนฺนาคโต.}

\csromangathameasure{“น ตสฺส อทฺทิฏฺฐมิธตฺถิ กิญฺจิ, \\ อโถ อวิญฺญาตมชานิตพฺพํ. \\ สพฺพํ อภิญฺญาสิ ยทตฺถิ เนยฺยํ, \\ ตถาคโต เตน สมนฺตจกฺขู”ติ. เอวํ ภควา \textbf{สมนฺตจกฺขุ}นาปิ จกฺขุมาติ น เม พฺยากาสิ จกฺขุมา.}
\csromangathagroup{\csromangathastanza{“น ตสฺส อทฺทิฏฺฐมิธตฺถิ กิญฺจิ, \\ อโถ อวิญฺญาตมชานิตพฺพํ. \\ สพฺพํ อภิญฺญาสิ ยทตฺถิ เนยฺยํ, \\ ตถาคโต เตน สมนฺตจกฺขู”ติ.\csromanspacer{} เอวํ ภควา \textbf{สมนฺตจกฺขุ}นาปิ จกฺขุมาติ น เม พฺยากาสิ จกฺขุมา.}}

\prose{\textbf{ยาวตติยญฺจ เทวีสิ, พฺยากโรตีติ เม สุตนฺ}ติ ยาวตติยํ พุทฺโธ สหธมฺมิกํ ปญฺหํ ปุฏฺโฐ พฺยากโรติ โน สํสาเรตีติ\footnote{สมฺปายตีติ (สฺยา)} เอวํ มยา อุคฺคหิตํ, เอวํ มยา อุปธาริตํ, เอวํ มยา อุปลกฺขิตํ.\csromanspacer{} \textbf{เทวีสี}ติ ภควา เจว อิสิ จาติ เทวีสิ, ยถา ราชา ปพฺพชิตา วุจฺจนฺติ ราชิสโย, พฺราหฺมณา ปพฺพชิตา วุจฺจนฺติ พฺราหฺมณิสโย, เอวเมว ภควา เจว อิสิ จาติ เทวีสิ.}

\prose{อถ วา ภควา ปพฺพชิโตติปิ \textbf{อิสิ,} มหนฺตํ สีลกฺขนฺธํ เอสี คเวสี ปริเยสีติปิ \textbf{อิสิ,} มหนฺตํ สมาธิกฺขนฺธํ ฯเปฯ มหนฺตํ ปญฺญากฺขนฺธํ.\csromanspacer{} มหนฺตํ วิมุตฺติกฺขนฺธํ.\csromanspacer{} มหนฺตํ วิมุตฺติญาณทสฺสนกฺขนฺธํ เอสี คเวสี ปริเยสีติปิ \textbf{อิสิ,} มหโต ตโมกายสฺส ปทาลนํ เอสี คเวสี ปริเยสีติปิ \textbf{อิสิ,} มหโต วิปลฺลาสสฺส ปเภทนํ เอสี คเวสี ปริเยสีติปิ \textbf{อิสิ,} มหโต ตณฺหาสลฺลสฺส อพฺพหนํ.\csromanspacer{} มหโต ทิฏฺฐิสงฺฆาฏสฺส วินิเวฐนํ.\csromanspacer{} มหโต มานทฺธชสฺส ปปาตนํ.\csromanspacer{} มหโต อภิสงฺขารสฺส วูปสมํ.\csromanspacer{} มหโต โอฆสฺส นิตฺถรณํ.\csromanspacer{} มหโต ภารสฺส นิกฺเขปนํ.\csromanspacer{} มหโต สํสารวฏฺฏสฺส อุปจฺเฉทํ.\csromanspacer{} มหโต สนฺตาปสฺส นิพฺพาปนํ.\csromanspacer{} มหโต ปริฬาหสฺส ปฏิปฺปสฺสทฺธิํ.\csromanspacer{} มหโต ธมฺมทฺธชสฺส อุสฺสาปนํ เอสี คเวสี ปริเยสีติปิ \textbf{อิสิ,} มหนฺเต สติปฏฺฐาเน.\csromanspacer{} มหนฺเต สมฺมปฺปธาเน.\csromanspacer{} มหนฺตานิ อินฺทฺริยานิ.\csromanspacer{} มหนฺตานิ พลานิ.\csromanspacer{} มหนฺเต โพชฺฌงฺเค.\csromanspacer{} มหนฺตํ อริยํ อฏฺฐงฺคิกํ มคฺคํ.\csromanspacer{} มหนฺตํ ปรมตฺถํ อมตํ นิพฺพานํ เอสี คเวสี ปริเยสีติปิ \textbf{อิสิ,} มเหสกฺเขหิ วา สตฺเตหิ เอสิโต คเวสิโต ปริเยสิโต}

\csromanheader{2}{15. โมฆราชมาณวปุจฺฉานิทฺเทส}
\prose{\csromanfolio{179}“กหํ พุทฺโธ, กหํ ภควา, กหํ เทวเทโว, กหํ นราสโภ”ติปิ อิสีติ ยาวตติยญฺจ เทวีสิ, พฺยากโรตีติ เม สุตํ.\csromanspacer{} เตนาห โส พฺราหฺมโณ\csromandash{}}

\prose{ทฺวาหํ สกฺกํ อปุจฺฉิสฺสํ, (อิจฺจายสฺมา โมฆราชา,)}

\prose{น เม พฺยากาสิ จกฺขุมา.}

\csromangathameasure{ยาวตติยญฺจ เทวีสิ, \\ พฺยากโรตีติ เม สุตนฺติ.}
\csromangathagroup{\csromangathastanza{ยาวตติยญฺจ เทวีสิ, \\ พฺยากโรตีติ เม สุตนฺติ.}}

\csromanhangingbegin
\hangingheaditem{86}{\textbf{อยํ โลโก ปโร โลโก},}
\hangingbody{\textbf{พฺรหฺมโลโก สเทวโก.}}
\hangingbody{\textbf{ทิฏฺฐิํ เต นาภิชานาติ,}}
\hangingbody{\textbf{โคตมสฺส ยสสฺสิโน}.}
\csromanhangingend

\prose{\textbf{อยํ โลโก ปโร โลโก}ติ \textbf{อยํ โลโก}ติ มนุสฺสโลโก.\csromanspacer{} \textbf{ปโร โลโก}ติ มนุสฺสโลกํ ฐเปตฺวา สพฺโพ ปโร โลโกติ \textbf{อยํ โลโก} ปโร โลโก.}

\prose{\textbf{พฺรหฺมโลโก สเทวโก}ติ สเทวโก โลโก สมารโก สพฺรหฺมโก สสฺสมณพฺราหฺมณี ปชา สเทวมนุสฺสาติ พฺรหฺมโลโก สเทวโก.}

\prose{\textbf{ทิฏฺฐิํ เต นาภิชานาตี}ติ ตุยฺหํ ทิฏฺฐิํ ขนฺติํ รุจิํ ลทฺธิํ อชฺฌาสยํ อธิปฺปายํ โลโก น ชานาติ, “อยํ เอวํทิฏฺฐิโก เอวํขนฺติโก เอวํรุจิโก เอวํลทฺธิโก เอวํอชฺฌาสโย เอวํอธิปฺปาโย”ติ น ชานาติ น ปสฺสติ น ทกฺขติ นาธิคจฺฉติ น วินฺทติ น ปฏิลภตีติ ทิฏฺฐิํ เต นาภิชานาติ.}

\prose{\textbf{โคตมสฺส ยสสฺสิโน}ติ ภควา ยสปฺปตฺโตติ ยสสฺสี.\csromanspacer{} อถ วา ภควา สกฺกโต ครุกโต มานิโต ปูชิโต อปจิโต ลาภี จีวรปิณฺฑปาตเสนาสนคิลาน- ปจฺจยเภสชฺชปริกฺขารานนฺติปิ ยสสฺสีติ \textbf{โคตมสฺส ยสสฺสิโน}.\csromanspacer{} เตนาห โส พฺราหฺมโณ\csromandash{}}

\prose{\textbf{อยํ โลโก ปโร โลโก},}

\prose{\textbf{พฺรหฺมโลโก สเทวโก.}}

\csromangathameasure{\textbf{ทิฏฺฐิํ เต นาภิชานาติ,} \\ \textbf{โคตมสฺส ยสสฺสิโน}ติ.}
\csromangathagroup{\csromangathastanza{\textbf{ทิฏฺฐิํ เต นาภิชานาติ,} \\ \textbf{โคตมสฺส ยสสฺสิโน}ติ.}}

\csromanhangingbegin
\hangingheaditem{87}{\csromanfolio{180}เอวํ อภิกฺกนฺต ทสฺสาวิํ, อตฺถิ ปเญฺหน อาคมํ.}
\hangingbody{กถํ โลกํ อเวกฺขนฺตํ, มจฺจุราชา น ปสฺสติ.}
\csromanhangingend

\prose{\textbf{เอวํ อภิกฺกนฺตทสฺสาวินฺ}ติ เอวํ อภิกฺกนฺตทสฺสาวิํ อคฺคทสฺสาวิํ เสฏฺฐทสฺสาวิํ วิเสฏฺฐทสฺสาวิํ ปาโมกฺขทสฺสาวิํ อุตฺตมทสฺสาวิํ ปรมทสฺสาวินฺติ เอวํ อภิกฺกนฺตทสฺสาวิํ.}

\prose{\textbf{อตฺถิ ปเญฺหน อาคมนฺ}ติ ปเญฺหน อตฺถิโก อาคโตมฺหิ ฯเปฯ วหสฺเสตํ ภารนฺติ เอวมฺปิ อตฺถิ ปเญฺหน อาคมํ.}

\prose{\textbf{กถํ โลกํ อเวกฺขนฺตนฺ}ติ กถํ โลกํ อเวกฺขนฺตํ ปจฺจเวกฺขนฺตํ ตุลยนฺตํ ตีรยนฺตํ วิภาวยนฺตํ วิภูตํ กโรนฺตนฺติ กถํ โลกํ อเวกฺขนฺตํ.}

\prose{\textbf{มจฺจุราชา น ปสฺสตี}ติ มจฺจุราชา น ปสฺสติ น ทกฺขติ นาธิคจฺฉติ น วินฺทติ น ปฏิลภตีติ มจฺจุราชา น ปสฺสติ.\csromanspacer{} เตนาห โส พฺราหฺมโณ\csromandash{}}

\csromangathameasure{เอวํ อภิกฺกนฺตทสฺสาวิํ, อตฺถิ ปเญฺหน อาคมํ. \\ กถํ โลกํ อเวกฺขนฺตํ, มจฺจุราชา น ปสฺสตีติ.}
\csromangathagroup{\csromangathastanza{เอวํ อภิกฺกนฺตทสฺสาวิํ, อตฺถิ ปเญฺหน อาคมํ. \\ กถํ โลกํ อเวกฺขนฺตํ, มจฺจุราชา น ปสฺสตีติ.}}

\csromanhangingbegin
\hangingheaditem{88}{สุญฺญโต โลกํ อเวกฺขสฺสุ, โมฆราช สทา สโต.}
\hangingbody{\textbf{อตฺตานุทิฏฺฐิํ อูหจฺจ, เอวํ มจฺจุตโร สิยา.}}
\hangingbody{เอวํ โลกํ อเวกฺขนฺตํ, มจฺจุราชา น ปสฺสติ.}
\csromanhangingend

\prose{\textbf{สุญฺญโต โลกํ อเวกฺขสฺสู}ติ \textbf{โลโก}ติ นิรย\textbf{โลโก} ติรจฺฉาน\textbf{โลโก} เปตฺติวิสย\textbf{โลโก} มนุสฺส\textbf{โลโก} เทว\textbf{โลโก} ขนฺธ\textbf{โลโก} ธาตุ\textbf{โลโก} อายตน\textbf{โลโก} อยํ \textbf{โลโก} ปโร \textbf{โลโก} พฺรหฺม\textbf{โลโก} สเทวโก\footnote{สเทวโก โลโก (ก)}.\csromanspacer{} อญฺญตโร ภิกฺขุ ภควนฺตํ เอตทโวจ “\textbf{โลโก} \textbf{โลโก}”ติ ภนฺเต วุจฺจติ, กิตฺตาวตา นุ โข ภนฺเต “\textbf{โลโก}”ติ วุจฺจตีติ.\csromanspacer{} ลุชฺชตีติ โข ภิกฺขุ ตสฺมา “\textbf{โลโก}”ติ วุจฺจติ.\csromanspacer{} กิญฺจ ลุชฺชติ.\csromanspacer{} จกฺขุ โข ภิกฺขุ ลุชฺชติ, รูปา ลุชฺชนฺติ, จกฺขุวิญฺญาณํ ลุชฺชติ, จกฺขุสมฺผสฺโส ลุชฺชติ, ยมฺปิทํ\footnote{ยมิทํ (ก) ปสฺส สํ 2. 278 ปิฏฺเฐ.} จกฺขุสมฺผสฺสปจฺจยา อุปฺปชฺชติ เวทยิตํ สุขํ วา ทุกฺขํ วา อทุกฺขมสุขํ วา, ตมฺปิ ลุชฺชติ.\csromanspacer{} โสตํ ลุชฺชติ, คนฺธา ลุชฺชนฺติ ฯเปฯ กาโย}

\vspace{\tipitakaparskip}
\csromancenter{โมฆราชมาณวปุจฺฉานิทฺเทส ลุชฺชติ, โผฏฺฐพฺพา ลุชฺชนฺติ.\csromanspacer{} มโน ลุชฺชติ, ธมฺมา ลุชฺชนฺติ, มโนวิญฺญาณํ ลุชฺชติ, มโนสมฺผสฺโส ลุชฺชติ, ยมฺปิทํ มโนสมฺผสฺสปจฺจยา อุปฺปชฺชติ เวทยิตํ สุขํ วา ทุกฺขํ วา อทุกฺขมสุขํ วา, ตมฺปิ ลุชฺชติ.\csromanspacer{} ลุชฺชตีติ โข ภิกฺขุ ตสฺมา “โลโก”ติ วุจฺจติ.}
\prose{\csromanfolio{181}\textbf{สุญฺญโต โลกํ อเวกฺขสฺสู}ติ ทฺวีหิ การเณหิ สุญฺญโต โลกํ อเวกฺขติ อวสิยปวตฺตสลฺลกฺขณวเสน\footnote{อวสฺสิยปวตฺต... (สฺยา)} วา ตุจฺฉสงฺขารสมนุปสฺสนาวเสน วา.\csromanspacer{} กถํ อวสิยปวตฺตสลฺลกฺขณวเสน สุญฺญโต โลกํ อเวกฺขติ\csromandash{} รูเป วโส น ลพฺภติ.\csromanspacer{} เวทนาย วโส น ลพฺภติ.\csromanspacer{} สญฺญาย วโส น ลพฺภติ.\csromanspacer{} สงฺขาเรสุ วโส น ลพฺภติ.\csromanspacer{} วิญฺญาเณ วโส น ลพฺภติ.\csromanspacer{} วุตฺตเญฺหตํ ภควตา\footnote{ปสฺส สํ 2. 55 ปิฏฺเฐสุ.}\csromandash{}รูปํ ภิกฺขเว อนตฺตา, รูปญฺจ หิทํ ภิกฺขเว อตฺตา อภวิสฺส, นยิทํ รูปํ อาพาธาย สํวตฺเตยฺย, ลพฺเภถ จ รูเป “เอวํ เม รูปํ โหตุ, เอวํ เม รูปํ มา อโหสี”ติ.\csromanspacer{} ยสฺมา จ โข ภิกฺขเว รูปํ อนตฺตา, ตสฺมา รูปํ อาพาธาย สํวตฺตติ.\csromanspacer{} น จ ลพฺภติ รูเป “เอวํ เม รูปํ โหตุ, เอวํ เม รูปํ มา อโหสี”ติ.}

\prose{เวทนา อนตฺตา, เวทนา จ หิทํ ภิกฺขเว อตฺตา อภวิสฺส, นยิทํ เวทนา อาพาธาย สํวตฺเตยฺย, ลพฺเภถ จ เวทนาย “เอวํ เม เวทนา โหตุ, เอวํ เม เวทนา มา อโหสี”ติ.\csromanspacer{} ยสฺมา จ โข ภิกฺขเว เวทนา อนตฺตา, ตสฺมา เวทนา อาพาธาย สํวตฺตติ, น จ ลพฺพติ เวทนาย “เอวํ เม เวทนา โหตุ, เอวํ เม เวทนา มา อโหสี”ติ.}

\prose{สญฺญา อนตฺตา, สญฺญา จ หิทํ ภิกฺขเว อตฺตา อภวิสฺส, นยิทํ สญฺญา อาพาธาย สํวตฺเตยฺย, ลพฺเภถ จ สญฺญาย “เอวํ เม สญฺญา โหตุ, เอวํ เม สญฺญา มา อโหสี”ติ.\csromanspacer{} ยสฺมา จ โข ภิกฺขเว สญฺญา อนตฺตา, ตสฺมา สญฺญา อาพาธาย สํวตฺตติ, น จ ลพฺภติ สญฺญาย “เอวํ เม สญฺญา โหตุ, เอวํ เม สญฺญา มา อโหสี”ติ.}

\prose{สงฺขารา อนตฺตา, สงฺขารา จ หิทํ ภิกฺขเว อตฺตา อภวิสฺสํสุ, นยิทํ สงฺขารา อาพาธาย สํวตฺเตยฺยุํ, ลพฺเภถ จ สงฺขาเรสุ “เอวํ เม สงฺขารา โหนฺตุ, เอวํ เม สงฺขารา มา อเหสุนฺ”ติ.\csromanspacer{} ยสฺมา จ โข \csromanfolio{182}ภิกฺขเว สงฺขารา อนตฺตา, ตสฺมา สงฺขารา อาพาธาย สํวตฺตนฺติ, น จ ลพฺภติ สงฺขาเรสุ “เอวํ เม สงฺขารา โหนฺตุ, เอวํ เม สงฺขารา มา อเหสุนฺ”ติ.}

\prose{วิญฺญาณํ อนตฺตา, วิญฺญาณญฺจ หิทํ ภิกฺขเว อตฺตา อภวิสฺส, นยิทํ วิญฺญาณํ อาพาธาย สํวตฺเตยฺย, ลพฺเภถ จ วิญฺญาเณ “เอวํ เม วิญฺญาณํ โหตุ, เอวํ เม วิญฺญาณํ มา อโหสี”ติ.\csromanspacer{} ยสฺมา จ โข ภิกฺขเว วิญฺญาณํ อนตฺตา, ตสฺมา วิญฺญาณํ อาพาธาย สํวตฺตติ, น จ ลพฺภติ วิญฺญาเณ “เอวํ เม วิญฺญาณํ โหตุ, เอวํ เม วิญฺญาณํ มา อโหสี”ติ.}

\prose{วุตฺตเญฺหตํ ภควตา\csromandash{}นายํ ภิกฺขเว กาโย ตุมฺหากํ นปิ อญฺเญสํ\footnote{ปเรสํ (สฺยา) ปสฺส สํ 1. 294 ปิฏฺเฐ.}, ปุราณมิทํ ภิกฺขเว กมฺมํ อภิสงฺขตํ อภิสญฺเจตยิตํ เวทนิยํ ทฏฺฐพฺพํ.\csromanspacer{} ตตฺร โข ภิกฺขเว สุตวา อริยสาวโก ปฏิจฺจสมุปฺปาทํเยว สาธุกํ โยนิโส มนสิ กโรติ “อิติ อิมสฺมิํ สติ อิทํ โหติ, อิมสฺสุปฺปาทา อิทํ อุปฺปชฺชติ.\csromanspacer{} อิมสฺมิํ อสติ อิทํ น โหติ, อิมสฺส นิโรธา อิทํ นิรุชฺฌติ.\csromanspacer{} ยทิทํ อวิชฺชาปจฺจยา สงฺขารา, สงฺขารปจฺจยา วิญฺญาณํ, วิญฺญาณปจฺจยา นามรูปํ, นามรูปปจฺจยา สฬายตนํ, สฬายตนปจฺจยา ผสฺโส, ผสฺสปจฺจยา เวทนา, เวทนาปจฺจยา ตณฺหา, ตณฺหาปจฺจยา อุปาทานํ, อุปาทานปจฺจยา ภโว, ภวปจฺจยา ชาติ, ชาติปจฺจยา ชรามรณํ โสกปริเทวทุกฺข- โทมนสฺสุปายาสา สมฺภวนฺติ.\csromanspacer{} เอวเมตสฺส เกวลสฺส ทุกฺขกฺขนฺธสฺส สมุทโย โหติ.}

\prose{อวิชฺชาย เตฺวว อเสสวิราคนิโรธา สงฺขารนิโรโธ, สงฺขารนิโรธา วิญฺญาณนิโรโธ ฯเปฯ ชาตินิโรธา ชรามรณํ โสกปริเทวทุกฺข- โทมนสฺสุปายาสา นิรุชฺฌนฺติ.\csromanspacer{} เอวเมตสฺส เกวลสฺส ทุกฺขกฺขนฺธสฺส นิโรโธ โหตี”ติ.\csromanspacer{} เอวํ อวสิยปวตฺตสลฺลกฺขณวเสน สุญฺญโต โลกํ อเวกฺขติ.}

\prose{กถํ ตุจฺฉสงฺขารสมนุปสฺสนาวเสน สุญฺญโต โลกํ อเวกฺขติ\csromandash{}รูเป สาโร น ลพฺภติ.\csromanspacer{} เวทนาย สาโร น ลพฺภติ.\csromanspacer{} สญฺญาย สาโร น ลพฺภติ.\csromanspacer{} สงฺขาเรสุ สาโร น ลพฺภติ.\csromanspacer{} วิญฺญาเณ สาโร น ลพฺภติ.\csromanspacer{} รูปํ อสฺสารํ นิสฺสารํ สาราปคตํ นิจฺจสารสาเรน วา สุขสารสาเรน วา อตฺตสารสาเรน วา นิจฺเจน วา ธุเวน วา สสฺสเตน วา}

\vspace{\tipitakaparskip}
\csromancenter{โมฆราชมาณวปุจฺฉานิทฺเทส อวิปริณามธมฺเมน วา.\csromanspacer{} เวทนา อสฺสารา นิสฺสารา สาราปคตา.\csromanspacer{} สญฺญา อสฺสารา นิสฺสารา สาราปคตา.\csromanspacer{} สงฺขารา อสฺสารา นิสฺสารา สาราปคตา.\csromanspacer{} วิญฺญาณํ อสฺสารํ นิสฺสารํ สาราปคตํ นิจฺจสารสาเรน วา สุขสารสาเรน วา อตฺตสารสาเรน วา นิจฺเจน วา ธุเวน วา สสฺสเตน วา อวิปริณามธมฺเมน วา.\csromanspacer{} ยถา นโฬ อสฺสาโร นิสฺสาโร สาราปคโต.\csromanspacer{} ยถา จ เอรณฺโฑ ฯเปฯ.\csromanspacer{} ยถา จ อุทุมฺพโร อสฺสาโร นิสฺสาโร สาราปคโต.\csromanspacer{} ยถา จ เสตกจฺโฉ\footnote{เสตวจฺโฉ (ก)} อสฺสาโร นิสฺสาโร สาราปคโต.\csromanspacer{} ยถา จ ปาลิภทฺทโก\footnote{ปาฬิภทฺทโก (ก)} อสฺสาโร นิสฺสาโร สาราปคโต.\csromanspacer{} ยถา จ เผณปิณฺโฑ\footnote{เผณุปิณฺโฑ (สฺยา)} อสฺสาโร นิสฺสาโร สาราปคโต.\csromanspacer{} ยถา จ อุทกปุพฺพุฬํ\footnote{ปุพฺพุลกํ (สฺยา)} อสฺสารํ นิสฺสารํ สาราปคตํ.\csromanspacer{} ยถา จ มรีจิ อสฺสารา นิสฺสารา สาราปคตา.\csromanspacer{} ยถา กทลิกฺขนฺโธ อสฺสาโร นิสฺสาโร สาราปคโต.\csromanspacer{} ยถา มายา อสฺสารา นิสฺสารา สาราปคตา.\csromanspacer{} เอวเมว รูปํ อสฺสารํ นิสฺสารํ สาราปคตํ นิจฺจสารสาเรน วา สุขสารสาเรน วา อตฺตสารสาเรน วา นิจฺเจน วา ธุเวน วา สสฺสเตน วา อวิปริณามธมฺเมน วา.\csromanspacer{} เวทนา อสฺสารา นิสฺสารา สาราปคตา.\csromanspacer{} สญฺญา อสฺสารา นิสฺสารา สาราปคตา.\csromanspacer{} สงฺขารา อสฺสารา นิสฺสารา สาราปคตา.\csromanspacer{} วิญฺญาณํ อสฺสารํ นิสฺสารํ สาราปคตํ นิจฺจสารสาเรน วา สุขสารสาเรน วา อตฺตสารสาเรน วา นิจฺเจน วา ธุเวน วา สสฺสเตน วา อวิปริณามธมฺเมน วา.\csromanspacer{} เอวํ ตุจฺฉสงฺขารสมนุปสฺสนาวเสน สุญฺญโต โลกํ อเวกฺขติ.\csromanspacer{} อิเมหิ ทฺวีหิ การเณหิ สุญฺญโต โลกํ อเวกฺขติ.}
\prose{\csromanfolio{183}อปิ จ ฉหากาเรหิ สุญฺญโต โลกํ อเวกฺขติ.\csromanspacer{} จกฺขุ สุญฺญํ\footnote{สฺยา...โปตฺถเก อิมสฺมิํ ฐาเน อญฺญถา ทิสฺสติ.} อตฺเตน วา อตฺตนิเยน วา นิจฺเจน วา ธุเวน วา สสฺสเตน วา อวิปริณามธมฺเมน วา.\csromanspacer{} โสตํ สุญฺญํ.\csromanspacer{} ฆานํ สุญฺญํ.\csromanspacer{} ชิวฺหา สุญฺญา.\csromanspacer{} กาโย สุญฺโญ.\csromanspacer{} มโน สุญฺโญ อตฺเตน วา อตฺตนิเยน วา นิจฺเจน วา ธุเวน วา สสฺสเตน วา อวิปริณามธมฺเมน วา.\csromanspacer{} รูปา สุญฺญา.\csromanspacer{} สทฺทา สุญฺญา.\csromanspacer{} คนฺธา สุญฺญา.\csromanspacer{} รสา สุญฺญา.\csromanspacer{} โผฏฺฐพฺพา สุญฺญา.\csromanspacer{} ธมฺมา สุญฺญา อตฺเตน วา อตฺตนิเยน วา นิจฺเจน วา ธุเวน วา สสฺสเตน วา อวิปริณามธมฺเมน วา.\csromanspacer{} จกฺขุวิญฺญาณํ สุญฺญํ ฯเปฯ.\csromanspacer{} มโนวิญฺญาณํ สุญฺญํ.\csromanspacer{} จกฺขุสมฺผสฺโส \csromanfolio{184}สุญฺโญ ฯเปฯ.\csromanspacer{} มโนสมฺผสฺโส สุญฺโญ.\csromanspacer{} จกฺขุสมฺผสฺสชา เวทนา สุญฺญา ฯเปฯ.\csromanspacer{} มโนสมฺผสฺสชา เวทนา สุญฺญา.\csromanspacer{} รูปสญฺญา สุญฺญา ฯเปฯ.\csromanspacer{} ธมฺมสญฺญา สุญฺญา.\csromanspacer{} รูปสญฺเจตนา สุญฺญา ฯเปฯ.\csromanspacer{} ธมฺมสญฺเจตนา สุญฺญา.\csromanspacer{} รูปตณฺหา สุญฺญา.\csromanspacer{} รูปวิตกฺโก สุญฺโญ.\csromanspacer{} รูปวิจาโร สุญฺโญ ฯเปฯ.\csromanspacer{} ธมฺมวิจาโร สุญฺโญ อตฺเตน วา อตฺตนิเยน วา นิจฺเจน วา ธุเวน วา สสฺสเตน วา อวิปริณามธมฺเมน วา.\csromanspacer{} เอวํ ฉหากาเรหิ สุญฺญโต โลกํ อเวกฺขติ.}

\prose{อปิ จ ทสหากาเรหิ สุญฺญโต โลกํ อเวกฺขติ, รูปํ ริตฺตโต ตุจฺฉโต สุญฺญโต อนตฺตโต อสารกโต วธกโต วิภวโต อฆมูลโต สาสวโต สงฺขตโต.\csromanspacer{} เวทนํ.\csromanspacer{} สญฺญํ.\csromanspacer{} สงฺขาเร.\csromanspacer{} วิญฺญาณํ.\csromanspacer{} จุติํ.\csromanspacer{} อุปปตฺติํ.\csromanspacer{} ปฏิสนฺธิํ.\csromanspacer{} ภวํ.\csromanspacer{} สํสารวฏฺฏํ ริตฺตโต ตุจฺฉโต สุญฺญโต อนตฺตโต อสารกโต วธกโต วิภวโต อฆมูลโต สาสวโต สงฺขตโต.\csromanspacer{} เอวํ ทสหากาเรหิ สุญฺญโต โลกํ อเวกฺขติ.}

\prose{อปิ จ ทฺวาทสหากาเรหิ สุญฺญโต โลกํ อเวกฺขติ.\csromanspacer{} รูปํ น สตฺโต, น ชีโว, น นโร, น มานโว, น อิตฺถี, น ปุริโส, น อตฺตา, น อตฺตนิยํ, นาหํ, น มม, น โกจิ, น กสฺสจิ.\csromanspacer{} เวทนา.\csromanspacer{} สญฺญา.\csromanspacer{} สงฺขารา.\csromanspacer{} วิญฺญาณํ น สตฺโต, น ชีโว, น นโร, น มานโว, น อิตฺถี, น ปุริโส, น อตฺตา, น อตฺตนิยํ, นาหํ, น มม, น โกจิ, น กสฺสจิ.\csromanspacer{} เอวํ ทฺวาทสหากาเรหิ สุญฺญโต โลกํ อเวกฺขติ.}

\prose{วุตฺตเญฺหตํ ภควตา\csromandash{}“ยํ ภิกฺขเว น ตุมฺหากํ, ตํ ปชหถ, ตํ โว ปหีนํ ทีฆรตฺตํ หิตาย สุขาย ภวิสฺสติ.\csromanspacer{} กิญฺจ ภิกฺขเว น ตุมฺหากํ.\csromanspacer{} รูปํ ภิกฺขเว น ตุมฺหากํ, ตํ ปชหถ.\csromanspacer{} ตํ โว ปหีนํ ทีฆรตฺตํ หิตาย สุขาย ภวิสฺสติ.\csromanspacer{} เวทนา ภิกฺขเว น ตุมฺหากํ, ตํ ปชหถ, สา โว ปหีนา ทีฆรตฺตํ หิตาย สุขาย ภวิสฺสติ.\csromanspacer{} สญฺญา ภิกฺขเว น ตุมฺหากํ, ตํ ปชหถ, สา โว ปหีนา ทีฆรตฺตํ หิตาย สุขาย ภวิสฺสติ.\csromanspacer{} สงฺขารา ภิกฺขเว น ตุมฺหากํ, เต ปชหถ, เต โว ปหีนา ทีฆรตฺตํ หิตาย สุขาย ภวิสฺสนฺติ.\csromanspacer{} วิญฺญาณํ ภิกฺขเว น ตุมฺหากํ, ตํ ปชหถ, ตํ โว ปหีนํ ทีฆรตฺตํ หิตาย สุขาย ภวิสฺสติ.}

\vspace{\tipitakaparskip}
\csromancenter{โมฆราชมาณวปุจฺฉานิทฺเทส เสยฺยถาปิ\footnote{ตํ กิํ มญฺญถ (สฺยา, ก) ปสฺส สํ 2. 28 ปิฏฺเฐ.} ภิกฺขเว ยํ อิมสฺมิํ เชตวเน ติณกฏฺฐสาขาปลาสํ, ตํ ชโน หเรยฺย วา ฑเหยฺย วา ยถาปจฺจยํ วา กเรยฺย, อปิ นุ ตุมฺหากํ เอวมสฺส ‘อเมฺห ชโน ตรติ วา ฑหติ วา ยถาปจฺจยํ วา กโรตี’ติ.\csromanspacer{} โน เหตํ ภนฺเต.\csromanspacer{} ตํ กิสฺส เหตุ, น หิ โน เอตํ ภนฺเต อตฺตา วา อตฺตนิยํ วาติ.\csromanspacer{} เอวเมว โข ภิกฺขเว ยํ น ตุมฺหากํ, ตํ ปชหถ, ตํ โว ปหีนํ ทีฆรตฺตํ หิตาย สุขาย ภวิสฺสติ.\csromanspacer{} กิญฺจ ภิกฺขเว น ตุมฺหากํ.\csromanspacer{} รูปํ ภิกฺขเว น ตุมฺหากํ, ตํ ปชหถ, ตํ โว ปหีนํ ทีฆรตฺตํ หิตาย สุขาย ภวิสฺสติ.\csromanspacer{} เวทนา.\csromanspacer{} สญฺญา.\csromanspacer{} สงฺขารา.\csromanspacer{} วิญฺญาณํ ภิกฺขเว น ตุมฺหากํ, ตํ ปชหถ, ตํ โว ปหีนํ ทีฆรตฺตํ หิตาย สุขาย ภวิสฺสตี”ติ.\csromanspacer{} เอวมฺปิ สุญฺญโต โลกํ อเวกฺขติ.}
\prose{\csromanfolio{185}อายสฺมา อานนฺโท ภควนฺตํ เอตทโวจ “สุญฺโญ\footnote{สุญฺญโต (ก) ปสฺส สํ 2. 280 ปิฏฺเฐ.} โลโก สุญฺโญ โลโก”ติ ภนฺเต วุจฺจติ, กิตฺตาวตา นุ โข ภนฺเต “สุญฺโญ โลโก”ติ วุจฺจตีติ.\csromanspacer{} ยสฺมา จ โข อานนฺท สุญฺญํ อตฺเตน วา อตฺตนิเยน วา, ตสฺมา “สุญฺโญ โลโก”ติ วุจฺจติ.\csromanspacer{} กิญฺจานนฺท สุญฺญํ อตฺเตน วา อตฺตนิเยน วา.\csromanspacer{} จกฺขุ โข อานนฺท สุญฺญํ อตฺเตน วา อตฺตนิเยน วา.\csromanspacer{} รูปา สุญฺญา.\csromanspacer{} จกฺขุวิญฺญาณํ สุญฺญํ.\csromanspacer{} จกฺขุสมฺผสฺโส สุญฺโญ, ยมฺปิทํ จกฺขุสมฺผสฺสปจฺจยา อุปฺปชฺชติ เวทยิตํ สุขํ วา ทุกฺขํ วา อทุกฺขมสุขํ วา.\csromanspacer{} ตมฺปิ สุญฺญํ อตฺเตน วา อตฺตนิเยน วา.\csromanspacer{} โสตํ สุญฺญํ.\csromanspacer{} สทฺทา สุญฺญา.\csromanspacer{} ฆานํ สุญฺญํ.\csromanspacer{} คนฺธา สุญฺญา.\csromanspacer{} ชิวฺหา สุญฺญา.\csromanspacer{} รสา สุญฺญา.\csromanspacer{} กาโย สุญฺโญ.\csromanspacer{} โผฏฺฐพฺพา สุญฺญา.\csromanspacer{} มโน สุญฺโญ.\csromanspacer{} ธมฺมา สุญฺญา.\csromanspacer{} มโนวิญฺญาณํ สุญฺญํ.\csromanspacer{} มโนสมฺผสฺโส สุญฺโญ, ยมฺปิทํ สุญฺญํ มโนสมฺผสฺสปจฺจยา อุปฺปชฺชติ เวทยิตํ สุขํ วา ทุขํ วา อทุกฺขมสุขํ วา, ตมฺปิ สุญฺญํ อตฺเตน วา อตฺตนิเยน วา.\csromanspacer{} ยสฺมา จ โข อานนฺท สุญฺญํ อตฺเตน วา อตฺตนิเยน วา, ตสฺมา “สุญฺโญ โลโก”ติ วุจฺจตีติ.\csromanspacer{} เอวมฺปิ สุญฺญโต โลกํ อเวกฺขติ.}

\csromangathameasure{“สุทฺธํ ธมฺมสมุปฺปาทํ, สุทฺธสงฺขารสนฺตติํ. \\ ปสฺสนฺตสฺส ยถาภูตํ, น ภยํ โหติ คามณิ. \\ ติณกฏฺฐสมํ โลกํ, ยทา ปญฺญาย ปสฺสติ. \\ นาญฺญํ ปตฺถยเต กิญฺจิ, อญฺญตฺรปฺปฏิสนฺธิยา”ติ.}
\csromangathagroup{\csromangathastanza{“สุทฺธํ ธมฺมสมุปฺปาทํ, สุทฺธสงฺขารสนฺตติํ. \\ ปสฺสนฺตสฺส ยถาภูตํ, น ภยํ โหติ คามณิ.}\csromangathastanza{\csromanfolio{186}ติณกฏฺฐสมํ โลกํ, ยทา ปญฺญาย ปสฺสติ. \\ นาญฺญํ\footnote{น อญฺญํ (สี, สฺยา, ก)} ปตฺถยเต กิญฺจิ, อญฺญตฺรปฺปฏิสนฺธิยา”ติ.}}

\prose{เอวมฺปิ สุญฺญโต โลกํ อเวกฺขติ.}

\prose{วุตฺตเญฺหตํ ภควตา\footnote{ปสฺส สํ 2. 401 ปิฏฺเฐ.}\csromandash{}เอวเมว โข ภิกฺขเว ภิกฺขุ รูปํ สมนฺเนสติ ยาวตา รูปสฺส คติ, เวทนํ สมนฺเนสติ ยาวตา เวทนาย คติ, สญฺญํ สมนฺเนสติ ยาวตา สญฺญาย คติ, สงฺขาเร สมนฺเนสติ ยาวตา สงฺขารานํ คติ, วิญฺญาณํ สมนฺเนสติ ยาวตา วิญฺญาณสฺส คติ, ตสฺส รูปํ\footnote{ตสฺส ภิกฺขุโน รูปํ (สฺยา)} สมนฺเนสโต ยาวตา รูปสฺส คติ.\csromanspacer{} เวทนํ ฯเปฯ.\csromanspacer{} สญฺญํ.\csromanspacer{} สงฺขาเร.\csromanspacer{} วิญฺญาณํ สมนฺเนสโต ยาวตา วิญฺญาณสฺส คติ, ยมฺปิสฺส ตํ โหติ “อหนฺ”ติ วา “มมนฺ”ติ วา “อสฺมี”ติ วา, ตมฺปิ ตสฺส น โหตีติ.\csromanspacer{} เอวมฺปิ สุญฺญโต โลกํ อเวกฺขติ.}

\prose{\textbf{สุญฺญโต โลกํ อเวกฺขสฺสู}ติ สุญฺญโต โลกํ อเวกฺขสฺสุ ปจฺจเวกฺขสฺสุ ทกฺขสฺสุ ตุเลหิ ตีเรหิ วิภาเวหิ วิภูตํ กโรหีติ สุญฺญโต โลกํ อเวกฺขสฺสุ.}

\prose{\textbf{โมฆราช สทา สโต}ติ \textbf{โมฆราชา}ติ ภควา ตํ พฺราหฺมณํ นาเมน อาลปติ.\csromanspacer{} \textbf{สทา}ติ สพฺพกาลํ ฯเปฯ สจฺฉิเม วโยขนฺเธ.\csromanspacer{} \textbf{สโต}ติ จตูหิ การเณหิ สโต.\csromanspacer{} กาเย กายานุปสฺสนาสติปฏฺฐานํ ภาเวนฺโต สโต ฯเปฯ โส วุจฺจติ สโตติ \textbf{โมฆราช สทา สโต}.}

\prose{\textbf{อตฺตานุทิฏฺฐิํ อูหจฺจา}ติ อตฺตานุทิฏฺฐิ วุจฺจติ วีสติวตฺถุกา สกฺกายทิฏฺฐิ.\csromanspacer{} อิธ อสฺสุตวา ปุถุชฺชโน อริยานํ อทสฺสาวี อริยธมฺมสฺส อโกวิโท อริยธมฺเม อวินีโต, สปฺปุริสานํ อทสฺสาวี สปฺปุริสธมฺมสฺส อโกวิโท สปฺปุริสธมฺเม อวินีโต รูปํ อตฺตโต สมนุปสฺสติ, รูปวนฺตํ วา อตฺตานํ, อตฺตนิ วา รูปํ, รูปสฺมิํ วา อตฺตานํ.\csromanspacer{} เวทนํ.\csromanspacer{} สญฺญํ.\csromanspacer{} สงฺขาเร.\csromanspacer{} วิญฺญาณํ อตฺตโต สมนุปสฺสติ, วิญฺญาณวนฺตํ วา อตฺตานํ, อตฺตนิ วา วิญฺญาณํ, วิญฺญาณสฺมิํ วา อตฺตานํ.\csromanspacer{} ยา เอวรูปา ทิฏฺฐิ ทิฏฺฐิคตํ ทิฏฺฐิคหนํ ทิฏฺฐิกนฺตาโร ทิฏฺฐิวิสูกายิกํ ทิฏฺฐิวิปฺผนฺทิตํ ทิฏฺฐิสํโยชนํ คาโห ปฏิคฺคาโห อภินิเวโส ปรามาโส กุมฺมคฺโค มิจฺฉาปโถ มิจฺฉตฺตํ ติตฺถายตนํ วิปริยาสคฺคาโห วิปรีตคฺคาโห วิปลฺลาสคฺคาโห มิจฺฉาคาโห}

\vspace{\tipitakaparskip}
\csromancenter{โมฆราชมาณวปุจฺฉานิทฺเทส อยาถาวกสฺมิํ ยาถาวกนฺติ คาโห ยาวตา ทฺวาสฏฺฐิ ทิฏฺฐิคตานิ, อยํ อตฺตานุทิฏฺฐิ.\csromanspacer{} \textbf{อตฺตานุทิฏฺฐิํ อูหจฺจา}ติ อตฺตานุทิฏฺฐิํ อูหจฺจ สมูหจฺจ\footnote{อุหจฺจ สมุหจฺจ (ก) สทฺทนีติยา ปน สเมติ.} อุทฺธริตฺวา สมุทฺธริตฺวา อุปฺปาฏยิตฺวา สมุปฺปาฏยิตฺวา ปชหิตฺวา วิโนเทตฺวา พฺยนฺตีกริตฺวา อนภาวํ คเมตฺวาติ อตฺตานุทิฏฺฐิํ อูหจฺจ.}
\prose{\csromanfolio{187}\textbf{เอวํ มจฺจุตโร สิยา}ติ เอวํ มจฺจุปิ ตเรยฺยาสิ, ชราปิ ตเรยฺยาสิ, มรณมฺปิ ตเรยฺยาสิ อุตฺตเรยฺยาสิ ปตเรยฺยาสิ สมติกฺกเมยฺยาสิ วีติวตฺเตยฺยาสีติ เอวํ มจฺจุตโร สิยา.}

\prose{\textbf{เอวํ โลกํ อเวกฺขนฺตนฺ}ติ เอวํ โลกํ อเวกฺขนฺตํ ปจฺจเวกฺขนฺตํ อุลยนฺตํ ตีรยนฺตํ วิภาวยนฺตํ วิภูตํ กโรนฺตนฺติ เอวํ โลกํ อเวกฺขนฺตํ.}

\prose{\textbf{มจฺจุ ราชา น ปสฺสตี}ติ มจฺจุปิ มจฺจุราชา, มาโรปิ มจฺจุราชา, มรณมฺปิ มจฺจุราชา.\csromanspacer{} \textbf{น ปสฺสตี}ติ มจฺจุราชา น ปสฺสติ น ทกฺขติ นาธิคจฺฉติ น วินฺทติ น ปฏิลภติ.\csromanspacer{} วุตฺตเญฺหตํ ภควตา\csromandash{} “เสยฺยถาปิ ภิกฺขเว อารญฺญิโก มิโค อรญฺเญ ปวเน จรมาโน วิสฺสตฺโถ คจฺฉติ, วิสฺสตฺโถ\footnote{วิสฺสฏฺโฐ (ก)} ติฏฺฐติ, วิสฺสตฺโถ นิสีทติ, วิสฺสตฺโถ เสยฺยํ กปฺเปติ.\csromanspacer{} ตํ กิสฺส เหตุ.\csromanspacer{} อนาปาถคโต\footnote{อนาปาตคโต (ก)} ภิกฺขเว ลุทฺทสฺส.\csromanspacer{} เอวเมว โข ภิกฺขเว ภิกฺขุ วิวิจฺเจว กาเมหิ วิวิจฺจ อกุสเลหิ ธมฺเมหิ สวิตกฺกํ สวิจารํ วิเวกชํ ปีติสุขํ ปฐมํ ฌานํ อุปสมฺปชฺช วิหรติ.\csromanspacer{} อยํ วุจฺจติ ภิกฺขเว ภิกฺขุ อนฺธมกาสิ\footnote{อนฺตมกาสิ (ก) ปสฺส ม 1. 214 ปิฏฺเฐ.}, มารํ อปทํ วธิตฺวา มารจกฺขุํ อทสฺสนํ คโต ปาปิมโต.}

\prose{ปุน จปรํ ภิกฺขเว ภิกฺขุ วิตกฺกวิจารานํ วูปสมา อชฺฌตฺตํ สมฺปสาทนํ เจตโส เอโกทิภาวํ อวิตกฺกํ อวิจารํ สมาธิชํ ปีติสุขํ ทุติยํ ฌานํ ฯเปฯ ตติยํ ฌานํ.\csromanspacer{} จตุตฺถํ ฌานํ อุปสมฺปชฺช วิหรติ.\csromanspacer{} อยํ วุจฺจติ ภิกฺขเว ภิกฺขุ อนฺธมกาสิ มารํ, อปทํ วธิตฺวา มารจกฺขุํ อทสฺสนํ คโต ปาปิมโต.}

\prose{ปุน จปรํ ภิกฺขเว ภิกฺขุ สพฺพโส รูปสญฺญานํ สมติกฺกมา ปฏิฆสญฺญานํ อตฺถงฺคมา นานตฺตสญฺญานํ อมนสิการา “อนนฺโต อากาโส”ติ อากาสานญฺจายตนํ อุปสมฺปชฺช วิหรติ.\csromanspacer{} อยํ วุจฺจติ ภิกฺขเว \csromanfolio{188}ภิกฺขุ อนฺธมกาสิ มารํ\textbf{,} อปทํ วธิตฺวา มารจกฺขุํ อทสฺสนํ คโต ปาปิมโต.}

\prose{ปุน จปรํ ภิกฺขเว ภิกฺขุ สพฺพโส อากาสานญฺจายตนํ สมติกฺกมฺม “อนนฺตํ วิญฺญาณนฺ”ติ วิญฺญาณญฺจายตนํ อุปสมฺปชฺช วิหรติ ฯเปฯ.}

\prose{ปุน จปรํ ภิกฺขเว สพฺพโส วิญฺญาณญฺจายตนํ สมติกฺกมฺม “นตฺถิ กิญฺจี”ติ อากิญฺจญฺญายตนํ อุปสมฺปชฺช วิหรติ ฯเปฯ.}

\prose{ปุน จปรํ ภิกฺขเว สพฺพโส อากิญฺจญฺญายตนํ สมติกฺกมฺม เนวสญฺญานาสญฺญายตนํ อุปสมฺปชฺช วิหรติ ฯเปฯ.}

\prose{ปุน จปรํ ภิกฺขเว สพฺพโส เนวสญฺญานาสญฺญายตนํ สมติกฺกมฺม สญฺญาเวทยิตนิโรธํ อุปสมฺปชฺช วิหรติ.\csromanspacer{} ปญฺญาย จสฺส ทิสฺวา อาสวา ปริกฺขีณา โหนฺติ.\csromanspacer{} อยํ วุจฺจติ ภิกฺขเว ภิกฺขุ อนฺธมกาสิ มารํ\textbf{,} อปทํ วธิตฺวา มารจกฺขุํ อทสฺสนํ คโต ปาปิมโต\textbf{,} ติณฺโณ โลเก วิสตฺติกนฺ”ติ.\csromanspacer{} โส วิสฺสตฺโถ คจฺฉติ\textbf{,} วิสฺสตฺโถ ติฏฺฐติ\textbf{,} วิสฺสตฺโถ นิสีทติ\textbf{,} วิสฺสตฺโถ เสยฺยํ กปฺเปติ.\csromanspacer{} ตํ กิสฺส เหตุ\textbf{,} อนาปาถคโต ภิกฺขุ ปาปิมโตติ มจฺจุราชา น ปสฺสติ.\csromanspacer{} เตนาห ภควา\csromandash{}}

\csromangathameasure{สุญฺญโต โลกํ อเวกฺขสฺสุ, โมฆราช สทา สโต. \\ อตฺตานุทิฏฺฐิํ อูหจฺจ, เอวํ มจฺจุตโร สิยา. \\ เอวํ โลกํ อเวกฺขนฺตํ, มจฺจุราชา น ปสฺสตีติ.}
\csromangathagroup{\csromangathastanza{สุญฺญโต โลกํ อเวกฺขสฺสุ, โมฆราช สทา สโต. \\ อตฺตานุทิฏฺฐิํ อูหจฺจ, เอวํ มจฺจุตโร สิยา. \\ เอวํ โลกํ อเวกฺขนฺตํ, มจฺจุราชา น ปสฺสตีติ.}}

\prose{สห คาถาปริโยสานา ฯเปฯ “สตฺถา เม ภนฺเต ภควา\textbf{,} สาวโกหมสฺมี”ติ.}

\vspace{\tipitakaparskip}
\csromancenter{โมฆราชมาณวปุจฺฉานิทฺเทโส ปนฺนรสโม.}
\csromanmatikaanchor{mtk.37}
\csromantocmarkref{subsection}{16. ปิงฺคิยมาณวปุจฺฉานิทฺเทส}{mtk.37}
\bookmark[dest={mtk.37},level=2]{16. ปิงฺคิยมาณวปุจฺฉานิทฺเทส}
\csromanheader{2}{16. ปิงฺคิยมาณวปุจฺฉานิทฺเทส}
\csromanhangingbegin
\hangingheaditem{89}{ชิณฺโณหมสฺมิ อพโล วีตวณฺโณ\footnote{วิวณฺโณ (สฺยา)}\textbf{, (อิจฺจายสฺมา ปิงฺคิโย,)}}
\hangingbody{\textbf{เนตฺตา น สุทฺธา สวนํ น ผาสุ.}}
\hangingbody{\textbf{มาหํ นสฺสํ โมมุโห อนฺตราว2,}}
\hangingbody{\textbf{อาจิกฺข ธมฺมํ ยมหํ วิชญฺญํ.}}
\hangingbody{ชาติชราย อิธ วิปฺปหานํ.}
\csromanhangingend

\csromanheader{2}{16. ปิงฺคิยมาณวปุจฺฉานิทฺเทส}
\prose{\csromanfolio{189}\textbf{ชิณฺโณหมสฺมิ อพโล วีตวณฺโณ}ติ \textbf{ชิณฺโณหมสฺมี}ติ ชิณฺโณ วุฑฺโฒ มหลฺลโก อทฺธคโต วโย-อนุปฺปตฺโต วีสวสฺสสติโก ชาติโย.\csromanspacer{} \textbf{อพโล}ติ ทุพฺพโล อปฺปพโล อปฺปถาโม.\csromanspacer{} \textbf{วีตวณฺโณ}ติ วีตวณฺโณ วิคตวณฺโณ วิคจฺฉิตวณฺโณ, ยา สา ปุริมา สุภา วณฺณนิภา, สา อนฺตรหิตา อาทีนโว ปาตุภูโตติ ชิณฺโณหมสฺมิ อพโล วีตวณฺโณ.}

\prose{\textbf{อิจฺจายสฺมา ปิงฺคิโย}ติ \textbf{อิจฺจา}ติ ปทสทฺธิ ฯเปฯ.\csromanspacer{} \textbf{อายสฺมา}ติ ปิยวจนํ ฯเปฯ.\csromanspacer{} \textbf{ปิงฺคิโย}ติ ตสฺส พฺราหฺมณสฺส นามํ ฯเปฯ อภิลาโปติ \textbf{อิจฺจา}ยสฺมา \textbf{ปิงฺคิโย}.}

\prose{\textbf{เนตฺตา น สุทฺธา สวนํ น ผาสู}ติ เนตฺตา อสุทฺธา อวิสุทฺธา อปริสุทฺธา อโวทาตา, โน ตถา จกฺขุนา รูเป ปสฺสามีติ เนตฺตา น สุทฺธา.\csromanspacer{} สฺ\textbf{อวนํ น ผาสู}ติ โสตํ อสุทฺธํ อวิสุทฺธํ อปริสุทฺธํ อโวทาตํ, โน ตถา โสเตน สทฺทํ สุโณมีติ เนตฺตา น สุทฺธา สวนํ น ผาสุ.}

\prose{\textbf{มาหํ นสฺสํ โมมุโห อนฺตราวา}ติ \textbf{มาหํ นสฺสนฺ}ติ มาหํ นสฺสํ มาหํ วินสฺสํ มาหํ ปนสฺสํ.\csromanspacer{} \textbf{โมมุโห}ติ โมหมุโห อวิชฺชาคโต อญฺญาณี อวิภาวี ทุปฺปญฺโญ.\csromanspacer{} \textbf{อนฺตราวา}ติ ตุยฺหํ ธมฺมํ ทิฏฺฐิํ ปฏิปทํ มคฺคํ อนญฺญาย อนธิคนฺตฺวา อวิทิตฺวา อปฺปฏิลภิตฺวา อผสฺสยิตฺวา อสจฺฉิกริตฺวา อนฺตราเยว กาลํ กเรยฺยนฺติ มาหํ นสฺสํ โมมุโห อนฺตราว.}

\prose{\textbf{อาจิกฺข ธมฺมํ ยมหํ วิชญฺญนฺ}ติ \textbf{ธมฺมนฺ}ติ อาทิกลฺยาณํ มชฺเฌกลฺยาณํ ปริโยสานกลฺยาณํ สาตฺถํ สพฺยญฺชนํ เกวลปริปุณฺณํ ปริสุทฺธํ พฺรหฺมจริยํ จตฺตาโร สติปฏฺฐาเน จตฺตาโร สมฺมปฺปทฺ.หาเน จตฺตาโร อิทฺธิปาเท ปญฺจินฺทฺริยานิ ปญฺจ พลานิ สตฺต โพชฺฌงฺเค อริยํ อฏฺฐงฺคิกํ มคฺคํ นิพฺพานญฺจ นิพฺพานคามินิญฺจ ปฏิปทํ อาจิกฺขาหิ เทเสหิ ปญฺญเปหิ ปฏฺฐเปหิ วิวราหิ วิภชาหิ อุตฺตานีกโรหิ ปกาเสหีติ อาจิกฺข ธมฺมํ.\csromanspacer{} \textbf{ยมหํ วิชญฺญนฺ}ติ ยมหํ ชาเนยฺยํ อาชาเนยฺยํ วิชาเนยฺยํ ปฏิวิชาเนยฺยํ ปฏิวิชฺเฌยฺยํ อธิคจฺเฉยฺยํ ผสฺเสยฺยํ สจฺฉิกเรยฺยนฺติ อาจิกฺข ธมฺมํ ยมหํ วิชญฺญํ.}

\prose{\textbf{ชาติชราย อิธ วิปฺปหานนฺ}ติ อิเธว ชาติชรามรณสฺส ปหานํ วูปสมํ ปฏินิสฺสคฺคํ ปฏิปฺปสฺสทฺธิํ อมตํ นิพฺพานนฺติ ชาติชราย อิธ วิปฺปหานํ.\csromanspacer{} เตนาห โส พฺราหฺมโณ\csromandash{}}

\prose{\csromanfolio{190}ชิณฺโณหมสฺมิ อพโล วีตวณฺโณ, (อิจฺจายสฺมา ปิงฺคิโย,)}

\prose{เนตฺตา น สุทฺธา สวนํ น ผาสุ.}

\csromangathameasure{มาหํ นสฺสํ โมมุโห อนฺตราว, \\ อาจิกฺข ธมฺมํ ยมหํ วิชญฺญํ. \\ ชาติชราย อิธ วิปฺปหานนฺติ.}
\csromangathagroup{\csromangathastanza{มาหํ นสฺสํ โมมุโห อนฺตราว, \\ อาจิกฺข ธมฺมํ ยมหํ วิชญฺญํ. \\ ชาติชราย อิธ วิปฺปหานนฺติ.}}

\csromanhangingbegin
\hangingheaditem{90}{\textbf{ทิสฺวาน รูเปสุ วิหญฺญมาเน}, (\textbf{ปิงฺคิยาติ ภควา},)}
\hangingbody{\textbf{รุปฺปนฺติ รูเปสุ ชนา ปมตฺตา.}}
\hangingbody{\textbf{ตสฺมา ตุวํ ปิงฺคิย อปฺปมตฺโต,}}
\hangingbody{ชหสฺสุ รูปํ อปุนพฺภวาย.}
\csromanhangingend

\prose{\textbf{ทิสฺวาน รูเปสุ วิหญฺญมาเน}ติ \textbf{รูปนฺ}ติ จตฺตาโร จ มหาภูตา จตุนฺนญฺจ มหาภูตานํ อุปาทาย รูปํ, สตฺตา รูปเหตุ รูปปฺปจฺจยา รูปการณา หญฺญนฺติ วิหญฺญนฺติ อุปหญฺญนฺติ อุปฆาติยนฺติ\footnote{อุปฆาตยนฺติ (สฺยา, ก)}, รูเป สติ วิวิธกมฺมการณา\footnote{วิวิธกมฺมกรณานิ (ก)} กาเรนฺติ, กสาหิปิ ตาเฬนฺติ, เวตฺเตหิปิ ตาเฬนฺติ, อฑฺฒทณฺฑเกหิปิ ตาเฬนฺติ, หตฺถมฺปิ ฉินฺทนฺติ, ปาทมฺปิ ฉินฺทนฺติ, หตฺถปาทมฺปิ ฉินฺทนฺติ, กณฺณมฺปิ ฉินฺทนฺติ, นาสมฺปิ ฉินฺทนฺติ, กณฺณนาสมฺปิ ฉินฺทนฺติ, พิลงฺคถาลิกมฺปิ กโรนฺติ, สงฺขมุณฺฑิกมฺปิ กโรนฺติ, ราหุมุขมฺปิ กโรนฺติ, โชติมาลิกมฺปิ กโรนฺติ, หตฺถปชฺโชติกมฺปิ กโรนฺติ, เอรกวตฺติกมฺปิ กโรนฺติ, จีรกวาสิกมฺปิ กโรนฺติ, เอเณยฺยกมฺปิ กโรนฺติ, พฬิสมํสิกมฺปิ กโรนฺติ, กหาปณิกมฺปิ กโรนฺติ, ขาราปตจฺฉิกมฺปิ\footnote{ขาราปฏิจฺฉิกมฺปิ (ก)} กโรนฺติ, ปลิฆปริวตฺติกมฺปิ กโรนฺติ, ปลาลปีฐกมฺปิ กโรนฺติ, ตตฺเตนปิ เตเลน โอสิญฺจนฺติ, สุนเขหิปิ ขาทาเปนฺติ, ชีวนฺตมฺปิ สูเล อุตฺตาเสนฺติ, อสินาปิ สีสํ ฉินฺทนฺติ, เอวํ สตฺตา รูปเหตุ รูปปฺปจฺจยา รูปการณา หญฺญนฺติ วิหญฺญนฺติ อุปหญฺญนฺติ อุปฆาติยนฺติ.\csromanspacer{} เอวํ หญฺญมาเน วิหญฺญมาเน อุปหญฺญมาเน อุปฆาติยมาเน ทิสฺวา ปสฺสิตฺวา ตุลยิตฺวา ตีรยิตฺวา วิภาวยิตฺวา วิภูตํ กตฺวาติ ทิสฺวาน รูเปสุ วิหญฺญมาเน.}

\prose{\textbf{ปิงฺคิยาติ ภควา}ติ \textbf{ปิงฺคิยาติ ภควา} ตํ พฺราหฺมณํ นาเมน อาลปติ.\csromanspacer{} \textbf{ภควา}ติ คารวาธิวจนเมตํ ฯเปฯ สจฺฉิกา ปญฺญตฺติ, ยทิทํ \textbf{ภควา}ติ \textbf{ปิงฺคิยาติ ภควา}.}

\csromanheader{2}{16. ปิงฺคิยมาณวปุจฺฉานิทฺเทส}
\prose{\csromanfolio{191}\textbf{รุปฺปนฺติ รูเปสุ ชนา ปมตฺตา}ติ \textbf{รุปฺปนฺตี}ติ รุปฺปนฺติ กุปฺปนฺติ ปีฬยนฺติ\footnote{ปีฬิยนฺติ (สฺยา, ก)} ฆฏฺฏยนฺติ, พฺยาธิตา โทมนสฺสิตา โหนฺติ.\csromanspacer{} จกฺขุโรเคน รุปฺปนฺติ กุปฺปนฺติ ปีฬยนฺติ ฆฏฺฏยนฺติ, พฺยาธิตา โทมนสฺสิตา โหนฺติ.\csromanspacer{} โสตโรเคน ฯเปฯ.\csromanspacer{} กายโรเคน ฯเปฯ.\csromanspacer{} qอํสมกสวาตาตปสรีสปสมฺผสฺเสหิ รุปฺปนฺติ กุปฺปนฺติ ปีฬยนฺติ ฆฏฺฏยนฺติ, พฺยาธิตา โทมนสฺสิตา โหนฺตีติ รุปฺปนฺติ รูเปสุ.}

\prose{อถ วา จกฺขุสฺมิํ หียมาเน หายมาเน ปริหายมาเน เวมาเน\footnote{วิหายมาเน (ก)} วิคจฺฉมาเน อนฺตารธายมาเน รุปฺปนฺติ ฯเปฯ โทมนสฺสิตา โหนฺติ.\csromanspacer{} โสตสฺมิํ.\csromanspacer{} ฆานสฺมิํ.\csromanspacer{} ชิวฺหาย.\csromanspacer{} กายสฺมิํ.\csromanspacer{} รูปสฺมิํ.\csromanspacer{} สทฺทสฺมิํ.\csromanspacer{} คนฺธสฺมิํ.\csromanspacer{} รสสฺมิํ.\csromanspacer{} โผฏฺฐพฺพสฺมิํ.\csromanspacer{} กุลสฺมิํ.\csromanspacer{} คณสฺมิํ.\csromanspacer{} อาวาสสฺมิํ.\csromanspacer{} ลาภสฺมิํ.\csromanspacer{} ยสสฺมิํ.\csromanspacer{} ปสํสาย.\csromanspacer{} สุขสฺมิํ.\csromanspacer{} จีวรสฺมิํ.\csromanspacer{} ปิณฺฑปาตสฺมิํ.\csromanspacer{} เสนาสนสฺมิํ.\csromanspacer{} คิลานปจฺจยเภสชฺชปริกฺขารสฺมิํ หียมาเน หายมาเน ปริหายมาเน เวมาเน วิคจฺฉมาเน อนฺตรธายมาเน รุปฺปนฺติ กุปฺปนฺติ ปีฬยนฺติ ฆฏฺฏยนฺติ, พฺยาธิตา โทมนสฺสิตา โหนฺตีติ เอวมฺปิ รุปฺปนฺติ รูเปสุ.}

\prose{\textbf{ชนา}ติ ขตฺติยา จ พฺราหฺมณา จ เวสฺสา จ สุทฺทา จ คหฏฺฐา จ ปพฺพชิตา จ เทวา จ มนุสฺสา จ.\csromanspacer{} \textbf{ปมตฺตา}ติ ปมาโท วตฺตพฺโพ กายทุจฺจริเตน วา วจีทุจฺจริเตน วา มโนทุจฺจริเตน วา ปญฺจสุ กามคุเณสุ จิตฺตสฺส โวสคฺโค โวสคฺคานุปฺปทานํ กุสลานํ วา ธมฺมานํ ภาวนาย อสกฺกจฺจกิริยตา อสาตจฺจกิริยตา อนฏฺฐิตกิริยตา โอลีนวุตฺติตา นิกฺขิตฺตจฺฉนฺทตา นิกฺขิตฺตธุรตา อนาเสวนา อภาวนา อพหุลีกมฺมํ\footnote{อพหุลิกมฺมํ (ก)} อนธิฏฺฐานํ อนนุโยโค ปมาโท, โย เอวรูโป ปมาโท ปมชฺชนา ปมชฺชิตตฺตํ, อยํ วุจฺจติ ปมาโท, อิมินา ปมาเทน สมนฺนาคตา ชนา ปมตฺตาติ รุปฺปนฺติ รูเปสุ ชนา ปมตฺตา.}

\prose{\textbf{ตสฺมา ตุวํ ปิงฺคิย อปฺปมตฺโต}ติ \textbf{ตสฺมา}ติ \textbf{ตสฺมา} ตํการณา ตํเหตุ ตปฺปจฺจยา ตํนิทานํ เอวํ อาทีนวํ สมฺปสฺสมาโน รูเปสูติ \textbf{ตสฺมา} ตุวํ ปิงฺคิย.\csromanspacer{} \textbf{อปฺปมตฺโต}ติ สกฺกจฺจการี สาตจฺจการี ฯเปฯ อปฺปมาโท กุสเลสุ ธมฺเมสูติ \textbf{ตสฺมา} ตุวํ ปิงฺคิย อปฺปมตฺโต.}

\prose{\textbf{ชหสฺสุ รูปํ อปุนพฺภวายา}ติ \textbf{รูปนฺ}ติ จตฺตาโร จ มหาภูตา จตุนฺนญฺจ มหาภูตานํ อุปาทาย รูปํ.\csromanspacer{} \textbf{ชหสฺสุ รูปนฺ}ติ ชหสฺสุ รูปํ, ปชหสฺสุ รูปํ, วิโนเทหิ รูปํ, พฺยนฺตีกโรหิ รูปํ, อนภาวํ คเมหิ รูปํ.\csromanspacer{} \textbf{อปุนพฺภ}}

\prose{\csromanfolio{192}\textbf{วายา}ติ ยถา เต รูปํ อิเธว นิรุชฺเฌยฺย, ปุน ปฏิสนฺธิโก ภโว น นิพฺพตฺเตยฺย กามธาตุยา วา รูปธาตุยา วา อรูปธาตุยา วา, กามภเว วา รูปภเว วา อรูปภเว วา, สญฺญาภเว วา อสญฺญาภเว วา เนวสญฺญานาสญฺญาภเว วา, เอกโวการภเว วา จตุโวการภเว วา ปญฺจโวการภเว วา, ปุน คติยา วา อุปปตฺติยา วา ปฏิสนฺธิยา วา ภเว วา สํสาเร วา วฏฺเฏ วา น ชเนยฺย น สญฺชเนยฺย น นิพฺพตฺเตยฺย นาภินิพฺพตฺเตยฺย อิเธว นิรุชฺเฌย วูปสเมยฺย อตฺถํ คจฺเฉยฺย ปฏิปฺปสฺสมฺเภยฺยาติ ชหสฺสุ รูปํ อปุนพฺภวาย.\csromanspacer{} เตนาห ภควา\csromandash{}}

\prose{ทิสฺวาน รูเปสุ วิหญฺญมาเน, (ปิงฺคิยาติ ภควา,)}

\prose{รุปฺปนฺติ รูเปสุ ชนา ปมตฺตา.}

\csromangathameasure{ตสฺมา ตุวํ ปิงฺคิย อปฺปมตฺโต, \\ ชหสฺสุ รูปํ อปุนพฺภ\textbf{วายา}ติ.}
\csromangathagroup{\csromangathastanza{ตสฺมา ตุวํ ปิงฺคิย อปฺปมตฺโต, \\ ชหสฺสุ รูปํ อปุนพฺภ\textbf{วายา}ติ.}}

\csromanhangingbegin
\hangingheaditem{91}{ทิสา จตสฺโส วิทิสา จตสฺโส,}
\hangingbody{\textbf{อุทฺธํ อโธ ทส ทิสา อิมาโย.}}
\hangingbody{\textbf{น ตุยฺหํ อทิฏฺฐํ อสุตํ อมุตํ,}}
\hangingbody{\textbf{อโถ อวิญฺญาตํ กิญฺจนมตฺถิ โลเก.}}
\hangingbody{\textbf{อาจิกฺข ธมฺมํ ยมหํ วิชญฺญํ,}}
\hangingbody{ชาติชราย อิธ วิปฺปหานํ.}
\csromanhangingend

\prose{\textbf{ทิสาจตสฺโส วิทิสา จตสฺโส, อุทฺธํ อโธ ทส ทิสา อิมาโย}ติ ทส ทิสา.}

\prose{\textbf{น ตุยฺหํ อทิฏฺฐํ อสุตํ อมุตํ, อโถ อวิญฺญาตํ กิญฺจนมตฺถิ โลเก}ติ น ตุยฺหํ อทิฏฺฐํ อสุตํ อมุตํ อวิญฺญาตํ กิญฺจิ อตฺตตฺโถ วา ปรตฺโถ วา อุภยตฺโถ วา ทิฏฺฐธมฺมิโก วา อตฺโถ สมฺปรายิโก วา อตฺโถ อุตฺตาโน วา อตฺโถ คมฺภีโร วา อตฺโถ คูโฬฺห วา อตฺโถ ปฏิจฺฉนฺโน วา อตฺโถ เนยฺโย วา อตฺโถ นีโต วา อตฺโถ อนวชฺโช วา อตฺโถ นิกฺกิเลโส วา อตฺโถ โวทาโน วา อตฺโถ ปรมตฺโถ วา อตฺโถ นตฺถิ น อตฺถิ น สํวิชฺชติ นุปลพฺภตีติ น ตุยฺหํ อทิฏฺฐํ อสุตํ อมุตํ อโถ อวิญฺญาตํ กิญฺจนมตฺถิ โลเก.}

\prose{\textbf{อาจิกฺข ธมฺมํ ยมหํ วิชญฺญนฺ}ติ \textbf{ธมฺมนฺ}ติ อาทิกลฺยาณํ ฯเปฯ นิพฺพานญฺจ นิพฺพานคามินิญฺจ ปฏิปทํ อาจิกฺขาหิ เทเสหิ ปญฺญเปหิ ปฏฺฐเปหิ วิวราหิ วิภชาหิ}

\vspace{\tipitakaparskip}
\csromancenter{ปิงฺคิยมาณวปุจฺฉานิทฺเทส อุตฺตานีกโรหิ ปกาเสหิ.\csromanspacer{} \textbf{ยมหํ วิชญฺญนฺ}ติ ยมหํ ชาเนยฺยํ อาชาเนยฺยํ วิชาเนยฺยํ ปฏิวิชาเนยฺยํ ปฏิวิชฺเฌยฺยํ อธิคจฺเฉยฺยํ ผสฺเสยฺยํ สจฺฉิกเรยฺยนฺติ อาจิกฺข ธมฺมํ ยมหํ วิชญฺญํ.}
\prose{\csromanfolio{193}\textbf{ชาติชราย อิธ วิปฺปหานนฺ}ติ อิเธว ชาติชรามรณสฺส ปหานํ วูปสมํ ปฏินิสฺสคฺคํ ปฏิปฺปสฺสทฺธิํ อมตํ นิพฺพานนฺติ ชาติชราย อิธ วิปฺปหานํ.\csromanspacer{} เตนาห โส พฺราหฺมโณ\csromandash{}}

\csromangathameasure{ทิสา จตสฺโส วิทิสา จตสฺโส, \\ อุทฺธํ อโธ ทส ทิสา อิมาโย. \\ น ตุยฺหํ อทิฏฺฐํ อสุตํ อมุตํ, \\ อโถ อวิญฺญาตํ กิญฺจนมตฺถิ โลเก. \\ อาจิกฺข ธมฺมํ ยมหํ วิชญฺญํ,}
\csromangathagroup{\csromangathastanza{ทิสา จตสฺโส วิทิสา จตสฺโส, \\ อุทฺธํ อโธ ทส ทิสา อิมาโย. \\ น ตุยฺหํ อทิฏฺฐํ อสุตํ อมุตํ, \\ อโถ อวิญฺญาตํ กิญฺจนมตฺถิ โลเก.}\csromangathastanza{อาจิกฺข ธมฺมํ ยมหํ วิชญฺญํ,}}

\prose{\textbf{ชาติชราย อิธ วิปฺปหานนฺ}ติ.}

\csromanhangingbegin
\hangingheaditem{92}{\textbf{ตณฺหาธิปนฺเน} มนุเช เปกฺขมาโน, (\textbf{ปิงฺคิยาติ ภควา},)}
\hangingbody{\textbf{สนฺตาปชาเต ชรสา ปเรเต.}}
\hangingbody{\textbf{ตสฺมา ตุวํ ปิงฺคิย อปฺปมตฺโต,}}
\hangingbody{ชหสฺสุ ตณฺหํ อปุนพฺภวาย.}
\csromanhangingend

\prose{\textbf{ตณาธิปนฺเน มนุเช เปกฺขมาโน}ติ \textbf{ตณฺหา}ติ รูป\textbf{ตณฺหา} ฯเปฯ ธมฺม\textbf{ตณฺหา}.\csromanspacer{} \textbf{ตณฺหาธิปนฺเน}ติ \textbf{ตณฺหา}ธิปนฺเน\footnote{ตณฺหาย อธิปนฺเน (ก)} \textbf{ตณฺหา}นุเค \textbf{ตณฺหา}นุคเต \textbf{ตณฺหา}นุสเฏ \textbf{ตณฺหา}ย ปนฺเน ปฏิปนฺเน อภิภูเต ปริยาทินฺนจิตฺเต.\csromanspacer{} \textbf{มนุเช}ติ สตฺตาธิวจนํ.\csromanspacer{} \textbf{เปกฺขมาโน}ติ เปกฺขมาโน ทกฺขมาโน โอโลกยมาโน นิชฺฌายมาโน อุปปริกฺขมาโนติ \textbf{ตณฺหา}ธิปนฺเน มนุเช เปกฺขมาโน.\csromanspacer{} \textbf{ปิงฺคิยาติ ภควา}ติ \textbf{ปิงฺคิยาติ ภควา} ตํ พฺราหฺมณํ นาเมน อาลปติ.\csromanspacer{} \textbf{ภควา}ติ คารวาธิวจนเมตํ ฯเปฯ สจฺฉิกา ปญฺญตฺติ, ยทิทํ \textbf{ภควา}ติ \textbf{ปิงฺคิยาติ ภควา}.}

\prose{\textbf{สนฺตาปชาเต ชรสา ปเรเต}ติ \textbf{สนฺตาปชาเต}ติ ชาติยา \textbf{สนฺตาปชาเต}, ชราย \textbf{สนฺตาปชาเต}, พฺยาธินา \textbf{สนฺตาปชาเต}, มรเณน \textbf{สนฺตาปชาเต}, โสกปริเทวทุกฺขโทมนสฺสุปายาเสหิ \textbf{สนฺตาปชาเต}, เนรยิเกน ทุกฺเขน \textbf{สนฺตาปชาเต} ฯเปฯ ทิฏฺฐิพฺยสเนน ทุกฺเขน \textbf{สนฺตาปชาเต} \csromanfolio{194}อีติชาเต อุปทฺทวชาเต อุปสคฺคชาเตติ สนฺตาปชาเต.\csromanspacer{} \textbf{ชรสา ปเรเต}ติ ชราย ผุฏฺเฐ ปเรเต สโมหิเต สมนฺนาคเต.\csromanspacer{} ชาติยา อนุคเต ชราย อนุสเฏ พฺยาธินา อภิภูเต มรเณน อพฺภาหเต อตาเณ อเลเณ อสรเณ อสรณีภูเตติ สนฺตาปชาเต ชรสา ปเรเต.}

\prose{\textbf{ตสฺมา ตุวํ ปิงฺคิย อปฺปมตฺโต}ติ \textbf{ตสฺมา}ติ \textbf{ตสฺมา} ตํการณา ตํเหตุ ตปฺปจฺจยา ตํนิทานา เอวํ อาทีนวํ สมฺปสฺสมาโน \textbf{ตณฺหา}ยาติ \textbf{ตสฺมา} ตุวํ ปิงฺคิย.\csromanspacer{} \textbf{อปฺปมตฺโต}ติ สกฺกจฺจการี ฯเปฯ อปฺปมาโท กุสเลสุ ธมฺเมสูติ \textbf{ตสฺมา} ตุวํ ปิงฺคิย อปฺปมตฺโต.}

\prose{\textbf{ชหสฺสุ ตณฺหํ อปุนพฺภวายา}ติ \textbf{ตณฺหา}ติ รูป\textbf{ตณฺหา} ฯเปฯ ธมฺม\textbf{ตณฺหา}.\csromanspacer{} \textbf{ชหสฺสุ ตณฺหนฺ}ติ ชหสฺสุ ตณฺหํ, ปชหสฺสุ ตณฺหํ, วิโนเทหิ ตณฺหํ, พฺยนฺตีกโรหิ ตณฺหํ, อนภาวํ คเมหิ ตณฺหํ.\csromanspacer{} \textbf{อปุนพฺภวายา}ติ ยถา เต ฯเปฯ ปุนปฏิสนฺธิโก ภโว น นิพฺพตฺเตยฺย กามธาตุยา วา รูปธาตุยา วา อรูปธาตุยา วา กามภเว วา รูปภเว วา อรูปภเว วา สญฺญาภเว วา อสญฺญาภเว วา เนวสญฺญานาสญฺญาภเว วา, เอกโวการภเว วา จตุโวการภเว วา ปญฺจโวการภเว วา, ปุนคติยา วา อุปปตฺติยา วา ปฏิสนฺธิยา วา ภเว วา สํสาเร วา วฏฺเฏ วา น ชเนยฺย น สญฺชเนยฺย น นิพฺพตฺเตยฺย นาภินิพฺพตฺเตยฺย อิเธว นิรุชฺเฌยฺย วูปสเมยฺย อตฺถํ คจฺเฉยฺย ปฏิปฺปสฺสมฺเภยฺยาติ ชหสฺสุ ตณฺหํ อปุนพฺภวาย.\csromanspacer{} เตนาห ภควา\csromandash{}}

\prose{ตณฺหาธิปนฺเน มนุเช เปกฺขมาโน, (ปิงฺคิยาติ ภควา,)}

\prose{สนฺตาปชาเต ชรสา ปเรเต.}

\prose{\textbf{ตสฺมา ตุวํ ปิงฺคิย อปฺปมตฺโต},}

\prose{\textbf{ชหสฺสุ ตณฺหํ อปุนพฺภวายา}ติ.}

\prose{สห คาถาปริโยสานา เย เต พฺราหฺมเณน สทฺธิํ เอกจฺฉนฺทา เอกปโยคา เอกาธิปฺปายา เอกวาสนวาสิตา, เตสํ อเนกปาณสหสฺสานํ วิรชํ วีตมลํ ธมฺมจกฺขุํ อุทปาทิ “ยํ กิญฺจิ สมุทยธมฺมํ, สพฺพํ ตํ นิโรธธมฺมนฺ”ติ.\csromanspacer{} ตสฺส จ พฺราหฺมณสฺส วิรชํ วีตมลํ ธมฺมจกฺขุํ อุทปาทิ “ยํ กิญฺจิ สมุทยธมฺมํ, สพฺพํ ตํ นิโรธธมฺมนฺ”ติ, สห ธมฺมจกฺขุสฺส ปฏิลาภา อชินชฏาวากจีรติทณฺฑกมณฺฑลุเกสา จ มสฺสู จ อนฺตรหิตา.}

\vspace{\tipitakaparskip}
\csromancenter{ปารายนตฺถุติคาถานิทฺเทส ภณฺฑุกาสายวตฺถวสโน สงฺฆาฏิปตฺตจีวรธโร อนฺวตฺถปฏิปตฺติยา ปญฺชลิโก ภควนฺตํ นมสฺสมาโน นิสินฺโน โหติ “สตฺถา เม ภนฺเต \textbf{ภควา}, สาวโกหมสฺมี”ติ.}
\csromancenter{ปิงฺคิยมาณวปุจฺฉานิทฺเทโส\footnote{สิงฺคิยปญฺหํ (ก)} โสฬสโม.}
\csromanmatikaanchor{mtk.38}
\csromantocmarkref{subsection}{17. ปารายนตฺถุติคาถานิทฺเทส}{mtk.38}
\bookmark[dest={mtk.38},level=2]{17. ปารายนตฺถุติคาถานิทฺเทส}
\csromanheader{2}{17. ปารายนตฺถุติคาถานิทฺเทส}
\proseitem{93}{\csromanfolio{195}\textbf{อิทมโวจ ภควา} \textbf{มคเธสุ วิหรนฺโต} ปาสาณเก เจติเย, ปริจารกโสฬสานํ\footnote{ปริจาริกโสฬสนฺนํ (สฺยา, ก)} \textbf{พฺราหฺมณานํ อชฺฌิฏฺโฐ ปุฏฺโฐ ปุฏฺโฐ ปญฺหํ พฺยากาสิ.}}

\prose{\textbf{อิทมโวจ ภควา}ติ อิทํ ปารายนํ อโวจ.\csromanspacer{} \textbf{ภควา}ติ คารวาธิวจนเมตํ ฯเปฯ สจฺฉิกา ปญฺญตฺติ, ยทิทํ \textbf{ภควา}ติ อิทมโวจ \textbf{ภควา}.\csromanspacer{} \textbf{มคเธสุ วิหรนฺโต}ติ มคธนามเก ชนปเท วิหรนฺโต อิริยโต วตฺเตนฺโต ปาเลนฺโต ยเปนฺโต ยาเปนฺโต.\csromanspacer{} \textbf{ปาสาณเก เจติเย}ติ ปาสาณกเจติยํ วุจฺจติ พุทฺธาสนนฺติ \textbf{มคเธสุ วิหรนฺโต} ปาสาณเก เจติเย.\csromanspacer{} \textbf{ปริจารกโสฬสานํ พฺราหฺมณานนฺ}ติ ปิงฺคิโย\footnote{สิงฺคิโย (ก)} พฺราหฺมโณ พาวริสฺส พฺราหฺมณสฺส ปทฺโธปทฺธจโร ปริจารโก\footnote{ปริจาริโก (สฺยา, ก)} สิสฺโส, ปิงฺคิเยน\footnote{เตน (ก)} เต โสฬสาติ เอวมฺปิ ปริจารกโสฬสานํ พฺราหฺมณานํ.\csromanspacer{} อถ วา เต โสฬส พฺราหฺมณา พุทฺธสฺส ภควโต ปทฺธา ปทฺธจรา ปริจารกา สิสฺสาติ เอวมฺปิ ปริจารกโสฬสานํ พฺราหฺมณานํ.}

\prose{\textbf{อชฺฌิฏฺโฐ ปุฏฺโฐ ปุฏฺโฐ ปญฺหํ พฺยากาสี}ติ \textbf{อชฺฌิฏฺโฐ}ติ \textbf{อชฺฌิฏฺโฐ} อชฺเฌสิโต.\csromanspacer{} \textbf{ปุฏฺโฐ ปุฏฺโฐ}ติ ปุฏฺโฐ ปุฏฺโฐ ปุจฺฉิโต ปุจฺฉิโต ยาจิโต ยาจิโต อชฺเฌสิโต อชฺเฌสิโต ปสาทิโต ปสาทิโต.\csromanspacer{} \textbf{ปญฺหํ พฺยากาสี}ติ ปญฺหํ พฺยากาสิ อาจิกฺขิ เทเสสิ ปญฺญเปสิ ปฏฺฐเปสิ วิวริ วิภชิ อุตฺตานี-อากาสิ ปกาเสสีติ \textbf{อชฺฌิฏฺโฐ} ปุฏฺโฐ ปุฏฺโฐ ปญฺหํ พฺยากาสิ.\csromanspacer{} เตเนตํ วุจฺจติ\csromandash{}}

\prose{\textbf{อิทมโวจ ภควา} \textbf{มคเธสุ วิหรนฺโต} ปาสาณเก เจติเย,}

\prose{ปริจารกโสฬสานํ พฺราหฺมณานํ \textbf{อชฺฌิฏฺโฐ} ปุฏฺโฐ ปุฏฺโฐ ปญฺหํ}

\prose{พฺยากาสีติ.}

\proseitem{94}{\csromanfolio{196}\textbf{เอกเมกสฺส เจปิ ปญฺหสฺส อตฺถมญฺญาย ธมฺมมญฺญาย ธมฺมานุธมฺมํ ปฏิปชฺเชยฺย, คจฺเฉยฺเยว ชรามรณสฺส ปารํ, ปารงฺคมนียา อิเม ธมฺมาติ, ตสฺมา อิมสฺส ธมฺมปริยายสฺส ปารายนนฺเตว อธิวจนํ.}}

\prose{\textbf{เอกเมกสฺส เจปิ ปญฺหสฺสา}ติ เอกเมกสฺส เจปิ อชิตปญฺหสฺส เอกเมกสฺส เจปิ ติสฺสเมตฺเตยฺยปญฺหสฺส เอกเมกสฺส เจปิ ปุณฺณกปญฺหสฺส เอกเมกสฺส เจปิ เมตฺตคูปญฺหสฺส เอกเมกสฺส เจปิ โธตกปญฺหสฺส เอกเมกสฺส เจปิ อุปสีวปญฺหสฺส เอกเมกสฺส เจปิ นนฺทกปญฺหสฺส เอกเมกสฺส เจปิ เหมกปญฺหสฺส เอกเมกสฺส เจปิ โตเทยฺยปญฺหสฺส เอกเมกสฺส เจปิ กปฺปปญฺหสฺส เอกเมกสฺส เจปิ ชตุกณฺณิกปญฺหสฺส เอกเมกสฺส เจปิ ภทฺราวุธปญฺหสฺส เอกเมกสฺส เจปิ อุทยปญฺหสฺส เอกเมกสฺส เจปิ โปสาลปญฺหสฺส เอกเมกสฺส เจปิ โมฆราชปญฺหสฺส เอกเมกสฺส เจปิ ปิงฺคิยปญฺหสฺสาติ เอกเมกสฺส เจปิ ปญฺหสฺส.}

\prose{\textbf{อตฺถมญฺญาย ธมฺมมญฺญายา}ติ เสฺวว ปโญฺห ธมฺโม, วิสชฺชนํ\footnote{วิสฺสชฺชนํ (ก)} อตฺโถติ อตฺถํ อญฺญาย ชานิตฺวา ตุลยิตฺวา ตีรยิตฺวา วิภาวยิตฺวา วิภูตํ กตฺวาติ อตฺถมญฺญาย.\csromanspacer{} \textbf{ธมฺมมญฺญายา}ติ ธมฺมํ อญฺญาย ชานิตฺวา ตุลยิตฺวา ตีรยิตฺวา วิภาวยิตฺวา วิภูตํ กตฺวาติ ธมฺมมญฺญายาติ อตฺถมญฺญาย ธมฺมมญฺญาย.\csromanspacer{} \textbf{ธมฺมานุธมฺมํ ปฏิปชฺเชยฺยา}ติ สมฺมาปฏิปทํ อนุโลมปฏิปทํ อปจฺจนีกปฏิปทํ อนฺวตฺถปฏิปทํ ธมฺมานุธมฺมปฏิปทํ ปฏิปชฺเชยฺยาติ ธมฺมานุธมฺมํ ปฏิปชฺเชยฺย.\csromanspacer{} \textbf{คจฺเฉยฺเยว ชรามรณสฺส ปารนฺ}ติ ชรามรณสฺส ปารํ วุจฺจติ อมตํ นิพฺพานํ.\csromanspacer{} โย โส สพฺพสงฺขารสมโถ สพฺพูปธิปฺปฏินิสฺสคฺโค ตณฺหกฺขโย วิราโค นิโรโธ นิพฺพานํ.\csromanspacer{} \textbf{คจฺเฉยฺเยว ชรามรณสฺส ปารนฺ}ติ ชรามรณสฺส ปารํ คจฺเฉยฺย, ปารํ อธิคจฺเฉยฺย, ปารํ อธิผสฺเสยฺย, ปารํ สจฺฉิกเรยฺยาติ คจฺเฉยฺเยว ชรามรณสฺส ปารํ.\csromanspacer{} \textbf{ปารงฺคมนียา อิเม ธมฺมา}ติ อิเม ธมฺมา ปารงฺคมนียา, ปารํ ปาเปนฺติ, ปารํ สมฺปาเปนฺติ, ปารํ สมนุปาเปนฺติ, ชรามรณสฺส ตรณาย\footnote{ตารณาย (สฺยา)} สํวตฺตนฺตีติ ปารงฺคมนียา อิเม ธมฺมาติ.}

\prose{\textbf{ตสฺมา อิมสฺส ธมฺมปริยายสฺสา}ติ \textbf{ตสฺมา}ติ \textbf{ตสฺมา} ตํการณา ตํเหตุ ตปฺปจฺจยา ตํนิทานาติ \textbf{ตสฺมา}.\csromanspacer{} \textbf{อิมสฺส ธมฺมปริยายสฺสา}ติ}

\vspace{\tipitakaparskip}
\csromancenter{ปารายนตฺถุติคาถานิทฺเทส อิมสฺส ปารายนสฺสาติ ตสฺมา อิมสฺส ธมฺมปริยายสฺส.\csromanspacer{} \textbf{ปารายนนฺเตว อธิวจนนฺ}ติ ปารํ วุจฺจติ อมตํ นิพฺพานํ ฯเปฯ นิโรโธ นิพฺพานํ.\csromanspacer{} อยนํ วุจฺจติ มคฺโค.\csromanspacer{} เสยฺยถิทํ, สมฺมาทิฏฺฐิ ฯเปฯ สมฺมาสมาธิ.\csromanspacer{} \textbf{อธิวจนนฺ}ติ สงฺขา สมญฺญา ปญฺญตฺติ โวหาโร นามํ นามกมฺมํ นามเธยฺยํ นิรุตฺติ พฺยญฺชนํ อภิลาโปติ ปารายนนฺเตว อธิวจนํ.\csromanspacer{} เตเนตํ วุจฺจติ\csromandash{}}
\prose{\csromanfolio{197}เอกเมกสฺส เจปิ ปญฺหสฺส, อตฺถมญฺญาย ธมฺมมญฺญาย}

\prose{ธมฺมานุธมฺมํ ปฏิปชฺเชยฺย, คจฺเฉยฺเยว ชรามรณสฺส}

\prose{ปารํ, ปารงฺคมนียา อิเม ธมฺมาติ, ตสฺมา อิมสฺส}

\prose{ธมฺมปริยายสฺส ปารายนนฺเตว อธิวจนนฺติ.}

\csromanhangingbegin
\hangingheaditem{95}{อชิโต ติสฺสเมตฺเตยฺโย, ปุณฺณโก อถ เมตฺตคู.}
\hangingbody{โธตโก อุปสีโว จ, นนฺโท จ อถ เหมโก.}
\csromanhangingend

\csromanhangingbegin
\hangingheaditem{96}{โตเทยฺย-กปฺปา ทุภโย, ชตุกณฺณี จ ปณฺฑิโต.}
\hangingbody{\textbf{ภทฺราวุโธ อุทโย จ, โปสาโล จาปิ พฺราหฺมโณ.}}
\hangingbody{โมฆราชา จ เมธาวี, ปิงฺคิโย จ มหา-อิสิ.}
\csromanhangingend

\csromanhangingbegin
\hangingheaditem{97}{เอเต พุทฺธํ อุปาคจฺฉุํ, สมฺปนฺนจรณํ อิสิํ.}
\hangingbody{ปุจฺฉนฺตา นิปุเณ ปเญฺห, พุทฺธเสฏฺฐํ อุปาคมุํ.}
\csromanhangingend

\prose{\textbf{เอเต พุทฺธํ อุปาคจฺฉุนฺ}ติ \textbf{เอเต}ติ โสฬส ปารายนิยา พฺราหฺมณา.\csromanspacer{} พุทฺโธติ โย โส ภควา สยมฺภู อนาจริยโก ปุพฺเพ อนนุสฺสุเตสุ ธมฺเมสุ สามํ สจฺจานิ อภิสมฺพุชฺฌิ, ตตฺถ จ สพฺพญฺญุตํ ปตฺโต พเลสุ จ วสีภาวํ.\csromanspacer{} พุทฺโธติ เกนฏฺเฐน พุทฺโธ, พุชฺฌิตา สจฺจานีติ พุทฺโธ, โพเธตา ปชายาติ พุทฺโธ, สพฺพญฺญุตาย พุทฺโธ, สพฺพทสฺสาวิตาย พุทฺโธ, อภิญฺเญยฺยตาย พุทฺโธ, วิสวิตาย พุทฺโธ, ขีณาสวสงฺขาเตน พุทฺโธ, นิรุปกฺกิเลสสงฺขาเตน พุทฺโธ, เอกนฺตวีตราโคติ พุทฺโธ, เอกนฺตวีตโทโสติ พุทฺโธ, เอกนฺตวีตโมโหติ พุทฺโธ, เอกนฺตนิกฺกิเลโสติ พุทฺโธ, เอกายนมคฺคํ คโตติ พุทฺโธ, เอโก อนุตฺตรํ สมฺมาสมฺโพธิํ อภิสมฺพุทฺโธติ พุทฺโธ, อพุทฺธิวิหตตฺตา พุทฺธิปฏิลาภาติ พุทฺโธ.\csromanspacer{} พุทฺโธติ เนตํ นามํ มาตรา กตํ, น ปิตรา กตํ, น ภาตรา กตํ, น ภคินิยา กตํ, น มิตฺตามจฺเจหิ กตํ, น ญาติสาโลหิเตหิ กตํ, น สมณพฺราหฺมเณหิ กตํ, น เทวตาหิ กตํ, วิโมกฺขนฺติกเมตํ พุทฺธานํ ภควนฺตานํ โพธิยา มูเล สห สพฺพญฺญุตญาณสฺส ปฏิลาภา สจฺฉิกา \csromanfolio{198}ปญฺญตฺติ, ยทิทํ พุทฺโธติ.\csromanspacer{} \textbf{เอเต พุทฺธํ อุปาคจฺฉุนฺ}ติ เอ\textbf{เต} พุทฺธํ อุปาคมิํสุ อุปสงฺกมิํสุ ปยิรุปาสิํสุ ปริปุจฺฉิํสุ ปริปญฺหิํสูติ เอ\textbf{เต} พุทฺธํ อุปาคจฺฉุํ.}

\prose{\textbf{สมฺปนฺนจรณํ อิสินฺ}ติ \textbf{จรณํ} วุจฺจติ สีลาจารนิพฺพตฺติ.\csromanspacer{} สีลสํวโรปิ \textbf{จรณํ}, อินฺทฺริยสํวโรปิ \textbf{จรณํ}, โภชเน มตฺตญฺญุตาปิ \textbf{จรณํ}, ชาคริยานุโยโคปิ \textbf{จรณํ}, สตฺตปิ สทฺธมฺมา \textbf{จรณํ}, จตฺตาริปิ ฌานานิ \textbf{จรณํ}.\csromanspacer{} \textbf{สมฺปนฺนจรณนฺ}ติ สมฺปนฺน\textbf{จรณํ} เสฏฺฐ\textbf{จรณํ} วิเสฏฺฐ\textbf{จรณํ}\footnote{วิสิฏฺฐจรณํ (ก)} ปาโมกฺข\textbf{จรณํ} อุตฺตม\textbf{จรณํ} ปวร\textbf{จรณํ}.\csromanspacer{} \textbf{อิสี}ติ อิสิ, ภควา มหนฺตํ สีลกฺขนฺธํ เอสี คเวสี ปริเยสีติ อิสิ ฯเปฯ มเหสกฺเขหิ วา สตฺเตหิ เอสิโต คเวสิโต ปริเยสิโต “กหํ พุทฺโธ กหํ ภควา กหํ เทวเทโว กหํ นราสโภ”ติ อิสีติ สมฺปนฺน\textbf{จรณํ} อิสิํ.}

\prose{\textbf{ปุจฺฉนฺตา นิปุเณ ปเญฺห}ติ \textbf{ปุจฺฉนฺตา}ติ \textbf{ปุจฺฉนฺตา} ยาจนฺตา อชฺเฌสนฺตา ปสาเทนฺตา.\csromanspacer{} \textbf{นิปุเณ ปเญฺห}ติ คมฺภีเร ทุทฺทเส ทุรนุโพเธ สนฺเต ปณีเต อตกฺกาวจเร นิปุเณ ปณฺฑิตเวทนีเย ปเญฺหติ \textbf{ปุจฺฉนฺตา} นิปุเณ ปเญฺห.}

\prose{\textbf{พุทฺธเสฏฺฐํ อุปาคมุนฺ}ติ พุทฺโธติ โย โส ภควา ฯเปฯ สจฺฉิกา ปญฺญตฺติ, ยทิทํ พุทฺโธติ.\csromanspacer{} \textbf{เสฏฺฐนฺ}ติ อคฺคํ เสฏฺฐํ วิเสฏฺฐํ ปาโมกฺขํ อุตฺตมํ ปวรํ พุทฺธํ อุปาคมุํ อุปาคมิํสุ อุปสงฺกมิํสุ ปยิรุปาสิํสุ ปริปุจฺฉิํสุ ปริปญฺหิํสูติ พุทฺธเสฏฺฐํ อุปาคมุํ.\csromanspacer{} เตเนตํ วุจฺจติ\csromandash{}}

\csromangathameasure{เอเต พุทฺธํ อุปาคจฺฉุํ, สมฺปนฺนจรณํ อิสิํ. \\ \textbf{ปุจฺฉนฺตา นิปุเณ ปเญฺห}, \textbf{พุทฺธเสฏฺฐํ อุปาคมุนฺ}ติ.}
\csromangathagroup{\csromangathastanza{เอเต พุทฺธํ อุปาคจฺฉุํ, สมฺปนฺนจรณํ อิสิํ. \\ \textbf{ปุจฺฉนฺตา นิปุเณ ปเญฺห}, \textbf{พุทฺธเสฏฺฐํ อุปาคมุนฺ}ติ.}}

\csromanhangingbegin
\hangingheaditem{98}{เตสํ พุทฺโธ ปพฺยากาสิ, \textbf{ปญฺหํ ปุฏฺโฐ} ยถาตถํ.}
\hangingbody{ปญฺหานํ เวยฺยากรเณน, โตเสสิ พฺราหฺมเณ มุนิ.}
\csromanhangingend

\prose{\textbf{เตสํ พุทฺโธ ปพฺยากาสี}ติ \textbf{เตสนฺ}ติ โสฬสานํ ปารายนิยานํ พฺราหฺมณานํ.\csromanspacer{} พุทฺโธติ โย โส ภควา ฯเปฯ สจฺฉิกา ปญฺญตฺติ, ยทิทํ พุทฺโธติ.\csromanspacer{} \textbf{ปพฺยากาสี}ติ เตสํ พุทฺโธ ปพฺยากาสิ อาจิกฺขิ เทเสสิ ปญฺญเปสิ ปฏฺฐเปสิ วิวริ วิภชิ อุตฺตานี-อกาสิ ปกาเสสีติ เตสํ พุทฺโธ ปพฺยากาสิ.}

\prose{\textbf{ปญฺหํ ปุฏฺโฐ ยถาตถนฺ}ติ \textbf{ปญฺหํ ปุฏฺโฐ}ติ \textbf{ปญฺหํ ปุฏฺโฐ} ปุจฺฉิโต ยาจิโต อชฺเฌสิโต ปสาทิโต.\csromanspacer{} \textbf{ยถาตถนฺ}ติ ยถา อาจิกฺขิตพฺพํ, ตถา}

\vspace{\tipitakaparskip}
\csromancenter{ปารายนตฺถุติคาถานิทฺเทส อาจิกฺขิ.\csromanspacer{} ยถา เทสิตพฺพํ, ตถา เทเสสิ.\csromanspacer{} ยถา ปญฺญเปตพฺพํ, ตถา ปญฺญเปสิ.\csromanspacer{} ยถา ปฏฺฐเปตพฺพํ, ตถา ปฏฺฐเปสิ.\csromanspacer{} ยถา วิวริตพฺพํ, ตถา วิวริ.\csromanspacer{} ยถา วิภชิตพฺพํ, ตถา วิภชิ.\csromanspacer{} ยถา อุตฺตานีกาตพฺพํ, ตถา อุตฺตานี-อกาสิ.\csromanspacer{} ยถา ปกาสิตพฺพํ, ตถา ปกาเสสีติ ปญฺหํ ปุฏฺโฐ ยถาตถํ.}
\prose{\csromanfolio{199}\textbf{ปญฺหานํ เวยฺยากรเณนา}ติ ปญฺหานํ เวยฺยากรเณน อาจิกฺขเนน เทสเนน ปญฺญปเนน ปฏฺฐปเนน วิวรเณน วิภชเนน อุตฺตานีกมฺเมน ปกาสเนนาติ ปญฺหานํ เวยฺยากรเณน.}

\prose{\textbf{โตเสสิ พฺราหฺมเณ มุนี}ติ \textbf{โตเสสี}ติ โตเสสิ วิโตเสสิ ปสาเทสิ อาราเธสิ อตฺตมเน อกาสิ.\csromanspacer{} \textbf{พฺราหฺมเณ}ติ โสฬส ปารายนิเย พฺราหฺมเณ.\csromanspacer{} \textbf{มุนี}ติ โมนํ วุจฺจติ ญาณํ ฯเปฯ สงฺคชาลมติจฺจ โส มุนีติ โตเสสิ พฺราหฺมเณ มุนิ.\csromanspacer{} เตเนตํ วุจฺจติ\csromandash{}}

\csromangathameasure{เตสํ พุทฺโธ ปพฺยากาสิ, ปญฺหํ ปุฏฺโฐ ยถาตถํ. \\ ปญฺหานํ เวยฺยากรเณน, โตเสสิ พฺราหฺมเณ มุนีติ.}
\csromangathagroup{\csromangathastanza{เตสํ พุทฺโธ ปพฺยากาสิ, ปญฺหํ ปุฏฺโฐ ยถาตถํ. \\ ปญฺหานํ เวยฺยากรเณน, โตเสสิ พฺราหฺมเณ มุนีติ.}}

\csromanhangingbegin
\hangingheaditem{99}{\textbf{เต โตสิตา จกฺขุมตา}, \textbf{พุทฺเธนาทิจฺจพนฺธุนา}.}
\hangingbody{พฺรหฺมจริยมจริํสุ, วรปญฺญสฺส สนฺติเก.}
\csromanhangingend

\prose{\textbf{เต โตสิตา จกฺขุมตา}ติ \textbf{เต}ติ โสฬส ปารายนิยา พฺราหฺมณา.\csromanspacer{} \textbf{โตสิตา}ติ โตสิตา วิโตสิตา ปสาทิตา อาราธิตา อตฺตมนา กตาติ \textbf{เต} โตสิตา.\csromanspacer{} \textbf{จกฺขุมตา}ติ ภควา ปญฺจหิ จกฺขูหิ จกฺขุมา, มํสจกฺขุนาปิ จกฺขุมา, ทิพฺพจกฺขุนาปิ จกฺขุมา, ปญฺญาจกฺขุนาปิ จกฺขุมา, พุทฺธจกฺขุนาปิ จกฺขุมา, สมนฺตจกฺขุนาปิ จกฺขุมา.\csromanspacer{} กถํ ภควา มํสจกฺขุนาปิ จกฺขุมา ฯเปฯ.\csromanspacer{} เอวํ ภควา สมนฺตจกฺขุนาปิ จกฺขุมาติ \textbf{เต} โตสิตา \textbf{จกฺขุมตา}.}

\prose{\textbf{พุทฺเธนาทิจฺจพนฺธุนา}ติ พุทฺโธติ โย โส ภควา ฯเปฯ สจฺฉิกา ปญฺญตฺติ, ยทิทํ พุทฺโธติ.\csromanspacer{} \textbf{อาทิจฺจพนฺธุนา}ติ \textbf{อาทิจฺโจ} วุจฺจติ สูริโย, โส โคตโม โคตฺ\textbf{เต}น, ภควาปิ โคตโม โคตฺ\textbf{เต}น, ภควา สูริยสฺส โคตฺตญาตโก โคตฺตพนฺธุ, ตสฺมา พุทฺโธ อาทิจฺจพนฺธูติ \textbf{พุทฺเธนาทิจฺจพนฺธุนา}.}

\prose{\csromanfolio{200}\textbf{พฺรหฺมจริยมจริํสู}ติ พฺรหฺมจริยํ วุจฺจติ อสทฺธมฺมสมาปตฺติยา อารติ วิรติ ปฏิวิรติ เวรมณี วิรมณํ อกิริยา อกรณํ อนชฺฌาปตฺติ เวลา-อนติกฺกโม เสตุฆาโต.\csromanspacer{} อปิ จ นิปฺปริยายวเสน พฺรหฺมจริยํ วุจฺจติ อริโย อฏฺฐงฺคิโก มคฺโค.\csromanspacer{} เสยฺยถิทํ, สมฺมาทิฏฺฐิ สมฺมาสงฺกปฺโป สมฺมาวาจา สมฺมากมฺมนฺโต สมฺมา-อาชีโว สมฺมาวายาโม สมฺมาสติ สมฺมาสมาธิ.\csromanspacer{} \textbf{พฺรหฺมจริยมจริํสู}ติ พฺรหฺมจริยํ จริํสุ อจริํสุ สมาทาย วตฺติํสูติ พฺรหฺมจริยมจริํสุ.}

\prose{\textbf{วรปญฺญสฺส สนฺติเก}ติ วรปญฺญสฺส อคฺคปญฺญสฺส เสฏฺฐปญฺญสฺส วิเสฏฺฐปญฺญสฺส ปาโมกฺขปญฺญสฺส อุตฺตมปญฺญสฺส ปวรปญฺญสฺส.\csromanspacer{} \textbf{สนฺติเก}ติ สนฺติเก สามนฺตา อาสนฺเน อวิทูเร อุปกฏฺเฐติ วรปญฺญสฺส สนฺติเก.\csromanspacer{} เตเนตํ วุจฺจติ\csromandash{}}

\csromangathameasure{เต โตสิตา จกฺขุมตา, พุทฺเธนาทิจฺจพนฺธุนา. \\ พฺรหฺมจริยมจริํสุ, วรปญฺญสฺส สนฺติเกติ.}
\csromangathagroup{\csromangathastanza{เต โตสิตา จกฺขุมตา, พุทฺเธนาทิจฺจพนฺธุนา. \\ พฺรหฺมจริยมจริํสุ, วรปญฺญสฺส สนฺติเกติ.}}

\csromanhangingbegin
\hangingheaditem{100}{เอกเมกสฺส ปญฺหสฺส, ยถา พุทฺเธน เทสิตํ.}
\hangingbody{ตถา โย ปฏิปชฺเชยฺย, คจฺเฉ ปารํ อปารโต.}
\csromanhangingend

\prose{\textbf{เอกเมกสฺส ปญฺหสฺสา}ติ เอกเมกสฺส อชิตปญฺหสฺส เอกเมกสฺส ติสฺสเมตฺเตยฺยปญฺหสฺส ฯเปฯ เอกเมกสฺส ปิงฺคิยปญฺหสฺสาติ เอกเมกสฺส ปญฺหสฺส.}

\prose{\textbf{ยถา พุทฺเธน เทสิตนฺ}ติ พุทฺโธติ โย โส ภควา สยมฺภู ฯเปฯ สจฺฉิกา ปญฺญตฺติ, ยทิทํ พุทฺโธติ.\csromanspacer{} \textbf{ยถา พุทฺเธน เทสิตนฺ}ติ ยถา พุทฺเธน อาจิกฺขิตํ เทสิตํ ปญฺญปิตํ ปฏฺฐปิตํ วิวริตํ วิภชิตํ\footnote{วิภตฺตํ (สฺยา)} อุตฺตานีกตํ ปกาสิตนฺติ ยถา พุทฺเธน เทสิตํ.}

\prose{\textbf{ตถา โย ปฏิปชฺเชยฺยา}ติ สมฺมาปฏิปทํ อนุโลมปฏิปทํ อปจฺจนีกปฏิปทํ อนฺวตฺถปฏิปทํ ธมฺมานุธมฺมปฏิปทํ ปฏิปชฺเชยฺยาติ ตถา โย ปฏิปชฺเชยฺย.}

\prose{\textbf{คจฺเฉ ปารํ อปารโต}ติ ปารํ วุจฺจติ อมตํ นิพฺพานํ ฯเปฯ นิโรโธ นิพฺพานํ.\csromanspacer{} อปารํ วุจฺจนฺติ กิเลสา จ ขนฺธา จ อภิสงฺขารา จ.\csromanspacer{} \textbf{คจฺเฉ ปารํ}}

\csromanheader{2}{17. ปารายนตฺถุติคาถานิทฺเทส}
\csromangathameasure{อปารโตติ อปารโต ปารํ คจฺเฉยฺย, ปารํ อธิคจฺเฉยฺย, ปารํ ผสฺเสยฺย, ปารํ สจฺฉิกเรยฺยาติ คจฺเฉ ปารํ อปารโต. เตเนตํ วุจฺจติ\csromandash{} เอกเมกสฺส ปญฺหสฺส, \\ ยถา พุทฺเธน เทสิตํ. ตถา โย ปฏิปชฺเชยฺย, \\ คจฺเฉ ปารํ อปารโตติ.}
\csromangathagroup{\csromangathastanza{\csromanfolio{201}อปารโตติ อปารโต ปารํ คจฺเฉยฺย, ปารํ อธิคจฺเฉยฺย, ปารํ ผสฺเสยฺย, ปารํ สจฺฉิกเรยฺยาติ คจฺเฉ ปารํ อปารโต.\csromanspacer{} เตเนตํ วุจฺจติ\csromandash{} เอกเมกสฺส ปญฺหสฺส, \\ ยถา พุทฺเธน เทสิตํ. ตถา โย ปฏิปชฺเชยฺย, \\ คจฺเฉ ปารํ อปารโตติ.}}

\csromanhangingbegin
\hangingheaditem{101}{อปารา ปารํ คจฺเฉยฺย, ภาเวนฺโต มคฺคมุตฺตมํ.}
\hangingbody{มคฺโค โส ปารํ คมนาย, \textbf{ตสฺมา} ปารายนํ อิติ.}
\csromanhangingend

\prose{\textbf{อปารา ปารํ คจฺเฉยฺยา}ติ อปารํ วุจฺจนฺติ กิเลสา จ ขนฺธา จ อภิสงฺขารา จ.\csromanspacer{} ปารํ วุจฺจติ อมตํ นิพฺพานํ ฯเปฯ ตณฺหกฺขโย วิราโค นิโรโธ นิพฺพานํ.\csromanspacer{} \textbf{อปารา ปารํ คจฺเฉยฺยา}ติ อปารา ปารํ คจฺเฉยฺย ปารํ อธิคจฺเฉยฺย, ปารํ ผสฺเสยฺย, ปารํ สจฺฉิกเรยฺยาติ อปารา ปารํ คจฺเฉยฺย.}

\prose{\textbf{ภาเวนฺโต มคฺคมุตฺตมนฺ}ติ มคฺคมุตฺตมํ วุจฺจติ อริโย อฏฺฐงฺคิโก มคฺโค, เสยฺยถิทํ, สมฺมาทิฏฺฐิ ฯเปฯ สมฺมาสมาธิ.\csromanspacer{} \textbf{มคฺคมุตฺตมนฺ}ติ มคฺคํ อคฺคํ เสฏฺฐํ วิเสฏฺฐํ ปาโมกฺขํ อุตฺตมํ ปวรํ.\csromanspacer{} \textbf{ภาเวนฺโต}ติ ภาเวนฺโต อาเสวนฺโต พหุลีกโรนฺโตติ ภาเวนฺโต มคฺคมุตฺตมํ.}

\csromanheader{2}{\textbf{มคฺโค โส ปารํ คมนายา}ติ\csromandash{}}
\csromangathameasure{มคฺโค ปนฺโถ ปโถ ปชฺโช, อญฺชสํ วฏุมา’ยนํ. \\ นาวา อุตฺตรเสตุ จ, กุลฺโล จ ภิสิ สงฺกโม.}
\csromangathagroup{\csromangathastanza{มคฺโค ปนฺโถ ปโถ ปชฺโช\footnote{อทฺโธ (ก)}, อญฺชสํ วฏุมา’ยนํ. \\ นาวา อุตฺตรเสตุ จ, กุลฺโล จ ภิสิ สงฺกโม\footnote{อทฺโธ (ก)}.}}

\prose{\textbf{ปารํ คมนายา}ติ ปารํ คมนาย ปารํ สมฺปาปนาย ปารํ สมนุปาปนาย ชรามรณสฺส ตรณายาติ มคฺโค โส ปารํ คมนาย.}

\prose{\textbf{ตสฺมา ปารายนํ อิตี}ติ \textbf{ตสฺมา}ติ \textbf{ตสฺมา} ตํการณา ตํเหตุ ตปฺปจฺจยา ตํนิทานา.\csromanspacer{} ปารํ วุจฺจติ อมตํ นิพฺพานํ ฯเปฯ นิโรโธ นิพฺพานํ.\csromanspacer{} อยนํ วุจฺจติ มคฺโค.\csromanspacer{} \textbf{อิตี}ติ ปทสนฺธิ ฯเปฯ ปทานุปุพฺพตาเปตํ อิตีติ \textbf{ตสฺมา} ปารายนํ อิติ.\csromanspacer{} เตเนตํ วุจฺจติ\csromandash{}}

\prose{อปารา ปารํ คจฺเฉยฺย, ภาเวนฺโต มคฺคมุตฺตมํ.}

\csromangathameasure{มคฺโค โส ปารํ คมนาย, ตสฺมา ปารายนํ อิตีติ.}
\csromangathagroup{\csromangathastanza{มคฺโค โส ปารํ คมนาย, ตสฺมา ปารายนํ อิตีติ.}}

\vspace{\tipitakaparskip}
\csromancenter{ปารายนตฺถุติคาถานิทฺเทโส สตฺตรสโม.}
\csromanmatikaanchor{mtk.39}
\csromantocmarkref{subsection}{18. ปารายนานุคีติคาถานิทฺเทส}{mtk.39}
\bookmark[dest={mtk.39},level=2]{18. ปารายนานุคีติคาถานิทฺเทส}
\csromanheader{2}{18. ปารายนานุคีติคาถานิทฺเทส}
\csromanhangingbegin
\hangingheaditem{102}{\csromanfolio{202}ปารายนมนุคายิสฺสํ, (อิจฺจฺ\textbf{อายสฺมา} \textbf{ปิงฺคิโย},)}
\hangingbody{\textbf{ยถา’ทฺทกฺขิ ตถา’กฺขาสิ, วิมโล ภูริเมธโส.}}
\hangingbody{นิกฺกาโม นิพฺพโน นาโค, กิสฺส เหตุ มุสา ภเณ.}
\csromanhangingend

\prose{\textbf{ปารายนมนุคายิสฺสนฺ}ติ คีตมนุคายิสฺสํ, กถิตมนุกถยิสฺสํ\footnote{กถิตมนุคายิสฺสํ (สฺยา) เอวํ สพฺพปเทสุ อนุคายิสฺสนฺติ อาคตํ.}, ภณิตมนุภณิสฺสํ, ลปิตมนุลปิสฺสํ, ภาสิตมนุภาสิสฺสนฺติ ปารายนมนุคายิสฺสํ.\csromanspacer{} \textbf{อิจฺจายสฺมา ปิงฺคิโย}ติ \textbf{อิจฺจา}ติ ปทสนฺธิ ฯเปฯ ปทานุปุพฺพตาเปตํ \textbf{อิจฺจา}ติ.\csromanspacer{} \textbf{อายสฺมา}ติ ปิยวจนํ ครุวจนํ สคารวสปฺปติสฺสาธิวจนเมตํ \textbf{อายสฺมา}ติ.\csromanspacer{} \textbf{ปิงฺคิโย}ติ ตสฺส เถรสฺส นามํ สงฺขา สมญฺญา ปญฺญตฺติ โวหาโร นามํ นามกมฺมํ นามเธยฺยํ นิรุตฺติ พฺยญฺชนํ อภิลาโปติ อิจฺจฺ\textbf{อายสฺมา} \textbf{ปิงฺคิโย}.}

\prose{\textbf{ยถา’ทฺทกฺขิ ตถา’กฺขาสี}ติ ยถา อทฺทกฺขิ, ตถา อกฺขาสิ อาจิกฺขิ เทเสสิ ปญฺญเปสิ ปฏฺฐเปสิ วิวริ วิภชิ อุตฺตานี-อกาสิ ปกาเสสิ. “สพฺเพ สงฺขารา อนฺ\textbf{อิจฺจา}”ติ ยถา อทฺทกฺขิ, ตถา อกฺขาสิ อาจิกฺขิ เทเสสิ ปญฺญเปสิ ปฏฺฐเปสิ วิวริ วิภชิ อุตฺตานี-อกาสิ ปกาเสสิ. “สพฺเพ สงฺขารา ทุกฺขา”ติ ฯเปฯ. “สพฺเพ ธมฺมา อนตฺตา”ติ. “ยํ กิญฺจิ สมุทยธมฺมํ, สพฺพํ ตํ นิโรธธมฺมนฺ”ติ ยถา อทฺทกฺขิ, ตถา อกฺขาสิ อาจิกฺขิ เทเสสิ ปญฺญเปสิ ปฏฺฐเปสิ วิวริ วิภชิ อุตฺตานี-อกาสิ ปกาเสสีติ ยถา’ทฺทกฺขิ ตถา’กฺขาสิ.}

\prose{\textbf{วิมโล ภูริเมธโส}ติ \textbf{วิมโล}ติ ราโค มลํ, โทโส มลํ, โมโห มลํ, โกโธ, อุปนาโห ฯเปฯ สพฺพากุสลาภิสงฺขารา มลา.\csromanspacer{} เต มลา พุทฺธสฺส ภควโต ปหีนา อุจฺฉินฺนมูลา ตาลาวตฺถุกตา อนภาวํกตา อายติํ อนุปฺปาทธมฺมา.\csromanspacer{} อมโล พุทฺโธ \textbf{วิมโล} นิมฺมโล มลาปคโต มลวิปฺปหีโน มลวิมุตฺโต สพฺพมลวีติวตฺโต, ภูริ วุจฺจติ ปถวี.\csromanspacer{} ภควา ตาย\footnote{ภควา อิมาย (สฺยา)} ปถวิสมาย ปญฺญาย วิปุลาย วิตฺถตาย สมนฺนาคโต, เมธา วุจฺจติ ปญฺญา, ยา ปญฺญา ปชานนา ฯเปฯ อโมโห ธมฺมวิจโย สมฺมาทิฏฺฐิ.\csromanspacer{} ภควา อิมาย เมธาย ปญฺญาย อุเปโต}

\vspace{\tipitakaparskip}
\csromancenter{ปารายนานุคีติคาถานิทฺเทส สมุเปโต อุปาคโต สมุปาคโต อุปปนฺโน สมุปปนฺโน สมนฺนาคโต, ตสฺมา พุทฺโธ สุเมธโสติ วิมโล ภูริเมธโส.}
\prose{\csromanfolio{203}\textbf{นิกฺกาโม นิพฺพโน นาโค}ติ \textbf{กามา}ติ อุทฺทานโต เทฺว \textbf{กามา} วตฺถุ\textbf{กามา} จ กิเลส\textbf{กามา} จ ฯเปฯ.\csromanspacer{} อิเม วุจฺจนฺติ วตฺถุ\textbf{กามา} ฯเปฯ.\csromanspacer{} อิเม วุจฺจนฺติ กิเลส\textbf{กามา}.\csromanspacer{} พุทฺธสฺส ภควโต วตฺถุ\textbf{กามา} ปริญฺญาตา, กิเลส\textbf{กามา} ปหีนา.\csromanspacer{} วตฺถุ\textbf{กามา}นํ ปริญฺญาตตฺตา กิเลส\textbf{กามา}นํ ปหีนตฺตา ภควา น กาเม กาเมติ, น กาเม อิจฺฉติ, น กาเม ปตฺเถติ, น กาเม ปิเหติ, น กาเม อภิชปฺปติ.\csromanspacer{} เย กาเม กาเมนฺติ, กาเม อิจฺฉนฺติ, กาเม ปตฺเถนฺติ, กาเม ปิเหนฺติ, กาเม อภิชปฺปนฺติ, เต กามกามิโน ราคราคิโน สญฺญสญฺญิโน.\csromanspacer{} ภควา น กาเม กาเมติ, น กาเม อิจฺฉติ, น กาเม ปตฺเถติ, น กาเม ปิเหติ, น กาเม อภิชปฺปติ.\csromanspacer{} ตสฺมา พุทฺโธ อกาโม นิกฺกาโม จตฺตกาโม วนฺตกาโม มุตฺตกาโม ปหีนกาโม ปฏินิสฺสฏฺฐกาโม วีตราโค วิคตราโค จตฺตราโค วนฺตราโค มุตฺตราโค ปหีนราโค ปฏินิสฺสฏฺฐราโค นิจฺฉาโต นิพฺพุโต สีติภูโต สุขปฺปฏิสํเวที พฺรหฺมภูเตน อตฺตนา วิหรตีติ นิกฺกาโม.}

\prose{\textbf{นิพฺพโน}ติ ราโค วนํ, โทโส วนํ, โมโห วนํ, โกโธ วนํ, อุปนาโห วนํ ฯเปฯ สพฺพากุสลาภิสงฺขารา วนา.\csromanspacer{} เต วนา พุทฺธสฺส ภควโต ปหีนา อุจฺฉินฺนมูลา ตาลาวตฺถุกตา อนภาวํกตา อายติํ อนุปฺปาทธมฺมา.\csromanspacer{} ตสฺมา พุทฺโธ อวโน วิวโน นิพฺพโน วนาปคโต วนวิปฺปหีโน วนวิมุตฺโต สพฺพวนวีติวตฺโตติ นิพฺพโน.\csromanspacer{} \textbf{นาโค}ติ นาโค, ภควา อาคุํ น กโรตีติ นาโค.\csromanspacer{} น คจฺฉตีติ นาโค.\csromanspacer{} น อาคจฺฉตีติ นาโค ฯเปฯ.\csromanspacer{} เอวํ ภควา น อาคจฺฉตีติ นาโคติ นิกฺกาโม นิพฺพโน นาโค.}

\prose{\textbf{กิสฺส เหตุ มุสา ภเณ}ติ \textbf{กิสฺส เหตู}ติ กิสฺส เหตุ กิํเหตุ กิํการณา กิํนิทานา กิํปจฺจยาติ กิสฺส เหตุ.\csromanspacer{} \textbf{มุสา ภเณ}ติ มุสา ภเณยฺย กเถยฺย ทีเปยฺย โวหเรยฺย.\csromanspacer{} \textbf{มุสา ภเณ}ติ โมสวชฺชํ ภเณยฺย.\csromanspacer{} มุสาวาทํ ภเณยฺย.\csromanspacer{} อนริยวาทํ ภเณยฺย.\csromanspacer{} อิเธกจฺโจ สภาคโต\footnote{สภคฺคโต (สฺยา)} วา ปริสาคโต\footnote{ปริสคฺคโต (สฺยา)} วา ญาติมชฺฌคโต วา ปูคมชฺฌคโต วา ราชกุลมชฺฌคโต วา อภินีโต สกฺขิปุฏฺโฐ “เอหมฺโภ\footnote{เอหิ โภ (สฺยา) ปสฺส ม 3. 95 ปิฏฺเฐ.} ปุริส ยํ ชานาสิ, ตํ วเทหี”ติ โส อชานํ วา อาห \csromanfolio{204}“ชานามี”ติ, ชานํ วา อาห “น ชานามี”ติ.\csromanspacer{} อปสฺสํ วา อาห “ปสฺสามี”ติ, ปสฺสํ วา อาห “น ปสฺสามี”ติ อิติ อตฺตเหตุ วา ปรเหตุ วา อามิสกิญฺจิกฺขเหตุ วา สมฺปชานมุสา ภาสติ.\csromanspacer{} อิทํ วุจฺจติ โมสวชฺชํ.}

\prose{อปิ จ ตีหากาเรหิ มุสาวาโท โหติ, ปุพฺเพวสฺส โหติ “มุสา ภณิสฺสนฺ”ติ.\csromanspacer{} ภณนฺตสฺส โหติ “มุสา ภณามี”ติ.\csromanspacer{} ภณิตสฺส โหติ “มุสา มยา ภณิตนฺ”ติ.\csromanspacer{} อิเมหิ ตีหากาเรหิ มุสาวาโท โหติ.\csromanspacer{}\csromanspacer{}อปิ จ จตูหากาเรหิ มุสาวาโท โหติ ปุพฺเพวสฺส โหติ “มุสา ภณิสฺสนฺ”ติ.\csromanspacer{} ภณนฺตสฺส โหติ “มุสา ภณามี”ติ.\csromanspacer{} ภณิตสฺส โหติ “มุสา มยา ภณิตนฺ”ติ วินิธาย ทิฏฺฐิํ.\csromanspacer{} อิเมหิ จตูหากาเรหิ มุสาวาโท โหติ.\csromanspacer{}\csromanspacer{}อปิ จ ปญฺจหากาเรหิ.\csromanspacer{} ฉหากาเรหิ.\csromanspacer{} สตฺตหากาเรหิ.\csromanspacer{} อฏฺฐหากาเรหิ มุสาวาโท โหติ, ปุพฺเพวสฺส โหติ “มุสา ภณิสฺสนฺ”ติ.\csromanspacer{} ภณนฺตสฺส โหติ “มุสา ภณามี”ติ.\csromanspacer{} ภณิตสฺส โหติ “มุสา มยา ภณิตนฺ”ติ วินิธาย ทิฏฺฐิํ, วินิธาย ขนฺติํ, วินิธาย รุจิํ, วินิธาย สญฺญํ, วินิธาย ภาวํ.\csromanspacer{} อิเมหิ อฏฺฐหากาเรหิ มุสาวาโท โหติ.\csromanspacer{} โมสวชฺชํ กิสฺส เหตุ มุสา ภเณยฺย กถายฺย ทีเปยฺย โวหเรยฺยาติ กิสฺส เหตุ มุสา ภเณ.\csromanspacer{} เตนาห เถโร ปิงฺคิโย\csromandash{}}

\prose{ปารายนมนุคายิสฺสํ, (อิจฺจายสฺมา ปิงฺคิโย,)}

\csromangathameasure{ยถา’ทฺทกฺขิ ตถา’กฺขาสิ, วิมโล ภูริเมธโส. \\ นิกฺกาโม นิพฺพโน นาโค, กิสฺส เหตุ มุสา ภเณติ.}
\csromangathagroup{\csromangathastanza{ยถา’ทฺทกฺขิ ตถา’กฺขาสิ, วิมโล ภูริเมธโส. \\ นิกฺกาโม นิพฺพโน นาโค, กิสฺส เหตุ มุสา ภเณติ.}}

\csromanhangingbegin
\hangingheaditem{103}{ปหีนมลโมหสฺส, มานมกฺขปฺปหายิโน.}
\hangingbody{หนฺทาหํ กิตฺตยิสฺสามิ, คิรํ วณฺณูปสํหิตํ.}
\csromanhangingend

\prose{\textbf{ปหีนมลโมหสฺสา}ติ \textbf{มลนฺ}ติ ราโค มลํ, โทโส มลํ, โมโห มลํ, มาโน มลํ, ทิฏฺฐิ มลํ, กิเลโส มลํ, สพฺพทุจฺจริตํ มลํ, สพฺพภวคามิกมฺมํ มลํ.}

\prose{\textbf{โมโห}ติ ยํ ทุกฺเข อญฺญาณํ ฯเปฯ อวิชฺชาลงฺคี โมโห อกุสลมูลํ.\csromanspacer{} อยํ วุจฺจติ โมโห.\csromanspacer{} มลญฺจ โมโห จ พุทฺธสฺส ภควโต ปหีนา อุจฺฉินฺนมูลา ตาลาวตฺถุกตา อนภาวํกตา}

\vspace{\tipitakaparskip}
\csromancenter{ปารายนานุคีติคาถานิทฺเทส อายติํ อนุปฺปาทธมฺมา.\csromanspacer{} ตสฺมา พุทฺโธ ปหีนมลโมโหติ ปหีนมลโมหสฺส.}
\prose{\csromanfolio{205}\textbf{มานมกฺขปฺปหายิโน}ติ \textbf{มาโน}ติ เอกวิเธน \textbf{มาโน}, ยา จิตฺตสฺส อุนฺนติ\footnote{อุณฺณติ (สฺยา, ก)}.\csromanspacer{} ทุวิเธน \textbf{มาโน}, อตฺตุกฺกํสน\textbf{มาโน}, ปรวมฺภน\textbf{มาโน}.\csromanspacer{} ติวิเธน \textbf{มาโน}, “เสยฺโยหมสฺมี”ติ \textbf{มาโน}, “สทิโสหมสฺมี”ติ \textbf{มาโน}, “หีโนหมสฺมี”ติ \textbf{มาโน}.\csromanspacer{} จตุพฺพิเธน \textbf{มาโน}, ลาเภน มานํ ชเนติ, ยเสน มานํ ชเนติ, ปสํสาย มานํ ชเนติ, สุเขน มานํ ชเนติ.\csromanspacer{} ปญฺจวิเธน \textbf{มาโน}, “ลาภิมฺหิ มนาปิกานํ รูปานนฺ”ติ มานํ ชเนติ, “ลาภิมฺหิ มนาปิกานํ สทฺทานํ.\csromanspacer{} คนฺธานํ.\csromanspacer{} รสานํ.\csromanspacer{} โผฏฺฐพฺพานนฺ”ติ มานํ ชเนติ.\csromanspacer{} ฉพฺพิเธน \textbf{มาโน}, จกฺขุสมฺปทาย มานํ ชเนติ, โสตสมฺปทาย.\csromanspacer{} ฆานสมฺปทาย.\csromanspacer{} ชิวฺหาสมฺปทาย.\csromanspacer{} กายสมฺปทาย.\csromanspacer{} มโนสมฺปทาย มานํ ชเนติ.\csromanspacer{} สตฺตวิเธน \textbf{มาโน}, \textbf{มาโน} อติ\textbf{มาโน} มานาติ\textbf{มาโน} โอ\textbf{มาโน} อว\textbf{มาโน} อสฺมิ\textbf{มาโน} มิจฺฉา\textbf{มาโน}.\csromanspacer{} อฏฺฐวิเธน \textbf{มาโน}, ลาเภน มานํ ชเนติ, อลาเภน โอมานํ ชเนติ, ยเสน มานํ ชเนติ, อยเสน โอมานํ ชเนติ, ปสํสาย มานํ ชเนติ, นินฺทาย โอมานํ ชเนติ, สุเขน มานํ ชเนติ, ทุกฺเขน โอมานํ ชเนติ.\csromanspacer{} นววิเธน \textbf{มาโน}, เสยฺยสฺส “เสยฺโยหมสฺมี”ติ \textbf{มาโน}, เสยฺยสฺส “สทิโสหมสฺมี”ติ \textbf{มาโน}, เสยฺยสฺส “หีโนหมสฺมี”ติ \textbf{มาโน}.\csromanspacer{} สทิสสฺส “เสยฺโยหมสฺมี”ติ \textbf{มาโน}, สทิสสฺส “สทิโสหมสฺมี”ติ \textbf{มาโน}, สทิสสฺส “หีโนหมสฺมี”ติ \textbf{มาโน}.\csromanspacer{} หีนสฺส “เสยฺโยหมสฺมี”ติ \textbf{มาโน}.\csromanspacer{} หีนสฺส “สทิโสหมสฺมี”ติ \textbf{มาโน}, หีนสฺส “หีโนหมสฺมี”ติ \textbf{มาโน}.\csromanspacer{} ทสวิเธน \textbf{มาโน}, อิเธกจฺโจ มานํ ชเนติ ชาติยา วา โคตฺเตน วา โกลปุตฺติเยน วา วณฺณโปกฺขรตาย วา ธเนน วา อชฺเฌเนน วา กมฺมายตเนน วา สิปฺปายตเนน วา วิชฺชาฏฺฐาเนน\footnote{วิชฺชาฐาเนน (ก)} วา สุเตน วา ปฏิภาเนน วา อญฺญตรญฺญตเรน วา วตฺถุนา.\csromanspacer{} โย เอวรูโป \textbf{มาโน} มญฺญนา มญฺญิตตฺตํ อุนฺนติ อุนฺนโม ธโช สมฺปคฺคาโห เกตุกมฺยตา จิตฺตสฺส.\csromanspacer{} อยํ วุจฺจติ \textbf{มาโน}.}

\prose{\textbf{มกฺโข}ติ โย มกฺโข มกฺขายนา มกฺขายิตตฺตํ นิฏฺฐุริยํ นิฏฺฐุริยกมฺมํ\footnote{นิตฺถุริยกมฺมํ (ก) ปสฺส อภิ 2. 371 ปิฏฺเฐ.}.\csromanspacer{} อยํ วุจฺจติ มกฺโข.\csromanspacer{} พุทฺธสฺส ภควโต \textbf{มาโน} จ มกฺโข จ ปหีนา \csromanfolio{206}อุจฺฉินฺนมูลา ตาลาวตฺถุกตา อนภาวํกตา อายติํ อนุปฺปาทธมฺมา.\csromanspacer{} ตสฺมา \textbf{พุทฺโธ} มานมกฺขปฺปหายีติ มานมกฺขปฺปหายิโน.}

\prose{\textbf{หนฺทาหํ กิตฺตยิสฺสามิ, คิรํ วณฺณูปสญฺหิตนฺ}ติ \textbf{หนฺทาหนฺ}ติ ปทสนฺธิ ปทสํสคฺโค ปทปาริปูรี อกฺขรสมวาโย พฺยญฺชนสิลิฏฺฐตา ปทานุปุพฺพตาเปตํ \textbf{หนฺทาหนฺ}ติ.\csromanspacer{} \textbf{กิตฺตยิสฺสามิ คิรํ วณฺณูปสํหิตนฺ}ติ วณฺเณน อุเปตํ สมุเปตํ อุปาคตํ สมุปาคตํ อุปปนฺนํ สมุปปนฺนํ สมนฺนาคตํ วาจํ คิรํ พฺยปฺปถํ อุทีรณํ\footnote{โอทีรณํ (สฺยา)} กิตฺตยิสฺสามิ เทเสสฺสามิ ปญฺญเปสฺสามิ ปฏฺฐเปสฺสามิ วิวริสฺสามิ วิภชิสฺสามิ อุตฺตานีกริสฺสามิ ปกาเสสฺสามีติ หนฺทาหํ กิตฺตยิสฺสามิ, คิรํ วณฺณูปสํหิตํ.\csromanspacer{} เตนาห เถโร ปิงฺคิโย\csromandash{}}

\csromangathameasure{ปหีนมลโมหสฺส, มานมกฺขปฺปหายิโน.}
\csromangathagroup{\csromangathastanza{ปหีนมลโมหสฺส, มานมกฺขปฺปหายิโน.}}

\prose{\textbf{หนฺทาหํ กิตฺตยิสฺสามิ, คิรํ วณฺณูปสญฺหิตนฺ}ติ.}

\csromanhangingbegin
\hangingheaditem{104}{ตโมนุโท \textbf{พุทฺโธ} \textbf{สมนฺตจกฺขุ},}
\hangingbody{\textbf{โลกนฺตคู สพฺพภวาติวตฺโต.}}
\hangingbody{\textbf{อนาสโว สพฺพทุกฺขปฺปหีโน,}}
\hangingbody{สจฺจวฺหโย พฺรหฺเม อุปาสิโต เม.}
\csromanhangingend

\prose{\textbf{ตโมนุโท พุทฺโธ สมนฺตจกฺขู}ติ \textbf{ตโมนุโท}ติ ราคตมํ โทสตมํ โมหตมํ มานตมํ ทิฏฺฐิตมํ กิเลสตมํ ทุจฺจริตตมํ อนฺธกรณํ อญฺญาณกรณํ ปญฺญานิโรธิกํ วิฆาตปกฺขิกํ อนิพฺพานสํวตฺตนิกํ นุทิ ปนุทิ ปชหิ วิโนเทสิ พฺยนฺตี-อกาสิ อนภาวํ คเมสิ.\csromanspacer{} \textbf{พุทฺโธ}ติ โย โส ภควา ฯเปฯ สจฺฉิกา ปญฺญตฺติ, ยทิทํ \textbf{พุทฺโธ}ติ.\csromanspacer{} \textbf{สมนฺตจกฺขุ} วุจฺจติ สพฺพญฺญุตญาณํ ฯเปฯ ตถาคโต เตน สมนฺตจกฺขูติ \textbf{ตโมนุโท} \textbf{พุทฺโธ} \textbf{สมนฺตจกฺขุ}.}

\prose{\textbf{โลกนฺตคู สพฺพภวาติวตฺโต}ติ \textbf{โลโก}ติ เอโก \textbf{โลโก}, ภว\textbf{โลโก}.\csromanspacer{} เทฺว โลกา, ภว\textbf{โลโก} จ สมฺภว\textbf{โลโก} จ.\csromanspacer{} สมฺปตฺติภว\textbf{โลโก} จ สมฺปตฺติสมฺภว\textbf{โลโก} จ.\csromanspacer{} วิปตฺติภว\textbf{โลโก} จ วิปตฺติสมฺภว\textbf{โลโก} จ\footnote{เทฺว โลกา สมฺปตฺติ จ ภวโลโก วิปตฺติ จ ภวโลโก (สฺยา)}.\csromanspacer{} ตโย โลกา.\csromanspacer{} ติสฺโส เวทนา.\csromanspacer{} จตฺตาโร โลกา, จตฺตาโร อาหารา.\csromanspacer{} ปญฺจ โลกา, ปญฺจุปาทานกฺขนฺธา.\csromanspacer{} ฉ โลกา, ฉ อชฺฌตฺติกานิ อายตนานิ.\csromanspacer{} สตฺต โลกา,}

\vspace{\tipitakaparskip}
\csromancenter{ปารายนานุคีติคาถานิทฺเทส สตฺตวิญฺญาณฏฺฐิติโย.\csromanspacer{} อฏฺฐ โลกา, อฏฺฐ โลกธมฺมา.\csromanspacer{} นว โลกา, นว สตฺตาวาสา.\csromanspacer{} ทส โลกา, ทส อายตนานิ.\csromanspacer{} ทฺวาทส โลกา, ทฺวาทสายตนานิ.\csromanspacer{} อฏฺฐารส โลกา, อฏฺฐารส ธาตุโย.\csromanspacer{} \textbf{โลกนฺตคู}ติ ภควา โลกสฺส อนฺตคโต อนฺตปฺปตฺโต โกฏิคโต โกฏิปฺปตฺโต ฯเปฯ นิพฺพานคโต นิพฺพานปฺปตฺโต.\csromanspacer{} โส วุตฺถวาโส จิณฺณจรโณ ฯเปฯ ชาติมรณสํสาโร, นตฺถิ ตสฺส ปุนพฺภโวติ โลกนฺตคู.}
\prose{\csromanfolio{207}\textbf{สพฺพภวาติวตฺโต}ติ \textbf{ภวา}ติ เทฺว \textbf{ภวา} กมฺมภโว จ ปฏิสนฺธิโก จ ปุนพฺภโว.\csromanspacer{} กตโม กมฺมภโว.\csromanspacer{} ปุญฺญาภิสงฺขาโร อปุญฺญาภิสงฺขาโร อาเนญฺชาภิสงฺขาโร, อยํ กมฺมภโว.\csromanspacer{} กตโม ปฏิสนฺธิโก ปุนพฺภโว.\csromanspacer{} ปฏิสนฺธิกา รูปา เวทนา สญฺญา สงฺขารา วิญฺญาณํ, อยํ ปฏิสนฺธิโก ปุนพฺภโว.\csromanspacer{} ภควา กมฺมภวญฺจ ปฏิสนฺธิกญฺจ ปุนพฺภวํ อติวตฺโต\footnote{อุปาติวตฺโต (ก)} อติกฺกนฺโต วีติวตฺโตติ โลกนฺตคู สพฺพ\textbf{ภวา}ติวตฺโต.}

\prose{\textbf{อนาสโว สพฺพทุกฺขปฺปหีโน}ติ \textbf{อนาสโว}ติ จตฺตาโร อาสวา กามาสโว \textbf{ภวา}สโว ทิฏฺฐาสโว อวิชฺชาสโว.\csromanspacer{} เต อาสวา พุทฺธสฺส ภควโต ปหีนา อุจฺฉินฺนมูลา ตาลาวตฺถุกตา อนภาวํกตา อายติํ อนุปฺปาทธมฺมา.\csromanspacer{} ตสฺมา พุทฺโธ \textbf{อนาสโว}.\csromanspacer{} \textbf{สพฺพทุกฺขปฺปหีโน}ติ สพฺพํ ตสฺส ปฏิสนฺธิกํ ชาติทุกฺขํ ชราทุกฺขํ พฺยาธิทุกฺขํ มรณทุกฺขํ โสกปริเทวทุกฺขโทมนสฺสุปายาสทุกฺขํ ฯเปฯ ทิฏฺฐิพฺยสนทุกฺขํ ปหีนํ สมุจฺฉินฺนํ วูปสนฺตํ ปฏิปฺปสฺสทฺธํ อภพฺพุปฺปตฺติกํ ญาณคฺคินา ทฑฺฒํ.\csromanspacer{} ตสฺมา พุทฺโธ สพฺพทุกฺขปฺปหีโนติ \textbf{อนาสโว} สพฺพทุกฺขปฺปหีโน.}

\prose{\textbf{สจฺจวฺหโย พฺรหฺเม อุปาสิโต เม}ติ \textbf{สจฺจวฺหโย}ติ \textbf{สจฺจวฺหโย} สทิสนาโม สทิสวฺหโย สจฺจสทิสวฺหโย.\csromanspacer{} วิปสฺสี ภควา สิขี ภควา เวสฺสภู ภควา กกุสนฺโธ ภควา โกณาคมโน ภควา กสฺสโป ภควา.\csromanspacer{} เต พุทฺธา ภควนฺโต สทิสนามา สทิสวฺหยา.\csromanspacer{} ภควาปิ สกฺยมุนิ เตสํ พุทฺธานํ ภควนฺตานํ สทิสนาโม สทิสวฺหโยติ ตสฺมา พุทฺโธ \textbf{สจฺจวฺหโย}.}

\prose{\textbf{พฺรหฺเม อุปาสิโต เม}ติ โส มยา ภควา อาสิโต อุปาสิโต ปยิรุปาสิโต ปริปุจฺฉิโต ปริปญฺหิโตติ \textbf{สจฺจวฺหโย} พฺรหฺเม อุปาสิโต เม.\csromanspacer{} เตนาห เถโร ปิงฺคิโย\csromandash{}}

\csromangathameasure{ตโมนุโท พุทฺโธ สมนฺตจกฺขุ, \\ โลกนฺตคู สพฺพภวาติวตฺโต. \\ อนาสโว สพฺพทุกฺขปฺปหีโน, \\ สจฺจวฺหโย พฺรหฺเม อุปาสิโต เมติ.}
\csromangathagroup{\csromangathastanza{\csromanfolio{208}ตโมนุโท พุทฺโธ สมนฺตจกฺขุ, \\ โลกนฺตคู สพฺพภวาติวตฺโต. \\ อนาสโว สพฺพทุกฺขปฺปหีโน, \\ สจฺจวฺหโย พฺรหฺเม อุปาสิโต เมติ.}}

\csromanhangingbegin
\hangingheaditem{105}{ทิโช ยถา กุพฺพนกํ ปหาย,}
\hangingbody{\textbf{พหุปฺผลํ กานนมาวเสยฺย.}}
\hangingbody{\textbf{เอวมหํ อปฺปทสฺเส ปหาย,}}
\hangingbody{\textbf{มโหทธิํ หํโสริว อชฺฌปตฺโต1}.}
\csromanhangingend

\prose{\textbf{ทิโช ยถา กุพฺพนกํ ปหาย, พหุปฺผลํ กานนมาวเสยฺยา}ติ ทิโช วุจฺจติ ปกฺขี.\csromanspacer{} กิํการณา ทิโช วุจฺจติ ปกฺขี, ทฺวิกฺขตฺตุํ ชายตีติ ทิโช มาตุกุจฺฉิมฺหา จ อณฺฑโกสมฺหา จ.\csromanspacer{} ตํการณา ทิโช วุจฺจติ ปกฺขีติ ทิโช.\csromanspacer{} \textbf{ยถา กุพฺพนกํ ปหายา}ติ ยถา ทิโช กุพฺพนกํ ปริตฺตวนกํ อปฺปผลํ อปฺปภกฺขํ อปฺโปทกํ ปหาย ชหิตฺวา อติกฺกมิตฺวา สมติกฺกมิตฺวา วีติวตฺเตตฺวา อญฺญํ พหุปฺผลํ พหุภกฺขํ พหูทกํ\footnote{พหุรุกฺขํ (สฺยา)} มหนฺตํ กานนํ วนสณฺฑํ อธิคจฺเฉยฺย วินฺเทยฺย ปฏิลเภยฺย ตสฺมิํ จ วนสณฺเฑ วาสํ กปฺเปยฺยาติ ทิโช ยถา กุพฺพนกํ ปหาย, พหุปฺผลํ กานนํ อาวเสยฺย.}

\prose{\textbf{เอวมหํ อปฺปทสฺเส ปหาย, มโหทธิํ หํโสริว อชฺฌปตฺโต}ติ \textbf{เอวนฺ}ติ โอปมฺมสมฺปฏิปาทนํ.\csromanspacer{} \textbf{อปฺปทสฺเส ปหายา}ติ โย จ พาวรี พฺราหฺมโณ เย จญฺเญ ตสฺส อาจริยา พุทฺธํ ภควนฺตํ อุปาทาย อปฺปทสฺสา ปริตฺตทสฺสา โถกทสฺสา โอมกทสฺสา ลามกทสฺสา ฉตุกฺกทสฺสา\footnote{ชตุกฺกทสฺสา (สฺยา), ชตุกทสฺสา (สี-ฏฺฐ)} วา.\csromanspacer{} เต อปฺปทสฺเส ปริตฺตทสฺเส โถกทสฺเส โอมกทสฺเส ลามกทสฺเส ฉตุกฺกทสฺเส ปหาย ปชหิตฺวา อติกฺกมิตฺวา สมติกฺกมิตฺวา วีติวตฺเตตฺวา พุทฺธํ ภควนฺตํ อปฺปมาณทสฺสํ อคฺคทสฺสํ เสฏฺฐทสฺสํ วิเสฏฺฐทสฺสํ ปาโมกฺขทสฺสํ อุตฺตมทสฺสํ ปวรทสฺสํ อสมํ อสมสมํ อปฺปฏิสมํ อปฺปฏิภาคํ อปฺปฏิปุคฺคลํ เทวาติเทวํ นราสภํ ปุริสสีหํ ปุริสนาคํ ปุริสาชญฺญํ ปุริสนิสภํ ปุริสโธรยฺหํ ทสพลธาริํ\footnote{ทสพลํ ตาทิํ (สฺยา)} อธิคจฺฉิํ วินฺทิํ ปฏิลภิํ.\csromanspacer{} ยถา จ หํโส}

\vspace{\tipitakaparskip}
\csromancenter{ปารายนานุคีติคาถานิทฺเทส มหนฺตํ มานสกํ\footnote{มานุสกตํ (สฺยา)} วา สรํ อโนตตฺตํ วา ทหํ มหาสมุทฺทํ วา อกฺโขภํ อมิโตทกํ ชลราสิํ อธิคจฺเฉยฺย วินฺเทยฺย ปฏิลเภยฺย, เอวเมว พุทฺธํ ภควนฺตํ อกฺโขภํ อมิตเตชํ ปภินฺนญาณํ วิวฏจกฺขุํ ปญฺญาปเภทกุสลํ อธิคตปฏิสมฺภิทํ จตุเวสารชฺชปฺปตฺตํ สุทฺธาธิมุตฺตํ เสตปจฺจตฺตํ อทฺวยภาณิํ ตาทิํ ตถาปฏิญฺญํ อปริตฺตํ มหนฺตํ คมฺภีรํ อปฺปเมยฺยํ ทุปฺปริโยคาหํ ปหูตรตนํ สาครสมํ ฉฬงฺคุเปกฺขาย สมนฺนาคตํ อตุลํ วิปุลํ อปฺปเมยฺยํ, ตํ ตาทิสํ ปวทตํ มคฺควาทินํ\footnote{ปวรมคฺควาทินํ (ก)} เมรุมิว นคานํ ครุฬมิว ทิชานํ สีหมิว มิคานํ อุทธิมิว อณฺณวานํ อธิคจฺฉิํ, ตํ สตฺถารํ ชินปวรํ มเหสินฺติ เอวมหํ อปฺปทสฺเส ปหาย, มโหทธิํ หํโสริว อชฺฌปตฺโต.\csromanspacer{} เตนาห เถโร ปิงฺคิโย\csromandash{}}
\vspace{0.5\baselineskip}
\csromangathameasure{ทิโช ยถา กุพฺพนกํ ปหาย, \\ พหุปฺผลํ กานนมาวเสยฺย. \\ เอวมหํ อปฺปทสฺเส ปหาย, \\ มโหทธิํ หํโสริว อชฺฌปตฺโตติ.}
\csromangathagroup{\csromangathastanza{\csromanfolio{209}ทิโช ยถา กุพฺพนกํ ปหาย, \\ พหุปฺผลํ กานนมาวเสยฺย. \\ เอวมหํ อปฺปทสฺเส ปหาย, \\ มโหทธิํ หํโสริว อชฺฌปตฺโตติ.}}

\csromanhangingbegin
\hangingheaditem{106}{เย เม ปุพฺเพ วิยากํสุ,}
\hangingbody{\textbf{หุรํ โคตมสาสนา, “อิจฺจาสิ อิติ ภวิสฺสติ.”}}
\hangingbody{\textbf{สพฺพํ ตํ อิติหีติหํ, สพฺพํ ตํ ตกฺกวฑฺฒนํ. (5)}}
\csromanhangingend

\prose{\textbf{เย เม ปุพฺเพ วิยากํสู}ติ \textbf{เย}ติ โย จ พาวรี พฺราหฺมโณ \textbf{เย} จญฺเญ ตสฺส อาจริยา, เต สกํ ทิฏฺฐิํ สกํ ขนฺติํ สกํ รุจิํ สกํ ลทฺธิํ สกํ อชฺฌาสยํ สกํ อธิปฺปายํ พฺยากํสุ อาจิกฺขิํสุ เทสยิํสุ ปญฺญปิํสุ ปฏฺฐปิํสุ วิวริํสุ วิภชิํสุ อุตฺตานี-อกํสุ ปกาเสสุนฺติ \textbf{เย} เม ปุพฺเพ วิยากํสุ.}

\prose{\textbf{หุรํ โคตมสาสนา}ติ หุรํ โคตมสาสนา ปรํ โคตมสาสนา ปุเร โคตมสาสนา ปฐมตรํ โคตมสาสนา พุทฺธสาสนา ชินสาสนา ตถาคตสาสนา\footnote{ตถาคตสาสนา เทวสาสนา (สฺยา, ก)} อรหนฺตสาสนาติ หุรํ โคตมสาสนา.}

\prose{\textbf{อิจฺจาสิ อิติ ภวิสฺสตี}ติ เอวํ กิร อาสิ, เอวํ กิร ภวิสฺสตีติ อิจฺจาสิ อิติ ภวิสฺสติ.}

\prose{\csromanfolio{210}\textbf{สพฺพํ ตํ อิติหีติหนฺ}ติ สพฺพํ ตํ อิติหีติหํ อิติกิราย ปรมฺปราย ปิฏกสมฺปทาย ตกฺกเหตุ นยเหตุ อาการปริวิตกฺเกน ทิฏฺฐินิชฺฌานกฺขนฺติยา น สามํ สยมภิญฺญาตํ น อตฺตปจฺจกฺขํ ธมฺมํ ยํ กถยิํสูติ สพฺพํ ตํ อิติหีติหํ.}

\prose{\textbf{สพฺพํ ตํ ตกฺกวฑฺฒนนฺ}ติ สพฺพํ ตํ ตกฺกวฑฺฒนํ วิตกฺกวฑฺฒนํ สงฺกปฺปวฑฺฒนํ กามวิตกฺกวฑฺฒนํ พฺยาปาทวิตกฺกวฑฺฒนํ วิหิํสาวิตกฺกวฑฺฒนํ ญาติวิตกฺกวฑฺฒนํ ชนปทวิตกฺกวฑฺฒนํ อมราวิตกฺกวฑฺฒนํ ปรานุทยตาปฏิสํยุตฺตวิตกฺกวฑฺฒนํ ลาภสกฺการสิโลก- ปฏิสํยุตฺตวิตกฺกวฑฺฒนํ อนวญฺญตฺติปฏิสํยุตฺตวิตกฺกวฑฺฒนนฺติ สพฺพํ ตํ ตกฺกวฑฺฒนํ.\csromanspacer{} เตนาห เถโร ปิงฺคิโย\csromandash{}}

\prose{เย เม ปุพฺเพ วิยากํสุ,}

\csromangathameasure{หุรํ โคตมสาสนา, “อิจฺจาสิ อิติ ภวิสฺสติ.” \\ สพฺพํ ตํ อิติหีติหํ, สพฺพํ ตํ ตกฺกวฑฺฒนนฺติ.}
\csromangathagroup{\csromangathastanza{หุรํ โคตมสาสนา, “อิจฺจาสิ อิติ ภวิสฺสติ.” \\ สพฺพํ ตํ อิติหีติหํ, สพฺพํ ตํ ตกฺกวฑฺฒนนฺติ.}}

\csromanhangingbegin
\hangingheaditem{107}{\textbf{เอโก ตโมนุทาสีโน}, ชุติมา โส ปภงฺกโร.}
\hangingbody{โคตโม ภูริปญฺญาโณ, โคตโม ภูริเมธโส.}
\csromanhangingend

\prose{\textbf{เอโก ตโมนุทาสีโน}ติ \textbf{เอโก}ติ ภควา ปพฺพชฺชสงฺขาเตน \textbf{เอโก}, อทุติยฏฺเฐน \textbf{เอโก}, ตณฺหาย ปหานฏฺเฐน \textbf{เอโก}, เอกนฺตวีตราโคติ \textbf{เอโก}, เอกนฺตวีตโทโสติ \textbf{เอโก}, เอกนฺตวีตโมโหติ \textbf{เอโก}, เอกนฺตนิกฺกิเลโสติ \textbf{เอโก}, เอกายนมคฺคํ คโตติ \textbf{เอโก}, \textbf{เอโก} อนุตฺตรํ สมฺมาสมฺโพธิํ อภิสมฺพุทฺโธติ \textbf{เอโก}.}

\prose{กถํ ภควา ปพฺพชฺชสงฺขาเตน \textbf{เอโก}, ภควา ทหโรว สมาโน สุสุ กาฬเกโส ภเทฺรน โยพฺพเนน สมนฺนาคโต ปฐเมน วยสา อกามกานํ มาตาปิตูนํ อสฺสุมุขานํ โรทนฺตานํ วิลปนฺตานํ ญาติสํฆํสพฺพํ ฆราวาสปลิโพธํ ฉินฺทิตฺวา ปุตฺตทารปลิโพธํ ฉินฺทิตฺวา ญาติปลิโพธํ ฉินฺทิตฺวา มิตฺตามจฺจปลิโพธํ ฉินฺทิตฺวา สนฺนิธิปลิโพธํ ฉินฺทิตฺวา เกสมสฺสุํ โอหาเรตฺวา กาสายานิ วตฺถานิ อจฺฉาเทตฺวา อคารสฺมา อนคาริยํ ปพฺพชิตฺวา อกิญฺจนภาวํ อุปคนฺตฺวา \textbf{เอโก} จรติ วิหรติ อิริยติ วตฺเตติ ปาเลติ ยเปติ ยาเปติ.\csromanspacer{} เอวํ ภควา ปพฺพชฺชสงฺขาเตน \textbf{เอโก}.}

\prose{กถํ ภควา อทุติยฏฺเฐน \textbf{เอโก}, เอวํ ปพฺพชิโต สมาโน \textbf{เอโก} อรญฺญวนปตฺถานิ ปนฺตานิ เสนาสนานิ ปฏิเสวติ อปฺปสทฺทานิ}

\vspace{\tipitakaparskip}
\csromancenter{ปารายนานุคีติคาถานิทฺเทส อปฺปนิคฺโฆสานิ วิชนวาตานิ มนุสฺสราหสฺเสยฺยกานิ\footnote{มนุสฺสราหเสยฺยกานิ (สฺยา, ก)} ปฏิสลฺลานสารุปฺปานิ\footnote{ปฏิสลฺลาณสารุปฺปานิ (ก)}, โส เอโก คจฺฉติ, เอโก ติฏฺฐติ, เอโก นิสีทติ, เอโก เสยฺยํ กปฺเปติ, เอโก คามํ ปิณฺฑาย ปวิสติ, เอโก อภิกฺกมติ, เอโก ปฏิกฺกมติ, เอโก รโห นิสีทติ, เอโก จงฺกมํ อธิฏฺฐาติ, เอโก จรติ วิหรติ อิริยติ วตฺเตติ ปาเลติ ยเปติ ยาเปติ เอวํ ภควา อทุติยฏฺเฐน เอโก.}
\prose{\csromanfolio{211}กถํ ภควา ตณฺหาย ปหานฏฺเฐน เอโก, โส เอวํ เอโก อทุติโย อปฺปมตฺโต อาตาปี ปหิตตฺโต วิหรนฺโต นชฺชา เนรญฺชราย ตีเร โพธิรุกฺขมูเล มหาปธานํ ปทหนฺโต มารํ สเสนํ กณฺหํ นมุจิํ ปมตฺตพนฺธุํ วิธมิตฺวา ตณฺหาชาลินิํ\footnote{ตณฺหํ ชาลินิํ (สฺยา)} วิสฏํ\footnote{สริตํ (สฺยา) ขุ 7. 362 ปิฏฺเฐปิ.} วิสตฺติกํ ปชหิ วิโนเทสิ พฺยนฺตี-อกาสิ อนภาวํ คเมสิ.}

\csromangathameasure{“ตณฺหาทุติโย ปุริโส, ทีฆมทฺธาน สํสรํ. \\ อิตฺถภาวญฺญถาภาวํ, สํสารํ นาติวตฺตติ. \\ เอตมาทีนวํ ญ ตฺวา, ตณฺหํ ทุกฺขสฺส สมฺภวํ. \\ วีตตโณฺห อนาทาโน, สโต ภิกฺขุ ปริพฺพเช”ติ.}
\csromangathagroup{\csromangathastanza{“ตณฺหาทุติโย ปุริโส, ทีฆมทฺธาน สํสรํ. \\ อิตฺถภาวญฺญถาภาวํ, สํสารํ นาติวตฺตติ.}\csromangathastanza{เอตมาทีนวํ\footnote{เอวมาทีนวํ (ก) ปสฺส ขุ 1. 201 ปิฏฺเฐ.} ญ ตฺวา, ตณฺหํ\footnote{ตณฺหา (สฺยา, ก) ขุ 7. 362 ปิฏฺเฐ.} ทุกฺขสฺส สมฺภวํ. \\ วีตตโณฺห อนาทาโน, สโต ภิกฺขุ ปริพฺพเช”ติ.}}

\prose{เอวํ ภควา ตณฺหาย ปหานฏฺเฐน เอโก.}

\prose{กถํ ภควา เอกนฺตวีตราโคติ เอโก, ราคสฺส ปหีนตฺตา เอกนฺตวีตราโคติ เอโก, โทสสฺส ปหีนตฺตา เอกนฺตวีตโทโสติ เอโก, โมหสฺส ปหีนตฺตา เอกนฺตาวีตโมโหติ เอโก, กิเลสานํ ปหีนตฺตา เอกนฺตนิกฺกิเลโสติ เอโก.}

\prose{กถํ ภควา เอกายนมคฺคํ คโตติ เอโก, เอกายนมคฺโค วุจฺจติ จตฺตาโร สติปฏฺฐานา ฯเปฯ อริโย อฏฺฐงฺคิโก มคฺโค.}

\csromangathameasure{“เอกายนํ ชาติขยนฺตทสฺสี, \\ มคฺคํ ปชานาติ หิตานุกมฺปี. \\ เอเตน มคฺเคน ตริํสุ ปุพฺเพ, \\ ตริสฺสนฺติ เย จ ตรนฺติ โอฆนฺ”ติ.}
\csromangathagroup{\csromangathastanza{\csromanfolio{212}“เอกายนํ ชาติขยนฺตทสฺสี, \\ มคฺคํ ปชานาติ หิตานุกมฺปี. \\ เอเตน มคฺเคน ตริํสุ\footnote{อตริํสุ (ก) ปสฺส สํ 3. 162 ปิฏฺเฐ, ขุ 7. 362 ปิฏฺเฐ จ.} ปุพฺเพ, \\ ตริสฺสนฺติ เย จ ตรนฺติ โอฆนฺ”ติ.}}

\prose{เอวํ ภควา เอกายนมคฺคํ คโตติ เอโก.}

\prose{กถํ ภควา เอโก อนุตฺตรํ สมฺมาสมฺโพธิํ อภิสมฺพุทฺโธติ เอโก.\csromanspacer{} โพธิ วุจฺจติ จตูสุ มคฺเคสุ ญาณํ, ปญฺญา ปญฺญินฺทฺริยํ ปญฺญาพลํ ธมฺมวิจยสมฺโพชฺฌงฺโค วีมํสา วิปสฺสนา สมฺมาทิฏฺฐิ, ภควา เตน โพธิญาเณน “สพฺเพ สงฺขารา อนิจฺจา”ติ พุชฺฌิ, “สพฺเพ สงฺขารา ทุกฺขา”ติ พุชฺฌิ ฯเปฯ “สพฺเพ ธมฺมา อนตฺตา”ติ พุชฺฌิ,. “ยํ กิญฺจิ สมุทยธมฺมํ, สพฺพํ ตํ นิโรธธมฺมนฺ”ติ พุชฺฌิ.\csromanspacer{} อถ วา ยํ พุชฺฌิตพฺพํ อนุพุชฺฌิตพฺพํ ปฏิพุชฺฌิตพฺพํ สมฺพุชฺฌิตพฺพํ อธิคนฺตพฺพํ ผสฺสิตพฺพํ สจฺฉิกาตพฺพํ, สพฺพํ ตํ เตน โพธิญาเณน พุชฺฌิ อนุพุชฺฌิ ปฏิพุชฺฌิ สมฺพุชฺฌิ อธิคจฺฉิ ผสฺเสสิ สจฺฉากาสิ.\csromanspacer{} เอวํ ภควา เอโก อนุตฺตรํ สมฺมาสมฺโพธิํ อภิสมฺพุทฺโธติ เอโก.}

\prose{\textbf{ตโมนุโท}ติ ภควา ราคตมํ โทสตมํ โมหตมํ ทิฏฺฐิตมํ กิเลสตมํ ทุจฺจริตตมํ อนฺธกรณํ อจกฺขุกรณํ อญฺญาณกรณํ ปญฺญานิโรธิกํ วิฆาตปกฺขิกํ อนิพฺพานสํวตฺตนิกํ นุทิ ปนุทิ ปชหิ วิโนเทสิ พฺยนฺตี-อกาสิ อนภาวํ คเมสิ.\csromanspacer{} \textbf{อาสีโน}ติ นิสินฺโน ภควา ปาสาณเก เจติเยติ อาสีโน\footnote{อาสิโน (ก)}.}

\csromangathameasure{นคสฺส ปสฺเส อาสีนํ, มุนิํ ทุกฺขสฺส ปารคุํ. \\ สาวกา ปยิรุปาสนฺติ, เตวิชฺชา มจฺจุหายิโนติ.}
\csromangathagroup{\csromangathastanza{นคสฺส ปสฺเส อาสีนํ, มุนิํ ทุกฺขสฺส ปารคุํ. \\ สาวกา ปยิรุปาสนฺติ, เตวิชฺชา มจฺจุหายิโนติ.}}

\prose{เอวมฺปิ ภควา อาสีโน.\csromanspacer{}\csromanspacer{}อถ วา ภควา สพฺโพสฺสุกฺกปฏิปฺปสฺสทฺธตฺตา อาสีโน, โส วุตฺถวาโส จิณฺณจรโณ ฯเปฯ ชาติมรณสํสาโร, นตฺถิ ตสฺส ปุนพฺภโวติ.\csromanspacer{} เอวมฺปิ ภควา อาสีโนติ เอโก ตโมนุทาสีโน.}

\prose{\textbf{ชุติมา โส ปภงฺกโร}ติ \textbf{ชุติมา}ติ \textbf{ชุติมา} มติมา ปณฺฑิโต ปญฺญวา พุทฺธิมา ญาณี วิภาวี เมธาวิ.\csromanspacer{} \textbf{ปภงฺกโร}ติ ปภงฺกโร อาโลกกโร โอภาสกโร ทีปงฺกโร ปทีปกโร อุชฺโชตกโร ปชฺโชตกโรติ \textbf{ชุติมา} โส ปภงฺกโร.}

\csromanheader{2}{18. ปารายนานุคีติคาถานิทฺเทส}
\prose{\csromanfolio{213}\textbf{โคตโม ภูริปญฺญาโณ}ติ \textbf{โคตโม ภูริปญฺญาโณ} ญาณปญฺญาโณ ปญฺญาธโช ปญฺญาเกตุ ปญฺญาธิปเตยฺ\textbf{โย} วิจยพหุโล ปวิจยพหุโล โอกฺขายนพหุโล สโมกฺขายนธมฺโม วิภูตวิหารี ตจฺจริโต ตพฺพหุโล ตคฺครุโก ตนฺนินฺโน ตปฺโปโณ ตปฺปพฺภาโร ตทธิมุตฺโต ตทธิปเตยฺ\textbf{โย}.}

\csromangathameasure{ธโช รถสฺส ปญฺญาณํ, ธูโม ปญฺญาณมคฺคิโน. \\ ราชา รฏฺฐสฺส ปญฺญาณํ, ภตฺตา ปญฺญาณมิตฺถิยาติ.}
\csromangathagroup{\csromangathastanza{ธโช รถสฺส ปญฺญาณํ, ธูโม\footnote{ธุโม (สฺยา)} ปญฺญาณมคฺคิโน. \\ ราชา รฏฺฐสฺส ปญฺญาณํ, ภตฺตา ปญฺญาณมิตฺถิยาติ.}}

\prose{เอวเมว \textbf{โคตโม ภูริปญฺญาโณ} ญาณปญฺญาโณ ปญฺญาธโช ปญฺญาเกตุ ปญฺญาธิปเตยฺ\textbf{โย} วิจยพหุโล ปวิจยพหุโล โอกฺขายนพหุโล สโมกฺขายนธมฺโม วิภูตวิหารี ตจฺจริโต ตพฺพหุโล ตคฺครูโก ตนฺนินฺโน ตปฺโปโณ ตปฺปพฺภาโร ตทธิมุตฺโต ตทธิปเตยฺ\textbf{โย}ติ \textbf{โคตโม ภูริปญฺญาโณ}.}

\prose{\textbf{โคตโม ภูริเมธโส}ติ ภูริ วุจฺจติ ปถวี, ภควา ตาย ปถวิสมาย ปญฺญาย วิปุลาย วิตฺถตาย สมนฺนาคโต.\csromanspacer{} เมธา วุจฺจติ ปญฺญา, ยา ปญฺญา ปชานนา ฯเปฯ อโมโห ธมฺมวิจ\textbf{โย} สมฺมาทิฏฺฐิ.\csromanspacer{} ภควา อิมาย เมธาย อุเปโต สมุเปโต อุปาคโต สมุปาคโต อุปปนฺโน สมุปปนฺโน สมนฺนาคโต, ตสฺมา พุทฺโธ สุเมธโสติ\footnote{ภูริเมธโสติ (สฺยา) เอวมุปริปิ.} \textbf{โคตโม ภูริเมธโส}.\csromanspacer{} เตนาห เถโร ปิงฺคิ\textbf{โย}\csromandash{}}

\csromangathameasure{เอโก ตโมนุทาสีโน, ชุติมา โส ปภงฺกโร. \\ \textbf{โคตโม ภูริปญฺญาโณ}, \textbf{โคตโม ภูริเมธโส}ติ.}
\csromangathagroup{\csromangathastanza{เอโก ตโมนุทาสีโน, ชุติมา โส ปภงฺกโร. \\ \textbf{โคตโม ภูริปญฺญาโณ}, \textbf{โคตโม ภูริเมธโส}ติ.}}

\csromanhangingbegin
\hangingheaditem{108}{โย เม ธมฺมมเทเสสิ, สนฺทิฏฺฐิกมกาลิกํ.}
\hangingbody{ตณฺหกฺขยมนีติกํ, ยสฺส นตฺถิ อุปมา กฺวจิ.}
\csromanhangingend

\prose{\textbf{โย เม ธมฺมมเทเสสี}ติ \textbf{โย}ติ \textbf{โย} โส ภควา สยมฺภู อนาจริยโก ปุพฺเพ อนนุสฺสุเตสุ ธมฺเมสุ สามํ สจฺจานิ อภิสมฺพุชฺฌิ, ตตฺถ จ สพฺพญฺญุตํ ปตฺโต พเลสุ จ วสีภาวํ.\csromanspacer{} \textbf{ธมฺมมเทเสสี}ติ \textbf{ธมฺมนฺ}ติ อาทิกลฺยาณํ มชฺเฌกลฺยาณํ ปริ\textbf{โย}สานกลฺยาณํ สาตฺถํ สพฺยญฺชนํ เกวลปริปุณฺณํ ปริสุทฺธํ พฺรหฺมจริยํ จตฺตาโร สติปฏฺฐเน ฯเปฯ อริยํ อฏฺฐงฺคิกํ มคฺคํ นิพฺพานญฺจ นิพฺพานคามินิญฺจ ปฏิปทํ อาจิกฺขิ เทเสสิ ปญฺญเปสิ \csromanfolio{214}ปฏฺฐเปสิ วิวริ วิภชิ อุตฺตานี-อกาสิ ปกาเสสฺ\textbf{อีติ} โย เม ธมฺมมเทเสสิ.}

\prose{\textbf{สนฺทิฏฺฐิกมกาลิกนฺ}ติ สนฺทิฏฺฐิกํ อกาลิกํ เอหิปสฺสิกํ โอปเนยฺยิกํ ปจฺจตฺตํ เวทิตพฺพํ วิญฺญูหฺ\textbf{อีติ} เอวํ สนฺทิฏฺฐิกํ.\csromanspacer{} อถ วา โย ทิฏฺเฐว ธมฺเม อริยํ อฏฺฐงฺคิกํ มคฺคํ ภาเวติ, ตสฺส มคฺคสฺส อนนฺตรา สมนนฺตรา อธิคจฺฉเตว ผลํ วินฺทติ ปฏิลภตฺ\textbf{อีติ} เอวมฺปิ สนฺทิฏฺฐิกมกาลิกํ.\csromanspacer{} ยถา มนุสฺสา กาลิกํ ธนํ ทตฺวา อนนฺตรา น ลภนฺติ กาลํ อาคเมนฺติ, เนวายํ ธมฺโม.\csromanspacer{} โย ทิฏฺเฐว ธมฺเม อริยํ อฏฺฐงฺคิกํ มคฺคํ ภาเวติ, ตสฺส มคฺคสฺส อนนฺตรา สมนนฺตรา อธิคจฺฉเตว ผลํ วินฺทติ ปฏิลภติ, น ปรตฺถ น ปรโลเก, เอวํ อกาลิกนฺติ สนฺทิฏฺฐิกมกาลิกํ.}

\prose{\textbf{ตณฺหกฺขยมนีติกนฺ}ติ \textbf{ตณฺหา}ติ รูป\textbf{ตณฺหา} ฯเปฯ ธมฺม\textbf{ตณฺหา}.\csromanspacer{} \textbf{ตณฺหกฺขยนฺ}ติ ตณฺหกฺขยํ ราคกฺขยํ โทสกฺขยํ โมหกฺขยํ คติกฺขยํ อุปปตฺติกฺขยํ ปฏิสนฺธิกฺขยํ ภวกฺขยํ สํสารกฺขยํ วฏฺฏกฺขยํ.\csromanspacer{} \textbf{อนีติกนฺ}ติ \textbf{อีติ}วุจฺจนฺติ กิเลสา จ ขนฺธา จ อภิสงฺขารา จ. ¢ติปฺปหานํ \textbf{อีติ}วูปสมํ \textbf{อีติ}ปฏินิสฺสคฺคํ \textbf{อีติ}ปฏิปฺปสฺสทฺธิํ อมตํ นิพฺพานนฺติ ตณฺหกฺขยมนฺ\textbf{อีติ}กํ.}

\prose{\textbf{ยสฺส นตฺถิ อุปมา กฺวจี}ติ \textbf{ยสฺสา}ติ นิพฺพานสฺส.\csromanspacer{} \textbf{นตฺถิ อุปมา}ติ อุปมา นตฺถิ, อุปนิธา นตฺถิ, สทิสํ นตฺถิ, ปฏิภาโค นตฺถิ น อตฺถิ น สํวิชฺชติ นุปลพฺภติ.\csromanspacer{} \textbf{กฺวจี}ติ กฺวจิ กิมฺหิจิ กตฺถจิ อชฺฌตฺตํ วา พหิทฺธา วา อชฺฌตฺตพหิทฺธา วาติ ยสฺส นตฺถิ อุปมา กฺวจิ.\csromanspacer{} เตนาห เถโร ปิงฺคิโย\csromandash{}}

\csromangathameasure{โย เม ธมฺมมเทเสสิ, สนฺทิฏฺฐิกมกาลิกํ. \\ ตณฺหกฺขยมนฺ\textbf{อีติ}กํ, ยสฺส นตฺถิ อุปมา กฺวจฺ\textbf{อีติ}.}
\csromangathagroup{\csromangathastanza{โย เม ธมฺมมเทเสสิ, สนฺทิฏฺฐิกมกาลิกํ. \\ ตณฺหกฺขยมนฺ\textbf{อีติ}กํ, ยสฺส นตฺถิ อุปมา กฺวจฺ\textbf{อีติ}.}}

\csromanhangingbegin
\hangingheaditem{109}{กิํ นุ ตมฺหา วิปฺปวสสิ, มุหุตฺตมปิ ปิงฺคิย.}
\hangingbody{โคตมา ภูริปญฺญาณา, โคตมา ภูริเมโส.}
\csromanhangingend

\prose{\textbf{กิํนุ ตมฺหา วิปฺปวสสี}ติ กิํนุ พุทฺธมฺหา วิปฺปวสสิ อเปสิ อปคจฺฉสิ\footnote{อปคจฺฉิ (สี)} วินา โหสฺ\textbf{อีติ} กิํนุ ตมฺหา วิปฺปวสสิ.}

\prose{\textbf{มุหุตฺตมปิ ปิงฺคิยา}ติ มุหุตฺตมฺปิ ขณมฺปิ ลยมฺปิ วยมฺปิ อทฺธมฺปฺ\textbf{อีติ} มุหุตฺตมปิ.\csromanspacer{} \textbf{ปิงฺคิยา}ติ พาวรี ตํ นตฺตารํ นาเมน อาลปติ.}

\csromanheader{2}{18. ปารายนานุคีติคาถานิทฺเทส}
\prose{\csromanfolio{215}\textbf{โคตมา ภูริปญฺญาณา}ติ \textbf{โคตมา ภูริปญฺญาณา} ญาณปญฺญาณา ปญฺญาธชา ปญฺญาเกตุมฺหา ปญฺญาธิปเตยฺยมฺหา วิจยพหุลา ปวิจยพหุลา โอกฺขายนพหุลา สโมกฺขายนธมฺมา วิภูตวิหาริมฺหา ตจฺจริตา ตพฺพหุลา ตคฺครุกา ตนฺนินฺนา ตปฺโปณา ตปฺปพฺภารา ตทธิมุตฺตา ตทธิปเตยฺยมฺหาติ \textbf{โคตมา ภูริปญฺญาณา}.}

\prose{\textbf{โคตมา ภูริเมธส}ติ ภูริ วุจฺจติ ปถวี, ภควา ตาย ปถวิสมาย ปญฺญาย วิปุลาย วิตฺถตาย สมนฺนาคโต.\csromanspacer{} เมธา วุจฺจติ ปญฺญา, ยา ปญฺญา ปชานนา ฯเปฯ อโมโห ธมฺมวิจโย สมฺมาทิฏฺฐิ.\csromanspacer{} ภควา อิมาย เมธาย ปญฺญาย อุเปโต สมุเปโต อุปาคโต สมุปาคโต อุปปนฺโน สมุปปนฺโน สมนฺนาคโต, ตสฺมา พุทฺโธ สุเมธโสติ \textbf{โคตมา ภูริเมธส}.\csromanspacer{} เตนาห โส พฺราหฺมโณ\csromandash{}}

\csromangathameasure{กิํนุ ตมฺหา วิปฺปวสิ, มุหุตฺตมปิ ปิงฺคิย. \\ โคตมา ภูริปญฺญาณา, โคตมา ภูริเมธสาติ.}
\csromangathagroup{\csromangathastanza{กิํนุ ตมฺหา วิปฺปวสิ, มุหุตฺตมปิ ปิงฺคิย. \\ โคตมา ภูริปญฺญาณา, โคตมา ภูริเมธสาติ.}}

\proseitem{110}{โย เต ธมฺมมเทเสสิ, สนฺทิฏฺฐิกมกาลิกํ.}

\prose{\textbf{ตณฺหกฺขยมนีติกํ, ยสฺส นตฺถิ อุปมา กฺวจิ. (9)}}

\prose{\textbf{โย เต ธมฺมมเทเสสี}ติ โย โส ภควา ฯเปฯ ตตฺถ จ สพฺพญฺญุตํ ปตฺโต พเลสุ จ วสีภาวํ.\csromanspacer{} \textbf{ธมฺมมเทเสสี}ติ \textbf{ธมฺมนฺ}ติ อาทิกลฺยาณํ มชฺเฌกลฺยาณํ ฯเปฯ นิพฺพานญฺจ นิพฺพานคามินิญฺจ ปฏิปทํ อาจิกฺขิ เทเสสิ ปญฺญเปสิ ปฏฺฐเปสิ วิวริ วิภชิ อุตฺตานี-อกาสิ ปกาเสสีติ โย เต ธมฺมมเทเสสิ.}

\prose{\textbf{สนฺทิฏฺฐิกมกาลิกนฺ}ติ สนฺทิฏฺฐิกํ อกาลิกํ เอหิปสฺสิกํ โอปเนยฺยิกํ ปจฺจตฺตํ เวทิตพฺพํ วิญฺญูหีติ เอวํ สนฺทิฏฺฐิกํ.\csromanspacer{} อถ วา โย ทิฏฺเฐว ธมฺเม อริยํ อฏฺฐงฺคิกํ มคฺคํ ภาเวติ, ตสฺส มคฺคสฺส อนนฺตรา สมนนฺตรา อธิคจฺฉเตว ผลํ วินฺทติ ปฏิลภตีติ เอวมฺปิ สนฺทิฏฺฐิกํ.\csromanspacer{} \textbf{อกาลิกนฺ}ติ ยถา มนุสฺสา กาลิกํ ธนํ ทตฺวา อนนฺตรา น ลภนฺติ, กาลํ อาคเมนฺติ, เนวายํ ธมฺโม.\csromanspacer{} โย ทิฏฺเฐว ธมฺเม อริยํ อฏฺฐงฺคิกํ มคฺคํ ภาเวติ.\csromanspacer{} ตสฺส มคฺคสฺส อนนฺตรา สมนนฺตรา อธิคจฺฉเตว ผลํ วินฺทติ ปฏิลภติ, น ปรตฺถ น ปรโลเก, เอวํ อกาลิกนฺติ สนฺทิฏฺฐิกมกาลิกํ.}

\prose{\csromanfolio{216}\textbf{ตณฺหกฺขยมนีติกนฺ}ติ \textbf{ตณฺหา}ติ รูป\textbf{ตณฺหา} ฯเปฯ ธมฺม\textbf{ตณฺหา}.\csromanspacer{} \textbf{ตณฺหกฺขยนฺ}ติ ตณฺหกฺขยํ ราคกฺขยํ โทสกฺขยํ โมหกฺขยํ คติกฺขยํ อุปปตฺติกฺขยํ ปฏิสนฺธิกฺขยํ ภวกฺขยํ สํสารกฺขยํ วฏฺฏกฺขยํ.\csromanspacer{} \textbf{อนีติกนฺ}ติ \textbf{อีติ} วุจฺจนฺติ กิเลสา จ ขนฺธา จ อภิสงฺขารา จ. ¢ติปฺปหานํ \textbf{อีติ}วูปสมํ \textbf{อีติ}ปฏินิสฺสคฺคํ \textbf{อีติ}ปฏิปฺปสฺสทฺธิํ อมตํ นิพฺพานนฺติ ตณฺหกฺขยมนฺ\textbf{อีติ}กํ.}

\prose{\textbf{ยสฺส นตฺถิ อุปมา กฺวจี}ติ \textbf{ยสฺสา}ติ นิพฺพานสฺส.\csromanspacer{} \textbf{นตฺถิ อุปมา}ติ อุปมา นตฺถิ, อุปนิธา นตฺถิ, สทิสํ นตฺถิ, ปฏิภาโค นตฺถิ น สติ น สํวิชฺชติ นุปลพฺภติ.\csromanspacer{} \textbf{กฺวจี}ติ กฺวจิ กิมฺหิจิ กตฺถจิ อชฺฌตฺตํ วา พหิทฺธา วา อชฺฌตฺตพหิทฺธา วาติ ยสฺส นตฺถิ อุปมา กฺวจิ.\csromanspacer{} เตนาห โส พฺราหฺมโณ\csromandash{}}

\csromangathameasure{โย เต ธมฺมมเทเสสิ, สนฺทิฏฺฐิกมกาลิกํ. \\ ตณฺหกฺขยมนฺ\textbf{อีติ}กํ, ยสฺส นตฺถิ อุปมา กฺวจฺ\textbf{อีติ}.}
\csromangathagroup{\csromangathastanza{โย เต ธมฺมมเทเสสิ, สนฺทิฏฺฐิกมกาลิกํ. \\ ตณฺหกฺขยมนฺ\textbf{อีติ}กํ, ยสฺส นตฺถิ อุปมา กฺวจฺ\textbf{อีติ}.}}

\proseitem{111}{นาหํ ตมฺหา วิปฺปวสามิ, มุหุตฺตมปิ พฺราหฺมณ.}

\prose{\textbf{โคตมา ภูริปญฺญาณา, โคตมา ภูริเมธสา. (10)}}

\prose{\textbf{นาหํ ตมฺหา วิปฺปวสามี}ติ นาหํ พุทฺธมฺหา วิปฺปวสามิ อเปมิ อปคจฺฉามิ วินา โหมฺ\textbf{อีติ} นาหํ ตมฺหา วิปฺปวสามิ.}

\prose{\textbf{มุหุตฺตมปิ พฺราหฺมณา}ติ มุหุตฺตมฺปิ ขณมฺปิ ลยมฺปิ วยมฺปิ อทฺธมฺปฺ\textbf{อีติ} มุหุตฺตมปิ.\csromanspacer{} \textbf{พฺราหฺมณา}ติ คารเวน มาตุลํ อาลปติ.}

\prose{\textbf{โคตมา ภูริปญฺญาณา}ติ \textbf{โคตมา ภูริปญฺญาณา} ญาณปญฺญาณา ปญฺญาธชา ปญฺญาเกตุมฺหา ปญฺญาธิปเตยฺยมฺหา วิจยพหุลา ปวิจยพหุลา โอกฺขายนพหุลา สโมกฺขายนธมฺมา วิภูตวิหาริมฺหา ตจฺจริตา ตพฺพหุลา ตคฺครุกา ตนฺนินฺนา ตปฺโปณา ตปฺปพฺภารา ตทธิมุตฺตา ตทธิปเตยฺยมฺหาติ \textbf{โคตมา ภูริปญฺญาณา}.}

\prose{\textbf{โคตมา ภูริเมธสา}ติ ภูริ วุจฺจติ ปถวี, ภควา ตาย ปถวิสมาย ปญฺญาย วิปุลาย วิตฺถตาย สมนฺนาคโต.\csromanspacer{} เมธา วุจฺจติ ปญฺญา, ยา ปญฺญา ปชานนา ฯเปฯ อโมโห ธมฺมวิจโย สมฺมาทิฏฺฐิ.\csromanspacer{} ภควา อิมาย เมธาย ปญฺญาย อุเปโต สมุเปโต อุปาคโต สมุปาคโต อุปปนฺโน สมุปปนฺโน สมนฺนาคโต, ตสฺมา พุทฺโธ สุเมธโสติ \textbf{โคตมา ภูริเมธสา}.\csromanspacer{} เตนาห เถโร ปิงฺคิโย\csromandash{}}

\csromanheader{2}{18. ปารายนานุคฺ\textbf{อีติ}คาถานิทฺเทส}
\csromangathameasure{นาหํ ตมฺหา วิปฺปวสามิ, มุหุตฺตมปิ พฺราหฺมณ. \\ โคตมา ภูริปญฺญาณา, โคตมา ภูริเมธสาติ.}
\csromangathagroup{\csromangathastanza{\csromanfolio{217}นาหํ ตมฺหา วิปฺปวสามิ, มุหุตฺตมปิ พฺราหฺมณ. \\ โคตมา ภูริปญฺญาณา, โคตมา ภูริเมธสาติ.}}

\proseitem{112}{โย เม ธมฺมมเทเสสิ, สนฺทิฏฺฐิกมกาลิกํ.}

\prose{\textbf{ตณฺหกฺขยมนีติกํ, ยสฺส นตฺถิ อุปมา กฺวจิ. (11)}}

\prose{\textbf{โย เม ธมฺมมเทเสสี}ติ โย โส ภควา สยมฺภู อนาจริยโก ปุพฺเพ อนนุสฺสุเตสุ ธมฺเมสุ สามํ สจฺจานิ อภิสมฺพุชฺฌิ, ตตฺถ จ สพฺพญฺญุตํ ปตฺโต พเลสุ จ วสีภาวํ.\csromanspacer{} \textbf{ธมฺมมเทเสสี}ติ \textbf{ธมฺมนฺ}ติ อาทิกลฺยาณํ มชฺเฌกลฺยาณํ ปริโยสานกลฺยาณํ สาตฺถํ สพฺยญฺชนํ เกวลปริปุณฺณํ ปริสุทฺธํ พฺรหฺมจริยํ จตฺตาโร สติปฏฺฐาเน จตฺตาโร สมฺมปฺปธาเน จตฺตาโร อิทฺธิปาเท ปญฺจินฺทฺริยานิ ปญฺจ พลานิ สตฺต โพชฺฌงฺเค อริยํ อฏฺฐงฺคิกํ มคฺคํ นิพฺพานญฺจ นิพฺพานคามินิญฺจ ปฏิปทํ อาจิกฺขิ เทเสสิ ปญฺญเปสิ ปฏฺฐเปสิ วิวริ วิภชิ อุตฺตานี-อกาสิ ปกาเสสฺ\textbf{อีติ} โย เม ธมฺมมเทเสสิ.}

\prose{\textbf{สนฺทิฏฺฐิกมกาลิกนฺ}ติ สนฺทิฏฺฐิกํ อกาลิกํ เอหิปสฺสิกํ โอปเนยฺยิกํ ปจฺจตฺตํ เวทิตพฺพํ วิญฺญูหฺ\textbf{อีติ} เอวํ สนฺทิฏฺฐิกํ.\csromanspacer{} อถ วา โย ทิฏฺเฐว ธมฺเม อริยํ อฏฺฐงฺคิกํ มคฺคํ ภาเวติ, ตสฺส มคฺคสฺส อนนฺตรา สมนนฺตรา อธิคจฺฉเตว ผลํ วินฺทติ ปฏิลภตฺ\textbf{อีติ} เอวมฺปิ สนฺทิฏฺฐิกํ.\csromanspacer{} \textbf{อกาลิกนฺ}ติ ยถา มนุสฺสา กาลิกํ ธนํ ทตฺวา อนนฺตรา น ลภนฺติ, กาลํ อาคเมนฺติ, เนวายํ ธมฺโม.\csromanspacer{} โย ทิฏฺเฐว ธมฺเม อริยํ อฏฺฐงฺคิกํ มคฺคํ ภาเวติ, ตสฺส มคฺคสฺส อนนฺตรา สมนนฺตรา อธิคจฺฉเตว ผลํ วินฺทติ ปฏิลภติ, น ปรตฺถ น ปรโลเก, เอวํ อกาลิกนฺติ สนฺทิฏฺฐิกมกาลิกํ.}

\prose{\textbf{ตณฺหกฺขยมนีติกนฺ}ติ \textbf{ตณฺหา}ติ รูป\textbf{ตณฺหา} ฯเปฯ ธมฺม\textbf{ตณฺหา}.\csromanspacer{} \textbf{ตณฺหกฺขยนฺ}ติ ตณฺหกฺขยํ ราคกฺขยํ โทสกฺขยํ โมหกฺขยํ คติกฺขยํ อุปปตฺติกฺขยํ ปฏิสนฺธิกฺขยํ ภวกฺขยํ สํสารกฺขยํ วฏฺฏกฺขยํ.\csromanspacer{} \textbf{อนีติกนฺ}ติ \textbf{อีติ} วุจฺจนฺติ กิเลสา จ ขนฺธา จ อภิสงฺขารา จ. ¢ติปฺปหานํ \textbf{อีติ}วูปสมํ \textbf{อีติ}ปฏิปฺปสฺสทฺธิํ อมตํ นิพฺพานนฺติ ตณฺหกฺขยมนฺ\textbf{อีติ}กํ.}

\prose{\textbf{ยสฺส นตฺถิ อุปมา กฺวจี}ติ \textbf{ยสฺสา}ติ นิพฺพานสฺส.\csromanspacer{} \textbf{นตฺถิ อุปมา}ติ อุปมา นตฺถิ, อุปนิธา นตฺถิ, สทิสํ นตฺถิ, ปฏิภาโค นตฺถิ น สติ น สํวิชฺชติ นุปลพฺภติ.\csromanspacer{} \textbf{กฺวจี}ติ กฺวจิ กิมฺหิจิ กตฺถจิ อชฺฌตฺตํ วา พหิทฺธา วา อชฺฌตฺตพหิทฺธา วาติ ยสฺส นตฺถิ อุปมา กฺวจิ.\csromanspacer{} เตนาห เถโร ปิงฺคิโย\csromandash{}}

\csromangathameasure{โย เม ธมฺมมเทเสสิ, สนฺทิฏฺฐิกมกาลิกํ. \\ ตณฺหกฺขยมนีติกํ, ยสฺส นตฺถิ อุปมา กฺวจีติ.}
\csromangathagroup{\csromangathastanza{\csromanfolio{218}โย เม ธมฺมมเทเสสิ, สนฺทิฏฺฐิกมกาลิกํ. \\ ตณฺหกฺขยมนีติกํ, ยสฺส นตฺถิ อุปมา กฺวจีติ.}}

\proseitem{113}{ปสฺสามิ นํ มนสา จกฺขุนาว,}

\prose{\textbf{รตฺตินฺทิวํ พฺราหฺมณ อปฺปมตฺโต.}}

\prose{\textbf{นมสฺสมาโน วิวเสมิ}\footnote{นมสฺสมาโนว วเสมิ (สี-ฏฺฐ) ...วิวสามิ (สฺยา)} \textbf{รตฺติํ,}}

\csromanheader{2}{\textbf{เตเนว มญฺญามิ อวิปฺปวาสํ. (12)}}
\prose{\textbf{ปสฺสามิ นํ มนสา จกฺขุนาวา}ติ ยถา จกฺขุมา ปุริโส อาโลเก รูปคตานิ ปสฺเสยฺย ทกฺเขยฺย โอโลเกยฺย นิชฺฌาเยยฺย อุปปริกฺเขยฺย, เอวเมวาหํ พุทฺธํ ภควนฺตํ มนสา ปสฺสามิ ทกฺขามิ โอโลเกมิ นิชฺฌายามิ อุปปริกฺขามีติ ปสฺสามิ นํ มนสา จกฺขุนาว.}

\prose{\textbf{รตฺตินฺทิวํ พฺราหฺมณ อปฺปมตฺโต}ติ รตฺติญฺจ ทิวา จ พุทฺธานุสฺสติํ มนสา ภาเวนฺโต อปฺปมตฺโตติ รตฺตินฺทิวํ พฺราหฺมณ อปฺปมตฺโต.}

\prose{\textbf{นมสฺสมาโน วิวเสมิ รตฺตินฺ}ติ \textbf{นมสฺสมาโน}ติ กาเยน วา \textbf{นมสฺสมาโน}, วาจาย วา \textbf{นมสฺสมาโน}, จิตฺเตน วา \textbf{นมสฺสมาโน}, อนฺวตฺถปฏิปตฺติยา วา \textbf{นมสฺสมาโน}, ธมฺมานุธมฺมปฏิปตฺติยา วา \textbf{นมสฺสมาโน} สกฺการมาโน ครุการมาโน มานยมาโน ปูชยมาโน รตฺตินฺทิวํ วิวเสมิ อตินาเมมิ อติกฺกเมมีติ \textbf{นมสฺสมาโน} วิวเสมิ รตฺติํ.}

\prose{\textbf{เตเนว มญฺญามิ อวิปฺปวาสนฺ}ติ ตาย พุทฺธานุสฺสติยา ภาเวนฺโต อวิปฺปวาโสติ ตํ มญฺญามิ, อวิปฺปวุฏฺโฐติ ตํ มญฺญามิ ชานามิ, เอวํ ชานามิ เอวํ อาชานามิ เอวํ วิชานามิ เอวํ ปฏิวิชานามิ เอวํ ปฏิวิชฺฌามีติ เตเนว มญฺญามิ อวิปฺปวาสํ.\csromanspacer{} เตนาห เถโร ปิงฺคิโย\csromandash{}}

\prose{ปสฺสามิ นํ มนสา จกฺขุนาว,}

\prose{\textbf{รตฺตินฺทิวํ พฺราหฺมณ อปฺปมตฺโต.}}

\prose{\textbf{นมสฺสมาโน วิวเสมิ} \textbf{รตฺติํ,}}

\prose{\textbf{เตเนว มญฺญามิ อวิปฺปวาสนฺ}ติ.}

\csromanheader{2}{18. ปารายนานุคีติคาถานิทฺเทส}
\proseitem{114}{\csromanfolio{219}สทฺธา จ ปีติ จ มโน สติ จ,}

\prose{\textbf{นา’เปนฺติ’เม โคตมสาสนมฺหา.}}

\prose{\textbf{ยํ ยํ ทิสํ วชติ ภูริปญฺโญ,}}

\csromanheader{2}{\textbf{ส เตน เตเนว นโตหมสฺมิ. (13)}}
\prose{\textbf{สทฺธา จ ปีติ จ มโน สติ จา}ติ \textbf{สทฺธา}ติ ยา จ ภควนฺตํ อารพฺภ \textbf{สทฺธา} สทฺทหนา\footnote{สทฺธหนา (ก)} โอกปฺปนา อภิปฺปสาโท \textbf{สทฺธา} สทฺธินฺทฺริยํ \textbf{สทฺธา}พลํ.\csromanspacer{} \textbf{ปีตี}ติ ยา ภควนฺตํ อารพฺภ ปีติ ปาโมชฺชํ\footnote{ปามุชฺชํ (สฺยา)} โมทนา อาโมทนา ปโมทนา หาโส ปหาโส วิตฺติ ตุฏฺฐิ โอทคฺยํ อตฺตมนตา จิตฺตสฺส.\csromanspacer{} \textbf{มโน}ติ ยญฺจ ภควนฺตํ อารพฺภ จิตฺตํ มโน มานสํ หทยํ ปณฺฑรํ มโน มนายตนํ มนินฺทฺริยํ วิญฺญาณํ วิญฺญาณกฺขนฺโธ ตชฺชา มโนวิญฺญาณธาตุ.\csromanspacer{} \textbf{สตี}ติ ยา ภควนฺตํ อารพฺภ สติ อนุสฺสติ สมฺมาสตีติ \textbf{สทฺธา} จ ปีติ จ มโน สติ จ.}

\prose{\textbf{นา’เปนฺติ’เม โคตมสาสนมฺหา}ติ อิเม จตฺตาโร ธมฺมา โคตมสาสนา พุทฺธสาสนา ชินสาสนา ตถาคตสาสนา อรหนฺตสาสนา นาเปนฺติ น คจฺฉนฺติ น วิชหนฺติ น วินาเสนฺตีติ นาเปนฺติเม โคตมสาสนมฺหา.}

\prose{\textbf{ยํ ยํ ทิสํ วชติ ภูริปญฺโญ}ติ \textbf{ยํ ยํ ทิสนฺ}ติ ปุรตฺถิมํ วา ทิสํ ปจฺฉิมํ วา ทิสํ ทกฺขิณํ วา ทิสํ อุตฺตรํ วา ทิสํ วชติ คจฺฉติ กมติ อภิกฺกมติ.\csromanspacer{} \textbf{ภูริปญฺโญ}ติ \textbf{ภูริปญฺโญ} มหาปญฺโญ ติกฺขปญฺโญ ปุถุปญฺโญ หาสปญฺโญ ชวนปญฺโญ นิพฺเพธิกปญฺโญ.\csromanspacer{} ภูริ วุจฺจติ ปถวี, ภควา ตาย ปถวิสมาย ปญฺญาย วิปุลาย วิตฺถตาย สมนฺนาคโตติ ยํ ยํ ทิสํ วชติ \textbf{ภูริปญฺโญ}.}

\prose{\textbf{ส เตน เตเนว นโตหมสฺมี}ติ โส เยน พุทฺโธ เตน เตเนว นโต ตนฺนินฺโน ตปฺโปโณ ตปฺปพฺภาโร ตทธิมุตฺโต ตทธิปเตยฺโยติ ส เตน เตเนว นโตหมสฺมิ.\csromanspacer{} เตนาห เถโร ปิงฺคิโย\csromandash{}}

\csromangathameasure{สทฺธา จ ปีติ จ \textbf{มโน} สติ จ, \\ นา’เปนฺติ’เม โคตมสาสนมฺหา. \\ ยํ ยํ ทิสํ วชติ ภูริปญฺโญ, \\ ส เตน เตเนว นโตหมสฺมีติ.}
\csromangathagroup{\csromangathastanza{\csromanfolio{220}สทฺธา จ ปีติ จ \textbf{มโน} สติ จ, \\ นา’เปนฺติ’เม โคตมสาสนมฺหา. \\ ยํ ยํ ทิสํ วชติ ภูริปญฺโญ, \\ ส เตน เตเนว นโตหมสฺมีติ.}}

\proseitem{115}{ชิณฺณสฺส เม ทุพฺพลถามกสฺส,}

\prose{\textbf{เตเนว กาโย น ปเลติ ตตฺถ.}}

\prose{\textbf{สงฺกปฺปยนฺตาย วชามิ นิจฺจํ,}}

\csromanheader{2}{\textbf{มโน หิ เม พฺราหฺมณ เตน ยุตฺโต. (14)}}
\prose{\textbf{ชิณฺณสฺส เม ทุพฺพลถามกสฺสา}ติ ชิณฺณสฺส วุฑฺฒสฺส มหลฺลกสฺส อทฺธคตสฺส วโย-อนุปฺปตฺตสฺส.\csromanspacer{} \textbf{ทุพฺพลถามกสฺสา}ติ ทุพฺพลถามกสฺส อปฺปถามกสฺส ปริตฺตถามกสฺสาติ ชิณฺณสฺส เม ทุพฺพลถามกสฺส.}

\prose{\textbf{เตเนว กาโย น ปเลติ ตตฺถา}ติ กาโย เยน พุทฺโธ, เตน น ปเลติ น วชติ น คจฺฉติ นาติกฺกมตีติ เตเนว กาโย น ปเลติ ตตฺถ.}

\prose{\textbf{สงฺกปฺปยนฺตาย วชามิ นิจฺจนฺ}ติ สงฺกปฺปคมเนน วิตกฺกคมเนน ญาณคมเนน ปญฺญาคมเนน พุทฺธิคมเนน วชามิ คจฺฉามิ อติกฺกมามีติ สงฺกปฺปยนฺตาย วชามิ นิจฺจํ.}

\prose{\textbf{มโน หิ เม พฺราหฺมณ เตน ยุตฺโต}ติ \textbf{มโน}ติ ยํ จิตฺตํ \textbf{มโน} มานสํ ฯเปฯ ตชฺชา \textbf{มโน}วิญฺญาณธาตุ.\csromanspacer{} \textbf{มโน หิ เม พฺราหฺมณ เตน ยุตฺโต}ติ \textbf{มโน} เยน พุทฺโธ เตน ยุตฺโต ปยุตฺโต สํยุตฺโตติ \textbf{มโน} หิ เม พฺราหฺมณ เตน ยุตฺโต.\csromanspacer{} เตนาห เถโร ปิงฺคิโย\csromandash{}}

\prose{ชิณฺณสฺส เม ทุพฺพลถามกสฺส,}

\prose{\textbf{เตเนว กาโย น ปเลติ ตตฺถ.}}

\csromangathameasure{\textbf{สงฺกปฺปยนฺตาย วชามิ นิจฺจํ,} \\ \textbf{มโน หิ เม พฺราหฺมณ เตน ยุตฺโต}ติ.}
\csromangathagroup{\csromangathastanza{\textbf{สงฺกปฺปยนฺตาย วชามิ นิจฺจํ,} \\ \textbf{มโน หิ เม พฺราหฺมณ เตน ยุตฺโต}ติ.}}

\csromanhangingbegin
\hangingheaditem{116}{ปงฺเก สยาโน ปริผนฺทมาโน, ทีปา ทีปํ อุปลฺลวิํ.}
\hangingbody{อถทฺทสาสิํ สมฺพุทฺธํ, โอฆติณฺณมนาสวํ.}
\csromanhangingend

\csromanheader{2}{(15)}
\csromanheader{2}{18. ปารายนานุคีติคาถานิทฺเทส}
\prose{\csromanfolio{221}\textbf{ปงฺเก สยาโน ปริผนฺทมาโน}ติ \textbf{ปงฺเก สยาโน}ติ กามปงฺเก กามกทฺทเม กามกิเลเส กามพฬิเส กามปริฬาเห กามปลิโพเธ เสมาโน สยมาโน วสมาโน อาวสมาโน ปริวสมาโน1ติ \textbf{ปงฺเก สยาโน}.\csromanspacer{} \textbf{ปริผนฺทมาโน}ติ ตณฺหาผนฺทนาย ผนฺทมาโน, ทิฏฺฐิผนฺทนาย ผนฺทมาโน, กิเลสผนฺทนาย ผนฺทมาโน, ปโยคผนฺทนาย ผนฺทมาโน, วิปากผนฺทนาย ผนฺทมาโน, มโนทุจฺจริตผนฺทนาย ผนฺทมาโน, รตฺโต ราเคน ผนฺทมาโน, ทุฏฺโฐ โทเสน ผนฺทมาโน, มูโฬฺห โมเหน ผนฺทมาโน, วินิพนฺโธ มาเนน ผนฺทมาโน, ปรามฏฺโฐ ทิฏฺฐิยา ผนฺทมาโน, วิกฺเขปคโต อุทฺธจฺเจน ผนฺทมาโน, อนิฏฺฐงฺคโต วิจิกิจฺฉาย ผนฺทมาโน, ถามคโต อนุสเยหิ ผนฺทมาโน, ลาเภน ผนฺทมาโน, อลาเภน ผนฺทมาโน, ยเสน ผนฺทมาโน, อยเสน ผนฺทมาโน, ปสํสาย ผนฺทมาโน, นินฺทาย ผนฺทมาโน, สุเขน ผนฺทมาโน, ทุกฺเขน ผนฺทมาโน, ชาติยา ผนฺทมาโน, ชราย ผนฺทมาโน, พฺยาธินา ผนฺทมาโน, มรเณน ผนฺทมาโน, โสกปริเทวทุกฺขโทมนสฺสุปายาเสหิ ผนฺทมาโน, เนรยิเกน ทุกฺเขน ผนฺทมาโน, ติรจฺฉานโยนิเกน ทุกฺเขน ผนฺทมาโน, เปตฺติวิสยิเกน ทุกฺเขน ผนฺทมาโน, มานุสิเกน ทุกฺเขน.\csromanspacer{} คพฺโภกฺกนฺติมูลเกน ทุกฺเขน.\csromanspacer{} คพฺภฏฺฐิติมูลเกน ทุกฺเขน.\csromanspacer{} คพฺภวุฏฺฐานมูลเกน ทุกฺเขน.\csromanspacer{} ชาตสฺสูปนิพนฺธเกน ทุกฺเขน.\csromanspacer{} ชาตสฺส ปราเธยฺยเกน ทุกฺเขน.\csromanspacer{} อตฺตูปกฺกเมน ทุกฺเขน.\csromanspacer{} ปรูปกฺกเมน ทุกฺเขน.\csromanspacer{} สงฺขารทุกฺเขน.\csromanspacer{} วิปริณามทุกฺเขน.\csromanspacer{} จกฺขุโรเคน ทุกฺเขน.\csromanspacer{} โสตโรเคน ทุกฺเขน.\csromanspacer{} ฆานโรเคน ทุกฺเขน.\csromanspacer{} ชิวฺหาโรเคน ทุกฺเขน.\csromanspacer{} กายโรเคน ทุกฺเขน.\csromanspacer{} สีสโรเคน ทุกฺเขน.\csromanspacer{} กณฺณโรเคน ทุกฺเขน.\csromanspacer{} มุขโรเคน ทุกฺเขน.\csromanspacer{} ทนฺตโรเคน ทุกฺเขน.\csromanspacer{} โอฏฺฐโรเคน ทุกฺเขน.\csromanspacer{} กาเสน.\csromanspacer{} สาเสน.\csromanspacer{} ปินาเสน.\csromanspacer{} qอาเหน\footnote{ทาเหน (ก) ขุ 7. 35 ปิฏฺเฐปิ.}.\csromanspacer{} ชเรน.\csromanspacer{} กุจฺฉิโรเคน.\csromanspacer{} มุจฺฉาย.\csromanspacer{} ปกฺขนฺทิกาย.\csromanspacer{} สูลาย.\csromanspacer{} วิสูจิกาย.\csromanspacer{} กุฏฺเฐน.\csromanspacer{} คณฺเฑน.\csromanspacer{} กิลาเสน.\csromanspacer{} โสเสน.\csromanspacer{} อปมาเรน.\csromanspacer{} ททฺทุยา.\csromanspacer{} กณฺฑุยา.\csromanspacer{} กจฺฉุยา.\csromanspacer{} รขสาย.\csromanspacer{} วิตจฺฉิกาย.\csromanspacer{} โลหิตปิตฺเตน\footnote{โลหิเตน. ปิตฺเตน (สฺยา, ก)}.\csromanspacer{} มธุเมเหน.\csromanspacer{} อํสาย.\csromanspacer{} ปิฬกาย.\csromanspacer{} ภคนฺทลาย\footnote{ภคนฺทลาย (สฺยา)}.\csromanspacer{} ปิตฺตสมุฏฺฐาเนน อาพาเธน.\csromanspacer{} เสมฺหสมุฏฺฐาเนน อาพาเธน.\csromanspacer{} วาตสมุฏฺฐาเนน อาพาเธน.\csromanspacer{} สนฺนิปาติเกน อาพาเธน.\csromanspacer{} อุตุปริณามเชน อาพาเธน.\csromanspacer{} วิสมปริหารเชน อาพาเธน.\csromanspacer{} โอปกฺกมิเกน อาพาเธน.\csromanspacer{} กมฺมวิปากเชน อาพาเธน.\csromanspacer{} สีเตน.\csromanspacer{} อุเณฺหน.}

\prose{\csromanfolio{222}ชิฆจฺฉาย.\csromanspacer{} ปิปาสาย.\csromanspacer{} อุจฺจาเรน.\csromanspacer{} ปสฺสาเวน.\csromanspacer{} qอํสมกสวาตาตปสรีสปสมฺผสฺเสน ทุกฺเขน.\csromanspacer{} มาตุมรเณน ทุกฺเขน.\csromanspacer{} ปิตุมรเณน ทุกฺเขน.\csromanspacer{} ปุตฺตมรเณน ทุกฺเขน.\csromanspacer{} ธีตุมรเณน ทุกฺเขน.\csromanspacer{} ญาติพฺยสเนน ทุกฺเขน.\csromanspacer{} โภคพฺยสเนน ทุกฺเขน.\csromanspacer{} โรคพฺยสเนน ทุกฺเขน.\csromanspacer{} สีลพฺยสเนน ทุกฺเขน.\csromanspacer{} ทิฏฺฐิพฺยเนน ทุกฺเขน ผนฺทมาโน ปริผนฺทมาโน ปเวธมาโน สมฺปเวธมาโนติ ปงฺเก สยาโน ปริผนฺทมาโน.}

\prose{\textbf{ทีปา ทีปํ อุปลฺลวินฺ}ติ สตฺถารโต สตฺถารํ ธมฺมกฺขานโต ธมฺมกฺขานํ คณโต คณํ ทิฏฺฐิยา ทิฏฺฐิํ ปฏิปทาย ปฏิปทํ มคฺคโต มคฺคํ ปลฺลวิํ อุปลฺลวิํ สมฺปลฺลวินฺติ ทีปา ทีปํ อุปลฺลวิํ.}

\prose{\textbf{อถทฺทสาสิํ สมฺพุทฺธนฺ}ติ \textbf{อถา}ติ ปทสนฺธิ ปทสํสคฺโค ปทปาริปูรี อกฺขรสมวาโย พฺยญฺชนสิลิฏฺฐตา ปทานุปุพฺพตาเปตํ \textbf{อถา}ติ.\csromanspacer{} \textbf{อทฺทสาสินฺ}ติ อทฺทสํ อทฺทกฺขิํ อปสฺสิํ ปฏิวิชฺฌิํ.\csromanspacer{} \textbf{พุทฺโธ}ติ โย โส ภควา สยมฺภู อนาจริยโก ฯเปฯ สจฺฉิกา ปญฺญตฺติ, ยทิทํ \textbf{พุทฺโธ}ติ อถทฺทสาสิํ สมฺพุทฺธํ.}

\prose{\textbf{โอฆติณฺณมนาสวนฺ}ติ \textbf{โอฆติณฺณนฺ}ติ ภควา กาโมฆํ ติณฺโณ, ภโวฆํ ติณฺโณ, ทิฏฺโฐฆํ ติณฺโณ, อวิชฺโชฆํ ติณฺโณ, สพฺพสํสารปถํ ติณฺโณ อุตฺติณฺโณ นิตฺถิณฺโณ อติกฺกนฺโต สมติกฺกนฺโต วีติวตฺโต, โส วุตฺถวาโส จิณฺณจรโณ ฯเปฯ ชาติมรณสํสาโร, นตฺถิ ตสฺส ปุนพฺภโวติ โอฆติณฺณํ.\csromanspacer{} \textbf{อนาสวนฺ}ติ จตฺตาโร อาสวา กามาสโว ภวาสโว ทิฏฺฐาสโว อวิชฺชาสโว, เต อาสวา พุทฺธสฺส ภควโต ปหีนา อุจฺฉินฺนมูลา ตาลาวตฺถุกตา อนภาวํกตา อายติํ อนุปฺปาทธมฺมา, ตสฺมา \textbf{พุทฺโธ} อนาสโวติ โอฆติณฺณมนาสวํ.\csromanspacer{} เตนาห เถโร ปิงฺคิโย\csromandash{}}

\csromangathameasure{ปงฺเก สยาโน ปริผนฺทมาโน, ทีปา ทีปํ อุปลฺลวิํ. \\ อถทฺทสาสิํ สมฺพุทฺธํ, โอฆติณฺณมนาสวนฺติ.}
\csromangathagroup{\csromangathastanza{ปงฺเก สยาโน ปริผนฺทมาโน, ทีปา ทีปํ อุปลฺลวิํ. \\ อถทฺทสาสิํ สมฺพุทฺธํ, โอฆติณฺณมนาสวนฺติ.}}

\proseitem{117}{ยฺ\textbf{อถา} อหู วกฺกลิ มุตฺตสทฺโธ,}

\prose{\textbf{ภทฺราวุโธ อาฬวิโคตโม จ.}}

\prose{\textbf{เอวเมว ตฺวมฺปิ ปมุญฺจสฺสุ สทฺธํ,}}

\csromanheader{2}{\textbf{คมิสฺสสิ ตฺวํ ปิงฺคิย มจฺจุเธยฺยสฺส ปารํ. (16)}}
\csromanheader{2}{18. ปารายนานุคีติคาถานิทฺเทส}
\prose{\csromanfolio{223}\textbf{ยถา อหู วกฺกลิ มุตฺตสทฺโธ, ภทฺราวุโธ อาฬวิโคตโม จา}ติ ยถา วกฺกลิตฺเถโร\footnote{วกฺกลิ (สฺยา)} สทฺโธ สทฺธาครุโก สทฺธาปุพฺพงฺคโม สทฺธาธิมุตฺโต สทฺธาธิปเตยฺโย อรหตฺตปฺปตฺโต, ยถา ภทฺราวุโธ เถโร สทฺโธ สทฺธาครุโก สทฺธาปุพฺพงฺคโม สทฺธาธิมุตฺโต สทฺธาธิปเตยฺโย อรหตฺตปฺปตฺโต, ยถา อาฬวิโคตโม เถโร สทฺโธ สทฺธาครุโก สทฺธาปุพฺพงฺคโม สทฺธาธิมุตฺโต สทฺธาธิปเตยฺโย อรหตฺตปฺปตฺโตติ ยถา อหู วกฺกลิ มุตฺตสทฺโธ, ภทฺราวุโธ อาฬวิโคตโม จ.}

\prose{\textbf{เอวเมว ตฺวมฺปิ ปมุญฺจสฺสุ สทฺธนฺ}ติ เอวเมว ตฺวํ สทฺธํ มุญฺจสฺสุ ปมุญฺจสฺสุ สมฺปมุญฺจสฺสุ อธิมุญฺจสฺสุ โอกปฺเปหิ. “สพฺเพ สงฺขารา อนิจฺจา”ติ สทฺธํ มุญฺจสฺสุ ปมุญฺจสฺสุ สมฺปมุญฺจสฺสุ อธิมุญฺจสฺสุ โอกปฺเปหิ. “สพฺเพ สงฺขารา ทุกฺขา”ติ. “สพฺเพ ธมฺมา อนตฺตา”ติ สทฺธํ มุญฺจสฺสุ ปมุญฺจสฺสุ สมฺปมุญฺจสฺสุ อธิมุญฺจสฺสุ โอกปฺเปหิ ฯเปฯ. “ยํ กิญฺจิ สมุทยธมฺมํ, สพฺพํ ตํ นิโรธธมฺมนฺ”ติ สทฺธํ มุญฺจสฺสุ ปมุญฺจสฺสุ สมฺปมุญฺจสฺสุ อธิมุญฺจสฺสุ โอกปฺเปหีติ เอวเมว ตฺวมฺปิ ปมุญฺจสฺสุ สทฺธํ.}

\prose{\textbf{คมิสฺสสิ ตฺวํ ปิงฺคิย มจฺจุเธยฺยสฺส ปารนฺ}ติ มจฺจุเธยฺยํ วุจฺจนฺติ กิเลสา จ ขนฺธา จ อภิสงฺขารา จ.\csromanspacer{} มจฺจุเธยฺยสฺส ปารํ วุจฺจติ อมตํ นิพฺพานํ, โย โส สพฺพสงฺขารสมโถ สพฺพูปธิปฏินิสฺสคฺโค ตณฺหกฺขโย วิราโค นิโรโธ นิพฺพานํ.\csromanspacer{} \textbf{คมิสฺสสิ ตฺวํ ปิงฺคิย มจฺจุเธยฺยสฺส ปารนฺ}ติ ตฺวํ ปารํ คมิสฺสสิ, ปารํ อธิคมิสฺสสิ, ปารํ ผสฺสิสฺสสิ, ปารํ สจฺฉิกริสฺสสีติ คมิสฺสสิ ตฺวํ ปิงฺคิย มจฺจุเธยฺยสฺส ปารํ.\csromanspacer{} เตนาห ภควา\csromandash{}}

\prose{ยถา อหู วกฺกลิ มุตฺตสทฺโธ,}

\prose{ภทฺราวุโธ อาฬวิโคตโม จ.}

\prose{เอวเมว ตฺวมฺปิ ปมุญฺจสฺสุ สทฺธํ,}

\prose{\textbf{คมิสฺสสิ ตฺวํ ปิงฺคิย มจฺจุเธยฺยสฺส ปารนฺ}ติ.}

\proseitem{118}{เอส ภิยฺโย ปสีทามิ, สุตฺวาน มุนิโน วโจ.}

\prose{\textbf{วิวฏฺฏจฺฉโท}\footnote{วิวฏจฺฉทโน (ก) สทฺทนีติปทมาลา โอโลเกตพฺพา.} \textbf{สมฺพุทฺโธ, อขิโล ปฏิภานวา}\footnote{ปฏิภาณวา (สฺยา)}\textbf{. (17)}}

\prose{\textbf{เอส ภิยฺโย ปสีทามี}ติ เอส ภิยฺโย ปสีทามิ, ภิยฺโย ภิยฺโย สทฺทหามิ, ภิยฺโย ภิยฺโย โอกปฺเปมิ, ภิยฺโย ภิยฺโย อธิมุจฺจามิ, “สพฺเพ สงฺขารา อนิจฺจา”ติ ภิยฺโย ภิยฺโย ปสีทามิ, ภิยฺโย ภิยฺโย \csromanfolio{224}สทฺทหามิ, ภิยฺโย ภิยฺโย โอกปฺเปมิ, ภิยฺโย ภิยฺโย อธิมุจฺจามิ, “สพฺเพ สงฺขารา ทุกฺขา”ติ ภิยฺโย ภิยฺโย ปสีทามิ ฯเปฯ “สพฺเพ ธมฺมา อนตฺตา”ติ ภิยฺโย ภิยฺโย ปสีทามิ ฯเปฯ\textbf{.} “ยํ กิญฺจิ สมุทยธมฺมํ, สพฺพํ ตํ นิโรธธมฺมนฺ”ติ ภิยฺโย ภิยฺโย ปสีทามิ, ภิยฺโย ภิยฺโย สทฺทหามิ, ภิยฺโย ภิยฺโย โอกปฺเปมิ, ภิยฺโย ภิยฺโย อธิมุจฺจามีติ เอส ภิยฺโย ปสีทามิ\textbf{.}}

\prose{\textbf{สุตฺวาน มุนิโน วโจ}ติ \textbf{มุนี}ติ โมนํ วุจฺจติ ญาณํ ฯเปฯ สงฺคชาลมติจฺจ โส มุนิ.\csromanspacer{} \textbf{สุตฺวาน มุนิโน วโจ}ติ ตุยฺหํ วจนํ พฺยปฺปถํ เทสนํ อนุสาสนํ อนุสิฏฺฐํ สุตฺวาน อุคฺคเหตฺวาน อุปธารยิตฺวาน อุปลกฺขยิตฺวานาติ สุตฺวาน มุนิโน วโจ.}

\prose{\textbf{วิวฏฺฏจฺฉโท สมฺพุทฺโธ}ติ \textbf{ฉทนนฺ}ติ ปญฺจ ฉทนานิ ตณฺหาฉทนํ ทิฏฺฐิฉทนํ กิเลสฉทนํ ทุจฺจริตฉทนํ อวิชฺชาฉทนํ, ตานิ ฉทนานิ พุทฺธสฺส ภควโต วิวฏานิ วิทฺธํสิตานิ สมุคฺฆาฏิตานิ ปหีนานิ สมุจฺฉินฺนานิ วูปสนฺตานิ ปฏิปฺปสฺสทฺธานิ อภพฺพุปฺปตฺติกานิ ญาณคฺคินา ทฑฺฒานิ, ตสฺมา \textbf{พุทฺโธ} วิวฏฺฏจฺฉโท.\csromanspacer{} \textbf{พุทฺโธ}ติ โย โส ภควา ฯเปฯ สจฺฉิกา ปญฺญตฺติ, ยทิทํ \textbf{พุทฺโธ}ติ วิวฏฺฏจฺฉโท สมฺพุทฺโธ.}

\prose{\textbf{อขิโล ปฏิภานวา}ติ \textbf{อขิโล}ติ ราโค ขิโล, โทโส ขิโล, โมโห ขิโล, โกโธ ขิโล, อุปนาโห ฯเปฯ สพฺพากุสลาภิสงฺขารา ขิลา, เต ขิลา พุทฺธสฺส ภควโต ปหีนา อุจฺฉินฺนมูลา ตาลาวตฺถุกตา อนภาวํกตา อายติํ อนุปฺปาทธมฺมา, ตสฺมา \textbf{พุทฺโธ} \textbf{อขิโล}.}

\prose{\textbf{ปฏิภานวา}ติ ตโย ปฏิภานวนฺโต ปริยตฺติ ปฏิภานวา ปริปุจฺฉาปฏิภานวา อธิคมปฏิภานวา.\csromanspacer{} กตโม ปริยตฺติปฏิภานวา, อิเธกจฺจสฺส พุทฺธวจนํ\footnote{ปกติยา (ก)} ปริยาปุตํ\footnote{ปริยาปุฏํ (สฺยา, ก)} โหติ สุตฺตํ เคยฺยํ เวยฺยากรณํ คาถา อุทานํ อิติวุตฺตกํ ชาตกํ อพฺภุตธมฺมํ เวทลฺลํ, ตสฺส ปริยตฺติํ นิสฺสาย ปฏิภาติ\footnote{ปฏิภายติ (ก)}, อยํ ปริยตฺติปฏิภานวา.}

\prose{กตโม ปริปุจฺฉาปฏิภานวา, อิเธกจฺโจ ปริปุจฺฉิตา โหติ อตฺเถ จ ญาเย จ ลกฺขเณ จ การเณ จ ฐานาฐาเน จ, ตสฺส ปริปุจฺฉํ นิสฺสาย ปฏิภาติ, อยํ ปริปุจฺฉาปฏิภานวา.}

\csromanheader{2}{18. ปารายนานุคีติคาถานิทฺเทส}
\prose{\csromanfolio{225}กตโม อธิคมปฏิภานวา, อิเธกจฺจสฺส อธิคตา โหนฺติ จตฺตาโร สติปฏฺฐานา จตฺตาโร สมฺมปฺปธานา จตฺตาโร อิทฺธิปาทา ปญฺจินฺทฺริยานิ ปญฺจ พลานิ สตฺต โพชฺฌงฺคา อริโย อฏฺฐงฺคิโก มคฺโค จตฺตาโร อริยมคฺคา จตฺตาริ สามญฺญผลานิ จตสฺโส ปฏิสมฺภิทาโย ฉ อภิญฺญาโย, ตสฺส อตฺโถ ญาโต ธมฺโม ญาโต นิรุตฺติ ญาตา, อตฺเถ ญาเต อตฺโถ ปฏิภาติ, ธมฺเม ญาเต ธมฺโม ปฏิภาติ, นิรุตฺติยา ญาตาย นิรุตฺติ ปฏิภาติ, อิเมสุ ตีสุ ญาเณสุ ญาณํ ปฏิภานปฏิสมฺภิทา\textbf{.}\csromanspacer{} ภควา อิมาย ปฏิภานปฏิสมฺภิทาย อุเปโต สมุเปโต อุปาคโต สมุปาคโต อุปปนฺโน สมุปปนฺโน สมนฺนาคโต, ตสฺมา พุทฺโธ ปฏิภานวา\textbf{.}\csromanspacer{} ยสฺส ปริยตฺติ นตฺถิ, ปริปุจฺฉา นตฺถิ, อธิคโม นตฺถิ, กิํ ตสฺส ปฏิภายิสฺสตีติ อขิโล ปฏิภานวา\textbf{.}\csromanspacer{} เตนาห เถโร ปิงฺคิโย\csromandash{}}

\csromangathameasure{เอส ภิยฺโย ปสีทามิ, สุตฺวาน มุนิโน วโจ. \\ วิวฏฺฏจฺฉโท สมฺพุทฺโธ, อขิโล ปฏิภานวาติ.}
\csromangathagroup{\csromangathastanza{เอส ภิยฺโย ปสีทามิ, สุตฺวาน มุนิโน วโจ. \\ วิวฏฺฏจฺฉโท สมฺพุทฺโธ, อขิโล ปฏิภานวาติ.}}

\proseitem{119}{อธิเทเว\footnote{อติเทเว (ก)} \textbf{อภิญฺญาย, สพฺพํ เวทิ ปโรปรํ}\footnote{ปโรวรํ (สี-ฏฺฐ)}\textbf{.}}

\prose{\textbf{ปญฺหานนฺตกโร สตฺถา, กงฺขีนํ ปฏิชานตํ. (18)}}

\prose{\textbf{อธิเทเว อภิญฺญายา}ติ \textbf{เทวา}ติ ตโย \textbf{เทวา} สมฺมุติ\textbf{เทวา}\footnote{สมฺมติเทวา (สฺยา) 85 ปิฏฺเฐปิ.} อุปปตฺติ\textbf{เทวา} วิสุทฺธิ\textbf{เทวา.}\csromanspacer{} กตเม สมฺมุติ\textbf{เทวา}, สมฺมุติ\textbf{เทวา} วุจฺจนฺติ ราชาโน\footnote{กตเม สมฺมติเทวา ราชาโน (สฺยา) เอวมุปริปิ.} จ ราชกุมารา จ เทวิโย จ, อิเม วุจฺจนฺติ สมฺมุติ\textbf{เทวา.}\csromanspacer{} กตเม อุปปตฺติ\textbf{เทวา}, อุปปตฺติ\textbf{เทวา} วุจฺจนฺติ จาตุมหาราชิกา \textbf{เทวา} ตาวติํสา \textbf{เทวา} ฯเปฯ พฺรหฺมกายิกา \textbf{เทวา}, เย จ \textbf{เทวา} ตทุตฺตริ, อิเม วุจฺจนฺติ อุปปตฺติ\textbf{เทวา.}\csromanspacer{} กตเม วิสุทฺธิ\textbf{เทวา}, วิสุทฺธิ\textbf{เทวา} วุจฺจนฺติ ตถาคตา ตถาคตสาวกา อรหนฺโต ขีณาสวา, เย จ ปจฺเจกสมฺพุทฺธา, อิเม วุจฺจนฺติ วิสุทฺธิ\textbf{เทวา.}\csromanspacer{} ภควา สมฺมุติเทเว อธิ\textbf{เทวา}ติ อภิญฺญาย, อุปปตฺติเทเว อธิ\textbf{เทวา}ติ อภิญฺญาย, วิสุทฺธิเทเว อธิ\textbf{เทวา}ติ อภิญฺญาย ชานิตฺวา ตุลยิตฺวา ตีรยิตฺวา วิภาวยิตฺวา วิภูตํ กตฺวาติ อธิเทเว อภิญฺญาย\textbf{.}}

\prose{\csromanfolio{226}\textbf{สพฺพํ เวทิ ปโรปรนฺ}ติ ภควา อตฺตโน จ ปเรสญฺจ อธิเทวกเร ธมฺเม เวทิ อญฺญาสิ อผสฺสิ ปฏิวิชฺฌิ.\csromanspacer{} กตเม อตฺตโน อธิเทวกรา ธมฺมา, สมฺมาปฏิปทา อนุโลมปฏิปทา อปจฺจนีกปฏิปทา อนฺวตฺถปฏิปทา ธมฺมานุธมฺมปฏิปทา สีเลสุ ปริปูรการิตา อินฺทฺริเยสุ คุตฺตทฺวารตา โภชเน มตฺตญฺญุตา ชาคริยานุโยโค สติสมฺปชญฺญํ จตฺตาโร สติปฏฺฐานา ฯเปฯ อริโย อฏฺฐงฺคิโก มคฺโค, อิเม วุจฺจนฺติ อตฺตโน อธิเทวกรา ธมฺมา.}

\prose{กตเม ปเรสํ อธิเทวกรา ธมฺมา, สมฺมาปฏิปทา ฯเปฯ อริโย อฏฺฐงฺคิโก มคฺโค, อิเม วุจฺจนฺติ ปเรสํ อธิเทวกรา ธมฺมา.\csromanspacer{} เอวํ ภควา อตฺตโน จ ปเรสญฺจ อธิเทวกเร ธมฺเม เวทิ อญฺญาสิ อผสฺสิ ปฏิวิชฺฌีติ สพฺพํ เวทิ ปโรปรํ.}

\prose{\textbf{ปญฺหานนฺตกโร สตฺถา}ติ ภควา ปารายนิกปญฺหานํ อนฺตกโร ปริยนฺตกโร ปริจฺเฉทกโร ปริวฏุมกโร.\csromanspacer{} สภิยปญฺหานํ\footnote{ปริสปญฺหานํ (สฺยา), ปิงฺคิยปญฺหานํ (ก)} อนฺตกโร ปริยนฺตกโร ปริจฺเฉทกโร ปริวฏุมกโร.\csromanspacer{} สกฺกปญฺหานํ.\csromanspacer{} สุยามปญฺหานํ.\csromanspacer{} ภิกฺขุปญฺหานํ.\csromanspacer{} ภิกฺขุนีปญฺหานํ.\csromanspacer{} อุปาสกปญฺหานํ.\csromanspacer{} อุปาสิกาปญฺหานํ.\csromanspacer{} ราชปญฺหานํ.\csromanspacer{} ขตฺติยปญฺหานํ.\csromanspacer{} พฺราหฺมณปญฺหานํ.\csromanspacer{} เวสฺสปญฺหานํ.\csromanspacer{} สุทฺทปญฺหานํ.\csromanspacer{} เทวปญฺหานํ.\csromanspacer{} พฺรหฺมปญฺหานํ อนฺตกโร ปริยนฺตกโร ปริจฺเฉทกโร ปริวฏุมกโรติ ปญฺหานนฺตกโร.\csromanspacer{} \textbf{สตฺถา}ติ ภควา สตฺถวาโห.\csromanspacer{} ยถา สตฺถวาโห สตฺเถ กนฺตารํ ตาเรติ, โจรกนฺตารํ ตาเรติ, วาฬกนฺตารํ ตาเรติ, ทุพฺภิกฺขกนฺตารํ ตาเรติ, นิรุทกกนฺตารํ ตาเรติ อุตฺตาเรติ นิตฺตาเรติ\footnote{นิตฺถาเรติ (สฺยา, ก)} ปตาเรติ, เขมนฺตภูมิํ สมฺปาเปติ.\csromanspacer{} เอวเมว ภควา สตฺถวาโห สตฺเต กนฺตารํ ตาเรติ.\csromanspacer{} ชาติกนฺตารํ ตาเรติ.\csromanspacer{} ชรากนฺตารํ.\csromanspacer{} พฺยาธิกนฺตารํ.\csromanspacer{} มรณกนฺตารํ.\csromanspacer{} โสกปริเทวทุกฺขโทมนสฺสุปายาสกนฺตารํ ตาเรติ, ราคกนฺตารํ ตาเรติ, โทสกนฺตารํ.\csromanspacer{} โมหกนฺตารํ.\csromanspacer{} มานกนฺตารํ.\csromanspacer{} ทิฏฺฐิกนฺตารํ.\csromanspacer{} กิเลสกนฺตารํ.\csromanspacer{} ทุจฺจริตกนฺตารํ ตาเรติ, ราคคหนํ ตาเรติ, โทสคหนํ.\csromanspacer{} โมหคหนํ.\csromanspacer{} ทิฏฺฐิคหนํ.\csromanspacer{} กิเลสคหนํ.\csromanspacer{} ทุจฺจริตคหนํ ตาเรติ อุตฺตาเรติ นิตฺตาเรติ ปตาเรติ, เขมนฺตํ อมตํ นิพฺพานํ สมฺปาเปตีติ เอวมฺปิ ภควา สตฺถวาโห.}

\prose{อถ วา ภควา เนตา วิเนตา อนุเนตา ปญฺญาเปตา นิชฺฌาเปตา เปกฺขตา ปสาเทตาติ เอวํ ภควา สตฺถวาโห.\csromanspacer{}\csromanspacer{}อถ วา ภควา}

\vspace{\tipitakaparskip}
\csromancenter{ปารายนานุคีติคาถานิทฺเทส อนุปฺปนฺนสฺส มคฺคสฺส อุปฺปาเทตา, อสญฺชาตสฺส มคฺคสฺส สญฺชเนตา, อนกฺขาตสฺส มคฺคสฺส อกฺขาตา, มคฺคญฺญู มคฺควิทู มคฺคโกวิโท มคฺคานุคา จ ปน เอตรหิ สาวกา วิหรนฺติ ปจฺฉา สมนฺนาคตาติ เอวมฺปิ ภควา สตฺถวาโหติ ปญฺหานนฺตกโร สตฺถา.}
\prose{\csromanfolio{227}\textbf{กงฺขีนํ ปฏิชานตนฺ}ติ สกงฺขา อาคนฺตฺวา นิกฺกงฺขา สมฺปชฺชนฺติ, สลฺเลขา อาคนฺตฺวา นิลฺเลขา สมฺปชฺชนฺติ, สเทฺวฬฺหกา อาคนฺตฺวา นิเทฺวฬฺหกา สมฺปชฺชนฺติ, สวิจิกิจฺฉา อาคนฺตฺวา นิพฺพิจิกิจฺฉา สมฺปชฺชนฺติ, สราคา อาคนฺตฺวา วีตราคา สมฺปชฺชนฺติ, สโทสา อาคนฺตฺวา วีตโทสา สมฺปชฺชนฺติ, สโมหา อาคนฺตฺวา วีตโมหา สมฺปชฺชนฺติ, สกิเลสา อาคนฺตฺวา นิกฺกิเลสา สมฺปชฺชนฺตีติ กงฺขีนํ ปฏิชานตํ.\csromanspacer{} เตนาห เถโร ปิงฺคิโย\csromandash{}}

\csromangathameasure{อธิเทเว อภิญฺญาย, สพฺพํ เวทิ ปโรปรํ. \\ ปญฺหานนฺตกโร สตฺถา, กงฺขีนํ ปฏิชานตนฺติ.}
\csromangathagroup{\csromangathastanza{อธิเทเว อภิญฺญาย, สพฺพํ เวทิ ปโรปรํ. \\ ปญฺหานนฺตกโร สตฺถา, กงฺขีนํ ปฏิชานตนฺติ.}}

\proseitem{120}{อสํหีรํ อสํกุปฺปํ,}

\prose{\textbf{ยสฺส นตฺถิ อุปมา กฺวจิ.}}

\prose{\textbf{อทฺธา คมิสฺสามิ น เมตฺถ กงฺขา,}}

\csromanheader{2}{\textbf{เอวํ มํ ธาเรหิ อธิมุตฺตจิตฺตํ. (19)}}
\prose{\textbf{อสํหีรํ อสํกุปฺปนฺ}ติ อสํหีรํ วุจฺจติ อมตํ นิพฺพานํ.\csromanspacer{} โย โส สพฺพสงฺขารสมโถ สพฺพูปธิปฏินิสฺสคฺโค ตณฺหกฺขโย วิราโค นิโรโธ นิพฺพานํ.\csromanspacer{} \textbf{อสํหีรนฺ}ติ ราเคน โทเสน โมเหน โกเธน อุปนาเหน มกฺเขน ปฬาเสน อิสฺสาย มจฺฉริเยน มายาย สาเฐยฺเยน ถมฺเภน สารมฺเภน มาเนน อติมาเนน มเทน ปมาเทน สพฺพกิเลเสหิ สพฺพทุจฺจริเตหิ สพฺพปริฬาเหหิ สพฺพาสเวหิ สพฺพทรเถหิ สพฺพสนฺตาเปหิ สพฺพากุสลาภิสงฺขาเรหิ อสํหาริยํ นิพฺพานํ นิจฺจํ ธุวํ สสฺสตํ อวิปริณามธมฺมนฺติ อสํหีรํ.}

\prose{\textbf{อสํกุปฺปนฺ}ติ อสํกุปฺปํ วุจฺจติ อมตํ นิพฺพานํ.\csromanspacer{} โย โส สพฺพสงฺขารสมโถ ฯเปฯ นิโรโธ นิพฺพานํ.\csromanspacer{} นิพฺพานสฺส\footnote{ยสฺส (สฺยา)} น อุปฺปาโท ปญฺญายติ, วโย นตฺถิ, น ตสฺส อญฺญถตฺตํ\footnote{ตสฺส อญฺญทตฺถุ (สฺยา)} ปญฺญายติ, นิพฺพานํ นิจฺจํ ธุวํ สสฺสตํ อวิปริณามธมฺมนฺติ อสํหีรํ อสํกุปฺปํ.}

\prose{\csromanfolio{228}\textbf{ยสฺส นตฺถิ อุปมา กฺวจี}ติ \textbf{ยสฺสา}ติ นิพฺพานสฺส.\csromanspacer{} \textbf{นตฺถิ อุปมา}ติ อุปมา นตฺถิ, อุปนิธา นตฺถิ, สทิสํ นตฺถิ, ปฏิภาโค นตฺถิ น สติ น สํวิชฺชติ นุปลพฺภติ.\csromanspacer{} \textbf{กฺวจี}ติ กฺวจิ กิมฺหิจิ กตฺถจิ อชฺฌตฺตํ วา พหิทฺธา วา อชฺฌตฺตพหิทฺธา วาติ ยสฺส นตฺถิ อุปมา กฺวจิ.}

\prose{\textbf{อทฺธา คมิสฺสามิ น เมตฺถ กงฺขา}ติ \textbf{อทฺธา}ติ เอกํสวจนํ นิสฺสํสยวจนํ นิกฺกงฺขวจนํ อเทฺวชฺฌวจนํ อเทฺวฬฺหกวจนํ นิโยควจนํ อปณฺณกวจนํ อวิรทฺธวจนํ อวตฺถาปนวจนเมตํ \textbf{อทฺธา}ติ.\csromanspacer{} \textbf{คมิสฺสามี}ติ คมิสฺสามิ อธิคมิสฺสามิ ผสฺสิสฺสามิ สจฺฉิกริสฺสามีติ \textbf{อทฺธา} คมิสฺสามิ.\csromanspacer{} \textbf{น เมตฺถ กงฺขา}ติ \textbf{เอตฺถา}ติ นิพฺพาเน กงฺขา นตฺถิ, วิจิกิจฺฉา นตฺถิ, เทฺวฬฺหกํ นตฺถิ, สํสโย นตฺถิ น อตฺถิ น สํวิชฺชติ นุปลพฺภติ, ปหีโน สมุจฺฉินฺโน วูปสนฺโต ปฏิปฺปสฺสทฺโธ อภพฺพุปฺปตฺติโก ญาณคฺคินา ทฑฺโฒติ \textbf{อทฺธา} คมิสฺสามิ น เมตฺถ กงฺขา.}

\prose{\textbf{เอวํ มํ ธาเรหิ อธิมุตฺตจิตฺตนฺ}ติ \textbf{เอวํ มํ ธาเรหี}ติ เอวํ มํ อุปลกฺเขหิ.\csromanspacer{} \textbf{อธิมุตฺตจิตฺตนฺ}ติ นิพฺพานนินฺนํ นิพฺพานโปณํ นิพฺพานปพฺภารํ นิพฺพานาธิมุตฺตนฺติ เอวํ มํ ธาเรหิ อธิมุตฺตจิตฺตนฺติ.\csromanspacer{} เตนาห เถโร ปิงฺคิโย\csromandash{}}

\csromangathameasure{อสํหีรํ อสํกุปฺปํ, \\ ยสฺส นตฺถิ อุปมา กฺวจิ.}
\csromangathagroup{\csromangathastanza{อสํหีรํ อสํกุปฺปํ, \\ ยสฺส นตฺถิ อุปมา กฺวจิ.}}

\prose{\textbf{อทฺธา คมิสฺสามิ น เมตฺถ กงฺขา},}

\prose{\textbf{เอวํ มํ ธาเรหิ อธิมุตฺตจิตฺตนฺ}ติ.}

\vspace{\tipitakaparskip}
\csromancenter{ปารายนานุคีติคาถานิทฺเทโส อฏฺฐารสโม.}
\csromancenter{\textbf{ปารายนวคฺโค สมตฺโต.}}
\csromanmatikaanchor{mtk.40}
\csromantocmarkref{subsection}{19. ขคฺควิสาณสุตฺตนิทฺเทส}{mtk.40}
\bookmark[dest={mtk.40},level=2]{19. ขคฺควิสาณสุตฺตนิทฺเทส}
\csromanheader{2}{19. ขคฺควิสาณสุตฺตนิทฺเทส}
\chapterhead{\textbf{ปฐมวคฺค}}
\csromanhangingbegin
\hangingheaditem{211}{\csromanfolio{229}สพฺเพสุ ภูเตสุ นิธาย ทณฺฑํ,}
\hangingbody{อวิเหฐยํ อญฺญตรมฺปิ เตสํ.}
\hangingbody{น ปุตฺตมิจฺเฉยฺย กุโต สหายํ,}
\hangingbody{เอโก จเร ขคฺควิสาณกปฺโป.}
\csromanhangingend

\prose{\textbf{สพฺเพสุ ภูเตสุ นิธาย ทณฺฑนฺติ สพฺเพสู}ติ สพฺเพน สพฺพํ สพฺพถา สพฺพํ อเสสํ นิสฺเสสํ ปริยาทิยนวจนเมตํ สพฺเพสูติ.\csromanspacer{} \textbf{ภูเตสู}ติ ภูตา วุจฺจนฺติ ตสา จ ถาวรา จ.\csromanspacer{} ตสาติ เยสํ ตสิตตณฺหา อปฺปหี\textbf{นา}, เยสญฺจ ภยเภรวา อปฺปหี\textbf{นา}.\csromanspacer{} กิํการณา วุจฺจนฺติ ตสา, เต ตสนฺติ อุตฺตสนฺติ ปริตสนฺติ ภายนฺติ สนฺตาสํ อาปชฺชนฺติ, ตํการณา วุจฺจนฺติ ตสา.\csromanspacer{} ถาวราติ เยสํ ตสิตตณฺหา ปหี\textbf{นา}, เยสญฺจ ภยเภรวา ปหี\textbf{นา}.\csromanspacer{} กิํการณา วุจฺจนฺติ ถาวรา, เต น ตสนฺติ น อุตฺตสนฺติ น ปริตสนฺติ น ภายนฺติ น สนฺตาสํ อาปชฺชนฺติ, ตํการณา วุจฺจนฺติ ถาวรา.\csromanspacer{} \textbf{ทณฺฑนฺ}ติ ตโย ทณฺฑา กายทณฺโฑ วจีทณฺโฑ มโนทณฺโฑ.\csromanspacer{} ติวิธํ กายทุจฺจริตํ กายทณฺโฑ, จตุพฺพิธํ วจีทุจฺจริตํ วจีทณฺโฑ, ติวิธํ มโนทุจฺจริตํ มโนทณฺโฑ.\csromanspacer{} \textbf{สพฺเพสุ ภูเตสุ นิธาย ทณฺฑนฺ}ติ สพฺเพสุ ภูเตสุ ทณฺฑํ นิธาย นิทหิตฺวา โอโรปยิตฺวา สโมโรปยิตฺวา นิกฺขิปิตฺวา ปฏิปฺปสฺสมฺภิตฺวาติ สพฺเพสุ ภูเตสุ นิธาย ทณฺฑํ.}

\prose{\textbf{อวิเหฐยํ อญฺญตรมฺปิ เตสนฺ}ติ เอกเมกมฺปิ สตฺตํ ปาณิ\textbf{นา} วา เลฑฺฑุ\textbf{นา} วา ทณฺเฑน วา สตฺเถน วา อนฺทุยา\footnote{อรุยา (สฺยา), อทฺทุยา (ก)} วา รชฺชุยา วา อวิเหฐยนฺโต, สพฺเพปิ สตฺเต ปาณิ\textbf{นา} วา เลฑฺฑุ\textbf{นา} วา ทณฺเฑน วา สตฺเถน วา อนฺทุยา วา รชฺชุยา วา อวิเหฐยนฺโตติ อวิเหฐยํ อญฺญตรมฺปิ เตสํ.}

\prose{\textbf{น ปุตฺตมิจฺเฉยฺย กุโต สหายนฺ}ติ \textbf{นา}ติ ปฏิกฺเขโป.\csromanspacer{} \textbf{ปุตฺตา}ติ จตฺตาโร ปุตฺตา อตฺรโช ปุตฺโต เขตฺตโช ปุตฺโต ทินฺนโก ปุตฺโต อนฺเตวาสิโก ปุตฺโต.\csromanspacer{} \textbf{สหายนฺ}ติ สหายา วุจฺจนฺติ เยหิ สห อาคมนํ ผาสุ, คมนํ ผาสุ, คม\textbf{นา}คมนํ ผาสุ, ฐานํ ผาสุ, นิสชฺชนํ ผาสุ, สยนํ\footnote{นิปชฺชนํ (สฺยา)} ผาสุ, อาลปนํ ผาสุ, สลฺลปนํ ผาสุ, อุลฺลปนํ ผาสุ, สมุลฺลปนํ \csromanfolio{230}ผาสุ.\csromanspacer{} \textbf{น ปุตฺตมิจฺเฉยฺย กุโต สหายนฺ}ติ ปุตฺตมฺปิ น อิจฺเฉยฺย น สาทิเยยฺย น ปตฺถเยยฺย น ปิหเยยฺย นาภิชปฺเปยฺย, กุโต มิตฺตํ วา สนฺทิฏฺฐํ วา สมฺภตฺตํ วา สหายํ วา อิจฺเฉยฺย\footnote{อิจฺฉิสฺสติ (สฺยา) เอวมีทิเสสุ ปเทสุ อนาคตวิภตฺติยา.} สาทิเยยฺย ปตฺถเยยฺย ปิหเยยฺย อภิชปฺเปยฺยาติ น ปุตฺตมิจฺเฉยฺย กุโต สหายํ.}

\prose{\textbf{เอโก จเร ขคฺควิสาณกปฺโป}ติ \textbf{เอโก}ติ โส ปจฺเจกสมฺพุทฺโธ ปพฺพชฺชาสงฺขาเตน \textbf{เอโก}, อทุติยฏฺเฐน \textbf{เอโก}, ตณฺหาย ปหานฏฺเฐน\footnote{ตณฺหาปหานฏฺเฐน (สฺยา) ขุ 7. 361 ปิฏฺเฐ.} \textbf{เอโก}, เอกนฺตวีตราโคติ \textbf{เอโก}, เอกนฺตวีตโทโสติ \textbf{เอโก}, เอกนฺตวีตโมโหติ \textbf{เอโก}, เอกนฺตนิกฺกิเลโสติ \textbf{เอโก}, เอกายนมคฺคํ คโตติ \textbf{เอโก}, \textbf{เอโก} อนุตฺตรํ ปจฺเจกสมฺโพธิํ อภิสมฺพุทฺโธติ \textbf{เอโก}.}

\prose{กถํ โส ปจฺเจกสมฺพุทฺโธ ปพฺพชฺชาสงฺขาเตน \textbf{เอโก}, โส ปจฺเจกสมฺพุทฺโธ สพฺพํ ฆราวาสปลิโพธํ ฉินฺทิตฺวา ปุตฺตทารปลิโพธํ ฉินฺทิตฺวา ญาติพลิโพธํ ฉินฺทิตฺวา สนฺนิธิปลิโพธํ ฉินฺทิตฺวา เกสมสฺสุํ โอหาเรตฺวา กาสายานิ วตฺถานิ อจฺฉาเทตฺวา อคารสฺมา อนคาริยํ ปพฺพชิตฺวา อกิญฺจนภาวํ อุปคนฺตฺวา \textbf{เอโก} จรติ วิหรติ อิริยติ วตฺเตติ ปาเลติ ยเปติ ยาเปตีติ เอวํ โส ปจฺเจกสมฺพุทฺโธ ปพฺพชฺชาสงฺขาเตน \textbf{เอโก}.}

\prose{กถํ โส ปจฺเจกสมฺพุทฺโธ อทุติยฏฺเฐน \textbf{เอโก}, โส เอวํ ปพฺพชิโต สมาโน \textbf{เอโก} อรญฺญวนปตฺถานิ ปนฺตานิ เสนาสนานิ ปฏิเสวติ อปฺปสทฺทานิ อปฺปนิคฺโฆสานิ วิชนวาตานิ มนุสฺสราหสฺเสยฺยกานิ ปฏิสลฺลานสารุปฺปานิ, โส \textbf{เอโก} คจฺฉติ, \textbf{เอโก} ติฏฺฐติ, \textbf{เอโก} นิสีทติ, \textbf{เอโก} เสยฺยํ กปฺเปติ, \textbf{เอโก} คามํ ปิณฺฑาย ปวิสติ, \textbf{เอโก} อภิกฺกมติ, \textbf{เอโก} ปฏิกฺกมติ, \textbf{เอโก} รโห นิสีทติ, \textbf{เอโก} จงฺกมํ อธิฏฺฐาติ, \textbf{เอโก} จรติ วิหรติ อิริยติ วตฺเตติ ปาเลติ ยเปติ ยาเปตีติ เอวํ โส ปจฺเจกสมฺพุทฺโธ อทุติยฏฺเฐน \textbf{เอโก}.}

\prose{กถํ โส ปจฺเจกสมฺพุทฺโธ ตณฺหาย ปหานฏฺเฐน \textbf{เอโก}, โส เอวํ \textbf{เอโก} อทุติโย อปฺปมตฺโต อาตาปี ปหิตตฺโต วิหรนฺโต มหาปธานํ ปทหนฺโต มารํ สเสนกํ นมุจิํ กณฺหํ ปมตฺตพนฺธุํ วิธเมตฺวา จ ตณฺหาชาลินิํ วิสริตํ วิสตฺติกํ ปชหิ วิโนเทสิ พฺยนฺตี-อกาสิ อนภาวํคเมสิ.}

\csromanheader{2}{19. ขคฺควิสาณสุตฺตนิทฺเทส}
\csromangathameasure{“ตณฺหาทุติโย ปุริโส, ทีฆมทฺธาน สํสรํ. \\ อิตฺถภาวญฺญถาภาวํ, สํสารํ นาติวตฺตติ. \\ เอตมาทีนวํ ญ ตฺวา, ตณฺหํ ทุกฺขสฺส สมฺภวํ. \\ วีตตโณฺห อนาทาโน, สโต ภิกฺขุ ปริพฺพเช”ติ.}
\csromangathagroup{\csromangathastanza{\csromanfolio{231}“ตณฺหาทุติโย ปุริโส, ทีฆมทฺธาน สํสรํ. \\ อิตฺถภาวญฺญถาภาวํ, สํสารํ นาติวตฺตติ.}\csromangathastanza{เอตมาทีนวํ ญ ตฺวา, ตณฺหํ ทุกฺขสฺส สมฺภวํ. \\ วีตตโณฺห อนาทาโน, สโต ภิกฺขุ ปริพฺพเช”ติ.}}

\prose{เอวํ โส ปจฺเจกสมฺพุทฺโธ ตณฺหาย ปหานฏฺเฐน เอโก.}

\prose{กถํ โส ปจฺเจกสมฺพุทฺโธ เอกนฺตวีตราโคติ เอโก, ราคสฺส ปหีนตฺตา เอกนฺตวีตราโคติ เอโก, โทสสฺส ปหีนตฺตา เอกนฺตวีตโทโสติ เอโก, โมหสฺส ปหีนตฺตา เอกนฺตวีตโมโหติ เอโก, กิเลสานํ ปหีนตฺตา เอกนฺตนิกฺกิเลโสติ เอโก, เอวํ โส ปจฺเจกสมฺพุทฺโธ เอกนฺตวีตราโคติ เอโก.}

\prose{กถํ โส ปจฺเจกสมฺพุทฺโธ เอกายนมคฺคํ คโตติ เอโก, เอกายนมคฺโค วุจฺจติ จตฺตาโร สติปฏฺฐานา จตฺตาโร สมฺมปฺปธานา จตฺตาโร อิทฺธิปาทา ปญฺจินฺทฺริยานิ ปญฺจ พลานิ สตฺต โพชฺฌงฺคา อริโย อฏฺฐงฺคิโก มคฺโค.}

\csromangathameasure{“เอกายนํ ชาติขยนฺตทสฺสี, \\ มคฺคํ ปชานาติ หิตานุกมฺปี. \\ เอเตน มคฺเคน ตริํสุ ปุพฺเพ, \\ ตริสฺสนฺติ เย จ ตรนฺติ โอฆนฺ”ติ.}
\csromangathagroup{\csromangathastanza{“เอกายนํ ชาติขยนฺตทสฺสี, \\ มคฺคํ ปชานาติ หิตานุกมฺปี. \\ เอเตน มคฺเคน ตริํสุ ปุพฺเพ, \\ ตริสฺสนฺติ เย จ ตรนฺติ โอฆนฺ”ติ.}}

\prose{เอวํ โส ปจฺเจกสมฺพุทฺโธ เอกายนมคฺคํ คโตติ เอโก.}

\prose{กถํ โส ปจฺเจกสมฺพุทฺโธ เอโก อนุตฺตรํ ปจฺเจกสมฺโพธิํ อภิสมฺพุทฺโธติ เอโก, โพธิ วุจฺจติ จตูสุ มคฺเคสุ ญาณํ.\csromanspacer{} ปญฺญา ปญฺญินฺทฺริยํ ปญฺญาพลํ ธมฺมวิจยสมฺโพชฺฌงฺโค วีมํสา วิปสฺสนา สมฺมาทิฏฺฐิ.\csromanspacer{} โส ปจฺเจกสมฺพุทฺโธ มคฺคปจฺเจกสมฺพุทฺโธ ญาณปจฺเจกสมฺพุทฺโธ “สพฺเพ สงฺขารา อนิจฺจา”ติ พุชฺฌิ, “สพฺเพ สงฺขารา ทุกฺขา”ติ พุชฺฌิ, “สพฺเพ ธมฺมา อนตฺตา”ติ พุชฺฌิ. “อวิชฺชาปจฺจยา สงฺขารา”ติ พุชฺฌิ, “สงฺขารปจฺจยา วิญฺญาณนฺ”ติ พุชฺฌิ, “วิญฺญาณปจฺจยา นามรูปนฺ”ติ พุชฺฌิ, “นามรูปปจฺจยา สฬายตนนฺ”ติ พุชฺฌิ, “สฬายตนปจฺจยา ผสฺโส”ติ พุชฺฌิ, “ผสฺสปจฺจยา เวทนา”ติ พุชฺฌิ, “เวทนาปจฺจยา ตณฺหา”ติ พุชฺฌิ, “ตณฺหาปจฺจยา อุปาทานนฺ”ติ พุชฺฌิ, “อุปาทานปจฺจยา ภโว”ติ พุชฺฌิ, “ภวปจฺจยา ชาตี”ติ พุชฺฌิ,}

\prose{\csromanfolio{232}“ชาติปจฺจยา ชรามรณนฺ”ติ พุชฺฌิ. “อวิชฺชานิโรธา สงฺขารนิโรโธ”ติ พุชฺฌิ, “สงฺขารนิโรธา วิญฺญาณนิโรโธ”ติ พุชฺฌิ, “วิญฺญาณนิโรธา นามรูปนิโรโธ”ติ พุชฺฌิ, “นามรูปนิโรธา สฬายตนนิโรโธ”ติ พุชฺฌิ, “สฬายตนนิโรธา ผสฺสนิโรโธ”ติ พุชฺฌิ, “ผสฺสนิโรธา เวทนานิโรโธ”ติ พุชฺฌิ, “เวทนานิโรธา ตณฺหานิโรโธ”ติ พุชฺฌิ, “ตณฺหานิโรธา อุปาทานนิโรโธ”ติ พุชฺฌิ, “อุปาทานนิโรธา ภวนิโรโธ”ติ พุชฺฌิ, “ภวนิโรธา ชาตินิโรโธ”ติ พุชฺฌิ, “ชาตินิโรธา ชรามรณนิโรโธ”ติ พุชฺฌิ. “อิทํ ทุกฺขนฺ”ติ พุชฺฌิ, “อยํ ทุกฺขสมุทโย”ติ พุชฺฌิ, “อยํ ทุกฺขนิโรโธ”ติ พุชฺฌิ, “อยํ ทุกฺขนิโรธคามินี ปฏิปทา”ติ พุชฺฌิ. “อิเม อาสวา”ติ พุชฺฌิ, “อยํ อาสวสมุทโย”ติ พุชฺฌิ ฯเปฯ “อยํ อาสวนิโรธคามินี ปฏิปทา”ติ พุชฺฌิ. “อิเม ธมฺมา อภิญฺเญยฺยา”ติ พุชฺฌิ, “อิเม ธมฺมา ปหาตพฺพา”ติ พุชฺฌิ, “อิเม ธมฺมา สจฺฉิกาตพฺพา”ติ พุชฺฌิ, “อิเม ธมฺมา ภาเวตพฺพา”ติ พุชฺฌิ.\csromanspacer{} ฉนฺนํ ผสฺสายตนานํ สมุทยญฺจ อตฺถงฺคมญฺจ อสฺสาทญฺจ อาทีนวญฺจ นิสฺสรณญฺจ พุชฺฌิ, ปญฺจนฺนํ อุปาทานกฺขนฺธานํ สมุทยญฺจ ฯเปฯ นิสฺสรณญฺจ พุชฺฌิ, จตุนฺนํ มหาภูตานํ สมุทยญฺจ อตฺถงฺคมญฺจ อสฺสาทญฺจ อาทีนวญฺจ นิสฺสรณญฺจ พุชฺฌิ, “ยํ กิญฺจิ สมุทยธมฺมํ, สพฺพํ ตํ นิโรธธมฺมนฺ”ติ พุชฺฌิ.}

\prose{อถ วา ยํ พุชฺฌิตพฺพํ อนุพุชฺฌิตพฺพํ ปฏิพุชฺฌิตพฺพํ สมฺพุชฺฌิตพฺพํ อธิคนฺตพฺพํ ผสฺสิตพฺพํ สจฺฉิกาตพฺพํ, สพฺพํ ตํ เตน ปจฺเจกโพธิญาเณน พุชฺฌิ อนุพุชฺฌิ ปฏิพุชฺฌิ สมฺพุชฺฌิ อธิคจฺฉิ ผสฺเสสิ สจฺฉากาสีติ เอวํ โส ปจฺเจกสมฺพุทฺโธ เอโก อนุตฺตรํ ปจฺเจกสมฺโพธิํ อภิสมฺพุทฺโธติ เอโก.}

\prose{\textbf{จเร}ติ อฏฺฐ จริยาโย อิริยาปถจริยา อายตนจริยา สติจริยา สมาธิจริยา ญาณจริยา มคฺคจริยา ปตฺติจริยา โลกตฺถจริยา.\csromanspacer{} \textbf{อิริยาปถจริยา}ติ จตูสุ อิริยาปเถสุ.\csromanspacer{} \textbf{อายตนจริยา}ติ ฉสุ อชฺฌตฺติกพาหิเรสุ อายตเนสุ.\csromanspacer{} \textbf{สติจริยา}ติ จตูสุ สติปฏฺฐาเนสุ.\csromanspacer{} \textbf{สมาธิจริยา}ติ จตูสุ ฌาเนสุ.\csromanspacer{} \textbf{ญาณจริยา}ติ จตูสุ อริยสจฺเจสุ.\csromanspacer{} \textbf{มคฺคจริยา}ติ จตูสุ อริยมคฺเคสุ.\csromanspacer{} \textbf{ปตฺติจริยา}ติ จตูสุ สามญฺญผเลสุ.\csromanspacer{} \textbf{โลกตฺถจริยา}ติ ตถาคเตสุ อรหนฺเตสุ สมฺมาสมฺพุทฺเธสุ ปเทสโต ปจฺเจกสมฺพุทฺเธสุ ปเทสโต สาวเกสุ.\csromanspacer{} \textbf{อิริยาปถจริยา} จ ปณิธิสมฺปนฺนานํ, อายตนจริยา จ อินฺทฺริเยสุ}

\csromanheader{2}{19. ขคฺควิสาณสุตฺตนิทฺเทส}
\prose{\csromanfolio{233}คุตฺตทฺวารานํ, สติจริยา จ อปฺปมาทวิหารีนํ, สมาธิจริยา จ อธิจิตฺตมนุยุตฺตานํ, ญาณจริยา จ พุทฺธิสมฺปนฺนานํ, มคฺคจริยา จ สมฺมาปฏิปนฺนานํ, ปตฺติจริยา จ อธิคตผลานํ, โลกตฺถจริยา จ ตถาคตานํ อรหนฺตานํ สมฺมาสมฺพุทฺธานํ ปเทสโต ปจฺเจกพุทฺธานํ ปเทสโต สาวกานํ, อิมา อฏฺฐ จริยาโย.\csromanspacer{} อปราปิ อฏฺฐ จริยาโย อธิมุจฺจนฺโต สทฺธาย จรติ, ปคฺคณฺหนฺโต วีริเยน จรติ, อุปฏฺฐาเปนฺโต สติยา จรติ, อวิกฺเขปํ กโรนฺโต สมาธินา จรติ, ปชานนฺโต ปญฺญาย จรติ, วิชานนฺโต วิญฺญาณจริยาย จรติ. “เอวํ ปฏิปนฺนสฺส กุสลา ธมฺมา อายาเปนฺตี”ติ อายตนจริยาย จรติ. “เอวํ ปฏิปนฺโน วิเสสมธิคจฺฉตี”ติ วิเสสจริยาย จรติ.\csromanspacer{} อิมา อฏฺฐ จริยาโย.}

\prose{อปราปิ อฏฺฐ จริยาโย ทสฺสนจริยา จ สมฺมาทิฏฺฐิยา, อภินิโรปนจริยา จ สมฺมาสงฺกปฺปสฺส, ปริคฺคหจริยา จ สมฺมาวาจาย, สมุฏฺฐานจริยา จ สมฺมากมฺมนฺตสฺส, โวทานจริยา จ สมฺมา-อาชีวสฺส, ปคฺคหจริยา จ สมฺมาวายามสฺส, อุปฏฺฐานจริยา จ สมฺมาสติยา, อวิกฺเขปจริยา จ สมฺมาสมาธิสฺส, อิมา อฏฺฐ จริยาโย.}

\prose{\textbf{ขคฺควิสาณกปฺโป}ติ ยถา ขคฺคสฺส นาม วิสาณํ เอกํ โหติ อทุติยํ, เอวเมว โส ปจฺเจกสมฺพุทฺโธ ตกฺกปฺโป ตสฺสทิโส ตปฺปฏิภาโค.\csromanspacer{} ยถา อติโลณํ วุจฺจติ โลณกปฺโป, อติติตฺตกํ วุจฺจติ ติตฺตกปฺโป, อติมธุรํ วุจฺจติ มธุรกปฺโป, อติ-อุณฺหํ วุจฺจติ อคฺคิกปฺโป, อติสีตลํ วุจฺจติ หิมกปฺโป, มหา-อุทกกฺขนฺโธ วุจฺจติ สมุทฺทกปฺโป, มหาภิญฺญาพลปฺปตฺโต สาวโก วุจฺจติ สตฺถุกปฺโปติ.\csromanspacer{} เอวเมว โส ปจฺเจกสมฺพุทฺโธ ตตฺถ ตกฺกปฺโป ตสฺสทิโส ตปฺปฏิภาโค เอโก อทุติโย มุตฺตพนฺธโน สมฺมา โลเก จรติ วิหรติ อิริยติ วตฺเตติ ปาเลติ ยเปติ ยาเปตีติ เอโก จเร ขคฺควิสาณกปฺโป.\csromanspacer{} เตนาห โส ปจฺเจกสมฺพุทฺโธ\csromandash{}}

\csromangathameasure{สพฺเพสุ ภูเตสุ นิธาย ทณฺฑํ, \\ อวิเหฐยํ อญฺญตรมฺปิ เตสํ. \\ น ปุตฺตมิจฺเฉยฺย กุโต สหายํ, \\ เอโก จเร ขคฺควิสาณกปฺโปติ.}
\csromangathagroup{\csromangathastanza{สพฺเพสุ ภูเตสุ นิธาย ทณฺฑํ, \\ อวิเหฐยํ อญฺญตรมฺปิ เตสํ. \\ น ปุตฺตมิจฺเฉยฺย กุโต สหายํ, \\ เอโก จเร ขคฺควิสาณกปฺโปติ.}}

\csromanhangingbegin
\hangingheaditem{122}{\csromanfolio{234}\textbf{สํสคฺคชาตสฺส ภวนฺติ สฺเนหา},}
\hangingbody{\textbf{สฺเนหนฺวยํ ทุกฺขมิทํ ปโหติ.}}
\hangingbody{\textbf{อาทีนวํ สฺเนหชํ เปกฺขมาโน,}}
\hangingbody{เอโก จเร ขคฺควิสาณกปฺโป.}
\csromanhangingend

\prose{\textbf{สํสคฺคชาตสฺส ภวนฺติ สฺเนหา}ติ \textbf{สํสคฺคา}ติ เทฺว \textbf{สํสคฺคา} ทสฺสนสํสคฺโค จ สวนสํสคฺโค จ.\csromanspacer{} กตโม ทสฺสนสํสคฺโค, อิเธกจฺโจ ปสฺสติ อิตฺถิํ วา กุมาริํ วา อภิรูปํ ทสฺสนียํ ปาสาทิกํ ปรมาย วณฺณโปกฺขรตาย สมนฺนาคตํ, ทิสฺวา ปสฺสิตฺวา อนุพฺยญฺชนโส นิมิตฺตํ คณฺหาติ เกสา วา โสภนา\footnote{โสภณา (สฺยา)}, มุขํ วา โสภนํ, อกฺขี วา โสภนา, กณฺณา วา โสภนา, นาสา วา โสภนา, โอฏฺฐา วา โสภนา, ทนฺตา วา โสภนา, มุขํ วา โสภนํ, คีวา วา โสภนา, ถนา วา โสภนา, อุรํ วา โสภนํ, อุทรํ วา โสภนํ, กฏิ วา โสภนา, อูรู วา โสภนา, ชงฺฆา วา โสภนา, หตฺถา วา โสภนา, ปาทา วา โสภนา, องฺคุลิโย วา โสภนา, นขา วา โสภนาติ ทิสฺวา ปสฺสิตฺวา อภินนฺทติ อภิวทติ อภิปตฺเถติ อนุปฺปาเทติ\footnote{อนุสฺสรติ (ก)} อนุพนฺธติ ราคพนฺธนํ, อยํ ทสฺสนสํสคฺโค.}

\prose{กตโม สวนสํสคฺโค, อิเธกจฺโจ สุณาติ “อสุกสฺมิํ นาม คาเม วา นิคเม วา อิตฺถี วา กุมารี วา อภิรูปา ทสฺสนียา ปาสาทิกา ปรมาย วณฺณโปกฺขรตาย สมนฺนาคตา”ติ สุตฺวา สุณิตฺวา อภินนฺทติ อภิวทติ อภิปตฺเถติ อนุปฺปาเทติ อนุพนฺธติ ราคพนฺธนํ, อยํ สวนสํสคฺโค.}

\prose{\textbf{สฺเนหา}ติ เทฺว สฺเนหา ตณฺหาสฺเนโห จ ทิฏฺฐิสฺเนโห จ.\csromanspacer{} กตโม ตณฺหาสฺเนโห, ยาวตา ตณฺหาสงฺขาเตน สีมกตํ มริยาทิกตํ โอธิกตํ ปริยนฺตกตํ ปริคฺคหิตํ มมายิตํ, อิทํ มม, เอตํ มม, เอตฺตกํ มม, เอตฺตาวตา มม, มม รูปา สทฺทา คนฺธา รสา โผฏฺฐพฺพา อตฺถรณา ปาวุรณา ทาสิทาสา อเชฬกา กุกฺกุฏสูกรา หตฺถิควาสฺสวฬวา เขตฺตํ วตฺถุ หิรญฺญํ สุวณฺณํ คามนิคมราชธานิโย รฏฺฐญฺจ ชนปโท จ โกโส จ โกฏฺฐาคารญฺจ เกวลมฺปิ มหาปถวิํ ตณฺหาวเสน มมายติ, ยาวตา อฏฺฐสตํ ตณฺหาวิจริตํ, อยํ ตณฺหาสฺเนโห.}

\csromanheader{2}{19. ขคฺควิสาณสุตฺตนิทฺเทส}
\prose{\csromanfolio{235}กตโม ทิฏฺฐิ\textbf{สฺเนโห}, วีสติวตฺถุกา สกฺกายทิฏฺฐิ, ทสวตฺถุกา มิจฺฉาทิฏฺฐิ, ทสวตฺถุกา อนฺตคฺคาหิกา ทิฏฺฐิ.\csromanspacer{} ยา เอวรูปา ทิฏฺฐิ ทิฏฺฐิคตํ ทิฏฺฐิคหนํ ทิฏฺฐิกนฺตาโร ทิฏฺฐิวิสูกายิกํ ทิฏฺฐิวิปฺผนฺทิตํ ทิฏฺฐิสํโยชนํ คาโห ปฏิคฺคาโห\footnote{ปติฏฺฐาโห (ก)} อภินิเวโส ปรามาโส กุมฺมคฺโค มิจฺฉาปโถ มิจฺฉตฺตํ ติตฺถายตนํ วิปริยาสคฺคาโห วิปรีตคฺคาโห วิปลฺลาสคฺคาโห มิจฺฉาคาโห อยาถาวกสฺมิํ “ยาถาวกนฺ”ติ คาโห ยาวตา ทฺวาสฏฺฐิ ทิฏฺฐิคตานิ.\csromanspacer{} อยํ ทิฏฺฐิ\textbf{สฺเนโห}.}

\prose{\textbf{สํสคฺคชาตสฺส ภวนฺติ สฺเนหา}ติ ทสฺสนสํสคฺคปจฺจยา จ สวนสํสคฺคปจฺจยา จ ตณฺหา\textbf{สฺเนโห} จ ทิฏฺฐิ\textbf{สฺเนโห} จ ภวนฺติ สมฺภวนฺติ ชายนฺติ สญฺชายนฺติ นิพฺพตฺตนฺติ อภินิพฺพตฺตนฺติ ปาตุภวนฺตีติ สํสคฺคชาตสฺส ภวนฺติ สฺเนหา.}

\prose{\textbf{สฺเนหนฺวยํ ทุกฺขมิทํ ปโหตี}ติ \textbf{สฺเนโห}ติ เทฺว สฺเนหา ตณฺหาสฺเนโย จ ทิฏฺฐิ\textbf{สฺเนโห} จ ฯเปฯ อยํ ตณฺหา\textbf{สฺเนโห} ฯเปฯ อยํ ทิฏฺฐิ\textbf{สฺเนโห}.\csromanspacer{} \textbf{ทุกฺขมิทํ ปโหตี}ติ อิเธกจฺโจ กาเยน ทุจฺจริตํ จรติ, วาจาย ทุจฺจริตํ จรติ, มนสา ทุจฺจริตํ จรติ, ปาณมฺปิ หนติ, อทินฺนมฺปิ อาทิยติ, สนฺธิมฺปิ ฉินฺทติ, นิลฺโลปมฺปิ\footnote{วิโลปมฺปิ (สฺยา) ปสฺส ขุ 7. 315 ปิฏฺเฐ.} หรติ, เอกาคาริกมฺปิ กโรติ, ปริปนฺเถปิ ติฏฺฐติ, ปรทารมฺปิ คจฺฉติ, มุสาปิ ภณติ, ตเมนํ คเหตฺวา รญฺโญ ทสฺเสนฺติ “อยํ เทว โจโร อาคุจารี, อิมสฺส ยํ อิจฺฉสิ, ตํ ทณฺฑํ ปเณหี”ติ.\csromanspacer{} ตเมนํ ราชา ตํ ปริภาสติ, โส ปริภาสปจฺจยาปิ ทุกฺขํ โทมนสฺสํ\footnote{ทุกฺขโทมนสฺสํ (สฺยา)} ปฏิสํเวเทติ.\csromanspacer{} เอตํ ภยํ ทุกฺขํ โทมนสฺสํ กุโต ชาตํ, ตสฺส สฺเนหปจฺจยา จ นนฺทิปจฺจยา จ ราคปจฺจยา จ นนฺทิราคปจฺจยา จ ชาตํ.}

\prose{เอตฺตเกนปิ ราชา น ตุสฺสติ, ตเมนํ ราชา พนฺธาเปติ อนฺทุพนฺธเนน วา รชฺชุพนฺธเนน วา สงฺขาลิกพนฺธเนน วา เวตฺตพนฺธเนน วา ลตาพนฺธเนน วา ปกฺเขปพนฺธเนน วา ปริกฺเขปพนฺธเนน วา คามพนฺธเนน วา นิคมพนฺธเนน วา รฏฺฐพนฺธเนน วา ชนปทพนฺธเนน วา อนฺตมโส สวจนียมฺปิ กโรติ “น เต ลพฺภา อิโต ปกฺกมิตุนฺ”ติ, โส พนฺธนปจฺจยาปิ ทุกฺขํ โทมนสฺสํ ปฏิสํเวเทติ.\csromanspacer{} เอตํ ภยํ ทุกฺขํ โทมนสฺสํ กุโต ชาตํ, ตสฺส สฺเนหปจฺจยา จ นนฺทิปจฺจยา จ ราคปจฺจยา จ นนฺทิราคปจฺจยา จ ชาตํ.}

\prose{\csromanfolio{236}เอตฺตเกนปิ ราชา น ตุสฺสติ, ตเมนํ ราชา ตสฺเสว\footnote{ตสฺส (สฺยา)} ธนํ อาหราเปติ สตํ วา สหสฺสํ วา สตสหสฺสํ วา, โส ธนชานิปจฺจยาปิ ทุกฺขํ โทมนสฺสํ ปฏิสํเวเทติ.\csromanspacer{} เอตํ ภยํ ทุกฺขํ โทมนสฺสํ กุโตชาตํ, ตสฺส สฺเนหปจฺจยา จ นนฺทิปจฺจยา จ ราคปจฺจยา จ นนฺทิราคปจฺจยา จ ชาตํ.}

\prose{เอตฺตเกนปิ ราชา น ตุสฺสติ, ตเมนํ ราชา วิวิธา กมฺมการณา\footnote{วิวิธานิ กมฺมกรณานิ (ก)} การาเปติ, กสาหิปิ ตาเฬติ, เวตฺเตหิปิ ตาเฬติ, อฑฺฒทณฺเฑหิปิ ตาเฬติ, หตฺถมฺปิ ฉินฺทติ, ปาทมฺปิ ฉินฺทติ, หตฺถปาทมฺปิ ฉินฺทติ, กณฺณมฺปิ ฉินฺทติ, นาสมฺปิ ฉินฺทติ, กณฺณนาสมฺปิ ฉินฺทติ, พิลงฺคถาลิกมฺปิ กโรติ, สงฺขมุณฺฑิกมฺปิ กโรติ, ราหุมุขมฺปิ กโรติ, โชติมาลิกมฺปิ กโรติ, หตฺถปชฺโชติกมฺปิ กโรติ, เอรกวตฺติกมฺปิ\footnote{เอรกวฏฺฏิกมฺปิ (สฺยา, ก) ปสฺส ม 1. 122 ปิฏฺเฐ.} กโรติ, จีรกวาสิกมฺปิ กโรติ, เอเณยฺยกมฺปิ กโรติ, พฬิสมํสิกมฺปิ กโรติ, กหาปณิกมฺปิ กโรติ, ขาราปตจฺฉิกมฺปิ\footnote{ขาราปฏิจฺฉิกมฺปิ (ก) เหฏฺฐา 190 ปิฏฺเฐ.} กโรติ, ปลิฆปริวตฺติกมฺปิ กโรติ, ปลาลปีฐกมฺปิ กโรติ, ตตฺเตนปิ เตเลน โอสิญฺจติ, สุนเขหิปิ ขาทาเปติ, ชีวนฺตมฺปิ สูเล อุตฺตาเสติ, อสินาปิ สีสํ ฉินฺทติ, โส กมฺมการณปจฺจยาปิ ทุกฺขํ โทมนสฺสํ ปฏิสํเวเทติ.\csromanspacer{} เอตํ ภยํ ทุกฺขํ โทมนสฺสํ กุโต ชาตํ, ตสฺส สฺเนหปจฺจยา จ นนฺทิปจฺจยา จ ราคปจฺจยา จ นนฺทิราคปจฺจยา จ ชาตํ.\csromanspacer{} ราชา อิเมสํ จตุนฺนํ ทณฺฑานํ อิสฺสโร.}

\prose{โส สเกน กมฺเมน กายสฺส เภทา ปรํ มรณา อปายํ ทุคฺคติํ วินิปาตํ นิรยํ อุปปชฺชติ.\csromanspacer{} ตเมนํ นิรยปาลา ปญฺจวิธพนฺธนํ นาม กมฺมการณํ กโรนฺติ\footnote{กาเรนฺติ (สฺยา, ก) ปสฺส ม 3. 221 ปิฏฺเฐ.}, ตตฺตํ อโยขิลํ หตฺเถ คเมนฺติ, ตตฺตํ อโยขิลํ ทุติเย หตฺเถ คเมนฺติ, ตตฺตํ อโยขิลํ ปาเท คเมนฺติ, ตตฺตํ อโยขิลํ ทุติเย ปาเท คเมนฺติ, ตตฺตํ อโยขิลํ มชฺเฌ อุรสฺมิํ คเมนฺติ.\csromanspacer{} โส ตตฺถ ทุกฺขา ติพฺพา\footnote{ติปฺปา (สฺยา)} ขรา กฏุกา เวทนา เวเทติ, น จ ตาว กาลํ กโรติ, ยาว น ตํ ปาปกมฺมํ พฺยนฺตีโหติ.\csromanspacer{} เอตํ ภยํ ทุกฺขํ โทมนสฺสํ กุโต ชาตํ, ตสฺส สฺเนหปจฺจยา จ นนฺทิปจฺจยา จ ราคปจฺจยา จ นนฺทิราคปจฺจยา จ ชาตํ.}

\csromanheader{2}{19. ขคฺควิสาณสุตฺตนิทฺเทส}
\prose{\csromanfolio{237}ตเมนํ นิรยปาลา สํเวเสตฺวา\footnote{สํเวสิตฺวา (สฺยา, ก) ปสฺส ม 3. 221 ปิฏฺเฐ.} กุฐารีหิ\footnote{กุธารีหิ (สฺยา, ก)} ตจฺฉนฺติ ฯเปฯ.\csromanspacer{} ตเมนํ นิรยปาลา อุทฺธํปาทํ อโธสิรํ คเหตฺวา วาสีหิ ตจฺฉนฺติ.\csromanspacer{} ตเมนํ นิรยปาลา รเถ โยเชตฺวา อาทิตฺตาย ปถวิยา สมฺปชฺชลิตาย สโชติภูตาย\footnote{สญฺโชติภูตาย (สฺยา)} สาเรนฺติปิ ปจฺจาสาเรนฺติปิ.\csromanspacer{} ตเมนํ นิรยปาลา มหนฺตํ องฺคารปพฺพตํ อาทิตฺตํ สมฺปชฺชลิตํ สโชติภูตํ อาโรเปนฺติปิ โอโรเปนฺติปิ.\csromanspacer{} ตเมนํ นิรยปาลา อุทฺธํปาทํ อโธสิรํ คเหตฺวา ตตฺตาย โลหกุมฺภิยา ปกฺขิปนฺติ อาทิตฺตาย สมฺปชฺชลิตาย สโชติภูตาย, โส ตตฺถ เผณุทฺเทหกํ ปจฺจติ, โส ตตฺถ เผณุทฺเทหกํ ปจฺจมาโน สกิมฺปิ อุทฺธํ คจฺฉติ, สกิมฺปิ อโธ คจฺฉติ, สกิมฺปิ ติริยํ คจฺฉติ.\csromanspacer{} โส ตตฺถ ทุกฺขา ติพฺพา ขรา กฏุกา เวทนา เวเทติ, น จ ตาว กาลํ กโรติ, ยาว น ตํ ปาปกมฺมํ พฺยนฺตีโหติ.\csromanspacer{} เอตํ ภยํ ทุกฺขํ โทมนสฺสํ กุโต ชาตํ, ตสฺส สฺเนหปจฺจยา จ นนฺทิปจฺจยา จ ราคปจฺจยา จ นนฺทิราคปจฺจยา จ ชาตํ.}

\prose{ตเมนํ นิรยปาลา มหานิรเย ปกฺขิปนฺติ, โส โข ปน มหานิรโย\csromandash{}}

\csromangathameasure{จตุกฺกณฺโณ จตุทฺวาโร, วิภตฺโต ภาคโส มิโต. \\ อโยปาการปริยนฺโต, อยสา ปฏิกุชฺชิโต. \\ ตสฺส อโยมยา ภูมิ, ชลิตา เตชสา ยุตา. \\ สมนฺตา โยชนสตํ, ผริตฺวา ติฏฺฐติ สพฺพทา. \\ กทริยาตปนา โฆรา, อจฺจิมนฺโต ทุราสทา. \\ โลมหํสนรูปา จ, เภสฺมา ปฏิภยา ทุขา. \\ ปุรตฺถิมาย จ ภิตฺติยา, อจฺจิกฺขนฺโธ สมุฏฺฐิโต. \\ qอหนฺโต ปาปกมฺมนฺเต, ปจฺฉิมาย ปฏิหญฺญติ. \\ ปจฺฉิมาย จ ภิตฺติยา, อจฺจิกฺขนฺโธ สมุฏฺฐิโต. \\ qอหนฺโต ปาปกมฺมนฺเต, ปุริมาย ปฏิหญฺญติ. \\ ทกฺขิณาย จ ภิตฺติยา, อจฺจิกฺขนฺโธ สมุฏฺฐิโต. \\ qอหนฺโต ปาปกมฺมนฺเต, อุตฺตราย ปฏิหญฺญติ. \\ อุตฺตราย จ ภิตฺติยา, อจฺจิกฺขนฺโธ สมุฏฺฐิโต. \\ qอหนฺโต ปาปกมฺมนฺเต, ทกฺขิณาย ปฏิหญฺญติ. \\ เหฏฺฐโต จ สมุฏฺฐาย, อจฺจิกฺขนฺโธ ภยานโก. \\ qอหนฺโต ปาปกมฺมนฺเต, ฉทนสฺมิํ ปฏิหญฺญติ. \\ ฉทนมฺหา สมุฏฺฐาย, อจฺจิกฺขนฺโธ ภยานโก. \\ qอหนฺโต ปาปกมฺมนฺเต, ภูมิยํ ปฏิหญฺญติ. \\ อโยกปาลมาทิตฺตํ, สนฺตตฺตํ ชลิตํ ยถา. \\ เอวํ อวีจินิรโย, เหฏฺฐา อุปริ ปสฺสโต. \\ ตตฺถ สตฺตา มหาลุทฺทา, มหากิพฺพิสการิโน. \\ อจฺจนฺตปาปกมฺมนฺตา, ปจฺจนฺติ น จ มิยฺยเร. \\ ชาตเวทสโม กาโย, เตสํ นิรยวาสินํ. \\ ปสฺส กมฺมานํ ทฬฺหตฺตํ, น ภสฺมา โหติ นปิ มสิ. \\ ปุรตฺถิเมนปิ ธาวนฺติ, ตโต ธาวนฺติ ปจฺฉิมํ. \\ อุตฺตเรนปิ ธาวนฺติ, ตโต ธาวนฺติ ทกฺขิณํ. \\ ยํ ยํ ทิสํ ปธาวนฺติ, ตํ ตํ ทฺวารํ ปิธียติ. \\ อภินิกฺขมิตาสา เต, สตฺตา โมกฺขคเวสิโน. \\ น เต ตโต นิกฺขมิตุํ, ลภนฺติ กมฺมปจฺจยา. \\ เตสญฺจ ปาปกมฺมนฺตํ, อวิปกฺกํ กตํ พหุนฺติ.}
\csromangathagroup{\csromangathastanza{จตุกฺกณฺโณ จตุทฺวาโร, วิภตฺโต ภาคโส มิโต. \\ อโยปาการปริยนฺโต, อยสา ปฏิกุชฺชิโต.}\csromangathastanza{ตสฺส อโยมยา ภูมิ, ชลิตา เตชสา ยุตา. \\ สมนฺตา โยชนสตํ, ผริตฺวา ติฏฺฐติ สพฺพทา.}\csromangathastanza{กทริยาตปนา\footnote{กทริยา ตาปนา (ก) ขุ 7. 317 ปิฏฺเฐปิ.} โฆรา, อจฺจิมนฺโต ทุราสทา. \\ โลมหํสนรูปา จ, เภสฺมา ปฏิภยา ทุขา.}\csromangathastanza{ปุรตฺถิมาย จ ภิตฺติยา, อจฺจิกฺขนฺโธ สมุฏฺฐิโต. \\ qอหนฺโต ปาปกมฺมนฺเต, ปจฺฉิมาย ปฏิหญฺญติ.}\csromangathastanza{ปจฺฉิมาย จ ภิตฺติยา, อจฺจิกฺขนฺโธ สมุฏฺฐิโต. \\ qอหนฺโต ปาปกมฺมนฺเต, ปุริมาย ปฏิหญฺญติ.}\csromangathastanza{ทกฺขิณาย จ ภิตฺติยา, อจฺจิกฺขนฺโธ สมุฏฺฐิโต. \\ qอหนฺโต ปาปกมฺมนฺเต, อุตฺตราย ปฏิหญฺญติ.}\csromangathastanza{\csromanfolio{238}อุตฺตราย จ ภิตฺติยา, อจฺจิกฺขนฺโธ สมุฏฺฐิโต. \\ qอหนฺโต ปาปกมฺมนฺเต, ทกฺขิณาย ปฏิหญฺญติ.}\csromangathastanza{เหฏฺฐโต จ สมุฏฺฐาย, อจฺจิกฺขนฺโธ ภยานโก. \\ qอหนฺโต ปาปกมฺมนฺเต, ฉทนสฺมิํ ปฏิหญฺญติ.}\csromangathastanza{ฉทนมฺหา สมุฏฺฐาย, อจฺจิกฺขนฺโธ ภยานโก. \\ qอหนฺโต ปาปกมฺมนฺเต, ภูมิยํ ปฏิหญฺญติ.}\csromangathastanza{อโยกปาลมาทิตฺตํ, สนฺตตฺตํ ชลิตํ ยถา. \\ เอวํ อวีจินิรโย, เหฏฺฐา อุปริ ปสฺสโต.}\csromangathastanza{ตตฺถ สตฺตา มหาลุทฺทา, มหากิพฺพิสการิโน. \\ อจฺจนฺตปาปกมฺมนฺตา, ปจฺจนฺติ น จ มิยฺยเร\footnote{มียเร (ก)}.}\csromangathastanza{ชาตเวทสโม กาโย, เตสํ นิรยวาสินํ. \\ ปสฺส กมฺมานํ ทฬฺหตฺตํ, น ภสฺมา โหติ นปิ มสิ.}\csromangathastanza{ปุรตฺถิเมนปิ ธาวนฺติ, ตโต ธาวนฺติ ปจฺฉิมํ. \\ อุตฺตเรนปิ ธาวนฺติ, ตโต ธาวนฺติ ทกฺขิณํ.}\csromangathastanza{ยํ ยํ ทิสํ ปธาวนฺติ, ตํ ตํ ทฺวารํ ปิธียติ. \\ อภินิกฺขมิตาสา เต, สตฺตา โมกฺขคเวสิโน.}\csromangathastanza{น เต ตโต นิกฺขมิตุํ, ลภนฺติ กมฺมปจฺจยา. \\ เตสญฺจ ปาปกมฺมนฺตํ, อวิปกฺกํ กตํ พหุนฺติ.}}

\prose{เอตํ ภยํ ทุกฺขํ โทมนสฺสํ กุโต ชาตํ, ตสฺส สฺเนหปจฺจยา จ นนฺทิปจฺจยา จ ราคปจฺจยา จ นนฺทิราคปจฺจยา จ ชาตํ.}

\prose{ยานิ จ เนรยิกานิ ทุกฺขานิ ยานิ จ ติรจฺฉานโยนิกานิ ทุกฺขานิ ยานิ จ เปตฺติวิสยิกานิ ทุกฺขานิ ยานิ จ มานุสิกานิ ทุกฺขานิ, ตานิ กุโต ชาตานิ กุโต สญฺชาตานิ กุโต นิพฺพตฺตานิ กุโต อภินิพฺพตฺตานิ กุโต ปาตุภูตานิ.\csromanspacer{} ตสฺส สฺเนหปจฺจยา จ นนฺทิปจฺจยา จ ราคปจฺจยา จ นนฺทิราคปจฺจยา จ ภวนฺติ สมฺภวนฺติ ชายนฺติ สญฺชายนฺติ นิพฺพตฺตนฺติ อภินิพฺพตฺตนฺติ ปาตุภวนฺตีติ สฺเนหนฺวยํ ทุกฺขมิทํ ปโหติ.}

\csromanheader{2}{19. ขคฺควิสาณสุตฺตนิทฺเทส}
\prose{\csromanfolio{239}\textbf{อาทีนวํ สฺเนหชํ เปกฺขมาโน}ติ \textbf{สฺเนโห}ติ เทฺว สฺเนหา ตณฺหา\textbf{สฺเนโห} จ ทิฏฺฐิ\textbf{สฺเนโห} จ ฯเปฯ อยํ ตณฺหา\textbf{สฺเนโห} ฯเปฯ อยํ ทิฏฺฐิ\textbf{สฺเนโห.}\csromanspacer{} \textbf{อาทีนวํ สฺเนหชํ เปกฺขมาโน}ติ ตณฺหา\textbf{สฺเนโห} จ ทิฏฺฐิ\textbf{สฺเนโห} จ อาทีนวํ สฺเนหชํ เปกฺขมาโน ทกฺขมาโน โอโลกยมาโน นิชฺฌายมาโน อุปปริกฺขมาโนติ อาทีนวํ สฺเนหชํ \textbf{เปกฺขมาโน,} เอโก จเร ขคฺควิสาณกปฺโป\textbf{.}\csromanspacer{} เตนาห โส ปจฺเจกสมฺพุทฺโธ\csromandash{}}

\csromangathameasure{สํสคฺคชาตสฺส ภวนฺติ สฺเนหา, \\ สฺเนหนฺวยํ ทุกฺขมิทํ ปโหติ\textbf{.}}
\csromangathagroup{\csromangathastanza{สํสคฺคชาตสฺส ภวนฺติ สฺเนหา, \\ สฺเนหนฺวยํ ทุกฺขมิทํ ปโหติ\textbf{.}}}

\prose{\textbf{อาทีนวํ สฺเนหชํ เปกฺขมาโน},}

\prose{เอโก จเร ขคฺควิสาณกปฺโปติ\textbf{.}}

\csromanhangingbegin
\hangingheaditem{123}{มิตฺเต สุหชฺเช อนุกมฺปมาโน,}
\hangingbody{\textbf{หาเปติ อตฺถํ ปฏิพทฺธจิตฺโต1.}}
\hangingbody{\textbf{เอตํ ภยํ สนฺถเว2} \textbf{เปกฺขมาโน,}}
\hangingbody{เอโก จเร ขคฺควิสาณกปฺโป\textbf{.}}
\csromanhangingend

\prose{\textbf{มิตฺเต สุหชฺเช อนุกมฺปมาโน, หาเปติ อตฺถํ ปฏิพทฺธจิตฺโต}ติ \textbf{มิตฺตา}ติ เทฺว \textbf{มิตฺตา} อคาริกมิตฺโต จ อนาคาริกมิตฺโต\footnote{ปพฺพชิตมิตฺโต (ก) เอวมุปริปิ.} จ\textbf{.}\csromanspacer{} กตโม อคาริกมิตฺโต, อิเธกจฺโจ ทุทฺททํ ททาติ, ทุจฺจชํ จชติ, ทุกฺกรํ กโรติ, ทุกฺขมํ ขมติ, คุยฺหมสฺส อาจิกฺขติ, คุยฺหมสฺส ปริคูหติ\footnote{ปริคุยฺหติ (สฺยา, ก)}, อาปทาสุ น วิชหติ, ชีวิตมฺปิสฺส อตฺถาย ปริจฺจตฺตํ โหติ, ขีเณ นาติมญฺญติ\textbf{.}\csromanspacer{} อยํ อคาริกมิตฺโต\textbf{.}}

\prose{กตโม อนาคาริกมิตฺโต, อิธ ภิกฺขุ ปิโย จ โหติ มนาโป จ ครุ จ ภาวนีโย จ วตฺตา จ วจนกฺขโม จ คมฺภีรญฺจ กถํ กตฺตา โน จ อฏฺฐาเน นิโยเชติ, อธิสีเล สมาทเปติ, จตุนฺนํ สติปฏฺฐานานํ ภาวนานุโยเค สมาทเปติ, จตุนฺนํ สมฺมปฺปธานานํ ฯเปฯ จตุนฺนํ อิทฺธิปาทานํ\textbf{.}\csromanspacer{} ปญฺจนฺนํ อินฺทฺริยานํ\textbf{.}\csromanspacer{} ปญฺจนฺนํ พลานํ\textbf{.}\csromanspacer{} สตฺตนฺนํ โพชฺฌงฺคานํ\textbf{.}\csromanspacer{} อริยสฺส อฏฺฐงฺคิกสฺส มคฺคสฺส ภาวนานุโยเค สมาทเปติ\textbf{.}\csromanspacer{} อยํ อนาคาริกมิตฺโต\textbf{.}}

\prose{\csromanfolio{240}สุหชฺชา วุจฺจนฺติ เยหิ สห อาคมนํ ผาสุ, คมนํ ผาสุ, ฐานํ ผาสุ, นิสชฺชนํ\footnote{นิสชฺชา (ก)} ผาสุ, สยนํ ผาสุ, อาลปนํ ผาสุ, สลฺลปนํ ผาสุ, อุลฺลปนํ ผาสุ, สมุลฺลปนํ ผาสุ.\csromanspacer{} \textbf{มิตฺเต สุหชฺเช อนุกมฺปมาโน หาเปติ อตฺถนฺ}ติ มิตฺเต จ สุหชฺเช จ สนฺทิฏฺเฐ จ สมฺภตฺเต จ สหาเย จ อนุกมฺปมาโน อนุเปกฺขมาโน อนุคยฺหมาโน อตฺตตฺถมฺปิ ปรตฺถมฺปิ อุภยตฺถมฺปิ หาเปติ, ทิฏฺฐธมฺมิกมฺปิ อตฺถํ หาเปติ, สมฺปรายิกมฺปิ อตฺถํ หาเปติ, ปรมตฺถมฺปิ หาเปติ ปหาเปติ ปริหาเปติ ปริธํเสติ ปริวชฺเชติ\footnote{ปริสชฺเชติ (สฺยา)} อนฺตรธาเปตีติ มิตฺเต สุหชฺเช อนุกมฺปมาโน หาเปติ อตฺถํ.}

\prose{\textbf{ปฏิพทฺธจิตฺโต}ติ ทฺวีหิ การเณหิ ปฏิพทฺธจิตฺโต โหติ อตฺตานํ วา นีจํ ฐเปนฺโต ปรํ อุจฺจํ ฐเปนฺโต ปฏิพทฺธจิตฺโต โหติ, อตฺตานํ วา อุจฺจํ ฐเปนฺโต ปรํ นีจํ ฐเปนฺโต ปฏิพทฺธจิตฺโต โหติ.\csromanspacer{} กถํ อตฺตานํ นีจํ ฐเปนฺโต ปรํ อุจฺจํ ฐเปนฺโต ปฏิพทฺธจิตฺโต โหติ, ตุเมฺห เม พหูปการา, อหํ ตุเมฺห นิสฺสาย ลภามิ จีวร ปิณฺฑปาต เสนาสนคิลาน ปจฺจย เภสชฺช ปริกฺขารํ, ยมฺปิ\footnote{เยปิ (ก) เอวํ 293 ปิฏฺเฐปิ.} เม อญฺเญ ทาตุํ วา กาตุํ วา มญฺญนฺติ ตุเมฺห นิสฺสาย ตุเมฺห สมฺปสฺสนฺตา, ยมฺปิ เม โปราณํ มาตาเปตฺติกํ นามโคตฺตํ, ตมฺปิ เม อนฺตรหิตํ, ตุเมฺหหิ อหํ ญายามิ “อสุกสฺส กุลุปโก อสุกาย กุลุปโก”ติ.\csromanspacer{} เอวํ อตฺตานํ นีจํ ฐเปนฺโต ปรํ อุจฺจํ ฐเปนฺโต ปฏิพทฺธจิตฺโต โหติ.}

\prose{กถํ อตฺตานํ อุจฺจํ ฐเปนฺโต ปรํ นีจํ ฐเปนฺโต ปฏิพทฺธจิตฺโต โหติ, อหํ ตุมฺหากํ พหูปกาโร, ตุเมฺห มํ อาคมฺม พุทฺธํ สรณํ คตา, ธมฺมํ สรณํ คตา, สํฆํ สรณํ คตา, ปาณาติปาตา ปฏิวิรตา, อทินฺนาทานา ปฏิวิรตา, กาเมสุมิจฺฉาจารา ปฏิวิรตา, มุสาวาทา ปฏิวิรตา, สุราเมรยมชฺชปมาทฏฺฐานา ปฏิวิรตา, อหํ ตุมฺหากํ อุทฺเทสํ เทมิ, ปริปุจฺฉํ เทมิ, อุโปสถํ อาจิกฺขามิ, นวกมฺมํ อธิฏฺฐามิ, อถ ปน ตุเมฺห มํ อุชฺฌิตฺวา\footnote{ปริจฺจชิตฺวา (สฺยา)} อญฺเญ สกฺกโรถ ครุํ กโรถ มาเนถ ปูเชถาติ เอวํ อตฺตานํ อุจฺจํ ฐเปนฺโต ปรํ นีจํ ฐเปนฺโต ปฏิพทฺธจิตฺโต โหตีติ มิตฺเต สุหชฺเช อนุกมฺปมาโน, หาเปติ อตฺถํ ปฏิพทฺธจิตฺโต.}

\prose{\textbf{เอตํ ภยํ สนฺถเว เปกฺขมาโน}ติ \textbf{ภยนฺ}ติ ชาติภยํ ชราภยํ พฺยาธิภยํ มรณภยํ ราชภยํ โจรภยํ อคฺคิภยํ อุทกภยํ}

\csromanheader{2}{19. ขคฺควิสาณสุตฺตนิทฺเทส}
\prose{\csromanfolio{241}อตฺตานุวาทภยํ ปรานุวาทภยํ ทณฺฑภยํ ทุคฺคติภยํ อูมิภยํ กุมฺภิลภยํ อาวฏฺฏภยํ สุสุมารภยํ\footnote{สุํสุมารภยํ (สฺยา)} อาชีวิกภยํ อสิโลกภยํ ปริสสารชฺชภยํ มทนภยํ ภยานกํ ฉมฺภิตตฺตํ โลมหํโส เจตโส อุพฺเพโค อุตฺราโส.\csromanspacer{} \textbf{สนฺถเว}ติ เทฺว สนฺถวา ตณฺหาสนฺถโว จ ทิฏฺฐิสนฺถโว จ ฯเปฯ อยํ ตณฺหาสนฺถโว ฯเปฯ อยํ ทิฏฺฐิสนฺถโว.\csromanspacer{} \textbf{เอตํ ภยํ สนฺถเว เปกฺขมาโน}ติ เอตํ ภยํ สนฺถเว เปกฺขมาโน ทกฺขมาโน โอโลกยมาโน นิชฺฌายมาโน อุปปริกฺขมาโนติ เอตํ ภยํ สนฺถเว เปกฺขมาโน, เอโก จเร ขคฺควิสาณกปฺโป.\csromanspacer{} เตนาห โส ปจฺเจกสมฺพุทฺโธ\csromandash{}}

\csromangathameasure{มิตฺเต สุหชฺเช อนุกมฺปมาโน, \\ หาเปติ อตฺถํ ปฏิพทฺธจิตฺโต. \\ \textbf{เอตํ ภยํ สนฺถเว เปกฺขมาโน}, \\ เอโก จเร ขคฺควิสาณกปฺโปติ.}
\csromangathagroup{\csromangathastanza{มิตฺเต สุหชฺเช อนุกมฺปมาโน, \\ หาเปติ อตฺถํ ปฏิพทฺธจิตฺโต. \\ \textbf{เอตํ ภยํ สนฺถเว เปกฺขมาโน}, \\ เอโก จเร ขคฺควิสาณกปฺโปติ.}}

\csromanhangingbegin
\hangingheaditem{124}{\textbf{วํโส วิสาโลว ยถา วิสตฺโต},}
\hangingbody{\textbf{ปุตฺเตสุ ทาเรสุ จ ยา อเปกฺขา.}}
\hangingbody{\textbf{วํสกฺกฬีโรว2} \textbf{อสชฺชมาโน,}}
\hangingbody{เอโก จเร ขคฺควิสาณกปฺโป.}
\csromanhangingend

\prose{\textbf{วํโส วิสาโลว ยถา วิสตฺโต}ติ วํโส วุจฺจติ เวฬุคุมฺโพ.\csromanspacer{} ยถา เวฬุคุมฺพสฺมิํ โปราณกา วํสา สตฺตา\footnote{เวฬุคุมฺพสฺมิํ กณฺฏกา ชฏิตา สํสิพฺพิตา (สฺยา)} วิสตฺตา อาสตฺตา ลคฺคา ลคฺคิตา ปลิพุทฺธา, เอวเมว วิสตฺติกา วุจฺจติ ตณฺหา.\csromanspacer{} โย ราโค สาราโค อนุนโย อนุโรโธ นนฺที นนฺทิราโค จิตฺตสฺส สาราโค อิจฺฉา มุจฺฉา อชฺโฌสานํ เคโธ ปลิเคโธ สงฺโค ปงฺโก เอชา มายา ชนิกา สญฺชนนี สิพฺพินี ชาลินี สริตา วิสตฺติกา สุตฺตํ วิสฏา อายูหนี ทุติยา ปณิธิ ภวเนตฺติ วนํ วนโถ สนฺถโว สิเนโห อเปกฺขา ปฏิพนฺธุ อาสา อาสีสนา อาสีสิตตฺตํ รูปาสา สทฺทาสา คนฺธาสา รสาสา โผฏฺฐพฺพาสา ลาภาสา ธนาสา ปุตฺตาสา ชีวิตาสา ชปฺปา ปชปฺปา อภิชปฺปา ชปฺปนา ชปฺปิตตฺตํ โลลุปฺปํ โลลุปฺปายนา โลลุปฺปายิตตฺตํ ปุจฺฉญฺชิกตา สาธุกมฺยตา อธมฺมราโค วิสมโลโภ นิกนฺติ นิกามนา ปตฺถนา ปิหนา สมฺปตฺถนา กามตณฺหา \csromanfolio{242}ภวตณฺหา วิภวตณฺหา รูปตณฺหา อรูปตณฺหา นิโรธตณฺหา รูปตณฺหา สทฺทตณฺหา คนฺธตณฺหา รสตณฺหา โผฏฺฐพฺพตณฺหา ธมฺมตณฺหา โอโฆ โยโค คนฺโถ อุปาทานํ อาวรณํ นีวรณํ ฉทนํ พนฺธนํ อุปกฺกิเลโส อนุสโย ปริยุฏฺฐานํ ลตา เววิจฺฉํ ทุกฺขมูลํ ทุกฺขนิทานํ ทุกฺขปฺปภโว มารปาโส มารพฬิสํ มารวิสโย มารนิวาโส มารพนฺธนํ ตณฺหานที ตณฺหาชาลํ ตณฺหาคทฺทุลํ ตณฺหาสมุทฺโท อภิชฺฌา โลโภ อกุสลมูลํ.}

\prose{\textbf{วิสตฺติกา}ติ เกนฏฺเฐน วิสตฺติกา, วิสาลาติ วิสตฺติกา, วิสตาติ วิสตฺติกา, วิสฏาติ วิสตฺติกา, วิสมาติ วิสตฺติกา, วิสกฺกตีติ วิสตฺติกา, วิสํหรตีติ วิสตฺติกา, วิสํวาทิกาติ วิสตฺติกา, วิสมูลาติ วิสตฺติกา, วิสผลาติ วิสตฺติกา, วิสปริโภคาติ วิสตฺติกา, วิสาลา วา ปน ตณฺหา รูเป สทฺเท คนฺเธ รเส โผฏฺฐพฺเพ กุเล คเณ อาวาเส ลาเภ ยเส ปสํสาย สุเข จีวเร ปิณฺฑปาเต เสนาสเน คิลานปจฺจยเภสชฺชปริกฺขาเร กามธาตุยา รูปธาตุยา อรูปธาตุยา กามภเว รูปภเว อรูปภเว สญฺญาภเว อสญฺญาภเว เนวสญฺญานาสญฺญาภเว เอกโวการภเว จตุโวการภเว ปญฺจโวการภเว อตีเต อนาคเต ปจฺจุปฺปนฺเน ทิฏฺฐสุตมุตวิญฺญาตพฺเพสุ ธมฺเมสุ วิสตา วิตฺถตาติ วิสตฺติกาติ วํโส วิสาโลว ยถา วิสตฺโต.}

\prose{\textbf{ปุตฺเตสุ ทาเรสุ จ ยา อเปกฺขา}ติ \textbf{ปุตฺตา}ติ จตฺตาโร \textbf{ปุตฺตา} อตฺรโช ปุตฺโต เขตฺตโช ปุตฺโต ทินฺนโก ปุตฺโต อนฺเตวาสิโก ปุตฺโต.\csromanspacer{} \textbf{ทารา} วุจฺจนฺติ ภริยาโย.\csromanspacer{} \textbf{อเปกฺขา} วุจฺจนฺติ ตณฺหา.\csromanspacer{} โย ราโค สาราโค ฯเปฯ อภิชฺฌา โลโภ อกุสลมูลนฺติ ปุตฺเตสุ ทาเรสุ จ ยา อเปกฺขา.}

\prose{\textbf{วํสกฺกฬีโรว อสชฺชมาโน}ติ วํโส วุจฺจติ เวฬุคุมฺโพ.\csromanspacer{} ยถา เวฬุคุมฺพสฺมิํ ตรุณกา กฬีรกา\footnote{ตรุณกฬีรา (สฺยา)} อสตฺตา อลคฺคา อคธิตา\footnote{อปลิเวฏฺฐิตา (สฺยา)} อปลิพุทฺธา นิกฺขนฺตา นิสฺสฏา วิปฺปมุตฺตา เอวเมว.\csromanspacer{} \textbf{สชฺชา}ติ เทฺว สชฺชนา ตณฺหาสชฺชนา จ ทิฏฺฐิสชฺชนา จ ฯเปฯ อยํ ตณฺหาสชฺชนา ฯเปฯ อยํ ทิฏฺฐิสชฺชนา.\csromanspacer{} ตสฺส}

\csromanheader{2}{19. ขคฺควิสาณสุตฺตนิทฺเทส}
\prose{\csromanfolio{243}ปจฺเจกสมฺพุทฺธสฺส ตณฺหาสชฺชนา ปหีนา\textbf{,} ทิฏฺฐิสชฺชนา ปฏินิสฺสฏฺฐา.\csromanspacer{} ตณฺหาสชฺชนาย ปหีนตฺตา ทิฏฺฐิสชฺชนาย ปฏินิสฺสฏฺฐตฺตา โส ปจฺเจกสมฺพุทฺโธ รูเป น สชฺชติ.\csromanspacer{} สทฺเท น สชฺชติ.\csromanspacer{} คนฺเธ น สชฺชติ.\csromanspacer{} รเส น สชฺชติ.\csromanspacer{} โผฏฺฐพฺเพ น สชฺชติ.\csromanspacer{} กุเล.\csromanspacer{} คเณ.\csromanspacer{} อาวาเส.\csromanspacer{} ลาเภ.\csromanspacer{} ยเส.\csromanspacer{} ปสํสาย.\csromanspacer{} สุเข.\csromanspacer{} จีวเร.\csromanspacer{} ปิณฺฑปาเต.\csromanspacer{} เสนาสเน.\csromanspacer{} คิลานปจฺจยเภสชฺชปริกฺขาเร.\csromanspacer{} กามธาตุยา.\csromanspacer{} รูปธาตุยา.\csromanspacer{} อรูปธาตุยา.\csromanspacer{} กามภเว.\csromanspacer{} รูปภเว.\csromanspacer{} อรูปภเว.\csromanspacer{} สญฺญาภเว.\csromanspacer{} อสญฺญาภเว.\csromanspacer{} เนวสญฺญานาสญฺญาภเว.\csromanspacer{} เอกโวการภเว.\csromanspacer{} จตุโวการภเว.\csromanspacer{} ปญฺจโวการภเว.\csromanspacer{} อตีเต.\csromanspacer{} อนาคเต.\csromanspacer{} ปจฺจุปฺปนฺเน.\csromanspacer{} ทิฏฺฐสุตมุตวิญฺญาตพฺเพสุ ธมฺเมสุ น สชฺชติ น คณฺหาติ น พชฺฌติ น ปลิพชฺฌติ น มุจฺฉติ นิกฺขนฺโต นิสฺสโฏ วิปฺปมุตฺโต วิสญฺญุตฺโต วิมริยาทิกเตน เจตสา วิหรตีติ วํสกฺกฬีโรว อสชฺชมาโน\textbf{,} เอโก จเร ขคฺควิสาณกปฺโป.\csromanspacer{} เตนาห โส ปจฺเจกสมฺพุทฺโธ\csromandash{}}

\csromangathameasure{วํโส วิสาโลว ยถา วิสตฺโต\textbf{,} \\ ปุตฺเตสุ ทาเรสุ จ ยา อเปกฺขา. \\ วํสกฺกฬีโรว อสชฺชมาโน\textbf{,} \\ เอโก จเร ขคฺควิสาณกปฺโปติ.}
\csromangathagroup{\csromangathastanza{วํโส วิสาโลว ยถา วิสตฺโต\textbf{,} \\ ปุตฺเตสุ ทาเรสุ จ ยา อเปกฺขา. \\ วํสกฺกฬีโรว อสชฺชมาโน\textbf{,} \\ เอโก จเร ขคฺควิสาณกปฺโปติ.}}

\csromanhangingbegin
\hangingheaditem{125}{มิโค อรญฺญมฺหิ ยถา อพทฺโธ\footnote{อพนฺโธ (สฺยา, ก)}\textbf{,}}
\hangingbody{\textbf{เยนิจฺฉกํ คจฺฉติ โคจราย.}}
\hangingbody{\textbf{วิญฺญู นโร เสริตํ เปกฺขมาโน.}}
\hangingbody{เอโก จเร ขคฺควิสาณกปฺโป.}
\csromanhangingend

\prose{\textbf{มิโค อรญฺญมฺหิ ยถา อพทฺโธ, เยนิจฺฉกํ คจฺฉติ โคจรายา}ติ \textbf{มิโค}ติ เทฺว มิคา เอณิ\textbf{มิโค} จ ปสท\textbf{มิโค} จ.\csromanspacer{} ยถา อารญฺญิโก\footnote{อารญฺญโก (สฺยา, ก)} \textbf{มิโค} อรญฺเญ ปวเน จรมาโน\footnote{อรญฺเญ วสมาโน (สฺยา)} วิสฺสตฺโถ คจฺฉติ\textbf{,} วิสฺสตฺโถ ติฏฺฐติ\textbf{,} วิสฺสตฺโถ นิสีทติ\textbf{,} วิสฺสตฺโถ เสยฺยํ กปฺเปติ.}

\prose{วุตฺตํเหตํ ภควตา\csromandash{}เสยฺยถาปิ ภิกฺขเว อารญฺญิโก \textbf{มิโค} อรญฺเญ ปวเน จรมาโน วิสฺสตฺโถ คจฺฉติ\textbf{,} วิสฺสตฺโถ ติฏฺฐติ\textbf{,} วิสฺสตฺโถ นิสีทติ\textbf{,} วิสฺสตฺโถ เสยฺยํ กปฺเปติ.\csromanspacer{} ตํ กิสฺส เหตุ\textbf{,} อนาปาถคโต ภิกฺขเว ลุทฺทสฺส.\csromanspacer{} เอวเมว โข ภิกฺขเว ภิกฺขุ \csromanfolio{244}วิวิจฺเจว กาเมหิ วิวิจฺจ อกุสเลหิ ธมฺเมหิ สวิตกฺกํ สวิจารํ วิเวกชํ ปีติสุขํ ปฐมํ ฌานํ อุปสมฺปชฺช วิหรติ.\csromanspacer{} อยํ วุจฺจติ ภิกฺขเว ภิกฺขุ อนฺธมกาสิ มารํ, อปทํ\footnote{ปรํ (ก) ม 1. 214 ปิฏฺเฐ.} วธิตฺวา มารจกฺขุํ อทสฺสนํ คโต ปาปิมโต.}

\prose{ปุน จปรํ ภิกฺขเว ภิกฺขุ วิตกฺกวิจารานํ วูปสมา อชฺฌตฺตํ สมฺปสาทนํ เจตโส เอโกทิภาวํ อวิตกฺกํ อวิจารํ สมาธิชํ ปีติสุขํ ทุติยํ ฌานํ อุปสมฺปชฺช วิหรติ.\csromanspacer{} อยํ วุจฺจติ ภิกฺขเว ภิกฺขุ อนฺธมกาสิ มารํ, อปทํ วธิตฺวา มารจกฺขุํ อทสฺสนํ คโต ปาปิมโต.}

\prose{ปุน จปรํ ภิกฺขเว ภิกฺขุ ปีติยา จ วิราคา อุเปกฺขโก จ วิหรติ, สโต จ สมฺปชาโน สุขญฺจ กาเยน ปฏิสํเวเทติ.\csromanspacer{} ยํ ตํ อริยา อาจิกฺขนฺติ “อุเปกฺขโก สติมา สุขวิหารี”ติ ตติยํ ฌานํ อุปสมฺปชฺช วิหรติ, อยํ วุจฺจติ ภิกฺขเว ภิกฺขุ อนฺธมกาสิ มารํ, อปทํ วธิตฺวา มารจกฺขุํ อทสฺสนํ คโต ปาปิมโต.}

\prose{ปุน จปรํ ภิกฺขเว ภิกฺขุ สุขสฺส จ ปหานา ทุกฺขสฺส จ ปหานา ปุพฺเพว โสมนสฺสโทมนสฺสานํ อตฺถงฺคมา อทุกฺขมสุขํ อุเปกฺขาสติปาริสุทฺธิํ จตุตฺถํ ฌานํ อุปสมฺปชฺช วิหรติ, อยํ วุจฺจติ ภิกฺขเว ภิกฺขุ อนฺธมกาสิ มารํ, อปทํ วธิตฺวา มารจกฺขุํ อทสฺสนํ คโต ปาปิมโต.}

\prose{ปุน จปรํ ภิกฺขเว ภิกฺขุ สพฺพโส รูปสญฺญานํ สมติกฺกมา ปฏิฆสญฺญานํ อตฺถงฺคมา นานตฺตสญฺญานํ อมนสิการา “อนนฺโต อากาโส”ติ อากาสานญฺจายตนํ อุปสมฺปชฺช วิหรติ, อยํ วุจฺจติ ภิกฺขเว ภิกฺขุ อนฺธมกาสิ มารํ, อปทํ วธิตฺวา มารจกฺขุํ อทสฺสนํ คโต ปาปิมโต.}

\prose{ปุน จปรํ ภิกฺขเว ภิกฺขุ สพฺพโส อากาสานญฺจายตนํ สมติกฺกมฺม “อนนฺตํ วิญฺญาณนฺ”ติ วิญฺญาณญฺจายตนํ อุปสมฺปชฺช วิหรติ ฯเปฯ.}

\prose{ปุน จปรํ ภิกฺขเว สพฺพโส วิญฺญาณญฺจายตนํ สมติกฺกมฺม “นตฺถิ กิญฺจี”ติ อากิญฺจญฺญายตนํ อุปสมฺปชฺช วิหรติ ฯเปฯ.}

\prose{ปุน จปรํ ภิกฺขเว สพฺพโส อากิญฺจญฺญายตนํ สมติกฺกมฺม เนวสญฺญานาสญฺญายตนํ อุปสมฺปชฺช วิหรติ ฯเปฯ.}

\csromanheader{2}{19. ขคฺควิสาณสุตฺตนิทฺเทส}
\prose{\csromanfolio{245}ปุน จปรํ ภิกฺขเว สพฺพโส เนวสญฺญานาสญฺญายตนํ สมติกฺกมฺม สญฺญาเวทยิตนิโรธํ อุปสมฺปชฺช วิหรติ.\csromanspacer{} ปญฺญาย จสฺส ทิสฺวา อาสวา ปริกฺขีณา โหนฺติ.\csromanspacer{} อยํ วุจฺจติ ภิกฺขเว ภิกฺขุ อนฺธมกาสิ มารํ, อปทํ วธิตฺวา มารจกฺขุํ อทสฺสนํ คโต ปาปิมโต, ติณฺโณ โลเก วิสตฺติกํ, โส วิสฺสตฺโถ คจฺฉติ, วิสฺสตฺโถ ติฏฺฐติ, วิสฺสตฺโถ นิสีทติ, วิสฺสตฺโถ เสยฺยํ กปฺเปติ.\csromanspacer{} ตํ กิสฺส เหตุ, อนาปาถคโต ภิกฺขเว ปาปิมโตติ มิโค อรญฺญมฺหิ ยถา อพทฺโธ, เยนิจฺฉกํ คจฺฉติ โคจราย.}

\prose{\textbf{วิญฺญู นโร เสริตํ เปกฺขมาโน}ติ \textbf{วิญฺญู}ติ ปณฺฑิโต ปญฺญวา พุทฺธิมา ญาณี วิภาวี เมธาวี.\csromanspacer{} \textbf{นโร}ติ สตฺโต มานโว โปโส ปุคฺคโล ชีโว ชาคุ ชนฺตุ อินฺทคุ มนุโช.\csromanspacer{} \textbf{เสรี}ติ เทฺว เสรี ธมฺโมปิ เสรี ปุคฺคโลปิ เสรี.\csromanspacer{} กตโม ธมฺโม เสรี, จตฺตาโร สติปฏฺฐานา จตฺตาโร สมฺมปฺปธานา จตฺตาโร อิทฺธิปาทา ปญฺจินฺทฺริยานิ ปญฺจ พลานิ สตฺต โพชฺฌงฺคา อริโย อฏฺฐงฺคิโก มคฺโค.\csromanspacer{} อยํ ธมฺโม เสรี.\csromanspacer{} กตโม ปุคฺคโล เสรี.\csromanspacer{} โย อิมินา เสรินา ธมฺเมน สมนฺนาคโต.\csromanspacer{} โส วุจฺจติ ปุคฺคโล เสรี.\csromanspacer{} \textbf{วีญฺญู นโร เสริตํ เปกฺขมาโน}ติ \textbf{วิญฺญู} นโร เสริตํ ธมฺมํ เปกฺขมาโน ทกฺขมาโน โอโลกยมาโน นิชฺฌายมาโน อุปปริกฺขมาโนติ \textbf{วิญฺญู} นโร เสริตํ เปกฺขมาโน, เอโก จเร ขคฺควิสาณกปฺโป.\csromanspacer{} เตนาห โส ปจฺเจกสมฺพุทฺโธ\csromandash{}}

\csromangathameasure{มิโค อรญฺญมฺหิ ยถา อพทฺโธ, \\ เยนิจฺฉกํ คจฺฉติ โคจราย.}
\csromangathagroup{\csromangathastanza{มิโค อรญฺญมฺหิ ยถา อพทฺโธ, \\ เยนิจฺฉกํ คจฺฉติ โคจราย.}}

\prose{\textbf{วิญฺญู นโร เสริตํ เปกฺขมาโน},}

\prose{เอโก จเร ขคฺควิสาณกปฺโปติ.}

\csromanhangingbegin
\hangingheaditem{126}{อามนฺตนา โหติ สหายมชฺเฌ,}
\hangingbody{\textbf{วาเส ฐาเน คมเน จาริกาย.}}
\hangingbody{\textbf{อนภิชฺฌิตํ เสริตํ เปกฺขมาโน,}}
\hangingbody{เอโก จเร ขคฺควิสาณกปฺโป.}
\csromanhangingend

\prose{\textbf{อามนฺตนา โหติ สหายมชฺเฌ, วาเส ฐาเน คมเน จาริกายา}ติ สหายา วุจฺจนฺติ เยหิ สห อาคมนํ ผาสุ, คมนํ ผาสุ, คมนาคมนํ ผาสุ, ฐานํ ผาสุ, นิสชฺชนํ ผาสุ, สยนํ ผาสุ, อาลปนํ \csromanfolio{246}ผาสุ, สลฺลปนํ ผาสุ, อุลฺลปนํ ผาสุ, สมุลฺลปนํ ผาสุ.\csromanspacer{} \textbf{อามนฺตนา โหติ สหายมชฺเฌ, วาเส ฐาเน คมเน จาริกายา}ติ สหายมชฺเฌ วาเสปิ ฐาเนปิ คมเนปิ จาริกายปิ อตฺตตฺถมนฺตนา ปรตฺถมนฺตนา อุภยตฺถมนฺตนา ทิฏฺฐธมฺมิกตฺถมนฺตนา สมฺปรายิกตฺถมนฺตนา ปรมตฺถมนฺตนา1ติ อามนฺตนา โหติ สหายมชฺเฌ, วาเส ฐาเน คมเน จาริกาย.}

\prose{\textbf{อนภิชฺฌิตํ เสริตํ เปกฺขมาโน}ติ อนภิชฺฌิตํ เอตํ วตฺถุ พาลานํ อสปฺปุริสานํ ติตฺถิยานํ ติตฺถิยสาวกานํ, ยทิทํ ภณฺฑุกาสายวตฺถวสนตา.\csromanspacer{} อภิชฺฌิตํ เอตํ วตฺถุ ปณฺฑิตานํ สปฺปุริสานํ พุทฺธสาวกานํ ปจฺเจกพุทฺธานํ, ยทิทํ ภณฺฑุกาสายวตฺถวสนตา.\csromanspacer{} \textbf{เสรี}ติ เทฺว เสรี ธมฺโมปิ เสรี ปุคฺคโลปิ เสรี.\csromanspacer{} กตโม ธมฺโม เสรี.\csromanspacer{} จตฺตาโร สติปฏฺฐานา ฯเปฯ อริโย อฏฺฐงฺคิโก มคฺโค, อยํ ธมฺโม เสรี.\csromanspacer{} กตโม ปุคฺคโล เสรี.\csromanspacer{} โย อิมินา เสรินา ธมฺเมน สมนฺนาคโต, โส วุจฺจติ ปุคฺคโล เสรี.\csromanspacer{} \textbf{อนภิชฺฌิตํ เสริตํ เปกฺขมาโน}ติ เสริตํ ธมฺมํ เปกฺขมาโน ทกฺขมาโน โอโลกยมาโน นิชฺฌายมาโน อุปปริกฺขมาโนติ อนภิชฺฌิตํ เสริตํ เปกฺขมาโน, เอโก จเร ขคฺควิสาณกปฺโป.\csromanspacer{} เตนาห โส ปจฺเจกสมฺพุทฺโธ\csromandash{}}

\csromangathameasure{อามนฺตนา โหติ สหายมชฺเฌ, \\ วาเส ฐาเน คมเน จาริกาย.}
\csromangathagroup{\csromangathastanza{อามนฺตนา โหติ สหายมชฺเฌ, \\ วาเส ฐาเน คมเน จาริกาย.}}

\prose{\textbf{อนภิชฺฌิตํ เสริตํ เปกฺขมาโน},}

\prose{เอโก จเร ขคฺควิสาณกปฺโปติ.}

\csromanhangingbegin
\hangingheaditem{127}{\textbf{ขิฑฺฑา}\footnote{ขิฏฺฏา (ก)}\textbf{ รตี โหติ สหายมชฺเฌ},}
\hangingbody{\textbf{ปุตฺเตสุ จ วิปุลํ โหติ เปมํ.}}
\hangingbody{\textbf{ปิยวิปฺปโยคํ3} \textbf{วิชิคุจฺฉมาโน,}}
\hangingbody{เอโก จเร ขคฺควิสาณกปฺโป.}
\csromanhangingend

\prose{\textbf{ขิฑฺฑา รตี โหติ สหายมชฺเฌ}ติ \textbf{ขิฑฺฑา}ติ เทฺว \textbf{ขิฑฺฑา} กายิกา จ ขิฏฺฏา วาจสิกา จ \textbf{ขิฑฺฑา}.\csromanspacer{} กตมา กายิกา \textbf{ขิฑฺฑา}.\csromanspacer{} หตฺถีหิปิ กีฬนฺติ, อสฺเสหิปิ กีฬนฺติ, รเถหิปิ กีฬนฺติ, ธนุหิปิ กีฬนฺติ, ถรูหิปิ กีฬนฺติ, อฏฺฐปเทหิปิ กีฬนฺติ, ทสปเทหิปิ กีฬนฺติ, อากาเสปิ กีฬนฺติ, ปริหารปเถปิ กีฬนฺติ, สนฺติกายปิ กีฬนฺติ, ขลิกายปิ กีฬนฺติ,}

\csromanheader{2}{19. ขคฺควิสาณสุตฺตนิทฺเทส}
\prose{\csromanfolio{247}ฆฏิกายปิ กีฬนฺติ, สลากหตฺเถนปิ กีฬนฺติ, อกฺเขนปิ กีฬนฺติ, ปงฺคจีเรนปิ กีฬนฺติ, วงฺกเกนปิ กีฬนฺติ, โมกฺขจิกายปิ กีฬนฺติ, จิงฺคุลเกนปิ กีฬนฺติ, ปตฺตาฬฺหเกนปิ กีฬนฺติ, รถเกนปิ กีฬนฺติ, ธนุเกนปิ กีฬนฺติ, อกฺขริกายปิ กีฬนฺติ, มเนสิกายปิ กีฬนฺติ, ยถาวชฺเชนปิ กีฬนฺติ, อยํ กายิกา ขิฑฺฑา.}

\prose{กตมา วาจสิกา ขิฑฺฑา.\csromanspacer{} มุขเภริกํ มุขาลมฺพรํ มุขฑิณฺฑิมกํ\footnote{มุขเทณฺฑิมกํ (สฺยา), มุขทินฺทิมกํ (ก)} มุขจลิมกํ มุขเกรกํ\footnote{มุขเภรุกํ (สฺยา)} มุขททฺทริกํ นาฏกํ ลาสํ คีตํ ทวกมฺมํ.\csromanspacer{} อยํ วาจสิกา ขิฑฺฑา.}

\prose{\textbf{รตี}ติ อนุกฺกณฺฐิตาธิวจนเมตํ รตีติ.\csromanspacer{} \textbf{สหายา} วุจฺจนฺติ เยหิ สห อาคมนํ ผาสุ, คมนํ ผาสุ, คมนาคมนํ ผาสุ, ฐานํ ผาสุ, นิสชฺชนํ ผาสุ, สยนํ ผาสุ, อาลปนํ ผาสุ, สลฺลปนํ ผาสุ, อุลฺลปนํ ผาสุ, สมุลฺลปนํ ผาสุ.\csromanspacer{} \textbf{ขิฑฺฑา รตี โหติ สหายมชฺเฌ}ติ ขิฑฺฑา จ รติ จ สหายมชฺเฌ โหตีติ ขิฑฺฑา รตี โหติ สหายมชฺเฌ.}

\prose{\textbf{ปุตฺเตสุ จ วิปุลํ โหติ เปมนฺ}ติ \textbf{ปุตฺตา}ติ จตฺตาโร \textbf{ปุตฺตา} อตฺรโช ปุตฺโต เขตฺตโช ปุตฺโต ทินฺนโก ปุตฺโต อนฺเตวาสิโก ปุตฺโต.\csromanspacer{} \textbf{ปุตฺเตสุ จ วิปุลํ โหติ เปมนฺ}ติ ปุตฺเตสุ จ อธิมตฺตํ โหติ เปมนฺติ ปุตฺเตสุ จ วิปุลํ โหติ เปมํ.}

\prose{\textbf{ปิยวิปฺปโยคํ วิชิคุจฺฉมาโน}ติ เทฺว ปิยา สตฺตา วา สงฺขารา วา.\csromanspacer{} กตเม สตฺตา ปิยา.\csromanspacer{} อิธ ย’สฺส เต โหนฺติ อตฺถกามา หิตกามา ผาสุกามา โยคกฺเขมกามา มาตา วา ปิตา วา ภาตา วา ภคินี วา ปุตฺโต วา ธีตา วา มิตฺตา วา อมจฺจา วา ญาตี วา สาโลหิตา วา.\csromanspacer{} อิเม สตฺตา ปิยา.}

\prose{กตเม สงฺขารา ปิยา.\csromanspacer{} มนาปิกา รูปา มนาปิกา สทฺทา มนาปิกา คนฺธา มนาปิกา รสา มนาปิกา โผฏฺฐพฺพา.\csromanspacer{} อิเม สงฺขารา ปิยา.\csromanspacer{} \textbf{ปิยวิปฺปโยคํ วิชิคุจฺฉมาโน}ติ ปิยานํ วิปฺปโยคํ วิชิคุจฺฉมาโน อฏฺฏิยมาโน หรายมาโนติ ปิยวิปฺปโยคํ วิชิคุจฺฉมาโน, เอโก จเร ขคฺควิสาณกปฺโป.\csromanspacer{} เตนาห โส ปจฺเจกสมฺพุทฺโธ\csromandash{}}

\csromangathameasure{ขิฑฺฑา รตี โหติ สหายมชฺเฌ, \\ ปุตฺเตสุ จ วิปุลํ โหติ เปมํ. \\ ปิยวิปฺปโยคํ วิชิคุจฺฉมาโน, \\ เอโก จเร ขคฺควิสาณกปฺโปติ.}
\csromangathagroup{\csromangathastanza{\csromanfolio{248}ขิฑฺฑา รตี โหติ สหายมชฺเฌ, \\ ปุตฺเตสุ จ วิปุลํ โหติ เปมํ. \\ ปิยวิปฺปโยคํ วิชิคุจฺฉมาโน, \\ เอโก จเร ขคฺควิสาณกปฺโปติ.}}

\csromanhangingbegin
\hangingheaditem{128}{จาตุทฺทิโส อปฺปฏิโฆ จ โหติ,}
\hangingbody{\textbf{สนฺตุสฺสมาโน อิตรีตเรน.}}
\hangingbody{\textbf{ปริสฺสยานํ สหิตา อฉมฺภี,}}
\hangingbody{เอโก จเร ขคฺควิสาณกปฺโป.}
\csromanhangingend

\prose{\textbf{จาตุทฺทิโส อปฺปฏิโฆ จ โหตี}ติ \textbf{จาตุทฺทิโส}ติ โส ปจฺเจกสมฺพุทฺโธ เมตฺตาสหคเตน เจตสา เอกํ ทิสํ ผริตฺวา วิหรติ.\csromanspacer{} ตถา ทุติยํ.\csromanspacer{} ตถา ตติยํ.\csromanspacer{} ตถา จตุตฺถํ.\csromanspacer{} อิติ อุทฺธมโธ ติริยํ สพฺพธิ สพฺพตฺตตาย สพฺพาวนฺตํ โลกํ เมตฺตาสหคเตน เจตสา วิปุเลน มหคฺคเตน อปฺปมาเณน อเวเรน อพฺยาปชฺเชน\footnote{อพฺยาปชฺเฌน (สฺยา) ปสฺส ที 3. 187 ปิฏฺเฐ.} ผริตฺวา วิหรติ.\csromanspacer{} กรุณาสหคเตน ฯเปฯ.\csromanspacer{} มุทิตาสหคเตน ฯเปฯ.\csromanspacer{} อุเปกฺขาสหคเตน เจตสา เอกํ ทิสํ ผริตฺวา วิหรติ.\csromanspacer{} ตถา ทุติยํ.\csromanspacer{} ตถา ตติยํ ฯเปฯ อพฺยาปชฺเชน ผริตฺวา วิหรติ.\csromanspacer{} \textbf{จาตุทฺทิโส อปฺปฏิโฆ จ โหตี}ติ เมตฺตาย ภาวิตตฺตา เย ปุรตฺถิมาย ทิสาย สตฺตา, เต อปฺปฏิกูลา\footnote{อปฺปฏิกุลา (พหูสุ)} โหนฺติ, เย ทกฺขิณาย ทิสาย สตฺตา, เต อปฺปฏิกูลา โหนฺติ.\csromanspacer{} เย ปจฺฉิมาย ทิสาย สตฺตา, เต อปฺปฏิกูลา โหนฺติ.\csromanspacer{} เย อุตฺตราย ทิสาย สตฺตา, เต อปฺปฏิกูลา โหนฺติ.\csromanspacer{} เย ปุรตฺถิมาย อนุทิสาย สตฺตา, เต อปฺปฏิกูลา โหนฺติ.\csromanspacer{} เย ทกฺขิณาย อนุทิสาย สตฺตา, เต อปฺปฏิกูลา โหนฺติ.\csromanspacer{} เย ปจฺฉิมาย อนุทิสาย สตฺตา, เต อปฺปฏิกูลา โหนฺติ.\csromanspacer{} เย อุตฺตราย อนุทิสาย สตฺตา, เต อปฺปฏิกูลา โหนฺติ.\csromanspacer{} เย เหฏฺฐิมาย ทิสาย สตฺตา, เต อปฺปฏิกูลา โหนฺติ.\csromanspacer{} เย อุปริมาย ทิสาย สตฺตา, เต อปฺปฏิกูลา โหนฺติ.\csromanspacer{} เย ทิสาสุ วิทิสาสุ สตฺตา, เต อปฺปฏิกูลา โหนฺติ.\csromanspacer{} กรุณาย ภาวิตตฺตา.\csromanspacer{} มุทิตาย ภาวิตตฺตา.\csromanspacer{} อุเปกฺขาย ภาวิตตฺตา เย ปุรตฺถิมาย ทิสาย สตฺตา, เต อปฺปฏิกูลา โหนฺติ ฯเปฯ.\csromanspacer{} เย ทิสาสุ วิทิสาสุ สตฺตา, เต อปฺปฏิกูลา โหนฺตีติ \textbf{จาตุทฺทิโส} อปฺปฏิโฆ จ โหติ.}

\csromanheader{2}{19. ขคฺควิสาณสุตฺตนิทฺเทส}
\prose{\csromanfolio{249}\textbf{สนฺตุสฺสมาโน อิตรีตเรนา}ติ โส ปจฺเจกสมฺพุทฺโธ สนฺตุฏฺโฐ โหติ อิตรีตเรน จีวเรน, อิตรีตรจีวรสนฺตุฏฺฐิยา จ วณฺณวาที, น จ จีวรเหตุ อเนสนํ อปฺปติรูปํ\footnote{อปฺปฏิรูปํ (สฺยา)} อาปชฺชติ.\csromanspacer{} อลทฺธา จ จีวรํ น ปริตสฺสติ, ลทฺธา จ จีวรํ อคธิโต อมุจฺฉิโต อนชฺโฌสนฺโน\footnote{อนชฺโฌปนฺโน (สฺยา)} อาทีนวทสฺสาวี นิสฺสรณปญฺโญ ปริภุญฺชติ.\csromanspacer{} ตาย จ ปน อิตรีตรจีวรสนฺตุฏฺฐิยา เนวตฺตานุกฺกํเสติ, น ปรํ วมฺเภติ, โย หิ ตตฺถ ทกฺโข อนลโส สมฺปชาโน ปติสฺสโต.\csromanspacer{} อยํ วุจฺจติ ปจฺเจกสมฺพุทฺโธ โปราเณ อคฺคญฺเญ อริยวํเส ฐิโต.\csromanspacer{} สนฺตุฏฺโฐ โหติ อิตรีตเรน ปิณฺฑปาเตน ฯเปฯ.\csromanspacer{} สนฺตุฏฺโฐ โหติ อิตรีตเรน เสนาสเนน ฯเปฯ.\csromanspacer{} สนฺตุฏฺโฐ โหติ อิตรีตเรน คิลานปจฺจยเภสชฺชปริกฺขาเรน.\csromanspacer{} อิตรีตรคิลานปจฺจยเภสชฺช- ปริกฺขารสนฺตุฏฺฐิยา จ วณฺณวาที, น จ คิลานปจฺจยเภสชฺช- ปริกฺขารเหตุ อเนสนํ อปฺปติรูปํ อาปชฺชติ.\csromanspacer{} อลทฺธา จ คิลานปจฺจยเภสชฺชปริกฺขารานํ น ปริตสฺสติ, ลทฺธา จ คิลานปจฺจยเภสชฺชปริกฺขารํ อคธิโต อมุจฺฉิโต อนชฺโฌสนฺโน อาทีนวทสฺสาวี นิสฺสรณปญฺโญ ปริภุญฺชติ.\csromanspacer{} ตาย จ อิตรีตรคิลานปจฺจยเภสชฺชปริกฺขารสนฺตุฏฺฐิยา เนวตฺตานุกฺกํเสติ, น ปรํ วมฺเภติ, โย หิ ตตฺถ ทกฺโข อนลโส สมฺปชาโน ปติสฺสโต.\csromanspacer{} อยํ วุจฺจติ ปจฺเจกสมฺพุทฺโธ โปราเณ อคฺคญฺเญ อริยวํเส ฐิโตติ สนฺตุสฺสมาโน อิตรีตเรน.}

\prose{\textbf{ปริสฺสยานํ สหิตา อฉมฺภี}ติ \textbf{ปริสฺสยา}ติ เทฺว \textbf{ปริสฺสยา} ปากฏ\textbf{ปริสฺสยา} จ ปฏิจฺฉนฺน\textbf{ปริสฺสยา} จ.\csromanspacer{} กตเม \textbf{ปากฏปริสฺสยา.}\csromanspacer{} สีหา พฺยคฺฆาทีปี อจฺฉา ตรจฺฉา โกกา มหิํสา\footnote{โคมหิสา (สฺยา) ขุ 7. 10 ปิฏฺเฐสุปิ.} หตฺถี อหี วิจฺฉิกา สตปที โจรา วา อสฺสุ มานวา วา กตกมฺมา วา อกตกมฺมา วา, จกฺขุโรโค โสตโรโค ฆานโรโค ชิวฺหาโรโค กายโรโค สีสโรโค กณฺณโรโค มุขโรโค ทนฺตโรโค กาโส สาโส ปินาโส ฑาโห ชโร กุจฺฉิโรโค มุจฺฉา ปกฺขนฺทิกา สูลา วิสูจิกา กุฏฺฐํ คณฺโฑ กิลาโส โสโส อปมาโร ททฺทุ กณฺฑุ กจฺฉุ รขสา วิตจฺฉิกา โลหิตปิตฺตํ มธุเมโห อํสา ปิฬกา ภคนฺทลา ปิตฺตสมุฏฺฐานา อาพาธา, เสมฺหสมุฏฺฐานา อาพาธา, วาตสมุฏฺฐานา อาพาธา, สนฺนิปาติกา อาพาธา, อุตุปริณามชา อาพาธา,}

\prose{\csromanfolio{250}วิสมปริหารชา อาพาธา, โอปกฺกมิกา อาพาธา กมฺมวิปากชา อาพาธา, สีตํ อุณฺหํ ชิฆจฺฉา ปิปาสา อุจฺจาโร ปสฺสาโว ฑํสมกสวาตาตปสรีสปสมฺผสฺสา อิติ วา, อิเม วุจฺจนฺติ ปากฏปริสฺสาย.}

\prose{กตเม \textbf{ปฏิจฺฉนฺนปริสฺสยา.}\csromanspacer{} กายทุจฺจริตํ วจีทุจฺจริตํ มโนทุจฺจริตํ กามจฺฉนฺทนีวรณํ พฺยาปาทนีวรณํ ถินมิทฺธนีวรณํ อุทฺธจฺจกุกฺกุจฺจนีวรณํ วิจิกิจฺฉานีวรณํ ราโค โทโส โมโห โกโธ อุปนาโห มกฺโข ปฬาโส อิสฺสา มจฺฉริยํ มายา สาเฐยฺยํ ถมฺโภ สารมฺโภ มาโน อติมาโน มโท ปมาโท สพฺเพ กิเลสา สพฺเพ ทุจฺจริตา สพฺเพ ทรถา สพฺเพ ปริฬาหา สพฺเพ สนฺตาปา สพฺพากุสลาภิสงฺขารา, อิเม วุจฺจนฺติ \textbf{ปฏิจฺฉนฺนปริสฺสยา.}}

\prose{\textbf{ปริสฺสยา}ติ เกนฏฺเฐน ปริสฺสยา, ปริสหนฺตีติ ปริสฺสยา, ปริหานาย สํวตฺตนฺตีติ ปริสฺสยา, ตตฺราสยาติ ปริสฺสยา.\csromanspacer{} กถํ ปริสหนฺตีติ ปริสฺสยา, เต ปริสฺสยา ตํ ปุคฺคลํ สหนฺติ ปริสหนฺติ อภิภวนฺติ อชฺโฌตฺถรนฺติ ปริยาทิยนฺติ มทฺทนฺติ, เอวํ ปริสหนฺตีติ ปริสฺสยา.}

\prose{กถํ ปริหานาย สํวตฺตนฺตีติ ปริสฺสยา, เต ปริสฺสยา กุสลานํ ธมฺมานํ อนฺตรายาย ปริหานาย สํวตฺตนฺติ, กตเมสํ กุสลานํ ธมฺมานํ, สมฺมาปฏิปทาย อนุโลมปฏิปทาย อปจฺจนีกปฏีปทาย อนฺวตฺถปฏิปทาย ธมฺมานุธมฺมปฏิปทาย สีเลสุ ปริปูรการิตาย อินฺทฺริเยสุ คุตฺตทฺวารตาย โภชเน มตฺตญฺญุตาย ชาคริยานุโยคสฺส สติสมฺปชญฺญสฺส จตุนฺนํ สติปฏฺฐานานํ ภาวนานุโยคสฺส จตุนฺนํ สมฺมปฺปธานานํ.\csromanspacer{} จตุนฺนํ อิทฺธิปาทานํ.\csromanspacer{} ปญฺจนฺนํ อินฺทฺริยานํ.\csromanspacer{} ปญฺจนฺนํ พลานํ.\csromanspacer{} สตฺตนฺนํ โพชฺฌงฺคานํ.\csromanspacer{} อริยสฺส อฏฺฐงฺคิกสฺส มคฺคสฺส ภาวนานุโยคสฺส อิเมสํ กุสลานํ ธมฺมานํ อนฺตรายาย ปริหานาย สํวตฺตนฺติ, เอวํ ปริหานาย สํวตฺตนฺตีติ ปริสฺสยา.}

\prose{กถํ ตตฺราสยาติ ปริสฺสยา, ตตฺเถเต\footnote{ตตฺร เต (ก)} ปาปกา อกุสลา ธมฺมา อุปฺปชฺชนฺติ อตฺตภาวสนฺนิสฺสยา, ยถา พิเล พิลาสยา ปาณา สยนฺติ, ทเก ทกาสยา\footnote{อุทเก อุทกาสยา (สฺยา)} ปาณา สยนฺติ, วเน วนาสยา ปาณา สยนฺติ, รุกฺเข รุกฺขาสยา ปาณา สยนฺติ.\csromanspacer{} เอวเมว ตตฺเถเต ปาปกา อกุสลา ธมฺมา อุปฺปชฺชนฺติ อตฺตภาวสนฺนิสฺสยาติ, เอวมฺปิ ตตฺรสยาติ ปริสฺสยา.}

\csromanheader{2}{19. ขคฺควิสาณสุตฺตนิทฺเทส}
\prose{\csromanfolio{251}วุตฺตํเหตํ ภควตา\csromandash{}สานฺเตวาสิโก ภิกฺขเว ภิกฺขุ สาจริยโก ทุกฺขํ น ผาสุ วิหรติ.\csromanspacer{} กถญฺจ ภิกฺขเว ภิกฺขุ สานฺเตวาสิโก สาจริยโก ทุกฺขํ น ผาสุ วิหรติ.\csromanspacer{} อิธ ภิกฺขเว จกฺขุนา รูปํ ทิสฺวา อุปฺปชฺชนฺติ เย ปาปกา อกุสลา ธมฺมา สรสงฺกปฺปา สญฺโญชนียา\footnote{สญฺโญชนิกา (ก)}, ตฺยสฺส อนฺโต วสนฺติ อนฺวาสฺสวนฺติ ปาปกา อกุสลา ธมฺมาติ ตสฺมา สานฺเตวาสิโก วุจฺจติ, เตน สมุทาจเรน สมุทาจรนฺติ นํ “ปาปกา อกุสลา ธมฺมา”ติ ตสฺมา “สาจริยโก”ติ วุจฺจติ.}

\prose{ปุน จปรํ ภิกฺขเว ภิกฺขุโน โสเตน สทฺทํ สุตฺวา ฯเปฯ ฆาเนน คนฺธํ ฆายิตฺวา ฯเปฯ ชิวฺหาย รสํ สายิตฺวา ฯเปฯ กาเยน โผฏฺฐพฺพํ ผุสิตฺวา ฯเปฯ มนสา ธมฺมํ วิญฺญาย อุปฺปชฺชนฺติ เย ปาปกา อกุสลา ธมฺมา สรสงฺกปฺปา สญฺโญชนียา.\csromanspacer{} ตฺยสฺส อนฺโต วสนฺติ อนฺวาสฺสวนฺติ ปาปกา อกุสลา ธมฺมาติ ตสฺมา “สานฺเตวาสิโก”ติ วุจฺจติ.\csromanspacer{} เตน สมุทาจเรน สมุทาจรนฺติ นํ “ปาปกา อกุสลา ธมฺมา”ติ ตสฺมา “สาจริยโก”ติ วุจฺจติ.\csromanspacer{} เอวํ โข ภิกฺขเว ภิกฺขุ สานฺเตวาสิโก สาจริยโก ทุกฺขํ น ผาสุ วิหรตีติ.\csromanspacer{} เอวมฺปิ ตตฺราสยาติ ปริสฺสยา.}

\prose{วุตฺตํ เหตํ ภควตา\csromandash{}ตโยเม ภิกฺขเว อนฺตรามลา อนฺตรา-อมิตฺตา อนฺตราสปตฺตา อนฺตราวธกา อนฺตราปจฺจตฺถิกา.\csromanspacer{} กตเม ตโย.\csromanspacer{} โลโภ ภิกฺขเว อนฺตรามโล\footnote{อนฺตรามลํ (สฺยา, ก) ปสฺส ขุ 1. 252 ปิฏฺเฐ.} อนฺตรา-อมิตฺโต อนฺตราสปตฺโต อนฺตราวธโก อนฺตราปจฺจตฺถิโก, โทโส ภิกฺขเว ฯเปฯ โมโห ภิกฺขเว อนฺตรามโล อนฺตรา- อมิตฺโต อนฺตราสปตฺโต อนฺตราวธโก อนฺตราปจฺจตฺถิโก.\csromanspacer{} อิเม โข ภิกฺขเว ตโย อนฺตรามลา อนฺตรา-อมิตฺตา อนฺตาราสปตฺตา อนฺตราวธกา อนฺตราปจฺจตฺถิกาติ.}

\csromangathameasure{อนตฺถชนโน โลโภ, โลโภ จิตฺตปฺปโกปโน. \\ ภยมนฺตรโต ชาตํ, ตํ ชโน นาวพุชฺฌติ. \\ ลุทฺโธ อตฺถํ น ชานาติ, ลุทฺโธ ธมฺมํ น ปสฺสติ. \\ อนฺธตมํ ตทา โหติ, ยํ โลโภ สหเต นรํ. \\ อนตฺถชนโน โทโส, โทโส จิตฺตปฺปโกปโน. \\ ภยมนฺตรโต ชาตํ, ตํ ชโน นาวพุชฺฌติ. \\ ทุฏฺโฐ อตฺถํ น ชานาติ, ทุฏฺโฐ ธมฺมํ น ปสฺสติ. \\ อนฺธตมํ ตทา โหติ, ยํ โทโส สหเต นรํ. \\ อนตฺถชนโน โมโห, โมโห จิตฺตปฺปโกปโน. \\ ภยมนฺตรโต ชาตํ, ตํ ชโน นาวพุชฺฌติ. \\ มูโฬฺห อตฺถํ น ชานาติ, มูโฬฺห ธมฺมํ น ปสฺสติ. \\ อนฺธตมํ ตทา โหติ, ยํ โมโห สหเต นรนฺติ.}
\csromangathagroup{\csromangathastanza{อนตฺถชนโน โลโภ, โลโภ จิตฺตปฺปโกปโน. \\ ภยมนฺตรโต ชาตํ, ตํ ชโน นาวพุชฺฌติ.}\csromangathastanza{ลุทฺโธ อตฺถํ น ชานาติ, ลุทฺโธ ธมฺมํ น ปสฺสติ. \\ อนฺธตมํ ตทา โหติ, ยํ โลโภ สหเต นรํ.}\csromangathastanza{อนตฺถชนโน โทโส, โทโส จิตฺตปฺปโกปโน. \\ ภยมนฺตรโต ชาตํ, ตํ ชโน นาวพุชฺฌติ.}\csromangathastanza{\csromanfolio{252}ทุฏฺโฐ อตฺถํ น ชานาติ, ทุฏฺโฐ ธมฺมํ น ปสฺสติ. \\ อนฺธตมํ ตทา โหติ, ยํ โทโส สหเต นรํ.}\csromangathastanza{อนตฺถชนโน โมโห, โมโห จิตฺตปฺปโกปโน. \\ ภยมนฺตรโต ชาตํ, ตํ ชโน นาวพุชฺฌติ.}\csromangathastanza{มูโฬฺห อตฺถํ น ชานาติ, มูโฬฺห ธมฺมํ น ปสฺสติ. \\ อนฺธตมํ ตทา โหติ, ยํ โมโห สหเต นรนฺติ.}}

\prose{เอวมฺปิ ตตฺราสยาติ ปริสฺสยา.}

\prose{วุตฺตเญฺหตํ ภควตา\csromandash{}ตโย โข มหาราช ปุริสสฺส ธมฺมา อชฺฌตฺตํ อุปฺปชฺชมานา อุปฺปชฺชนฺติ อหิตาย ทุกฺขาย อผาสุวิหาราย.\csromanspacer{} กตเม ตโย, โลโภ โข มหาราช ปุริสสฺส ธมฺโม อชฺฌตฺตํ อุปฺปชฺชมาโน อุปฺปชฺชติ อหิตาย ทุกฺขาย อผาสุวิหาราย, โทโส โข มหาราช ฯเปฯ โมโห โข มหาราช ปุริสสฺส ธมฺโม อชฺฌตฺตํ อุปฺปชฺชมาโน อุปฺปชฺชติ อหิตาย ทุกฺขาย อผาสุวิหาราย, อิเม โข มหาราช ตโย ปุริสสฺส ธมฺมา อชฺฌตฺตํ อุปฺปชฺชมานา อุปฺปชฺชนฺติ อหิตาย ทุกฺขาย อผาสุวิหาราย.}

\csromangathameasure{โลโภ โทโส จ โมโห จ, ปุริสํ ปาปเจตสํ. \\ หิํสนฺติ อตฺตสมฺภูตา, ตจสารํว สมฺผลนฺติ.}
\csromangathagroup{\csromangathastanza{โลโภ โทโส จ โมโห จ, ปุริสํ ปาปเจตสํ. \\ หิํสนฺติ อตฺตสมฺภูตา, ตจสารํว สมฺผลนฺติ\footnote{สผลนฺติ (ก) ขุ 1. 226 ปิฏฺเฐ.}.}}

\prose{เอวมฺปิ ตตฺราสยาติ ปริสฺสยา.}

\csromanheader{2}{วุตฺตเญฺหตํ ภควตา\csromandash{}}
\csromangathameasure{ราโค จ โทโส จ อิโตนิทานา, \\ อรตี รตี โลมหํโส อิโตชา. \\ อิโต สมุฏฺฐาย มโนวิตกฺกา, \\ กุมารกา ธงฺกมิโวสฺสชนฺตีติ.}
\csromangathagroup{\csromangathastanza{ราโค จ โทโส จ อิโตนิทานา, \\ อรตี รตี โลมหํโส อิโตชา. \\ อิโต สมุฏฺฐาย มโนวิตกฺกา, \\ กุมารกา ธงฺกมิโวสฺสชนฺตีติ\footnote{ธงฺกมิโวสฺสชฺชนฺตีติ (สฺยา, ก) ปสฺส ขุ 1. 320 ปิฏฺเฐ.}.}}

\prose{เอวมฺปิ ตตฺราสยาติ ปริสฺสยา.}

\prose{\textbf{ปริสฺสยานํ สหิตา}ติ ปริสฺสเย สหิตา อาราธิตา อชฺโฌตฺถริตา ปริยาทิตา ปฏินิสฺสตาติ ปริสฺสยานํ สหิตา.\csromanspacer{} \textbf{อฉมฺภี}ติ โส}

\csromanheader{2}{19. ขคฺควิสาณสุตฺตนิทฺเทส}
\prose{\csromanfolio{253}ปจฺเจกสมฺพุทฺโธ อภีรู อจฺฉมฺภี อนุตฺราสี อปลายี ปหีนภยเภรโว วิคตโลมหํโส วิหรตีติ ปริสฺสยานํ สหิตา อฉมฺภี, เอโก จเร ขคฺควิสาณกปฺโป.\csromanspacer{} เตนาห โส ปจฺเจกสมฺพุทฺโธ\csromandash{}}

\csromangathameasure{จาตุทฺทิโส อปฺปฏิโฆ จ โหติ, \\ สนฺตุสฺสมาโน อิตรีตเรน. \\ ปริสฺสยานํ สหิตา อฉมฺภี, \\ เอโก จเร ขคฺควิสาณกปฺโปติ.}
\csromangathagroup{\csromangathastanza{จาตุทฺทิโส อปฺปฏิโฆ จ โหติ, \\ สนฺตุสฺสมาโน อิตรีตเรน. \\ ปริสฺสยานํ สหิตา อฉมฺภี, \\ เอโก จเร ขคฺควิสาณกปฺโปติ.}}

\proseitem{129}{\textbf{ทุสฺสงฺคหา ปพฺพชิตาปิ เอเก},}

\prose{\textbf{อโถ คหฏฺฐา ฆรมาวสนฺตา.}}

\prose{\textbf{อปฺโปสฺสุกฺโก ปรปุตฺเตสุ หุตฺวา,}}

\csromanheader{2}{\textbf{เอโก จเร ขคฺควิสาณกปฺโป. (9)}}
\prose{\textbf{ทุสฺสงฺคหา ปพฺพชิตาปิ เอเก}ติ ปพฺพชิตาปิ อิเธกจฺเจ นิสฺสเยปิ ทิยฺยมาเน อุทฺเทเสปิ ทิยฺยมาเน ปริปุจฺฉายปิ\footnote{ปริปุจฺเฉปิ (ก)} ทิยฺยมาเน จีวเรปิ ทิยฺยมาเน ปตฺเตปิ ทิยฺยมาเน โลหถาลเกปิ ทิยฺยมาเน ธมฺมกรเณปิ\footnote{ธมฺมกรเกปิ (สฺยา)} ทิยฺยมาเน ปริสฺสาวเนปิ ทิยฺยมาเน ถวิเกปิ ทิยฺยมาเน อุปาหเนปิ ทิยฺยมาเน กายพนฺธเนปิ ทิยฺยมาเน น สุณนฺติ, น โสตํ โอทหนฺติ, น อญฺญา จิตฺตํ อุปฏฺฐเปนฺติ, อนสฺสวา อวจนกรา ปฏิโลมวุตฺติโน อญฺเญเนว มุขํ กโรนฺตีติ ทุสฺสงฺคหา ปพฺพชิตาปิ เอเก.}

\prose{\textbf{อโถ คหฏฺฐา ฆรมาวสนฺตา}ติ คหฏฺฐาปิ อิเธกจฺเจ หตฺถิมฺหิปิ ทิยฺยมาเน ฯเปฯ.\csromanspacer{} รเถปิ.\csromanspacer{} เขตฺเตปิ.\csromanspacer{} วตฺถุมฺหิปิ.\csromanspacer{} หิรญฺเญปิ.\csromanspacer{} สุวณฺเณปิ ทิยฺยมาเน คาเมปิ.\csromanspacer{} นิคเมปิ.\csromanspacer{} นคเรปิ.\csromanspacer{} รฏฺเฐปิ.\csromanspacer{} ชนปเทปิ ทิยฺยมาเน น สุณนฺติ น โสตํ โอทหนฺติ, น อญฺญาจิตฺตํ อุปฏฺฐเปนฺติ, อนสฺสวา อวจนกรา ปฏิโลมวุตฺติโน อญฺเญเนว มุขํ กโรนฺตีติ อโถ คหฏฺฐา ฆรมาวสนฺตา.}

\prose{\textbf{อปฺโปสฺสุกฺโก ปรปุตฺเตสุ หุตฺวา}ติ อตฺตานํ ฐเปตฺวา สพฺเพ อิมสฺมิํ อตฺเถ ปรปุตฺตา เตสุ ปรปุตฺเตสุ อปฺโปสฺสุกฺโก หุตฺวา อพฺยาวโฏ หุตฺวา อนเปกฺโข หุตฺวาติ อปฺโปสฺสุกฺโก ปรปุตฺเตสุ หุตฺวา, เอโก จเร ขคฺควิสาณกปฺโป.\csromanspacer{} เตนาห โส ปจฺเจกสมฺพุทฺโธ\csromandash{}}

\csromangathameasure{ทุสฺสงฺคหา ปพฺพชิตาปิ เอเก, \\ อโถ คหฏฺฐา ฆรมาวสนฺตา. \\ อปฺโปสฺสุกฺโก ปรปุตฺเตสุ หุตฺวา, \\ เอโก จเร ขคฺควิสาณกปฺโปติ.}
\csromangathagroup{\csromangathastanza{\csromanfolio{254}ทุสฺสงฺคหา ปพฺพชิตาปิ เอเก, \\ อโถ คหฏฺฐา ฆรมาวสนฺตา. \\ อปฺโปสฺสุกฺโก ปรปุตฺเตสุ หุตฺวา, \\ เอโก จเร ขคฺควิสาณกปฺโปติ.}}

\proseitem{130}{โอโรปยิตฺวา\footnote{โวโรปยิตฺวา (สฺยา)} \textbf{คิหิพฺยญฺชนานิ,}}

\prose{\textbf{สญฺฉินฺนปตฺโต}\footnote{สํสีนปตฺโต (สี-ฏฺฐ)}\textbf{ ยถา โกวิฬาโร}.}

\prose{\textbf{เฉตฺวาน วีโร คิหิพนฺธนานิ,}}

\csromanheader{2}{\textbf{เอโก จเร ขคฺควิสาณกปฺโป. (10)}}
\prose{\textbf{โอโรปยิตฺวา คิหิพฺยญฺชนานี}ติ คิหิพฺยญฺชนานิ วุจฺจนฺติ เกสา จ มสฺสู จ มาลา จ คนฺธญฺจ วิเลปนญฺจ อาภรณญฺจ ปิลนฺธนญฺจ วตฺถญฺจ ปารุปนญฺจ เวฐนญฺจ อุจฺฉาทนํ ปริมทฺทนํ นฺหาปนํ\footnote{นหาปนํ (สฺยา)} สมฺพาหนํ อาทาสํ อญฺชนํ มาลาคนฺธวิเลปนํ มุขจุณฺณํ มุขเลปนํ หตฺถพนฺธํ สิขาพนฺธํ ทณฺฑํ นาฬิกํ\footnote{นาลิกํ (ก) ปสฺส ที 1. 7 ปิฏฺเฐ.} ขคฺคํ ฉตฺตํ จิตฺรูปาหนํ อุณฺหีสํ มณิํ วาลพีชนิํ โอทาตานิ วตฺถานิ ทีฆทสานิ\footnote{ทีฆรสฺสานิ (สฺยา)} อิติ วา.\csromanspacer{} \textbf{โอโรปยิตฺวา คิหิพฺยญฺชนานี}ติ คิหิพฺยญฺชนานิ โอโรปยิตฺวา สโมโรปยิตฺวา นิกฺขิปิตฺวา ปฏิปสฺสมฺภยิตฺวาติ โอโรปยิตฺวา คิหิพฺยญฺชนานิ.}

\prose{\textbf{สญฺฉินฺนปตฺโต ยถา โกวิฬาโร}ติ ยถา โกวิฬารสฺส ปตฺตานิ ฉินฺนานิ สญฺฉินฺนานิ ปติตานิ ปริปติตานิ, เอวเมว ตสฺส ปจฺเจกสมฺพุทฺธสฺส คิหิพฺยญฺชนานิ ฉินฺนานิ สญฺฉินฺนานิ ปติตานีติ สญฺฉินฺนปตฺโต \textbf{ยถา โกวิฬาโร.}}

\prose{\textbf{เฉตฺวาน วีโร คิหิพนฺธนานี}ติ \textbf{วีโร}ติ วีริยวาติ \textbf{วีโร}, ปหูติ \textbf{วีโร}, วิสวีติ \textbf{วีโร}, อลมตฺโตติ \textbf{วีโร}, สูโรติ \textbf{วีโร}, วิกฺกนฺโต อภิรู อจฺฉมฺภี อนุตฺราสี อปลายี ปหีนภยเภรโวติ \textbf{วีโร}, วิคตโลมหํโสติ \textbf{วีโร}.}

\csromangathameasure{วิรโต อิธ สพฺพปาเกหิ, \\ นิรยทุกฺขํ อติจฺจ วีริยวาโส. \\ โส วีริยวา ปธานวา, \\ ธีโร ตาทิ ปวุจฺจเต ตถตฺตา.}
\csromangathagroup{\csromangathastanza{วิรโต อิธ\footnote{อารโต อิเธว (สฺยา) ปสฺส ขุ 1. 360 ปิฏฺเฐ.} สพฺพปาเกหิ, \\ นิรยทุกฺขํ อติจฺจ วีริยวาโส. \\ โส วีริยวา ปธานวา, \\ ธีโร\footnote{อารโต อิเธว (สฺยา) ปสฺส ขุ 1. 360 ปิฏฺเฐ.} ตาทิ ปวุจฺจเต ตถตฺตา.}}

\csromanheader{2}{19. ขคฺควิสาณสุตฺตนิทฺเทส}
\prose{\csromanfolio{255}คิหิพนฺธนานิ วุจฺจนฺติ ปุตฺตา จ ภริยา จ ทาสา จ ทาสี จ อเชฬกา จ กุกฺกุฏสูกรา จ หตฺถิควาสฺสวฬวา จ เขตฺตญฺจ วตฺถุ จ หิรญฺญญฺจ สุวณฺณญฺจ คามนิคมราชธานิโย จ รฏฺฐญฺจ ชนปโท จ โกโส จ โกฏฺฐาคารญฺจ ยํ กิญฺจิ รชนียวตฺถุ.}

\prose{\textbf{เฉตฺวาน วีโร คิหิพนฺธนานี}ติ โส ปจฺเจกสมฺพุทฺโธ วีโร คิหิพนฺธนานิ ฉินฺทิตฺวา สมุจฺฉินฺทิตฺวา ชหิตฺวา วิโนเทตฺวา พฺยนฺตีกริตฺวา อนภาวํ คเมตฺวาติ เฉตฺวาน วีโร คิหิพนฺธนานิ, เอโก จเร ขคฺควิสาณกปฺโป.\csromanspacer{} เตนาห โส ปจฺเจกสมฺพุทฺโธ\csromandash{}}

\csromangathameasure{โอโรปยิตฺวา คิหิพฺยญฺชนานิ, \\ สญฺฉินฺนปตฺโต ยถา โกวิฬาโร. \\ เฉตฺวาน วีโร คิหิพนฺธนานิ, \\ เอโก จเร ขคฺควิสาณกปฺโปติ.}
\csromangathagroup{\csromangathastanza{โอโรปยิตฺวา คิหิพฺยญฺชนานิ, \\ สญฺฉินฺนปตฺโต ยถา โกวิฬาโร. \\ เฉตฺวาน วีโร คิหิพนฺธนานิ, \\ เอโก จเร ขคฺควิสาณกปฺโปติ.}}

\vspace{\tipitakaparskip}
\csromancenter{ปฐโม วคฺโค.}
\chapterhead{\textbf{ทุติยวคฺค}}
\csromanhangingbegin
\hangingheaditem{131}{สเจ ลเภถ นิปกํ สหายํ,}
\hangingbody{\textbf{สทฺธิํ จรํ สาธุวิหาริ ธีรํ.}}
\hangingbody{\textbf{อภิภุยฺย สพฺพานิ ปริสฺสยานิ,}}
\hangingbody{จเรยฺย เตน’ตฺตมโน สตีมา.}
\csromanhangingend

\prose{\textbf{สเจ ลเภถ นิปกํ สหายนฺ}ติ สเจ นิปกํ ปณฺฑิตํ ปญฺญวนฺตํ พุทฺธิมนฺตํ ญาณิํ วิภาวิํ เมธาวิํ สหายํ ลเภยฺย ปฏิลเภยฺย อธิคจฺเฉยฺย วินฺเทยฺยาติ สเจ ลเภถ นิปกํ สหายํ.}

\prose{\textbf{สทฺธิํจรํ สาธุวิหาริ ธีรนฺ}ติ \textbf{สทฺธิํจรนฺ}ติ เอกโต จรํ.\csromanspacer{} \textbf{สาธุวิหารินฺ}ติ ปฐเมนปิ ฌาเนน สาธุวิหาริํ.\csromanspacer{} ทุติเยนปิ ฌาเนน.\csromanspacer{} ตติเยนปิ ฌาเนน.\csromanspacer{} จตุตฺเถนปิ ฌาเนน สาธุวิหาริํ.\csromanspacer{} เมตฺตายปิ เจโตวิมุตฺติยา สาธุวิหาริํ.\csromanspacer{} กรุณายปิ.\csromanspacer{} มุทิตายปิ.\csromanspacer{} อุเปกฺขายปิ เจโตวิมุตฺติยา สาธุวิหาริํ.\csromanspacer{} อากาสานญฺจายตนสมาปตฺติยาปิ สาธุวิหาริํ.\csromanspacer{} วิญฺญาณญฺจายตนสมาปตฺติยาปิ ฯเปฯ.\csromanspacer{} อากิญฺจญฺญายตนสมาปตฺติยาปิ ฯเปฯ. \csromanfolio{256}เนวสญฺญานาสญฺญายตนสมาปตฺติยาปิ สาธุวิหาริํ นิโรธสมาปตฺติยาปิ สาธุวิหาริํ.\csromanspacer{} ผลสมาปตฺติยาปิ สาธุวิหาริํ.\csromanspacer{} \textbf{ธีรนฺ}ติ ธีรํ ปณฺฑิตํ ปญฺญวนฺตํ พุทฺธิมนฺตํ ญาณิํ วิภาวิํ เมธาวินฺติ สทฺธิํจรํ สาธุวิหาริ ธีรํ.}

\prose{\textbf{อภิภุยฺย สพฺพานิ ปริสฺสยานี}ติ \textbf{ปริสฺสยา}ติ เทฺว \textbf{ปริสฺสยา} ปากฏ\textbf{ปริสฺสยา} จ ปฏิจฺฉนฺน\textbf{ปริสฺสยา} จ ฯเปฯ อิเม วุจฺจนฺติ ปากฏ\textbf{ปริสฺสยา} ฯเปฯ อิเม วุจฺจนฺติ ปฏิจฺฉนฺน\textbf{ปริสฺสยา} ฯเปฯ.\csromanspacer{} เอวมฺปิ ตตฺราสยาติ \textbf{ปริสฺสยา}.\csromanspacer{} \textbf{อภิภุยฺย สพฺพานิ ปริสฺสยานี}ติ สพฺเพ ปริสฺสเย อภิภุยฺย อภิภวิตฺวา อชฺโฌตฺถริตฺวา ปริยาทิยิตฺวา มทฺทิตฺวาติ อภิภุยฺย สพฺพานิ \textbf{ปริสฺสยา}นิ.}

\prose{\textbf{จเรยฺย เตน’ตฺตมโน สตีมา}ติ โส ปจฺเจกสมฺพุทฺโธ เตน นิปเกน ปณฺฑิเตน ปญฺญวนฺเตน พุทฺธิมนฺเตน ญาณินา วิภาวินา เมธาวินา สหาเยน สทฺธิํ อตฺตมโน ตุฏฺฐมโน หฏฺฐมโน ปหฏฺฐมโน อุทคฺคมโน มุทิตมโน จเรยฺย วิหเรยฺย อิริเยยฺย วตฺเตยฺย ปาเลยฺย ยเปยฺย ยาเปยฺยาติ จเรยฺย เตนตฺตมโน.\csromanspacer{} \textbf{สตีมา}ติ โส ปจฺเจกสมฺพุทฺโธ สติมา โหติ ปรเมน สติเนปกฺเกน สมนฺนาคโต จิรกตมฺปิ จิรภาสิตมฺปิ สริตา อนุสฺสริตาติ จเรยฺย เตน’ตฺตมโน สตีมา.\csromanspacer{} เตนาห โส ปจฺเจกสมฺพุทฺโธ\csromandash{}}

\csromangathameasure{สเจ ลเภถ นิปกํ สหายํ, \\ \textbf{สทฺธิํจรํ สาธุวิหาริ ธีรํ.}}
\csromangathagroup{\csromangathastanza{สเจ ลเภถ นิปกํ สหายํ, \\ \textbf{สทฺธิํจรํ สาธุวิหาริ ธีรํ.}}}

\prose{อภิภุยฺย สพฺพานิ \textbf{ปริสฺสยา}นิ,}

\prose{จเรยฺย เตนตฺตมโน สตีมาติ.}

\csromanhangingbegin
\hangingheaditem{132}{โน เจ ลเภถ นิปกํ สหายํ,}
\hangingbody{\textbf{สทฺธิํจรํ สาธุวิหาริ ธีรํ.}}
\hangingbody{\textbf{ราชาว รฏฺฐํ วิชิตํ ปหาย,}}
\hangingbody{เอโก จเร ขคฺควิสาณกปฺโป.}
\csromanhangingend

\prose{\textbf{โน เจ ลเภถ นิปกํ สหายนฺ}ติ โน เจ นิปกํ ปณฺฑิตํ ปญฺญวนฺตํ พุทฺธิมนฺตํ ญาณิํ วิภาวิํ เมธาวิํ สหายํ ลเภยฺย ปฏิลเภยฺย อธิคจฺเฉยฺย วินฺเทยฺยาติ โน เจ ลเภถ นิปกํ สหายํ.}

\csromanheader{2}{19. ขคฺควิสาณสุตฺตนิทฺเทส}
\prose{\csromanfolio{257}\textbf{สทฺธิํจรํ สาธุวิหาริ ธีรนฺ}ติ \textbf{สทฺธิํจรนฺ}ติ เอกโต จรํ.\csromanspacer{} \textbf{สาธุวิหารินฺ}ติ ปฐเมนปิ ฌาเนน สาธุวิหาริํ ฯเปฯ นิโรธสมาปตฺติยาปิ สาธุวิหาริํ ผลสมาปตฺติยาปิ สาธุวิหาริํ.\csromanspacer{} \textbf{ธีรนฺ}ติ ธีรํ ปณฺฑิตํ ปญฺญวนฺตํ พุทฺธิมนฺตํ ญาณิํ วิภาวิํ เมธาวินฺติ สทฺธิํจรํ สาธุวิหาริ ธีรํ.}

\prose{\textbf{ราชาว รฏฺฐํ วิชิตํ ปหายา}ติ ราชา ขตฺติโย มุทฺธาภิสิตฺโต วิชิตสงฺคาโม นิหตปจฺจามิตฺโต ลฺ\textbf{อทฺธา}ธิปฺปาโย ปริปุณฺณโกสโกฏฺฐาคาโร รฏฺฐญฺจ ชนปทญฺจ โกสญฺจ โกฏฺฐาคารญฺจ ปหูตหิรญฺญสุวณฺณํ นครญฺจ ปริจฺจชิตฺวา เกสมสฺสุํ โอหาเรตฺวา กาสายานิ วตฺถานิ อจฺฉาเทตฺวา อคารสฺมา อนคาริยํ ปพฺพชิตฺวา อกิญฺจนภาวํ อุปคนฺตฺวา เอโก จรติ วิหรติ อิริยติ วตฺเตติ ปาเลติ ยเปติ ยาเปติ.\csromanspacer{} เอวํ ปจฺเจกสมฺพุทฺโธปิ สพฺพํ ฆราวาสปลิโพธํ ฉินฺทิตฺวา ปุตฺตทารปลิโพธํ ฉินฺทิตฺวา ญาติปลิโพธํ ฉินฺทิตฺวา มิตฺตามจฺจปลิโพธํ ฉินฺทิตฺวา เกสมสฺสุํ โอหาเรตฺวา กาสายานิ วตฺถานิ อจฺฉาเทตฺวา อคารสฺมา อนคาริยํ ปพฺพชิตฺวา อกิญฺจนภาวํ อุปคนฺตฺวา เอโก จรติ วิหรติ อิริยติ วตฺเตติ ปาเลติ ยเปติ ยาเปตีติ ราชาว รฏฺฐํ วิชิตํ ปหาย, เอโก จเร ขคฺควิสาณกปฺโป.\csromanspacer{} เตนาห โส ปจฺเจกสมฺพุทฺโธ\csromandash{}}

\prose{โน เจ ลเภถ นิปกํ สหายํ,}

\prose{สทฺธิํจรํ สาธุวิหาริ ธีรํ.}

\prose{ราชาว รฏฺฐํ วิชิตํ ปหาย,}

\prose{เอโก จเร ขคฺควิสาณกปฺโปติ.}

\csromanhangingbegin
\hangingheaditem{133}{อทฺธา ปสํสาม สหายสมฺปทํ,}
\hangingbody{\textbf{เสฏฺฐา สมา เสวิตพฺพา สหายา.}}
\hangingbody{\textbf{เอเต อลทฺธา อนวชฺชโภชี,}}
\hangingbody{เอโก จเร ขคฺควิสาณกปฺโป.}
\csromanhangingend

\prose{\textbf{อทฺธา ปสํสาม สหายสมฺปทนฺ}ติ \textbf{อทฺธา}ติ เอกํสวจนํ นิสฺสํสยวจนํ นิกฺกงฺขวจนํ อเทฺวชฺฌวจนํ อเทฺวฬฺหกวจนํ นิโยควจนํ อปณฺณกวจนํ อวิรทฺธวจนํ อวตฺถาปนวจนเมตํ \textbf{อทฺธา}ติ.\csromanspacer{} \textbf{สหายสมฺปทนฺ}ติ สหายสมฺปทา วุจฺจติ โย โส สหาโย อเสกฺเขน สีลกฺขนฺเธน สมนฺนาคโต โหติ.\csromanspacer{} อเสกฺเขน สมาธิกฺขนฺเธน.\csromanspacer{} อเสกฺเขน \csromanfolio{258}ปญฺญากฺขนฺเธน.\csromanspacer{} อเสกฺเขน วิมุตฺติกฺขนฺเธน.\csromanspacer{} อเสกฺเขน วิมุตฺติญาณทสฺสนกฺขนฺเธน สมนฺนาคโต โหติ.\csromanspacer{} \textbf{อทฺธา ปสํสาม สหายสมฺปทนฺ}ติ สหายสมฺปทํ ปสํสาม โถเมม กิตฺเตม วณฺเณมาติ อทฺธา ปสํสาม สหายสมฺปทํ.}

\prose{\textbf{เสฏฺฐา สมา เสวิตพฺพา สหายา}ติ เสฏฺฐา โหนฺติ สหายา สีเลน สมาธินา ปญฺญาย วิมุตฺติยา วิมุตฺติญาณทสฺสเนน, สมา สทิสา โหนฺติ สหาย สีเลน สมาธินา ปญฺญาย วิมุตฺติยา วิมุตฺติญาณทสฺสเนน, เสฏฺฐา วา สหายา สทิสา วา สหายา เสวิตพฺพา ภชิตพฺพา ปยิรุปาสิตพฺพา ปริปุจฺฉิตพฺพา ปริปญฺหิตพฺพาติ เสฏฺฐา สมา เสวิตพฺพา สหายา.}

\prose{\textbf{เอเต อลทฺธา อนวชฺชโภชี}ติ อตฺถิ ปุคฺคโล สาวชฺชโภชี อตฺถิ ปุคฺคโล อนวชฺชโภชีติ.\csromanspacer{} กตโม จ ปุคฺคโล สาวชฺชโภชี.\csromanspacer{} อิเธกจฺโจ ปุคฺคโล กุหนาย ลปนาย เนมิตฺติกตาย นิปฺเปสิกตาย ลาเภน ลาภํ นิชิคีสนตาย\footnote{นิชิคิํสนตาย (สฺยา)} ทารุทาเนน เวฬุทาเนน ปตฺตทาเนน ปุปฺผทาเนน ผลทาเนน สินานทาเนน จุณฺณทาเนน มตฺติกาทาเนน ทนฺตกฏฺฐทาเนน มุโขทกทาเนน จาฏุกมฺยตาย\footnote{ปาตุกมฺยตาย (สฺยา) ขุ 7. 290 ปิฏฺเฐปิ.} มุคฺคสูปฺยตาย\footnote{มุคฺคสูปตาย (สฺยา)} ปาริภฏฺยตาย ปีฐมทฺทิกตาย วตฺถุวิชฺชาย ติรจฺฉานวิชฺชาย องฺควิชฺชาย นกฺขตฺตวิชฺชาย ทูตคมเนน ปหิณคมเนน ชงฺฆเปสนิเยน เวชฺชกมฺเมน นวกมฺเมน\footnote{ทูตกมฺเมน (สฺยา, ก) ขุ 7. 290 ปิฏฺเฐ.} ปิณฺฑปฏิปิณฺฑเกน ทานานุปฺปทาเนน อธมฺเมน วิสเมน ลทฺธา ลภิตฺวา อธิคนฺตฺวา วินฺทิตฺวา ปฏิลภิตฺวา ชีวิกํ\footnote{ชีวิตํ (ก)} กปฺเปติ.\csromanspacer{} อยํ วุจฺจติ ปุคฺคโล สาวชฺชโภชี.}

\prose{กตโม จ ปุคฺคโล อนวชฺชโภชี.\csromanspacer{} อิเธกจฺโจ ปุคฺคโล น กุหนาย น ลปนาย น เนมิตฺติกตาย น นิปฺเปสิกตาย น ลาเภน ลาภํ นิชิคีสนตาย น ทารุทาเนน น เวฬุทาเนน น ปตฺตทาเนน น ปุปฺผทาเนน น ผลทาเนน น สินานทาเนน น จุณฺณทาเนน น มตฺติกาทาเนน น ทนฺตกฏฺฐทาเนน น มุโขทกทาเนน น จาฏุกมฺยตาย น มุคฺคสูปฺยตาย น ปาริภฏฺยตาย น ปีฐมทฺทิกตาย น วตฺถุวิชฺชาย น ติรจฺฉานวิชฺชาย}

\csromanheader{2}{19. ขคฺควิสาณสุตฺตนิทฺเทส}
\prose{\csromanfolio{259}น องฺควิชฺชาย น นกฺขตฺตวิชฺชาย น ทูตคมเนน น ปหิณคมเนน น ชงฺฆเปสนิเยน น เวชฺชกมฺเมน น นวกมฺเมน น ปิณฺฑปฏิปิณฺฑเกน น ทานานุปฺปทาเนน ธมฺเมน สเมน ลทฺธา ลภิตฺวา อธิคนฺตฺวา วินฺทิตฺวา ปฏิลภิตฺวา ชีวิกํ กปฺเปติ.\csromanspacer{} อยํ วุจฺจติ ปุคฺคโล อนวชฺชโภชี.}

\prose{\textbf{เอเต อลทฺธา อนวชฺชโภชี}ติ เอเต อนวชฺชโภชี อลทฺธา อลภิตฺวา อนธิคนฺตฺวา อวินฺทิตฺวา อปฺปฏิลภิตฺวาติ เอเต อลทฺธา อนวชฺชโภชี, เอโก จเร ขคฺควิสาณกปฺโป.\csromanspacer{} เตนาห โส ปจฺเจกสมฺพุทฺโธ\csromandash{}}

\csromangathameasure{อทฺธา ปสํสาม สหายสมฺปทํ, \\ เสฏฺฐา สมา เสวิตพฺพา สหายา. \\ \textbf{เอเต อลทฺธา อนวชฺชโภชี}, \\ เอโก จเร ขคฺควิสาณกปฺโปติ.}
\csromangathagroup{\csromangathastanza{อทฺธา ปสํสาม สหายสมฺปทํ, \\ เสฏฺฐา สมา เสวิตพฺพา สหายา. \\ \textbf{เอเต อลทฺธา อนวชฺชโภชี}, \\ เอโก จเร ขคฺควิสาณกปฺโปติ.}}

\proseitem{134}{ทิสฺวา สุวณฺณสฺส ปภสฺสรานิ,}

\nitthitam{\textbf{กมฺมารปุตฺเตน สุนิฏฺฐิตานิ.}}
\csromangathameasure{\textbf{สํฆฏฺฏยนฺตานิ} \textbf{ทุเว ภุชสฺมิํ,} \\ เอโก จเร ขคฺควิสาณกปฺโป.}
\csromangathagroup{\csromangathastanza{\textbf{สํฆฏฺฏยนฺตานิ}\footnote{สํฆฏฺฏมานานิ (ขุ 1. 286 ปิฏฺเฐ.)} \textbf{ทุเว ภุชสฺมิํ,} \\ เอโก จเร ขคฺควิสาณกปฺโป.}}

\prose{\textbf{ทิสฺวา สุวณฺณสฺส ปภสฺสรานี}ติ ทิสฺวา ปสฺสิตฺวา ตุลยิตฺวา ตีรยิตฺวา วิภาวยิตฺวา วิภูตํ กตฺวา.\csromanspacer{} \textbf{สุวณฺณสฺสา}ติ ชาตรูปสฺส.\csromanspacer{} \textbf{ปภสฺสรานี}ติ ปริสุทฺธานิ ปริโยทาตานีติ ทิสฺวา สุวณฺณสฺส ปภสฺสรานิ.}

\prose{\textbf{กมฺมารปุตฺเตน สุนิฏฺฐิตานี}ติ กมฺมารปุตฺโต วุจฺจติ สุวณฺณกาโร.\csromanspacer{} \textbf{กมฺมารปุตฺเตน สุนิฏฺฐิตานี}ติ กมฺมารปุตฺเตน สุนิฏฺฐิตานิ สุกตานิ สุปริกมฺมกตานีติ กมฺมารปุตฺเตน สุนิฏฺฐิตานิ.}

\prose{\textbf{สํฆฏฺฏยนฺตานิ ทุเว ภุชสฺมินฺ}ติ ภุโช วุจฺจติ หตฺโถ.\csromanspacer{} ยถา เอกสฺมิํ หตฺเถ เทฺว นูปุรานิ\footnote{ธุวรานิ (สฺยา)} ฆฏฺเฏนฺติ\footnote{ฆเฏนฺติ (สฺยา)}, เอวเมว สตฺตา ตณฺหาวเสน ทิฏฺฐิวเสน นิรเย ฆฏฺเฏนฺติ, ติรจฺฉานโยนิยํ ฆฏฺเฏนฺติ, เปตฺติวิสเย ฆฏฺเฏนฺติ, มนุสฺสาโลเก ฆฏฺเฏนฺติ, เทวโลเก ฆฏฺเฏนฺติ, คติยา คติํ}

\prose{\csromanfolio{260}อุปปตฺติยา อุปปตฺติํ ปฏิสนฺธิยา ปฏิสนฺธิํ ภเวน ภวํ สํสาเรน สํสารํ วฏฺเฏน วฏฺฏํ ฆฏฺเฏนฺติ สํฆฏฺเฏนฺติ สํฆฏฺเฏนฺตา จรนฺติ วิหรนฺติ อิริยนฺติ วตฺเตนฺติ ปาเลนฺติ ยเปนฺติ ยาเปนฺตีติ สํฆฏฺฏยนฺตานิ ทุเว ภุชสฺมิํ, เอโก จเร ขคฺควิสาณกปฺโป.\csromanspacer{} เตนาห โส ปจฺเจกสมฺพุทฺโธ\csromandash{}}

\prose{ทิสฺวา สุวณฺณสฺส ปภสฺสรานิ,}

\nitthitam{กมฺมารปุตฺเตน สุนิฏฺฐิตานิ.}
\csromangathameasure{สํฆฏฺฏยนฺตานิ ทุเว ภุชสฺมิํ, \\ เอโก จเร ขคฺควิสาณกปฺโป.}
\csromangathagroup{\csromangathastanza{สํฆฏฺฏยนฺตานิ ทุเว ภุชสฺมิํ, \\ เอโก จเร ขคฺควิสาณกปฺโป.}}

\csromanhangingbegin
\hangingheaditem{135}{เอวํ ทุตีเยน สห มมสฺส,}
\hangingbody{\textbf{วาจาภิลาโป อภิสชฺชนา วา.}}
\hangingbody{\textbf{เอตํ ภยํ อายติํ เปกฺขมาโน,}}
\hangingbody{เอโก จเร ขคฺควิสาณกปฺโป.}
\csromanhangingend

\prose{\textbf{เอวํ ทุตีเยน สห มมสฺสา}ติ \textbf{ตณฺหา}ทุติโย วา โหติ ปุคฺคลทุติโย วา.\csromanspacer{} กถํ \textbf{ตณฺหา}ทุติโย โหติ, \textbf{ตณฺหา}ติ รูป\textbf{ตณฺหา} ฯเปฯ ธมฺม\textbf{ตณฺหา}.\csromanspacer{} ยสฺเสสา \textbf{ตณฺหา} อปฺปหีนา, โส วุจฺจติ \textbf{ตณฺหา}ทุติโย.}

\csromangathameasure{ตณฺหาทุติโย ปุริโส, ทีฆมทฺธาน สํสรํ. \\ อิตฺถภาวญฺญถาภาวํ, สํสารํ นาติวตฺตตีติ.}
\csromangathagroup{\csromangathastanza{ตณฺหาทุติโย ปุริโส, ทีฆมทฺธาน สํสรํ. \\ อิตฺถภาวญฺญถาภาวํ, สํสารํ นาติวตฺตตีติ.}}

\prose{เอวํ \textbf{ตณฺหา}ทุติโย วา โหติ.}

\prose{กถํ ปุคฺคลทุติโย โหติ, อิเธกจฺโจ น อตฺถเหตุ\footnote{อตฺตเหตุ (สฺยา)} น การณเหตุ อุทฺธโต อวูปสนฺตจิตฺโต เอกสฺส วา ทุติโย โหติ ทฺวินฺนํ วา ตติโย โหติ, ติณฺณํ วา จตุตฺโถ โหติ.\csromanspacer{} ตตฺถ พหุํ สมฺผปฺปลาปํ ปลปติ.\csromanspacer{} เสยฺยถิทํ, ราชกถํ โจรกถํ มหามตฺตกถํ เสนากถํ ภยกถํ ยุทฺธกถํ อนฺนกถํ ปานกถํ วตฺถกถํ สยนกถํ มาลากถํ คนฺธกถํ ญาติกถํ ยานกถํ คามกถํ นิคมกถํ นครกถํ ชนปทกถํ อิตฺถิกถํ\footnote{อิตฺถิกถํ ปุริสกถํ (สฺยา, ก) ที 1. 7, 62 ปิฏฺเฐ; สํ 3. 368 ปิฏฺเฐ จ ปสฺสิตพฺพํ.} สูรกถํ วิสิขากถํ กุมฺภฏฺฐานกถํ ปุพฺพเปตกถํ}

\csromanheader{2}{19. ขคฺควิสาณสุตฺตนิทฺเทส}
\prose{\csromanfolio{261}นานตฺตกถํ โลกกฺขายิกํ สมุทฺทกฺขายิกํ อิติภวาภวกถํ กเถติ.\csromanspacer{} เอวํ ปุคฺคลทุติโย โหตีติ เอวํ ทุตีเยน สห มมสฺส.}

\prose{\textbf{วาจาภิลาโป อภิสชฺชนา วา}ติ วาจาภิลาโป วุจฺจติ พาตฺติํสติรจฺฉานกถา.\csromanspacer{} เสยฺยถิทํ, ราชกถํ ฯเปฯ อิติภวาภวกถํ.\csromanspacer{} \textbf{อภิสชฺชนา วา}ติ เทฺว สชฺชนา, ตณฺหาสชฺชนา จ ทิฏฺฐิสชฺชนา จ ฯเปฯ อยํ ตณฺหาสชฺชนา ฯเปฯ อยํ ทิฏฺฐิสชฺชนาติ วาจาภิลาโป อภิสชฺชนา วา.}

\prose{\textbf{เอตํ ภยํ อายติํ เปกฺขมาโน}ติ \textbf{ภยนฺ}ติ ชาติภยํ ชราภยํ พฺยาธิภยํ มรณภยํ ราชภยํ โจรภยํ อคฺคิภยํ อุทกภยํ อตฺตานุวาทภยํ ปรานุวาทภยํ ทณฺฑภยํ ทุคฺคติภยํ อูมิภยํ กุมฺภิลภยํ อาวฏฺฏภยํ สุสุมารภยํ อาชีวิกภยํ อสิโลกภยํ ปริสสารชฺชภยํ มทนภยํ ภยานกํ ฉมฺภิตตฺตํ โลมหํโส เจตโส อุพฺเพโค อุตฺราโส.\csromanspacer{}\csromanspacer{}\textbf{เอตํ ภยํ อายติํ เปกฺขมาโน}ติ เอตํ ภยํ อายติํ เปกฺขมาโน ทกฺขมาโน โอโลกยมาโน นิชฺฌายมาโน อุปปริกฺขมาโนติ เอตํ ภยํ อายติํ เปกฺขมาโน, เอโก จเร ขคฺควิสาณกปฺโป.\csromanspacer{} เตนาห โส ปจฺเจกสมฺพุทฺโธ\csromandash{}}

\csromangathameasure{เอวํ ทุตีเยน สห มมสฺส, \\ \textbf{วาจาภิลาโป อภิสชฺชนา วา}.}
\csromangathagroup{\csromangathastanza{เอวํ ทุตีเยน สห มมสฺส, \\ \textbf{วาจาภิลาโป อภิสชฺชนา วา}.}}

\prose{\textbf{เอตํ ภยํ อายติํ เปกฺขมาโน},}

\prose{เอโก จเร ขคฺควิสาณกปฺโป.}

\csromanhangingbegin
\hangingheaditem{136}{\textbf{กามา หิ จิตฺรา มธุรา มโนรมา},}
\hangingbody{\textbf{วิรูปรูเปน มเถนฺติ จิตฺตํ.}}
\hangingbody{\textbf{อาทีนวํ กามคุเณสุ ทิสฺวา,}}
\hangingbody{เอโก จเร ขคฺควิสาณกปฺโป.}
\csromanhangingend

\prose{\textbf{กามา หิ จิตฺรา มธุรา มโนรมา}ติ \textbf{กามา}ติ อุทฺทานโต เทฺว \textbf{กามา} วตฺถุ\textbf{กามา} จ กิเลส\textbf{กามา} จ ฯเปฯ.\csromanspacer{} อิเม วุจฺจนฺติ วตฺถุ\textbf{กามา} ฯเปฯ.\csromanspacer{} อิเม วุจฺจนฺติ กิเลส\textbf{กามา}.\csromanspacer{} \textbf{จิตฺรา}ติ นานาวณฺณา รูปา นานาวณฺณา สทฺทา นานาวณฺณา คนฺธา นานาวณฺณา รสา นานาวณฺณา โผฏฺฐพฺพา อิฏฺฐา กนฺตา มนาปา ปิยรูปา กามูปสํหิตา รชนียา.\csromanspacer{} \textbf{มธุรา}ติ วุตฺตํ เหตํ ภควตา\footnote{ม 1. 120 ปิฏฺเฐ; ม 3. 276 ปิฏฺเฐ จ ปสฺสิตพฺพํ.} \csromanfolio{262}ปญฺจิเม ภิกฺขเว กามคุณา.\csromanspacer{} กตเม ปญฺจ, จกฺขุวิญฺเญยฺยา รูปา อิฏฺฐา กนฺตา มนาปา ปิยรูปา กามูปสํหิตา รชนียา.\csromanspacer{} โสตวิญฺเญยฺยา สทฺทา ฯเปฯ.\csromanspacer{} ฆานวิญฺเญยฺยา คนฺธา.\csromanspacer{} ชิวฺหาวิญฺเญยฺยา รสา.\csromanspacer{} กายวิญฺเญยฺยา โผฏฺฐพฺพา อิฏฺฐา กนฺตา มนาปา ปิยรูปา กามูปสํหิตา รชนียา.\csromanspacer{} อิเม โข ภิกฺขเว ปญฺจ กามคุณา.\csromanspacer{} ยํ โข ภิกฺขเว อิเม ปญฺจ กามคุเณ ปฏิจฺจ อุปฺปชฺชติ สุขํ โสมนสฺสํ, อิทํ วุจฺจติ กามสุขํ มิฬฺหสุขํ\footnote{มีฬฺหสุขํ (ปสฺส-ม 3. 276 ปิฏฺเฐ.)} ปุถุชฺชนสุขํ อนริยสุขํ “น เสวิตพฺพํ น ภาเวตพฺพํ น พหุลีกาตพฺพํ ภายิตพฺพํ เอตสฺส สุขสฺสา”ติ วทามีติ กามา หิ จิตฺรา มธุรา.\csromanspacer{} \textbf{มโนรมา}ติ \textbf{มโน}ติ ยํ จิตฺตํ ฯเปฯ ตชฺชา \textbf{มโน}วิญฺญาณธาตุ, \textbf{มโน} รเมนฺติ โถเมนฺติ โตเสนฺติ ปหาเสนฺตีติ กามา หิ จิตฺรา มธุรา \textbf{มโน}รมา.}

\prose{\textbf{วิรูปรูเปน มเถนฺติ จิตฺตนฺ}ติ นานาวณฺเณหิ รูเปหิ ฯเปฯ นานาวณฺเณหิ โผฏฺฐพฺเพหิ จิตฺตํ มเถนฺติ โตเสนฺติ ปหาเสนฺตีติ วิรูปรูเปน มเถนฺติ จิตฺตํ.}

\prose{\textbf{อาทีนวํ กามคุเณสุ ทิสฺวา}ติ วุตฺตํ เหตํ ภควตา\csromandash{}โก จ ภิกฺขเว กามานํ อาทีนโว.\csromanspacer{} อิธ ภิกฺขเว กุลปุตฺโต เยน สิปฺปฏฺฐาเนน ชีวิกํ กปฺเปติ, ยทิ มุทฺทาย, ยทิ คณนาย, ยทิ สงฺขาเนน, ยทิ กสิยา, ยทิ วณิชฺชาย, ยทิ โครกฺเขน, ยทิ อิสฺสตฺเถน,\footnote{อิสฺสฏฺเฐน (สฺยา), อิสฺสตฺเตน (ก) ปสฺส ม 1. 120 ปิฏฺเฐ.} ยทิ ราชโปริเสน, ยทิ สิปฺปญฺญตเรน, สีตสฺส ปุรกฺขโต อุณฺหสฺส ปุรกฺขโต ฑํสมกสวาตาตปสรีสปสมฺผสฺเสหิ สมฺผสฺสมาโน\footnote{ริสฺสมาโน (ม 1. 120 ปิฏฺเฐ.)} ขุปฺปิปาสาย มียมาโน, อยํ ภิกฺขเว กามานํ อาทีนโว สนฺทิฏฺฐิโก ทุกฺขกฺขนฺโธ กามเหตุ กามนิทานํ กามาธิกรณํ กามานเมว เหตุ.}

\prose{ตสฺส เจ ภิกฺขเว กุลปุตฺตสฺส เอวํ อุฏฺฐหโต ฆฏโต วายมโต เต โภคา นาภินิปฺผชฺชนฺติ, โส โสจติ กิลมติ ปริเทวติ, อุรตฺตาฬิํ กนฺทติ, สมฺโมหํ อาปชฺชติ “โมฆํ วต เม อุฏฺฐานํ, อผโล วต เม วายาโม”ติ, อยมฺปิ ภิกฺขเว กามานํ อาทีนโว สนฺทิฏฺฐิโก ทุกฺขกฺขนฺโธ กามเหตุ กามนิทานํ กามาธิกรณํ กามานเมว เหตุ.}

\csromanheader{2}{19. ขคฺควิสาณสุตฺตนิทฺเทส}
\prose{\csromanfolio{263}ตสฺส เจ ภิกฺขเว กุลปุตฺตสฺส เอวํ อุฏฺฐหโต ฆฏโต วายมโต เต โภคา อภินิปฺผชฺชนฺติ, โส เตสํ โภคานํ อารกฺขาธิกรณํ ทุกฺขํ โทมนสฺสํ\footnote{ทุกฺขโทมนสฺสํ (สฺยา, ก)} ปฏิสํเวเทติ “กินฺติ เม โภเค เนว ราชาโน หเรยฺยุํ, น โจรา หเรยฺยุํ, น อคฺคิ ฑเหยฺย, น อุทกํ วเหยฺย, น อปฺปิยา ทายาทา หเรยฺยุนฺ”ติ.\csromanspacer{} ตสฺส เอวํ อารกฺขโต โคปยโต เต โภเค ราชาโน วา หรนฺติ, โจรา วา หรนฺติ, อคฺคิ วา ฑหติ, อุทกํ วา วหติ, อปฺปิยา วา ทายาทา หรนฺติ, โส โสจติ ฯเปฯ สมฺโมหํ อาปชฺชติ “ยมฺปิ เม อโหสิ, ตมฺปิ โน นตฺถี”ติ.\csromanspacer{} อยมฺปิ ภิกฺขเว กามานํ อาทีนโว สนฺทิฏฺฐิโก ทุกฺขกฺขนฺโธ กามเหตุ กามนิทานํ กามาธิกรณํ กามนเมว เหตุ.}

\prose{ปุน จปรํ ภิกฺขเว กามเหตุ กามนิทานํ กามาธิกรณํ กามานเมว เหตุ ราชาโนปิ ราชูหิ วิวทนฺติ, ขตฺติยาปิ ขตฺติเยหิ วิวทนฺติ, พฺราหฺมณาปิ พฺราหฺมเณหิ วิวทนฺติ, คหปตีปิ คหปตีหิ วิวทนฺติ, มาตาปิ ปุตฺเตน วิวทติ, ปุตฺโตปิ มาตรา วิวทติ, ปิตาปิ ปุตฺเตน วิวทติ, ปุตฺโตปิ ปิตรา วิวทติ, ภาตาปิ ภคินิยา วิวทติ, ภคินีปิ ภาตรา วิวทติ, สหาโยปิ สหาเยน วิวทติ.\csromanspacer{} เต ตตฺถ กลหวิคฺคหวิวาทาปนฺนา อญฺญมญฺญํ ปาณีหิปิ อุปกฺกมนฺติ, เลฑฺฑูหิปิ อุปกฺกมนฺติ, ทณฺเฑหิปิ อุปกฺกมนฺติ, สตฺเถหิปิ อุปกฺกมนฺติ, เต ตตฺถ มรณมฺปิ นิคจฺฉนฺติ มรณมตฺตมฺปิ ทุกฺขํ.\csromanspacer{} อยมฺปิ ภิกฺขเว กามานํ อาทีนโว สนฺทิฏฺฐิโก ทุกฺขกฺขนฺโธ กามเหตุ กามนิทานํ กามาธิกรณํ กามานเมว เหตุ.}

\prose{ปุน จปรํ ภิกฺขเว กามเหตุ กามนิทานํ กามาธิกรณํ กามานเมว เหตุ อสิจมฺมํ คเหตฺวา ธนุกลาปํ สนฺนยฺหิตฺวา อุภโตพฺยูฬฺหํ\footnote{อุภโตวิยูฬฺหํ (สฺยา) ปสฺส ม 1. 121 ปิฏฺเฐ.} สงฺคามํ ปกฺขนฺทนฺติ อุสูสุปิ ขิปฺปมาเนสุ สตฺตีสุปิ ขิปฺปมานาสุ อสีสุปิ วิชฺโชตลนฺเตสุ, เต ตตฺถ อุสูหิปิ วิชฺฌนฺติ, สตฺตีหิปิ\footnote{สตฺติยาปิ (ม 1. 121, 128 ปิฏฺเฐ.)} วิชฺฌนฺติ, อสินาปิ สีสํ ฉินฺทนฺติ.\csromanspacer{} เต ตตฺถ มรณมฺปิ นิคจฺฉนฺติ มรณมตฺตมฺปิ ทุกฺขํ.\csromanspacer{} อยมฺปิ ภิกฺขเว กามานํ อาทีนโว สนฺทิฏฺฐิโก ทุกฺขกฺขนฺโธ กามเหตุ กามนิทานํ กามาธิกรณํ กามานเมว เหตุ.}

\prose{ปุน จปรํ ภิกฺขเว กามเหตุ กามนิทานํ กามาธิกรณํ กามานเมว เหตุ อสิจมฺมํ คเหตฺวา ธนุกลาปํ สนฺนยฺหิตฺวา อทฺทาวเลปนา\footnote{อทฺธาวเลปนา (สฺยา)}}

\prose{\csromanfolio{264}อุปการิโย ปกฺขนฺทนฺติ อุสูสุปิ ขิปฺปมาเนสุ สตฺตีสุปิ ขิปฺปมานาสุ อสีสุปิ วิชฺโชตลนฺเตสุ.\csromanspacer{} เต ตตฺถ อุสูหิปิ วิชฺฌนฺติ, สตฺตีหิปิ วิชฺฌนฺติ, ฉกณกายปิ\footnote{ฉกณฏิยาปิ (สฺยา), ปกฺกฏฺฐิยา (สี-ฏฺฐ) ปสฺส ม 1. 121 ปิฏฺเฐ.} โอสิญฺจนฺติ.\csromanspacer{} อภิวคฺเคนปิ โอมทฺทนฺติ, อสินาปิ สีสํ ฉินฺทนฺติ.\csromanspacer{} เต ตตฺถ มรณมฺปิ นิคจฺฉนฺติ มรณมตฺตมฺปิ ทุกฺขํ.\csromanspacer{} อยมฺปิ ภิกฺขเว กามานํ อาทีนโว สนฺทิฏฺฐิโก ทุกฺขกฺขนฺโธ กามเหตุ กามนิทานํ กามาธิกรณํ กามานเมว เหตุ.}

\prose{ปุน จปรํ ภิกฺขเว กามเหตุ กามนิทานํ กามาธิกรณํ กามานเมว เหตุ สนฺธิมฺปิ ฉินฺทนฺติ, นิลฺโลปมฺปิ หรนฺติ, เอกาคาริกมฺปิ กโรนฺติ, ปริปนฺเถปิ ติฏฺฐนฺติ, ปรทารมฺปิ คจฺฉนฺติ.\csromanspacer{} ตเมนํ ราชาโน คเหตฺวา วิวิธา กมฺมการณา กาเรนฺติ, กสาหิปิ ตาเฬนฺติ, วตฺเตหิปิ ตาเฬนฺติ, อฑฺฒทณฺฑเกหิปิ ตาเฬนฺติ, หตฺถมฺปิ ฉินฺทนฺติ ฯเปฯ อสินาปิ สีสํ ฉินฺทนฺติ.\csromanspacer{} เต ตตฺถ มรณมฺปิ นิคจฺฉนฺติ มรณมตฺตมฺปิ ทุกฺขํ.\csromanspacer{} อยมฺปิ ภิกฺขเว กามานํ อาทีนโว สนฺทิฏฺฐิโก ทุกฺขกฺขนฺโธ กามเหตุ กามนิทานํ กามาธิกรณํ กามานเมว เหตุ.}

\prose{ปุน จปรํ ภิกฺขเว กามเหตุ กามนิทานํ กามาธิกรณํ กามานเมว เหตุ กาเยน ทุจฺจริตํ จรนฺติ, วาจาย ทุจฺจริตํ จรนฺติ, มนสา ทุจฺจริตํ จรนฺติ.\csromanspacer{} เต กาเยน ทุจฺจริตํ จริตฺวา วาจาย ทุจฺจริตํ จริตฺวา มนสา ทุจฺจริตํ จริตฺวา กายสฺส เภทา ปรํ มรณา อปายํ ทุคฺคติํ วินิปาตํ นิรยํ อุปปชฺชนฺติ.\csromanspacer{} อยมฺปิ ภิกฺขเว กามานํ อาทีนโว สมฺปรายิโก ทุกฺขกฺขนฺโธ กามเหตุ กามนิทานํ กามาธิกรณํ กามานเมว เหตุ.}

\prose{\textbf{อาทีนวํ กามคุเณสุ ทิสฺวา}ติ กามคุเณสุ อาทีนวํ ทิสฺวา ปสฺสิตฺวา ตุลยิตฺวา ตีรยิตฺวา วิภาวยิตฺวา วิภูตํ กตฺวาติ อาทีนวํ กามคุเณสุ ทิสฺวา, เอโก จเร ขคฺควิสาณกปฺโป.\csromanspacer{} เตนาห โส ปจฺเจกสมฺพุทฺโธ\csromandash{}}

\csromangathameasure{กามา หิ จิตฺรา มธุรา มโนรมา, \\ วิรูปรูเปน มเถนฺติ จิตฺตํ. \\ \textbf{อาทีนวํ กามคุเณสุ ทิสฺวา}, \\ เอโก จเร ขคฺควิสาณกปฺโปติ.}
\csromangathagroup{\csromangathastanza{กามา หิ จิตฺรา มธุรา มโนรมา, \\ วิรูปรูเปน มเถนฺติ จิตฺตํ. \\ \textbf{อาทีนวํ กามคุเณสุ ทิสฺวา}, \\ เอโก จเร ขคฺควิสาณกปฺโปติ.}}

\csromanheader{2}{19. ขคฺควิสาณสุตฺตนิทฺเทส}
\csromanhangingbegin
\hangingheaditem{137}{\csromanfolio{265}¢ตี จ คณฺโฑ จ อุปทฺทโว จ,}
\hangingbody{\textbf{โรโค จ สลฺลญฺจ ภยญฺจ เมตํ.}}
\hangingbody{\textbf{เอตํ ภยํ กามคุเณสุ ทิสฺวา,}}
\hangingbody{เอโก จเร ขคฺควิสาณกปฺโป.}
\csromanhangingend

\prose{\textbf{¢ตี จ คณฺโฑ จ อุปทฺทโว จ, โรโค จ สลฺลญฺจ ภยญฺจ เมตนฺ}ติ วุตฺตํ เหตํ ภควตา\csromandash{} “ภยนฺ”ติ ภิกฺขเว กามานเมตํ อธิวจนํ, “ทุกฺขนฺ”ติ ภิกฺขเว กามานเมตํ อธิวจนํ, “โรโค”ติ ภิกฺขเว กามานเมตํ อธิวจนํ, “คณฺโฑ”ติ ภิกฺขเว กามานเมตํ อธิวจนํ. “สลฺลนฺ”ติ ภิกฺขเว กามานเมตํ อธิวจนํ, “สงฺโค”ติ ภิกฺขเว กามานเมตํ อธิวจนํ, “ปงฺโก”ติ ภิกฺขเว กามานเมตํ อธิวจนํ, “คพฺโภ”ติ ภิกฺขเว กามานเมตํ อธิวจนํ.\csromanspacer{} กสฺมา จ ภิกฺขเว “ภยนฺ”ติ กามานเมตํ อธิวจนํ.\csromanspacer{} ยสฺมา จ กามราครตฺตายํ ภิกฺขเว ฉนฺทราควินิพทฺโธ\footnote{ฉนฺทราควินิพนฺโธ (สฺยา, ก) ปสฺส อํ 3. 114 ปิฏฺเฐ.} ทิฏฺฐธมฺมิกาปิ ภยา น ปริมุจฺจติ, สมฺปรายิกาปิ ภยา น ปริมุจฺจติ, ตสฺมา “ภยนฺ”ติ กามานเมตํ อธิวจนํ.\csromanspacer{} กสฺมา จ ภิกฺขเว “ทุกฺขนฺ”ติ. “โรโค”ติ. “คณฺโฑ”ติ. “สลฺลนฺ”ติ. “สงฺโค”ติ. “ปงฺโก”ติ. “คพฺโภ”ติ กามานเมตํ อธิวจนํ.\csromanspacer{} ยสฺมา จ กามราครตฺตายํ ภิกฺขเว ฉนฺทราควินิพทฺโธ ทิฏฺฐธมฺมิกาปิ คพฺภา น ปริมุจฺจติ, สมฺปรายิกาปิ คพฺภา น ปริมุจฺจติ, ตสฺมา “คพฺโภ”ติ กามานเมตํ อธิวจนนฺติ.}

\csromangathameasure{ภยํ ทุกฺขญฺจ โรโค จ, คณฺโฑ สลฺลญฺจ สงฺโค จ. \\ ปงฺโก คพฺโภ จ อุภยํ, เอเต กามา ปวุจฺจนฺติ. ยตฺถ สตฺโต ปุถุชฺชโน. \\ โอติณฺโณ สาตรูเปน, ปุน คพฺภาย คจฺฉติ. \\ ยโต จ ภิกฺขุ อาตาปี, สมฺปชญฺญํ น ริจฺจติ. \\ โส อิมํ ปลิปถํ ทุคฺคํ, อติกฺกมฺม ตถาวิโธ.}
\csromangathagroup{\csromangathastanza{ภยํ ทุกฺขญฺจ โรโค จ, คณฺโฑ สลฺลญฺจ สงฺโค จ. \\ ปงฺโก คพฺโภ จ อุภยํ, เอเต กามา ปวุจฺจนฺติ.\csromanspacer{} ยตฺถ สตฺโต ปุถุชฺชโน.}\csromangathastanza{โอติณฺโณ สาตรูเปน, ปุน คพฺภาย คจฺฉติ. \\ ยโต จ ภิกฺขุ อาตาปี, สมฺปชญฺญํ น ริจฺจติ\footnote{น ริญฺจติ (สฺยา, ก) สํ 2. 408 ปิฏฺเฐปิ.}. \\ โส อิมํ ปลิปถํ ทุคฺคํ, อติกฺกมฺม ตถาวิโธ.}}

\prose{ปชํ ชาติชรูเปตํ, ผนฺทมานํ อเวกฺขตีติ\csromandash{}}

\prose{¢ตี จ คณฺโฑ จ อุปทฺทโว จ,}

\prose{\textbf{โรโค จ สลฺลญฺจ ภยญฺจ เมตํ.}}

\prose{\csromanfolio{266}\textbf{เอตํ ภยํ กามคุเณสุ ทิสฺวา}ติ เอตํ ภยํ กามคุเณสุ ทิสฺวา ปสฺสิตฺวา ตุลยิตฺวา ตีรยิตฺวา วิภาวยิตฺวา วิภูตํ กตฺวาติ เอตํ ภยํ กามคุเณสุ ทิสฺวา, เอโก จเร ขคฺควิสาณกปฺโป.\csromanspacer{} เตนาห โส ปจฺเจกสมฺพุทฺโธ\csromandash{}}

\csromangathameasure{¢ตี จ คณฺโฑ จ อุปทฺทโว จ, \\ โรโค จ สลฺลญฺจ ภยญฺจ เมตํ.}
\csromangathagroup{\csromangathastanza{¢ตี จ คณฺโฑ จ อุปทฺทโว จ, \\ โรโค จ สลฺลญฺจ ภยญฺจ เมตํ.}}

\prose{\textbf{เอตํ ภยํ กามคุเณสุ ทิสฺวา},}

\prose{เอโก จเร ขคฺควิสาณกปฺโปติ.}

\csromanhangingbegin
\hangingheaditem{138}{สีตญฺจ อุณฺหญฺจ ขุทํ ปิปาสํ,}
\hangingbody{\textbf{วาตาตเป ฑํสสรีสเป1} \textbf{จ.}}
\hangingbody{\textbf{สพฺพานิเป’ตานิ อภิสมฺภวิตฺวา,}}
\hangingbody{เอโก จเร ขคฺควิสาณกปฺโป.}
\csromanhangingend

\prose{\textbf{สีตญฺจ อุณฺหญฺจ ขุทํ ปิปาสนฺ}ติ \textbf{สีตนฺ}ติ ทฺวีหิ การเณหิ สีตํ โหติ, อพฺภนฺตรธาตุปฺปโกปวเสน วา สีตํ โหติ, พหิทฺธา อุตุวเสน วา สีตํ โหติ.\csromanspacer{} \textbf{อุณฺหนฺ}ติ ทฺวีหิ การเณหิ อุณฺหํ โหติ, อพฺภนฺตรธาตุปฺปโกปวเสน วา อุณฺหํ โหติ, พหิทฺธา อุตุวเสน วา อุณฺหํ โหติ.\csromanspacer{} \textbf{ขุทา}\footnote{ขุทฺทา (สฺยา, ก)} วุจฺจติ ฉาตโก.\csromanspacer{} \textbf{ปิปาสา} วุจฺจติ อุทกปิปาสาติ สีตญฺจ อุณฺหญฺจ ขุทํ ปิปาสํ.}

\prose{\textbf{วาตาตเป ฑํสสรีสเป จา}ติ \textbf{วาตา}ติ ปุรตฺถิมา \textbf{วาตา} ปจฺฉิมา \textbf{วาตา} อุตฺตรา \textbf{วาตา} ทกฺขิณา \textbf{วาตา} สรชา \textbf{วาตา} อรชา \textbf{วาตา} สีตา \textbf{วาตา} อุณฺหา \textbf{วาตา} ปริตฺตา \textbf{วาตา} อธิมตฺตา \textbf{วาตา} เวรมฺภ\textbf{วาตา} ปกฺข\textbf{วาตา} สุปณฺณ\textbf{วาตา} ตาลปณฺณ\textbf{วาตา}\footnote{ตาลวณฺฑวาตา (ก) 109 ปิฏฺเฐ นตฺถิ ปาฐนานตฺตํ.} วิธูปน\textbf{วาตา}.\csromanspacer{} \textbf{อาตโป} วุจฺจติ สูริยสนฺตาโป.\csromanspacer{} q\textbf{อํสา} วุจฺจนฺติ ปิงฺคลมกฺขิกา.\csromanspacer{} \textbf{สรีสปา} วุจฺจนฺติ อหีติ \textbf{วาตา}ตเป ฑํสสรีสเป \textbf{จ.}}

\prose{\textbf{สพฺพานิเป’ตานิ อภิสมฺภวิตฺวา}ติ อภิภวิตฺวา อชฺโฌตฺถริตฺวา ปริยาทิยิตฺวา มทฺทิตฺวาติ สพฺพานิเป’ตานิ อภิสมฺภวิตฺวา, เอโก จเร ขคฺควิสาณกปฺโป.\csromanspacer{} เตนาห โส ปจฺเจกสมฺพุทฺโธ\csromandash{}}

\csromanheader{2}{19. ขคฺควิสาณสุตฺตนิทฺเทส}
\csromangathameasure{สีตญฺจ อุณฺหญฺจ ขุทํ ปิปาสํ, \\ วาตาตเป ฑํสสรีสเป จ. \\ สพฺพานิเป’ตานิ อภิสมฺภวิตฺวา, \\ เอโก จเร ขคฺควิสาณกปฺโปติ.}
\csromangathagroup{\csromangathastanza{\csromanfolio{267}สีตญฺจ อุณฺหญฺจ ขุทํ ปิปาสํ, \\ วาตาตเป ฑํสสรีสเป จ. \\ สพฺพานิเป’ตานิ อภิสมฺภวิตฺวา, \\ เอโก จเร ขคฺควิสาณกปฺโปติ.}}

\proseitem{139}{\textbf{นาโคว ยูถานิ วิวชฺชยิตฺวา},}

\prose{\textbf{สญฺชาตขนฺโธ ปทุมี อุฬาโร.}}

\prose{\textbf{ยถาภิรนฺตํ วิหเร}\footnote{วีหรํ (ขุ 1. 287 ปิฏฺเฐ.)} \textbf{อรญฺเญ,}}

\csromanheader{2}{\textbf{เอโก จเร ขคฺควิสาณกปฺโป. (9)}}
\prose{\textbf{นาโคว ยูถานิ วิวชฺชยิตฺวา}ติ นาโค วุจฺจติ หตฺถินาโค, ปจฺเจกสมฺพุทฺโธปิ นาโค.\csromanspacer{} กิํการณา ปจฺเจกสมฺพุทฺโธ นาโค.\csromanspacer{} อาคุํ น กโรตีติ นาโค.\csromanspacer{} น คจฺฉตีติ นาโค.\csromanspacer{} น อาคจฺฉตีติ นาโค.\csromanspacer{} กถํ โส ปจฺเจกสมฺพุทฺโธ อาคุํ น กโรตีติ นาโค.\csromanspacer{} อาคุ วุจฺจติ ปาปกา อกุสลา ธมฺมา สํกิเลสิกา โปโนภวิกา สทรา ทุกฺขวิปากา อายติํ ชาติชรามรณิยา.}

\csromangathameasure{อาคุํ น กโรติ กิญฺจิ โลเก, \\ สพฺพสํโยเค วิสชฺช พนฺธนานิ. \\ สพฺพตฺถ น สชฺชตี วิมุตฺโต, \\ นาโค ตาทิ ปวุจฺจเต ตถตฺตา.}
\csromangathagroup{\csromangathastanza{อาคุํ น กโรติ กิญฺจิ โลเก, \\ สพฺพสํโยเค วิสชฺช พนฺธนานิ. \\ สพฺพตฺถ น สชฺชตี วิมุตฺโต, \\ นาโค ตาทิ ปวุจฺจเต ตถตฺตา.}}

\prose{เอวํ โส ปจฺเจกสมฺพุทฺโธ อาคุํ น กโรตีติ นาโค.}

\prose{กถํ โส ปจฺเจกสมฺพุทฺโธ น คจฺฉตีติ นาโค, โส ปจฺเจกสมฺพุทฺโธ น ฉนฺทาคติํ คจฺฉติ, น โทสาคติํ คจฺฉติ, น โมหาคติํ คจฺฉติ, น ภยาคติํ คจฺฉติ, น ราควเสน คจฺฉติ, น โทสวเสน คจฺฉติ, น โมหวเสน คจฺฉติ, น มานวเสน คจฺฉติ, น ทิฏฺฐิวเสน คจฺฉติ, น อุทฺธจฺจวเสน คจฺฉติ, น วิจิกิจฺฉาวเสน คจฺฉติ, น อนุสยวเสน คจฺฉติ, น วคฺเคหิ กปฺเปหิ ยายติ นียติ\footnote{นิยฺยติ (สฺยา, ก)} วุยฺหติ สํหรียติ.\csromanspacer{} เอวํ โส ปจฺเจกสมฺพุทฺโธ น คจฺฉตีติ นาโค.}

\prose{\csromanfolio{268}กถํ โส ปจฺเจกสมฺพุทฺโธ น อาคจฺฉตีติ นาโค, โสตาปตฺติ มคฺเคน เย กิเลสา ปหีนา, เต กิเลเส น ปุเนติ น ปจฺเจติ น ปจฺจาคจฺฉติ.\csromanspacer{} สกทาคามิมคฺเคน ฯเปฯ อนาคามิมคฺเคน ฯเปฯ อรหตฺตมคฺเคน เย กิเลสา ปหีนา, เต กิเลเส น ปุเนติ น ปจฺเจติ น ปจฺจาคจฺฉติ.\csromanspacer{} เอวํ โส ปจฺเจกสมฺพุทฺโธ น อาคจฺฉตีติ นาโค.}

\prose{\textbf{นาโคว ยูถานิ วิวชฺชยิตฺวา}ติ ยถา โส หตฺถินาโค ยูถานิ วิวชฺเชตฺวา ปริวชฺเชตฺวา อภินิวชฺเชตฺวา เอโกว อรญฺญวนมชฺโฌคาเหตฺวา\footnote{อรญฺเญ วนมชฺฌสฺส อชฺโฌคาเหตฺวา (สฺยา)} จรติ วิหรติ อิริยติ วตฺเตติ ปาเลติ ยเปติ ยาเปติ.\csromanspacer{} ปจฺเจกสมฺพุทฺโธปิ คณํ วิวชฺเชตฺวา ปริวชฺเชตฺวา อภิวชฺเชตฺวา เอโก\footnote{เอโก จเร ขคฺควิสาณกปฺโป (สฺยา)} อรญฺญวนปตฺถานิ ปนฺตานิ เสนาสนานิ ปฏิเสวติ อปฺปสทฺทานิ อปฺปนิคฺโฆสานิ วิชนวาตานิ มนุสฺสราหสฺเสยฺยกานิ ปฏิสลฺลานสารุปฺปานิ.\csromanspacer{} โส เอโก คจฺฉติ, เอโก ติฏฺฐติ, เอโก นิสีทติ, เอโก เสยฺยํ กปฺเปติ, เอโก คามํ ปิณฺฑาย ปวิสติ, เอโก ปฏิกฺกมติ, เอโก รโห นิสีทติ, เอโก จงฺกมํ อธิฏฺฐาติ, เอโก จรติ วิหรติ อิริยติ วตฺเตติ ปาเลติ ยเปติ ยาเปตีติ นาโคว ยูถานิ วิวชฺชยิตฺวา.}

\prose{\textbf{สญฺชาตขนฺโธ ปทุมี อุฬาโร}ติ ยถา โส หตฺถินาโค สญฺชาตกฺขนฺโธ สตฺตรตโน วา โหติ อฏฺฐรตโน วา.\csromanspacer{} ปจฺเจกสมฺพุทฺโธปิ สญฺชาตกฺขนฺโธ อเสกฺเขน สีลกฺขนฺเธน, อเสกฺเขน สมาธิกฺขนฺเธน, อเสกฺเขน ปญฺญากฺขนฺเธน, อเสกฺเขน วิมุตฺติกฺขนฺเธน, อเสกฺเขน วิมุตฺติญาณทสฺสนกฺขนฺเธน.\csromanspacer{} ยถา โส หตฺถินาโค ปทุมี, ปจฺเจกสมฺพุทฺโธปิ สตฺตหิ โพชฺฌงฺคปุปฺเผหิ ปทุมี สติสมฺโพชฺฌงฺคปุปฺเผน ธมฺมวิจยสมฺโพชฺฌงฺคปุปฺเผน วีริยสมฺโพชฺฌงฺคปุปฺเผน ปีติสมฺโพชฺฌงฺคปุปฺเผน ปสฺสทฺธิสมฺโพชฺฌงฺคปุปฺเผน สมาธิสมฺโพชฺฌงฺคปุปฺเผน อุเปกฺขาสมฺโพชฺฌงฺคปุปฺเผน.\csromanspacer{} ยถา โส หตฺถินาโค อุฬาโร ถาเมน พเลน ชเวน สูเรน, ปจฺเจกสมฺพุทฺโธปิ อุฬาโร สีเลน สมาธินา ปญฺญาย วิมุตฺติยา วิมุตฺติญาณทสฺสเนนาติ สญฺชาตขนฺโธ ปทุมี อุฬาโร.}

\prose{\textbf{ยถาภิรนฺตํ วิหเร อรญฺเญ}ติ ยถา โส หตฺถินาโค ยถาภิรนฺตํ อรญฺเญ วิหรติ.\csromanspacer{} ปจฺเจกสมฺพุทฺโธปิ ยถาภิรนฺตํ อรญฺเญ วิหรติ.}

\csromanheader{2}{19. ขคฺควิสาณสุตฺตนิทฺเทส}
\prose{\csromanfolio{269}ปฐเมนปิ ฌาเนน ยถาภิรนฺตํ อรญฺเญ วิหรติ.\csromanspacer{} ทุติเยนปิ ฌาเนน.\csromanspacer{} ตติเยนปิ ฌาเนน.\csromanspacer{} จตุตฺเถนปิ ฌาเนน ยถาภิรนฺตํ อรญฺเญ วิหรติ.\csromanspacer{} เมตฺตายปิ เจโตวิมุตฺติยา ยถาภิรนฺตํ อรญฺเญ วิหรติ.\csromanspacer{} กรุณายปิ เจโตวิมุตฺติยา.\csromanspacer{} มุทิตายปิ เจโตวิมุตฺติยา.\csromanspacer{} อุเปกฺขายปิ เจโตวิมุตฺติยา ยถาภิรนฺตํ อรญฺเญ วิหรติ.\csromanspacer{} อากาสานญฺจายตนมาปตฺติยาปิ ยถาภิรนฺตํ อรญฺเญ วิหรติ.\csromanspacer{} วิญฺญาณญฺจายตนสมาปตฺติยาปิ.\csromanspacer{} อากิญฺจญฺญายตนสมาปตฺติยาปิ.\csromanspacer{} เนวสญฺญานาสญฺญายตนสมาปตฺติยาปิ.\csromanspacer{} นิโรธสมาปตฺติยาปิ.\csromanspacer{} ผลสมาปตฺติยาปิ ยถาภิรนฺตํ อรญฺเญ วิหรตีติ ยถาภิรนฺตํ วิหเร อรญฺเญ, เอโก จเร ขคฺควิสาณกปฺโป.\csromanspacer{} เตนาห โส ปจฺเจกสมฺพุทฺโธ\csromandash{}}

\csromangathameasure{นาโคว ยูถานิ วิวชฺชยิตฺวา, \\ สญฺชาตขนฺโธ ปทุมี อุฬาโร. \\ ยถาภิรนฺตํ วิหเร อรญฺเญ, \\ เอโก จเร ขคฺควิสาณกปฺโปติ.}
\csromangathagroup{\csromangathastanza{นาโคว ยูถานิ วิวชฺชยิตฺวา, \\ สญฺชาตขนฺโธ ปทุมี อุฬาโร. \\ ยถาภิรนฺตํ วิหเร อรญฺเญ, \\ เอโก จเร ขคฺควิสาณกปฺโปติ.}}

\proseitem{140}{อฏฺฐาน’ตํ สงฺคณิการตสฺส,}

\prose{\textbf{ยํ ผสฺสเย}\footnote{ผุสฺสเย (สฺยา, ก)} \textbf{สามยิกํ}\footnote{อาสามายิกํ (ก)} \textbf{วิมุตฺติํ.}}

\prose{\textbf{อาทิจฺจพนฺธุสฺส วโจ นิสมฺม,}}

\csromanheader{2}{\textbf{เอโก จเร ขคฺควิสาณกปฺโป. (10)}}
\prose{\textbf{อฏฺฐาน’ตํ สงฺคณิการตสฺส, ยํ ผสฺสเย สามยิกํ วิมุตฺตินฺ}ติ วุตฺตํ เหตํ ภควตา\csromandash{}ยาวตานนฺท\footnote{โส วตานนฺท (ม 3. 152 ปิฏฺเฐ.)} ภิกฺขุ สงฺคณิการาโม สงฺคณิกรโต สงฺคณิการามตํ อนุยุตฺโต คณาราโม คณรโต คณสมฺมุทิโต (คณารามตํ อนุยุตฺโต)\footnote{( ) นตฺถิ ม 3. 152 ปิฏฺเฐ.}, ยํ ตํ เนกฺขมฺมสุขํ ปวิเวกสุขํ อุปสมสุขํ สมฺโพธิสุขํ, ตสฺส สุขสฺส นิกามลาภี ภวิสฺสติ อกิจฺฉลาภี อกสิรลาภีติ เนตํ ฐานํ วิชฺชติ.\csromanspacer{} โย จ โข โส อานนฺท ภิกฺขุ เอโก คณสฺมา วูปกฏฺโฐ วิหรติ, ตสฺเสตํ ภิกฺขุโน ปาฏิกงฺขํ.\csromanspacer{} ยํ ตํ เนกฺขมฺมสุขํ ปวิเวกสุขํ อุปสมสุขํ สมฺโพธิสุขํ, ตสฺส สุขสฺส นิกามลาภี ภวิสฺสติ อกิจฺฉลาภี อกสิรลาภีติ ฐานเมตํ วิชฺชติ.\csromanspacer{} ยาวตานนฺท ภิกฺขุ สงฺคณิการาโม สงฺคณิกรโต}

\prose{\csromanfolio{270}สงฺคณิการามตํ อนุยุตฺโต คณาราโม คณรโต คณสมฺมุทิโต (คณารามตํ อนุยุตฺโต,) สามายิกํ\footnote{สามยิกํ (สฺยา, ก)} วา กนฺตํ เจโตวิมุตฺติํ อุปสมฺปชฺช วิหริสฺสติ, อสามายิกํ วา อกุปฺปนฺติ เนตํ ฐานํ วิชฺชติ.\csromanspacer{} โย จ โข โส อานนฺท ภิกฺขุ เอโก คณสฺมา วูปกฏฺโฐ วิหรติ, ตสฺเสตํ ภิกฺขุโน ปาฏิกงฺขํ สามายิกํ วา กนฺตํ เจโตวิมุตฺติํ อุปสมฺปชฺช วิหริสฺสติ, อสามายิกํ วา อกุปฺปนฺติ ฐานเมตํ วิชฺชตีติ อฏฺฐาน’ตํ สงฺคณิการตสฺส, ยํ ผสฺสเย สามยิกํ วิมุตฺติํ.}

\prose{\textbf{อาทิจฺจพนฺธุสฺส วโจ นิสมฺมา}ติ อาทิจฺโจ วุจฺจติ สูริโย, โส โคตโม โคตฺเตน.\csromanspacer{} ปจฺเจกสมฺพุทฺโธปิ โคตโม โคตฺเตน.\csromanspacer{} โส ปจฺเจกสมฺพุทฺโธ สูริยสฺส โคตฺตญาตโก โคตฺตพนฺธุ, ตสฺมา ปจฺเจกสมฺพุทฺโธ อาทิจฺจพนฺธุ.\csromanspacer{} \textbf{อาทิจฺจพนฺธุสฺส วโจ นิสมฺมา}ติ อาทิจฺจพนฺธุสฺส วจนํ พฺยปฺปถํ เทสนํ อนุสาสนํ อนุสิฏฺฐํ สุตฺวา สุณิตฺวา อุคฺคเหตฺวา อุปธารยิตฺวา อุปลกฺขยิตฺวาติ อาทิจฺจพนฺธุสฺส วโจ นิสมฺม, เอโก จเร ขคฺควิสาณกปฺโป.\csromanspacer{} เตนาห โส ปจฺเจกสมฺพุทฺโธ\csromandash{}}

\csromangathameasure{อฏฺฐาน’ตํ สงฺคณิการตสฺส, \\ ยํ ผสฺสเย สามยิกํ วิมุตฺติํ. \\ อาทิจฺจพนฺธุสฺส วโจ นิสมฺม, \\ เอโก จเร ขคฺควิสาณกปฺโปติ.}
\csromangathagroup{\csromangathastanza{อฏฺฐาน’ตํ สงฺคณิการตสฺส, \\ ยํ ผสฺสเย สามยิกํ วิมุตฺติํ. \\ อาทิจฺจพนฺธุสฺส วโจ นิสมฺม, \\ เอโก จเร ขคฺควิสาณกปฺโปติ.}}

\vspace{\tipitakaparskip}
\csromancenter{ทุติโย วคฺโค.}
\chapterhead{\textbf{ตติยวคฺค}}
\csromanhangingbegin
\hangingheaditem{141}{\textbf{ทิฏฺฐีวิสูกานิ อุปาติวตฺโต},}
\hangingbody{\textbf{ปตฺโต นิยามํ ปฏิลทฺธมคฺโค.}}
\hangingbody{\textbf{อุปฺปนฺนญาโณมฺหิ อนญฺญเนยฺโย,}}
\hangingbody{เอโก จเร ขคฺควิสาณกปฺโป.}
\csromanhangingend

\prose{\textbf{ทิฏฺฐีวิสูกานิ อุปาติวตฺโต}ติ ทิฏฺฐิวิสูกานิ วุจฺจนฺติ วีสติ วตฺถุกา สกฺกายทิฏฺฐี.\csromanspacer{} อิธ อสฺสุตวา ปุถุชฺชโน อริยานํ อทสฺสาวี อริยธมฺมสฺส}

\csromanheader{2}{19. ขคฺควิสาณสุตฺตนิทฺเทส}
\prose{\csromanfolio{271}อโกวิโท อริยธมฺเม อวินีโต สปฺปุริสานํ อทสฺสาวี สปฺปุริสธมฺมสฺส อโกวิโท สปฺปุริสธมฺเม อวินีโต รูปํ อตฺตโต สมนุปสฺสติ รูปวนฺตํ วา อตฺตานํ, อตฺตนิ วา รูปํ, รูปสฺมิํ วา อตฺตานํ.\csromanspacer{} เวทนํ.\csromanspacer{} สญฺญํ.\csromanspacer{} สงฺขาเร.\csromanspacer{} วิญฺญาณํ อตฺตโต สมนุปสฺสติ, วิญฺญาณวนฺตํ วา อตฺตานํ, อตฺตนิ วา วิญฺญาณํ, วิญฺญาณสฺมิํ วา อตฺตนํ.\csromanspacer{} ยา เอวรูปา ทิฏฺฐิ ทิฏฺฐิคตํ ทิฏฺฐิคหนํ ทิฏฺฐิกนฺตาโร ทิฏฺฐิวิสูกายิกํ ทิฏฺฐิวิปฺผนฺทิตํ ทิฏฺฐิสํโยชนํ คาโห ปฏิคฺคาโห อภินิเวโส ปรามาโส กุมฺมคฺโค มิจฺฉาปโถ มิจฺฉตฺตํ ติตฺถายตนํ วิปริเยสคฺคาโห วิปรีตคฺคาโห วิปลฺลาสคฺคาโห มิจฺฉาคาโห อยาถาวกสฺมิํ “ยาถาวกนฺ”ติ คาโห, ยาวตา ทฺวาสฏฺฐิ ทิฏฺฐิคตานิ, อิมานิ ทิฏฺฐิวิสูกานิ.\csromanspacer{} \textbf{ทิฏฺฐีวิสูกานิ อุปาติวตฺโต}ติ ทิฏฺฐิวิสูกานิ อุปาติวตฺโต อติกฺกนฺโต สมติกฺกนฺโต วีติวตฺโตติ ทิฏฺฐีวิสูกานิ อุปาติวตฺโต.}

\prose{\textbf{ปตฺโต นิยามํ ปฏิลทฺธมคฺโค}ติ นิยามา วุจฺจนฺติ จตฺตาโร มคฺคา.\csromanspacer{} อริโย อฏฺฐงฺคิโก มคฺโค.\csromanspacer{} เสยฺยถิทํ, สมฺมาทิฏฺฐิสมฺมาสงฺกปฺโป สมฺมาวาจา สมฺมากมฺมนฺโต สมฺมา-อาชีโว สมฺมาวายาโม สมฺมาสติ สมฺมาสมาธิ.\csromanspacer{} จตูหิ อริยมคฺเคหิ สมนฺนาคโต นิยามํ ปตฺโต สมฺปตฺโต อธิคโต ผสฺสิโต สจฺฉิกโตติ ปตฺโต นิยามํ.\csromanspacer{} \textbf{ปฏิลทฺธมคฺโค}ติ ลทฺธมคฺโค ปฏิลทฺธมคฺโค อธิคตมคฺโค ผสฺสิตมคฺโค สจฺฉิกตมคฺโคติ ปตฺโต นิยามํ ปฏิลทฺธมคฺโค.}

\prose{\textbf{อุปฺปนฺนญาโณมฺหิ อนญฺญเนยฺโย}ติ ตสฺส ปจฺเจกสมฺพุทฺธสฺส ญาณํ อุปฺปนฺนํ สมุปฺปนฺนํ นิพฺพตฺตํ อภินิพฺพตฺตํ ปาตุภูตํ “สพฺเพ สงฺขารา อนิจฺจา”ติ ญาณํ อุปฺปนฺนํ สมุปฺปนฺนํ นิพฺพตฺตํ อภินิพฺพตฺตํ ปาตุภูตํ. “สพฺเพ สงฺขารา ทุกฺขา”ติ ฯเปฯ. “สพฺเพ ธมฺมา อนตฺตา”ติ ฯเปฯ. “ยํ กิญฺจิ สมุทยธมฺมํ, สพฺพํ ตํ นิโรธธมฺมนฺ”ติ ญาณํ อุปฺปนฺนํ สมุปฺปนฺนํ นิพฺพตฺตํ อภินิพฺพตฺตํ ปาตุภูตนฺติ อุปฺปนฺนญาโณมฺหิ.\csromanspacer{} \textbf{อนญฺญเนยฺโย}ติ โส ปจฺเจกสมฺพุทฺโธ น ปรเนยฺโย น ปรปฺปตฺติโย น ปรปฺปจฺจโย น ปรปฏิพทฺธคู, ยถาภูตํ\footnote{ตํ (ก)} ชานาติ ปสฺสติ อสมฺมูโฬฺห สมฺปชาโน ปฏิสฺสโต. “สพฺเพ สงฺขารา อนิจฺจา”ติ น ปรเนยฺโย น ปรปฺปตฺติโย น ปรปฺปจฺจโย น ปรปฏิพทฺธคู, ยถาภูตํ1 ชานาติ ปสฺสติ อสมฺมูโฬฺห สมฺปชาโน ปฏิสฺสโต. “สพฺเพ สงฺขารา ทุกฺขา”ติ ฯเปฯ. “สพฺเพ ธมฺมา อนตฺตา”ติ ฯเปฯ. “ยํ กิญฺจิ สมุทยธมฺมํ, สพฺพํ ตํ นิโรธธมฺมนฺ”ติ น ปรเนยฺโย}

\prose{\csromanfolio{272}น ปรปฺปตฺติโย น ปรปฺปจฺจโย น ปรปฏิพทฺธคู, ยถาภูตํ\footnote{นิราสโย (สี-ฏฺฐ) ขุ 1. 288 ปิฏฺเฐปิ.} ชานาติ ปสฺสติ อสมฺมูโฬฺห สมฺปชาโน ปฏิสฺสโตติ อุปฺปนฺนญาโณมฺหิ อนญฺญเนยฺโย, เอโก จเร ขคฺควิสาณกปฺโป.\csromanspacer{} เตนาห โส ปจฺเจกสมฺพุทฺโธ\csromandash{}}

\csromangathameasure{ทีฏฺฐีวิสูกานิ อุปาติวตฺโต, \\ ปตฺโต นิยามํ ปฏิลทฺธมคฺโค. \\ อุปฺปนฺนญาโณมฺหิ อนญฺญเนยฺโย, \\ เอโก จเร ขคฺควิสาณกปฺโปติ.}
\csromangathagroup{\csromangathastanza{ทีฏฺฐีวิสูกานิ อุปาติวตฺโต, \\ ปตฺโต นิยามํ ปฏิลทฺธมคฺโค. \\ อุปฺปนฺนญาโณมฺหิ อนญฺญเนยฺโย, \\ เอโก จเร ขคฺควิสาณกปฺโปติ.}}

\csromanhangingbegin
\hangingheaditem{142}{\textbf{นิลฺโลลุโป นิกฺกุโห นิปฺปิปาโส},}
\hangingbody{\textbf{นิมฺมกฺโข นิทฺธนฺตกสาวโมโห.}}
\hangingbody{\textbf{นิราสโส1} \textbf{สพฺพโลเก ภวิตฺวา,}}
\hangingbody{เอโก จเร ขคฺควิสาณกปฺโป.}
\csromanhangingend

\prose{\textbf{นิลฺโลลุโป นิกฺกุโห นิปฺปิปาโส}ติ \textbf{โลลุปฺปํ} วุจฺจติ ตณฺหา.\csromanspacer{} โย ราโค สาราโค ฯเปฯ อภิชฺฌา โลโภ อกุสลมูลํ.\csromanspacer{} สา โลลุปฺปา ตณฺหา ตสฺส ปจฺเจกสมฺพุทฺธสฺส ปหีนา อุจฺฉินฺนมูลา ตาลาวตฺถุกตา อนภาวํกตา อายติํ อนุปฺปาทธมฺมา.\csromanspacer{} ตสฺมา ปจฺเจกสมฺพุทฺโธ นิลฺโลลุโป.}

\prose{\textbf{นิกฺกุโห}ติ ตีณิ กุหนวตฺถูนิ ปจฺจยปฏิเสวนสงฺขาตํ กุหนวตฺถุ อิริยาปถสงฺขาตํ กุหนวตฺถุ สามนฺตชปฺปนสงฺขาตํ กุหนวตฺถุ.}

\prose{กตมํ ปจฺจยปฏิเสวนสงฺขาตํ กุหนวตฺถุ, อิธ คหปติกา ภิกฺขุํ\footnote{ภิกฺขู (ก) ขุ 7. 172 ปิฏฺเฐ.} นิมนฺเตนฺติ จีวรปิณฺฑปาตเสนาสนคิลานปจฺจยเภสชฺช- ปริกฺขาเรหิ, โส ปาปิจฺโฉ อิจฺฉาปกโต อตฺถิโก จีวรปิณฺฑปาตเสนาสนคิลานปจฺจยเภสชฺชปริกฺขารานํ ภิยฺโยกมฺยตํ อุปาทาย จีวรํ ปจฺจกฺขาติ, ปิณฺฑปาตํ ปจฺจกฺขาติ, เสนาสนํ ปจฺจกฺขาติ, คิลานปจฺจยเภสชฺชปริกฺขารํ ปจฺจกฺขาติ, โส เอวมาห “กิํ สมณสฺส มหคฺเฆน จีวเรน, เอตํ สารุปฺปํ, ยํ สมโณ สุสานา วา สงฺการกูฏา วา ปาปณิกา วา นนฺตกานิ อุจฺจินิตฺวา สงฺฆาฏิกํ กริตฺวา ธาเรยฺย.\csromanspacer{} กิํ สมณสฺส มหคฺเฆน ปิณฺฑปาเตน, เอตํ สารุปฺปํ ยํ สมโณ อุญฺฉาจริยาย ปิณฺฑิยาโลเปน ชีวิกํ กปฺเปยฺย.\csromanspacer{} กิํ สมณสฺส มหคฺเฆน เสนาสเนน, เอตํ สารุปฺปํ, ยํ สมโณ}

\csromanheader{2}{19. ขคฺควิสาณสุตฺตนิทฺเทส}
\prose{\csromanfolio{273}รุกฺขมูลิโก วา อสฺส โสสานิโก วา อพฺโภกาสิโก วา, กิํ สมณสฺส มหคฺเฆน คิลานปจฺจยเภสชฺชปริกฺขาเรน, เอตํ สารุปฺปํ, ยํ สมโณ ปูติมุตฺเตน วา หริตกีขณฺเฑน วา โอสธํ กเรยฺยา”ติ.\csromanspacer{} ตทุปาทาย ลูขํ จีวรํ ธาเรติ, ลูขํ ปิณฺฑปาตํ ปริภุญฺชติ, ลูขํ เสนาสนํ ปฏิเสวติ, ลูขํ คิลานปจฺจยเภสชฺชปริกฺขารํ ปฏิเสวติ.\csromanspacer{} ตเมนํ คหปติกา เอวํ ชานนฺติ “อยํ สมโณ อปฺปิจฺโฉ สนฺตุฏฺโฐ ปวิวิตฺโต อสํสฏฺโฐ อารทฺธวีริโย ธุตวาโท”ติ ภิยฺโย ภิยฺโย นิมนฺเตนฺติ จีวรปิณฺฑปาตเสนาสนคิลาน ปจฺจย เภสชฺช ปริกฺขาเรหิ, โส เอวมาห “ติณฺณํ สมฺมุขีภาวา สทฺโธ กุลปุตฺโต พหุํ ปุญฺญํ ปสวติ, สทฺธาย สมฺมุขีภาวา สทฺโธ กุลปุตฺโต พหุํ ปุญฺญํ ปสวติ, เทยฺยธมฺมสฺส สมฺมุขีภาวา สทฺโธ กุลปุตฺโต พหุํ ปุญฺญํ ปสวติ, ทกฺขิเณยฺยานํ สมฺมุขีภาวา สทฺโธ กุลปุตฺโต พหุํ ปุญฺญํ ปสวติ.\csromanspacer{} ตุมฺหากญฺเจวายํ สทฺธา อตฺถิ, เทยฺยธมฺโม จ สํวิชฺชติ, อหญฺจ ปฏิคฺคาหโก.\csromanspacer{} สเจหํ น ปฏิคฺคเหสฺสามิ, เอวํ ตุเมฺห ปุญฺเญน ปริพาหิรา ภวิสฺสถ\footnote{ภวิสฺสนฺติ (ขุ 7. 173 ปิฏฺเฐ)}, น มยฺหํ อิมินา อตฺโถ, อปิ จ ตุมฺหากํเยว อนุกมฺปาย ปฏิคฺคณฺหามี”ติ ตทุปาทาย พหุมฺปิ จีวรํ ปฏิคฺคณฺหาติ, พหุมฺปิ ปิณฺฑปาตํ ปฏิคฺคณฺหาติ, พหุมฺปิ เสนาสนํ ปฏิคฺคณฺหาติ, พหุมฺปิ คิลานปจฺจยเภสชฺชปริกฺขารํ ปฏิคฺคณฺหาติ.\csromanspacer{} ยา เอวรูปา ภากุฏิกา ภากุฏิยํ กุหนา กุหายนา กุหิตตฺตํ, อิทํ ปจฺจยปฏิเสวนสงฺขาตํ กุหนวตฺถุ.}

\prose{กตมํ อิริยาปถสงฺขาตํ กุหนวตฺถุ, อิเธกจฺโจ ปาปิจฺโฉ อิจฺฉาปกโต สมฺภาวนาธิปฺปาโย “เอวํ มํ ชโน สมฺภาเวสฺสตี”ติ คมนํ สณฺฐเปติ, ฐานํ สณฺฐเปติ, นิสชฺชํ\footnote{นิสชฺชนํ (ก)} สณฺฐเปติ, สยนํ สณฺฐเปติ.\csromanspacer{} ปณิธาย คจฺฉติ, ปณิธาย ติฏฺฐติ, ปณิธาย นิสีทติ, ปณิธาย เสยฺยํ กปฺเปติ.\csromanspacer{} สมาหิโต วิย คจฺฉติ, สมาหิโต วิย ติฏฺฐติ, สมาหิโต วิย นิสีทติ, สมาหิโต วิย เสยฺยํ กปฺเปติ.\csromanspacer{} อาปาถกชฺฌายีว\footnote{อาปาตกฌายี จ (ก) 4. กุหายิตตฺตํ (สฺยา, ก), วิสุทฺธิมคฺคฏฺฐกถา โอโลเกตพฺพา.} โหติ.\csromanspacer{} ยา เอวรูปา อิริยาปถสฺส อาฐปนา ฐปนา สณฺฐปนา ภากุฏิกา ภากุฏิยํ กุหนา กุหายนา กุหิตตฺตํ4.\csromanspacer{} อิทํ อิริยาปถสงฺขาตํ กุหนวตฺถุ.}

\prose{\csromanfolio{274}กตมํ สามนฺตชปฺปนสงฺขาตํ กุหนวตฺถุ, อิเธกจฺโจ ปาปิจฺโฉ อิจฺฉาปกโต สมฺภาวนาธิปฺปาโย “เอวํ มํ ชโน สมฺภาเวสฺสตี”ติ อริยธมฺเม สนฺนิสฺสิตวาจํ ภาสติ, “โย เอวรูปํ จีวรํ ธาเรติ, โส สมโณ มเหสกฺโข”ติ ภณติ, “โย เอวรูปํ ปตฺตํ ธาเรติ, โลหถาลกํ ธาเรติ, ธมฺมกรณํ ธาเรติ.\csromanspacer{} ปริสฺสาวนํ ธาเรติ, กุญฺจิกํ ธาเรติ, อุปาหนํ ธาเรติ, กายพนฺธนํ ธาเรติ, อาโยคํ\footnote{อาโยคพนฺธนํ (สฺยา, ก) ปสฺส ขุ 7. 173 ปิฏฺเฐ.} ธาเรติ, โส สมโณ มเหสกฺโข”ติ ภณติ. “ยสฺส เอวรูโป อุปชฺฌาโย, โส สมโณ มเหสกฺโข”ติ ภณติ. “ยสฺส เอวรูโป อาจริโย.\csromanspacer{} เอวรูปา สมานุปชฺฌายกา.\csromanspacer{} สมานาจริยกา.\csromanspacer{} มิตฺตา.\csromanspacer{} สนฺทิฏฺฐา.\csromanspacer{} สมฺภตฺตา.\csromanspacer{} สหายา, โส สมโณ มเหสกฺโข”ติ ภณติ. “โย เอวรูเป วิหาเร วสติ.\csromanspacer{} อฑฺฒโยเค วสติ.\csromanspacer{} ปาสาเท วสติ.\csromanspacer{} หมฺมิเย วสติ.\csromanspacer{} คุหายํ วสติ.\csromanspacer{} เลเณ วสติ.\csromanspacer{} กุฏิยํ วสติ.\csromanspacer{} กูฏาคาเร วสติ.\csromanspacer{} อฏฺเฏ วสติ.\csromanspacer{} มาเฬ วสติ.\csromanspacer{} อุทฺทณฺเฑ\footnote{อุฏฺฏณฺเฑ (ก)} วสติ.\csromanspacer{} อุปฏฺฐานสาลายํ วสติ.\csromanspacer{} มณฺฑเป วสติ.\csromanspacer{} รุกฺขมูเล วสติ.\csromanspacer{} โส สมโณ มเหสกฺโข”ติ ภณติ.}

\prose{อถ วา โกรชิกโกรชิโก ภากุฏิกภากุฏิโก กุหกกุหโก ลปกลปโก มุขสมฺภาวิโก “อยํ สมโณ อิมาสํ เอวรูปานํ สนฺตานํ วิหารสมาปตฺตีนํ ลาภี”ติ ตาทิสํ คมฺภีรํ คูฬฺหํ นิปุณํ ปฏิจฺฉนฺนํ โลกุตฺตรํ สุญฺญตาปฏิสญฺญุตฺตํ กถํ กเถติ.\csromanspacer{} ยา เอวรูปา ภากุฏิกา ภากุฏิยํ กุหนา กุหายนา กุหิตตฺตํ, อิทํ สามนฺตชปฺปนสงฺขาตํ กุหนวตฺถุ.\csromanspacer{} ตสฺส ปจฺเจกสมฺพุทฺธสฺส อิมานิ ตีณิ กุหนวตฺถุนิ ปหีนานิ สมุจฺฉินฺนานิ วูปสนฺตานิ ปฏิปฺปสฺสทฺธานิ อภพฺพุปฺปตฺติกานิ ญาณคฺคินา ทฑฺฒานิ.\csromanspacer{} ตสฺมา โส ปจฺเจกสมฺพุทฺโธ นิกฺกุโห.}

\prose{\textbf{นิปฺปิปาโส}ติ \textbf{ปิปาสา} วุจฺจติ ตณฺหา.\csromanspacer{} โย ราโค สาราโค ฯเปฯ อภิชฺฌา โลโภ อกุสลมูลํ.\csromanspacer{} สา \textbf{ปิปาสา} ตณฺหา ตสฺส ปจฺเจกสมฺพุทฺธสฺส ปหีนา อุจฺฉินฺนมูลา ตาลาวตฺถุกตา อนภาวํกตา อายติํ อนุปฺปาทธมฺมา.\csromanspacer{} ตสฺมา ปจฺเจกสมฺพุทฺโธ นิปฺปิปาโสติ นิลฺโลลุโป นิกฺกุโห นิปฺปิปาโส.}

\csromanheader{2}{19. ขคฺควิสาณสุตฺตนิทฺเทส}
\prose{\csromanfolio{275}\textbf{นิมฺมกฺโข นิทฺธนฺตกสาวโมโห}ติ \textbf{มกฺโข}ติ โย \textbf{มกฺโข} มกฺขายนา\footnote{มกฺขิยนา (ก) ปสฺส อภิ 2. 371 ปิฏฺเฐ.} มกฺขายิตตฺตํ นิฏฺฐุริยํ นิฏฺฐุริยกมฺมํ.\csromanspacer{} \textbf{กสาโว}ติ ราโค กสาโว, โทโส กสาโว, โมโห กสาโว, โกโธ.\csromanspacer{} อุปนาโห.\csromanspacer{} มกฺโข.\csromanspacer{} ปฬาโส ฯเปฯ สพฺพากุสลาภิสงฺขารา กสาวา.\csromanspacer{} \textbf{โมโห}ติ ทุกฺเข อญฺญาณํ, ทุกฺขสมุทเย อญฺญาณํ, ทุกฺขนิโรเธ อญฺญาณํ, ทุกฺขนิโรธคามินิยา ปฏิปทาย อญฺญาณํ, ปุพฺพนฺเต อญฺญาณํ, อปรนฺเต อญฺญาณํ, ปุพฺพนฺตาปรนฺเต อญฺญาณํ.\csromanspacer{} อิทปฺปจฺจยตาปฏิจฺจสมุปฺปนฺเนยุ ธมฺเมสุ อญฺญาณํ.\csromanspacer{} ยํ เอวรูปํ อญฺญาณํ อทสฺสนํ อนภิสมโย อนนุโพโธ อปฺปฏิเวโธ อสํคาหนา อปริโยคาหนา อสมเปกฺขนา อปจฺจเวกฺขณา อปจฺจกฺขกมฺมํ ทุมฺเมชฺฌํ พาลฺยํ อสมฺปชญฺญํ โมโห ปโมโห สมฺโมโห อวิชฺชา อวิชฺโชโฆ อวิชฺชาโยโค อวิชฺชานุสโย อวิชฺชาปริยุฏฺฐานํ อวิชฺชาลงฺคี โมโห อกุสลมูลํ.\csromanspacer{} ตสฺส ปจฺเจกสมฺพุทฺธสฺส \textbf{มกฺโข} จ กสาโว จ โมโห จ วนฺตา สํวนฺตา นิทฺธนฺตา ปหีนา สมุจฺฉินฺนา วูปสนฺตา ปฏิปฺปสฺสทฺธา อภพฺพุปฺปตฺติกา ญาณคฺคินา ทฑฺฒาติ โส ปจฺเจกสมฺพุทฺโธ นิมฺ\textbf{มกฺโข} นิทฺธนฺตกสาวโมโห.}

\prose{\textbf{นิราสโส สพฺพโลเก ภวิตฺวา}ติ \textbf{อาสา} วุจฺจติ ตณฺหา.\csromanspacer{} โย ราโค สาราโค ฯเปฯ อภิชฺฌา โลโภ อกุสลมูลํ.\csromanspacer{} \textbf{สพฺพโลเก}ติ สพฺพ-อปายโลเก สพฺพมนุสฺสโลเก สพฺพเทวโลเก สพฺพขนฺธโลเก สพฺพธาตุโลเก สพฺพ- อายตนโลเก.\csromanspacer{} \textbf{นิราสโส สพฺพโลเก ภวิตฺวา}ติ สพฺพโลเก นิราสโส ภวิตฺวา นิตฺตโณฺห ภวิตฺวา นิปฺปิปาโส ภวิตฺวาติ นิราสโส สพฺพโลเก ภวิตฺวา, เอโก จเร ขคฺควิสาณกปฺโป.\csromanspacer{} เตนาห โส ปจฺเจกสมฺพุทฺโธ\csromandash{}}

\prose{นิลฺโลลุโป นิกฺกุโห นิปฺปิปาโส,}

\prose{\textbf{นิมฺมกฺโข นิทฺธนฺตกสาวโมโห}.}

\csromangathameasure{\textbf{นิราสโส สพฺพโลเก ภวิตฺวา}, \\ เอโก จเร ขคฺควิสาณกปฺโปติ.}
\csromangathagroup{\csromangathastanza{\textbf{นิราสโส สพฺพโลเก ภวิตฺวา}, \\ เอโก จเร ขคฺควิสาณกปฺโปติ.}}

\csromanhangingbegin
\hangingheaditem{143}{ปาปํ สหายํ ปริวชฺชเยถ,}
\hangingbody{\textbf{อนตฺถทสฺสิํ วิสเม นิวิฏฺฐํ.}}
\hangingbody{\textbf{สยํ น เสเว ปสุตํ ปมตฺตํ,}}
\hangingbody{เอโก จเร ขคฺควิสาณกปฺโป.}
\csromanhangingend

\prose{\csromanfolio{276}\textbf{ปาปํ สหายํ ปริวชฺชเยถา}ติ ปาปสหาโย วุจฺจติ โย โส สหาโย ทสวตฺถุกาย มิจฺฉาทิฏฺฐิยา สมนฺนาคโต “นตฺถิ ทินฺนํ, นตฺถิ ยิฏฺฐํ, นตฺถิ หุตํ, นตฺถิ สุกตทุกฺกฏานํ กมฺมานํ ผลํ วิปาโก, นตฺถิ อยํ โลโก, นตฺถิ ปโร โลโก, นตฺถิ มาตา, นตฺถิ ปิตา, นตฺถิ สตฺตา โอปปาติกา, นตฺถิ โลเก สมณพฺราหฺมณา สมฺมคฺคตา สมฺมาปฏิปนฺนา.\csromanspacer{} เย อิมญฺจ โลกํ ปรญฺจ โลกํ สยํ อภิญฺญา สจฺฉิกตฺวา ปเวเทนฺตี”ติ.\csromanspacer{} อยํ ปาปสหาโย, ปาปํ สหยํ ปริวชฺชเยถาติ ปาปํ สหายํ วชฺเชยฺย ปริวชฺเชยฺยาติ ปาปํ สหายํ ปริวชฺชเยถ.}

\prose{\textbf{อนตฺถทสฺสิํ วิสเม นิวิฏฺฐนฺ}ติ อนตฺถทสฺสี วุจฺจติ โย โส สหาโย ทสวตฺถุกาย มิจฺฉาทิฏฺฐิยา สมนฺนาคโต “นตฺถิ ทินฺนํ, นตฺถิ ยิฏฺฐํ ฯเปฯ.\csromanspacer{} เย อิมญฺจ โลกํ ปรญฺจ โลกํ สยํ อภิญฺญา สจฺฉิกตฺวา ปเวเทนฺตี”ติ.\csromanspacer{} \textbf{วิสเม นิวิฏฺฐนฺ}ติ วิสเม กายกมฺเม นิวิฏฺฐํ.\csromanspacer{} วิสเม วจีกมฺเม นิวิฏฺฐํ.\csromanspacer{} วิสเม มโนกมฺเม นิวิฏฺฐํ.\csromanspacer{} วิสเม ปาณาติปาเต นิวิฏฺฐํ.\csromanspacer{} วิสเม อทินฺนาทาเน นิวิฏฺฐํ.\csromanspacer{} วิสเม กาเมสุมิจฺฉาจาเร นิวิฏฺฐํ.\csromanspacer{} วิสเม มุสาวาเท นิวิฏฺฐํ.\csromanspacer{} วิสมาย ปิสุณาย วาจาย\footnote{ปิสุณวาจาย (ก)} นิวิฏฺฐํ.\csromanspacer{} วิสมาย ผรุสาย วาจาย\footnote{ผรุสวาจาย (ก)} นิวิฏฺฐํ.\csromanspacer{} วิสเม สมฺผปฺปลาเป นิวิฏฺฐํ.\csromanspacer{} วิสมาย อภิชฺฌาย นิวิฏฺฐํ.\csromanspacer{} วิสเม พฺยาปาเท นิวิฏฺฐํ.\csromanspacer{} วิสมาย มิจฺฉาทิฏฺฐิยา นิวิฏฺฐํ.\csromanspacer{} วิสเมสุ สงฺขาเรสุ นิวิฏฺฐํ.\csromanspacer{} วิสเมสุ ปญฺจสุ กามคุเณสุ นิวิฏฺฐํ.\csromanspacer{} วิสเมสุ ปญฺจสุ นีวรเณสุ นิวิฏฺฐํ วินิวิฏฺฐํ สตฺตํ อลฺลีนํ อุปคตํ อชฺโฌสิตํ อธิมุตฺตนฺติ อนตฺถทสฺสิํ วิสเม นิวิฏฺฐํ.}

\prose{\textbf{สยํ น เสเว ปสุตํ ปมตฺตนฺ}ติ \textbf{ปสุตนฺ}ติ โยปิ กาเม เอสติ คเวสติ ปริเยสติ ตจฺจริโต ตพฺพหุโล ตคฺครุโก ตนฺนินฺโน ตปฺโปโณ ตปฺปพฺภาโร ตทธิมุตฺโต ตทธิปเตยฺโย, โสปิ กามปฺปสุโต.\csromanspacer{} โยปิ ตณฺหาวเสน รูเป ปริเยสติ, สทฺเท.\csromanspacer{} คนฺเธ.\csromanspacer{} รเส.\csromanspacer{} โผฏฺฐพฺเพ ปริเยสติ ตจฺจริโต ตพฺพหุโล ตคฺครุโก ตนฺนินฺโน ตปฺโปโณ ตปฺปพฺภาโร ตทธิมุตฺโต ตทธิปเตยฺโย, โสปิ กามปฺปสุโต.\csromanspacer{} โยปิ ตณฺหาวเสน รูเป ปฏิลภติ, สทฺเท.\csromanspacer{} คนฺเธ.\csromanspacer{} รเส.\csromanspacer{} โผฏฺฐพฺเพ ปฏิลภติ ตจฺจริโต ตพฺพหุโล ตคฺครุโก ตนฺนินฺโน ตปฺโปโณ ตปฺปพฺภาโร ตทธิมุตฺโต ตทธิปเตยฺโย, โสปิ กามปฺปสุโต.\csromanspacer{} โยปิ ตณฺหาวเสน รูเป ปริภุญฺชติ, สทฺเท.\csromanspacer{} คนฺเธ.\csromanspacer{} รเส.}

\csromanheader{2}{19. ขคฺควิสาณสุตฺตนิทฺเทส}
\prose{\csromanfolio{277}โผฏฺฐพฺเพ ปริภุญฺชติ ตจฺจริโต ตพฺพหุโล ตคฺครุโก ตนฺนินฺโน ตปฺโปโณ ตปฺปพฺภาโร ตทธิมุตฺโต ตทธิปเตยฺโย, โสปิ กามปฺปสุโต.\csromanspacer{} ยถา กลหการโก กลหปฺปสุโต กมฺมการโก กมฺมปฺปสุโต โคจเร จรนฺโต โคจรปฺปสุโต ฌายี ฌานปฺปสุโต, เอวเมว โย กาเม เอสติ คเวสติ ปริเยสติ ตจฺจริโต ตพฺพหุโล ตคฺครุโก ตนฺนินฺโน ตปฺโปโณ ตปฺปพฺภาโร ตทธิมุตฺโต ตทธิปเตยฺโย, โสปิ กามปฺปสุโต.\csromanspacer{} โยปิ ตณฺหาวเสน รูเป ปริเยสติ ฯเปฯ.\csromanspacer{} โยปิ ตณฺหาวเสน รูเป ปฏิลภติ ฯเปฯ.\csromanspacer{} โยปิ ตณฺหาวเสน รูเป ปริภุญฺชติ, สทฺเท.\csromanspacer{} คนฺเธ.\csromanspacer{} รเส.\csromanspacer{} โผฏฺฐพฺเพ ปริภุญฺชติ ตจฺจริโต ตพฺพหุโล ตคฺครุโก ตนฺนินฺโน ตปฺโปโณ ตปฺปพฺภาโร ตทธิมุตฺโต ตทธิปเตยฺโย, โสปิ กามปฺปสุโต.\csromanspacer{} \textbf{ปมตฺตนฺ}ติ ปมาโท วตฺตพฺโพ กายทุจฺจริเต วา วจีทุจฺจริเต วา มโนทุจฺจริเต วา ปญฺจสุ กามคุเณสุ วา จิตฺตสฺส โวสคฺโค โวสคฺคานุปฺปทานํ กุสลานํ ธมฺมานํ ภาวนาย อสกฺกจฺจกิริยตา อสาตจฺจกิริยตา อนฏฺฐิตกิริยตา โอลีนวุตฺติตา นิกฺขิตฺตจฺฉนฺทตา นิกฺขิตฺตธุรตา อนาเสวนา อภาวนา อพหุลีกมฺมํ อนธิฏฺฐานํ อนนุโยโค, โย เอวรูโป มาโท ปมชฺชนา ปมชฺชิตตฺตํ, อยํ วุจฺจติ ปมาโท.}

\prose{\textbf{สยํ น เสเว ปสุตํ ปมตฺตนฺ}ติ ปสุตํ น เสเวยฺย ปมตฺตญฺจ สยํ น เสเวยฺย สามํ น เสเวยฺย น นิเสเวยฺย น สํเสเวยฺย น ปริสํเสเวยฺย น อาจเรยฺย น สมาจเรยฺย น สมาทาย วตฺเตยฺยาติ สยํ น เสเว ปสุตํ ปมตฺตํ, เอโก จเร ขคฺควิสาณกปฺโป.\csromanspacer{} เตนาห โส ปจฺเจกสมฺพุทฺโธ\csromandash{}}

\csromangathameasure{ปาปํ สหายํ ปริวชฺชเยถ, \\ อนตฺถทสฺสิํ วิสเม นิวิฏฺฐํ.}
\csromangathagroup{\csromangathastanza{ปาปํ สหายํ ปริวชฺชเยถ, \\ อนตฺถทสฺสิํ วิสเม นิวิฏฺฐํ.}}

\prose{สยํ น เสเว ปสุตํ ปมตฺตํ,}

\prose{เอโก จเร ขคฺควิสาณกปฺโปติ.}

\csromanhangingbegin
\hangingheaditem{144}{พหุสฺสุตํ ธมฺมธรํ ภเชถ,}
\hangingbody{\textbf{มิตฺตํ อุฬารํ ปฏิภานวนฺตํ.}}
\hangingbody{\textbf{อญฺญาย อตฺถานิ วิเนยฺย กงฺขํ,}}
\hangingbody{เอโก จเร ขคฺควิสาณกปฺโป.}
\csromanhangingend

\prose{\csromanfolio{278}\textbf{พหุสฺสุตํ ธมฺมธรํ ภเชถา}ติ พหุสฺสุโต โหติ มิตฺโต สุตธโร สุตสนฺนิจโย, เย เต ธมฺมา อาทิกลฺยาณา มชฺเฌกลฺยาณา ปริโยสานกลฺยาณา สาตฺถํ สพฺยญฺชนํ\footnote{สตฺถา สพฺยญฺชนา (สฺยา) ปสฺส ม 3. 62 ปิฏฺเฐ.} เกวลปริปุณฺณํ ปริสุทฺธํ พฺรหฺมจริยํ อภิวทนฺติ, ตถารูปาสฺส ธมฺมา พหุสฺสุตา โหนฺติ ธาตา\footnote{ธตา (สฺยา)} วจสา ปริจิตา มนสานุเปกฺขิตา ทิฏฺฐิยา สุปฺปฏิวิทฺธา.\csromanspacer{} \textbf{ธมฺมธรนฺ}ติ ธมฺมํ ธาเรนฺตํ\footnote{ธาเรติ (ก) 4. ปริยาปุฏํ (สฺยา, ก) ปสฺส ที 3. 164ปิฏฺเฐ. 5. ฉฬภิญฺญาโย (สฺยา)} สุตฺตํ เคยฺยํ เวยฺยากรณํ คาถํ อุทานํ อิติวุตฺตกํ ชาตกํ อพฺภุตธมฺมํ เวทลฺลํ.\csromanspacer{} \textbf{พหุสฺสุตํ ธมฺมธรํภเชถา}ติ พหุสฺสุตญฺจ ธมฺมธรญฺจ มิตฺตํ ภเชยฺย สํภเชยฺย เสเวยฺย นิเสเวยฺย สํเสเวยฺย ปฏิเสเวยฺยาติ พหุสฺสุตํ ธมฺมธรํ ภเชถ.}

\prose{\textbf{มิตฺตํ อุฬารํ ปฏิภานวนฺตนฺ}ติ อุฬาโร โหติ มิตฺโต สีเลน สมาธินา ปญฺญาย วิมุตฺติยา วิมุตฺติญาณทสฺสเนน.\csromanspacer{} \textbf{ปฏิภานวนฺตนฺ}ติ ตโย ปฏิภานวนฺโต ปริยตฺติปฏิภานวา ปริปุจฺฉาปฏิภานวา อธิคมปฏิภานวา.\csromanspacer{} กตโม ปริยตฺติปฏิภานวา, อิเธกจฺจสฺส พุทฺธวจนํ ปริยาปุตํ4 โหติ สุตฺตํ เคยฺยํ เวยฺยากรณํ คาถา อุทานํ อิติวุตฺตกํ ชาตกํ อพฺภุตธมฺมํ เวทลฺลํ, ตสฺส ปริยตฺติํ นิสฺสาย ปฏิภาติ, อยํ ปริยตฺติปฏิภานวา.}

\prose{กตโม ปริปุจฺฉาปฏิภานวา, อิเธกจฺโจ ปริปุจฺฉิโตปิ โหติ อตฺเถ จ ญาเย จ ลกฺขเณ จ การเณ จ ฐานาฏฺฐาเน จ, ตสฺส ปริปุจฺฉํ นิสฺสาย ปฏิภาติ, อยํ ปริปุจฺฉาปฏิภานวา.}

\prose{กตโม อธิคมปฏิภานวา, อิเธกจฺจสฺส อธิคตา โหนฺติ จตฺตาโร สติปฏฺฐานา จตฺตาโร สมฺมปฺปธานา จตฺตาโร อิทฺธิปาทา ปญฺจินฺทฺริยานิ ปญฺจ พลานิ สตฺต โพชฺฌงฺคา อริโย อฏฺฐงฺคิโก มคฺโค จตฺตาโร อริยมคฺคา จตฺตาริ สามญฺญผลานิ จตสฺโส ปฏิสมฺภิ ทาโย ฉ อภิญฺญาโย5.\csromanspacer{} ตสฺส อตฺโถ ญาโต, ธมฺโม ญาโต, นิรุตฺติ ญาตา.\csromanspacer{} อตฺเถ ญาเต อตฺโถ ปฏิภาติ, ธมฺเม ญาเต ธมฺโม ปฏิภาติ, นิรุตฺติยา ญาตาย นิรุตฺติ ปฏิภาติ, อิเมสุ ตีสุ ญาณํ ปฏิภานปฏิสมฺภิทา.\csromanspacer{} โส ปจฺเจกสมฺพุทฺโธ อิมาย ปฏิภานปฏิสมฺภิทาย อุเปโต สมุเปโต อุปาคโต สมุปาคโต อุปปนฺโน สมุปปนฺโน สมนฺนาคโต, ตสฺมา ปจฺเจกสมฺพุทฺโธ}

\csromanheader{2}{19. ขคฺควิสาณสุตฺตนิทฺเทส}
\prose{\csromanfolio{279}ปฏิภานวา.\csromanspacer{} ยสฺส ปริยตฺติ นตฺถิ ปริปุจฺฉา นตฺถิ อธิคโม นตฺถิ, กิํ ตสฺส ปฏิภายิสฺสตีติ มิตฺตํ อุฬารํ ปฏิภานวนฺตํ.}

\prose{\textbf{อญฺญาย อตฺถานิ วิเนยฺย กงฺขนฺ}ติ อตฺตตฺถํ อญฺญาย ปรตฺถํ อญฺญาย อุภยตฺถํ อญฺญาย ทิฏฺฐธมฺมิกตฺถํ อญฺญาย สมฺปรายิกตฺถํ อญฺญาย ปรมตฺถํ อญฺญาย อภิญฺญาย ชานิตฺวา ตุลยิตฺวา ตีรยิตฺวา วิภาวยิตฺวา วิภูตํ กตฺวา กงฺขํ วิเนยฺย ปฏิวิเนยฺย ปชเหยฺย วิโนเทยฺย พฺยนฺตีกเรยฺย อนภาวํ คเมยฺยาติ อญฺญาย อตฺถานิ วิเนยฺย กงฺขํ, เอโก จเร ขคฺควิสาณกปฺโป.\csromanspacer{} เตนาห โส ปจฺเจกสมฺพุทฺโธ\csromandash{}}

\csromangathameasure{พหุสฺสุตํ ธมฺมธรํ ภเชถ, \\ มิตฺตํ อุฬารํ ปฏิภานวนฺตํ.}
\csromangathagroup{\csromangathastanza{พหุสฺสุตํ ธมฺมธรํ ภเชถ, \\ มิตฺตํ อุฬารํ ปฏิภานวนฺตํ.}}

\prose{อญฺญาย อตฺถานิ วิเนยฺย กงฺขํ,}

\prose{เอโก จเร ขคฺควิสาณกปฺโปติ.}

\csromanhangingbegin
\hangingheaditem{145}{\textbf{ขิฑฺฑํ รติํ}\footnote{รตี (สฺยา)}\textbf{ กามสุขญฺจ โลเก},}
\hangingbody{\textbf{อนลงฺกริตฺวา อนเปกฺขมาโน.}}
\hangingbody{\textbf{วิภูสฏฺฐานา2} \textbf{วิรโต สจฺจวาที,}}
\hangingbody{เอโก จเร ขคฺควิสาณกปฺโป.}
\csromanhangingend

\prose{\textbf{ขิฑฺฑํรติํ กามสุขญฺจ โลเก}ติ \textbf{ขิฑฺฑา}ติ เทฺว \textbf{ขิฑฺฑา}, กายิกา \textbf{ขิฑฺฑา} จ วาจสิกา \textbf{ขิฑฺฑา} จ ฯเปฯ อยํ กายิกา \textbf{ขิฑฺฑา} ฯเปฯ อยํ วาจสิกา \textbf{ขิฑฺฑา}.\csromanspacer{} \textbf{รตี}ติ อนุกฺกณฺฐิตาธิวจนเมตํ รตีติ.\csromanspacer{} \textbf{กามสุขนฺ}ติ วุตฺตํ เหตํ ภควตา\footnote{ปสฺส ม 2. 233 ปิฏฺเฐ.}\csromandash{}ปญฺจิเม ภิกฺขเว กามคุณา.\csromanspacer{} กตเม ปญฺจ, จกฺขุวิญฺเญยฺยา รูปา อิฏฺฐา กนฺตา มนาปา ปิยรูปา กามูปสํหิตา รชนียา, โสตวิญฺเญยฺยา สทฺทา ฯเปฯ ฆานวิญฺเญยฺยา คนฺธา.\csromanspacer{} ชิวฺหาวิญฺเญยฺยา รสา.\csromanspacer{} กายวิญฺเญยฺยา โผฏฺฐพฺพา อิฏฺฐา กนฺตา มนาปา ปิยรูปา กามูปสํหิตา รชนียา, อิเม โข ภิกฺขเว ปญฺจ กามคุณา.\csromanspacer{} ยํ โข ภิกฺขเว อิเม ปญฺจ กามคุเณ ปฏิจฺจ อุปฺปชฺชติ สุขํ โสมนสฺสํ, อิทํ วุจฺจติ กามสุขํ.\csromanspacer{} \textbf{โลเก}ติ มนุสฺสโลเกติ ขิฑฺฑํ รติํ กามสุขญฺจ โลเก.}

\prose{\csromanfolio{280}\textbf{อนลงฺกริตฺวา อนเปกฺขมาโน}ติ ขิฑฺฑญฺจ รติญฺจ กามสุขญฺจ โลเก อนลงฺกริตฺวา อนเปกฺโข หุตฺวา ปชหิตฺวา วิโนเทตฺวา พฺยนฺตีกริตฺวา อนภาวํ คเมตฺวาติ อนลงฺกริตฺวา อนเปกฺขมาโน.}

\prose{\textbf{วิภูสฏฺฐานา วิรโต สจฺจวาที}ติ \textbf{วิภูสา}ติ เทฺว \textbf{วิภูสา}, อตฺถิ อคาริก\textbf{วิภูสา}\footnote{อคาริกสฺส วิภูสา (ก) เอวมุปริปิ.}, อตฺถิ อนาคาริก\textbf{วิภูสา}.\csromanspacer{} กตมา อคาริก\textbf{วิภูสา}, เกสา จ มสฺสู จ มาลาคนฺธญฺจ วิเลปนญฺจ อาภรณญฺจ ปิลนฺธนญฺจ วตฺถญฺจ ปารุปนญฺจ\footnote{ปสาธนญฺจ (สฺยา), สยนาสนญฺจ (ก) 254 ปิฏฺเฐ นตฺถิ ปาฐนานตฺตํ.} เวฐนญฺจ อุจฺฉาทนํ ปริมทฺทนํ นฺหาปนํ สมฺพาหนํ อาทาสํ อญฺชนํ มาลาคนฺธวิเลปนํ มุขจุณฺณํ มุขเลปนํ หตฺถพนฺธํ สิขาพนฺธํ ทณฺฑํ นาฬิกํ ขคฺคํ ฉตฺตํ จิตฺรุปาหนํ อุณฺหีสํ มณิํ วาฬพีชนิํ โอทาตานิ วตฺถานิ ทีฆทสานิ อิติ วา, อยํ อคาริก\textbf{วิภูสา}.}

\prose{กตมา อนาคาริก\textbf{วิภูสา}, จีวรมณฺฑนา ปตฺตมณฺฑนา เสนาสนมณฺฑนา อิมสฺส วา ปูติกายสฺส พาหิรานํ วา ปริกฺขารานํ มณฺฑนา วิภูสนา เกฬนา ปริเกฬนา คทฺธิกตา คทฺธิกตฺตํ\footnote{เคธิกตา เคธิกตฺตํ (สฺยา) ปสฺส อภิ 2. 365 ปิฏฺเฐ.} จปลตา\footnote{จปลนา (สฺยา)} จาปลฺยํ, อยํ อนาคาริก\textbf{วิภูสา}.}

\prose{\textbf{สจฺจวาที}ติ โส ปจฺเจกสมฺพุทฺโธ สจฺจวาที สจฺจสนฺโธ เถโต ปจฺจยิโก อวิสํวาทโก โลกสฺส วิภูสฏฺฐานา อารโต วิรโต ปฏิวิรโต นิกฺขนฺโต นิสฺสโฏ วิปฺปมุตฺโต วิสญฺญุตฺโต วิมริยาทิกเตน เจตสา วิหรตีติ วิภูสฏฺฐานา วิรโต สจฺจวาที, เอโก จเร ขคฺควิสาณกปฺโป.\csromanspacer{} เตนาห โส ปจฺเจกสมฺพุทฺโธ\csromandash{}}

\csromangathameasure{ขิฑฺฑํ รติํ กามสุขญฺจ โลเก, \\ \textbf{อนลงฺกริตฺวา อนเปกฺขมาโน}.}
\csromangathagroup{\csromangathastanza{ขิฑฺฑํ รติํ กามสุขญฺจ โลเก, \\ \textbf{อนลงฺกริตฺวา อนเปกฺขมาโน}.}}

\prose{\textbf{วิภูสฏฺฐานา วิรโต สจฺจวาที},}

\prose{เอโก จเร ขคฺควิสาณกปฺโปติ.}

\csromanhangingbegin
\hangingheaditem{146}{ปุตฺตญฺจ ทารํ ปิตรญฺจ มาตรํ,}
\hangingbody{\textbf{ธนานิ ธญฺญานิ จ พนฺธวานิ.}}
\hangingbody{\textbf{หิตฺวาน กามานิ ยโถธิกานิ,}}
\hangingbody{เอโก จเร ขคฺควิสาณกปฺโป.}
\csromanhangingend

\csromanheader{2}{19. ขคฺควิสาณสุตฺตนิทฺเทส}
\prose{\csromanfolio{281}\textbf{ปุตฺตญฺจ ทารํ ปิตรญฺจ มาตรนฺ}ติ \textbf{ปุตฺตา}ติ จตฺตาโร \textbf{ปุตฺตา} อตฺรโช ปุตฺโต เขตฺตโช\footnote{เขตฺรโช (สฺยา, ก)} ปุตฺโต ทินฺนโก ปุตฺโต อนฺเตวาสิโก ปุตฺโต.\csromanspacer{} \textbf{ทารา} วุจฺจนฺติ ภริยาโย.\csromanspacer{} \textbf{ปิตา}ติ โย โส ชนโก.\csromanspacer{} \textbf{มาตา}ติ ยา สา ชนิกาติ ปุตฺตญฺจ ทารํ ปิตรญฺจ มาตรํ.}

\prose{\textbf{ธนานิ ธญฺญานิ จ พนฺธวานี}ติ ธนานิ วุจฺจนฺติ หิรญฺญํ สุวณฺณํ มุตฺตา มณิ เวฬุริโย สงฺโข สิลา ปวาฬํ รชตํ ชาตรูปํ โลหิตงฺโค\footnote{โลหิตโก (?)} มสารคลฺลํ.\csromanspacer{} ธญฺญานิ วุจฺจนฺติ ปุพฺพณฺณํ อปรณฺณํ.\csromanspacer{} \textbf{ปุพฺพณฺณํ} นาม สาลิ วีหิ ยโว โคธุโม กงฺคุ วรโก กุทฺรูสโก\footnote{กุทฺรุสโก (สฺยา)}.\csromanspacer{} \textbf{อปรณฺณํ} นาม สูเปยฺยํ.\csromanspacer{} \textbf{พนฺธวานี}ติ จตฺตาโร พนฺธวา ญาติพนฺธวาปิ พนฺธุ โคตฺตพนฺธวาปิ พนฺธุ มิตฺตพนฺธวาปิ พนฺธุ สิปฺปพนฺธวาปิ พนฺธูติ ธนานิ ธญฺญานิ จ พนฺธวานิ.}

\prose{\textbf{หิตฺวาน กามานิ ยโถธิกานี}ติ \textbf{กามา}ติ อุทฺทานโต เทฺว \textbf{กามา} วตฺถุ\textbf{กามา} จ กิเลส\textbf{กามา} จ ฯเปฯ.\csromanspacer{} อิเม วุจฺจนฺติ วตฺถุ\textbf{กามา} ฯเปฯ.\csromanspacer{} อิเม วุจฺจนฺติ กิเลส\textbf{กามา}.\csromanspacer{} \textbf{หิตฺวาน กามานี}ติ วตฺถุกาเม ปริชานิตฺวา กิเลสกาเม ปหายา ปชหิตฺวา วิโนเทตฺวา พฺยนฺตีกริตฺวา อนภาวํ คเมตฺวา.\csromanspacer{} \textbf{หิตฺวาน กามานิ ยโถธิกานี}ติ โสตาปตฺติมคฺเคน เย กิเลสา ปหีนา, เต กิเลเส น ปุเนติ น ปจฺเจติ น ปจฺจาคจฺฉติ.\csromanspacer{} สกทาคามิมคฺเคน เย กิเลสา ปหีนา.\csromanspacer{} อนาคามิมคฺเคน เย กิเลสา ปหีนา.\csromanspacer{} อรหตฺตมคฺเคน เย กิเลสา ปหีนา, เต กิเลเส น ปุเนติ น ปจฺเจติ น ปจฺจาคจฺฉตีติ หิตฺวาน \textbf{กามา}นิ ยโถธิกานิ, เอโก จเร ขคฺควิสาณกปฺโป.\csromanspacer{} เตนาห โส ปจฺเจกสมฺพุทฺโธ\csromandash{}}

\prose{ปุตฺตญฺจ ทารํ ปิตรญฺจ มาตรํ,}

\prose{ธนานิ ธญฺญานิ จ พนฺธวานิ.}

\csromangathameasure{หิตฺวาน \textbf{กามา}นิ ยโถธิกานิ, \\ เอโก จเร ขคฺควิสาณกปฺโปติ.}
\csromangathagroup{\csromangathastanza{หิตฺวาน \textbf{กามา}นิ ยโถธิกานิ, \\ เอโก จเร ขคฺควิสาณกปฺโปติ.}}

\csromanhangingbegin
\hangingheaditem{147}{\csromanfolio{282}สงฺโค เอโส ปริตฺตเมตฺถ โสขฺยํ\textbf{,}}
\hangingbody{\textbf{อปฺปสฺสาโท ทุกฺขเมตฺถ ภิยฺโย.}}
\hangingbody{\textbf{คโฬ เอโส อิติ ญ ตฺวา มติมา1,}}
\hangingbody{เอโก จเร ขคฺควิสาณกปฺโป.}
\csromanhangingend

\prose{\textbf{สงฺโค เอโส ปริตฺตเมตฺถ โสขฺยนฺ}ติ สงฺโคติ วา พฬิสนฺติ วา อามิสนฺติ วา ลคฺคนนฺติ วา ปลิโพโธติ วา ปญฺจนฺเนตํ กามคุณานํ อธิวจนํ.\csromanspacer{} \textbf{ปริตฺตเมตฺถ โสขฺยนฺ}ติ วุตฺตํ เหตํ ภควตา\csromandash{}“ปญฺจิเม ภิกฺขเว กามคุณา.\csromanspacer{} กตเม ปญฺจ\textbf{,} จกฺขุวิญฺเญยฺยา รูปา อิฏฺฐา กนฺตา มนาปา ปิยรูปา กามูปสํหิตา รชนียา ฯเปฯ กายวิญฺเญยฺยา โผฏฺฐพฺพา อิฏฺฐา กนฺตา มนาปา ปิยรูปา กามูปสํหิตา รชนียา\textbf{,} อิเม โข ภิกฺขเว ปญฺจ กามคุณา.\csromanspacer{} ยํ โข ภิกฺขเว อิเม ปญฺจ กามคุเณ ปฏิจฺจ อุปฺปชฺชติ สุขํ โสมนสฺสํ\textbf{,} อิทํ วุจฺจติ กามสุขํ.\csromanspacer{} อปฺปกํ เอตํ สุขํ\textbf{,} ปริตฺตกํ เอตํ สุขํ\textbf{,} โถกกํ เอตํ สุขํ\textbf{,} โอมกํ เอตํ สุขํ\textbf{,} ลามกํ เอตํ สุขํ\textbf{,} ฉตุกฺกํ เอตํ สุขนฺ”ติ สงฺโค เอโส ปริตฺตเมตฺถ โสขฺยํ.}

\prose{\textbf{อปฺปสฺสาโท ทุกฺขเมตฺถ ภิยฺโย}ติ อปฺปสฺสาทา กามา วุตฺตา ภควตา\footnote{ปสฺส ม 1. 185 ปิฏฺเฐ.} พหุทุกฺขา พหุปายาสา\footnote{พหูปายาสา (สฺยา)}\textbf{,} อาทีนโว เอตฺถ ภิยฺโย.\csromanspacer{} อฏฺฐิกงฺกลูปมา กามา วุตฺตา ภควตา.\csromanspacer{} มํสเปสูปมา กามา วุตฺตา ภควตา.\csromanspacer{} ติณุกฺกูปมา กามา วุตฺตา ภควตา.\csromanspacer{} องฺคารกาสูปมา กามา วุตฺตา ภควตา.\csromanspacer{} สุปินกูปมา กามา วุตฺตา ภควตา.\csromanspacer{} ยาจิตกูปมา กามา วุตฺตา ภควตา.\csromanspacer{} รุกฺขผลูปมา กามา วุตฺตา ภควตา.\csromanspacer{} อสิสูนูปมา กามา วุตฺตา ภควตา\textbf{,} สตฺติสูลูปมา กามา วุตฺตา ภควตา.\csromanspacer{} สปฺปสิรูปมา กามา วุตฺตา ภควตา พหุทุกฺขา พหุปายาสา.\csromanspacer{} อาทีนโว เอตฺถ ภิยฺโยติ อปฺปสฺสาโท ทุกฺขเมตฺถ ภิยฺโย.}

\prose{\textbf{คโฬ เอโส อิติ ญ ตฺวา มติมา}ติ คโฬติ วา พฬิสนฺติ วา อามิสนฺติ วา ลคฺคนนฺติ วา พนฺธนนฺติ วา ปลิโพโธติ วา ปญฺจนฺเนตํ กามคุณานํ อธิวจนํ.\csromanspacer{} \textbf{อิตี}ติ ปทสนฺธิ ปทสํสคฺโค ปทปาริปูรี อกฺขรสมวาโย พฺยญฺชนสิลิฏฺฐตา ปทานุปุพฺพตาเปตํ อิตีติ.\csromanspacer{} \textbf{มติมา}ติ ปณฺฑิโต ปญฺญวา พุทฺธิมา ญาณี วิภาวี เมธาวี.\csromanspacer{} \textbf{คโฬ เอโส อิติ ญ ตฺวา มติมา}ติ มติมา คโฬติ ญ ตฺวา พฬิสนฺติ ญ ตฺวา อามิสนฺติ}

\csromanheader{2}{19. ขคฺควิสาณสุตฺตนิทฺเทส}
\prose{\csromanfolio{283}ญตฺวา ลคฺคนนฺติ ญ ตฺวา พนฺธนนฺติ ญ ตฺวา ปลิโพโธติ ญ ตฺวา ชานิตฺวา ตุลยิตฺวา ตีรยิตฺวา วิภาวยิตฺวา วิภูตํ กตฺวาติ คโฬ เอโส อิติ ญ ตฺวา มติมา, เอโก จเร ขคฺควิสาณกปฺโป.\csromanspacer{} เตนาห โส ปจฺเจกสมฺพุทฺโธ\csromandash{}}

\csromangathameasure{สงฺโค เอโส ปริตฺตเมตฺถ โสขฺยํ, \\ อปฺปสฺสาโท ทุกฺขเมตฺถ ภิยฺโย. \\ คโฬ เอโส อิติ ญตฺวา มติมา, \\ เอโก จเร ขคฺควิสาณกปฺโปติ.}
\csromangathagroup{\csromangathastanza{สงฺโค เอโส ปริตฺตเมตฺถ โสขฺยํ, \\ อปฺปสฺสาโท ทุกฺขเมตฺถ ภิยฺโย. \\ คโฬ เอโส อิติ ญตฺวา มติมา, \\ เอโก จเร ขคฺควิสาณกปฺโปติ.}}

\csromanhangingbegin
\hangingheaditem{148}{สนฺทาลยิตฺวาน สญฺโญชนานิ,}
\hangingbody{\textbf{ชาลํว เภตฺวา สลิลมฺพุจารี.}}
\hangingbody{\textbf{อคฺคีว ทฑฺฒํ อนิวตฺตมาโน,}}
\hangingbody{เอโก จเร ขคฺควิสาณกปฺโป.}
\csromanhangingend

\prose{\textbf{สนฺทาลยิตฺวาน สญฺโญชนานี}ติ ทส สญฺโญชนานิ, กามราคสญฺโญชนํ ปฏิฆสญฺโญชนํ มานสญฺโญชนํ ทิฏฺฐิสญฺโญชนํ วิจิกิจฺฉาสญฺโญชนํ สีลพฺพตปรามาสสญฺโญชนํ ภวราคสญฺโญชนํ อิสฺสาสญฺโญชนํ มจฺฉริยสญฺโญชนํ อวิชฺชาสญฺโญชนํ.\csromanspacer{} \textbf{สนฺทาลยิตฺวาน สญฺโญชนานี}ติ ทส สญฺโญชนานิ ทาลยิตฺวา สนฺทาลยิตฺวา ปชหิตฺวา วิโนเทตฺวา พฺยนฺตีกริตฺวา อนภาวํ คเมตฺวาติ สนฺทาลยิตฺวาน สญฺโญชนานิ.}

\prose{\textbf{ชาลํว เภตฺวา สลิลมฺพุจารี}ติ ชาลํ วุจฺจติ สุตฺตชาลํ.\csromanspacer{} สลิลํ วุจฺจติ อุทกํ, อมฺพุจารี วุจฺจติ มจฺโฉ, ยถา มจฺโฉ ชาลํ ภินฺทิตฺวา ปภินฺทิตฺวา ทาลยิตฺวา ปทาลยิตฺวา สมฺปทาลยิตฺวา จรติ วิหรติ อิริยติ วตฺเตติ ปาเลติ ยเปติ ยาเปติ, เอวเมว เทฺว ชาลา ตณฺหาชาลญฺจ ทิฏฺฐิชาลญฺจ ฯเปฯ อิทํ ตณฺหาชาลํ ฯเปฯ อิทํ ทิฏฺฐิชาลํ.\csromanspacer{} ตสฺส ปจฺเจกสมฺพุทฺธสฺส ตณฺหาชาลํ ปหีนํ, ทิฏฺฐิชาลํ ปฏินิสฺสฏฺฐํ.\csromanspacer{} ตณฺหาชาลสฺส ปหีนตฺตา ทิฏฺฐิชาลสฺส ปฏินิสฺสฏฺฐตฺตา โส ปจฺเจกสมฺพุทฺโธ รูเป น สชฺชติ, สทฺเท น สชฺชติ, คนฺเธ น สชฺชติ ฯเปฯ ทิฏฺฐสุตมุตวิญฺญาตพฺเพสุ ธมฺเมสุ น สชฺชติ น คณฺหาติ น พชฺฌติ น ปลิพชฺฌติ นิกฺขนฺโต นิสฺสโฏ วิปฺปมุตฺโต วิสญฺญุตฺโต วิมริยาทิกเตน เจตสา วิหรตีติ ชาลํว เภตฺวา สลิลมฺพุจารี.}

\prose{\csromanfolio{284}\textbf{อคฺคีว ทฑฺฒํ อนิวตฺตมาโน}ติ ยถา อคฺคิ ติณกฏฺฐุปาทานํ ทหนฺโต คจฺฉติ อนิวตฺตนฺโต\textbf{,} เอวเมว ตสฺส ปจฺเจกสมฺพุทฺธสฺส โสตาปตฺติมคฺเคน เย กิเลสา ปหีนา\textbf{,} เต กิเลเส น ปุเนติ น ปจฺเจติ น ปจฺจาคจฺฉติ.\csromanspacer{} สกทาคามิมคฺเคน.\csromanspacer{} อนาคามิมคฺเคน.\csromanspacer{} อรหตฺตมคฺเคน เย กิเลสา ปหีนา\textbf{,} เต กิเลเส น ปุเนติ น ปจฺเจติ น ปจฺจาคจฺฉตีติ อคฺคีว ทฑฺฒํ อนิวตฺตมาโน\textbf{,} เอโก จเร ขคฺควิสาณกปฺโป.\csromanspacer{} เตนาห โส ปจฺเจกสมฺพุทฺโธ\csromandash{}}

\csromangathameasure{สนฺทาลยิตฺวาน สญฺโญชนานิ\textbf{,} \\ ชาลํว เภตฺวา สลิลมฺพุจารี. \\ \textbf{อคฺคีว ทฑฺฒํ อนิวตฺตมาโน,} \\ เอโก จเร ขคฺควิสาณกปฺโปติ.}
\csromangathagroup{\csromangathastanza{สนฺทาลยิตฺวาน สญฺโญชนานิ\textbf{,} \\ ชาลํว เภตฺวา สลิลมฺพุจารี. \\ \textbf{อคฺคีว ทฑฺฒํ อนิวตฺตมาโน,} \\ เอโก จเร ขคฺควิสาณกปฺโปติ.}}

\proseitem{149}{\textbf{โอกฺขิตฺตจกฺขุ น จ ปาทโลโล,}}

\prose{\textbf{คุตฺตินฺทฺริโย รกฺขิตมานสาโน.}}

\prose{\textbf{อนวสฺสุโต อปริฑยฺหมาโน}\footnote{อปริทยฺหมาโน (ก)}\textbf{,}}

\csromanheader{2}{\textbf{เอโก จเร ขคฺควิสาณกปฺโป. (9)}}
\prose{\textbf{โอกฺขิตฺตจกฺขุ น จ ปาทโลโล}ติ กถํ ขิตฺตจกฺขุ โหติ.\csromanspacer{} อิเธกจฺโจ ภิกฺขุ\footnote{นตฺถิ สฺยา-โปตฺถเก, ขุ 7. 285 ปิฏฺเฐ จ.} จกฺขุโลโล จกฺขุโลลิเยน สมนฺนาคโต โหติ\textbf{,} “อทิฏฺฐํ ทกฺขิตพฺพํ\textbf{,} ทิฏฺฐํ สมติกฺกมิตพฺพนฺ”ติ อาราเมน อารามํ อุยฺยาเนน อุยฺยานํ คาเมน คามํ นิคเมน นิคมํ นคเรน นครํ รฏฺเฐน รฏฺฐํ ชนปเทน ชนปทํ ทีฆจาริกํ อนวฏฺฐิตจาริกํ\footnote{อนวตฺถจาริกํ (สฺยา)} อนุยุตฺโต โหติ รูปทสฺสนาย\textbf{,} เอวํ ขิตฺตจกฺขุ โหติ.}

\prose{อถ วา ภิกฺขุ อนฺตรฆรํ ปวิฏฺโฐ วีถิํ ปฏิปนฺโน อสํวุโต คจฺฉติ\textbf{,} หตฺถิํ โอโลเกนฺโต อสฺสํ โอโลเกนฺโต รถํ โอโลเกนฺโต ปตฺติํ โอโลเกนฺโต กุมารเก โอโลเกนฺโต กุมาริกาโย โอโลเกนฺโต อิตฺถิโย โอโลเกนฺโต ปุริเส โอโลเกนฺโต อนฺตราปณํ โอโลเกนฺโต ฆรมุขานิ โอโลเกนฺโต อุทฺธํ โอโลเกนฺโต อโธ โอโลเกนฺโต ทิสาวิทิสํ วิเปกฺขมาโน\footnote{เปกฺขมาโน (สฺยา, ก)} คจฺฉติ\textbf{,} เอวมฺปิ ขิตฺตจกฺขุ โหติ.}

\csromanheader{2}{19. ขคฺควิสาณสุตฺตนิทฺเทส}
\prose{\csromanfolio{285}อถ วา ภิกฺขุ จกฺขุนา รูปํ ทิสฺวา นิมิตฺตคฺคาหี โหติ อนุพฺยญฺชนคฺคาหี, ยตฺวาธิกรณเมนํ จกฺขุนฺทฺริยํ อสํวุตํ วิหรนฺตํ อภิชฺฌาโทมนสฺสา ปาปกา อกุสลา ธมฺมา อนฺวาสฺสเวยฺยุํ, ตสฺส สํวราย น ปฏิปชฺชติ, น รกฺขติ จกฺขุนฺทฺริยํ, จกฺขุนฺทฺริเย น สํวรํ อาปชฺชติ, เอวมฺปิ ขิตฺตจกฺขุ โหติ.}

\prose{ยถา วา ปเนเก โภนฺโต สมณพฺราหฺมณา สทฺธาเทยฺยานิ โภชนานิ ภุญฺชิตฺวา เต เอวรูปํ วิสูกทสฺสนํ อนุยุตฺตา วิหรนฺติ.\csromanspacer{} เสยฺยถิทํ, นจฺจํ คีตํ วาทิตํ เปกฺขํ อกฺขานํ ปาณิสฺสรํ เวตาฬํ กุมฺภถูณํ โสภนกํ\footnote{โสภนครกํ (สฺยา), โสภนกรณํ (ก)} จณฺฑาลํ วํสํ โธวนํ หตฺถิยุทฺธํ อสฺสยุทฺธํ มหิํสยุทฺธํ\footnote{มหิสยุทฺธํ (สฺยา)} อุสภยุทฺธํ อชยุทฺธํ เมณฺฑยุทฺธํ กุกฺกุฏยุทฺธํ วฏฺฏกยุทฺธํ ทณฺฑยุทฺธํ มุฏฺฐิยุทฺธํ นิพฺพุทฺธํ อุยฺโยธิกํ พลคฺคํ เสนาพฺยูหํ อนีกทสฺสนํ อิติ วา, อิติ เอวรูปํ วิสูกทสฺสนํ อนุยุตฺโต โหติ, เอวมฺปิ ขิตฺตจกฺขุ โหติ.}

\prose{กถํ โอกฺขิตฺตจกฺขุ โหติ.\csromanspacer{} อิเธกจฺโจ ภิกฺขุ น จกฺขุโลโล น จกฺขุโลลิเยน สมนฺนาคโต โหติ, “อทิฏฺฐํ ทกฺขิตพฺพํ, ทิฏฺฐํ สมติกฺกมิตพฺพนฺ”ติ น อาราเมน อารามํ น อุยฺยาเนน อุยฺยานํ น คาเมน คามํ น นิคเมน นิคมํ น นคเรน นครํ น รฏฺเฐน รฏฺฐํ น ชนปเทน ชนปทํ ทีฆจาริกํ อนวฏฺฐิตจารีกํ อนุยุตฺโต โหติ รูปทสฺสนาย, เอวํ โอกฺขิตฺตจกฺขุ โหติ.}

\prose{อถ วา ภิกฺขุ อนฺตรฆรํ ปวิฏฺโฐ วีถิํ ปฏิปนฺโน สํวุโต คจฺฉติ, น หตฺถิํ โอโลเกนฺโต น อสฺสํ โอโลเกนฺโต น รถํ โอโลเกนฺโต น ปตฺติํ โอโลเกนฺโต น กุมารเก โอโลเกนฺโต น กุมาริกาโย โอโลเกนฺโต น อิตฺถิโย โอโลเกนฺโต น ปุริเส โอโลเกนฺโต น อนฺตราปณํ โอโลเกนฺโต น ฆรมุขานิ โอโลเกนฺโต น อุทฺธํ โอโลเกนฺโต น อโธ โอโลเกนฺโต น ทิสาวิทิสํ วิเปกฺขมาโน คจฺฉติ, เอวมฺปิ โอกฺขิตฺตจกฺขุ โหติ.}

\prose{อถ วา ภิกฺขุ จกฺขุนา รูปํ ทิสฺวา น นิมิตฺตคฺคาหี โหติ นานุพฺยญฺชนคฺคาหี, ยตฺวาธิกรณเมนํ จกฺขุนฺทฺริยํ อสํวุตํ วิหรนฺตํ อภิชฺฌาโทมนสฺสา ปาปกา อกุสลา ธมฺมา อนฺวาสฺสเวยฺยุํ, ตสฺส สํวราย ปฏิปชฺชติ, รกฺขติ จกฺขุนฺทฺริยํ, จกฺขุนฺทฺริเย สํวรํ อาปชฺชติ, เอวมฺปิ โอกฺขิตฺตจกฺขุ โหติ.}

\prose{\csromanfolio{286}ยถา วา ปเนเก โภนฺโต สมณพฺราหฺมณา สทฺธาเทยฺยานิ โภชนานิ ภุญฺชิตฺวา เต เอวรูปํ วิสูกทสฺสนํ อนุยุตฺตา วิหรนฺติ.\csromanspacer{} เสยฺยถิทํ, นจฺจํ คีตํ วาทิตํ ฯเปฯ อนีกทสฺสนํ อิติ วา, อิติ เอวรูปา วิสูกทสฺสนา ปฏิวิรโต, เอวมฺปิ โอกฺขิตฺตจกฺขุ โหติ.}

\prose{\textbf{น จ ปาทโลโล}ติ กถํ ปาทโลโล โหติ.\csromanspacer{} อิเธกจฺโจ ภิกฺขุ ปาทโลโล ปาทโลลิเยน สมนฺนาคโต โหติ, อาราเมน อารามํ อุยฺยาเนน อุยฺยานํ คาเมน คามํ นิคเมน นิคมํ นคเรน นครํ รฏฺเฐน รฏฺฐํ ชนปเทน ชนปทํ ทีฆจาริกํ อนวฏฺฐิตจาริกํ อนุยุตฺโต โหติ, เอวมฺปิ ปาทโลโล โหติ.}

\prose{อถ วา ภิกฺขุ อนฺโตสํฆาราเม\footnote{อนฺโตปิ สํฆาราเม (ก)} ปาทโลโล ปาทโลลิเยน สมนฺนาคโต โหติ, น อตฺถเหตุ น การณเหตุ อุทฺธโต อวูปสนฺตจิตฺโต ปริเวณโต ปริเวณํ คจฺฉติ, วิหารโต วิหารํ คจฺฉติ, อฑฺฒโยคโต อฑฺฒโยคํ คจฺฉติ, ปาสาทโต ปาสาทํ คจฺฉติ, หมฺมิยโต หมฺมิยํ คจฺฉติ, คุหโต คุหํ คจฺฉติ, เลณโต เลณํ คจฺฉติ, กุฏิยา กุฏิํ คจฺฉติ, กูฏาคารโต กูฏาคารํ คจฺฉติ, อฏฺฏโต อฏฺฏํ คจฺฉติ, มาฬโต มาฬํ คจฺฉติ, อุทฺทณฺฑโต อุทฺทณฺฑํ คจฺฉติ\footnote{อุทฺทณฺฑํ คจฺฉติ, อุทฺโธสิตโต อุทฺโธสิตํ คจฺฉติ (สฺยา) ปสฺส ขุ 7. 51 ปิฏฺเฐ.}, อุปฏฺฐานสาลโต อุปฏฺฐานสาลํ คจฺฉติ, มณฺฑปโต มณฺฑปํ คจฺฉติ, รุกฺขมูลโต รุกฺขมูลํ คจฺฉติ.\csromanspacer{} ยตฺถ วา ปน ภิกฺขู นิสีทนฺติ วา คจฺฉนฺติ วา, ตตฺถ เอกสฺส วา ทุติโย โหติ, ทฺวินฺนํ วา ตติโย โหติ, ติณฺณํ วา จตุตฺโถ โหติ, ตตฺถ พหุํ สมฺผปฺปลาปํ ปลปติ.\csromanspacer{} เสยฺยถิทํ, ราชกถํ โจรกถํ ฯเปฯ อิติภวาภวกถํ กเถติ, เอวมฺปิ ปาทโลโล โหติ.}

\prose{\textbf{น จ ปาทโลโล}ติ โส ปจฺเจกสมฺพุทฺโธ ปาทโลลิยา อารโต วิรโต ปฏิวิรโต นิกฺขนฺโต นิสฺสโฏ วิปฺปมุตฺโต วิสญฺญุตฺโต วิมริยาทิกเตน เจตสา, ปฏิสลฺลานาราโม โหติ ปฏิสลฺลานรโต อชฺฌตฺตํ เจโตสมถมนุยุตฺโต อนิรากตชฺฌาโน วิปสฺสนาย สมนฺนาคโต พฺรูเหตา สุญฺญาคารํ ฌายี ฌานรโต เอกตฺตมนุยุตฺโต สทตฺถครุโกติ โอกฺขิตฺตจกฺขุ น จ ปาทโลโล.}

\csromanheader{2}{19. ขคฺควิสาณสุตฺตนิทฺเทส}
\prose{\csromanfolio{287}\textbf{คุตฺตินฺทฺริโย รกฺขิตมานสาโน}ติ \textbf{คุตฺตินฺทฺริโย}ติ โส ปจฺเจกสมฺพุทฺโธ จกฺขุนา รูปํ ทิสฺวา น นิมิตฺตคฺคาหี โหติ นานุพฺยญฺชนคฺคาหี, ยตฺวาธิกรณเมนํ จกฺขุนฺทฺริยํ อสํวุตํ วิหรนฺตํ อภิชฺฌาโทมนสฺสา ปาปกา อกุสลา ธมฺมา อนฺวาสฺสเวยฺยุํ, ตสฺส สํวราย ปฏิปชฺชติ, รกฺขติ จกฺขุนฺทฺริยํ, จกฺขุนฺทฺริเย สํวรํ อาปชฺชติ, โสเตน สทฺทํ สุตฺวา ฯเปฯ ฆาเนน คนฺธํ ฆายิตฺวา.\csromanspacer{} ชิวฺหาย รสํ สายิตฺวา.\csromanspacer{} กาเยน โผฏฺฐพฺพํ ผุสิตฺวา.\csromanspacer{} มนสา ธมฺมํ วิญฺญาย น นิมิตฺตคฺคาหี โหติ นานุพฺยญฺชนคฺคาหี, ยตฺวาธิกรณเมนํ มนินฺทฺริยํ อสํวุตํ วิหรนฺตํ อภิชฺฌาโทมนสฺสา ปาปกา อกุสลา ธมฺมา อนฺวาสฺสเวยฺยุํ, ตสฺส สํวราย ปฏิปชฺชติ, รกฺขติ มนินฺทฺริยํ, มนินฺทฺริเย สํวรํ อาปชฺชตีติ \textbf{คุตฺตินฺทฺริโย}.\csromanspacer{} \textbf{รกฺขิตมานสาโน}ติ โคปิตมานสาโนติ \textbf{คุตฺตินฺทฺริโย} รกฺขิตมานสาโน.}

\prose{\textbf{อนวสฺสุโต ปริฑยฺหมาโน}ติ วุตฺตํ เหตํ อายสฺมตา มหาโมคฺคลฺลาเนน\csromandash{}อวสฺสุตปริยายญฺจ\footnote{ปสฺส สํ 2. 390 ปิฏฺเฐ.} โว อาวุโส เทเสสฺสามิ อนวสฺสุตปริยายญฺจ, ตํ สุณาถ สาธุกํ มนสิกโรถ, ภาสิสฺสามีติ. “เอวมาวุโส”ติ โข เต ภิกฺขู อายสฺมโต มหาโมคฺคลฺลานสฺส ปจฺจสฺโสสุํ.\csromanspacer{} อายสฺมา มหาโมคฺคลฺลาโน\footnote{มหาโมคฺคลาโน (ก)} เอตทโวจ\csromandash{}}

\prose{กถญฺจาวุโส อวสฺสุโต โหติ, อิธาวุโส ภิกฺขุ จกฺขุนา รูปํ ทิสฺวา ปิยรูเป รูเป อธิมุจฺจติ, อปฺปิยรูเป รูเป พฺยาปชฺชติ, อนุปฏฺฐิตกายสฺสติ\footnote{อนุปฏฺฐิตกายสติ (สฺยา, ก)} จ วิหรติ ปริตฺตเจตโส, ตญฺจ เจโตวิมุตฺติํ ปญฺญาวิมุตฺติํ ยถาภูตํ นปฺปชานาติ, ยตฺถสฺส\footnote{ตตฺถ เย (ก) สํ 2. 390 ปิฏฺเฐ.} เต อุปฺปนฺนา ปาปกา อกุสลา ธมฺมา อปริเสสา นิรุชฺฌนฺติ.\csromanspacer{} โสเตน สทฺทํ สุตฺวา ฯเปฯ.\csromanspacer{} มนสา ธมฺมํ วิญฺญาย ปิยรูเป ธมฺเม อธิมุจฺจติ, อปฺปิยรูเป ธมฺเม พฺยาปชฺชติ, อนุปฏฺฐิตกายสฺสติ จ วิหรติ ปริตฺตเจตโส, ตญฺจ เจโตวิมุตฺติํ ปญฺญาวิมุตฺติํ ยถาภูตํ นปฺปชานาติ, ยตฺถสฺส เต อุปฺปนฺนา ปาปกา อกุสลา ธมฺมา อปริเสสา นิรุชฺฌนฺติ.\csromanspacer{} อยํ วุจฺจตาวุโส ภิกฺขุ อวสฺสุโต จกฺขุวิญฺเญยฺเยสุ รูเปสุ ฯเปฯ อวสฺสุโต มโนวิญฺเญยฺเยสุ ธมฺเมสุ เอวํวิหาริํ จาวุโส ภิกฺขุํ จกฺขุโต เจปิ นํ มาโร อุปสงฺกมติ, ลเภเถว\footnote{ลภเตว (สฺยา, ก) เอวมุปริปิ.} มาโร}

\prose{\csromanfolio{288}โอตารํ, ลเภถ\footnote{ลภติ (สฺยา, ก) เอวมุปริปิ.} มาโร อารมฺมณํ.\csromanspacer{} โสตโต เจปิ นํ ฯเปฯ.\csromanspacer{} มนโต เจปิ นํ มาโร อุปสงฺกมติ, ลเภเถว มาโร โอตารํ, ลเภถ มาโร อารมฺมณํ.}

\prose{เสยฺยถาปิ อาวุโส นฬาคารํ วา ติณาคารํ วา สุกฺขํ โกฬาปํ\footnote{โกลาปํ (สฺยา) สํ 2. 391 ปิฏฺเฐปิ.} เตโรวสฺสิกํ.\csromanspacer{} ปุรตฺถิมาย เจปิ นํ ทิสาย ปุริโส อาทิตฺตาย ติณุกฺกาย อุปสงฺกเมยฺย, ลเภเถว อคฺคิ โอตารํ, เลเภถ อคฺคิ อารมฺมณํ.\csromanspacer{} ปจฺฉิมาย เจปิ นํ ทิสาย ฯเปฯ.\csromanspacer{} อุตฺตราย เจปิ นํ ทิสาย.\csromanspacer{} ทกฺขิณาย เจปิ นํ ทิสาย.\csromanspacer{} เหฏฺฐิมโต\footnote{ปจฺฉโต (สฺยา), เหฏฺฐิมาย (ก)} เจปิ นํ ทิสาย.\csromanspacer{} อุปริมโต\footnote{อุปริโต (สฺยา), อุปริมาย (ก)} เจปิ นํ ทิสาย.\csromanspacer{} ยโต กุโตจิ เจปิ นํ ปุริโส อาทิตฺตาย ติณุกฺกาย อุปสงฺกเมยฺย, ลเภเถว อคฺคิ โอตารํ.\csromanspacer{} ลเภถ อคฺคิ อารมฺมณํ.\csromanspacer{} เอวเมว โข อาวุโส เอวํวิหาริํ ภิกฺขุํ จกฺขุโต เจปิ นํ มาโร อุปสงฺกมติ.\csromanspacer{} ลเภเถว มาโร โอตารํ, ลเภถ มาโร อารมฺมณํ.\csromanspacer{} โสตโต เจปิ นํ ฯเปฯ.\csromanspacer{} มนโต เจปิ นํ มาโร อุปสงฺกมติ, ลเภเถว มาโร โอตารํ, ลเภถ มาโร อารมฺมณํ.}

\prose{เอวํวิหาริํ จาวุโส ภิกฺขุํ รูปา อธิภํสุ\footnote{อภิภวิํสุ (สฺยา), อภิภํสุ (ก) เอวมุปริปิ.}, น ภิกฺขุ รูเป อธิโภสิ.\csromanspacer{} สทฺทา ภิกฺขุํ อธิภํสุ, น ภิกฺขุ สทฺเท อธิโภสิ.\csromanspacer{} คนฺธา ภิกฺขุํ อธิภํสุ, น ภิกฺขุ คนฺเธ อธิโภสิ.\csromanspacer{} รสา ภิกฺขุํ อธิภํสุ, น ภิกฺขุ รเส อธิโภสิ.\csromanspacer{} โผฏฺฐพฺพา ภิกฺขุํ อธิภํสุ, น ภิกฺขุ โผฏฺฐพฺเพ อธิโภสิ.\csromanspacer{} ธมฺมา ภิกฺขุํ อธิภํสุ, น ภิกฺขุ ธมฺเม อธิโภสิ.\csromanspacer{} อยํ วุจฺจตาวุโส ภิกฺขุ รูปาธิภูโต สทฺทาธิภูโต คนฺธาธิภูโต รสาธิภูโต โผฏฺฐพฺพาธิภูโต ธมฺมาธิภูโต อธิภูโต อนธิภู, อธิภํสุ นํ ปาปกา อกุสลา ธมฺมา สํกิเลสิกา โปโนภวิกา สทรา ทุกฺขวิปากา อายติํ ชาติชรามรณิยา, เอวํ โข อาวุโส อวสฺสุโต โหติ.}

\prose{กถญฺจาวุโส อนวสฺสุโต โหติ.\csromanspacer{} อิธาวุโส ภิกฺขุ จกฺขุนา รูปํ ทิสฺวา ปิยรูเป รูเป นาธิมุจฺจติ, อปฺปิยรูเป รูเป น พฺยาปชฺชติ, อุปฏฺฐิตกายสฺสติ จ วิหรติ อปฺปมาณเจตโส.\csromanspacer{} ตญฺจ เจโตวิมุตฺติํ}

\csromanheader{2}{19. ขคฺควิสาณสุตฺตนิทฺเทส}
\prose{\csromanfolio{289}ปญฺญาวิมุตฺติํ ยถาภูตํ ปชานาติ, ยตฺถสฺส เต อุปฺปนฺนา ปาปกา อกุสลา ธมฺมา อปริเสสา นิรุชฺฌนฺติ.\csromanspacer{} โสเตน สทฺทํ สุตฺวา ฯเปฯ.\csromanspacer{} มนสา ธมฺมํ วิญฺญาย ปิยรูเป ธมฺเม นาธิมุจฺจติ, อปฺปิยรูเป ธมฺเม น พฺยาปชฺชติ, อุปฏฺฐิตกายสฺสติ จ วิหรติ อปฺปมาณเจตโส.\csromanspacer{} ตญฺจ เจโตวิมุตฺติํ ปญฺญาวิมุตฺติํ ยถาภูตํ ปชานาติ.\csromanspacer{} ยตฺถสฺส เต อุปฺปนฺนา ปาปกา อกุสลา ธมฺมา อปริเสสา นิรุชฺฌนฺติ.\csromanspacer{} อยํ วุจฺจตาวุโส ภิกฺขุ อนวสฺสุโต จกฺขุวิญฺเญยฺเยสุ รูเปสุ ฯเปฯ อนวสฺสุโต มโนวิญฺเญยฺเยสุ ธมฺเมสุ.\csromanspacer{} เอวํวิหาริํ จาวุโส ภิกฺขุํ จกฺขุโต เจปิ นํ มาโร อุปสงฺกมติ, เนว ลเภถ มาโร โอตารํ, น ลเภถ มาโร อารมฺมณํ.\csromanspacer{} โสตโต เจปิ นํ ฯเปฯ.\csromanspacer{} มนโต เจปิ นํ มาโร อุปสงฺกมติ, เนว ลเภถ มาโร โอตารํ, น ลเภถ มาโร อารมฺมณํ.}

\prose{เสยฺยถาปิ อาวุโส กูฏาคารา วา กูฏาคารสาลา\footnote{สนฺถาคารสาลา (สฺยา) ปสฺส สํ 2. 392 ปิฏฺเฐ.} วา พหลมตฺติกา อทฺทาวเลปนา\footnote{อลฺลาวเลปนา (สฺยา)}.\csromanspacer{} ปุรตฺถิมาย เจปิ นํ ทิสาย ปุริโส อาทิตฺตาย ติณุกฺกาย อุปสงฺกเมยฺย, เนว ลเภถ อคฺคิ โอตารํ, น ลเภถ อคฺคิ อารมฺมณํ.\csromanspacer{} ปจฺฉิมาย เจปิ นํ ทิสาย.\csromanspacer{} อุตฺตราย เจปิ นํ ทิสาย.\csromanspacer{} ทกฺขิณาย เจปิ นํ ทิสาย.\csromanspacer{} เหฏฺฐิมโต เจปิ นํ ทิสาย.\csromanspacer{} อุปริมโต เจปิ นํ ทิสาย.\csromanspacer{} ยโต กุโตจิ เจปิ นํ ปุริโส อาทิตฺตาย ติณุกฺกาย อุปสงฺกเมยฺย, เนว ลเภถ อคฺคิ โอตารํ, น ลเภถ อคฺคิ อารมฺมณํ.\csromanspacer{} เอวเมว โข อาวุโส เอวํวิหาริํ ภิกฺขุํ จกฺขุโต เจปิ นํ มาโร อุปสงฺกมติ, เนว ลเภถ มาโร โอตารํ, น ลเภถ มาโร อารมฺมณํ.\csromanspacer{} โสตโต เจปิ นํ ฯเปฯ.\csromanspacer{} มนโต เจปิ นํ มาโร อุปสงฺกมติ, เนว ลเภถ มาโร โอตารํ, น ลเภถ มาโร อารมฺมณํ.}

\prose{เอวํวิหารี จาวุโส ภิกฺขุ รูเป อธิโภสิ\footnote{อภิภวิํสุ (สฺยา), อภิโภสิ (ก) ปสฺส สํ 2. 392 ปิฏฺเฐ.}, น รูปา ภิกฺขุํ อธิภํสุ.\csromanspacer{} สทฺเท ภิกฺขุ อธิโภสิ, น สทฺทา ภิกฺขุํ อธิภํสุ.\csromanspacer{} คนฺเธ ภิกฺขุ อธิโภสิ, น คนฺธา ภิกฺขุํ อธิภํสุ.\csromanspacer{} รเส ภิกฺขุ อธิโภสิ, น รสา ภิกฺขุํ อธิภํสุ.\csromanspacer{} โผฏฺฐพฺเพ ภิกฺขุ อธิโภสิ, น โผฏฺฐพฺพา ภิกฺขุํ อธิภํสุ.\csromanspacer{} ธมฺเม ภิกฺขุ อธิโภสิ, น ธมฺมา ภิกฺขุํ อธิภํสุ.\csromanspacer{} อยํ}

\prose{\csromanfolio{290}วุจฺจตาวุโส ภิกฺขุ รูปาธิภู สทฺทาธิภู คนฺธาธิภู รสาธิภู โผฏฺฐพฺพาธิภู ธมฺมาธิภู อธิภู อนธิภูโต\footnote{อนภิภูโต เกหิ จิ กิเลเสหิ (ก) ปสฺส สํ 2. 392 ปิฏฺเฐ.} อธิโภสิ เต ปาปเก อกุสเล ธมฺเม สํกิเลสิเก โปโนภวิเก สทเร ทุกฺขวิปาเก อายติํ ชาติชรามรณิเย.\csromanspacer{} เอวํ โข อาวุโส อนวสฺสุโต โหตีติ อนวสฺสุโต.}

\prose{\textbf{อปริฑยฺหมาโน}ติ ราคเชน ปริฬาเหน อปริฑยฺหมาโน, โทสเชน ปริฬาเหน อปริฑยฺหมาโน, โมหเชน ปริฬาเหน อปริฑยฺหมาโนติ อนวสฺสุโต อปริฑยฺหมาโน, เอโก จเร ขคฺควิสาณกปฺโป.\csromanspacer{} เตนาห โส ปจฺเจกสมฺพุทฺโธ\csromandash{}}

\csromangathameasure{โอกฺขิตฺตจกฺขุ น จ ปาทโลโล, \\ คุตฺตินฺทฺริโย รกฺขิตมานสาโน. \\ อนวสฺสุโต อปริฑยฺหมาโน, \\ เอโก จเร ขคฺควิสาณกปฺโปติ.}
\csromangathagroup{\csromangathastanza{โอกฺขิตฺตจกฺขุ น จ ปาทโลโล, \\ คุตฺตินฺทฺริโย รกฺขิตมานสาโน. \\ อนวสฺสุโต อปริฑยฺหมาโน, \\ เอโก จเร ขคฺควิสาณกปฺโปติ.}}

\proseitem{150}{โอหารยิตฺวา คิหิพฺยญฺชนานิ,}

\prose{\textbf{สญฺฉนฺนปตฺโต ยถา ปาริฉตฺตโก.}}

\prose{\textbf{กาสายวตฺโถ อภินิกฺขมิตฺวา,}}

\csromanheader{2}{\textbf{เอโก จเร ขคฺควิสาณกปฺโป. (10)}}
\prose{\textbf{โอหารยิตฺวา คิหิพฺยญฺชนานี}ติ คิหิพฺยญฺชนานิ วุจฺจนฺติ เกสา จ มสฺสู จ ฯเปฯ ทีฆทสานิ อิติ วา.\csromanspacer{} \textbf{โอหารยิตฺวา คิหิพฺยญฺชนานี}ติ คิหิพฺยญฺชนานิ โอโรปยิตฺวา สโมโรปยิตฺวา นิกฺขิปิตฺวา ปฏิปฺปสฺสมฺภิตฺวาติ โอหารยิตฺวา คิหิพฺยญฺชนานิ.}

\prose{\textbf{สญฺฉนฺนปตฺโต ยถา ปาริฉตฺตโก}ติ ยถา โส ปาริจฺฉตฺตโก โกวิฬาโร พหลปตฺตปลาโส\footnote{สาขปตฺตปลาโส (ก)} สนฺทจฺฉาโย\footnote{สณฺฑจฺฉาโย (สฺยา), สนฺตจฺฉาโย (ก) ปสฺส ม 1. 108 ปิฏฺเฐ.}.\csromanspacer{} เอวเมว โส ปจฺเจกสมฺพุทฺโธ ปริปุณฺณปตฺตจีวรธโรติ สญฺฉนฺนปตฺโต ยถา ปาริฉตฺตโก.}

\csromanheader{2}{19. ขคฺควิสาณสุตฺตนิทฺเทส}
\prose{\csromanfolio{291}\textbf{กาสายวตฺโถ อภินิกฺขมิตฺวา}ติ โส ปจฺเจกสมฺพุทฺโธ สพฺพํ ฆราวาสปลิโพธํ ฉินฺทิตฺวา ปุตฺตทารปลิโพธํ ฉินฺทิตฺวา ญาติปลิโพธํ ฉินฺทิตฺวา มิตฺตามจฺจปลิโพธํ ฉินฺทิตฺวา สนฺนิธิปลิโพธํ ฉินฺทิตฺวา เกสมสฺสุํ โอหาเรตฺวา\footnote{โอหารยิตฺวา (สฺยา, ก)} กาสายานิ วตฺถานิ อจฺฉาเทตฺวา อคารสฺมา อนคาริยํ ปพฺพชิตฺวา อกิญฺจนภาวํ อุปคนฺตฺวา เอโก จรติ วิหรติ อิริยติ วตฺเตติ ปาเลติ ยเปติ ยาเปตีติ กาสายวตฺโถ อภินิกฺขมิตฺวา, เอโก จเร ขคฺควิสาณกปฺโป.\csromanspacer{} เตนาห โส ปจฺเจกสมฺพุทฺโธ\csromandash{}}

\csromangathameasure{โอหารยิตฺวา คิหิพฺยญฺชนานิ, \\ สญฺฉนฺนปตฺโต ยถา ปาริฉตฺตโก.}
\csromangathagroup{\csromangathastanza{โอหารยิตฺวา คิหิพฺยญฺชนานิ, \\ สญฺฉนฺนปตฺโต ยถา ปาริฉตฺตโก.}}

\prose{\textbf{กาสายวตฺโถ อภินิกฺขมิตฺวา},}

\prose{เอโก จเร ขคฺควิสาณกปฺโปติ.}

\proseitem{151}{\textbf{รเสสุ เคธํ อกรํ อโลโล},}

\prose{\textbf{อนญฺญโปสี สปทานจารี.}}

\prose{\textbf{กุเล กุเล อปฺปฏิพทฺธจิตฺโต,}}

\csromanheader{2}{\textbf{เอโก จเร ขคฺควิสาณกปฺโป. (11)}}
\prose{\textbf{รเสสุ เคธํ อกรํ อโลโล}ติ \textbf{รโส}ติ มูล\textbf{รโส} ขนฺธ\textbf{รโส} ตจ\textbf{รโส} ปตฺต\textbf{รโส} ปุปฺผ\textbf{รโส} ผล\textbf{รโส} อมฺพิลํ มธุรํ ติตฺตกํ กฏุกํ โลณิกํ ขาริกํ ลมฺพิกํ\footnote{ลมฺพิลํ (สฺยา, ก) ปสฺส อายตนวิภงฺเค 73 ปิฏฺเฐ.} กสาโว สาทุ อสาทุ สีตํ อุณฺหํ.\csromanspacer{} สนฺเตเก สมณพฺราหฺมณา รสคิทฺธา, เต ชิวฺหคฺเคน รสคฺคานิ ปริเยสนฺตา อาหิณฺฑนฺติ, เต อมฺพิลํ ลภิตฺวา อนมฺพิลํ ปริเยสนฺติ, อนมฺพิลํ ลภิตฺวา อมฺพิลํ ปริเยสนฺติ, มธุรํ ลภิตฺวา อมธุรํ ปริเยสนฺติ, อมธุรํ ลภิตฺวา มธุรํ ปริเยสนฺติ, ติตฺตกํ ลภิตฺวา อติตฺตกํ ปริเยสนฺติ, อติกฺกกํ ลภิตฺวา ติตฺตกํ ปริเยสนฺติ, กฏุกํ ลภิตฺวา อกฏุกํ ปริเยสนฺติ, อกฏุกํ ลภิตฺวา กฏุกํ ปริเยสนฺติ, โลณิกํ ลภิตฺวา อโลณิกํ ปริเยสนฺติ, อโลณิกํ ลภิตฺวา โลณิกํ ปริเยสนฺติ, ขาริกํ ลภิตฺวา อขาริกํ ปริเยสนฺติ, อขาริกํ ลภิตฺวา ขาริกํ ปริเยสนฺติ, กสาวํ ลภิตฺวา อกสาวํ ปริเยสนฺติ, อกสาวํ ลภิตฺวา กสาวํ ปริเยสนฺติ, ลมฺพิกํ ลภิตฺวา อลมฺพิกํ ปริเยสนฺติ, อลมฺพิกํ ลภิตฺวา ลมฺพิกํ ปริเยสนฺติ, สาธุํ ลภิตฺวา อสาทุํ ปริเยสนฺติ, อสาทุํ}

\prose{\csromanfolio{292}ลภิตฺวา สาทุํ ปริเยสนฺติ, สีตํ ลภิตฺวา อุณฺหํ ปริเยสนฺติ, อุณฺหํ ลภิตฺวา สีตํ ปริเยสนฺติ.\csromanspacer{} เต ยํ ยํ ลภนฺติ เตน เตน น ตุสฺสนฺติ, อปราปรํ ปริเยสนฺติ มนาปิเกสุ รเสสุ รตฺตา คิทฺธา คถิตา มุจฺฉิตา อชฺโฌสนฺนา ลคฺคา ลคฺคิตา ปลิพุทฺธา.\csromanspacer{} สา รสตณฺหา ตสฺส ปจฺเจกสมฺพุทฺธสฺส ปหีนา อุจฺฉินฺนมูลา ตาลาวตฺถุกตา อนภาวํกตา อายติํ อนุปฺปาทมฺโม, ตสฺมา โส ปจฺเจกสมฺพุทฺโธ ปฏิสงฺขา โยนิโส อาหารํ อาหาเรติ “เนว ทวาย, น มทาย, น มณฺฑนาย, น วิภูสนาย, ยาวเทว อิมสฺส กายสฺส ฐิติยา, ยาปนาย, วิหิํสูปรติยา, พฺรหฺมจริยานุคฺคหาย, อิติ ปุราณญฺจ เวทนํ ปฏิหงฺขามิ, นวญฺจ เวทนํ น อุปฺปาเทสฺสามิ, ยาตฺรา จ เม ภวิสฺสติ อนวชฺชตา จ ผาสุวิหาโร จา”ติ.}

\prose{ยถา วณํ อาลิมฺเปยฺย ยาวเทว อารุหณตฺถาย\footnote{โรปนตฺถาย (สฺยา)}, ยถา วา อกฺขํ อพฺภญฺเชยฺย ยาวเทว ภารสฺส นิตฺถรณตฺถาย, ยถา ปุตฺตมํสํ อาหารํ อาหาเรยฺย ยาวเทว กนฺตารสฺส นิตฺถรณตฺถาย.\csromanspacer{} เอวเมว โส ปจฺเจกสมฺพุทฺโธ ปฏิสงฺขา โยนิโส อาหารํ อาหาเรติ “เนว ทวาย, น มทาย, น มณฺฑนาย, น วิภูสนาย, ยาวเทว อิมสฺส กายสฺส ฐิติยา, ยาปนาย, วิหิํสูปรติยา, พฺรหฺมจริยานุคฺคหาย, อิติ ปุราณญฺจ เวทนํ ปฏิหงฺขามิ, นวญฺจ เวทนํ น อุปฺปาเทสฺสามิ, ยาตฺรา จ เม ภวิสฺสติ อนวชฺชตา จ ผาสุวิหาโร จา”ติ.\csromanspacer{} รสตณฺหาย อารโต วิรโต ปฏิวิรโต นิกฺขนฺโต นิสฺสโฏ วิปฺปมุตฺโต วิสญฺญุตฺโต วิมริยาทิกเตน เจตสา วิหรตีติ รเสสุ เคธํ อกรํ.}

\prose{\textbf{อโลโล}ติ \textbf{โลลุปฺปํ} วุจฺจติ ตณฺหา, โย ราโค สาราโค ฯเปฯ อภิชฺฌา โลโภ อกุสลมูลํ.\csromanspacer{} สา โลลุปฺปา ตณฺหา ตสฺส ปจฺเจกสมฺพุทฺธสฺส ปหีนา อุจฺฉินฺนมูลา ตาลาวตฺถุกตา อนภาวํกตา อายติํ อนุปฺปาทธมฺมา, ตสฺมา ปจฺเจกสมฺพุทฺโธ อโลโลติ รเสสุ เคธํ อกรํ อโลโล.}

\prose{\textbf{อนญฺญโปสี สปทานจารี}ติ \textbf{อนญฺญโปสี}ติ โส ปจฺเจกสมฺพุทฺโธ อตฺตานญฺเญว โปเสติ น ปรนฺติ.}

\csromanheader{2}{19. ขคฺควิสาณสุตฺตนิทฺเทส}
\csromangathameasure{อนญฺญโปสิ’มญฺญาตํ, ทนฺตํ สาเร ปติฏฺฐิตํ. \\ ขีณาสวํ วนฺตโทสํ, ตมหํ พฺรูมิ พฺราหฺมณนฺติ.}
\csromangathagroup{\csromangathastanza{\csromanfolio{293}อนญฺญโปสิ’มญฺญาตํ, ทนฺตํ สาเร ปติฏฺฐิตํ\footnote{สาเรสุ สุปติฏฺฐิตํ (สฺยา, ก) ปสฺส ขุ 1. 81 ปิฏฺเฐ.}. \\ ขีณาสวํ วนฺตโทสํ, ตมหํ พฺรูมิ พฺราหฺมณนฺติ.}}

\prose{\textbf{อนญฺญโปสี สปทานจารี}ติ โส ปจฺเจกสมฺพุทฺโธ ปุพฺพณฺหสมยํ\footnote{ปุพฺพนฺหสมยํ (ก)} นิวาเสตฺวา ปตฺตจีวรมาทาย คามํ วา นิคมํ วา ปิณฺฑาย ปวิสติ, รกฺขิเตเนว กาเยน รกฺขิตาย วาจาย รกฺขิเตน จิตฺเตน อุปฏฺฐิตาย สติยา สํวุเตหิ อินฺทฺริเยหิ โอกฺขิตฺตจกฺขุ อิริยาปถสมฺปนฺโน กุลา กุลํ อนติกฺกมนฺโต ปิณฺฑาย จรตีติ อนญฺญโปสี สปทานจารี.}

\prose{\textbf{กุเล กุเล อปฺปฏิพทฺธจิตฺโต}ติ ทฺวีหิ การเณหิ ปฏิพทฺธจิตฺโต โหติ, อตฺตานํ วา นีจํ ฐเปนฺโต ปรํ อุจฺจํ ฐเปนฺโต ปฏิพทฺธจิตฺโต โหติ, อตฺตานํ วา อุจฺจํ ฐเปนฺโต ปรํ นีจํ ฐเปนฺโต ปฏิพทฺธจิตฺโต โหติ.\csromanspacer{} กถํ อตฺตานํ นีจํ ฐเปนฺโต ปรํ อุจฺจํ ฐเปนฺโต ปฏิพทฺธจิตฺโต โหติ, ตุเมฺห เม พหูปการา, อหํ ตุเมฺห นิสฺสาย ลภามิ จีวรปิณฺฑปาตเสนาสนคิลานปจฺจยเภสชฺชปริกฺขารํ.\csromanspacer{} ยมฺปิ เม อญฺเญ ทาตุํ วา กาตุํ วา มญฺญนฺติ ตุเมฺห นิสฺสาย ตุเมฺห ปสฺสนฺตา.\csromanspacer{} ยมฺปิ เม โปราณํ มาตาเปตฺติกํ นามโคตฺตํ, ตมฺปิ เม อนฺตรหิตํ ตุเมฺหหิ อหํ ญายามิ “อสุกสฺส กุลุปโก, อสุกาย กุลุปโก”ติ, เอวํ อตฺตานํ นีจํ ฐเปนฺโต ปรํ อุจฺจํ ฐเปนฺโต ปฏิพทฺธจิตฺโต โหติ.}

\prose{กถํ อตฺตานํ อุจฺจํ ฐเปนฺโต ปรํ นีจํ ฐเปนฺโต ปฏิพทฺธจิตฺโต โหติ, อหํ ตุมฺหากํ พหูปกาโร, ตุเมฺห มํ อาคมฺม พุทฺธํ สรณํ คตา, ธมฺมํ สรณํ คตา, สํฆํ สรณํ คตา, ปาณาติปาตา ปฏิวิรตา, อทินฺนาทานา ปฏิวิรตา, กาเมสุมิจฺฉาจารา ปฏิวิรตา, มุสาวาทา ปฏิวิรตา, สุราเมรยมชฺชปมาทฏฺฐานา ปฏิวิรตา.\csromanspacer{} ตุมฺหากํ อหํ อุทฺเทสํ เทมิ, ปริปุจฺฉํ เทมิ, อุโปสถํ อาจิกฺขามิ, นวกมฺมํ อธิฏฺฐามิ.\csromanspacer{} อถ จ ปน ตุเมฺห มํ อุชฺฌิตฺวา อญฺเญ สกฺกโรถ ครุํ กโรถ มาเนถ ปูเชถาติ เอวํ อตฺตานํ อุจฺจํ ฐเปนฺโต ปรํ นีจํ ฐเปนฺโต ปฏิพทฺธจิตฺโต โหติ.}

\prose{\textbf{กุเล กุเล อปฺปฏิพทฺธจิตฺโต}ติ โส ปจฺเจกสมฺพุทฺโธ กุลปลิโพเธน อปฺปฏิพทฺธจิตฺโต โหติ, คณปลิโพเธน อปฺปฏิพทฺธจิตฺโต โหติ, อาวาสปลิโพเธน อปฺปฏิพทฺธจิตฺโต โหติ, จีวรปลิโพเธน อปฺปฏิพทฺธจิตฺโต โหติ, ปิณฺฑปาตปลิโพเธน อปฺปฏิพทฺธจิตฺโต โหติ,}

\prose{\csromanfolio{294}เสนาสนปลิโพเธน อปฺปฏิพทฺธจิตฺโต โหติ\textbf{,} คิลานปจฺจยเภสชฺชปริกฺขารปลิโพเธน อปฺปฏิพทฺธจิตฺโต โหตีติ กุเล กุเล อปฺปฏิพทฺธจิตฺโต โหติ\textbf{,} เอโก จเร ขคฺควิสาณกปฺโป.\csromanspacer{} เตนาห โส ปจฺเจกสมฺพุทฺโธ\csromandash{}}

\csromangathameasure{รเสสุ เคธํ อกรํ อโลโล\textbf{,} \\ อนญฺญโปสี สปทานจารี. \\ กุเล กุเล อปฺปฏิพทฺธจิตฺโต\textbf{,} \\ เอโก จเร ขคฺควิสาณกปฺโปติ.}
\csromangathagroup{\csromangathastanza{รเสสุ เคธํ อกรํ อโลโล\textbf{,} \\ อนญฺญโปสี สปทานจารี. \\ กุเล กุเล อปฺปฏิพทฺธจิตฺโต\textbf{,} \\ เอโก จเร ขคฺควิสาณกปฺโปติ.}}

\vspace{\tipitakaparskip}
\csromancenter{ตติโย วคฺโค.}
\chapterhead{\textbf{จตุตฺถวคฺค}}
\csromanhangingbegin
\hangingheaditem{152}{\textbf{ปหาย ปญฺจาวรณานิ เจตโส,}}
\hangingbody{\textbf{อุปกฺกิเลเส พฺยปนุชฺช สพฺเพ.}}
\hangingbody{\textbf{อนิสฺสิโต เฉตฺว1} \textbf{สิเนหโทสํ2,}}
\hangingbody{เอโก จเร ขคฺควิสาณกปฺโป.}
\csromanhangingend

\prose{\textbf{ปหาย ปญฺจาวรณานิ เจตโส}ติ โส ปจฺเจกสมฺพุทฺโธ กามจฺฉนฺทนีวรณํ ปหาย ปชหิตฺวา วิโนเทตฺวา พฺยนฺตีกริตฺวา อนภาวํ คเมตฺวา\textbf{,} พฺยาปาทนีวรณํ.\csromanspacer{} ถินมิทฺธนีวรณํ.\csromanspacer{} อุทฺธจฺจกุกฺกุจฺจนีวรณํ.\csromanspacer{} วิจิกิจฺฉานีวรณํ ปหาย ปชหิตฺวา วิโนเทตฺวา พฺยนฺตีกริตฺวา อนภาวํ คเมตฺวา\textbf{,} วิวิจฺเจว กาเมหิ\textbf{,} วิวิจฺจ อกุสเลหิ ธมฺเมหิ สวิตกฺกํ สวิจารํ วิเวกชํ ปีติสุขํ ปฐมํ ฌานํ อุปสมฺปชฺช วิหรตีติ ปหาย ปญฺจาวรณานิ เจตโส.}

\prose{\textbf{อุปกฺกิเลเส พฺยปนุชฺช สพฺเพ}ติ ราโค จิตฺตสฺส อุปกฺกิเลโส\textbf{,} โทโส จิตฺตสฺส อุปกฺกิเลโส\textbf{,} โมโห จิตฺตสฺส อุปกฺกิเลโส\textbf{,} โกโธ.\csromanspacer{} อุปนาโห ฯเปฯ สพฺพากุสลาภิสงฺขารา จิตฺตสฺส อุปกฺกิเลสา.\csromanspacer{} \textbf{อุปกฺกิเลเส พฺยปนุชฺช สพฺเพ}ติ สพฺเพ จิตฺตสฺส อุปกฺกิเลเส พฺยปนุชฺช ปนุทิตฺวา ปชหิตฺวา วิโนเทตฺวา พฺยนฺตีกริตฺวา อนภาวํ คเมตฺวาติ อุปกฺกิเลเส พฺยปนุชฺช สพฺเพ.}

\csromanheader{2}{19. ขคฺควิสาณสุตฺตนิทฺเทส}
\prose{\csromanfolio{295}\textbf{อนิสฺสิโต เฉตฺว สิเนหโทสนฺ}ติ \textbf{อนิสฺสิโต}ติ เทฺว นิสฺสยา ตณฺหานิสฺสโย จ ทิฏฺฐินิสฺสโย จ ฯเปฯ อยํ ตณฺหานิสฺสโย ฯเปฯ อยํ ทิฏฺฐินิสฺสโย.\csromanspacer{} \textbf{สิเนโห}ติ เทฺว สฺเนหา ตณฺหาสฺเนโห จ ทิฏฺฐิสฺเนโห จ ฯเปฯ อยํ ตณฺหาสฺเนโห ฯเปฯ อยํ ทิฏฺฐิสฺเนโห.\csromanspacer{} \textbf{โทโส}ติ โย จิตฺตสฺส อาฆาโต ปฏิฆาโต ปฏิฆํ ปฏิวิโรโธ โกโป ปโกโป สมฺปโกโป โทโส ปโทโส สมฺปโทโส จิตฺตสฺส พฺยาปตฺติ มโนปโทโส โกโธ กุชฺฌนา กุชฺฌิตตฺตํ โทโส ทุสฺสนา ทุสฺสิตตฺตํ พฺยาปตฺติ พฺยาปชฺชนา พฺยาปชฺชิตตฺตํ จณฺฑิกฺกํ อสุโรโป อนตฺตมนตา จิตฺตสฺส.\csromanspacer{} \textbf{อนิสฺสิโต เฉตฺว สิเนหโทสนฺ}ติ โส ปจฺเจกสมฺพุทฺโธ ตณฺหาสฺเนหญฺจ ทิฏฺฐิสฺเนหญฺจ โทสญฺจ เฉตฺวา อุจฺฉินฺทิตฺวา สมุจฺฉินฺทิตฺวา ปชหิตฺวา วิโนเทตฺวา พฺยนฺตีกริตฺวา อนภาวํ คเมตฺวา จกฺขุํ \textbf{อนิสฺสิโต}, โสตํ \textbf{อนิสฺสิโต} ฯเปฯ.\csromanspacer{} ทิฏฺฐสุตมุตวิญฺญาตพฺเพ ธมฺเม \textbf{อนิสฺสิโต} อนลฺลีโน อนุปคโต อนชฺโฌสิโต อนธิมุตฺโต นิกฺขนฺโต นิสฺสโฏ วิปฺปมุตฺโต วิสญฺญุตฺโต วิมริยาทิกเตน เจตสา วิหรตีติ \textbf{อนิสฺสิโต} เฉตฺว สิเนหโทสํ, เอโก จเร ขคฺควิสาณกปฺโป.\csromanspacer{} เตนาห โส ปจฺเจกสมฺพุทฺโธ\csromandash{}}

\csromangathameasure{ปหาย ปญฺจาวรณานิ เจตโส, \\ อุปกฺกิเลเส พฺยปนุชฺช สพฺเพ. \\ อนิสฺสิโต เฉตฺว สิเนหโทสํ, \\ เอโก จเร ขคฺควิสาณกปฺโปติ.}
\csromangathagroup{\csromangathastanza{ปหาย ปญฺจาวรณานิ เจตโส, \\ อุปกฺกิเลเส พฺยปนุชฺช สพฺเพ. \\ อนิสฺสิโต เฉตฺว สิเนหโทสํ, \\ เอโก จเร ขคฺควิสาณกปฺโปติ.}}

\csromanhangingbegin
\hangingheaditem{153}{วิปิฏฺฐิกตฺวาน สุขํ ทุขญฺจ,}
\hangingbody{\textbf{ปุพฺเพว จ โสมนสฺสโทมนสฺสํ.}}
\hangingbody{\textbf{ลทฺธานุเปกฺขํ สมถํ วิสุทฺธํ,}}
\hangingbody{เอโก จเร ขคฺควิสาณกปฺโป.}
\csromanhangingend

\prose{\textbf{วิปิฏฺฐิกตฺวาน สุขํ ทุขญฺจ, ปุพฺเพว จ โสมนสฺสโทมนสฺสนฺ}ติ โส ปจฺเจกสมฺพุทฺโธ สุขสฺส จ ปหานา ทุกฺขสฺส จ ปหานา ปุพฺเพว โสมนสฺสโทมนสฺสานํ อตฺถงฺคมา อทุกฺขมสุขํ อุเปกฺขาสติปาริสุทฺธิํ จตุตฺถํ ฌานํ อุปสมฺปชฺช วิหรตีติ วิปิฏฺฐิกตฺวาน สุขํ ทุขญฺจ, ปุพฺเพว จ โสมนสฺสโทมนสฺสํ.}

\prose{\csromanfolio{296}\textbf{ลทฺธานุเปกฺขํ สมถํ วิสุทฺธนฺ}ติ \textbf{อุเปกฺขา}ติ ยา จตุตฺถชฺฌาเน \textbf{อุเปกฺขา} อุเปกฺขนา อชฺฌุเปกฺขนา จิตฺตสมตา จิตฺตปฺปสฺสทฺธตา\footnote{จิตฺตวิสฏตา (ก) ปสฺส ขุ 7. 402 ปิฏฺเฐ.} มชฺฌตฺตตา จิตฺตสฺส.\csromanspacer{} \textbf{สมโถ}ติ ยา จิตฺตสฺส ฐิติ สณฺฐิติ อวฏฺฐิติ อวิสาหาโร\footnote{อวิสํหาโร (ก) ปสฺส อภิ 1. 19 ปิฏฺเฐ.} อวิกฺเขโป อวิสาหฏมานสตา\footnote{อวิสํหฏมานสตา (ก)} สมโถ สมาธินฺทฺริยํ สมาธิพลํ สมฺมาสมาธิ, จตุตฺถชฺฌาเน \textbf{อุเปกฺขา} จ สมโถ จ สุทฺธา โหนฺติ วิสุทฺธา ปริโยทาตา อนงฺคณา วิคตูปกฺกิเลสา มุทุภูตา กมฺมนิยา ฐิตา อาเนญฺชปฺปตฺตา.\csromanspacer{} \textbf{ลทฺธานุเปกฺขํ สมถํ วิสุทฺธนฺ}ติ จตุตฺถชฺฌานํ อุเปกฺขญฺจ สมถญฺจ ลทฺธา ลภิตฺวา วินฺทิตฺวา ปฏิลภิตฺวาติ ลทฺธานุเปกฺขํ สมถํ วิสุทฺธํ, เอโก จเร ขคฺควิสาณกปฺโป.\csromanspacer{} เตนาห โส ปจฺเจกสมฺพุทฺโธ\csromandash{}}

\csromangathameasure{วิปิฏฺฐิกตฺวาน สุขํ ทุขญฺจ, \\ ปุพฺเพว จ โสมนสฺสโทมนสฺสํ.}
\csromangathagroup{\csromangathastanza{วิปิฏฺฐิกตฺวาน สุขํ ทุขญฺจ, \\ ปุพฺเพว จ โสมนสฺสโทมนสฺสํ.}}

\prose{\textbf{ลทฺธานุเปกฺขํ สมถํ วิสุทฺธ}ํ,}

\prose{เอโก จเร ขคฺควิสาณกปฺโปติ.}

\proseitem{154}{อารทฺธวิริโย ปรมตฺถปตฺติยา,}

\prose{\textbf{อลีนจิตฺโต อกุสีตวุตฺติ.}}

\prose{\textbf{ทฬฺหนิกฺกโม ถามพลูปปนฺโน,}}

\csromanheader{2}{\textbf{เอโก จเร ขคฺควิสาณกปฺโป.(3)}}
\prose{\textbf{อารทฺธวิริโย ปรมตฺตปตฺติยา}ติ \textbf{ปรมตฺถํ} วุจฺจติ อมตํ นิพฺพานํ.\csromanspacer{} โย โส สพฺพสงฺขารสมโถ สพฺพูปธิปฏินิสฺสคฺโค ตณฺหกฺขโย วิราโค นิโรโธ นิพฺพานํ.\csromanspacer{} ปรมตฺถสฺส ปตฺติยา ลาภาย ปฏิลาภาย อธิคมาย ผสฺสนาย สจฺฉิกิริยาย อารทฺธวีริโย วิหรติ, อกุสลานํ ธมฺมานํ ปหานาย กุสลานํ ธมฺมานํ สมฺปทาย ถามวา\footnote{ถามสา (ก)} ทฬฺหปรกฺกโม อนิกฺขิตฺตธุโร กุสเลสุ ธมฺเมสูติ อารทฺธวิริโย ปรมตฺถปตฺติยา.}

\prose{\textbf{อลีนจิตฺโต อกุสีตวุตฺตี}ติ โส ปจฺเจกสมฺพุทฺโธ อนุปฺปนฺนานํ ปาปกานํ อกุสลานํ ธมฺมานํ อนุปฺปาทาย ฉนฺทํ ชเนติ วายมติ, วีริยํ อารภติ, จิตฺตํ ปคฺคณฺหาติ ปทหติ.\csromanspacer{} อุปฺปนฺนานํ ปาปกานํ อกุสลานํ ธมฺมานํ ปหานาย ฯเปฯ.\csromanspacer{} อนุปฺปนฺนานํ กุสลานํ ธมฺมานํ อุปฺปาทาย ฯเปฯ.\csromanspacer{} อุปฺปนฺนานํ}

\csromanheader{2}{19. ขคฺควิสาณสุตฺตนิทฺเทส}
\prose{\csromanfolio{297}กุสลานํ ธมฺมานํ ฐิติยา อสมฺโมสาย ภิยฺโยภาวาย เวปุลฺลาย ภาวนาย ปาริปูริยา\footnote{ภาวนาปาริปูริยา (สฺยา)} ฉนฺทํ ชเนติ วายมติ, วีริยํ อารภติ, จิตฺตํ ปคฺคณฺหาติ ปทหตีติ, เอวํ อลีนจิตฺโต อกุสีตวุตฺติ.}

\prose{อถ วา “กามํ ตโจ จ นฺหารุ จ อฏฺฐิ จ อวสิสฺสตุ\footnote{อวสุสฺสตุ (สฺยา) ม 2. 146 ปิฏฺเฐ; สํ 1. 267 ปิฏฺเฐ จ ปสฺสิตพฺพํ.}, สรีเร อุปสุสฺสตุ มํสโลหิตํ.\csromanspacer{} ยํ ตํ ปุริสถาเมน ปุริสพเลน ปุริสวีริเยน ปุริสปรกฺกเมน ปตฺตพฺพํ, น ตํ อปาปุณิตฺวา วีริยสฺส สณฺฐานํ\footnote{วิริยสฺส ฐานํ (สฺยา)} ภวิสฺสตี”ติ จิตฺตํ ปคฺคณฺหาติ ปทหติ, เอวมฺปิ อลีนจิตฺโต อกุสีตวุตฺติ.}

\csromangathameasure{“นาสิสฺสํ น ปิวิสฺสามิ, วิหารโต น นิกฺขเม. \\ นปิ ปสฺสํ นิปาเตสฺสํ, ตณฺหาสลฺเล อนูหเต”ติ.}
\csromangathagroup{\csromangathastanza{“นาสิสฺสํ น ปิวิสฺสามิ, วิหารโต น นิกฺขเม. \\ นปิ ปสฺสํ นิปาเตสฺสํ, ตณฺหาสลฺเล อนูหเต”ติ.}}

\prose{จิตฺตํ ปคฺคณฺหาติ ปทหติ, เอวมฺปิ อลีนจิตฺโต อกุสีตวุตฺติ.}

\prose{“น ตาวาหํ อิมํ ปลฺลงฺกํ ภินฺทิสฺสามิ, ยาว เม น อนุปาทาย อาสเวหิ จิตฺตํ วิมุจฺจิสฺสตี”ติ จิตฺตํ ปคฺคณฺหาติ ปทหติ, เอวมฺปิ อลีนจิตฺโต อกุสีตวุตฺติ.อ ตาวาหํ อิมมฺหา อาสนา วุฏฺฐหิสฺสามิ, ยาว เม น อนุปาทาย อาสเวหิ จิตฺตํ วิมุจฺจิสฺสาตี”ติ จิตฺตํ ปคฺคณฺหาติ ปทหติ, เอวมฺปิ อลีนจิตฺโต อกุสีตวุตฺติ.}

\prose{“น ตาวาหํ อิมมฺหา จงฺกมา โอโรหิสฺสามิ.\csromanspacer{} วิหารา นิกฺขมิสฺสามิ.\csromanspacer{} อฑฺฒโยคา นิกฺขมิสฺสามิ.\csromanspacer{} ปาสาทา นิกฺขมิสฺสามิ.\csromanspacer{} หมฺมิยา นิกฺขมิสฺสามิ.\csromanspacer{} คุหาย นิกฺขมิสฺสามิ.\csromanspacer{} เลณา นิกฺขมิสฺสามิ.\csromanspacer{} กุฏิยา นิกฺขมิสฺสามิ.\csromanspacer{} กูฏาคารา นิกฺขมิสฺสามิ.\csromanspacer{} อฏฺฏา นิกฺขมิสฺสามิ.\csromanspacer{} มาฬา นิกฺขมิสฺสามิ.\csromanspacer{} อุทฺทณฺฑา นิกฺขมิสฺสามิ.\csromanspacer{} อุปฏฺฐานสาลาย นิกฺขมิสฺสามิ.\csromanspacer{} มณฺฑปา นิกฺขมิสฺสามิ.\csromanspacer{} รุกฺขมูลา นิกฺขมิสฺสามิ, ยาว เม น อนุปาทาย อาสเวหิ จิตฺตํ วิมุจฺจิสฺสตี”ติ จิตฺตํ ปคฺคณฺหาติ ปทหติ, เอวมฺปิ อลีนจิตฺโต อกุสีตวุตฺติ.}

\prose{“อิมสฺมิํ เยว ปุพฺพณฺหสมเย\footnote{ปุพฺพณฺหสมยํ (ก) ขุ 7. 51 ปิฏฺเฐปิ.} อริยธมฺมํ อาหริสฺสามิ สมาหริสฺสามิ อธิคจฺฉิสฺสามิ ผสฺสยิสฺสามิ สจฺฉิกริสฺสามีติ จิตฺตํ ปคฺคณฺหาติ ปทหติ, เอวมฺปิ อลีนจิตฺโต อกุสีตวุตฺติ.\csromanspacer{} อิมสฺมิํ เยว มชฺฌนฺหิกสมเย ฯเปฯ.}

\prose{\csromanfolio{298}สายนฺหสมเย.\csromanspacer{} ปุเรภตฺตํ.\csromanspacer{} ปจฺฉาภตฺตํ.\csromanspacer{} ปุริมยามํ.\csromanspacer{} มชฺฌิมยามํ.\csromanspacer{} ปจฺฉิมยามํ.\csromanspacer{} กาเฬ.\csromanspacer{} ชุเณฺห.\csromanspacer{} วสฺเส.\csromanspacer{} เหมนฺเต.\csromanspacer{} คิเมฺห.\csromanspacer{} ปุริเม วโยขนฺเธ.\csromanspacer{} มชฺฌิเม วโยขนฺเธ.\csromanspacer{} ปจฺฉิเม วโยขนฺเธ อริยธมฺมํ อาหริสฺสามิ สมาหริสฺสามิ อธิคจฺฉิสฺสามิ ผสฺสยิสฺสามิ สจฺฉิกริสฺสามี”ติ จิตฺตํ ปคฺคณฺหาติ ปทหติ, เอวมฺปิ อลีนจิตฺโต อกุสีตวุตฺติ.}

\prose{\textbf{ทฬฺหนิกฺกโม ถามพลูปปนฺโน}ติ โส ปจฺเจกสมฺพุทฺโธ ทฬฺหสมาทาโน อโหสิ กุสเลสุ ธมฺเมสุ อวฏฺฐิตสมาทาโน กายสุจริเต วจีสุจริเต มโนสุจริเต ทานสํวิภาเค สีลสมาทาเน อุโปสถุปวาเส มตฺเตยฺยตาย\footnote{เมตฺเตยฺยตาย (สฺยา, ก) โมคฺคลฺลานพฺยากรเณ 4. 39 สุตฺตํ โอโลเกตพฺพํ.} เปตฺเตยฺยตาย สามญฺญตาย พฺรหฺมญฺญตาย กุเลเชฏฺฐาปจายิตาย อญฺญตรญฺญตเรสุ อธิกุสเลสุ ธมฺเมสูติ ทฬฺหนิกฺกโม.\csromanspacer{} \textbf{ถามพลูปปนฺโน}ติ โส ปจฺเจกสมฺพุทฺโธ ถาเมน จ พเลน จ วีริเยน จ ปรกฺกเมน จ ปญฺญาย จ อุเปโต โหติ สมุเปโต อุปาคโต สมุปาคโต อุปปนฺโน สมุปปนฺโน สมนฺนาคโตติ ทฬฺหนิกฺกโม ถามพลูปปนฺโน, เอโก จเร ขคฺควิสาณกปฺโป.\csromanspacer{} เตนาห โส ปจฺเจกสมฺพุทฺโธ\csromandash{}}

\csromangathameasure{อารทฺธวิริโย ปรมตฺถปตฺติยา, \\ อลีนจิตฺโต อกุสีตวุตฺติ.}
\csromangathagroup{\csromangathastanza{อารทฺธวิริโย ปรมตฺถปตฺติยา, \\ อลีนจิตฺโต อกุสีตวุตฺติ.}}

\prose{\textbf{ทฬฺหนิกฺกโม ถามพลูปปนฺโน},}

\prose{เอโก จเร ขคฺควิสาณกปฺโปติ.}

\csromanhangingbegin
\hangingheaditem{155}{\textbf{ปฏิสลฺลานํ ฌานมริญฺจมาโน},}
\hangingbody{\textbf{ธมฺเมสุ นิจฺจํ อนุธมฺมจารี.}}
\hangingbody{\textbf{อาทีนวํ สมฺมสิตา ภเวสุ,}}
\hangingbody{เอโก จเร ขคฺควิสาณกปฺโป.}
\csromanhangingend

\prose{\textbf{ปฏิสลฺลานํ ฌานมริญฺจมาโน}ติ โส ปจฺเจกสมฺพุทฺโธ ปฏิสลฺลานาราโม โหติ ปฏิสลฺลานรโต อชฺฌตฺตํ เจโตสมถมนุยุตฺโต อนิรากตชฺฌาโน วิปสฺสนาย สมนฺนาคโต พฺรูเหตา สุญฺญาคารํ ฌายี ฌานรโต เอกตฺตมนุยุตฺโต สทตฺถครุโกติ ปฏิสลฺลานํ.\csromanspacer{} \textbf{ฌานมริญฺจมาโน}ติ โส ปจฺเจกสมฺพุทฺโธ ทฺวีหิ การเณหิ ฌานํ น ริญฺจติ อนุปฺปนฺนสฺส วา ปฐมสฺส ฌานสฺส อุปฺปาทาย ยุตฺโต ปยุตฺโต}

\csromanheader{2}{19. ขคฺควิสาณสุตฺตนิทฺเทส}
\prose{\csromanfolio{299}สํยุตฺโต อายุตฺโต สมายุตฺโต, อนุปฺปนฺนสฺส วา ทุติยสฺส ฌานสฺส.\csromanspacer{} อนุปฺปนฺนสฺส วา ตติยสฺส ฌานสฺส.\csromanspacer{} อนุปฺปนฺนสฺส วา จตุตฺถสฺส ฌานสฺส อุปฺปาทาย ยุตฺโต ปยุตฺโต สํยุตฺโต อายุตฺโต สมายุตฺโตติ เอวมฺปิ ฌานํ น ริญฺจติ.}

\prose{อถ วา อุปฺปนฺนํ วา ปฐมํ ฌานํ อาเสวติ ภาเวติ พหุลีกโรติ, อุปฺปนฺนํ วา ทุติยํ ฌานํ.\csromanspacer{} อุปฺปนฺนํ วา ตติยํ ฌานํ.\csromanspacer{} อุปฺปนฺนํ วา จตุตฺถํ ฌานํ อาเสวติ ภาเวติ พหุลีกโรติ, เอวมฺปิ ฌานํ น ริญฺจตีติ ปฏิสลฺลานํ ฌานมริญฺจมาโน.}

\prose{\textbf{ธมฺเมสุ นิจฺจํ อนุธมฺมจารี}ติ ธมฺมา วุจฺจนฺติ จตฺตาโร สติปฏฺฐานา ฯเปฯ อริโย อฏฺฐงฺคิโก มคฺโค.\csromanspacer{} กตเม อนุธมฺมา, สมฺมาปฏิปทา อปจฺจนีกปฏิปทา อนฺวตฺถปฏิปทา ธมฺมานุธมฺมปฏิปทา สีเลสุ ปริปูรการิตา อินฺทฺริเยสุ คุตฺตทฺวารตา โภชเน มตฺตญฺญุตา ชาคริยานุโยโค สติสมฺปชญฺญํ, อิเม วุจฺจนฺติ อนุธมฺมา.\csromanspacer{} \textbf{ธมฺเมสุ นิจฺจํ อนุธมฺมจารี}ติ ธมฺเมสุ นิจฺจกาลํ ธุวกาลํ สตตํ สมิตํ อโวกิณฺณํ โปงฺขานุโปงฺขํ อุทกูมิกชาตํ\footnote{อุทกุมฺมิชาตํ (สฺยา), อุทกุมฺมิกชาตํ (ก) 43 ปิฏฺเฐปิ.} อวีจิ สนฺตติ สหิตํ ผสฺสิตํ ปุเรภตฺตํ ปจฺฉาภตฺตํ ปุริมยามํ มชฺฌิมยามํ ปจฺฉิมยามํ กาเฬ ชุเณฺห วสฺเส เหมนฺเต คิเมฺห ปุริเม วโยขนฺเธ มชฺฌิเม วโยขนฺเธ ปจฺฉิเม วโยขนฺเธ จรติ วิหรติ อิริยติ วตฺเตติ ปาเลติ ยเปติ ยาเปตีติ ธมฺเมสุ นิจฺจํ อนุธมฺมจารี.}

\prose{\textbf{อาทีนวํ สมฺมสิตา ภเวสู}ติ “สพฺเพ สงฺขารา อนิจฺจา”ติ อาทีนวํ สมฺมสิตา ภเวสุ, “สพฺเพ สงฺขารา ทุกฺขา”ติ. “สพฺเพ ธมฺมา อนตฺตา”ติ ฯเปฯ. “ยํ กิญฺจิ สมุทยธมฺมํ, สพฺพํ ตํ นิโรธธมฺมนฺ”ติ อาทีนวํ สมฺมสิตา ภเวสูติ อาทีนวํ สมฺมสิตา ภเวสุ, เอโก จเร ขคฺควิสาณกปฺโป.\csromanspacer{} เตนาห โส ปจฺเจกสมฺพุทฺโธ\csromandash{}}

\csromangathameasure{ปฏิสลฺลานํ ฌานมริญฺจมาโน,}
\csromangathagroup{\csromangathastanza{ปฏิสลฺลานํ ฌานมริญฺจมาโน,}}

\prose{\textbf{ธมฺเมสุ นิจฺจํ อนุธมฺมจารี}.}

\csromangathameasure{อาทีนวํ สมฺมสิตา ภเวสุ, \\ เอโก จเร ขคฺควิสาณกปฺโปติ.}
\csromangathagroup{\csromangathastanza{อาทีนวํ สมฺมสิตา ภเวสุ, \\ เอโก จเร ขคฺควิสาณกปฺโปติ.}}

\csromanhangingbegin
\hangingheaditem{156}{\csromanfolio{300}\textbf{ตณฺหกฺขยํ ปตฺถยมปฺปมตฺโต},}
\hangingbody{\textbf{อเนฬมูโค1} \textbf{สุตวา สตีมา.}}
\hangingbody{\textbf{สงฺขาตธมฺโม นิยโต ปธานวา,}}
\hangingbody{เอโก จเร ขคฺควิสาณกปฺโป.}
\csromanhangingend

\prose{\textbf{ตณฺหกฺขยํ ปตฺถยมปฺปมตฺโต}ติ \textbf{ตณฺหา}ติ รูป\textbf{ตณฺหา} ฯเปฯ ธมฺม\textbf{ตณฺหา}.\csromanspacer{} \textbf{ตณฺหกฺขยนฺ}ติ ราคกฺขยํ โทสกฺขยํ โมหกฺขยํ คติกฺขยํ อุปปตฺติกฺขยํ ปฏิสนฺธิกฺขยํ ภวกฺขยํ สํสารกฺขยํ วฏฺฏกฺขยํ ปตฺถยนฺโต อิจฺฉนฺโต สาทิยนฺโต ปิหยนฺโต อภิชปฺปนฺโตติ ตณฺหกฺขยํ ปตฺถยํ.\csromanspacer{} \textbf{อปฺปมตฺโต}ติ โส ปจฺเจกสมฺพุทฺโธ สกฺกจฺจการี สาตจฺจการี ฯเปฯ อปฺปมาโท กุสเลสุ ธมฺเมสูติ ตณฺหกฺขยํ ปตฺถยมปฺปมตฺโต.}

\prose{\textbf{อเนฬมูโค สุตวา สตีมา}ติ \textbf{อเนฬมูโค}ติ โส ปจฺเจกสมฺพุทฺโธ ปณฺฑิโต ปญฺญวา พุทฺธิมา ญาณี วิภาวี เมธาวี.\csromanspacer{} \textbf{สุตวา}ติ โส ปจฺเจกสมฺพุทฺโธ พหุสฺสุโต โหติ สุตธโร สุตสนฺนิจฺจโย, เย เต ธมฺมา อาทิกลฺยาณา มชฺเฌกลฺยาณา ปริโยสานกลฺยาณา สาตฺถํ สพฺยญฺชนํ เกวลปริปุณฺณํ ปริสุทฺธํ พฺรหฺมจริยํ อภิวทนฺติ, ตถารูปาสฺส ธมฺมา พหุสฺสุตา โหนฺติ ธาตา วจสา ปริจิตา มนสานุเปกฺขิตา ทิฏฺฐิยา สุปฺปฏิวิทฺธา.\csromanspacer{} \textbf{สตีมา}ติ โส ปจฺเจกสมฺพุทฺโธ สติมา โหติ ปรเมน สติเนปกฺเกน สมนฺนาคตตฺตา จิรกตมฺปิ จิรภาสิตมฺปิ สริตา อนุสฺสริตาติ \textbf{อเนฬมูโค} \textbf{สุตวา สตีมา.}}

\prose{\textbf{สงฺขาตธมฺโม นิยโต ปธานวา}ติ สงฺขาตธมฺโม วุจฺจติ ญาณํ.\csromanspacer{} ยา ปญฺญา ปชานนา ฯเปฯ อโมโห ธมฺมวิจโย สมฺมาทิฏฺฐิ.\csromanspacer{} \textbf{สงฺขาตธมฺโม}ติ โส ปจฺเจกสมฺพุทฺโธ สงฺขาตธมฺโม ญาตธมฺโม ตุลิตธมฺโม ตีริตธมฺโม วิภูตธมฺโม วิภาวิตธมฺโม “สพฺเพ สงฺขารา อนิจฺจา”ติ สงฺขาตธมฺโม ฯเปฯ. “ยํ กิญฺจิ สมุทยธมฺมํ, สพฺพํ ตํ นิโรธธมฺมนฺ”ติ สงฺขาตธมฺโม ญาตธมฺโม ตุลิตธมฺโม ตีริตธมฺโม วิภูตธมฺโม วิภาวิตธมฺโม.\csromanspacer{} อถ วา ตสฺส ปจฺเจกสมฺพุทฺธสฺส จ ขนฺธา สํขิตฺตา, ธาตุโย สํขิตฺตา, อายตนานิ สํขิตฺตานิ, คติโย สํขิตฺตา, อุปปตฺติโย สํขิตฺตา, ปฏิสนฺธิโย สํขิตฺตา, ภวา สํขิตฺตา, สํสารา สํขิตฺตา, วฏฺฏา สํขิตฺตา.\csromanspacer{} อถ วา โส ปจฺเจกสมฺพุทฺโธ ขนฺธปริยนฺเต ฐิโต,}

\csromanheader{2}{19. ขคฺควิสาณสุตฺตนิทฺเทส}
\prose{\csromanfolio{301}ธาตุปริยนฺเต ฐิโต\textbf{,} อายตนปริยนฺเต ฐิโต\textbf{,} คติปริยนฺเต ฐิโต\textbf{,} อุปปตฺติปริยนฺเต ฐิโต\textbf{,} ปฏิสนฺธิปริยนฺเต ฐิโต\textbf{,} ภวปริยนฺเต ฐิโต\textbf{,} สํสารปริยนฺเต ฐิโต\textbf{,} วฏฺฏปริยนฺเต ฐิโต\textbf{,} สงฺขารปริยนฺเต ฐิโต\textbf{,} อนฺติมภเว\footnote{อนฺติเม ภเว (ก) 86 ปิฏฺเฐปิ.} ฐิโต\textbf{,} อนฺติมสมุสฺสเย ฐิโต อนฺติมเทหธโร ปจฺเจกสมฺพุทฺโธ.}

\csromangathameasure{ตสฺสายํ ปจฺฉิมโก ภโว, จริโมยํ สมุสฺสโย. \\ ชาติมรณสํสาโร, นตฺถิ ตสฺส ปุนพฺภโวติ.}
\csromangathagroup{\csromangathastanza{ตสฺสายํ ปจฺฉิมโก ภโว, จริโมยํ สมุสฺสโย. \\ ชาติมรณสํสาโร\footnote{ชาติชรามรณสํสาโร (สฺยา) เอวมีทิเสสุ ฐาเนสุ.}, นตฺถิ ตสฺส ปุนพฺภโวติ.}}

\prose{ตํการณา ปจฺเจกสมฺพุทฺโธ สงฺขาตธมฺโม.\csromanspacer{} \textbf{นิยโต}ติ นิยามา วุจฺจนฺติ จตฺตาโร อริยมคฺคา.\csromanspacer{} จตูหิ อริยมคฺเคหิ สมนฺนาคโตติ นิยโต\textbf{,} นิยามํ ปตฺโต สมฺปตฺโต อธิคโต ผสฺสิโต สจฺฉิกโต ปตฺโต นิยามํ.\csromanspacer{} \textbf{ปธานวา}ติ ปธานํ วุจฺจติ วีริยํ\textbf{,} โส เจตโส วีริยารมฺโภ นิกฺกโม ปรกฺกโม อุยฺยาโม วายาโม อุสฺสาโห อุสฺโสฬฺหี ถาโม ธิติ อสิถิลปรกฺกโม อนิกฺขิตฺตจฺฉนฺทตา อนิกฺขิตฺตธุรตา ธุรสมฺปคฺคาโห วีริยํ วีริยินฺทฺริยํ วีริยพลํ สมฺมาวายาโม.\csromanspacer{} โส ปจฺเจกสมฺพุทฺโธ อิมินา ปธาเนน อุเปโต สมุเปโต อุปาคโต สมุปาคโต อุปปนฺโน สมุปปนฺโน สมนฺนาคโต\textbf{,} ตสฺมา โส ปจฺเจกสมฺพุทฺโธ ปธานวาติ สงฺขาตธมฺโม นิยโต ปธานวา\textbf{,} เอโก จเร ขคฺควิสาณกปฺโป.\csromanspacer{} เตนาห โส ปจฺเจกสมฺพุทฺโธ\csromandash{}}

\csromangathameasure{ตณฺหกฺขยํ ปตฺถยมปฺปมตฺโต\textbf{,} \\ อเนฬมูโค สุตวา สตีมา. \\ สงฺขาตธมฺโม นิยโต ปธานวา\textbf{,} \\ เอโก จเร ขคฺควิสาณกปฺโปติ.}
\csromangathagroup{\csromangathastanza{ตณฺหกฺขยํ ปตฺถยมปฺปมตฺโต\textbf{,} \\ อเนฬมูโค สุตวา สตีมา. \\ สงฺขาตธมฺโม นิยโต ปธานวา\textbf{,} \\ เอโก จเร ขคฺควิสาณกปฺโปติ.}}

\csromanhangingbegin
\hangingheaditem{157}{\textbf{สีโหว สทฺเทสุ อสนฺตสนฺโต,}}
\hangingbody{\textbf{วาโตว ชาลมฺหิ อสชฺชมาโน.}}
\hangingbody{\textbf{ปทุมํว โตเยน อลิมฺปมาโน3,}}
\hangingbody{เอโก จเร ขคฺควิสาณกปฺโป.}
\csromanhangingend

\prose{\textbf{สีโหว สทฺเทสุ อสนฺตสนฺโต}ติ ยถา สีโห มิคราชา สทฺเทสุ อสนฺตาสี อปริสนฺตาสี อนุตฺราสี อนุพฺพิคฺโค อนุสฺสงฺกี\footnote{อนุสฺสุกี (สฺยา)} \csromanfolio{302}อนุตฺราโส อภีรู อจฺฉมฺภี อนุตฺราสี อปลายี.\csromanspacer{} ปจฺเจกสมฺพุทฺโธปิ สทฺเทสุ อสนฺตาสี อปริสนฺตาสี อนุตฺราสี อนุพฺพิคฺโค อนุสฺสงฺกี อนุตฺราโส อภีรู อจฺฉมฺภี อนุตฺราสี อปลายี ปหีนภยเภรโว วิคตโลมหํโส วิหรตีติ สีโหว สทฺเทสุ อสนฺตสนฺโต.}

\prose{\textbf{วาโตว ชาลมฺหิ อสชฺชมาโน}ติ \textbf{วาโต}ติ ปุรตฺถิมาวาตา ปจฺฉิมา วาตา อุตฺตารา วาตา ทกฺขิณา วาตา สรชา วาตา อรชา วาตา สีตา วาตา อุณฺหา วาตา ปริตฺตา วาตา อธิมตฺตา วาตา เวรมฺภวาตา ปกฺขวาตา\footnote{ปกฺขิวาตา (สฺยา) 109, 266 ปิฏฺเฐ, นตฺถิ ปาฐนานตฺตํ.} สุปณฺณวาตา ตาลปณฺณวาตา วิธูปนวาตา.\csromanspacer{} ชาลํ วุจฺจติ สุตฺตชาลํ.\csromanspacer{} ยถา \textbf{วาโต} ชาลมฺหิ น สชฺชติ น คณฺหาติ น พชฺฌติ น ปลิพชฺฌติ, เอวเมว เทฺว ชาลา ตณฺหาชาลญฺจ ทิฏฺฐิชาลญฺจ ฯเปฯ อิทํ ตณฺหาชาลํ ฯเปฯ อิทํ ทิฏฺฐิชาลํ.\csromanspacer{} ตสฺส ปจฺเจกสมฺพุทฺธสฺส ตณฺหาชาลํ ปหีนํ, ทิฏฺฐิชาลํ ปฏินิสฺสฏฺฐํ, ตณฺหาชาลสฺส ปหีนตฺตา ทิฏฺฐิชาลสฺส ปฏินิสฺสฏฺฐตฺตา โส ปจฺเจกสมฺพุทฺโธ รูเป น สชฺชติ, สทฺเท น สชฺชติ ฯเปฯ ทิฏฺฐสุตมุตวิญฺญาตพฺเพสุ ธมฺเมสุ น สชฺชติ น คณฺหาติ น พชฺฌติ น ปลิพชฺฌติ นิกฺขนฺโต นิสฺสโฏ วิปฺปมุตฺโต วิสญฺญุตฺโต วิมริยาทิกเตน เจตสา วิหรตีติ \textbf{วาโต}ว ชาลมฺหิ อสชฺชมาโน.}

\prose{\textbf{ปทุมํว โตเยน อลิมฺปมาโน}ติ ปทุมํ วุจฺจติ ปทุมปุปฺผํ.\csromanspacer{} โตยํ วุจฺจติ อุทกํ.\csromanspacer{} ยถา ปทุมปุปฺผํ โตเยน น ลิปฺปติ น ปลิปฺปติ น อุปลิปฺปติ, อลิตฺตํ อปลิตฺตํ อนุปลิตฺตํ, เอวเมว เทฺว เลปา ตณฺหาเลโป จ ทิฏฺฐิเลโป จ ฯเปฯ อยํ ตณฺหาเลโป ฯเปฯ อยํ ทิฏฺฐิเลโป.\csromanspacer{} ตสฺส ปจฺเจกสมฺพุทฺธสฺส ตณฺหาเลโป ปหีโน, ทิฏฺฐิเลโป ปฏินิสฺสฏฺโฐ, ตณฺหาเลปสฺส ปหีนตฺตา ทิฏฺฐิเลปสฺส ปฏินิสฺสฏฺฐตฺตา โส ปจฺเจกสมฺพุทฺโธ รูเป น ลิมฺปติ, สทฺเท น ลิมฺปติ ฯเปฯ ทิฏฺฐสุตมุตวิญฺญาตพฺเพสุ ธมฺเมสุ น ลิมฺปติ น ปลิมฺปติ น อุปลิมฺปติ, อลิตฺโต อปลิตฺโต อนุปลิตฺโต นิกฺขนฺโต นิสฺสโฏ วิปฺปมุตฺโต วิสญฺญุตฺโต วิมริยาทิกเตน เจตสา วิหรตีติ ปทุมํว โตเยน อลิมฺปมาโน, เอโก จเร ขคฺควิสาณกปฺโป.\csromanspacer{} เตนาห โส ปจฺเจกสมฺพุทฺโธ\csromandash{}}

\csromanheader{2}{19. ขคฺควิสาณสุตฺตนิทฺเทส}
\csromangathameasure{สีโหว สทฺเทสุ อสนฺตสนฺโต, \\ วาโตว ชาลมฺหิ อสชฺชมาโน. \\ ปทุมํว โตเยน อลิมฺปมาโน, \\ เอโก จเร ขคฺควิสาณกปฺโปติ.}
\csromangathagroup{\csromangathastanza{\csromanfolio{303}สีโหว สทฺเทสุ อสนฺตสนฺโต, \\ วาโตว ชาลมฺหิ อสชฺชมาโน. \\ ปทุมํว โตเยน อลิมฺปมาโน, \\ เอโก จเร ขคฺควิสาณกปฺโปติ.}}

\csromanhangingbegin
\hangingheaditem{158}{สีโห ยถา ทาฐพลี ปสยฺห,}
\hangingbody{\textbf{ราชา มิคานํ อภิภุยฺย จารี.}}
\hangingbody{\textbf{เสเวถ ปนฺตานิ เสนาสนานิ,}}
\hangingbody{เอโก จเร ขคฺควิสาณกปฺโป.}
\csromanhangingend

\prose{\textbf{สีโห ยถา ทาฐพลี ปสยฺห, ราชา มิคานํ อภิภุยฺย จารี}ติ ยถา สีโห มิคราชา ทาฐพลี ทาฐาวุโธ สพฺเพ ติรจฺฉานคเต ปาเณ อภิภุยฺย อภิภวิตฺวา อชฺโฌตฺถริตฺวา ปริยาทิยิตฺวา มทฺทิตฺวา จรติ วิหรติ อิริยติ วตฺเตติ ปาเลติ ยเปติ ยาเปติ, ปจฺเจกสมฺพุทฺโธปิ ปญฺญาพลี ปญฺญาวุโธ สพฺพปาณภูเต ปุคฺคเล ปญฺญาย อภิภุยฺย อภิภวิตฺวา อชฺโฌตฺถริตฺวา ปริยาทิยิตฺวา มทฺทิตฺวา จรติ วิหรติ อิริยติ วตฺเตติ ปาเลติ ยเปติ ยาเปตีติ สีโห ยถา ทาฐพลี ปสยฺห, ราชา มิคานํ อภิภุยฺย จารี.}

\prose{\textbf{เสเวถ ปนฺตานิ เสนาสนานี}ติ ยถา สีโห มิคราชา อรญฺญวนมชฺโฌคาเหตฺวา จรติ วิหรติ อิริยติ วตฺเตติ ปาเลติ ยเปติ ยาเปติ, ปจฺเจกสมฺพุทฺโธปิ อรญฺญวนปตฺถานิ ปนฺตานิ เสนาสนานิ ปฏิเสวติ อปฺปสทฺทานิ อปฺปนิคฺโฆสานิ วิชนวาตานิ มนุสฺสราหสฺเสยฺยกานิ ปฏิสลฺลานสารุปฺปานิ, โส เอโก คจฺฉติ, เอโก ติฏฺฐติ, เอโก นิสีทติ, เอโก เสยฺยํ กปฺเปติ, เอโก คามํ ปิณฺฑาย ปวิสติ, เอโก ปฏิกฺกมติ, เอโก รโห นิสีทติ, เอโก จงฺกมํ อธิฏฺฐาติ, เอโก จรติ วิหรติ อิริยติ วตฺเตติ ปาเลติ ยเปติ ยาเปตีติ เสเวถ ปนฺตานิ เสนาสนานิ, เอโก จเร ขคฺควิสาณกปฺโป.\csromanspacer{} เตนาห โส ปจฺเจกสมฺพุทฺโธ\csromandash{}}

\prose{สีโห ยถา ทาฐพลี ปสยฺห,}

\prose{ราชา มิคานํ อภิภูยฺย จารี.}

\csromangathameasure{\textbf{เสเวถ ปนฺตานิ เสนาสนานิ,} \\ เอโก จเร ขคฺควิสาณกปฺโปติ.}
\csromangathagroup{\csromangathastanza{\textbf{เสเวถ ปนฺตานิ เสนาสนานิ,} \\ เอโก จเร ขคฺควิสาณกปฺโปติ.}}

\csromanhangingbegin
\hangingheaditem{159}{\csromanfolio{304}เมตฺตํ อุเปกฺขํ กรุณํ วิมุตฺติํ,}
\hangingbody{\textbf{อาเสวมาโน มุทิตญฺจ กาเล.}}
\hangingbody{\textbf{สพฺเพน โลเกน อวิรุชฺฌมาโน,}}
\hangingbody{เอโก จเร ขคฺควิสาณกปฺโป.}
\csromanhangingend

\prose{\textbf{เมตฺตํ อุเปกฺขํ กรุณํ วิมุตฺติํ, อาเสวมาโน มุทิตญฺจ กาเล}ติ โส ปจฺเจกสมฺพุทฺโธ เมตฺตาสหคเตน เจตสา เอกํ ทิสํ ผริตฺวา วิหรติ.\csromanspacer{} ตถา ทุติยํ.\csromanspacer{} ตถา ตติยํ.\csromanspacer{} ตถา จตุตฺถํ.\csromanspacer{} อิติ อุทฺธมโธ ติริยํ สพฺพธิ สพฺพตฺตตาย สพฺพาวนฺตํ โลกํ เมตฺตาสหคเตน เจตสา วิปุเลน มหคฺคเตน อปฺปมาเณน อเวเรน อพฺยาปชฺเชน ผริตฺวา วิหรติ.\csromanspacer{} กรุณาสหคเตน เจตสา ฯเปฯ.\csromanspacer{} มุทิตาสหคเตน เจตสา ฯเปฯ.\csromanspacer{} อุเปกฺขาสหคเตน เจตสา วิปุเลน มหคฺคเตน อปฺปมาเณน อเวเรน อพฺยาปชฺเชน ผริตฺวา วิหรตีติ เมตฺตํ อุเปกฺขํ กรุณํ วิมุตฺติํ, อาเสวมาโน มุทิตญฺจ กาเล.}

\prose{\textbf{สพฺเพน โลเกน อวิรุชฺฌมาโน}ติ เมตฺตาย ภาวิตตฺตา เย ปุรตฺถิมาย ทิสาย สตฺตา, เต อปฺปฏิกูลา โหนฺติ.\csromanspacer{} เย ปจฺฉิมาย ทิสาย สตฺตา.\csromanspacer{} เย อุตฺตราย ทิสาย สตฺตา.\csromanspacer{} เย ทกฺขิณาย ทิสาย สตฺตา.\csromanspacer{} เย ปุรตฺถิมาย อนุทิสาย สตฺตา.\csromanspacer{} เย ปจฺฉิมาย อนุทิสาย สตฺตา.\csromanspacer{} เย อุตฺตราย อนุทิสาย สตฺตา.\csromanspacer{} เย ทกฺขิณาย อนุทิสาย สตฺตา.\csromanspacer{} เย เหฏฺฐิมาย\footnote{อโธคมาย (สฺยา)} ทิสาย สตฺตา.\csromanspacer{} เย อุปริมาย ทิสาย สตฺตา.\csromanspacer{} เย ทสสุ ทิสาสุ สตฺตา, เต อปฺปฏิกูลา โหนฺติ.\csromanspacer{} กรุณาย ภาวิตตฺตา.\csromanspacer{} มุทิตาย ภาวิตตฺตา.\csromanspacer{} อุเปกฺขาย ภาวิตตฺตา เย ปุรตฺถิมาย ทิสาย สตฺตา ฯเปฯ เย ทสสุ ทิสาสุ สตฺตา, เต อปฺปฏิกูลา โหนฺติ.\csromanspacer{} \textbf{สพฺเพน โลเกน อวิรุชฺฌมาโน}ติ สพฺเพน โลเกน อวิรุชฺฌมาโน อปฺปฏิวิรุชฺฌมาโน อนาฆาติยมาโน\footnote{อฆฏฺฏิยมาโน (สฺยา)} อปฺปฏิหญฺญมาโนติ สพฺเพน โลเกน อวิรุชฺฌมาโน, เอโก จเร ขคฺควิสาณกปฺโป.\csromanspacer{} เตนาห โส ปจฺเจกสมฺพุทฺโธ\csromandash{}}

\prose{เมตฺตํ อุเปกฺขํ กรุณํ วิมุตฺติํ,}

\prose{\textbf{อาเสวมาโน มุทิตญฺจ กาเล.}}

\csromangathameasure{\textbf{สพฺเพน โลเกน อวิรุชฺฌมาโน,} \\ เอโก จเร ขคฺควิสาณกปฺโปติ.}
\csromangathagroup{\csromangathastanza{\textbf{สพฺเพน โลเกน อวิรุชฺฌมาโน,} \\ เอโก จเร ขคฺควิสาณกปฺโปติ.}}

\csromanheader{2}{19. ขคฺควิสาณสุตฺตนิทฺเทส}
\proseitem{160}{\csromanfolio{305}ราคญฺจ โทสญฺจ ปหาย โมหํ,}

\prose{\textbf{สนฺทาลยิตฺวาน สญฺโญชนานิ}\footnote{สํโยชนานิ (ก)}\textbf{.}}

\prose{\textbf{อสนฺตสํ ชีวิตสงฺขยมฺหิ,}}

\csromanheader{2}{\textbf{เอโก จเร ขคฺควิสาณกปฺโป. (9)}}
\prose{\textbf{ราคญฺจ โทสญฺจ ปหาย โมหนฺ}ติ \textbf{ราโค}ติ โย \textbf{ราโค} สา\textbf{ราโค} ฯเปฯ อภิชฺฌา โลโภ อกุสลมูลํ\textbf{.}\csromanspacer{} \textbf{โทโส}ติ โย จิตฺตสฺส อาฆาโต ฯเปฯ จณฺฑิกฺกํ อสุโรโป อนตฺตมนตา จิตฺตสฺส\textbf{.}\csromanspacer{} \textbf{โมโห}ติ ทุกฺเข อญฺญาณํ ฯเปฯ อวิชฺชาลงฺคี โมโห อกุสลมูลํ\textbf{.}\csromanspacer{} \textbf{ราคญฺจ โทสญฺจ ปหาย โมหนฺ}ติ โส ปจฺเจกสมฺพุทฺโธ ราคญฺจ โทสญฺจ โมหญฺจ ปหาย ปชหิตฺวา วิโนเทตฺวา พฺยนฺตีกริตฺวา อนภาวํ คเมตฺวาติ ราคญฺจ โทสญฺจ ปหาย โมหํ\textbf{.}}

\prose{\textbf{สนฺทาลยิตฺวาน สญฺโญชนานี}ติ ทส สญฺโญชนานิ กามราคสญฺโญชนํ ปฏิฆสญฺโญชนํ ฯเปฯ อวิชฺชาสญฺโญชนํ\textbf{.}\csromanspacer{} \textbf{สนฺทาลยิตฺวาน สญฺโญชนานี}ติ ทส สญฺโญชนานิ สนฺทาลยิตฺวา ปทาลยิตฺวา สมฺปทาลยิตฺวา ปชหิตฺวา วิโนเทตฺวา พฺยนฺตีกริตฺวา อนภาวํ คเมตฺวาติ สนฺทาลยิตฺวาน สญฺโญชนานิ\textbf{.}}

\prose{\textbf{อสนฺตสํ ชีวิตสงฺขยมฺหี}ติ โส ปจฺเจกสมฺพุทฺโธ ชีวิตปริโยสาเน อสนฺตาสี อนุตฺราสี อนุพฺพิคฺโค อนุสฺสงฺกี อนุตฺราโส อภีรู อจฺฉมฺภี อนุตฺราสี อปลายี ปหีนภยเภรโว วิคตโลมหํโส วิหรตีติ อสนฺตสํ ชีวิตสงฺขยมฺหิ, เอโก จเร ขคฺควิสาณกปฺโป\textbf{.}\csromanspacer{} เตนาห โส ปจฺเจกสมฺพุทฺโธ\csromandash{}}

\prose{ราคญฺจ โทสญฺจ ปหาย โมหํ,}

\prose{\textbf{สนฺทาลยิตฺวาน สญฺโญชนานิ.}}

\csromangathameasure{\textbf{อสนฺตสํ ชีวิตสงฺขยมฺหิ,} \\ เอโก จเร ขคฺควิสาณกปฺโปติ\textbf{.}}
\csromangathagroup{\csromangathastanza{\textbf{อสนฺตสํ ชีวิตสงฺขยมฺหิ,} \\ เอโก จเร ขคฺควิสาณกปฺโปติ\textbf{.}}}

\proseitem{161}{ภชนฺติ เสวนฺติ จ การณตฺถา,}

\csromangathameasure{\textbf{นิกฺการณา ทุลฺลภา อชฺช มิตฺตา,} \\ \textbf{อตฺตตฺถปญฺญา} \textbf{อสุจี มนุสฺสา,}}
\csromangathagroup{\csromangathastanza{\textbf{นิกฺการณา ทุลฺลภา อชฺช มิตฺตา,} \\ \textbf{อตฺตตฺถปญฺญา}\footnote{อตฺตฏฺฐปญฺญา (ก)} \textbf{อสุจี มนุสฺสา,}}}

\csromanheader{2}{\textbf{เอโก จเร ขคฺควิสาณกปฺโป. (10)}}
\prose{\csromanfolio{306}\textbf{ภชนฺติ เสวนฺติ จ การณตฺถา}ติ อตฺตตฺถการณา ปรตฺถการณา อุภยตฺถการณา ทิฏฺฐธมฺมิกตฺถการณา สมฺปรายิกตฺถการณา ปรมตฺถการณา ภชนฺติ สมฺภชนฺติ เสวนฺติ นิเสวนฺติ สํเสวนฺติ ปฏิเสวนฺตีติ ภชนฺติ เสวนฺติ จ การณตฺถา.}

\prose{\textbf{นิกฺการณา ทุลฺลภา อชฺช มิตฺตา}ติ เทฺว มิตฺตา อคาริกมิตฺโต จ อนาคาริกมิตฺโต จ ฯเปฯ อยํ อคาริกมิตฺโต ฯเปฯ อยํ อนาคาริกมิตฺโต.\csromanspacer{} \textbf{นิกฺการณา ทุลฺลภา อชฺช มิตฺตา}ติ อิเม เทฺว มิตฺตา อการณา นิกฺการณา อเหตู อปฺปจฺจยา ทุลฺลภา (ทุลฺลทฺธา สุทุลฺลทฺธา)1ติ นิกฺการณา ทุลฺลภา อชฺช มิตฺตา.}

\prose{\textbf{อตฺตตฺถปญฺญา อสุจี มนุสฺสา}ติ \textbf{อตฺตตฺถปญฺญา}ติ อตฺตโน อตฺถาย อตฺตโน เหตุ อตฺตโน ปจฺจยา อตฺตโน การณา ภชนฺติ สมฺภชนฺติ เสวนฺติ นิเสวนฺติ สํเสวนฺติ ปฏิเสวนฺติ อาจรนฺติ สมาจรนฺติ ปยิรุปาสนฺติ ปริปุจฺฉนฺติ ปริปญฺหนฺตีติ \textbf{อตฺตตฺถปญฺญา}.\csromanspacer{} \textbf{อสุจี มนุสฺสา}ติ อสุจินา กายกมฺเมน สมนฺนาคตาติ อสุจี มนุสฺสา, อสุจินา วจีกมฺเมน สมนฺนาคตาติ อสุจี มนุสฺสา, อสุจินา มโนกมฺเมน สมนฺนาคตาติ อสุจี มนุสฺสา, อสุจินา ปาณาติปาเตน.\csromanspacer{} อสุจินา อทินฺนาทาเนน ฯเปฯ.\csromanspacer{} อสุจินา กาเมสุมิจฺฉาจาเรน.\csromanspacer{} อสุจินา มุสาวาเทน.\csromanspacer{} อสุจิยา ปิสุณาย วาจาย สมนฺนาคตา.\csromanspacer{} อสุจิยา ผรุสาย วาจาย สมนฺนาคตา.\csromanspacer{} อสุจินา สมฺผปฺปลาเปน สมนฺนาคตา.\csromanspacer{} อสุจิยา อภิชฺฌาย สมนฺนาคตา อสุจินา พฺยาปาเทน สมนฺนาคตาติ อสุจี มนุสฺสา, อสุจิยา มิจฺฉาทิฏฺฐิยา สมนฺนาคตาติ อสุจี มนุสฺสา, อสุจิยา เจตนาย สมนฺนาคตาติ อสุจี มนุสฺสา, อสุจิยา ปตฺถนาย สมนฺนาคตาติ อสุจี มนุสฺสา, อสุจินา ปณิธินา สมนฺนาคตาติ อสุจี มนุสฺสา, อสุจี หีนา นิหีนา โอมกา ลามกา ฉตุกฺกา ปริตฺตาติ \textbf{อตฺตตฺถปญฺญา} อสุจี มนุสฺสา.}

\prose{\textbf{เอโก จเร ขคฺควิสาณกปฺโป}ติ \textbf{เอโก}ติ โส ปจฺเจกสมฺพุทฺโธ ปพฺพชฺชาสงฺขาเตน \textbf{เอโก} ฯเปฯ.\csromanspacer{} \textbf{จเร}ติ อฏฺฐ จริยาโย ฯเปฯ.\csromanspacer{} \textbf{ขคฺควิสาณกปฺโป}ติ ยถา ขคฺคสฺส นาม วิสาณํ เอกํ โหติ อทุติยํ ฯเปฯ \textbf{เอโก} จเร ขคฺควิสาณกปฺโป.\csromanspacer{} เตนาห โส ปจฺเจกสมฺพุทฺโธ\csromandash{}}

\csromanheader{2}{19. ขคฺควิสาณสุตฺตนิทฺเทส}
\csromangathameasure{ภชนฺติ เสวนฺติ จ การณตฺถา, \\ นิกฺการณา ทุลฺลภา อชฺช มิตฺตา. \\ อตฺตตฺถปญฺญา อสุจี มนุสฺสา, \\ เอโก จเร ขคฺควิสาณกปฺโปติ. จตุตฺโถ วคฺโค. ขคฺควิสาณสุตฺตนิทฺเทโส เอกูนวีสติโม. \\ อชิโต ติสฺสเมตฺเตยฺโย, ปุณฺณโก อถ เมตฺตคู. \\ โธตโก อุปสีโว จ, นนฺโท จ อถ เหมโก. \\ โตเทยฺยกปฺปา ทุภโย, ชตุกณฺณี จ ปณฺฑิโต. \\ ภทฺราวุโธ อุทโย จ, โปสาโล จาปิ พฺราหฺมโณ. \\ โมฆราชา จ เมธาวี, ปิงฺคิโย จ มหา-อิสิ. \\ โสฬสานํ ปเนเตสํ, พฺราหฺมณานํว สาสนํ. \\ ปารายนานํ นิทฺเทสา, ตตฺถกา จ ภวนฺติ หิ. \\ ขคฺควิสาณสุตฺตานํ, นิทฺเทสาปิ ตเถว จ. \\ นิทฺเทสา ทุวิธา เญยฺยา, ปริปุณฺณา สุลกฺขิตาติ.}
\csromangathagroup{\csromangathastanza{\csromanfolio{307}ภชนฺติ เสวนฺติ จ การณตฺถา, \\ นิกฺการณา ทุลฺลภา อชฺช มิตฺตา. \\ อตฺตตฺถปญฺญา อสุจี มนุสฺสา, \\ เอโก จเร ขคฺควิสาณกปฺโปติ.\csromanspacer{} จตุตฺโถ วคฺโค.\csromanspacer{} ขคฺควิสาณสุตฺตนิทฺเทโส เอกูนวีสติโม. \\ อชิโต ติสฺสเมตฺเตยฺโย, ปุณฺณโก อถ เมตฺตคู. \\ โธตโก อุปสีโว จ, นนฺโท จ อถ เหมโก.}\csromangathastanza{โตเทยฺยกปฺปา ทุภโย, ชตุกณฺณี จ ปณฺฑิโต. \\ ภทฺราวุโธ อุทโย จ, โปสาโล จาปิ พฺราหฺมโณ.}\csromangathastanza{โมฆราชา จ เมธาวี, ปิงฺคิโย จ มหา-อิสิ. \\ โสฬสานํ\footnote{โสฬสนฺนํ (สฺยา, ก)} ปเนเตสํ, พฺราหฺมณานํว สาสนํ.}\csromangathastanza{ปารายนานํ นิทฺเทสา, ตตฺถกา จ ภวนฺติ หิ\footnote{วา (ก)}. \\ ขคฺควิสาณสุตฺตานํ, นิทฺเทสาปิ ตเถว จ. \\ นิทฺเทสา ทุวิธา เญยฺยา, ปริปุณฺณา สุลกฺขิตาติ.}}

\nitthitam{\textbf{จูฬนิทฺเทสปาฬิ นิฏฺฐิตา.}}
\orphannote{( ) เอตฺถนฺตเร ปาฐา นตฺถิ สฺยา-โปตฺถเก.}

\orphannote{อวเสมาโน ปริเสมาโน (สฺยา)}

\orphannote{สีมิกตํ มริยทิกตํ (สฺยา)}

\orphannote{อุภยตฺถมนฺตนา (ก)}

\orphannote{( ) นตฺถิ สฺยา-โปตฺถเก.}

